% =====================================================================
% NP -- Rigidity of eigenvalues of generalized Wigner matrices: a self-contained proof.
% Version 1.1. Main file; the sections are in sec/*.tex.
% =====================================================================
\documentclass[11pt]{article}
\usepackage{amsmath,amssymb,amsthm,mathrsfs}
\usepackage[margin=1in]{geometry}
\usepackage{enumitem}
\usepackage{alltt}
\usepackage[colorlinks=true,linkcolor=blue,citecolor=blue,urlcolor=blue,hypertexnames=false]{hyperref}
\numberwithin{equation}{section}
\setcounter{tocdepth}{2}

% ---------------- theorem-like environments ----------------
% Part I: definitions, remarks, conventions have their own counter (per section).
\newtheorem{theorem}{Theorem}[section]
\newtheorem{lemma}[theorem]{Lemma}
\newtheorem{proposition}[theorem]{Proposition}
\newtheorem{corollary}[theorem]{Corollary}
\theoremstyle{definition}
\newtheorem{definition}{Definition}[section]
\newtheorem{remark}{Remark}[section]
\newtheorem{convention}[definition]{Convention}
% Parts II and III: one counter for all statements of a section;
% switched by \useSharedCounters at the beginning of Part II.
\newcommand{\useSharedCounters}{%
  \theoremstyle{definition}%
  \let\definition\relax\let\enddefinition\relax
  \let\remark\relax\let\endremark\relax
  \newtheorem{definition}[theorem]{Definition}%
  \newtheorem{remark}[theorem]{Remark}%
  \theoremstyle{plain}}
% lettered auxiliary statements of Part III: numbered <section>.A, <section>.B, ...
\theoremstyle{plain}
\newtheorem{auxlemma}{Lemma}[section]
\renewcommand{\theauxlemma}{\thesection.\Alph{auxlemma}}
\theoremstyle{definition}
\newtheorem{auxdefinition}[auxlemma]{Definition}
\newtheorem{auxremark}[auxlemma]{Remark}
% unnumbered restatement of a theorem of Part I inside Part III
\theoremstyle{plain}
\newcommand{\restatedname}{Theorem}
\newtheorem*{restated}{\restatedname}

% ---------------- macros ----------------
\newcommand{\R}{\mathbb{R}}
\newcommand{\C}{\mathbb{C}}
\newcommand{\Z}{\mathbb{Z}}
\newcommand{\kk}{\Bbbk}
\newcommand{\PP}{\mathbb{P}}
\newcommand{\EE}{\mathbb{E}}
\newcommand{\re}{\operatorname{Re}}
\newcommand{\im}{\operatorname{Im}}
\newcommand{\tr}{\operatorname{Tr}}
\newcommand{\msc}{m_{sc}}
\newcommand{\rhosc}{\varrho_{sc}}
\newcommand{\nsc}{n_{sc}}
\newcommand{\fn}{\mathfrak{n}}
\newcommand{\fa}{\mathfrak{a}}
\newcommand{\fL}{\mathfrak{L}}
\newcommand{\cG}{\mathcal{G}}
\newcommand{\bS}{\mathbf{S}}
\newcommand{\bD}{\mathbf{D}}
\newcommand{\bone}{\mathbf{1}}
\newcommand{\be}{\mathbf{e}}
\newcommand{\Ld}{\Lambda_d}
\newcommand{\Lo}{\Lambda_o}
\newcommand{\vth}{\vartheta}
\newcommand{\dpm}{\delta_\pm}
\newcommand{\Emax}{E_{\max}}
\newcommand{\infra}{\textup{[INFRA]}}
\newcommand{\plan}{\noindent\textit{Proof plan.}\ }
\newcommand{\ol}[1]{\overline{#1}}
\newcommand{\avg}[1]{[#1]}
\newcommand{\norm}[1]{\|#1\|}
\newcommand{\abs}[1]{|#1|}
\newcommand{\wh}{\widehat}
\newcommand{\wt}{\widetilde}
\newcommand{\eps}{\varepsilon}
\newcommand{\rd}{\mathrm{d}}
\newcommand{\ii}{\mathrm{i}}

\newcommand{\E}{\mathbb{E}}
\renewcommand{\P}{\mathbb{P}}
\newcommand{\K}{\Bbbk}
\newcommand{\Log}{\operatorname{Log}}
\newcommand{\e}{\varepsilon}
\newcommand{\eg}{\varepsilon_{\mathrm g}}
\newcommand{\Par}{\mathsf{Par}}
\newcommand{\fA}{\mathfrak{A}}
\newcommand{\cR}{\mathcal{R}}
\newcommand{\br}[1]{\big[#1\big]_{1}^{2}}

\newcommand{\cH}{\mathcal{H}}
\newcommand{\cS}{\mathcal{S}}
\newcommand{\cV}{\mathcal{V}}
\newcommand{\cP}{\mathcal{P}}
\newcommand{\cX}{\mathcal{X}}
\newcommand{\bG}{\mathbf{G}}
\newcommand{\Tr}{\operatorname{Tr}}
\newcommand{\diag}{\operatorname{diag}}
\newcommand{\Ent}{\operatorname{Ent}}
\newcommand{\Lip}{\operatorname{Lip}}
\newcommand{\dv}{\operatorname{div}}
\newcommand{\Sig}{\Sigma_N}
\newcommand{\cSig}{\overline{\Sigma}_N}
\newcommand{\INFRA}{\textup{[INFRA]}}
\newcommand{\la}{\langle}
\newcommand{\ra}{\rangle}

\sloppy
\emergencystretch=4em
\allowdisplaybreaks

\begin{document}
\title{Rigidity of eigenvalues of generalized Wigner matrices:\\ a self-contained proof}
\author{Jun Yin \quad Claude\\[1.5ex] \normalsize Version 1.2 --- written as the source of a Lean~4 formalization;\\ \normalsize based on L.~Erd\H{o}s, H.-T.~Yau, J.~Yin, arXiv:1007.4652}
\date{}
\maketitle

\begin{abstract}
We give a self-contained proof of the four main results of L.~Erd\H{o}s, H.-T.~Yau and J.~Yin, \emph{Rigidity of eigenvalues of generalized Wigner matrices}
(arXiv:1007.4652): the strong local semicircle law (Theorem~\ref{thm:21}), the rigidity of the eigenvalues (Theorem~\ref{thm:22}), the strong local ergodicity of
Dyson Brownian motion (Theorem~\ref{thm:23}) and the universality of the extreme eigenvalues (Theorem~\ref{thm:24}), for real symmetric and complex Hermitian
generalized Wigner matrices with subexponentially decaying entries. All statements are quantitative and hold for every finite $N$ beyond an explicit threshold;
the only results used without proof are standard facts available in Lean's Mathlib, listed in Appendix~\ref{app:mathlib}. Several proofs differ from those of the
original; the differences in the statements are listed in Appendix~\ref{app:diff}.
The text is organised in three parts: Part~I (Sections~\ref{sec:intro}--\ref{sec:export}) contains the model, the statements and the proofs of Theorems~\ref{thm:21} and~\ref{thm:22};
Part~II (Section~\ref{sec:edge}) proves Theorem~\ref{thm:24}; Part~III (Sections~\ref{sec:C1}--\ref{sec:C9}) proves Theorem~\ref{thm:23}.
\end{abstract}

\tableofcontents

% ================= PART I =================
\clearpage
\part[The model and the main results; the local semicircle law and rigidity (Theorems 2.1, 2.2)]{The model and the main results;\\ the local semicircle law and rigidity (Theorems 2.1 and 2.2)}
\section{Introduction and conventions}\label{sec:intro}

\subsection{The results and the scope of this paper}
A \emph{generalized Wigner matrix} is an $N\times N$ Hermitian or real symmetric random matrix $H=(h_{ij})$ whose entries are independent up to the symmetry
constraint, centred, with variances $\sigma_{ij}^2=\EE|h_{ij}|^2$ that may depend on $(i,j)$ but sum to one in every row. The paper
of Erd\H{o}s, Yau and Yin, \emph{Rigidity of eigenvalues of generalized Wigner matrices} (below: \emph{the original}), established four results for such
matrices, under a spectral gap condition on the matrix of variances and a subexponential decay condition on the entries:
\begin{itemize}[leftmargin=2em,itemsep=1pt]
\item the \emph{strong local semicircle law} (Theorem~\ref{thm:21}): the Stieltjes transform $m(z)$ of the empirical spectral measure and the individual entries $G_{ij}(z)$
of the resolvent are close to their deterministic limits, down to the optimal scale $\im z\gg N^{-1}$, with the optimal error $(N\im z)^{-1}$ for $m$, uniformly in
the energy up to the spectral edges, together with the bound $2+N^{-2/3+o(1)}$ on the norm of $H$;
\item the \emph{rigidity of eigenvalues} (Theorem~\ref{thm:22}): the $j$-th eigenvalue is within $N^{-2/3+o(1)}\min(j,N+1-j)^{-1/3}$ of its classical location $\gamma_j$;
\item the \emph{strong local ergodicity of Dyson Brownian motion} (Theorem~\ref{thm:23}): after a time $t\gg N^{-1}$ the local eigenvalue statistics in the bulk, averaged over a small
energy window, coincide with those of the Gaussian ensemble, with an explicit power-law error;
\item the \emph{universality of the extreme eigenvalues} (Theorem~\ref{thm:24}): the distributions of the largest (and of the smallest) eigenvalue of two generalized Wigner
matrices with matching first two moments agree on the scale $N^{-2/3}$.
\end{itemize}
The present paper, called NP below, proves these four theorems again. Its purpose is different from that of the original: NP is written as the mathematical
source of a formalization in Lean~4 with Mathlib. This has three consequences.

\emph{(1) Self-containedness.} No result is quoted from another paper or book. Whatever is not available in Mathlib is proved here. A statement marked \infra\ (``infrastructure'') is a standard result of analysis, linear algebra or probability
which is \emph{not} available in Mathlib and which NP therefore proves in full (for instance the Helffer--Sj\"ostrand formula, the joint eigenvalue density of the Gaussian ensembles,
the Pr\'ekopa--Leindler inequality); the marker carries no other meaning, and it indicates the places where the formalization will have to add material to Mathlib.
Appendix~\ref{app:mathlib} lists the facts that \emph{are} taken from Mathlib. In particular the inputs which the original takes from earlier papers
(the Schur complement identities, the large deviation bounds for quadratic forms, the properties of $\msc$, the a priori bound on the norm, the Helffer--Sj\"ostrand
argument for the counting function, the logarithmic Sobolev inequality and the local relaxation flow) are all proved, in most cases by an argument chosen for
its shortness and for the absence of analytic prerequisites.

\emph{(2) Finite-$N$ statements with explicit quantifiers.} There are no limits and no ``with high probability'': every probability is bounded by an explicit expression,
every constant has a stated dependence, and the order of the quantifiers is fixed once and for all (Convention~\ref{conv:const}).

\emph{(3) New proofs where the old ones are long.} The statements are those of the original, up to the corrections listed in Appendix~\ref{app:diff}; several proofs are new.

NP has three parts. Part~I (Sections~\ref{sec:intro}--\ref{sec:export}) contains the model, the four statements, and the proofs of Theorems~\ref{thm:21}
and~\ref{thm:22}. Part~II (Section~\ref{sec:edge}) proves Theorem~\ref{thm:24} and Part~III (Sections~\ref{sec:C1}--\ref{sec:C9}) proves Theorem~\ref{thm:23}; both use Theorems~\ref{thm:21}, \ref{thm:22} and the auxiliary results
collected in Section~\ref{sec:export}, and nothing else from Part~I. The appendices contain the bookkeeping of the exponents of Part~I (Appendix~\ref{app:ledger}), its dependency graph
(Appendix~\ref{app:dep}), the list of differences from the original (Appendix~\ref{app:diff}) and the list of facts taken from Mathlib (Appendix~\ref{app:mathlib}).

\subsection{Strategy of Part~I}
\begin{enumerate}[leftmargin=2em,itemsep=2pt]
\item \emph{One level parameter.} The original works with three coupled exponents (in the factors $(\log N)^{\ell}$, $(\log N)^{L}$ and in an iteration index).
Here a single real parameter $K$, the \emph{level}, with $(\log N)^3\le K\le N^{c_0}$, is fixed. Every error factor is an explicit power $K^{C}$ and every exceptional
probability is at most $N^{C}e^{-K}$. The main internal result is Theorem~\ref{thm:K}; Theorems~\ref{thm:21} and~\ref{thm:22}, with the original powers of $\log N$, follow
from it by the substitution $K=(\log N)^{\phi L}$ with a small constant $\phi$ (Proposition~\ref{prop:reduction}, proved in Section~\ref{sec:main}).
\item \emph{Probabilistic input} (Section~\ref{sec:prob}). Large deviation bounds for linear, quadratic and bilinear forms in independent variables with unequal, possibly
vanishing variances. They are derived from Hoeffding's inequality after truncation, with a deterministic decoupling of the off-diagonal quadratic form.
\item \emph{Entrywise (weak) law} (Section~\ref{sec:weak}). Schur complement identities, the large deviation bounds conditionally on a minor, stability of the self-consistent
equation (the vector part on $\be^\perp$ without any factor $\log N$; the scalar part is a quadratic inequality), and a \emph{deterministic} continuity argument in
$\eta=\im z$ along a finite lattice.
\item \emph{Fluctuation averaging} (Section~\ref{sec:FA}). For the index set $L$ of a $p$-fold product one conditions on the minor $H^{(L)}$; the block of $G$ on $L$ is the inverse of an
$|L|\times|L|$ matrix $M$ whose entry $M_{ab}$ depends only on the rows $a$ and $b$ of $H$. A finite Neumann expansion in the off-diagonal part of $M$ exhibits the exact dependence
structure, and the terms in which a lone index is not visited have zero expectation. This replaces the thirty pages of Section~7 of the original.
\item \emph{Strong law} (Section~\ref{sec:strong}). A moment inequality that improves itself in one step (no iteration), followed by a dichotomy and a continuity argument at the level
of events; a second pass gives the improved bound outside the spectrum.
\item \emph{Norm bound} (Section~\ref{sec:norm}). A crude bound $\norm H\le K^{C}$ from a grid net and Hoeffding's inequality; eigenvalues between $2+K^{C_3}N^{-2/3}$ and the crude bound
are excluded by the averaged law outside the spectrum.
\item \emph{Rigidity} (Sections~\ref{sec:det} and~\ref{sec:rig}). The Helffer--Sj\"ostrand formula, proved from the divergence theorem on rectangles (Lemma~\ref{lem:HS}); a deterministic
bound on the counting function in terms of $m-\msc$ (Lemma~\ref{lem:count}); and a deterministic conversion lemma which turns the bound on the counting function
into the bound on $\lambda_j-\gamma_j$ with the factor $\min(j,N+1-j)^{-1/3}$, at both edges at once (Lemma~\ref{lem:conv}).
\end{enumerate}
The semicircle law itself enters only through the algebraic function $\msc$: its equation, its size (Lemmas~\ref{lem:msc}, \ref{lem:mscsize}), and the fact that it is the
Stieltjes transform of $\rhosc$ (Lemma~\ref{lem:stieltjes}), which is needed only for the counting function.

\subsection{Conventions}\label{sec:conv}
\begin{convention}[Constants]\label{conv:const}
The \emph{model parameters} are $\vth>0$, $\delta_->0$, $\delta_+>0$, $C_0>0$. By the monotonicity of the assumptions
(Lemma~\ref{lem:mono}(ii)) we assume in all proofs that $\vth\le1$, $\dpm\le1$, $C_0\ge1$, and we put
\[
 \fa:=1/\vth\ \ (\ge1),\qquad \delta_{\min}:=\min(\delta_-,\delta_+).
\]
Letters $C,c$ denote positive constants that depend only on $(\vth,\delta_-,\delta_+,C_0)$; their values may change from
line to line, and they never depend on $N$, $K$, $z$, the ensemble, the field $\kk$, or indices. An \emph{absolute} constant is a fixed real number.
Named constants ($c_*$, $C_*$, $C_s$, $\lambda_{\rm st}$, $c_F$, $C_1,\dots,C_7$, \dots) keep their value throughout Part~I. Parts~II and~III use a few
constants with the same names ($C_2,C_3,C_4$ and $\Psi$ in Section~\ref{sec:edge}; $c_0,c_1$, $\Psi$ and others in Part~III); these are defined locally in the section where they occur,
are unrelated to the constants of Part~I, and this is recalled there. The letter $\delta$ without index is used only for the exponents of Theorems~\ref{thm:23} and~\ref{thm:24}
(as in the original); the two are unrelated, and each is local to its theorem and to its proof.
``For $N\ge N_0$'' means: for all $N$ above a threshold that depends only on the model parameters.
Quantifiers are always in the order: model parameters $\to$ constants $\to$ $N_0$ $\to$ ($N$, $\kk$, ensemble) $\to$ $K$ (or $L$) $\to$ $z$.
\end{convention}
\begin{convention}[Probability]\label{conv:prob}
All probabilities are explicit. An event is called \emph{$K$-likely} if its complement has probability at most $N^{C}e^{-K}$
with a constant $C$ as in Convention~\ref{conv:const}. A union of at most $N^{C'}$ complements of $K$-likely events has probability
at most $N^{C+C'}e^{-K}$, and
\[
 N^{C}e^{-K}\le e^{-K/2}\qquad\text{whenever }K\ge(\log N)^3\text{ and }(\log N)^2\ge2C,
\]
because then $C\log N\le\frac12(\log N)^3\le K/2$. Events indexed by uncountably many $z$
are always handled as: a finite lattice of $z$'s, a union bound, and a deterministic Lipschitz extension; the final statements
assert the existence of a measurable event of large probability on which a bound holds for all $z$.
\end{convention}
\begin{convention}[Spectral parameter]\label{conv:z}
$z=E+\ii\eta$ with $E\in\R$, $\eta>0$; $\kappa:=\big||E|-2\big|$; $\varkappa:=\sqrt{\kappa+\eta}$. We write $\C_+:=\{z\in\C:\im z>0\}$.
\end{convention}
\begin{convention}[General notation]\label{conv:not}
$\log$ is the natural logarithm. For $x\in\C^n$, $|x|$ is the Euclidean norm, $\norm x_\infty$ the maximum norm, and $\langle x,y\rangle=\sum_i\ol{x_i}y_i$;
for a matrix $A$, $\norm A$ is the operator norm with respect to the Euclidean norm. $\bone(\cdot)$ is an indicator, $\#S$ the cardinality of a finite set $S$,
$(t)_+:=\max(t,0)$, and $|I|$ is the length of an interval $I$. For a complex number $\zeta\notin(-\infty,0]$, $\sqrt\zeta$ is the principal square root (the one with positive real part).
Integrals without a measure are with respect to Lebesgue measure.
\end{convention}

\section{Model, assumptions and main results}\label{sec:main}

\subsection{The ensemble}
Let $\kk\in\{\R,\C\}$ ($\kk=\R$: real symmetric case, $\beta=1$; $\kk=\C$: complex Hermitian case, $\beta=2$), $N\ge3$, and
$\mathcal P_N:=\{(i,j): 1\le i\le j\le N\}$.

\begin{definition}[Ensemble]\label{def:ens}
An \emph{ensemble} (of size $N$ over $\kk$) is a family $\nu=(\nu_{ij})_{(i,j)\in\mathcal P_N}$ of Borel probability measures on $\kk$.
The probability space is $\Omega:=\kk^{\mathcal P_N}$ with the product $\sigma$-algebra and the product law $\PP=\PP_\nu:=\bigotimes_{(i,j)\in\mathcal P_N}\nu_{ij}$; $\EE$ is the expectation with respect to $\PP$.
For $x\in\Omega$ the matrix $H=H(x)=(h_{ij})_{i,j=1}^N$ is
\[
 h_{ij}:=x_{ij}\ (i<j),\qquad h_{ji}:=\ol{x_{ij}}\ (i<j),\qquad h_{ii}:=\re x_{ii}.
\]
Thus $H(x)$ is Hermitian (symmetric if $\kk=\R$) at \emph{every} point $x\in\Omega$, the entries $(h_{ij})_{i\le j}$ are independent, and
$h_{ij}$ has law $\nu_{ij}$ for $i<j$ (also for $i=j$ if $\nu_{ii}(\R)=1$).
We put
\begin{equation}\label{eq:aver}
 \sigma_{ij}^2:=\int_{\kk}|y|^2\,\nu_{ij}(\rd y)\in[0,\infty]\ \ (i\le j),\qquad \sigma_{ji}:=\sigma_{ij},\qquad B:=(\sigma_{ij}^2)_{i,j=1}^N .
\end{equation}
\end{definition}

\begin{definition}[Assumptions]\label{def:ass}
Let $\vth,\delta_-,\delta_+,C_0>0$. The ensemble $\nu$ belongs to the class $\mathcal E_N(\vth,\delta_-,\delta_+,C_0)$ (over $\kk$) if:
\begin{description}[leftmargin=2.2em,itemsep=1pt]
\item[(R)] $\nu_{ii}(\R)=1$ for all $i$ (real diagonal).
\item[(D)] \emph{Subexponential decay:} for all $(i,j)\in\mathcal P_N$ and all real $x\ge1$,
\begin{equation}\label{eq:subexp}
 \nu_{ij}\big(\{y\in\kk:\ |y|> x\,\sigma_{ij}\}\big)\le \vth^{-1}\exp(-x^{\vth}).
\end{equation}
\item[(Z)] \emph{Centring (zero mean):} $\int y\,\nu_{ij}(\rd y)=0$ for all $(i,j)\in\mathcal P_N$ (the integral exists because $\sigma_{ij}<\infty$, by \textup{(A)}).
\item[(A)] $\sum_{i=1}^N\sigma_{ij}^2=1$ for every $j$.
\item[(B)] For every $v\in\R^N$ with $\sum_iv_i=0$:\quad $(-1+\delta_-)\,|v|^2\le \langle v,Bv\rangle\le(1-\delta_+)\,|v|^2$.
\item[(C)] $\sigma_{ij}^2\le C_0/N$ for all $i,j$.
\end{description}
The labels \textup{(A), (B), (C)} are those of the original; the conditions \textup{(D)} and \textup{(Z)} are (2.17) and (2.1) of the original.
\end{definition}
\noindent Vanishing variances $\sigma_{ij}=0$ are allowed; then $\nu_{ij}$ is the point mass at $0$ (trivially, or by \eqref{eq:subexp}: $\nu_{ij}(\{|y|>0\})\le\vth^{-1}e^{-x^\vth}$ for all $x\ge1$) and $h_{ij}=0$ almost surely. No statement in NP divides by $\sigma_{ij}$,
and every event comparing $|h_{ij}|$ with a multiple of $\sigma_{ij}$ uses a strict inequality, as in \eqref{eq:subexp}.
All moments of $h_{ij}$ are finite: if $\vth\le1$ (which may be assumed, by Lemma~\ref{lem:mono}(ii)), then $\EE|h_{ij}|^q\le(C_\vth q)^{q/\vth}\sigma_{ij}^q$ for every real $q\ge1$, with $C_\vth=90/\vth$ (Lemma~\ref{lem:moments}(a)); for $\vth>1$ the same bound holds with $\vth$ replaced by $1$.
Since $B$ is symmetric, (A) says that all row sums and all column sums of $B$ equal $1$.

Elementary properties of the class (equivalence of (B) with the spectral condition (2.3) of the original, monotonicity in the parameters, behaviour under $H\mapsto aH$) are collected in Lemma~\ref{lem:mono}.

\subsection{Green function, semicircle law, control parameters}\label{sec:green}
For $z=E+\ii\eta$, $\eta>0$, and for $H$ Hermitian,
\begin{equation}\label{eq:green}
 G(z):=(H-z)^{-1},\qquad m(z):=\frac1N\tr G(z)=\frac1N\sum_{j=1}^NG_{jj}(z).
\end{equation}
(If $\kk=\R$, $H$ is regarded as a complex matrix.) The Stieltjes transform of the semicircle law is defined \emph{algebraically} by
\begin{equation}\label{eq:defmsc}
 \msc(z):=\tfrac12\big(-z+\sqrt{z-2}\,\sqrt{z+2}\,\big)\qquad(\im z>0),
\end{equation}
with the principal branch of both square roots; by Lemma~\ref{lem:msc} it is the unique solution of
$\msc(z)+\frac1{z+\msc(z)}=0$ with $\im\msc(z)>0$, and by Lemma~\ref{lem:stieltjes} it is the Stieltjes transform of $\rhosc$. Further,
\begin{equation}\label{eq:sc}
 \rhosc(x):=\frac{\sqrt{(4-x^2)_+}}{2\pi},\quad \nsc(E):=\int_{-\infty}^E\rhosc(x)\,\rd x,\quad
 \gamma_j:=\inf\{x\in\R:\ N\nsc(x)\ge j\}\ \ (1\le j\le N),
\end{equation}
so that $N\nsc(\gamma_j)=j$ and $-2<\gamma_1<\dots<\gamma_{N-1}<\gamma_N=2$ (Lemma~\ref{lem:nsc}). The eigenvalues of $H$, counted with multiplicity, are
$\lambda_1\le\lambda_2\le\dots\le\lambda_N$, and the empirical counting function is
\[
 \fn(E):=\frac1N\#\{j:\lambda_j\le E\}\qquad(E\in\R).
\]
We set
\begin{equation}\label{eq:defLambda}
 v_k:=G_{kk}-\msc,\quad \Ld:=\max_k|G_{kk}-\msc|,\quad \Lo:=\max_{k\ne l}|G_{kl}|,\quad \Lambda:=|m-\msc|,
\end{equation}
all depending on $z$. For $L>0$,
\begin{equation}\label{eq:defS}
 \bS_L:=\big\{z=E+\ii\eta:\ |E|\le5,\ N^{-1}(\log N)^{10L}<\eta\le10\big\}.
\end{equation}
\emph{Measurability.} For fixed $z$ the functions $x\mapsto\lambda_j(H(x))$, $G_{ij}(z)$, $m(z)$, $\Lambda(z)$, $\Ld(z)$, $\Lo(z)$ are continuous on $\Omega$
(Lemma~\ref{lem:weyl}(d),(e)), hence Borel measurable. The counting function $x\mapsto\fn(E)$ is \emph{not} continuous on $\Omega$ (it is integer valued up to the factor $1/N$);
it is a Borel function, because $\{x:\fn(E)\ge k/N\}=\{x:\lambda_k(H(x))\le E\}$ is closed, and $\sup_{E\in J}|\fn(E)-\nsc(E)|$, for $J=\R$ or $J$ a compact interval,
is a Borel function on $\Omega$, being equal to a supremum over a countable set of $E$'s (Lemma~\ref{lem:weyl}(f)).
For a possibly non-measurable set $S\subset\Omega$, $\PP^*(S):=\inf\{\PP(S'):S\subset S',\ S'\text{ measurable}\}$ denotes the outer probability.

\subsection{The four main theorems}

\begin{theorem}[Strong local semicircle law]\label{thm:21}
For all $\vth,\delta_-,\delta_+,C_0>0$ there exist constants $A_0>1$, $C>0$, $c>0$, $\phi\in(0,1)$ and $N_0$, depending only on
$(\vth,\delta_-,\delta_+,C_0)$, such that for every $N\ge N_0$, every $\kk\in\{\R,\C\}$, every ensemble
$\nu\in\mathcal E_N(\vth,\delta_-,\delta_+,C_0)$ and every real $L$ with
\begin{equation}\label{eq:Lbound}
 A_0\log\log N\le L\le\frac{\log(10N)}{10\log\log N}
\end{equation}
the following hold.
\begin{enumerate}[label=\textup{(\roman*)},leftmargin=2.2em]
\item
\begin{equation}\label{eq:Lambdafinal}
 \PP^*\Big(\bigcup_{z\in\bS_L}\Big\{\Lambda(z)\ge\frac{(\log N)^{4L}}{N\eta}\Big\}\Big)\le C\exp\big[-c(\log N)^{\phi L}\big].
\end{equation}
\item
\begin{equation}\label{eq:Lambdaodfinal}
 \PP^*\Big(\bigcup_{z\in\bS_L}\Big\{\Ld(z)+\Lo(z)\ge(\log N)^{4L}\sqrt{\frac{\im\msc(z)}{N\eta}}+\frac{(\log N)^{4L}}{N\eta}\Big\}\Big)\le C\exp\big[-c(\log N)^{\phi L}\big].
\end{equation}
\item
\begin{equation}\label{eq:443}
 \PP\Big(\max_{1\le j\le N}|\lambda_j|\ge2+N^{-2/3}(\log N)^{9L}\Big)\le C\exp\big[-c(\log N)^{\phi L}\big].
\end{equation}
\end{enumerate}
In \textup{(i)} and \textup{(ii)}, $\eta=\im z$. More precisely, there is one measurable event of probability at least $1-C\exp[-c(\log N)^{\phi L}]$ on which none of the events in
\eqref{eq:Lambdafinal}, \eqref{eq:Lambdaodfinal}, \eqref{eq:443} occurs, for any $z\in\bS_L$.
\end{theorem}

\begin{theorem}[Rigidity of eigenvalues]\label{thm:22}
For all $\vth,\delta_-,\delta_+,C_0>0$ there exist $A_0>1$, $C,c>0$, $\phi\in(0,1)$ and $N_0$, depending only on
$(\vth,\delta_-,\delta_+,C_0)$, such that for every $N\ge N_0$, $\kk\in\{\R,\C\}$, $\nu\in\mathcal E_N(\vth,\delta_-,\delta_+,C_0)$ and $L$ satisfying \eqref{eq:Lbound}:
\begin{equation}\label{eq:rigidity}
 \PP\Big(\exists j\in\{1,\dots,N\}:\ |\lambda_j-\gamma_j|\ge(\log N)^L\big[\min(j,N-j+1)\big]^{-1/3}N^{-2/3}\Big)\le C\exp\big[-c(\log N)^{\phi L}\big],
\end{equation}
\begin{equation}\label{eq:nn}
 \PP\Big(\sup_{|E|\le5}\big|\fn(E)-\nsc(E)\big|\ge\frac{(\log N)^L}{N}\Big)\le C\exp\big[-c(\log N)^{\phi L}\big].
\end{equation}
\end{theorem}
\noindent In \eqref{eq:nn} the supremum itself is used; it need not be attained, since $\fn$ is only right-continuous in $E$. The constants $(A_0,C,c,\phi,N_0)$ can be taken
the same in Theorems~\ref{thm:21} and~\ref{thm:22} (Proposition~\ref{prop:reduction}).

\medskip
\noindent\textbf{Setting of Theorem~\ref{thm:23}.} Let $H$ be the matrix of an ensemble $\nu$ of size $N$ over $\kk$, and let $V$ be a Gaussian matrix independent of $H$, defined on a common probability space:
$V$ is Hermitian (symmetric if $\kk=\R$), the entries $(V_{ij})_{i\le j}$ are independent, $V_{ii}$ is a real centred Gaussian variable with variance $2/(\beta N)$, and for $i<j$:
$V_{ij}$ is a real centred Gaussian with variance $1/N$ if $\kk=\R$, and $V_{ij}=(\xi_{ij}+\ii\zeta_{ij})/\sqrt2$ with $\xi_{ij},\zeta_{ij}$ independent real centred Gaussians of variance $1/N$ if $\kk=\C$.
(This is the GOE for $\beta=1$ and the GUE for $\beta=2$; $\EE|V_{ij}|^2=1/N$ for $i<j$.) For $t\ge0$ put
\[
 H_t:=e^{-t/2}H+(1-e^{-t})^{1/2}V .
\]
The matrix $H_t$ has the law at time $t$ of the matrix Ornstein--Uhlenbeck flow started from $H$ (Dyson Brownian motion at the level of eigenvalues); NP uses only this formula, which is (2.29) of the original.
For a Hermitian matrix $A$ write $\lambda(A):=(\lambda_1(A),\dots,\lambda_N(A))\in\R^N$ for its ordered eigenvalues.
For a bounded Borel function $O:\R^n\to\R$, $1\le n\le N$, $E'\in\R$, $\varrho>0$ and $x\in\R^N$ put
\[
 \mathcal X^{O}_{E',\varrho}(x):=\sum_{\substack{i_1,\dots,i_n\in\{1,\dots,N\}\\ \text{pairwise distinct}}}O\big(N\varrho(x_{i_1}-E'),\dots,N\varrho(x_{i_n}-E')\big),\qquad C_{N,n}:=\frac{N^n(N-n)!}{N!}.
\]
For $O:\R^n\to\R$ let $\norm O_\infty:=\sup|O|$ and $\operatorname{Lip}(O):=\sup_{u\ne v}|O(u)-O(v)|/|u-v|$ (Euclidean norm on $\R^n$).

\begin{theorem}[Strong local ergodicity of Dyson Brownian motion]\label{thm:23}
Let $\vth,\delta_-,\delta_+,C_0>0$, $\eps'>0$, $\delta>0$, $0<c\le2$, $n\ge1$ and $\rho_O>0$.
There is a constant $C$ depending only on $(\eps',\delta,c,n,\rho_O,\vth,\delta_-,\delta_+,C_0)$ such that for all $N\ge\max(n,3)$,
$\kk\in\{\R,\C\}$, $\nu\in\mathcal E_N(\vth,\delta_-,\delta_+,C_0)$, every Lipschitz function $O:\R^n\to\R$ with support in $[-\rho_O,\rho_O]^n$, and all
$E\in[-2+c,2-c]$, $b\in(0,c/2)$ and $t\ge N^{-1+\delta+\eps'}$, with $\varrho:=\rhosc(E)$:
\begin{multline}\label{eq:abstrthm2}
 \Big|\int_{E-b}^{E+b}\frac{\rd E'}{2b}\;C_{N,n}\Big(\EE\,\mathcal X^{O}_{E',\varrho}\big(\lambda(H_t)\big)-\EE\,\mathcal X^{O}_{E',\varrho}\big(\lambda(V)\big)\Big)\Big|\\
 \le C\,\big(\norm O_\infty+\operatorname{Lip}(O)\big)\,N^{2\eps'}\Big[b^{-1}N^{-1+\eps'}+b^{-1/2}N^{-\delta/2}\Big].
\end{multline}
\end{theorem}
\noindent\emph{Comments on the formulation.}
(a) \emph{Correlation functions.} For a Borel probability measure $P$ on $\R^N$ let $P^{\rm sym}:=\frac1{N!}\sum_{\pi}P\circ\pi^{-1}$ be its symmetrisation (the sum is over the
$N!$ permutations of the coordinates) and let $p^{(n)}[P]$ be the marginal of $P^{\rm sym}$ on the first $n$ coordinates. For every bounded Borel $F:\R^n\to\R$,
\begin{align*}
 N^n\int_{\R^n}F\,\rd p^{(n)}[P]&=\frac{N^n}{N!}\sum_\pi\int_{\R^N}F(x_{\pi(1)},\dots,x_{\pi(n)})\,P(\rd x)\\
 &=C_{N,n}\int_{\R^N}\ \sum_{\substack{i_1,\dots,i_n\\ \text{pairwise distinct}}}F(x_{i_1},\dots,x_{i_n})\,P(\rd x),
\end{align*}
because each $n$-tuple of pairwise distinct indices equals $(\pi(1),\dots,\pi(n))$ for exactly $(N-n)!$ permutations $\pi$. Hence, with
$p^{(n)}_{t,N}:=p^{(n)}[\text{law of }\lambda(H_t)]$ and $p^{(n)}_{V,N}:=p^{(n)}[\text{law of }\lambda(V)]$, the quantity inside the absolute value in \eqref{eq:abstrthm2} is
\[
 \int_{E-b}^{E+b}\frac{\rd E'}{2b}\,N^n\!\int_{\R^n}O\big(N\varrho(x_1-E'),\dots,N\varrho(x_n-E')\big)\,\big(p^{(n)}_{t,N}-p^{(n)}_{V,N}\big)(\rd x_1\cdots\rd x_n).
\]
When $p^{(n)}_{t,N}$ and $p^{(n)}_{V,N}$ have densities, the change of variables $\alpha_i=N\varrho(x_i-E')$ turns this into the left side of (2.36) of the original
(with $N^n\,\rd x=\varrho^{-n}\rd\alpha$). The formulation \eqref{eq:abstrthm2} needs no densities. (The integrand is a bounded continuous function of $E'$, since $O$ is bounded and continuous.)

(b) \emph{The reference measure.} The law of $\lambda(V)$ is the Gaussian $\beta$-ensemble: Part~III proves that it has the density
$Z_{N,\beta}^{-1}\prod_{i<j}(x_j-x_i)^\beta\exp\big(-\frac{\beta N}4\sum_ix_i^2\big)$ on $\{x_1<\dots<x_N\}$, so that $p^{(n)}_{V,N}$ is the $n$-point correlation measure $p^{(n)}_{\mu,N}$ of the equilibrium measure $\mu\propto\exp\big(-\beta N\big[\sum_i\frac{x_i^2}4-\frac1N\sum_{i<j}\log|x_j-x_i|\big]\big)$ that appears in (2.36) of the original.
Theorem~\ref{thm:23} does not use this identification.

(c) \emph{The class of observables.} The theorem is stated for Lipschitz $O$, with a constant that depends on $O$ only through $\rho_O$ and, linearly, through $\norm O_\infty+\operatorname{Lip}(O)$.
Since the left side of \eqref{eq:abstrthm2} is linear in $O$, this is equivalent to the statement for observables normalised by $\norm O_\infty+\operatorname{Lip}(O)\le1$.
The original states the theorem for continuous $O$ with a constant depending on $O$; see Appendix~\ref{app:diff}, item PD-21. A version for H\"older continuous $O$, with a weaker power of $N$, is a corollary proved in Part~III.

(d) For $\kk=\R$ the diagonal entries of $V$ have variance $2/N$, so that $H_t$ does not satisfy \textup{(A)} exactly; Part~III deals with this by a rescaling (cf.\ Lemma~\ref{lem:mono}(v)).

\begin{theorem}[Universality of extreme eigenvalues]\label{thm:24}
For every $\vth>0$ there exist $\eps>0$ and $\delta>0$ such that for all $\delta_-,\delta_+,C_0>0$ there is $N_0$ with the following property.
Let $N\ge N_0$, $\kk\in\{\R,\C\}$, and let $\nu^{\mathbf v},\nu^{\mathbf w}\in\mathcal E_N(\vth,\delta_-,\delta_+,C_0)$ be two ensembles over $\kk$ whose first two moments match:
\begin{equation}\label{eq:2m}
 \int\ol y^{\,l}y^{u}\,\nu^{\mathbf v}_{ij}(\rd y)=\int\ol y^{\,l}y^{u}\,\nu^{\mathbf w}_{ij}(\rd y)\qquad\text{for all }(i,j)\in\mathcal P_N\text{ and all integers }l,u\ge0\text{ with }l+u\le2 .
\end{equation}
Write $\PP^{\mathbf v}:=\PP_{\nu^{\mathbf v}}$, $\PP^{\mathbf w}:=\PP_{\nu^{\mathbf w}}$. Then for all $s\in\R$
\begin{equation}\label{eq:tw}
 \PP^{\mathbf v}\big(N^{2/3}(\lambda_N-2)\le s-N^{-\eps}\big)-N^{-\delta}\le\PP^{\mathbf w}\big(N^{2/3}(\lambda_N-2)\le s\big)\le
 \PP^{\mathbf v}\big(N^{2/3}(\lambda_N-2)\le s+N^{-\eps}\big)+N^{-\delta},
\end{equation}
and
\begin{equation}\label{eq:tw1}
 \PP^{\mathbf v}\big(N^{2/3}(\lambda_1+2)\ge s+N^{-\eps}\big)-N^{-\delta}\le\PP^{\mathbf w}\big(N^{2/3}(\lambda_1+2)\ge s\big)\le
 \PP^{\mathbf v}\big(N^{2/3}(\lambda_1+2)\ge s-N^{-\eps}\big)+N^{-\delta}.
\end{equation}
\end{theorem}
\noindent Here $N_0$ does not depend on $s$. By \eqref{eq:2m} with $(l,u)=(1,1)$ the two ensembles have the same matrix of variances $B$.
Theorems~\ref{thm:23} and~\ref{thm:24} are proved in Parts~III and~II.

\subsection{The level form of Theorems 2.1 and 2.2}
\begin{lemma}[Elementary properties of the class]\label{lem:mono}
Let $\nu$ be an ensemble of size $N$ over $\kk$ with $\sigma_{ij}<\infty$ for all $i,j$.
\begin{enumerate}[label=\textup{(\roman*)},leftmargin=2.2em,itemsep=1pt]
\item Assume \textup{(A)}, and let $\delta_->0$, $\delta_+>0$. Then \textup{(B)} is equivalent to the spectral condition \textup{(2.3)} of the original:
$1$ is a simple eigenvalue of $B$ and $\operatorname{Spec}(B)\subset[-1+\delta_-,1-\delta_+]\cup\{1\}$.
\item \emph{Monotonicity:} $\mathcal E_N(\vth,\delta_-,\delta_+,C_0)\subset\mathcal E_N(\vth',\delta_-',\delta_+',C_0')$ if
$0<\vth'\le\vth$, $0<\delta_\pm'\le\delta_\pm$, $C_0'\ge C_0$. In particular
$\mathcal E_N(\vth,\delta_-,\delta_+,C_0)\subset\mathcal E_N\big(\min(\vth,1),\min(\delta_-,1),\min(\delta_+,1),\max(C_0,1)\big)$.
\item If $\mathcal E_N(\vth,\delta_-,\delta_+,C_0)\ne\emptyset$ then $C_0\ge1$ and (for $N\ge2$) $\delta_-+\delta_+\le2$.
\item \emph{Reflection:} if $\nu\in\mathcal E_N(\vth,\delta_-,\delta_+,C_0)$ then the ensemble $\nu^-$, $\nu^-_{ij}(S):=\nu_{ij}(-S)$, is in the same class, with the same $B$;
$H(-x)=-H(x)$, and the law of $H$ under $\PP_{\nu^-}$ is the law of $-H$ under $\PP_\nu$.
\item \emph{Scaling}: let $a\in\R$ and let $\nu^{a}$ be the ensemble $\nu^a_{ij}:=$ image of $\nu_{ij}$ under $y\mapsto ay$. Then $H(ax)=aH(x)$, the law of $H$ under $\PP_{\nu^a}$ is the law of $aH$ under $\PP_\nu$,
the variances of $\nu^a$ are $a^2\sigma_{ij}^2$, and each of the conditions \textup{(R)}, \textup{(Z)} and \textup{(D)} for $\nu$ implies the same condition for $\nu^a$,
\emph{with the same $\vth$} in \textup{(D)}. The same holds for the image of a single $\nu_{ij}$ under $y\mapsto ay$ with $a\in\kk$ $(i<j)$.
\end{enumerate}
\end{lemma}
\begin{proof}
Recall that $B$ is a real symmetric matrix with nonnegative entries. Let $\be:=N^{-1/2}(1,\dots,1)$ and $\be^\perp:=\{v\in\R^N:\sum_iv_i=0\}$.

(i) By (A) and the symmetry of $B$, $(B\be)_i=N^{-1/2}\sum_j\sigma_{ij}^2=N^{-1/2}$, i.e.\ $B\be=\be$. Consequently $\langle\be,Bv\rangle=\langle B\be,v\rangle=\langle\be,v\rangle$, so $B$ maps $\be^\perp$ into itself.

\emph{\textup{(B)} implies the spectral condition.} Let $\mu\in\R$ be an eigenvalue of $B$ and $u\in\R^N\setminus\{0\}$ a real eigenvector, $Bu=\mu u$. Write $u=a\be+v$ with $a\in\R$, $v\in\be^\perp$. Then $a\be+Bv=Bu=\mu a\be+\mu v$ with $Bv\in\be^\perp$, hence
$a(1-\mu)=0$ and $Bv=\mu v$. If $\mu\ne1$, then $a=0$, $v=u\ne0$ and $\mu|v|^2=\langle v,Bv\rangle\in[(-1+\delta_-)|v|^2,(1-\delta_+)|v|^2]$ by (B), so $\mu\in[-1+\delta_-,1-\delta_+]$.
If $\mu=1$, then $|v|^2=\langle v,Bv\rangle\le(1-\delta_+)|v|^2$, so $v=0$ because $\delta_+>0$; thus the eigenspace of the eigenvalue $1$ is $\R\be$, and since $B$ is symmetric (hence diagonalisable), $1$ is a simple eigenvalue.

\emph{The spectral condition implies \textup{(B)}.} By the spectral theorem there is an orthonormal basis $u_1,\dots,u_N$ of $\R^N$ with $Bu_\alpha=\mu_\alpha u_\alpha$. Since $1$ is a simple eigenvalue and $B\be=\be$,
there is exactly one index $\alpha_0$ with $\mu_{\alpha_0}=1$, and $u_{\alpha_0}=\pm\be$. For $\alpha\ne\alpha_0$ we have $\mu_\alpha\ne1$, hence $\mu_\alpha\in[-1+\delta_-,1-\delta_+]$.
If $v\in\be^\perp$, then $v=\sum_{\alpha\ne\alpha_0}c_\alpha u_\alpha$ with $c_\alpha=\langle u_\alpha,v\rangle$, $|v|^2=\sum_{\alpha\ne\alpha_0}c_\alpha^2$ and
$\langle v,Bv\rangle=\sum_{\alpha\ne\alpha_0}\mu_\alpha c_\alpha^2\in[(-1+\delta_-)|v|^2,(1-\delta_+)|v|^2]$.

(ii) Conditions (R), (Z), (A) do not involve the parameters. (C) and (B) are weakened when $C_0$ is increased and $\delta_\pm$ are decreased. For \eqref{eq:subexp}: if $x\ge1$ and $0<\vth'\le\vth$ then $x^{\vth'}\le x^{\vth}$ and $\vth^{-1}\le\vth'^{-1}$, so
$\vth^{-1}\exp(-x^\vth)\le\vth'^{-1}\exp(-x^{\vth'})$. The second statement is the case $\vth'=\min(\vth,1)$, $\delta'_\pm=\min(\delta_\pm,1)$, $C_0'=\max(C_0,1)$.

(iii) By (A) and (C), $1=\sum_{i}\sigma_{ij}^2\le N\cdot C_0/N=C_0$. If $N\ge2$ there is $v\in\be^\perp$ with $|v|=1$, and (B) gives $-1+\delta_-\le\langle v,Bv\rangle\le1-\delta_+$.

(iv) $H(-x)=-H(x)$ is clear from Definition~\ref{def:ens}. The map $x\mapsto-x$ sends the product measure $\PP_\nu$ to the product of the images, which is $\PP_{\nu^-}$; so for every Borel set $S$ of matrices
$\PP_{\nu^-}(H\in S)=\PP_\nu(H(-x)\in S)=\PP_\nu(-H\in S)$. Since $|-y|=|y|$, the measures $\nu^-_{ij}$ have the same second moments as $\nu_{ij}$, so $B$ and the conditions (A), (B), (C) are unchanged, and
\eqref{eq:subexp} is unchanged; $\nu^-_{ii}(\R)=\nu_{ii}(\R)=1$ and $\int y\,\nu^-_{ij}(\rd y)=-\int y\,\nu_{ij}(\rd y)=0$.

(v) Let $\mu$ be a Borel probability measure on $\kk$ with $s^2:=\int|y|^2\rd\mu<\infty$, let $a\in\kk$ and let $\mu^a$ be the image of $\mu$ under $y\mapsto ay$. Then $\int|y|^2\rd\mu^a=\int|ay|^2\rd\mu=|a|^2s^2$, so the standard deviation of $\mu^a$ is $|a|s$,
and $\int y\,\rd\mu^a=a\int y\,\rd\mu$. Suppose $\mu(\{|y|>xs\})\le\vth^{-1}e^{-x^\vth}$ for all $x\ge1$. If $a\ne0$ then for $x\ge1$
\[
 \mu^a\big(\{y:|y|>x|a|s\}\big)=\mu\big(\{y:|ay|>x|a|s\}\big)=\mu\big(\{y:|y|>xs\}\big)\le\vth^{-1}e^{-x^\vth};
\]
if $a=0$ then $\mu^a$ is the point mass at $0$ and $\mu^a(\{|y|>0\})=0$. So $\mu^a$ satisfies the same bound, with the same $\vth$, relative to its own standard deviation $|a|s$. If $a\in\R$ and $\mu(\R)=1$ then $\mu^a(\R)=1$.
Applying this to every $\nu_{ij}$ with a real $a$ gives the statements on $\nu^a$; $H(ax)=aH(x)$ holds for real $a$ by Definition~\ref{def:ens}, and the statement on the laws follows as in (iv).
\end{proof}
\noindent Under $H\mapsto aH$ the conditions (A), (B), (C) are of course not preserved ($B$ becomes $a^2B$); this is why the rescaling in Part~III has to be combined with a direct verification of (A)--(C).


Part~I proves the following statement, in which the single parameter $K$ replaces the powers of $\log N$.
Let (recall $\fa=1/\vth\ge1$ after the normalisation of Convention~\ref{conv:const})
\begin{equation}\label{eq:ledger0}
 C_1:=18\fa+20,\quad C_2:=12\fa+19,\quad C_3:=24\fa+36,\quad C_5:=18\fa+21,\quad C_4:=30\fa+44,\quad c_0:=\frac1{8C_3}.
\end{equation}
These values are fixed by the exponent ledger in Appendix~\ref{app:ledger}. The reduction below (Proposition~\ref{prop:reduction}) uses only that $C_1,\dots,C_5\ge1$ and $c_0>0$ are finite numbers depending on $\vth$ alone;
the proof of parts (iv), (v) in Section~\ref{sec:rig} uses, in addition, only the relations listed in \eqref{eq:ledgerrel} there. For $K\ge1$ let
\begin{equation}\label{eq:DK}
 \bD_K:=\big\{z=E+\ii\eta:\ |E|\le5,\ K^{C_1}N^{-1}\le\eta\le10\big\}.
\end{equation}

\begin{theorem}[Local law and rigidity at level $K$]\label{thm:K}
For all $\vth,\delta_-,\delta_+,C_0>0$, normalised as in Convention~\ref{conv:const}, there is $N_0$ such that for every $N\ge N_0$, $\kk\in\{\R,\C\}$,
$\nu\in\mathcal E_N(\vth,\delta_-,\delta_+,C_0)$ and every real $K$ with
\begin{equation}\label{eq:Krange}
 (\log N)^3\le K\le N^{c_0}
\end{equation}
there is a measurable event $\Xi_K\subset\Omega$ with $\PP(\Xi_K^c)\le e^{-K/2}$ on which:
\begin{enumerate}[label=\textup{(\roman*)},leftmargin=2.2em,itemsep=1pt]
\item for all $z\in\bD_K$:\quad $\Lambda(z)\le K^{C_2}/(N\eta)$;
\item for all $z\in\bD_K$:\quad $\Ld(z)+\Lo(z)\le K^{C_2}\big(\sqrt{\im\msc(z)/(N\eta)}+1/(N\eta)\big)$;
\item $\max_j|\lambda_j|\le2+K^{C_3}N^{-2/3}$;
\item $\sup_{E\in\R}|\fn(E)-\nsc(E)|\le K^{C_5}/N$;
\item for all $j\in\{1,\dots,N\}$:\quad $|\lambda_j-\gamma_j|\le K^{C_4}\,[\min(j,N+1-j)]^{-1/3}N^{-2/3}$.
\end{enumerate}
\end{theorem}
\noindent Parts (i), (ii) are proved in Section~\ref{sec:strong}, part (iii) in Section~\ref{sec:norm}, where the event $\Xi_K$ is defined; parts (iv) and (v) are deterministic consequences of (i) and (iii),
proved in Section~\ref{sec:rig} (Proposition~\ref{prop:rig}).

\begin{proposition}[Reduction]\label{prop:reduction}
Theorem~\ref{thm:K} implies Theorems~\ref{thm:21} and~\ref{thm:22}, with the same constants in both theorems: if $(\vth,\delta_-,\delta_+,C_0)$ are given and
$\bar\vth:=\min(\vth,1)$, $\bar\delta_\pm:=\min(\delta_\pm,1)$, $\bar C_0:=\max(C_0,1)$ are the normalised parameters, $C_1,\dots,C_5,c_0$ are defined by \eqref{eq:ledger0} with $\fa=1/\bar\vth$, then one can take
\[
 \phi:=\min\Big\{\frac{10}{C_1},\ \frac{2}{C_2},\ \frac{4}{C_3},\ \frac1{2C_4},\ \frac1{2C_5},\ 5c_0\Big\},\qquad A_0:=\frac3\phi,\qquad C:=1,\qquad c:=\frac12,
\]
and $N_0:=\max\{16,\ N_0'\}$, where $N_0'$ is the threshold of Theorem~\ref{thm:K} for the normalised parameters.
\end{proposition}
\begin{proof}
Let $(\vth,\delta_-,\delta_+,C_0)$ be given. By Lemma~\ref{lem:mono}(ii), $\mathcal E_N(\vth,\delta_-,\delta_+,C_0)\subset\mathcal E_N(\bar\vth,\bar\delta_-,\bar\delta_+,\bar C_0)$, and the normalised parameters
satisfy Convention~\ref{conv:const}. So Theorem~\ref{thm:K} applies to every $\nu\in\mathcal E_N(\vth,\delta_-,\delta_+,C_0)$ with the constants of the normalised parameters, which depend only on $(\vth,\delta_-,\delta_+,C_0)$.
Note that $C_1,\dots,C_5\ge1$, so $\phi\le10/C_1<1$ (as $C_1\ge38$), $\phi>0$ and $A_0=3/\phi>1$.

Let $N\ge N_0$ and let $L$ satisfy \eqref{eq:Lbound}. Since $N\ge16>e^e$ we have $\log N>e$ and $\ell:=\log\log N>1$; in particular $L\ge A_0\ell>0$. Put
\[
 K:=(\log N)^{\phi L}=\exp(\phi L\ell).
\]
\emph{Step 1: $K$ satisfies \eqref{eq:Krange}.} From $L\ge A_0\ell$ and $\ell\ge1$: $\phi L\ell\ge\phi A_0\ell^2=3\ell^2\ge3\ell$, so $K\ge e^{3\ell}=(\log N)^3$.
From the upper bound in \eqref{eq:Lbound}: $10L\ell\le\log(10N)$, so $(\log N)^{10L}=\exp(10L\ell)\le10N$ and
\[
 K=\big((\log N)^{10L}\big)^{\phi/10}\le(10N)^{\phi/10}\le N^{\phi/5}\le N^{c_0},
\]
where we used $10N\le N^2$ (as $N\ge10$) and $\phi\le5c_0$.

\emph{Step 2: comparison of the powers.} Since $\log N>1$ and $L>0$, the function $s\mapsto(\log N)^{sL}$ is strictly increasing. By the definition of $\phi$,
$\phi C_1\le10$, $\phi C_2\le2$, $\phi C_3\le4$, $\phi C_4\le\frac12$, $\phi C_5\le\frac12$. Hence
\begin{gather*}
 K^{C_1}\le(\log N)^{10L},\qquad K^{C_2}\le(\log N)^{2L}<(\log N)^{4L},\qquad K^{C_3}\le(\log N)^{4L}<(\log N)^{9L},\\
 K^{C_4}\le(\log N)^{L/2}<(\log N)^{L},\qquad K^{C_5}\le(\log N)^{L/2}<(\log N)^{L}.
\end{gather*}
In particular, if $z=E+\ii\eta\in\bS_L$ then $|E|\le5$ and $10\ge\eta>N^{-1}(\log N)^{10L}\ge K^{C_1}N^{-1}$, so $\bS_L\subset\bD_K$.

\emph{Step 3: conclusion.} Let $\Xi_K$ be the event of Theorem~\ref{thm:K} for this $K$. At every point of $\Xi_K$ we have, by Theorem~\ref{thm:K}(i)--(v) and Step 2:
\begin{itemize}[leftmargin=2em,itemsep=1pt]
\item for all $z\in\bS_L$: $\Lambda(z)\le K^{C_2}/(N\eta)<(\log N)^{4L}/(N\eta)$;
\item for all $z\in\bS_L$: $\Ld(z)+\Lo(z)\le K^{C_2}\big(\sqrt{\im\msc/(N\eta)}+1/(N\eta)\big)<(\log N)^{4L}\sqrt{\im\msc/(N\eta)}+(\log N)^{4L}/(N\eta)$ (strict because $1/(N\eta)>0$ and $\im\msc\ge0$);
\item $\max_j|\lambda_j|\le2+K^{C_3}N^{-2/3}<2+N^{-2/3}(\log N)^{9L}$;
\item for all $j$: $|\lambda_j-\gamma_j|\le K^{C_4}[\min(j,N+1-j)]^{-1/3}N^{-2/3}<(\log N)^{L}[\min(j,N-j+1)]^{-1/3}N^{-2/3}$;
\item $\sup_{|E|\le5}|\fn(E)-\nsc(E)|\le\sup_{E\in\R}|\fn(E)-\nsc(E)|\le K^{C_5}/N<(\log N)^L/N$.
\end{itemize}
Therefore each of the sets whose probability is estimated in \eqref{eq:Lambdafinal}, \eqref{eq:Lambdaodfinal}, \eqref{eq:443}, \eqref{eq:rigidity}, \eqref{eq:nn} is contained in the measurable set $\Xi_K^c$.
(The sets in \eqref{eq:443}, \eqref{eq:rigidity} and \eqref{eq:nn} are themselves measurable, by Lemma~\ref{lem:weyl}(d),(f); for the unions over $z\in\bS_L$ in \eqref{eq:Lambdafinal}, \eqref{eq:Lambdaodfinal} we use the outer probability.)
Its (outer) probability is thus at most
\[
 \PP(\Xi_K^c)\le e^{-K/2}=\exp\big[-\tfrac12(\log N)^{\phi L}\big],
\]
which is the asserted bound with $C=1$, $c=\frac12$. The last sentence of Theorem~\ref{thm:21} holds with the event $\Xi_K$.
\end{proof}

\begin{remark}\label{rem:Kstronger}
Theorem~\ref{thm:K} is slightly stronger than Theorems~\ref{thm:21}--\ref{thm:22}: the lower bound $L\ge A_0\log\log N$ is used only through $K\ge(\log N)^3$,
the estimate \textup{(iv)} holds for all real $E$, the level $K$ may be as large as a small power of $N$ (which gives error bounds $N^{\eps}$ with probability $1-\exp(-\frac12N^{\eps/C})$, see Section~\ref{sec:export}),
and the domain may be taken larger ($|E|\le\Emax+1$, see Theorem~\ref{thm:weak} and Section~\ref{sec:export}).
At the upper endpoint of \eqref{eq:Lbound}, $(\log N)^{10L}=10N$ and the set $\bS_L$ is empty; then \textup{(i)}, \textup{(ii)} of Theorem~\ref{thm:21} hold trivially.
\end{remark}

\section{Deterministic toolbox}\label{sec:det}
Everything in this section is deterministic: $H$ is a fixed Hermitian $N\times N$ matrix (with entries in $\kk\subset\C$) and $B$ a fixed matrix satisfying (A), (B), (C).

\subsection{The function $\msc$ and the semicircle law}
We use the following facts about the principal square root $\sqrt\zeta$, $\zeta\in\C\setminus(-\infty,0]$ (the unique $s$ with $s^2=\zeta$, $\re s>0$).
Let $\mathbb Q_1:=\{s\in\C:\re s>0,\ \im s>0\}$ be the open first quadrant.
\begin{itemize}[leftmargin=2em,itemsep=1pt]
\item[(S1)] If $\im\zeta>0$ then $\sqrt\zeta\in\mathbb Q_1$. Indeed, for $s=a+\ii b=\sqrt\zeta$ we have $a>0$ and $2ab=\im\zeta>0$, so $b>0$.
\item[(S2)] If $\im\zeta>0$ then $\sqrt{-\bar\zeta}=\ii\,\ol{\sqrt\zeta}$. Indeed, with $s=a+\ii b=\sqrt\zeta\in\mathbb Q_1$: $(\ii\bar s)^2=-\bar s^2=-\bar\zeta$ and $\ii\bar s=b+\ii a$ has positive real part.
\item[(S3)] If $s,s'\in\mathbb Q_1$ then $\im(ss')=\re s\,\im s'+\im s\,\re s'>0$.
\item[(S4)] $\zeta\mapsto\sqrt\zeta$ is holomorphic on $\C\setminus(-\infty,0]$ (Mathlib).
\end{itemize}

\begin{lemma}[Algebra of $\msc$]\label{lem:msc} \infra\
Let $z=E+\ii\eta$ with $\eta>0$, $w:=\sqrt{z-2}\sqrt{z+2}$ (principal roots) and $m:=\msc(z)=(-z+w)/2$. Then:
\begin{enumerate}[label=\textup{(\alph*)},leftmargin=2.2em,itemsep=1pt]
\item $m^2+zm+1=0$, $\im m>0$, $m\ne0$, $m=-1/(z+m)$, $-z-m=1/m$; and $m$ is the only root of $X^2+zX+1$ with positive imaginary part (the other root is $1/m$).
\item $\im m\,(|z+m|^2-1)=\eta$; hence $|m|<1$, $|m|^2=\frac{\im m}{\im m+\eta}$, $1-|m|^2=\frac{\eta}{\im m+\eta}$ and $\im m=\eta\frac{|m|^2}{1-|m|^2}$.
\item $\dfrac1{|z|+1}\le|m|\le1$ and $\im m\ge\dfrac{\eta}{(|z|+1)^2}$.
\item $w=z+2m$, $w^2=z^2-4$ and $1-m^2=-m\,w$; hence $|1-m^2|=|m|\,|z-2|^{1/2}|z+2|^{1/2}$.
\item With $P:=|z-2||z+2|-|z|^2+4$ and $Q:=|z-2||z+2|+|z|^2-4$:
\[
 (\im w)^2=\eta^2+\tfrac12P,\qquad \im m=\tfrac12(\im w-\eta)=\dfrac{P}{4(\im w+\eta)},\qquad PQ=16\eta^2 ;
\]
in particular $P>0$ and $Q>0$.
\item \emph{Lipschitz:} for $\im z,\im z'\ge\eta_0>0$:\quad $|\msc(z)-\msc(z')|\le\frac{\eta_0+1}{\eta_0}|z-z'|$.
\item \emph{Far regime:} if $|E|\ge4$ then $|m|\le1/3$, $\im m\le\eta/8$ and $|1-m^2|\ge8/9$.
\item \emph{Symmetry:} $\msc(-\bar z)=-\ol{\msc(z)}$; hence $|\msc|$, $\im\msc$ and $|1-\msc^2|$ take the same values at $E+\ii\eta$ and at $-E+\ii\eta$.
\item $\msc$ is holomorphic on $\C_+$.
\end{enumerate}
\end{lemma}
\begin{proof}
Let $p:=\sqrt{z-2}$, $q:=\sqrt{z+2}$. Since $\im(z\pm2)=\eta>0$, (S1) gives $p,q\in\mathbb Q_1$, and $w=pq$ satisfies $w^2=p^2q^2=(z-2)(z+2)=z^2-4$.

(a) We compute
\[
 m^2+zm+1=\tfrac14(z^2-2zw+w^2)+\tfrac12(-z^2+zw)+1=\tfrac14(w^2-z^2+4)=0 .
\]
Next we show $b:=\im w>\eta$. By (S3), $b>0$. Write $w=a+\ii b$. Taking real and imaginary parts of $w^2=z^2-4$: $a^2-b^2=E^2-\eta^2-4$ and $ab=E\eta$.
Suppose $b\le\eta$. Then $|a|=|E|\eta/b\ge|E|$, so $a^2-b^2\ge E^2-\eta^2>E^2-\eta^2-4$, a contradiction. Hence $b>\eta$ and $\im m=\frac12(b-\eta)>0$.
From $m(m+z)=-1$ we get $m\ne0$, $z+m\ne0$, $m=-1/(z+m)$ and $-z-m=1/m$. Finally $(X-m)(X-1/m)=X^2-(m+1/m)X+1=X^2+zX+1$, so the roots are $m$ and $1/m$, and
$\im(1/m)=-\im m/|m|^2<0$.

(b) By (a), $m=-1/(z+m)=-\ol{(z+m)}/|z+m|^2$, so $\im m=\im(z+m)/|z+m|^2=(\eta+\im m)/|z+m|^2$, which is the first identity. Since $\im m>0$ and $\eta>0$ it gives $|z+m|>1$, hence
$|m|=1/|z+m|<1$ and $|m|^2=|z+m|^{-2}=\im m/(\im m+\eta)$. The other two formulas are rearrangements.

(c) $|m|\le1$ by (b), and $|m|=1/|z+m|\ge1/(|z|+|m|)\ge1/(|z|+1)$. By (b), $\im m=\eta/(|z+m|^2-1)\ge\eta/((|z|+1)^2-1)\ge\eta/(|z|+1)^2$.

(d) $w=z+2m$ by the definition of $m$. By (a), $1=-m^2-zm$, so $1-m^2=-2m^2-zm=-m(2m+z)=-mw$. Taking absolute values, $|w|=|p||q|=|z-2|^{1/2}|z+2|^{1/2}$.

(e) For any complex number $w=a+\ii b$: $b^2=\frac12(|w|^2-\re(w^2))$. Here $|w|^2=|w^2|=|z-2||z+2|$ and $\re(w^2)=\re(z^2-4)=E^2-\eta^2-4=|z|^2-4-2\eta^2$. Hence
$(\im w)^2=\frac12(|z-2||z+2|-|z|^2+4)+\eta^2=\eta^2+\frac12P$. By the proof of (a), $\im w>\eta>0$, so $P=2((\im w)^2-\eta^2)>0$ and
\[
 \im m=\tfrac12(\im w-\eta)=\frac{(\im w)^2-\eta^2}{2(\im w+\eta)}=\frac{P}{4(\im w+\eta)} .
\]
Finally $PQ=|z^2-4|^2-(|z|^2-4)^2=\big(|z|^4-8\re(z^2)+16\big)-\big(|z|^4-8|z|^2+16\big)=8(|z|^2-\re z^2)=16\eta^2$, and $Q=16\eta^2/P>0$.

(f) Let $m=\msc(z)$, $m'=\msc(z')$. By (a), $1/(z+m)=-m$ and $1/(z'+m')=-m'$, so
\[
 m-m'=\frac1{z'+m'}-\frac1{z+m}=\frac{(z-z')+(m-m')}{(z+m)(z'+m')}=\big((z-z')+(m-m')\big)\,mm',
\]
i.e.\ $(m-m')(1-mm')=(z-z')mm'$. By (b) and $0<\im m\le|m|<1$: $1-|m|^2=\frac{\im z}{\im z+\im m}\ge\frac{\im z}{\im z+1}\ge\frac{\eta_0}{\eta_0+1}$, since $t\mapsto t/(t+1)$ is increasing; the same holds for $m'$.
Hence $|1-mm'|\ge1-|m||m'|\ge1-\max(|m|^2,|m'|^2)\ge\frac{\eta_0}{\eta_0+1}$, and $|mm'|\le1$. This gives the claim.

(g) If $|E|\ge4$ then $|z+m|\ge|z|-|m|\ge|E|-1\ge3$, so $|m|=1/|z+m|\le\frac13$; by (b), $\im m=\eta/(|z+m|^2-1)\le\eta/8$; and $|1-m^2|\ge1-|m|^2\ge\frac89$.

(h) $-\bar z\pm2=-\ol{(z\mp2)}$ and $\im(z\mp2)>0$, so by (S2): $\sqrt{-\bar z-2}=\ii\,\ol{\sqrt{z+2}}$ and $\sqrt{-\bar z+2}=\ii\,\ol{\sqrt{z-2}}$. The product is $-\ol{w}$, hence
$\msc(-\bar z)=\frac12(\bar z-\bar w)=-\ol{\msc(z)}$. The map $\mu\mapsto-\bar\mu$ preserves $|\mu|$, $\im\mu$ and $|1-\mu^2|$ (as $1-(-\bar\mu)^2=\ol{1-\mu^2}$).

(i) For $z\in\C_+$ the numbers $z\pm2$ lie in $\C_+\subset\C\setminus(-\infty,0]$; apply (S4).
\end{proof}

\begin{lemma}[Size of $\msc$ in the near regime]\label{lem:mscsize} \infra\
Let
\[
 c_*:=\tfrac1{150},\qquad C_*:=4 .
\]
For $|E|\le5$, $0<\eta\le10$ (recall $\kappa=||E|-2|$, $\varkappa=\sqrt{\kappa+\eta}$):
\begin{align}
 &\tfrac1{13}\le|\msc|\le1,\qquad c_*\varkappa\le|1-\msc^2|\le C_*\varkappa,\label{eq:alpha}\\
 &c_*\varkappa\le\im\msc\le C_*\varkappa\ \ \text{if }|E|\le2\text{ or }\kappa\le\eta;\qquad
 c_*\frac{\eta}{\varkappa}\le\im\msc\le C_*\frac{\eta}{\varkappa}\ \ \text{if }|E|>2\text{ and }\kappa\ge\eta.\label{eq:imm}
\end{align}
In particular $c_*\eta\le\im\msc\le\varkappa\le C_*\varkappa\le C_*^2c_*^{-1}|1-\msc^2|$ on this domain.
More precisely, on this domain:
$\frac{2^{1/4}}{13}\varkappa\le|1-\msc^2|\le149^{1/4}\varkappa$;
$\ \frac1{81}\varkappa\le\im\msc\le\varkappa$ if $|E|\le2$;
$\ \frac1{92}\varkappa\le\im\msc\le\frac1{\sqrt2}\varkappa$ if $|E|>2$, $\kappa<\eta$;
$\ \frac1{27}\frac\eta\varkappa\le\im\msc\le\frac1{\sqrt2}\frac\eta\varkappa$ if $|E|>2$, $\kappa\ge\eta$.
\end{lemma}
\begin{proof}
Write $m=\msc(z)$. By Lemma~\ref{lem:msc}(h) the quantities $|m|$, $\im m$, $|1-m^2|$ are unchanged under $E\mapsto-E$, and so are $\kappa$ and $\varkappa$. Hence we may assume $0\le E\le5$; then $\kappa=|E-2|$.

\emph{Proof of \eqref{eq:alpha}.} $|z|\le\sqrt{125}<11.2$, so $|m|\ge1/(|z|+1)\ge\frac1{13}$ by Lemma~\ref{lem:msc}(c), and $|m|\le1$. Since $E\ge0$: $2\le|z+2|\le|7+10\ii|=\sqrt{149}$, and $|z-2|=\sqrt{\kappa^2+\eta^2}$ satisfies
$(\kappa+\eta)/\sqrt2\le|z-2|\le\kappa+\eta$. By Lemma~\ref{lem:msc}(d), $|1-m^2|=|m|\,|z-2|^{1/2}|z+2|^{1/2}$, so
\[
 \tfrac{2^{1/4}}{13}\,\varkappa=\tfrac1{13}\big(\tfrac{\kappa+\eta}{\sqrt2}\cdot2\big)^{1/2}\le|1-m^2|\le\big((\kappa+\eta)\sqrt{149}\big)^{1/2}=149^{1/4}\varkappa,
\]
and $2^{1/4}/13>0.09>c_*$, $149^{1/4}<3.5<C_*$.

\emph{Preparations for \eqref{eq:imm}.} Let $R:=|z-2||z+2|\ge0$ and $S:=|z|^2-4=E^2+\eta^2-4$, so that $P=R-S$ and $Q=R+S$ in the notation of Lemma~\ref{lem:msc}(e). That lemma gives $P>0$, $Q>0$,
\begin{equation}\tag{F1}\label{eq:a1-F1}
 R^2-S^2=PQ=16\eta^2,\qquad\text{hence}\qquad \max(|S|,4\eta)\le R=\sqrt{S^2+16\eta^2}\le|S|+4\eta,
\end{equation}
and, with $\im w=\sqrt{\eta^2+P/2}$,
\begin{equation}\tag{F2}\label{eq:a1-F2}
 \im m=\tfrac12\Big(\sqrt{\eta^2+\tfrac12P}-\eta\Big)=\frac{P}{4\big(\sqrt{\eta^2+\frac12P}+\eta\big)} .
\end{equation}
By the first expression, for fixed $\eta$ the right side of \eqref{eq:a1-F2} is an increasing function of $P\in(0,\infty)$. Since $\sqrt{P/2}\le\sqrt{\eta^2+P/2}\le\eta+\sqrt{P/2}$, the second expression gives
\begin{equation}\tag{F3}\label{eq:a1-F3}
 \frac{P}{4\big(2\eta+\sqrt{P/2}\big)}\le\im m\le\frac{\sqrt P}{2\sqrt2},
\end{equation}
and the left side of \eqref{eq:a1-F3} is also increasing in $P$.

\emph{Case A: $E>2$ and $\kappa\ge\eta$.} Here $E=2+\kappa$, $0<\kappa\le3$, $\eta\le\kappa$, and $S=(E-2)(E+2)+\eta^2=\kappa(4+\kappa)+\eta^2$. So $S\ge4\kappa>0$ and, using $\eta^2\le\kappa^2\le3\kappa$, $S\le7\kappa+3\kappa=10\kappa$.
By \eqref{eq:a1-F1}, $2S\le Q=R+S\le2S+4\eta$, hence $8\kappa\le Q\le24\kappa$ and
\[
 \frac{2\eta^2}{3\kappa}\le P=\frac{16\eta^2}{Q}\le\frac{2\eta^2}{\kappa} .
\]
By \eqref{eq:a1-F3}, $\im m\le\sqrt{2\eta^2/\kappa}/(2\sqrt2)=\eta/(2\sqrt\kappa)$; and since $\sqrt{P/2}\le\eta/\sqrt\kappa$,
\[
 \im m\ge\frac{P}{4(2\eta+\eta/\sqrt\kappa)}=\frac{P\sqrt\kappa}{4\eta(2\sqrt\kappa+1)}\ge\frac{\eta}{6\sqrt\kappa\,(2\sqrt\kappa+1)}\ge\frac{\eta}{6(2\sqrt3+1)\sqrt\kappa}\ge\frac{\eta}{27\sqrt\kappa} .
\]
As $\kappa\le\varkappa^2=\kappa+\eta\le2\kappa$, i.e.\ $\varkappa/\sqrt2\le\sqrt\kappa\le\varkappa$, we get
\[
 \frac1{27}\,\frac\eta\varkappa\le\im m\le\frac1{\sqrt2}\,\frac\eta\varkappa\qquad(\text{Case A}).
\]

\emph{Case B: $E>2$ and $\kappa<\eta$.} Again $S=\kappa(4+\kappa)+\eta^2>0$, and now $S\le7\kappa+10\eta\le17\eta$ (using $\kappa\le3$, $\eta\le10$, $\kappa<\eta$). By \eqref{eq:a1-F1}, $4\eta\le R\le Q\le2S+4\eta\le38\eta$, hence
$\frac8{19}\eta\le P\le4\eta$. By \eqref{eq:a1-F3}, $\im m\le\sqrt{4\eta}/(2\sqrt2)=\sqrt{\eta/2}\le\varkappa/\sqrt2$. For the lower bound we use the monotonicity of the left side of \eqref{eq:a1-F3} in $P$ and $\eta\le10$:
\[
 \im m\ge\frac{\frac8{19}\eta}{4\big(2\eta+\sqrt{4\eta/19}\big)}=\frac{2\sqrt\eta}{19\big(2\sqrt\eta+2/\sqrt{19}\big)}\ge\frac{2\sqrt\eta}{19\,(2\sqrt{10}+0.46)}\ge\frac{\sqrt\eta}{65} .
\]
Since $\varkappa^2=\kappa+\eta<2\eta$, $\sqrt\eta>\varkappa/\sqrt2$ and
\[
 \frac1{92}\,\varkappa\le\im m\le\frac1{\sqrt2}\,\varkappa\qquad(\text{Case B}).
\]

\emph{Case C: $0\le E\le2$.} Here $\kappa=2-E\in[0,2]$ and $E^2-4=-\kappa(4-\kappa)$, so $S=\eta^2-\kappa(4-\kappa)$, which may have either sign.

\emph{Upper bound.} By \eqref{eq:a1-F1}, $P=R-S\le|S|-S+4\eta=2\max(-S,0)+4\eta\le2\kappa(4-\kappa)+4\eta\le8(\kappa+\eta)$, so by \eqref{eq:a1-F3} $\im m\le\sqrt{8\varkappa^2}/(2\sqrt2)=\varkappa$.

\emph{Lower bound on $P$.} We claim $P\ge\frac13\varkappa^2$.
If $S\le0$: by \eqref{eq:a1-F1}, $P=R+|S|\ge R\ge4\eta$ and $P\ge2|S|=2\kappa(4-\kappa)-2\eta^2\ge4\kappa-2\eta^2$. If moreover $\eta\le1$, then $P\ge4\kappa-2\eta$, and adding half of $P\ge4\eta$ gives $\frac32P\ge4\kappa$; so
$P\ge\max(\frac83\kappa,4\eta)\ge\frac12(\frac83\kappa+4\eta)\ge\frac43(\kappa+\eta)$. If $\eta>1$, then $\kappa\le2<2\eta$ and $P\ge4\eta>\frac43(\kappa+\eta)$.
If $S>0$: then $\eta^2>\kappa(4-\kappa)\ge2\kappa$, which implies $\kappa<\eta$ (if $\eta\le2$ then $\kappa<\eta^2/2\le\eta$; if $\eta>2$ then $\kappa\le2<\eta$). By \eqref{eq:a1-F1}, $Q\le2S+4\eta\le2\eta^2+4\eta$, so
$P=16\eta^2/Q\ge8\eta/(\eta+2)\ge\frac23\eta>\frac13(\kappa+\eta)$, using $\eta\le10$. This proves the claim.

\emph{Lower bound on $\im m$.} By the monotonicity of the left side of \eqref{eq:a1-F3} and the claim,
$\im m\ge\frac{\varkappa^2/3}{4(2\eta+\varkappa/\sqrt6)}$. Since $\eta\le\varkappa^2$ and $\eta\le10$ we have $\eta=\sqrt\eta\sqrt\eta\le\sqrt{10}\,\varkappa$, so $2\eta+\varkappa/\sqrt6\le(2\sqrt{10}+0.41)\varkappa\le6.74\,\varkappa$ and
\[
 \frac1{81}\,\varkappa\le\im m\le\varkappa\qquad(\text{Case C}).
\]

\emph{Conclusion.} If $|E|\le2$, the first statement of \eqref{eq:imm} is Case C. If $|E|>2$ and $\kappa<\eta$ it is Case B. If $|E|>2$ and $\kappa=\eta$, Case A applies and $\eta/\varkappa=\varkappa/2$ (as $\varkappa^2=2\eta$), so
$\frac1{54}\varkappa\le\im m\le\frac1{2\sqrt2}\varkappa$. The second statement of \eqref{eq:imm} is Case A. All the lower constants are $\ge\frac1{150}=c_*$ and all the upper constants are $\le4=C_*$.

For the last assertion: $\im m\ge\eta/(|z|+1)^2$ by Lemma~\ref{lem:msc}(c), and $(|z|+1)^2\le(\sqrt{125}+1)^2<150$, so $\im m\ge c_*\eta$. In Cases B and C, $\im m\le\varkappa$; in Case A, $\im m\le\eta/\varkappa\le\varkappa$ because $\eta\le\varkappa^2$.
Finally $\varkappa\le c_*^{-1}|1-m^2|$ by \eqref{eq:alpha} and $C_*\ge1$.
\end{proof}

\begin{lemma}[Stieltjes representation]\label{lem:stieltjes} \infra\
For $\im z>0$:\quad $\displaystyle\msc(z)=\int_{-2}^2\frac{\rhosc(x)}{x-z}\,\rd x$. Moreover $\displaystyle\int_{-2}^2\rhosc(x)\,\rd x=1$.
\end{lemma}
\begin{proof}
Let $m:=\msc(z)$. By Lemma~\ref{lem:msc}(a),(b): $0<|m|<1$ and $z=-(m+1/m)$. The integrand is continuous on $[-2,2]$ because $x-z\ne0$ for real $x$.
The function $\theta\mapsto2\cos\theta$ is continuously differentiable and maps $[0,\pi]$ onto $[-2,2]$, with $\rhosc(2\cos\theta)=\frac1{2\pi}\sqrt{4\sin^2\theta}=\frac1\pi\sin\theta$ for $\theta\in[0,\pi]$.
The substitution $x=2\cos\theta$, $\rd x=-2\sin\theta\,\rd\theta$, gives for every continuous $g:[-2,2]\to\C$
\begin{equation}\tag{$\Theta$}\label{eq:a1-Theta}
 \int_{-2}^2g(x)\rhosc(x)\,\rd x=\frac2\pi\int_0^\pi g(2\cos\theta)\sin^2\theta\,\rd\theta .
\end{equation}
For $g=1$ this is $\frac2\pi\int_0^\pi\sin^2\theta\,\rd\theta=\frac1\pi\int_0^\pi(1-\cos2\theta)\,\rd\theta=1$, which is the second claim. For $g(x)=1/(x-z)$ we have
\[
 2\cos\theta-z=2\cos\theta+m+\frac1m=\frac{1+2m\cos\theta+m^2}{m},\qquad 1+2m\cos\theta+m^2=(1+me^{\ii\theta})(1+me^{-\ii\theta})\ne0,
\]
so
\[
 \int_{-2}^2\frac{\rhosc(x)}{x-z}\,\rd x=\frac{2m}{\pi}\int_0^\pi\sin\theta\cdot\frac{\sin\theta}{1+2m\cos\theta+m^2}\,\rd\theta .
\]
\emph{Claim:} for every real $\theta$,
\[
 \frac{\sin\theta}{1+2m\cos\theta+m^2}=\sum_{k=0}^\infty(-m)^k\sin((k+1)\theta),
\]
the series being absolutely convergent with $|(-m)^k\sin((k+1)\theta)|\le|m|^k$. Indeed, by the geometric series ($|me^{\pm\ii\theta}|<1$) and $\sin((k+1)\theta)=\frac1{2\ii}(e^{\ii(k+1)\theta}-e^{-\ii(k+1)\theta})$,
\begin{align*}
 \sum_{k=0}^\infty(-m)^k\sin((k+1)\theta)&=\frac1{2\ii}\Big(\frac{e^{\ii\theta}}{1+me^{\ii\theta}}-\frac{e^{-\ii\theta}}{1+me^{-\ii\theta}}\Big)\\
 &=\frac1{2\ii}\cdot\frac{e^{\ii\theta}(1+me^{-\ii\theta})-e^{-\ii\theta}(1+me^{\ii\theta})}{(1+me^{\ii\theta})(1+me^{-\ii\theta})}=\frac{\sin\theta}{1+2m\cos\theta+m^2}.
\end{align*}
Since $\sum_k|m|^k<\infty$, the series may be integrated term by term against $\sin\theta$ on $[0,\pi]$ (dominated convergence). For integers $k\ge0$,
\[
 \int_0^\pi\sin\theta\,\sin((k+1)\theta)\,\rd\theta=\frac12\int_0^\pi\big(\cos(k\theta)-\cos((k+2)\theta)\big)\rd\theta=\begin{cases}\pi/2,&k=0,\\0,&k\ge1,\end{cases}
\]
because $\int_0^\pi\cos(j\theta)\,\rd\theta=0$ for every integer $j\ge1$. Therefore the integral equals $\frac{2m}\pi\cdot\frac\pi2=m$.
\end{proof}

\begin{lemma}[Semicircle distribution function and classical locations]\label{lem:nsc} \infra\
\begin{enumerate}[label=\textup{(\alph*)},leftmargin=2.2em,itemsep=1pt]
\item $\rhosc$ is continuous and even, $0\le\rhosc\le\frac1\pi$, $\rhosc>0$ exactly on $(-2,2)$, and $\rhosc$ is nonincreasing on $[0,\infty)$. The function $\nsc$ is continuously differentiable with $\nsc'=\rhosc$; it is
nondecreasing and $\frac1\pi$-Lipschitz, $\nsc(E)=0$ for $E\le-2$, $\nsc(E)=1$ for $E\ge2$, $\nsc$ is strictly increasing on $[-2,2]$, $\nsc(-E)=1-\nsc(E)$ for all $E$, and $\nsc(0)=\frac12$.
\item \emph{Edge:} for $0\le u\le2$:\quad $\dfrac{\sqrt u}{\sqrt2\,\pi}\le\rhosc(-2+u)\le\dfrac{\sqrt u}{\pi}$\quad and\quad $\dfrac{\sqrt2}{3\pi}u^{3/2}\le\nsc(-2+u)\le\dfrac{2}{3\pi}u^{3/2}$.
The two upper bounds hold for all $u\ge0$.
\item \emph{Classical locations:} for $1\le j\le N$ the infimum defining $\gamma_j$ in \eqref{eq:sc} is attained, $\gamma_j\in(-2,2]$, $N\nsc(\gamma_j)=j$, $\gamma_N=2$, $\gamma_j<\gamma_{j+1}$, and $\gamma_{N-j}=-\gamma_j$ for $1\le j\le N-1$. Moreover
\begin{align*}
 2.8\,(j/N)^{2/3}&\le\gamma_j+2\le3.6\,(j/N)^{2/3}&&\text{if }j\le N/2,\\
 2.8\,\big(\tfrac{N-j}N\big)^{2/3}&\le2-\gamma_j\le3.6\,\big(\tfrac{N-j}N\big)^{2/3}&&\text{if }j\ge N/2.
\end{align*}
\item \emph{Bulk}: for $0<c\le2$ and $|x|\le2-c/4$:\quad $\rhosc(x)\ge\frac1{2\pi}\sqrt{c-c^2/16}$.
\item \emph{Counting the $\gamma_j$}: for every bounded interval $I\subset\R$ (open, closed or half-open):\quad $\#\{j\in\{1,\dots,N\}:\ \gamma_j\in I\}\le N|I|/\pi+1$.
\item \emph{Density versus distribution function:} for all $x\in\R$:\quad $\rhosc(x)\ge\frac13\min\big(\nsc(x),1-\nsc(x)\big)^{1/3}$. Consequently, for all real $a\le b$,
\[
 (b-a)\,\min\big(\nsc(a),\,1-\nsc(b)\big)^{1/3}\le3\,\big(\nsc(b)-\nsc(a)\big).
\]
\end{enumerate}
\end{lemma}
\begin{proof}
(a) The statements on $\rhosc$ are read off from \eqref{eq:sc} ($\sqrt{(4-x^2)_+}\le2$). Since $\rhosc$ is continuous with compact support, $\nsc$ is well defined and continuously differentiable with $\nsc'=\rhosc$ (fundamental theorem of calculus);
it is nondecreasing because $\rhosc\ge0$, and $|\nsc(E)-\nsc(E')|=|\int_{E'}^E\rhosc|\le|E-E'|/\pi$. $\nsc(E)=0$ for $E\le-2$ because $\rhosc=0$ on $(-\infty,-2]$; $\nsc(E)=\int_{-2}^2\rhosc=1$ for $E\ge2$ by Lemma~\ref{lem:stieltjes}.
For $-2\le a<b\le2$, $\nsc(b)-\nsc(a)=\int_a^b\rhosc>0$ since $\rhosc$ is continuous and positive on $(a,b)$. Since $\rhosc$ is even, $\nsc(-E)=\int_{-\infty}^{-E}\rhosc=\int_E^\infty\rhosc=1-\nsc(E)$; for $E=0$ this gives $\nsc(0)=\frac12$.

(b) For $0\le u\le4$, $\rhosc(-2+u)=\frac1{2\pi}\sqrt{u(4-u)}$. If $0\le u\le2$ then $2\le4-u\le4$, which gives $\frac{\sqrt{2u}}{2\pi}\le\rhosc(-2+u)\le\frac{2\sqrt u}{2\pi}$. The upper bound $\rhosc(-2+u)\le\sqrt u/\pi$ holds for all $u\ge0$ (for $u\ge4$ the left side is $0$).
Integrating over $[0,u]$, $\nsc(-2+u)=\int_0^u\rhosc(-2+s)\,\rd s$ and $\int_0^u\sqrt s\,\rd s=\frac23u^{3/2}$ give the bounds on $\nsc$.

(c) Fix $j\in\{1,\dots,N\}$ and let $A_j:=\{x\in\R:N\nsc(x)\ge j\}$. By (a), $2\in A_j$, $A_j\subset(-2,\infty)$ (as $\nsc=0<j/N$ on $(-\infty,-2]$), $A_j$ is closed (continuity) and an up-set (monotonicity). So $\gamma_j=\min A_j\in(-2,2]$ and $\nsc(\gamma_j)\ge j/N$.
For $x<\gamma_j$ we have $\nsc(x)<j/N$; letting $x\uparrow\gamma_j$, continuity gives $\nsc(\gamma_j)\le j/N$. Hence $N\nsc(\gamma_j)=j$.
Since $\nsc(x)<1$ for $x<2$ (strict monotonicity on $[-2,2]$) and $\nsc(2)=1$, $\gamma_N=2$. Since $\nsc$ is nondecreasing and $\nsc(\gamma_j)=j/N<(j+1)/N=\nsc(\gamma_{j+1})$, we get $\gamma_j<\gamma_{j+1}$.
For $1\le j\le N-1$: $\gamma_j\in(-2,2)$ and $\nsc(-\gamma_j)=1-j/N=\nsc(\gamma_{N-j})$; as $\nsc$ is injective on $[-2,2]$, $-\gamma_j=\gamma_{N-j}$.

Let $j\le N/2$. Then $\nsc(0)=\frac12\ge j/N$, so $\gamma_j\le0$ and $u:=\gamma_j+2\in(0,2]$. By (b), $\frac{\sqrt2}{3\pi}u^{3/2}\le\frac jN\le\frac2{3\pi}u^{3/2}$, i.e.
\[
 \Big(\frac{3\pi}2\Big)^{2/3}\Big(\frac jN\Big)^{2/3}\le u\le\Big(\frac{3\pi}{\sqrt2}\Big)^{2/3}\Big(\frac jN\Big)^{2/3},\qquad\Big(\frac{3\pi}2\Big)^{2/3}>2.8,\quad\Big(\frac{3\pi}{\sqrt2}\Big)^{2/3}<3.6 .
\]
Let $j\ge N/2$. If $j=N$ both sides of the claim vanish. If $j<N$, then $N-j\le N/2$ and $2-\gamma_j=2+\gamma_{N-j}$, so the claim follows from the case just proved.

(d) If $|x|\le2-c/4$ then $4-x^2\ge4-(2-c/4)^2=c-c^2/16>0$.

(e) If no $\gamma_j$ lies in $I$ there is nothing to prove. Otherwise let $j_1\le j_2$ be the smallest and the largest index with $\gamma_j\in I$. By (c) and the Lipschitz bound in (a),
$\frac{j_2-j_1}N=\nsc(\gamma_{j_2})-\nsc(\gamma_{j_1})\le\frac{\gamma_{j_2}-\gamma_{j_1}}\pi\le\frac{|I|}\pi$, and the number of indices is at most $j_2-j_1+1$.

(f) Let $\tau(x):=\min(\nsc(x),1-\nsc(x))$. If $|x|\ge2$ then $\tau(x)=0$ and the inequality is trivial. Let $-2<x\le0$ and $u:=x+2\in(0,2]$. Then $\nsc(x)\le\frac12$, so $\tau(x)=\nsc(x)\le\frac2{3\pi}u^{3/2}$ by (b), i.e.\ $\sqrt u\ge(\frac{3\pi}2\tau(x))^{1/3}$, and by (b) again
\[
 \rhosc(x)\ge\frac{\sqrt u}{\sqrt2\,\pi}\ge\frac{(3\pi/2)^{1/3}}{\sqrt2\,\pi}\,\tau(x)^{1/3}\ge\tfrac13\,\tau(x)^{1/3},
\]
since $(3\pi/2)^{1/3}/(\sqrt2\pi)>0.37$. For $0<x<2$ use $\rhosc(x)=\rhosc(-x)$ and $\tau(x)=\tau(-x)$.
For the consequence, let $a\le b$ and $\tau_0:=\min(\nsc(a),1-\nsc(b))$. For $x\in[a,b]$, $\nsc(a)\le\nsc(x)\le\nsc(b)$, hence $\tau(x)\ge\tau_0$ and $\rhosc(x)\ge\frac13\tau_0^{1/3}$. Integrating over $[a,b]$:
$\nsc(b)-\nsc(a)\ge\frac13(b-a)\tau_0^{1/3}$.
\end{proof}

\subsection{Resolvent identities}\label{sec:resid}

\paragraph{Notation.} Vectors $x\in\C^N$ are columns, $\langle x,y\rangle:=\sum_i\ol{x_i}y_i$, $|x|:=\langle x,x\rangle^{1/2}$, $\norm x_\infty:=\max_i|x_i|$, and
$\norm M:=\sup_{|x|=1}|Mx|$ is the operator norm of a matrix $M$; thus $|M_{ij}|=|\langle e_i,Me_j\rangle|\le\norm M$ for the standard basis $(e_i)$.
$M^*$ is the conjugate transpose. If $M$ is indexed by a finite set $I$ and $P,Q\subset I$, then $M_{PQ}:=(M_{pq})_{p\in P,q\in Q}$.
A matrix indexed by the empty set is the \emph{empty matrix}; by convention it is invertible, equal to its inverse, its determinant is $1$, its trace is $0$, and sums over an empty index set are $0$.

For $T\subset\{1,\dots,N\}$ let $H^{(T)}:=(h_{ij})_{i,j\notin T}$ be the minor of $H$ with rows and columns in $T$ removed. It is indexed by $T^c:=\{1,\dots,N\}\setminus T$
(indices are \emph{not} relabelled) and it is Hermitian. For $\im z>0$,
\[
 G^{(T)}(z):=(H^{(T)}-z)^{-1},\qquad m^{(T)}(z):=\frac1N\sum_{k\notin T}G^{(T)}_{kk}(z)\quad(\text{normalised by }N,\text{ not by }N-|T|),
\]
and we write $(iT)$ for $\{i\}\cup T$, $(ij)$ for $\{i,j\}$, and $\sum^{(T)}_k$ for the sum over $k\in T^c$. $H^{(\emptyset)}=H$, $G^{(\emptyset)}=G$.
The matrix $G^{(T)}$ exists by Lemma~\ref{lem:basic}(a) applied to $H^{(T)}$. Throughout, $z=E+\ii\eta$ and $z'=E'+\ii\eta'$ with $\eta,\eta'>0$.

\begin{lemma}[Basic properties]\label{lem:basic}
Let $H$ be a Hermitian $N\times N$ matrix, $N\ge1$, and $\eta>0$. Then:
\begin{enumerate}[label=\textup{(\alph*)},leftmargin=2.2em,itemsep=1pt]
\item $|(H-z)x|\ge\eta|x|$ for all $x\in\C^N$; $H-z$ is invertible, $\norm{G}\le\eta^{-1}$, $|G_{ij}|\le\eta^{-1}$, $G(z)^*=(H-\bar z)^{-1}$, and $\im G_{ii}>0$ (in particular $G_{ii}\ne0$).
If $H$ is real symmetric then $G_{ij}=G_{ji}$.
\item \emph{Ward identity:} $\sum_j|G_{ij}|^2=\sum_j|G_{ji}|^2=\im G_{ii}/\eta$.
\item \emph{Resolvent identity:} $G(z)-G(z')=(z-z')G(z)G(z')$; hence $|G_{ij}(z)-G_{ij}(z')|\le|z-z'|/(\eta\eta')$ and $|m(z)-m(z')|\le|z-z'|/(\eta\eta')$.
Consequently, if $\eta,\eta'\ge\eta_0>0$, each of the three functions $F\in\{\Lambda,\Ld,\Lo\}$ of \eqref{eq:defLambda} satisfies $|F(z)-F(z')|\le\big(\eta_0^{-2}+\eta_0^{-1}+1\big)|z-z'|$.
\item \emph{Spectral form:} let $\lambda_1\le\dots\le\lambda_N$ be the eigenvalues of $H$ and $u_1,\dots,u_N$ an orthonormal basis of $\C^N$ with $Hu_\alpha=\lambda_\alpha u_\alpha$. Then
\[
 G_{ij}(z)=\sum_{\alpha=1}^N\frac{u_\alpha(i)\ol{u_\alpha(j)}}{\lambda_\alpha-z},\qquad\tr G(z)=\sum_{\alpha=1}^N\frac1{\lambda_\alpha-z},\qquad
 \im m(z)=\frac1N\sum_{\alpha=1}^N\frac{\eta}{(\lambda_\alpha-E)^2+\eta^2},
\]
and $\im G_{ii}(z)=\sum_\alpha|u_\alpha(i)|^2\,\eta/((\lambda_\alpha-E)^2+\eta^2)$.
\item \emph{Counting:} $\#\{\alpha:|\lambda_\alpha-E|\le\eta\}\le2\eta\,\im\tr G(E+\ii\eta)=2N\eta\,\im m(E+\ii\eta)$. In particular, if $\im m(E+\ii\eta)<\frac1{2N\eta}$ then $H$ has no eigenvalue in $[E-\eta,E+\eta]$.
\item For fixed $E$, the functions $\eta\mapsto\eta\,\im m(E+\ii\eta)$ and $\eta\mapsto\eta\,\im G_{ii}(E+\ii\eta)$ are nondecreasing on $(0,\infty)$, and $\eta\mapsto\eta^{-1}\im m(E+\ii\eta)$, $\eta\mapsto\eta^{-1}\im G_{ii}(E+\ii\eta)$ are nonincreasing.
\item $\norm{G(z)}=\max_\alpha|\lambda_\alpha-z|^{-1}=\operatorname{dist}(z,\operatorname{Spec}H)^{-1}$ and $G(z)+z^{-1}=z^{-1}G(z)H$.
\item \emph{Two matrices}: if $A,B$ are Hermitian $N\times N$ matrices and $G_A:=(A-z)^{-1}$, $G_B:=(B-z)^{-1}$, then
\[
 G_A-G_B=G_A\,(B-A)\,G_B,\quad\ \norm{G_A-G_B}\le\eta^{-2}\norm{A-B},\quad\ |\tr G_A-\tr G_B|\le N\eta^{-2}\norm{A-B}.
\]
\end{enumerate}
Items \textup{(a)--(c), (g), (h)} hold verbatim for every minor $H^{(T)}$, $T\ne\{1,\dots,N\}$; items \textup{(d)--(f)} hold for $H^{(T)}$ with $\tr G^{(T)}=Nm^{(T)}$ in place of $\tr G=Nm$
and the $N-|T|$ eigenvalues and eigenvectors of $H^{(T)}$.
\end{lemma}
\begin{proof}
(a) Since $H=H^*$, $\langle x,Hx\rangle$ is real, so $\im\langle x,(H-z)x\rangle=-\eta|x|^2$ and, by the Cauchy--Schwarz inequality,
\[
 |x|\,|(H-z)x|\ \ge\ |\langle x,(H-z)x\rangle|\ \ge\ \eta|x|^2 .
\]
Hence $|(H-z)x|\ge\eta|x|$; the linear map $H-z$ is injective, therefore bijective on $\C^N$, and putting $x=Gy$ gives $|Gy|\le\eta^{-1}|y|$, i.e.\ $\norm G\le\eta^{-1}$ and $|G_{ij}|\le\norm G\le\eta^{-1}$.
The same argument with $-\eta$ shows that $H-\bar z$ is invertible, and $G^*=((H-z)^*)^{-1}=(H-\bar z)^{-1}$. If $H$ is real symmetric, $H-z$ is a symmetric matrix and so is its inverse.
The claim $\im G_{ii}>0$ is proved with (b).

(b) From $(H-\bar z)-(H-z)=(z-\bar z)=2\ii\eta$ we get, multiplying by $G$ on the left and $G^*$ on the right, respectively by $G^*$ on the left and $G$ on the right,
\[
 G-G^*=2\ii\eta\,GG^*,\qquad G-G^*=2\ii\eta\,G^*G .
\]
The $(i,i)$ entry of $G-G^*$ is $G_{ii}-\ol{G_{ii}}=2\ii\im G_{ii}$, that of $GG^*$ is $\sum_j|G_{ij}|^2$ and that of $G^*G$ is $\sum_j|G_{ji}|^2$. This is (b).
Since $G$ is invertible, its $i$-th row is not zero, so $\im G_{ii}=\eta\sum_j|G_{ij}|^2>0$, which completes (a).

(c) $G(z)\big[(H-z')-(H-z)\big]G(z')=G(z)-G(z')$ and the bracket equals $(z-z')$. By (a), $|G_{ij}(z)-G_{ij}(z')|\le|z-z'|\,\norm{G(z)}\norm{G(z')}\le|z-z'|/(\eta\eta')$, and $m$ is the average of the $G_{ii}$.
For the last claim, $|\msc(z)-\msc(z')|\le(1+\eta_0^{-1})|z-z'|$ by Lemma~\ref{lem:msc}(f); so each of the functions $G_{ii}-\msc$, $m-\msc$, $G_{ij}$ $(i\ne j)$ changes by at most
$(\eta_0^{-2}+\eta_0^{-1}+1)|z-z'|$ between $z$ and $z'$, and $\big||a|-|b|\big|\le|a-b|$, $|\max_ka_k-\max_kb_k|\le\max_k|a_k-b_k|$.

(d) By the spectral theorem for Hermitian matrices there is an orthonormal eigenbasis $(u_\alpha)$; let $U$ be the unitary matrix with columns $u_\alpha$ and $D:=\operatorname{diag}(\lambda_\alpha)$, so $H=UDU^*$.
Since $\im(\lambda_\alpha-z)=-\eta\ne0$, the matrix $D-z$ is invertible and $(H-z)\,U(D-z)^{-1}U^*=U(D-z)(D-z)^{-1}U^*=1$. Hence $G=U(D-z)^{-1}U^*$, which is the formula for $G_{ij}$;
$\tr G=\tr\big((D-z)^{-1}U^*U\big)=\sum_\alpha(\lambda_\alpha-z)^{-1}$; and $\im\frac1{\lambda-z}=\im\frac{\lambda-E+\ii\eta}{|\lambda-z|^2}=\frac{\eta}{(\lambda-E)^2+\eta^2}$.

(e) Every $\alpha$ with $|\lambda_\alpha-E|\le\eta$ contributes at least $\eta/(2\eta^2)=1/(2\eta)$ to the sum $\im\tr G=\sum_\alpha\eta/((\lambda_\alpha-E)^2+\eta^2)$, and all terms are positive. Hence
$\#\{\alpha:|\lambda_\alpha-E|\le\eta\}\le2\eta\,\im\tr G$. If $\im m<1/(2N\eta)$, the right side is $<1$, so the set is empty.

(f) By (d), $\eta\,\im G_{ii}=\sum_\alpha|u_\alpha(i)|^2\frac{\eta^2}{(\lambda_\alpha-E)^2+\eta^2}$ and $\eta^{-1}\im G_{ii}=\sum_\alpha|u_\alpha(i)|^2\frac{1}{(\lambda_\alpha-E)^2+\eta^2}$; for $a\ge0$ the function $\eta\mapsto\eta^2/(a+\eta^2)=1-a/(a+\eta^2)$ is nondecreasing and
$\eta\mapsto1/(a+\eta^2)$ is nonincreasing. Average over $i$ for $m$.

(g) By (d), $G=U(D-z)^{-1}U^*$ with $U$ unitary, so $\norm G=\norm{(D-z)^{-1}}=\max_\alpha|\lambda_\alpha-z|^{-1}$, and $\min_\alpha|\lambda_\alpha-z|$ is the distance from $z$ to the spectrum.
Further $GH=G(H-z)+zG=1+zG$; divide by $z\ne0$.

(h) $G_A\big[(B-z)-(A-z)\big]G_B=G_A-G_B$. Then use $\norm{G_A},\norm{G_B}\le\eta^{-1}$ from (a), and $|\tr M|\le N\max_i|M_{ii}|\le N\norm M$.

The statement on minors needs no proof: $H^{(T)}$ is a Hermitian matrix of size $N-|T|\ge1$.
\end{proof}

\begin{lemma}[Schur complement in block form]\label{lem:schur} \infra\
Let $H$ be a Hermitian $N\times N$ matrix, $\eta>0$, $\emptyset\ne A\subset\{1,\dots,N\}$, and define the $A\times A$ matrix
\begin{equation}\label{eq:MA}
 M^A_{ab}:=h_{ab}-z\delta_{ab}-\sum_{k,l\notin A}h_{ak}G^{(A)}_{kl}h_{lb}\qquad(a,b\in A).
\end{equation}
Then $M^A$ is invertible and $(G_{ab})_{a,b\in A}=(M^A)^{-1}$. Moreover, for $a\in A$ and $k\notin A$,
\[
 G_{ak}=-\sum_{b\in A}\sum_{l\notin A}G_{ab}\,h_{bl}\,G^{(A)}_{lk},\qquad G_{ka}=-\sum_{b\in A}\sum_{l\notin A}G^{(A)}_{kl}\,h_{lb}\,G_{ba}.
\]
Consequently:
\begin{enumerate}[label=\textup{(I\arabic*)},leftmargin=2.6em,itemsep=1pt]
\item $\dfrac1{G_{ii}}=h_{ii}-z-\sum^{(i)}_{k,l}h_{ik}G^{(i)}_{kl}h_{li}$.
\item For $i\ne j$:
\[
 G_{ij}=-G_{ii}\sum^{(i)}_kh_{ik}G^{(i)}_{kj}=-G_{ii}G^{(i)}_{jj}\Big(h_{ij}-\sum^{(ij)}_{k,l}h_{ik}G^{(ij)}_{kl}h_{lj}\Big),
 \qquad G_{ji}=-G_{ii}\sum^{(i)}_kG^{(i)}_{jk}h_{ki}.
\]
\item For $A\ne\{1,\dots,N\}$ the matrix $G_{AA}=(G_{cd})_{c,d\in A}$ is invertible and for $i,j\notin A$:
$G^{(A)}_{ij}=G_{ij}-\sum_{a,b\in A}G_{ia}\big(G_{AA}^{-1}\big)_{ab}G_{bj}$; for $A=\{k\}$: $G^{(k)}_{ij}=G_{ij}-G_{ik}G_{kj}/G_{kk}$.
\item For $k\in A$, $A':=A\setminus\{k\}\ne\emptyset$: the matrix $M^A_{A'A'}:=(M^A_{ab})_{a,b\in A'}$ is invertible, $(M^A_{A'A'})^{-1}=(G^{(k)}_{ab})_{a,b\in A'}$, and
\begin{equation}\label{eq:1overG}
 \frac1{G_{kk}}=M^A_{kk}-\sum_{a,b\in A'}M^A_{ka}\,G^{(k)}_{ab}\,M^A_{bk}.
\end{equation}
Moreover $M^A_{aa}=1/G^{(A\setminus a)}_{aa}$ for every $a\in A$.
\item $\tr G-\tr G^{(i)}=(G^2)_{ii}/G_{ii}$ and $|(G^2)_{ii}|\le\im G_{ii}/\eta\le|G_{ii}|/\eta$; hence $|m-m^{(i)}|\le\frac1{N\eta}$ and, for every $T\subset\{1,\dots,N\}$,
\[
 |m^{(T)}-m|\le\frac{|T|}{N\eta},\qquad\text{in particular}\quad\im m^{(T)}\le\im m+\frac{|T|}{N\eta}.
\]
\item \emph{Determinants}: $\det(H-z)=\det(H^{(A)}-z)\cdot\det M^A=\det(H^{(A)}-z)/\det(G_{AA})$; for $A=\{a\}$:
\[
 \det(H-z)=\frac{\det(H^{(a)}-z)}{G_{aa}} .
\]
\end{enumerate}
The same identities hold for every minor $H^{(T)}$ in place of $H$ (with $A\subset T^c$, and with $G^{(T)}$, $G^{(T\cup A)}$, $m^{(T)}$, $\tr G^{(T)}$ in place of $G$, $G^{(A)}$, $m$, $\tr G$);
for instance $1/G^{(T)}_{ii}=h_{ii}-z-\sum^{(iT)}_{k,l}h_{ik}G^{(iT)}_{kl}h_{li}$ for $i\notin T$.
\end{lemma}
\begin{proof}
\emph{Step 0: the block inverse formula.} Let $I=A\sqcup A^c$ be a finite index set with $A,A^c\ne\emptyset$, and let $\mathcal M$ be an invertible $I\times I$ complex matrix, written in blocks as
\[
 \mathcal M=\begin{pmatrix}P&Q\\R&S\end{pmatrix},\qquad\mathcal M^{-1}=\begin{pmatrix}X&Y\\Z&W\end{pmatrix}\qquad(P,X:\ A\times A;\quad S,W:\ A^c\times A^c),
\]
and assume that $S$ is invertible. Then $P-QS^{-1}R$ is invertible and
\begin{equation}\tag{BI}\label{eq:BI}
 X=(P-QS^{-1}R)^{-1},\qquad Y=-XQS^{-1},\qquad Z=-S^{-1}RX,\qquad W=S^{-1}+S^{-1}RXQS^{-1}.
\end{equation}
Indeed, $\mathcal M\mathcal M^{-1}=1$ gives $PX+QZ=1$ and $RX+SZ=0$; the second relation gives $Z=-S^{-1}RX$ and then the first gives $(P-QS^{-1}R)X=1$.
For square matrices a one-sided inverse is the inverse, which proves the first and third formulas. $\mathcal M^{-1}\mathcal M=1$ gives $XQ+YS=0$ and $ZQ+WS=1$, hence $Y=-XQS^{-1}$ and
$W=S^{-1}-ZQS^{-1}=S^{-1}+S^{-1}RXQS^{-1}$. (In Mathlib: \texttt{Matrix.invOf\_fromBlocks\textsubscript{22}\_eq}.)
Moreover, without assuming that $\mathcal M$ is invertible,
\begin{equation}\tag{BD}\label{eq:BD}
 \det\mathcal M=\det S\cdot\det(P-QS^{-1}R),
\end{equation}
because
$\begin{pmatrix}P&Q\\R&S\end{pmatrix}=\begin{pmatrix}1&QS^{-1}\\0&1\end{pmatrix}\begin{pmatrix}P-QS^{-1}R&0\\0&S\end{pmatrix}\begin{pmatrix}1&0\\S^{-1}R&1\end{pmatrix}$
and the determinant of a block triangular matrix is the product of the determinants of its diagonal blocks (Mathlib: \texttt{Matrix.det\_fromBlocks\textsubscript{22}}).
Both the inverse and the determinant are unchanged under a simultaneous relabelling of rows and columns, so it does not matter how $I$ is ordered.

\emph{Step 1: $G_{AA}=(M^A)^{-1}$ and the off-diagonal blocks.} If $A=\{1,\dots,N\}$, then $M^A=H-z$ by definition (empty sum) and there is nothing to prove.
Otherwise apply Step~0 to $\mathcal M:=H-z$, which is invertible by Lemma~\ref{lem:basic}(a), with
$P=H_{AA}-z$, $Q=H_{AA^c}$, $R=H_{A^cA}$ and $S=H^{(A)}-z$. The matrix $S$ is invertible by Lemma~\ref{lem:basic}(a) applied to the Hermitian matrix $H^{(A)}$, and $S^{-1}=G^{(A)}$.
By definition \eqref{eq:MA}, $P-QS^{-1}R=M^A$. So \eqref{eq:BI} gives: $M^A$ is invertible, $G_{AA}=X=(M^A)^{-1}$, $G_{AA^c}=Y=-G_{AA}H_{AA^c}G^{(A)}$ and $G_{A^cA}=Z=-G^{(A)}H_{A^cA}G_{AA}$.
The last two are the displayed formulas for $G_{ak}$, $G_{ka}$.

\emph{Step 2: \textup{(I1)} and the first and third formulas of \textup{(I2)}.} Take $A=\{i\}$. Then $M^A$ is the $1\times1$ matrix with entry $h_{ii}-z-\sum^{(i)}_{k,l}h_{ik}G^{(i)}_{kl}h_{li}$, and Step~1 says that it is the inverse of $G_{ii}$: this is (I1).
The off-diagonal blocks of Step~1 read, for $j\ne i$: $G_{ij}=-G_{ii}\sum^{(i)}_lh_{il}G^{(i)}_{lj}$ and $G_{ji}=-\sum^{(i)}_lG^{(i)}_{jl}h_{li}G_{ii}$.

\emph{Step 3: a consistency property of $M^A$.} Let $\emptyset\ne A'\subset A$ and $T:=A\setminus A'$. Then
\begin{equation}\tag{$\ast$}\label{eq:MAcons}
 M^A_{ab}=M^{A'}\big[H^{(T)}\big]_{ab}\qquad(a,b\in A'),
\end{equation}
where $M^{A'}[H^{(T)}]$ denotes the matrix \eqref{eq:MA} formed with the Hermitian matrix $H^{(T)}$ (indexed by $T^c$) and its subset of indices $A'\subset T^c$.
Indeed, the complement of $A'$ in $T^c$ is $A^c$, and the minor of $H^{(T)}$ with the rows and columns in $A'$ removed is $H^{(T\cup A')}=H^{(A)}$; so
$M^{A'}[H^{(T)}]_{ab}=h_{ab}-z\delta_{ab}-\sum_{k,l\notin A}h_{ak}G^{(A)}_{kl}h_{lb}$, which is $M^A_{ab}$.

\emph{Step 4: the second formula of \textup{(I2)}.} Let $i\ne j$, $A:=\{i,j\}$ and $M:=M^A$. By Step~1, $M$ is an invertible $2\times2$ matrix with inverse $G_{AA}$, so by the formula for the inverse of a $2\times2$ matrix
\[
 G_{ii}=\frac{M_{jj}}{\det M},\qquad G_{ij}=-\frac{M_{ij}}{\det M}.
\]
By \eqref{eq:MAcons} with $A'=\{j\}$, $T=\{i\}$, and by (I1) applied to the Hermitian matrix $H^{(i)}$ and its index $j$: $M_{jj}=M^{\{j\}}[H^{(i)}]_{jj}=1/G^{(i)}_{jj}$; in particular $M_{jj}\ne0$.
Hence $\det M=M_{jj}/G_{ii}$ and $G_{ij}=-G_{ii}M_{ij}/M_{jj}=-G_{ii}G^{(i)}_{jj}M_{ij}$, where $M_{ij}=h_{ij}-\sum^{(ij)}_{k,l}h_{ik}G^{(ij)}_{kl}h_{lj}$ by \eqref{eq:MA}.

\emph{Step 5: \textup{(I3)}.} Let $A\ne\{1,\dots,N\}$. Apply Step~0 to the invertible matrix $\mathcal M:=G$, whose inverse is $H-z$, with the roles of the two blocks exchanged:
the block $G_{AA}$ plays the role of $S$ (it is invertible by Step~1), $G_{A^cA^c}$ the role of $P$. The first formula of \eqref{eq:BI} says that $G_{A^cA^c}-G_{A^cA}(G_{AA})^{-1}G_{AA^c}$ is invertible and
\[
 \big(G_{A^cA^c}-G_{A^cA}(G_{AA})^{-1}G_{AA^c}\big)^{-1}=(G^{-1})_{A^cA^c}=(H-z)_{A^cA^c}=H^{(A)}-z .
\]
Taking inverses, $G^{(A)}=G_{A^cA^c}-G_{A^cA}(G_{AA})^{-1}G_{AA^c}$, which is (I3). For $A=\{k\}$, $G_{AA}^{-1}=1/G_{kk}$.

\emph{Step 6: \textup{(I4)}.} Let $k\in A$ and $A'=A\setminus\{k\}\ne\emptyset$. By \eqref{eq:MAcons} with $T=\{k\}$, $M^A_{A'A'}=M^{A'}[H^{(k)}]$, and by Step~1 applied to the Hermitian matrix $H^{(k)}$ and $A'$ this matrix
is invertible with inverse $(G^{(k)}_{ab})_{a,b\in A'}$. Now apply Step~0 to $\mathcal M:=M^A$ (invertible, with inverse $G_{AA}$, by Step~1) with the blocks $\{k\}$ and $A'$: $S=M^A_{A'A'}$ is invertible, so
\[
 G_{kk}=\big((M^A)^{-1}\big)_{kk}=\Big(M^A_{kk}-\sum_{a,b\in A'}M^A_{ka}\big((M^A_{A'A'})^{-1}\big)_{ab}M^A_{bk}\Big)^{-1},
\]
which is \eqref{eq:1overG}. Finally, for $a\in A$, \eqref{eq:MAcons} with $A'=\{a\}$, $T=A\setminus\{a\}$ and (I1) for $H^{(T)}$ give $M^A_{aa}=1/G^{(A\setminus a)}_{aa}$ (for $A=\{a\}$ this is (I1) itself).

\emph{Step 7: \textup{(I5)}.} \emph{The identity.} If $N=1$ then $\tr G^{(i)}=0$ and $(G^2)_{ii}/G_{ii}=G_{ii}=\tr G$. If $N\ge2$, then by (I3) with $A=\{i\}$, $G^{(i)}_{kk}=G_{kk}-G_{ki}G_{ik}/G_{ii}$ for $k\ne i$; summing over $k\ne i$,
\[
 \tr G^{(i)}=\tr G-G_{ii}-\frac1{G_{ii}}\sum_{k\ne i}G_{ik}G_{ki}=\tr G-\frac1{G_{ii}}\sum_{k}G_{ik}G_{ki}=\tr G-\frac{(G^2)_{ii}}{G_{ii}} .
\]
\emph{The bounds} (for every $N\ge1$). By the Cauchy--Schwarz inequality and the Ward identity (Lemma~\ref{lem:basic}(b)),
$|(G^2)_{ii}|\le\big(\sum_k|G_{ik}|^2\big)^{1/2}\big(\sum_k|G_{ki}|^2\big)^{1/2}=\im G_{ii}/\eta\le|G_{ii}|/\eta$. Hence $|\tr G-\tr G^{(i)}|\le\eta^{-1}$ and $|m-m^{(i)}|\le(N\eta)^{-1}$.
For a general $T=\{t_1,\dots,t_r\}$ (with $r=|T|$ distinct elements; $T=\{1,\dots,N\}$ is allowed) put $T_s:=\{t_1,\dots,t_s\}$, $T_0:=\emptyset$. For $1\le s\le r$ the matrix $H^{(T_{s-1})}$ is Hermitian (of size $N-s+1\ge1$) and $H^{(T_s)}$ is its minor with the index $t_s$ removed, so the bound just proved gives
$|\tr G^{(T_{s-1})}-\tr G^{(T_s)}|\le\eta^{-1}$. Summing over $s$: $|\tr G-\tr G^{(T)}|\le r/\eta$, i.e.\ $|m-m^{(T)}|\le|T|/(N\eta)$.

\emph{Step 8: \textup{(I6)}.} If $A=\{1,\dots,N\}$ the claim is $\det(H-z)=\det M^A$, which is true as $M^A=H-z$. Otherwise \eqref{eq:BD} with the blocks of Step~1 gives
$\det(H-z)=\det(H^{(A)}-z)\det M^A$, and $\det M^A=1/\det G_{AA}$ by Step~1. For $A=\{a\}$, $\det G_{AA}=G_{aa}$.

\emph{Minors.} For $T\ne\{1,\dots,N\}$ the matrix $H^{(T)}$ is Hermitian, and for $A\subset T^c$ its minor with the indices $A$ removed is $H^{(T\cup A)}$. So every statement applies to $H^{(T)}$ as described.
\end{proof}

\begin{remark}\label{rem:rows}
Two structural facts used in Sections~\ref{sec:prob} and~\ref{sec:FA} can be read off \eqref{eq:MA}: (i) $M^A_{ab}$ is a polynomial in the entries $h_{ab}$, $(h_{ak})_{k\notin A}$, $(h_{lb})_{l\notin A}$ with coefficients $-z\delta_{ab}$ and $G^{(A)}_{kl}$
that depend only on the minor $H^{(A)}$; (ii) by (I4), $M^A_{aa}=1/G^{(A\setminus a)}_{aa}$, so $|1/M^A_{aa}|\le\eta^{-1}$ by Lemma~\ref{lem:basic}(a).
\end{remark}

\begin{lemma}[Perturbation of minors]\label{lem:minor}
Let $A\subset\{1,\dots,N\}$, $|A|=r\ge1$, and $g>0$. If $|G_{aa}|\ge g$ for all $a\in A$ and $\max_{a\ne b\in A}|G_{ab}|\le g/(2r)$, then every entry of $(G_{AA})^{-1}$ has modulus at most $2/g$, and for all $i,j\notin A$:
\[
 |G^{(A)}_{ij}-G_{ij}|\le\frac{2r^2}g\max_{a\in A}|G_{ia}|\max_{b\in A}|G_{bj}|.
\]
\end{lemma}
\begin{proof}
By Lemma~\ref{lem:schur}, $G_{AA}$ is invertible; let $X:=(G_{AA})^{-1}$ and $\mu:=\max_{a,b\in A}|X_{ab}|$. Writing the identity $G_{AA}X=1$ entrywise and isolating the diagonal entry of $G_{AA}$,
\[
 G_{aa}X_{ab}=\delta_{ab}-\sum_{c\in A,\,c\ne a}G_{ac}X_{cb}\qquad(a,b\in A).
\]
Hence $g|X_{ab}|\le1+(r-1)\frac{g}{2r}\mu\le1+\frac g2\mu$ for all $a,b$, so $g\mu\le1+\frac g2\mu$ and $\mu\le2/g$.
By (I3), $|G^{(A)}_{ij}-G_{ij}|\le\sum_{a,b\in A}|G_{ia}|\,|X_{ab}|\,|G_{bj}|\le r^2\frac2g\max_a|G_{ia}|\max_b|G_{bj}|$.
\end{proof}

\subsection{The self-consistent equation and its stability}\label{sec:stab}
In this subsection $B=(\sigma_{ij}^2)_{i,j=1}^N$ is a real symmetric matrix with nonnegative entries satisfying (A), (B), (C) of Definition~\ref{def:ass}; $H$ is an arbitrary Hermitian $N\times N$ matrix
(no relation between $H$ and $B$ is assumed), $N\ge2$, $z=E+\ii\eta$ with $\eta>0$, and $m:=\msc(z)$. By (A) and the symmetry of $B$,
\begin{equation}\tag{A$'$}\label{eq:Arow}
 \sum_{j=1}^N\sigma_{ij}^2=1\quad\text{for every }i,\qquad\text{i.e.}\quad B\bone=\bone,\qquad\bone:=(1,\dots,1)^{T}.
\end{equation}
\begin{definition}\label{def:XA}
For $i\in\{1,\dots,N\}$ (all quantities depend on $z$):
\begin{align*}
 Z_i&:=\sum^{(i)}_{k,l}h_{ik}G^{(i)}_{kl}h_{li}-\sum^{(i)}_k\sigma_{ik}^2G^{(i)}_{kk}, & X_i&:=h_{ii}-Z_i,\\
 A_i&:=\sigma_{ii}^2G_{ii}+\sum_{j\ne i}\sigma_{ij}^2\frac{G_{ij}G_{ji}}{G_{ii}}, & \Upsilon_i&:=A_i+X_i .
\end{align*}
For a vector $u\in\C^N$, $\avg u:=\frac1N\sum_iu_i$; for vectors $u,u'$ the products $uu'$, $u^2$, \dots\ are taken entrywise. Further $\Upsilon_{\max}:=\max_i|\Upsilon_i|$.
\end{definition}
\noindent($A_i$ is well defined since $G_{ii}\ne0$, Lemma~\ref{lem:basic}(a). In the probabilistic setting $X_i=Q_i(1/G_{ii})$, where $Q_i=1-P_i$ and $P_i$ is the expectation in the $i$-th row, see Section~\ref{sec:FA}; this is not used in the present section.)

\begin{lemma}[Exact self-consistent equations]\label{lem:sce}
Assume \textup{(A)}. With $v_i=G_{ii}-\msc$, $m=\msc$, $r_i:=v_i^2/G_{ii}$:
\begin{align}
 \frac1{G_{ii}}&=-z-\sum_j\sigma_{ij}^2G_{jj}+\Upsilon_i, &
 v_i&=m\,G_{ii}\Big(\sum_j\sigma_{ij}^2v_j-\Upsilon_i\Big),\label{eq:sce1}\\
 v_i-m^2\sum_j\sigma^2_{ij}v_j&=-m^2\Upsilon_i+r_i, &
 (1-m^2)\avg v&=-m^2\avg\Upsilon+\avg r,\label{eq:sce2}
\end{align}
and $A_i=\sigma_{ii}^2G_{ii}+G_{ii}\sum_{j\ne i}\sigma_{ij}^2\big(\sum^{(i)}_kh_{ik}G^{(i)}_{kj}\big)\big(\sum^{(i)}_lG^{(i)}_{jl}h_{li}\big)$.
\end{lemma}
\begin{proof}
\emph{First equation of \eqref{eq:sce1}.} By (I1) and the definitions of $Z_i$, $X_i$,
\[
 \frac1{G_{ii}}=h_{ii}-z-Z_i-\sum^{(i)}_k\sigma_{ik}^2G^{(i)}_{kk}=-z+X_i-\sum^{(i)}_k\sigma_{ik}^2G^{(i)}_{kk}.
\]
By (I3) with $A=\{i\}$, $G^{(i)}_{kk}=G_{kk}-G_{ki}G_{ik}/G_{ii}$ for $k\ne i$, so
\[
 \sum^{(i)}_k\sigma_{ik}^2G^{(i)}_{kk}=\sum_{k}\sigma_{ik}^2G_{kk}-\sigma_{ii}^2G_{ii}-\sum_{k\ne i}\sigma_{ik}^2\frac{G_{ik}G_{ki}}{G_{ii}}=\sum_k\sigma_{ik}^2G_{kk}-A_i .
\]
Inserting this gives $1/G_{ii}=-z-\sum_k\sigma_{ik}^2G_{kk}+A_i+X_i$.

\emph{Second equation of \eqref{eq:sce1}.} Put $t_i:=\sum_j\sigma_{ij}^2v_j-\Upsilon_i$. By \eqref{eq:Arow}, $\sum_j\sigma_{ij}^2G_{jj}=m+\sum_j\sigma_{ij}^2v_j$, and by Lemma~\ref{lem:msc}(a), $m\ne0$ and $-z-m=1/m$. So the first equation reads
$1/G_{ii}=1/m-t_i$. Multiplying by $mG_{ii}$: $m=G_{ii}-mG_{ii}t_i$, i.e.\ $v_i=mG_{ii}t_i$.

\emph{First equation of \eqref{eq:sce2}.} From $v_i=mG_{ii}t_i$ and $G_{ii}=m+v_i$: $v_i=m^2t_i+mv_it_i$, and $mt_i=v_i/G_{ii}$, so $mv_it_i=v_i^2/G_{ii}=r_i$. Thus $v_i=m^2t_i+r_i$, which is the claim.

\emph{Second equation of \eqref{eq:sce2}.} Average the first one over $i$ and use (A): $\frac1N\sum_i\sum_j\sigma_{ij}^2v_j=\frac1N\sum_jv_j\sum_i\sigma_{ij}^2=\avg v$.

\emph{Formula for $A_i$.} For $j\ne i$, the first and third formulas of (I2) give $G_{ij}G_{ji}=G_{ii}^2\big(\sum^{(i)}_kh_{ik}G^{(i)}_{kj}\big)\big(\sum^{(i)}_lG^{(i)}_{jl}h_{li}\big)$; divide by $G_{ii}$.
\end{proof}

\begin{lemma}[Vector stability on $\be^\perp$, no logarithm]\label{lem:vecstab}
Assume \textup{(A), (B), (C)} and let $\delta_{\min}:=\min(\delta_-,\delta_+)$. Let $\zeta\in\C$, $|\zeta|\le1$, and $u\in\C^N$ with $\sum_iu_i=0$. There is a unique $d\in\C^N$ with $\sum_id_i=0$ and $(1-\zeta B)d=u$, and
\[
 |d|\le\delta_{\min}^{-1}|u|,\qquad\norm d_\infty\le C_s\norm u_\infty,\qquad C_s:=1+\frac{\sqrt{C_0}}{\delta_{\min}}.
\]
\end{lemma}
\begin{proof}
Let $V:=\{x\in\C^N:\sum_ix_i=0\}$.

\emph{Step 1: $B$ maps $V$ into $V$.} For $x\in V$, by (A), $\sum_i(Bx)_i=\sum_jx_j\sum_i\sigma_{ij}^2=\sum_jx_j=0$.

\emph{Step 2: the quadratic form on $V$.} Let $x\in V$ and write $x=a+\ii b$ with $a,b\in\R^N$; then $\sum_ia_i=\sum_ib_i=0$. Since $B$ is real and symmetric, $\langle a,Bb\rangle=\langle b,Ba\rangle$ and
\[
 \langle x,Bx\rangle=\langle a,Ba\rangle+\langle b,Bb\rangle+\ii\big(\langle a,Bb\rangle-\langle b,Ba\rangle\big)=\langle a,Ba\rangle+\langle b,Bb\rangle .
\]
Applying (B) to $a$ and to $b$, and $|x|^2=|a|^2+|b|^2$: $(-1+\delta_-)|x|^2\le\langle x,Bx\rangle\le(1-\delta_+)|x|^2$, hence
\[
 |\langle x,Bx\rangle|\le(1-\delta_{\min})|x|^2\qquad(x\in V).
\]
(If $\delta_{\min}>1$ this forces $V=\{0\}$, which is impossible for $N\ge2$; so $\delta_{\min}\le1$.)

\emph{Step 3: coercivity, existence and uniqueness.} For $x\in V$, by the Cauchy--Schwarz inequality, Step~2 and $|\zeta|\le1$,
\[
 |x|\,|(1-\zeta B)x|\ \ge\ |\langle x,(1-\zeta B)x\rangle|\ \ge\ |x|^2-|\zeta|\,|\langle x,Bx\rangle|\ \ge\ \delta_{\min}|x|^2 .
\]
So $|(1-\zeta B)x|\ge\delta_{\min}|x|$ on $V$. By Step~1, $1-\zeta B$ maps the finite-dimensional space $V$ into itself; it is injective on $V$, hence bijective on $V$. This gives existence and uniqueness of $d\in V$ with $(1-\zeta B)d=u$, and $|d|\le\delta_{\min}^{-1}|u|$.

\emph{Step 4: the bound in the maximum norm.} From $d=u+\zeta Bd$, the Cauchy--Schwarz inequality, (C) and \eqref{eq:Arow},
\[
 |(Bd)_i|=\Big|\sum_j\sigma_{ij}^2d_j\Big|\le\Big(\sum_j\sigma_{ij}^4\Big)^{1/2}|d|\le\Big(\frac{C_0}N\sum_j\sigma_{ij}^2\Big)^{1/2}|d|=\sqrt{\frac{C_0}{N}}\,|d| ,
\]
and $|d|\le\delta_{\min}^{-1}|u|\le\delta_{\min}^{-1}\sqrt N\norm u_\infty$. Hence $\norm d_\infty\le\norm u_\infty+\norm{Bd}_\infty\le(1+\sqrt{C_0}/\delta_{\min})\norm u_\infty$.
\end{proof}
\noindent Neither a Neumann series nor the spectral theorem for $B$ is used; the lemma holds for all $\delta_\pm,C_0>0$, without the normalisation of Convention~\ref{conv:const}.
Note that $C_s\ge1$ always, and $C_s\ge2$ under that normalisation ($C_0\ge1$, $\delta_{\min}\le1$).

\begin{lemma}[Stability, near regime]\label{lem:stabnear}
Assume \textup{(A), (B), (C)}. Let $|E|\le5$, $0<\eta\le10$, and
\begin{equation}\label{eq:lambda1}
 \lambda_{\rm st}:=\frac1{576\,C_s}.
\end{equation}
Put $w:=\avg v$ (so $|w|=\Lambda$), $w_0:=\msc(1-\msc^2)$, $d_i:=v_i-w$, $D:=\max_i|v_i-w|$, $\Upsilon_{\max}:=\max_i|\Upsilon_i|$. If $\Ld\le\lambda_{\rm st}$, then $\frac3{52}\le|G_{ii}|\le\frac54$ for all $i$,
$\Lambda\le\Ld$, $D\le2\Ld\le2\lambda_{\rm st}$, and
\begin{align}
 D&\le4C_s\Upsilon_{\max}+144C_s\Lambda^2,\qquad\Ld\le\Lambda+D,\label{eq:Dbound}\\
 |w^2-w_0w|&\le|\avg\Upsilon|+D^2+18\Ld^3\ \le\ |\avg\Upsilon|+40C_s^2\,\Upsilon_{\max}^2+(90+72C_s)\Lambda^3,\label{eq:quad}
\end{align}
and $\frac1{13}c_*\varkappa\le|w_0|\le C_*\varkappa$, with $c_*,C_*$ from Lemma~\ref{lem:mscsize}. Moreover the following identity holds (whenever $\eta>0$, without any assumption on $E$, $\eta$ or $\Ld$):
\[
 w^2-w_0w=m^3\avg\Upsilon-\avg{d^2}+\avg{v^3/G},\qquad\text{where }\ \avg{v^3/G}:=\frac1N\sum_i\frac{v_i^3}{G_{ii}} .
\]
\end{lemma}
\noindent(Note $(90+72C_s)\lambda_{\rm st}\le\frac{162}{576}<\frac3{10}$, so the cubic term is at most $\frac3{10}\Lambda^2$ when $\Ld\le\lambda_{\rm st}$.)
\begin{proof}
Throughout, $m=\msc(z)$, and $\frac1{13}\le|m|\le1$ by \eqref{eq:alpha}.

\emph{Step 1: elementary bounds.} Since $C_s\ge1$, $\lambda_{\rm st}\le\frac1{576}\le\frac1{52}$. So $|G_{ii}|\ge|m|-|v_i|\ge\frac1{13}-\frac1{52}=\frac3{52}$ and $|G_{ii}|\le1+\frac1{52}\le\frac54$. Consequently
\[
 |r_i|=\frac{|v_i|^2}{|G_{ii}|}\le\frac{52}3\Ld^2\le18\Ld^2,\qquad\Big|\frac{v_i^3}{G_{ii}}\Big|\le18\Ld^3 .
\]
Further $\Lambda=|\avg v|\le\max_i|v_i|=\Ld$, $|d_i|\le|v_i|+|w|\le2\Ld$, so $D\le2\Ld\le2\lambda_{\rm st}$; and $|v_i|\le|w|+|d_i|$ gives $\Ld\le\Lambda+D$.

\emph{Step 2: the equation for $d$.} Let $y:=-m^2\Upsilon+r\in\C^N$ and $d:=v-w\bone$, so $\sum_id_i=0$. The first equation of \eqref{eq:sce2} is $(1-m^2B)v=y$. By \eqref{eq:Arow}, $B\bone=\bone$, hence
\[
 (1-m^2)w\bone+(1-m^2B)d=y .
\]
By Step~1 of the proof of Lemma~\ref{lem:vecstab}, $\sum_i((1-m^2B)d)_i=0$. Taking the average of both sides gives $(1-m^2)w=\avg y$ (this is the second equation of \eqref{eq:sce2}), and subtracting,
\[
 (1-m^2B)d=y-\avg y\bone=:u,\qquad\sum_iu_i=0,\qquad\norm u_\infty\le2\norm y_\infty\le2\big(\Upsilon_{\max}+18\Ld^2\big),
\]
where we used $|m|\le1$ and Step~1. By Lemma~\ref{lem:vecstab} with $\zeta:=m^2$ ($|\zeta|\le1$; $d$ is \emph{the} solution by uniqueness),
\[
 D=\norm d_\infty\le2C_s\Upsilon_{\max}+36C_s\Ld^2\le2C_s\Upsilon_{\max}+72C_s\Lambda^2+72C_sD^2,
\]
using $\Ld^2\le(\Lambda+D)^2\le2\Lambda^2+2D^2$. Since $72C_sD\le144C_s\lambda_{\rm st}=\frac14$, we get $\frac34D\le2C_s\Upsilon_{\max}+72C_s\Lambda^2$, hence $D\le\frac83C_s\Upsilon_{\max}+96C_s\Lambda^2$, which implies \eqref{eq:Dbound}.

\emph{Step 3: the identity for $w$.} Since $G_{ii}=m+v_i$,
\[
 \frac{v_i^2}{m}-\frac{v_i^3}{mG_{ii}}=\frac{v_i^2(G_{ii}-v_i)}{mG_{ii}}=\frac{v_i^2}{G_{ii}}=r_i,\qquad\text{so}\qquad\avg r=\frac1m\avg{v^2}-\frac1m\avg{v^3/G}.
\]
Since $\avg d=0$, $\avg{v^2}=\avg{(w\bone+d)^2}=w^2+\avg{d^2}$. Multiplying the second equation of \eqref{eq:sce2} by $-m$ therefore gives
\[
 -m(1-m^2)w=m^3\avg\Upsilon-w^2-\avg{d^2}+\avg{v^3/G},
\]
which is the asserted identity because $w_0=m(1-m^2)$. Only $\eta>0$ and (A) were used here.

\emph{Step 4: \eqref{eq:quad}.} By Step~3, $|m|\le1$, $|\avg{d^2}|\le D^2$ and Step~1: $|w^2-w_0w|\le|\avg\Upsilon|+D^2+18\Ld^3$. Next, $(a+b)^3\le4(a^3+b^3)$ for $a,b\ge0$, so by Step~1
\[
 18\Ld^3\le72\Lambda^3+72D^3\le72\Lambda^3+144\lambda_{\rm st}D^2=72\Lambda^3+\frac{D^2}{4C_s} .
\]
By \eqref{eq:Dbound}, $(a+b)^2\le2a^2+2b^2$ and $\Lambda\le\Ld\le\lambda_{\rm st}$,
\[
 D^2\le32C_s^2\Upsilon_{\max}^2+2\cdot144^2C_s^2\Lambda^4\le32C_s^2\Upsilon_{\max}^2+\frac{2\cdot144^2}{576}C_s\Lambda^3=32C_s^2\Upsilon_{\max}^2+72C_s\Lambda^3 .
\]
Hence, using $C_s\ge1$ in the form $32C_s^2+8C_s\le40C_s^2$,
\[
\begin{aligned}
 D^2+18\Ld^3&\le\Big(1+\frac1{4C_s}\Big)D^2+72\Lambda^3\le(32C_s^2+8C_s)\Upsilon_{\max}^2+(72C_s+18+72)\Lambda^3\\
 &\le40C_s^2\Upsilon_{\max}^2+(90+72C_s)\Lambda^3 .
\end{aligned}
\]

\emph{Step 5: size of $w_0$.} $|w_0|=|m|\,|1-m^2|$, and \eqref{eq:alpha} gives $\frac1{13}c_*\varkappa\le|w_0|\le C_*\varkappa$.
\end{proof}

\begin{lemma}[Scalar quadratic inequality]\label{lem:scalar}
Let $w,w_0\in\C$, $R\ge0$ with $|w^2-w_0w|\le R$. Then \textup{(i)} $|w|\le|w_0|+\sqrt R$; \textup{(ii)} if $|w|\le|w_0|/2$ then $|w_0|\,|w|\le2R$, i.e.\ $|w|\le2R/|w_0|$ when $w_0\ne0$ (and $w=0$ when $w_0=0$).
Moreover, for reals $x,b,c\ge0$: $x^2\le bx+c$ implies $x\le b+\sqrt c$.
\end{lemma}
\begin{proof}
\emph{The real statement.} If $x>b+\sqrt c$, then $x-b>\sqrt c\ge0$ and $x>\sqrt c$, so $x^2-bx=x(x-b)>c$, contradicting $x^2\le bx+c$.

(i) $|w|^2=|w^2|\le|w_0w|+|w^2-w_0w|\le|w_0|\,|w|+R$; apply the real statement with $x=|w|$, $b=|w_0|$, $c=R$.

(ii) If $|w|\le|w_0|/2$ then $|w-w_0|\ge|w_0|-|w|\ge|w_0|/2$, so $\frac12|w_0|\,|w|\le|w|\,|w-w_0|=|w^2-w_0w|\le R$.
\end{proof}

\begin{lemma}[Stability, far regime]\label{lem:stabfar}
Assume \textup{(A)}. Let $|E|\ge4$, $\eta>0$ and suppose $\max_i|G_{ii}|\le1$. Then, for every $i$,
\[
 \Ld\le\tfrac12\Upsilon_{\max},\qquad\Lambda\le\tfrac18|\avg\Upsilon|+\tfrac98\,\Ld^2\,(|z|+1+\Upsilon_{\max}),\qquad\frac1{|G_{ii}|}\le|z|+1+\Upsilon_{\max}.
\]
\end{lemma}
\begin{proof}
By Lemma~\ref{lem:msc}(g), $|m|\le\frac13$ and $|1-m^2|\ge\frac89$. By the second equation of \eqref{eq:sce1}, $|G_{ii}|\le1$ and \eqref{eq:Arow},
\[
 |v_i|\le|m|\,|G_{ii}|\Big(\sum_j\sigma_{ij}^2|v_j|+|\Upsilon_i|\Big)\le\tfrac13\big(\Ld+\Upsilon_{\max}\big)\qquad\text{for every }i,
\]
so $\Ld\le\frac13\Ld+\frac13\Upsilon_{\max}$, i.e.\ $\Ld\le\frac12\Upsilon_{\max}$. By the first equation of \eqref{eq:sce1}, $|G_{jj}|\le1$ and \eqref{eq:Arow},
\[
 \frac1{|G_{ii}|}\le|z|+\sum_j\sigma_{ij}^2|G_{jj}|+|\Upsilon_i|\le|z|+1+\Upsilon_{\max}.
\]
Therefore $|r_i|=|v_i|^2/|G_{ii}|\le\Ld^2(|z|+1+\Upsilon_{\max})$, and the second equation of \eqref{eq:sce2} gives
\[
 \tfrac89\Lambda\le|1-m^2|\,|\avg v|\le|m|^2|\avg\Upsilon|+|\avg r|\le\tfrac19|\avg\Upsilon|+\Ld^2(|z|+1+\Upsilon_{\max}).
\]
Multiply by $\frac98$.
\end{proof}

\subsection{Eigenvalues: continuity and a net for the operator norm}\label{sec:eig}
For a Hermitian $n\times n$ matrix $A$ ($n\ge1$; a real symmetric matrix is regarded as a complex Hermitian one) we denote by
$\lambda_1(A)\le\lambda_2(A)\le\dots\le\lambda_n(A)$ its eigenvalues, i.e.\ the roots of the characteristic polynomial $\det(t-A)$ counted with multiplicity, in nondecreasing order.
By the spectral theorem (Mathlib) they are real and there is an orthonormal basis $u_1,\dots,u_n$ of $\C^n$ with $Au_\alpha=\lambda_\alpha(A)u_\alpha$; conversely, if $(u'_\alpha)$ is any orthonormal basis of eigenvectors of $A$ with eigenvalues $\mu_\alpha$, then
$\det(t-A)=\prod_\alpha(t-\mu_\alpha)$ (conjugation by the unitary matrix with columns $u'_\alpha$), so $(\mu_\alpha)$ is a rearrangement of $(\lambda_\alpha(A))$.
Recall $\Omega=\kk^{\mathcal P_N}$ and the map $x\mapsto H(x)$ of Definition~\ref{def:ens}; on $\Omega$ we use the Euclidean norm $|x|^2=\sum_{(i,j)\in\mathcal P_N}|x_{ij}|^2$. The product $\sigma$-algebra of $\Omega$ is its Borel $\sigma$-algebra (finite product of separable metric spaces).

\begin{lemma}[Weyl inequality and measurability]\label{lem:weyl} \infra\
Let $A,A'$ be Hermitian $n\times n$ matrices. Then:
\begin{enumerate}[label=\textup{(\alph*)},leftmargin=2.2em,itemsep=1pt]
\item \emph{Rayleigh bounds:} $\lambda_1(A)|x|^2\le\langle x,Ax\rangle\le\lambda_n(A)|x|^2$ for all $x\in\C^n$, and
$\norm A=\max_k|\lambda_k(A)|=\max\big(\lambda_n(A),-\lambda_1(A)\big)=\max_{|x|=1}|\langle x,Ax\rangle|$.
\item \emph{Weyl's inequality:} $|\lambda_k(A)-\lambda_k(A')|\le\norm{A-A'}$ for every $k\in\{1,\dots,n\}$.
\item \emph{Reflection}: $\lambda_k(-A)=-\lambda_{n+1-k}(A)$ for every $k$; in particular $\lambda_1(A)=-\lambda_n(-A)$.
\item \emph{Lipschitz continuity on $\Omega$:} $x\mapsto H(x)$ is $\R$-linear, $\norm{H(x)-H(x')}\le\sqrt2\,|x-x'|$, and $|\lambda_k(H(x))-\lambda_k(H(x'))|\le\sqrt2\,|x-x'|$ for all $x,x'\in\Omega$ and all $k$.
\item \emph{Resolvents:} for fixed $z=E+\ii\eta$, $\eta>0$, and $T\subsetneq\{1,\dots,N\}$, the functions $x\mapsto G^{(T)}_{ij}(z)$, $m^{(T)}(z)$, $\Lambda(z)$, $\Ld(z)$, $\Lo(z)$ are Lipschitz on $\Omega$ with constant $\sqrt2\,\eta^{-2}$;
the functions $(x,z)\mapsto G^{(T)}_{ij}(z)$ are continuous on $\Omega\times\{\im z>0\}$; and for fixed $z$ the functions $1/G_{ii}$, $M^A_{ab}$, $Z_i$, $X_i$, $A_i$, $\Upsilon_i$ of \eqref{eq:MA} and Definition~\ref{def:XA} are continuous on $\Omega$.
\item \emph{Counting function:} for every $E\in\R$ and every $j$ the set $\{x:\lambda_j(H(x))\le E\}$ is closed, and $x\mapsto\fn(E)=\frac1N\#\{j:\lambda_j(H(x))\le E\}$ is a Borel function. For $J=\R$ or $J=[a,b]$ $(a<b)$,
\[
 \sup_{E\in J}|\fn(E)-\nsc(E)|=\sup_{E\in J_0}|\fn(E)-\nsc(E)|,\qquad J_0:=(J\cap\mathbb Q)\cup\{b\}\ \ (J_0:=\mathbb Q\text{ if }J=\R),
\]
at every point of $\Omega$, and this supremum is a Borel function on $\Omega$.
\end{enumerate}
In particular all the functions in \textup{(d)--(f)} are Borel measurable on $\Omega$.
\end{lemma}
\begin{proof}
Fix an orthonormal eigenbasis $(u_\alpha)$ of $A$ with $Au_\alpha=\lambda_\alpha u_\alpha$, $\lambda_\alpha:=\lambda_\alpha(A)$.

(a) For $x=\sum_\alpha c_\alpha u_\alpha$ we have $|x|^2=\sum_\alpha|c_\alpha|^2$, $\langle x,Ax\rangle=\sum_\alpha\lambda_\alpha|c_\alpha|^2$ and $|Ax|^2=\sum_\alpha\lambda_\alpha^2|c_\alpha|^2$.
The first gives the Rayleigh bounds. With $\rho:=\max_\alpha|\lambda_\alpha|=\max(\lambda_n,-\lambda_1)$ we get $|Ax|\le\rho|x|$ and $|\langle x,Ax\rangle|\le\rho|x|^2$, with equality in both for $x=u_\alpha$ when $|\lambda_\alpha|=\rho$.
Hence $\norm A=\rho=\max_{|x|=1}|\langle x,Ax\rangle|$.

(b) \emph{One-sided min--max.} We claim: for $t\in\R$ and $1\le k\le n$,
\begin{equation}\tag{MM}\label{eq:MM}
 \lambda_k(A)\le t\iff\text{there is a subspace }V\subset\C^n\text{ with }\dim V=k\text{ and }\langle x,Ax\rangle\le t|x|^2\text{ for all }x\in V.
\end{equation}
``$\Rightarrow$'': take $V:=\operatorname{span}(u_1,\dots,u_k)$; for $x=\sum_{\alpha\le k}c_\alpha u_\alpha$, $\langle x,Ax\rangle=\sum_{\alpha\le k}\lambda_\alpha|c_\alpha|^2\le\lambda_k|x|^2\le t|x|^2$.
``$\Leftarrow$'': let $W:=\operatorname{span}(u_k,\dots,u_n)$, of dimension $n-k+1$. Since $\dim V+\dim W=n+1>n$, there is $x\ne0$ in $V\cap W$. Writing $x=\sum_{\alpha\ge k}c_\alpha u_\alpha$,
$\lambda_k|x|^2\le\sum_{\alpha\ge k}\lambda_\alpha|c_\alpha|^2=\langle x,Ax\rangle\le t|x|^2$, so $\lambda_k\le t$.

\emph{Weyl.} Let $\eps:=\norm{A-A'}$ and $V:=\operatorname{span}(u_1,\dots,u_k)$. For $x\in V$,
$\langle x,A'x\rangle=\langle x,Ax\rangle+\langle x,(A'-A)x\rangle\le(\lambda_k(A)+\eps)|x|^2$. By \eqref{eq:MM} for $A'$ with $t=\lambda_k(A)+\eps$: $\lambda_k(A')\le\lambda_k(A)+\eps$. Exchanging $A$ and $A'$ gives the other inequality.

(c) $(u_\alpha)$ is an orthonormal eigenbasis of $-A$ with eigenvalues $-\lambda_\alpha$. By the remark before the lemma, the eigenvalues of $-A$ counted with multiplicity are the numbers $-\lambda_1\ge-\lambda_2\ge\dots\ge-\lambda_n$;
in nondecreasing order the $k$-th of them is $-\lambda_{n+1-k}$.

(d) Linearity over $\R$ is clear from the definition of $H(x)$, so $H(x)-H(x')=H(y)$ with $y:=x-x'$. For any matrix $M$, $|Mx|^2=\sum_i|\sum_jM_{ij}x_j|^2\le\sum_{i,j}|M_{ij}|^2\,|x|^2$ by the Cauchy--Schwarz inequality. Hence
\[
 \norm{H(y)}^2\le\sum_{i,j}|H(y)_{ij}|^2=\sum_i(\re y_{ii})^2+2\sum_{i<j}|y_{ij}|^2\le2|y|^2 .
\]
The eigenvalue bound follows from (b).

(e) For $T\subsetneq\{1,\dots,N\}$ and vectors $\xi$ supported on $T^c$, $|(H^{(T)}(x)-H^{(T)}(x'))\xi|\le|(H(x)-H(x'))\xi|$ (the left side keeps only the coordinates in $T^c$), so $\norm{H^{(T)}(x)-H^{(T)}(x')}\le\sqrt2|x-x'|$ by (d).
By Lemma~\ref{lem:basic}(h) applied to the Hermitian matrices $H^{(T)}(x)$, $H^{(T)}(x')$, the matrix $G^{(T)}(z)$ changes by at most $\sqrt2\eta^{-2}|x-x'|$ in operator norm, hence each entry does.
$m^{(T)}$ is $N^{-1}$ times a sum of at most $N$ diagonal entries; $\msc(z)$ does not depend on $x$; and $||a|-|b||\le|a-b|$, $|\max_ka_k-\max_kb_k|\le\max_k|a_k-b_k|$. This gives the Lipschitz bounds.
Joint continuity: by Lemma~\ref{lem:basic}(c),(h), $|G^{(T)}_{ij}(z;x)-G^{(T)}_{ij}(z';x')|\le\sqrt2\eta^{-2}|x-x'|+|z-z'|/(\eta\eta')$.
The entries $h_{ij}$ are continuous functions of $x$; $1/G_{ii}$ is continuous because $G_{ii}$ is continuous and never zero (Lemma~\ref{lem:basic}(a)); $M^A_{ab}$, $Z_i$, $X_i$, $A_i$, $\Upsilon_i$ are finite sums of products of the functions $h_{kl}$, $G^{(T)}_{kl}$, $1/G_{ii}$ and constants.

(f) By (d) the function $x\mapsto\lambda_j(H(x))$ is continuous, so $\{\lambda_j(H(x))\le E\}$ is closed and $\fn(E)=\frac1N\sum_j\bone(\lambda_j\le E)$ is Borel.
Fix $x$ and put $\varphi(E):=|\fn(E)-\nsc(E)|$. The function $\fn$ is right-continuous: if $\eps>0$ is smaller than $\lambda_j-E$ for all $j$ with $\lambda_j>E$, then $\fn=\fn(E)$ on $[E,E+\eps)$.
The function $\nsc$ is continuous (Lemma~\ref{lem:nsc}). So $\varphi$ is right-continuous, and for every $E\in J$ that is not the right endpoint of $J$ there are $q_l\in J\cap\mathbb Q$, $q_l\downarrow E$, with $\varphi(q_l)\to\varphi(E)$.
Hence $\sup_J\varphi=\sup_{J_0}\varphi$. A countable supremum of Borel functions is Borel.

Finally, continuous functions on $\Omega$ are Borel.
\end{proof}

\begin{lemma}[Grid net for the norm]\label{lem:net} \infra\
Let $n:=\beta N$, where $\beta:=\dim_\R\kk\in\{1,2\}$, and identify $\kk^N$ with $\R^{n}$ (real and imaginary parts of the coordinates; this identification preserves the Euclidean norm).
Let $\delta_n:=1/(4\sqrt n)$ and $\mathcal N:=\{y\in(\delta_n\Z)^n:\ |y|\le1\}\subset\kk^N$. Then
\[
 |\mathcal N|\le(9\sqrt n)^n\le e^{3N\log N}\quad\text{for }N\ge13,
\]
and for every Hermitian $N\times N$ matrix $A$ with entries in $\kk$ (real symmetric if $\kk=\R$):\quad$\norm A\le2\max_{y\in\mathcal N}|\langle y,Ay\rangle|$.
\end{lemma}
\begin{proof}
\emph{Cardinality.} If $y\in\mathcal N$ then every real coordinate is $y_i=\delta_nk_i$ with $k_i\in\Z$, $|k_i|\le1/\delta_n=4\sqrt n$; there are at most $8\sqrt n+1\le9\sqrt n$ such integers. So $|\mathcal N|\le(9\sqrt n)^n$.
Since $n\le2N$, $\ n\log(9\sqrt n)\le2N\big(\log9+\tfrac12\log(2N)\big)=N\big(\log162+\log N\big)\le3N\log N$ as soon as $N^2\ge162$, i.e.\ $N\ge13$.

\emph{A maximiser in $\kk^N$.} We claim that there is $x\in\kk^N$ with $|x|=1$ and $|\langle x,Ax\rangle|=\norm A$. By Lemma~\ref{lem:weyl}(a) there is a unit vector $u\in\C^N$ and $\lambda\in\R$ with $Au=\lambda u$, $|\lambda|=\norm A$.
If $\kk=\C$ take $x:=u$. If $\kk=\R$, write $u=a+\ii b$ with $a,b\in\R^N$; since $A$ and $\lambda$ are real, $Aa=\lambda a$ and $Ab=\lambda b$, and $a\ne0$ or $b\ne0$; take $x:=a/|a|$ or $x:=b/|b|$.

\emph{Rounding.} Let $x_1,\dots,x_n$ be the real coordinates of $x$ and put $y_i:=\delta_n\operatorname{sgn}(x_i)\lfloor|x_i|/\delta_n\rfloor$. Then $y_i\in\delta_n\Z$, $|y_i|\le|x_i|$ and $|x_i-y_i|<\delta_n$, so
$|y|\le|x|=1$, $y\in\mathcal N$, and $|x-y|\le\sqrt n\,\delta_n=\frac14$. Now
\[
 \langle x,Ax\rangle-\langle y,Ay\rangle=\langle x-y,Ax\rangle+\langle y,A(x-y)\rangle,
\]
so $\big|\langle x,Ax\rangle-\langle y,Ay\rangle\big|\le\norm A\,|x-y|\,(|x|+|y|)\le\tfrac12\norm A$.
Therefore $\norm A=|\langle x,Ax\rangle|\le|\langle y,Ay\rangle|+\frac12\norm A$, i.e.\ $\norm A\le2|\langle y,Ay\rangle|\le2\max_{\mathcal N}|\langle\cdot,A\,\cdot\rangle|$.
\end{proof}

\subsection{Helffer--Sj\"ostrand formula and the counting function}\label{sec:HS}
Throughout this subsection $\mathsf s:\R\to[0,1]$ denotes the quintic step
\[
 \mathsf s(t):=0\ \ (t\le0),\qquad \mathsf s(t):=10t^3-15t^4+6t^5\ \ (0\le t\le1),\qquad \mathsf s(t):=1\ \ (t\ge1).
\]
On $[0,1]$: $\mathsf s'(t)=30t^2(1-t)^2$ and $\mathsf s''(t)=60t(1-t)(1-2t)$; both vanish at $t=0$ and $t=1$, and $\mathsf s(1)=1$. Hence $\mathsf s\in C^2(\R)$, $\mathsf s$ is nondecreasing, $0\le\mathsf s'\le\mathsf s'(\frac12)=\frac{15}8$,
and, since $\mathsf s'$ increases on $[0,\frac12]$ and decreases on $[\frac12,1]$,
\begin{equation}\tag{Q}\label{eq:a1-quintic}
 \int_\R|\mathsf s'|=1,\qquad\int_\R|\mathsf s''|=2\,\mathsf s'(\tfrac12)=\tfrac{15}4 .
\end{equation}

\begin{lemma}[Helffer--Sj\"ostrand formula]\label{lem:HS} \infra\
Let $f\in C^2(\R)$ be real valued with compact support and let $\chi\in C^1(\R)$ be even with $\chi=1$ on $[-1,1]$, $\chi=0$ outside $[-2,2]$ and $0\le\chi\le1$. Put
\[
 g(x,y):=\ii yf''(x)\chi(y)+\ii\big(f(x)+\ii yf'(x)\big)\chi'(y)\qquad(x,y\in\R).
\]
Then for every $\lambda\in\R$
\begin{equation}\label{eq:HS}
 f(\lambda)=\frac1{2\pi}\int_{\R}\int_{\R}\frac{\ii yf''(x)\chi(y)+\ii\big(f(x)+\ii yf'(x)\big)\chi'(y)}{\lambda-x-\ii y}\,\rd x\,\rd y,
\end{equation}
the integral being absolutely convergent: for all $\lambda,x\in\R$ and $y\ne0$,
\begin{equation}\tag{HS$'$}\label{eq:a1-HSdom}
 \Big|\frac{g(x,y)}{\lambda-x-\ii y}\Big|\le\Gamma(x,y):=|f''(x)|\,\bone(|y|\le2)+|\chi'(y)|\big(|f(x)|+2|f'(x)|\big),\qquad\int_\R\int_\R\Gamma\,\rd x\,\rd y<\infty .
\end{equation}
Moreover $g(x,-y)=\ol{g(x,y)}$.
\end{lemma}
\begin{proof}
\emph{The function $\tilde f$.} Define $\tilde f:\R^2\to\C$ by $\tilde f(x,y):=\big(f(x)+\ii yf'(x)\big)\chi(y)$. Since $f\in C^2$ and $\chi\in C^1$, $\tilde f$ is continuously differentiable on $\R^2$, with
\[
 \partial_x\tilde f=\big(f'(x)+\ii yf''(x)\big)\chi(y),\qquad\partial_y\tilde f=\ii f'(x)\chi(y)+\big(f(x)+\ii yf'(x)\big)\chi'(y),
\]
so that
\begin{equation}\tag{HS1}\label{eq:a1-HS1}
 \partial_x\tilde f+\ii\,\partial_y\tilde f=\ii yf''(x)\chi(y)+\ii\big(f(x)+\ii yf'(x)\big)\chi'(y)=g(x,y).
\end{equation}
Fix $\rho>0$ with $\operatorname{supp}f\subset(-\rho,\rho)$. Then $\tilde f(x,y)=0$ if $|x|\ge\rho$ or $|y|\ge2$, and $\tilde f(x,y)=f(x)+\ii yf'(x)$ if $|y|\le1$.
Since $\chi$ is even, $\chi'$ is odd, and since $f$ is real, $\ol{g(x,y)}=-\ii yf''(x)\chi(y)-\ii(f(x)-\ii yf'(x))\chi'(y)=g(x,-y)$.

\emph{Absolute convergence.} $|\lambda-x-\ii y|\ge|y|$, so $|\ii yf''(x)\chi(y)|/|\lambda-x-\ii y|\le|f''(x)|\bone(|y|\le2)$. The function $\chi'$ vanishes for $|y|<1$ and for $|y|>2$; for $1\le|y|\le2$ we have $|\lambda-x-\ii y|\ge1$ and
$|f(x)+\ii yf'(x)|\le|f(x)|+2|f'(x)|$. This proves the pointwise bound in \eqref{eq:a1-HSdom}; $\Gamma$ is bounded with compact support, hence integrable.

\emph{The kernel.} Let $h(x,y):=(\lambda-x-\ii y)^{-1}$ for $(x,y)\ne(\lambda,0)$. It is smooth there, with $\partial_xh=h^2$ and $\partial_yh=\ii h^2$, hence
\begin{equation}\tag{HS2}\label{eq:a1-HS2}
 \partial_xh+\ii\,\partial_yh=0 .
\end{equation}
Let $F:=\tilde fh$ on $\{y\ne0\}$. By \eqref{eq:a1-HS1}, \eqref{eq:a1-HS2}: $\partial_xF+\ii\partial_yF=gh$, and $F$ is continuously differentiable on $\{y\ne0\}$.

\emph{Divergence theorem.} Let $0<\eps<1$. On the rectangle $[-\rho,\rho]\times[\eps,3]$ the divergence theorem (Mathlib), applied to the continuously differentiable vector field $(F,\ii F)$, gives
\[
 \int_\eps^3\!\!\int_{-\rho}^{\rho}g\,h\,\rd x\,\rd y=\int_\eps^3\big(F(\rho,y)-F(-\rho,y)\big)\rd y+\ii\int_{-\rho}^{\rho}\big(F(x,3)-F(x,\eps)\big)\rd x=-\ii\int_{-\rho}^\rho F(x,\eps)\,\rd x,
\]
because $\tilde f$ vanishes for $|x|=\rho$ and for $y=3$. Since $\chi(\eps)=1$ and $g=0$ outside $[-\rho,\rho]\times[-2,2]$,
\[
 \int_{y>\eps}\!\int_\R\frac{g(x,y)}{\lambda-x-\ii y}\,\rd x\,\rd y=-\ii\int_\R\frac{f(x)+\ii\eps f'(x)}{\lambda-x-\ii\eps}\,\rd x .
\]
In the same way, on $[-\rho,\rho]\times[-3,-\eps]$,
\[
 \int_{y<-\eps}\!\int_\R\frac{g(x,y)}{\lambda-x-\ii y}\,\rd x\,\rd y=\ii\int_{-\rho}^\rho\big(F(x,-\eps)-F(x,-3)\big)\rd x=\ii\int_\R\frac{f(x)-\ii\eps f'(x)}{\lambda-x+\ii\eps}\,\rd x .
\]
With $a:=\lambda-x$ we have $\frac1{a+\ii\eps}-\frac1{a-\ii\eps}=\frac{-2\ii\eps}{a^2+\eps^2}$ and $\frac1{a+\ii\eps}+\frac1{a-\ii\eps}=\frac{2a}{a^2+\eps^2}$. Adding the two identities and dividing by $2\pi$:
\begin{equation}\tag{HS3}\label{eq:a1-HS3}
 \frac1{2\pi}\int_{|y|>\eps}\!\int_\R\frac{g(x,y)}{\lambda-x-\ii y}\,\rd x\,\rd y=\frac1\pi\int_\R f(x)\frac{\eps}{(\lambda-x)^2+\eps^2}\,\rd x+\frac1\pi\int_\R f'(x)\frac{\eps(\lambda-x)}{(\lambda-x)^2+\eps^2}\,\rd x .
\end{equation}

\emph{The limit $\eps\downarrow0$.} By \eqref{eq:a1-HSdom} and dominated convergence the left side of \eqref{eq:a1-HS3} converges to the right side of \eqref{eq:HS}.
In the first term on the right substitute $x=\lambda-\eps s$: it equals $\frac1\pi\int_\R f(\lambda-\eps s)\frac{\rd s}{1+s^2}$, which converges to $f(\lambda)\cdot\frac1\pi\int_\R\frac{\rd s}{1+s^2}=f(\lambda)$ by dominated convergence ($f$ is bounded and continuous).
In the second term, $\big|\frac{\eps a}{a^2+\eps^2}\big|\le\frac12$ for all real $a$, and $\frac{\eps a}{a^2+\eps^2}\to0$ as $\eps\downarrow0$ for every fixed $a$ (it vanishes for $a=0$ and is at most $\eps/|a|$ otherwise); since $f'$ is integrable, the second term
tends to $0$ by dominated convergence. This proves \eqref{eq:HS}.
\end{proof}

\begin{lemma}[Counting function from the Stieltjes transform]\label{lem:count}
Let $\lambda_1\le\dots\le\lambda_N$ be real, $m(z):=\frac1N\sum_j(\lambda_j-z)^{-1}$, $\fn(E):=\frac1N\#\{j:\lambda_j\le E\}$. Let $E_1<E_2$ with $E_2-E_1\le8$, $0<\tilde\eta\le\frac12$, $U\ge1$, and assume
\begin{equation}\label{eq:HShyp}
 |m(x+\ii y)-\msc(x+\ii y)|\le\frac{U}{Ny}\qquad\text{for all }x\in[E_1-\tilde\eta,E_2+\tilde\eta],\ y\in[\tilde\eta,2].
\end{equation}
Then, with the absolute constant $C_6:=10$,
\[
 \Big|\big(\fn(E_2)-\fn(E_1)\big)-\big(\nsc(E_2)-\nsc(E_1)\big)\Big|\le\frac{C_6}{N}\Big(U\log(2/\tilde\eta)+N\tilde\eta\Big).
\]
\end{lemma}
\begin{proof}
\emph{Step 0: test functions.} Let $\chi(y):=\mathsf s(2-|y|)$. Then $\chi$ is even, $\chi=1$ on $[-1,1]$, $\chi=0$ outside $[-2,2]$, $0\le\chi\le1$, and $\chi\in C^1(\R)$ (near $y=0$ it is constant); for $y>0$, $\chi'(y)=-\mathsf s'(2-y)\le0$, so
\begin{equation}\tag{T$\chi$}\label{eq:a1-chi}
 \int_1^2|\chi'(y)|\,\rd y=1,\qquad \chi'(y)=0\ \text{ for }0<y<1\text{ and for }y>2 .
\end{equation}
Let
\[
 f(x):=\mathsf s\Big(\frac{x-E_1+\tilde\eta}{\tilde\eta}\Big)\,\mathsf s\Big(\frac{E_2+\tilde\eta-x}{\tilde\eta}\Big).
\]
The second factor equals $1$ for $x\le E_2$ and the first equals $1$ for $x\ge E_1$; since $E_1<E_2$, $f$ coincides with the first factor on $(-\infty,E_2]$ and with the second on $[E_1,\infty)$. Hence $f\in C^2(\R)$ is real, $0\le f\le1$,
$f=1$ on $[E_1,E_2]$, $f=0$ outside $(E_1-\tilde\eta,E_2+\tilde\eta)$, $f'$ and $f''$ vanish outside $I_1\cup I_2$ with $I_1:=[E_1-\tilde\eta,E_1]$, $I_2:=[E_2,E_2+\tilde\eta]$, and by \eqref{eq:a1-quintic} (substituting $t=(x-E_1+\tilde\eta)/\tilde\eta$, resp.\ $t=(E_2+\tilde\eta-x)/\tilde\eta$)
\begin{equation}\tag{T$f$}\label{eq:a1-fbounds}
 \int_\R|f|\le E_2-E_1+2\tilde\eta\le9,\qquad\int_\R|f'|=2,\qquad\int_\R|f''|=\frac{2}{\tilde\eta}\cdot\frac{15}4=\frac{15}{2\tilde\eta}.
\end{equation}

\emph{Step 1: a formula for $D:=\frac1N\sum_jf(\lambda_j)-\int_\R f\rhosc$.} For $z\in\C\setminus\R$ let $m(z)$ be defined by the same formula as above and let $\mathfrak s(z):=\int_{-2}^2\rhosc(\lambda)(\lambda-z)^{-1}\rd\lambda$. Then $m(\bar z)=\ol{m(z)}$, $\mathfrak s(\bar z)=\ol{\mathfrak s(z)}$,
and $\mathfrak s=\msc$ on $\C_+$ by Lemma~\ref{lem:stieltjes}. Put $\Delta(z):=m(z)-\mathfrak s(z)$; thus $\Delta=m-\msc$ on $\C_+$ and $\Delta(\bar z)=\ol{\Delta(z)}$.
Applying \eqref{eq:HS} with $\lambda=\lambda_j$ and averaging over $j$,
\[
 \frac1N\sum_jf(\lambda_j)=\frac1{2\pi}\int_\R\int_\R g(x,y)\,m(x+\ii y)\,\rd x\,\rd y .
\]
Multiplying \eqref{eq:HS} by $\rhosc(\lambda)$ and integrating over $\lambda$, Fubini's theorem (justified by \eqref{eq:a1-HSdom}: the integrand is dominated by $\Gamma(x,y)\rhosc(\lambda)$, which is integrable on $\R^3$) gives
$\int f\rhosc=\frac1{2\pi}\int\!\!\int g(x,y)\mathfrak s(x+\ii y)\,\rd x\,\rd y$. Hence
\[
 D=\frac1{2\pi}\int_\R\int_\R g(x,y)\,\Delta(x+\ii y)\,\rd x\,\rd y .
\]
By Lemma~\ref{lem:HS}, $g(x,-y)\Delta(x-\ii y)=\ol{g(x,y)\Delta(x+\ii y)}$, so the contributions of $y<0$ and $y>0$ are complex conjugates and $D=\frac1\pi\int_{y>0}\int\re\big[g\,\Delta\big]\rd x\,\rd y$.
Since $g=\ii yf''\chi+\ii f\chi'-yf'\chi'$ with $f,\chi$ real, $\re[g\Delta]=-yf''\chi\im\Delta-f\chi'\im\Delta-yf'\chi'\re\Delta$. By \eqref{eq:a1-chi} we obtain $D=T_1+T_2+T_3$ with
\begin{align*}
 T_1&:=-\frac1\pi\int_1^2\!\!\int_\R\chi'(y)\Big(f(x)\im\Delta(x+\ii y)+yf'(x)\re\Delta(x+\ii y)\Big)\rd x\,\rd y,\\
 T_2&:=-\frac1\pi\int_0^{\tilde\eta}\!\!\int_\R yf''(x)\im\Delta(x+\ii y)\,\rd x\,\rd y,\qquad
 T_3:=-\frac1\pi\int_{\tilde\eta}^2\!\!\int_\R yf''(x)\chi(y)\im\Delta(x+\ii y)\,\rd x\,\rd y
\end{align*}
(recall $\chi=1$ on $[0,\tilde\eta]\subset[0,1]$). In all three terms $x$ ranges over $\operatorname{supp}f\subset[E_1-\tilde\eta,E_2+\tilde\eta]$, so the hypothesis \eqref{eq:HShyp} is available whenever $\tilde\eta\le y\le2$.

\emph{Step 2: $T_1$.} For $1\le y\le2$, \eqref{eq:HShyp} gives $|\Delta|\le U/(Ny)\le U/N$ and $y|\Delta|\le U/N$. By \eqref{eq:a1-chi}, \eqref{eq:a1-fbounds},
\[
 |T_1|\le\frac1\pi\int_1^2|\chi'(y)|\rd y\cdot\frac UN\Big(\int|f|+\int|f'|\Big)\le\frac{11}{\pi}\cdot\frac UN .
\]

\emph{Step 3: $T_2$.} For $y>0$ and real $x$, $\im m(x+\ii y)=\frac1N\sum_j\frac{y}{(\lambda_j-x)^2+y^2}>0$, and $y\,\im m(x+\ii y)=\frac1N\sum_j\frac{y^2}{(\lambda_j-x)^2+y^2}$ is nondecreasing in $y$. Also $0<\im\msc\le|\msc|\le1$ (Lemma~\ref{lem:msc}).
Hence for $0<y\le\tilde\eta$ and $x\in[E_1-\tilde\eta,E_2+\tilde\eta]$, using \eqref{eq:HShyp} at $x+\ii\tilde\eta$,
\[
 y|\im\Delta(x+\ii y)|\le y\,\im m(x+\ii y)+y\,\im\msc(x+\ii y)\le\tilde\eta\,\im m(x+\ii\tilde\eta)+\tilde\eta\le\tilde\eta\Big(1+\frac{U}{N\tilde\eta}\Big)+\tilde\eta=2\tilde\eta+\frac UN .
\]
By \eqref{eq:a1-fbounds},
\[
 |T_2|\le\frac1\pi\cdot\tilde\eta\cdot\frac{15}{2\tilde\eta}\Big(2\tilde\eta+\frac UN\Big)=\frac{15}{\pi}\tilde\eta+\frac{15}{2\pi}\cdot\frac UN .
\]

\emph{Step 4: $T_3$.} The function $\Delta$ is holomorphic on $\C_+$: $m$ is a rational function with real poles and $\msc$ is holomorphic by Lemma~\ref{lem:msc}(i). Therefore $\partial_x\Delta=\Delta'$ and $\partial_y\Delta=\ii\Delta'$ as functions of $(x,y)$, $y>0$, and
\begin{equation}\tag{CR}\label{eq:a1-CR}
 \partial_x\im\Delta=\im\Delta'=-\re(\ii\Delta')=-\partial_y\re\Delta .
\end{equation}
Fix $y\in[\tilde\eta,2]$. Integrating by parts in $x$ ($f'$ has compact support and $x\mapsto\im\Delta(x+\ii y)$ is $C^1$) and using \eqref{eq:a1-CR},
\[
 \int_\R f''(x)\im\Delta(x+\ii y)\,\rd x=-\int_\R f'(x)\,\partial_x\im\Delta(x+\ii y)\,\rd x=\int_\R f'(x)\,\partial_y\re\Delta(x+\ii y)\,\rd x .
\]
Insert this into $T_3$ and exchange the order of integration (the integrand $yf'(x)\chi(y)\partial_y\re\Delta(x+\ii y)$ is continuous on the compact set $\operatorname{supp}f'\times[\tilde\eta,2]$). For fixed $x$, integration by parts in $y$, with $\chi(2)=0$ and $\chi(\tilde\eta)=1$, gives
\[
 \int_{\tilde\eta}^2y\chi(y)\,\partial_y\re\Delta(x+\ii y)\,\rd y=-\tilde\eta\,\re\Delta(x+\ii\tilde\eta)-\int_{\tilde\eta}^2\big(\chi(y)+y\chi'(y)\big)\re\Delta(x+\ii y)\,\rd y .
\]
Hence
\[
 T_3=\frac1\pi\int_\R f'(x)\Big[\tilde\eta\,\re\Delta(x+\ii\tilde\eta)+\int_{\tilde\eta}^2\big(\chi(y)+y\chi'(y)\big)\re\Delta(x+\ii y)\,\rd y\Big]\rd x .
\]
For $x\in\operatorname{supp}f'\subset I_1\cup I_2$ and $y\in[\tilde\eta,2]$, \eqref{eq:HShyp} gives $|\re\Delta(x+\ii y)|\le U/(Ny)$. So $|\tilde\eta\re\Delta(x+\ii\tilde\eta)|\le U/N$, $\int_{\tilde\eta}^2\chi(y)\frac{U}{Ny}\rd y\le\frac UN\log(2/\tilde\eta)$ and, by \eqref{eq:a1-chi},
$\int_{\tilde\eta}^2y|\chi'(y)|\frac U{Ny}\rd y=\frac UN$. With $\int|f'|=2$,
\[
 |T_3|\le\frac2\pi\cdot\frac UN\big(2+\log(2/\tilde\eta)\big).
\]

\emph{Step 5: the smooth statistic.} Adding Steps 2--4 and using $\log(2/\tilde\eta)\ge\log4>1.38$:
\[
 |D|\le\frac UN\Big(\frac{11}\pi+\frac{15}{2\pi}+\frac4\pi+\frac2\pi\log(2/\tilde\eta)\Big)+\frac{15}\pi\tilde\eta\le5.9\,\frac UN\log(2/\tilde\eta)+4.8\,\tilde\eta ,
\]
since $\frac{22.5}{\pi}<7.17<5.2\cdot1.38$ and $5.2+\frac2\pi<5.9$.

\emph{Step 6: removing the smoothing.} Since $\fn(E_2)-\fn(E_1)=\frac1N\#\{j:E_1<\lambda_j\le E_2\}$ and $f=1$ on $[E_1,E_2]$, $0\le f\le1$, $f=0$ outside $I_1\cup[E_1,E_2]\cup I_2$,
\[
 \Big|\frac1N\sum_jf(\lambda_j)-\big(\fn(E_2)-\fn(E_1)\big)\Big|\le\frac1N\#\{j:\lambda_j\in I_1\}+\frac1N\#\{j:\lambda_j\in I_2\}.
\]
For $E\in\{E_1,E_2\}$: every $j$ with $|\lambda_j-E|\le\tilde\eta$ contributes at least $\frac1{2\tilde\eta}$ to $N\im m(E+\ii\tilde\eta)=\sum_j\frac{\tilde\eta}{(\lambda_j-E)^2+\tilde\eta^2}$, so by \eqref{eq:HShyp} at $E+\ii\tilde\eta$
\[
 \#\{j:|\lambda_j-E|\le\tilde\eta\}\le2N\tilde\eta\,\im m(E+\ii\tilde\eta)\le2N\tilde\eta\Big(1+\frac{U}{N\tilde\eta}\Big)=2N\tilde\eta+2U .
\]
As $I_1\subset[E_1-\tilde\eta,E_1+\tilde\eta]$ and $I_2\subset[E_2-\tilde\eta,E_2+\tilde\eta]$, the right side of the previous display is at most $4\tilde\eta+4U/N$.
Similarly, by $0\le\rhosc\le\frac1\pi$ (Lemma~\ref{lem:nsc}(a)),
\[
 \Big|\int f\rhosc-\big(\nsc(E_2)-\nsc(E_1)\big)\Big|=\int_{I_1\cup I_2}f\rhosc\le\frac{2\tilde\eta}\pi .
\]
Altogether, using $4\le2.9\log(2/\tilde\eta)$,
\begin{align*}
 \Big|\big(\fn(E_2)-\fn(E_1)\big)-\big(\nsc(E_2)-\nsc(E_1)\big)\Big|&\le|D|+4\tilde\eta+\frac{4U}N+\frac{2\tilde\eta}\pi\le8.8\,\frac UN\log(2/\tilde\eta)+9.5\,\tilde\eta\\
 &\le\frac{10}N\Big(U\log(2/\tilde\eta)+N\tilde\eta\Big).\qedhere
\end{align*}
\end{proof}
\begin{remark}\label{rem:count}
Lemma~\ref{lem:count} is deterministic. It replaces Lemma~5.1 of the original (quoted there from another paper), whose hypothesis did not match its use; here the hypothesis is exactly what Theorem~\ref{thm:K}(i) provides.
The lemma is applied in Section~\ref{sec:rig} with $E_1=-3$, to the left of the spectrum.
\end{remark}

\section{Probabilistic toolbox}\label{sec:prob}
Throughout Sections~\ref{sec:prob} and~\ref{sec:weak} the model parameters are normalised as in Convention~\ref{conv:const}: $0<\vth\le1$, $0<\dpm\le1$, $C_0\ge1$, $\fa=1/\vth\ge1$.
All random variables are complex valued unless said otherwise, $\log$ is the natural logarithm, and $\bone(\cdot)$ denotes an indicator function.
We use two named constants that depend only on $\vth$:
\[
 C_\vth:=90\,\fa,\qquad\qquad\kappa_\vth:=\log\fa+4\fa\log(4C_\vth)\ \ (\ge23).
\]
(The bound $\kappa_\vth\ge23$ holds because $\fa\ge1$ and $4\log360>23$.)
\subsection{Moments and truncation}
\begin{lemma}[Moments]\label{lem:moments}
Let $0<\vth\le1$, let $\mu$ be a Borel probability measure on $\C$ and $s\ge0$ with $\mu(\{y:|y|>xs\})\le\vth^{-1}e^{-x^\vth}$ for all $x\ge1$. Then:
\begin{enumerate}[label=\textup{(\alph*)},leftmargin=2.2em,itemsep=1pt]
\item for every real $q\ge1$:\quad $\int|y|^q\,\mu(\rd y)\le(C_\vth q)^{q/\vth}s^q$, with $C_\vth=90/\vth$; if $s=0$ then $\mu=\delta_0$;
\item \emph{(truncation)} for every $T\ge1$:
\begin{align*}
 \int|y|\,\bone(|y|>Ts)\,\mu(\rd y)&\le\Big(\int|y|^2\rd\mu\Big)^{1/2}\vth^{-1/2}e^{-T^\vth/2},\\
 \int|y|^2\,\bone(|y|>Ts)\,\mu(\rd y)&\le(4C_\vth)^{2\fa}\,s^2\,\vth^{-1/2}e^{-T^\vth/2}.
\end{align*}
\end{enumerate}
In particular, for $\nu\in\mathcal E_N(\vth,\dpm,C_0)$ and all $i,j$: $\EE|h_{ij}|^q\le(C_\vth q)^{\fa q}\sigma_{ij}^q$ for every real $q\ge1$, and for $T\ge1$,
\[
 \PP\big(\exists\,i,j:\ |h_{ij}|>T\sigma_{ij}\big)\le\PP\big(\exists\,(i,j)\in\mathcal P_N:\ |x_{ij}|>T\sigma_{ij}\big)\le N^2\vth^{-1}e^{-T^\vth}.
\]
\end{lemma}
\begin{proof}
(a) \emph{The case $s=0$.} The hypothesis gives $\mu(|y|>0)\le\vth^{-1}e^{-x^\vth}$ for every $x\ge1$; letting $x\to\infty$, $\mu(|y|>0)=0$, i.e.\ $\mu=\delta_0$, and both sides of the moment bound vanish.

\emph{The case $s>0$.} By the layer-cake formula (Mathlib: \texttt{lintegral\_rpow\_eq\_lintegral\_meas\_lt\_mul}) applied to the function $y\mapsto|y|/s$, and since $\mu(\cdot)\le1$,
\[
 \int\Big(\frac{|y|}s\Big)^q\mu(\rd y)=\int_0^\infty qx^{q-1}\mu(|y|>xs)\,\rd x\le\int_0^1qx^{q-1}\rd x+\frac q\vth\int_1^\infty x^{q-1}e^{-x^\vth}\rd x .
\]
The first integral is $1$. In the second we substitute $x=u^{\fa}$ (so $x^\vth=u$, $\rd x=\fa u^{\fa-1}\rd u$) and put $t:=\fa q\ge1$:
\begin{align*}
 \frac q\vth\int_1^\infty x^{q-1}e^{-x^\vth}\rd x&=q\fa^2\int_1^\infty u^{t-1}e^{-u}\rd u\le q\fa^2\int_1^\infty u^{t}e^{-u/2}\,e^{-u/2}\rd u\\
 &\le q\fa^2(2t)^t\int_0^\infty e^{-u/2}\rd u=2q\fa^2(2t)^t .
\end{align*}
Here we used $u^{t-1}\le u^t$ for $u\ge1$ and $\sup_{u>0}u^te^{-u/2}=(2t/e)^t\le(2t)^t$ (the function $u\mapsto t\log u-u/2$ is maximal at $u=2t$).
Now $q\le2^q\le2^{t}$ and $\fa^2\le e^{2\fa}\le e^{2t}$, so with $X:=(4e^2t)^t\ge1$,
\[
 \int\Big(\frac{|y|}s\Big)^q\rd\mu\le1+2\cdot2^te^{2t}(2t)^t=1+2X\le3X\le3^tX=(12e^2\,\fa q)^{\fa q}\le(C_\vth q)^{\fa q},
\]
because $12e^2<90$.

(b) Let $R:=\{|y|>Ts\}$; by hypothesis (with $x=T\ge1$) $\mu(R)\le\vth^{-1}e^{-T^\vth}$. By the Cauchy--Schwarz inequality,
$\int|y|\bone_R\rd\mu\le(\int|y|^2\rd\mu)^{1/2}\mu(R)^{1/2}$ and $\int|y|^2\bone_R\rd\mu\le(\int|y|^4\rd\mu)^{1/2}\mu(R)^{1/2}$, and $\int|y|^4\rd\mu\le(4C_\vth)^{4\fa}s^4$ by (a).

\emph{The ensemble.} Let $(i,j)\in\mathcal P_N$. By Definition~\ref{def:ass}, the measure $\nu_{ij}$ satisfies the hypothesis with $s=\sigma_{ij}$ (this is \eqref{eq:subexp}), and by Definition~\ref{def:ens}
$|h_{ij}|=|h_{ji}|\le|x_{ij}|$ at every point of $\Omega$ (for $i<j$ with equality; for $i=j$, $h_{ii}=\re x_{ii}$), where $x_{ij}$ has law $\nu_{ij}$. So (a) gives the moment bound, and
$\{\exists i,j:|h_{ij}|>T\sigma_{ij}\}\subset\bigcup_{(i,j)\in\mathcal P_N}\{|x_{ij}|>T\sigma_{ij}\}$, a union of $N(N+1)/2\le N^2$ events of probability at most $\vth^{-1}e^{-T^\vth}$.
\end{proof}
\begin{remark}\label{rem:laws}
For later use we record what Definitions~\ref{def:ens} and~\ref{def:ass} say about a single entry. Let $i\ne j$. The random variable $h_{ij}$ is the coordinate $x_{\min(i,j),\max(i,j)}$ or its complex conjugate. Hence
$\EE h_{ij}=0$, $\EE|h_{ij}|^2=\sigma_{ij}^2$ (by \eqref{eq:aver}), and $\PP(|h_{ij}|>x\sigma_{ij})\le\vth^{-1}e^{-x^\vth}$ for $x\ge1$. If $\sigma_{ij}=0$ then $h_{ij}=0$ almost surely.
For $i=j$, $h_{ii}=x_{ii}$ almost surely by \textup{(R)}, and the same three properties hold.
\end{remark}
\subsection{Large deviation bounds}
Throughout this subsection $X_1,\dots,X_n$ ($n\ge1$) are independent complex random variables, defined on some probability space, with
\begin{equation}\label{eq:LDhyp}
 \EE X_i=0,\qquad\EE|X_i|^2=s_i^2\ (s_i\ge0),\qquad\PP(|X_i|>xs_i)\le\vth^{-1}e^{-x^\vth}\ \ (x\ge1),
\end{equation}
with $0<\vth\le1$, $\fa=1/\vth$. (If $s_i=0$ then $X_i=0$ almost surely, by Lemma~\ref{lem:moments}(a). If $X_i$ satisfies \eqref{eq:LDhyp}, so does $\ol{X_i}$.) For $\kappa\ge1$ define the \emph{loss factors}
\begin{equation}\label{eq:loss}
 \ell(\kappa):=5\,\kappa^{\fa+\frac12},\qquad\ell_2(\kappa,n):=2\,\ell(\kappa)^2\sqrt{2\lceil\log_2n\rceil+2}.
\end{equation}
The only probabilistic input is Hoeffding's inequality, in the following form.

\medskip\noindent\textbf{(H)} \emph{Let $\xi_1,\dots,\xi_n$ be independent real random variables with $\EE\xi_i=0$ and $|\xi_i|\le c_i$ almost surely, where $c_i\ge0$ are constants, and let $V:=\sum_ic_i^2$.
Then for every $\kappa\ge0$:\quad $\PP\big(|\sum_i\xi_i|>\sqrt{2\kappa V}\big)\le2e^{-\kappa}$.}

\smallskip\noindent Indeed, if $V=0$ then every $\xi_i$ vanishes almost surely and the event has probability $0$. If $V>0$: by Hoeffding's lemma a centred random variable with values in $[-c_i,c_i]$ has a sub-Gaussian moment generating function with
variance proxy $c_i^2$; a sum of independent such variables has variance proxy $V$; and the Chernoff bound gives $\PP(\sum_i\xi_i\ge t)\le\exp(-t^2/(2V))$ for $t\ge0$
(Mathlib: \texttt{hasSubgaussianMGF\_of\_mem\_Icc\_of\_integral\_eq\_zero}, \texttt{HasSubgaussianMGF.measure\_sum\_ge\_le\_of\_iIndepFun}). Apply this to $(\xi_i)$ and to $(-\xi_i)$ with $t=\sqrt{2\kappa V}$.

\begin{lemma}[Large deviations, tail form]\label{lem:LD} \infra\
Let $0<\vth\le1$ and $\kappa_\vth=\log\fa+4\fa\log(4C_\vth)$. For all $n\ge1$, all $X$ as in \eqref{eq:LDhyp}, all deterministic $a_i,b_i,B_{ij}\in\C$ and all $\kappa\ge\kappa_\vth+\log n$:
\begin{align}
 &\PP\Big(\Big|\sum_ia_iX_i\Big|>\ell(\kappa)\Big(\sum_i|a_i|^2s_i^2\Big)^{1/2}\Big)\le(4+n\vth^{-1})e^{-\kappa},\label{eq:LD1}\\
 &\PP\Big(\Big|\sum_ib_i(|X_i|^2-s_i^2)\Big|>\ell(\kappa)^2\Big(\sum_i|b_i|^2s_i^4\Big)^{1/2}\Big)\le(4+n\vth^{-1})e^{-\kappa},\label{eq:LD2}\\
 &\PP\Big(\Big|\sum_{i\ne j}\ol{X_i}B_{ij}X_j\Big|>\tfrac12\ell_2(\kappa,n)\Big(\sum_{i\ne j}s_i^2|B_{ij}|^2s_j^2\Big)^{1/2}\Big)\le20\,n^2\lceil\log_2 2n\rceil\vth^{-1}e^{-\kappa}.\label{eq:LD3}
\end{align}
If $Y_1,\dots,Y_{n'}$, $1\le n'\le n$, are random variables on the same probability space such that the $n+n'$ variables $X_1,\dots,X_n,Y_1,\dots,Y_{n'}$ are independent and
$Y$ satisfies \eqref{eq:LDhyp} with variances $t_j^2$, then for all $B_{ij}\in\C$ $(1\le i\le n,\ 1\le j\le n')$
\begin{equation}\label{eq:LD4}
 \PP\Big(\Big|\sum_{i,j}X_iB_{ij}Y_j\Big|>\ell(\kappa)^2\Big(\sum_{i,j}s_i^2|B_{ij}|^2t_j^2\Big)^{1/2}\Big)\le10\,n^2\vth^{-1}e^{-\kappa}.
\end{equation}
The same bounds hold with $X_i$ replaced by $\ol{X_i}$ in any position (i.e.\ for $\sum_ia_i\ol{X_i}$, for $\sum_{i\ne j}X_iB_{ij}\ol{X_j}$, and for $\sum_{i,j}\ol{X_i}B_{ij}Y_j$, $\sum_{i,j}X_iB_{ij}\ol{Y_j}$, $\sum_{i,j}\ol{X_i}B_{ij}\ol{Y_j}$).
All inequalities inside the probabilities are strict; in particular the bounds are meaningful (and true) when some or all $s_i$ vanish.
\end{lemma}
\begin{proof}
Fix $\kappa\ge\kappa_\vth+\log n$ and put $T:=\kappa^{\fa}$, so that $T\ge1$ and $T^\vth=\kappa$.

\emph{Step 1: truncation and recentring of one variable.} Let
\[
 \mathcal T:=\{|X_i|\le Ts_i\ \text{for all }i\},\qquad\PP(\mathcal T^c)\le\sum_i\PP(|X_i|>Ts_i)\le n\vth^{-1}e^{-\kappa}
\]
by \eqref{eq:LDhyp} with $x=T$. For each $i$ define
\[
 X_i^\le:=X_i\bone(|X_i|\le Ts_i),\qquad\mu_i:=\EE X_i^\le,\qquad Y_i:=X_i^\le-\mu_i,
\]
\[
 W_i^\le:=|X_i|^2\bone(|X_i|\le Ts_i),\qquad\rho_i:=\EE\big[|X_i|^2\bone(|X_i|>Ts_i)\big],\qquad V_i:=W_i^\le-\EE W_i^\le=W_i^\le-s_i^2+\rho_i .
\]
Since $\EE X_i=0$ we have $\mu_i=-\EE[X_i\bone(|X_i|>Ts_i)]$, so by Lemma~\ref{lem:moments}(b) (applied to the law of $X_i$, with $s=s_i$ and $\int|y|^2\rd\mu=s_i^2$)
\begin{equation}\tag{T1}\label{eq:T1}
 |\mu_i|\le s_i\,\vth^{-1/2}e^{-\kappa/2},\qquad0\le\rho_i\le(4C_\vth)^{2\fa}s_i^2\,\vth^{-1/2}e^{-\kappa/2}.
\end{equation}
Moreover, at every point, $|X_i^\le|\le Ts_i$, hence $|\mu_i|\le Ts_i$ and
\begin{equation}\tag{T2}\label{eq:T2}
 \EE Y_i=0,\quad|Y_i|\le2Ts_i;\qquad\EE V_i=0,\quad|V_i|\le T^2s_i^2 .
\end{equation}
(For the last bound: $W_i^\le$ and its expectation both lie in $[0,T^2s_i^2]$.) The families $(Y_i)_i$ and $(V_i)_i$ are independent, because $Y_i$ and $V_i$ are Borel functions of $X_i$. Finally, on the event $\mathcal T$,
\begin{equation}\tag{T3}\label{eq:T3}
 X_i=Y_i+\mu_i\qquad\text{and}\qquad|X_i|^2-s_i^2=V_i-\rho_i\qquad\text{for all }i .
\end{equation}

\emph{Step 2: proof of \eqref{eq:LD1}.} Let $S:=(\sum_i|a_i|^2s_i^2)^{1/2}$. The real random variables $\xi_i:=\re(a_iY_i)$ are independent, centred and bounded by $c_i:=2T|a_i|s_i$ by \eqref{eq:T2}, and $\sum_ic_i^2=4T^2S^2$.
By (H), $|\sum_i\re(a_iY_i)|\le\sqrt{8\kappa}\,TS$ outside an event of probability at most $2e^{-\kappa}$; the same holds for $\im(a_iY_i)$. Hence there is an event $\mathcal H_1$ with $\PP(\mathcal H_1^c)\le4e^{-\kappa}$ on which
\[
 \Big|\sum_ia_iY_i\Big|\le\sqrt{2\cdot8\kappa}\,TS=4T\sqrt\kappa\,S .
\]
By \eqref{eq:T1} and the Cauchy--Schwarz inequality, $|\sum_ia_i\mu_i|\le\sum_i|a_i|s_i\vth^{-1/2}e^{-\kappa/2}\le\sqrt n\,S\,\vth^{-1/2}e^{-\kappa/2}\le S$, because $\kappa\ge\log n+\log\fa$.
Therefore, by \eqref{eq:T3}, on $\mathcal T\cap\mathcal H_1$,
\[
 \Big|\sum_ia_iX_i\Big|\le(4T\sqrt\kappa+1)S\le5\kappa^{\fa+\frac12}S=\ell(\kappa)S,
\]
and $\PP((\mathcal T\cap\mathcal H_1)^c)\le(4+n\vth^{-1})e^{-\kappa}$. This is \eqref{eq:LD1}.

\emph{Step 3: proof of \eqref{eq:LD2}.} Let $S_2:=(\sum_i|b_i|^2s_i^4)^{1/2}$. The real variables $\re(b_iV_i)=\re(b_i)V_i$ are independent, centred and bounded by $|b_i|T^2s_i^2$ by \eqref{eq:T2}; the sum of the squares of these bounds is $T^4S_2^2$.
By (H) applied to the real and to the imaginary parts there is an event $\mathcal H_2$ with $\PP(\mathcal H_2^c)\le4e^{-\kappa}$ on which $|\sum_ib_iV_i|\le\sqrt{2\cdot2\kappa}\,T^2S_2=2T^2\sqrt\kappa\,S_2$.
By \eqref{eq:T1}, $|\sum_ib_i\rho_i|\le\sqrt n\,S_2\,(4C_\vth)^{2\fa}\vth^{-1/2}e^{-\kappa/2}\le S_2$, because $\kappa\ge\log n+\log\fa+4\fa\log(4C_\vth)=\log n+\kappa_\vth$. So by \eqref{eq:T3}, on $\mathcal T\cap\mathcal H_2$,
\[
 \Big|\sum_ib_i(|X_i|^2-s_i^2)\Big|\le(2T^2\sqrt\kappa+1)S_2\le3\kappa^{2\fa+\frac12}S_2\le\ell(\kappa)^2S_2 .
\]

\emph{Step 4: proof of \eqref{eq:LD4}.} Let $\mathsf P_X$ and $\mathsf P_Y$ be the laws of the random vectors $X=(X_1,\dots,X_n)\in\C^n$ and $Y=(Y_1,\dots,Y_{n'})\in\C^{n'}$.
Since the $n+n'$ variables are independent, the law of $(X,Y)$ on $\C^n\times\C^{n'}$ is the product measure $\mathsf P_X\otimes\mathsf P_Y$, and under $\mathsf P_X$ the coordinate functions $x\mapsto x_i$ are independent and have the laws of the $X_i$; so
they satisfy \eqref{eq:LDhyp} (Mathlib: \texttt{iIndepFun\_iff\_map\_fun\_eq\_pi\_map}, \texttt{iIndepFun.indepFun\_finset}). For $y\in\C^{n'}$ put
\[
 c_i(y):=\sum_jB_{ij}y_j,\qquad R(y):=\Big(\sum_is_i^2|c_i(y)|^2\Big)^{1/2},\qquad r_i:=\Big(\sum_j|B_{ij}|^2t_j^2\Big)^{1/2},
\]
and consider the Borel sets
\begin{align*}
 \mathcal E_1&:=\Big\{(x,y)\in\C^n\times\C^{n'}:\ \Big|\sum_ix_ic_i(y)\Big|>\ell(\kappa)R(y)\Big\},\\
 \mathcal E_2&:=\big\{y\in\C^{n'}:\ |c_i(y)|>\ell(\kappa)r_i\text{ for some }i\big\}.
\end{align*}
For every fixed $y$, \eqref{eq:LD1} with the deterministic coefficients $a_i=c_i(y)$ gives $\mathsf P_X(\{x:(x,y)\in\mathcal E_1\})\le(4+n\vth^{-1})e^{-\kappa}$. By Tonelli's theorem for the indicator function of $\mathcal E_1$,
\[
 \PP\big((X,Y)\in\mathcal E_1\big)=\int_{\C^{n'}}\mathsf P_X\big(\{x:(x,y)\in\mathcal E_1\}\big)\,\mathsf P_Y(\rd y)\le(4+n\vth^{-1})e^{-\kappa}.
\]
By \eqref{eq:LD1} for the family $Y$ ($n'$ variables; $\kappa\ge\kappa_\vth+\log n'$ since $n'\le n$) and the coefficients $(B_{ij})_j$, and a union bound over $i$: $\PP(Y\in\mathcal E_2)\le n(4+n'\vth^{-1})e^{-\kappa}$.
If $(X,Y)\notin\mathcal E_1$ and $Y\notin\mathcal E_2$, then
\[
 \Big|\sum_{i,j}X_iB_{ij}Y_j\Big|=\Big|\sum_iX_ic_i(Y)\Big|\le\ell(\kappa)R(Y)\le\ell(\kappa)\Big(\sum_is_i^2\ell(\kappa)^2r_i^2\Big)^{1/2}=\ell(\kappa)^2\Big(\sum_{i,j}s_i^2|B_{ij}|^2t_j^2\Big)^{1/2}.
\]
The probability of the exceptional event is at most $\big(4+n\vth^{-1}+4n+nn'\vth^{-1}\big)e^{-\kappa}\le10n^2\vth^{-1}e^{-\kappa}$, since $n'\le n$, $n\ge1$ and $\vth^{-1}\ge1$.

\emph{Step 5: proof of \eqref{eq:LD3} (deterministic decoupling).} If $n=1$ the sum is empty and the event is empty. Let $n\ge2$ and $D:=\lceil\log_2n\rceil\ge1$. Relabel the indices as $0,1,\dots,n-1$. Since $n\le2^D$, every index has a unique binary expansion
\[
 i=\sum_{d=1}^D\epsilon_d(i)\,2^{d-1},\qquad\epsilon_d(i)\in\{0,1\}.
\]
For $i\ne j$ let $d(i,j):=\max\{d:\epsilon_d(i)\ne\epsilon_d(j)\}$ be the position of the most significant digit in which $i$ and $j$ differ (the set is not empty because the expansion is unique).
For $d\in\{1,\dots,D\}$ and $\beta\in\{0,1\}$ put
\[
 I_{d,\beta}:=\{i:\epsilon_d(i)=\beta\},\qquad\mathcal P_{d,\beta}:=\{(i,j):\ i\ne j,\ d(i,j)=d,\ \epsilon_d(i)=\beta\}.
\]
Then: (i) the $2D$ sets $\mathcal P_{d,\beta}$ are pairwise disjoint and their union is $\{(i,j):i\ne j\}$, because a pair $(i,j)$ with $i\ne j$ belongs to $\mathcal P_{d,\beta}$ if and only if $d=d(i,j)$ and $\beta=\epsilon_d(i)$;
(ii) if $(i,j)\in\mathcal P_{d,\beta}$ then $i\in I_{d,\beta}$ and $j\in I_{d,1-\beta}$, because $\epsilon_d(j)\ne\epsilon_d(i)=\beta$; (iii) $I_{d,0}\cap I_{d,1}=\emptyset$. Consequently
\begin{equation}\tag{D}\label{eq:decoup}
 \sum_{i\ne j}\ol{X_i}B_{ij}X_j=\sum_{d=1}^D\sum_{\beta\in\{0,1\}}F_{d,\beta},\qquad F_{d,\beta}:=\sum_{i\in I_{d,\beta}}\ \sum_{j\in I_{d,1-\beta}}\ol{X_i}\,B^{(d,\beta)}_{ij}\,X_j,
\end{equation}
where $B^{(d,\beta)}_{ij}:=B_{ij}\bone\big((i,j)\in\mathcal P_{d,\beta}\big)$.
Fix $(d,\beta)$ and let $n_1:=|I_{d,\beta}|$, $n_2:=|I_{d,1-\beta}|$, so $n_1+n_2=n$. If $n_1=0$ or $n_2=0$ then $F_{d,\beta}=0$. Otherwise the families $(\ol{X_i})_{i\in I_{d,\beta}}$ and $(X_j)_{j\in I_{d,1-\beta}}$ are indexed by disjoint sets, so the $n$ variables together are independent,
and each family satisfies \eqref{eq:LDhyp}. We apply \eqref{eq:LD4} with ``$n$'' replaced by $\max(n_1,n_2)$: the larger family plays the role of $X$ and the smaller one the role of $Y$ (the bilinear form $\sum_{i,j}u_iB_{ij}v_j$ is unchanged when the two families are exchanged and the matrix is transposed);
this is allowed because $\kappa\ge\kappa_\vth+\log n\ge\kappa_\vth+\log\max(n_1,n_2)$. So outside an event of probability at most $10\max(n_1,n_2)^2\vth^{-1}e^{-\kappa}\le10n^2\vth^{-1}e^{-\kappa}$,
\[
 |F_{d,\beta}|\le\ell(\kappa)^2\Big(\sum_{(i,j)\in\mathcal P_{d,\beta}}s_i^2|B_{ij}|^2s_j^2\Big)^{1/2}.
\]
By a union bound over the $2D$ pairs $(d,\beta)$, by \eqref{eq:decoup}, by the Cauchy--Schwarz inequality for the sum over $(d,\beta)$ and by the partition property (i), outside an event of probability at most $20Dn^2\vth^{-1}e^{-\kappa}$,
\[
 \Big|\sum_{i\ne j}\ol{X_i}B_{ij}X_j\Big|\le\ell(\kappa)^2\sqrt{2D}\,\Big(\sum_{d,\beta}\ \sum_{(i,j)\in\mathcal P_{d,\beta}}s_i^2|B_{ij}|^2s_j^2\Big)^{1/2}=\ell(\kappa)^2\sqrt{2D}\,\Big(\sum_{i\ne j}s_i^2|B_{ij}|^2s_j^2\Big)^{1/2}.
\]
Since $\ell(\kappa)^2\sqrt{2D}\le\ell(\kappa)^2\sqrt{2D+2}=\frac12\ell_2(\kappa,n)$ and $D\le\lceil\log_22n\rceil$, this is \eqref{eq:LD3}.

\emph{Conjugates.} The family $(\ol{X_i})$ satisfies \eqref{eq:LDhyp} with the same $s_i$, and so does any family obtained by conjugating some of the variables; independence is not affected. So \eqref{eq:LD1}, \eqref{eq:LD4} hold with conjugates in any position, and
\eqref{eq:LD3} for $\sum_{i\ne j}X_iB_{ij}\ol{X_j}$ is \eqref{eq:LD3} for the family $(\ol{X_i})$.
\end{proof}
\begin{remark}\label{rem:LDn}
\textup{(i)} In \eqref{eq:LD1}, \eqref{eq:LD2} the lower bound on $\kappa$ is used only for the recentring; the term $\log n$ compensates the factor $\sqrt n$ lost in $\sum_i|a_i||\mu_i|$.
\textup{(ii)} Any rule that assigns to each pair $i\ne j$ a digit in which $i$ and $j$ differ would do in Step~5; we use the most significant one.
\textup{(iii)} The loss factor $\ell_2(\kappa,n)$ is nondecreasing in $n$ and in $\kappa$, and $\ell(\kappa)^2\le\frac1{2\sqrt2}\ell_2(\kappa,n)$.
\end{remark}

\begin{lemma}[Large deviations, moment form]\label{lem:LDmom} \infra\ 
In the setting of Lemma~\ref{lem:LD}, for all $n\ge1$, all real $q\ge2$ and all deterministic $a_i,B_{ij}\in\C$:
\begin{align}
 \Big\|\sum_ia_iX_i\Big\|_q&\le C_{\rm LD}\,(q+\log n)^{\fa+\frac12}\Big(\sum_i|a_i|^2s_i^2\Big)^{1/2},\label{eq:LDm1}\\
 \Big\|\sum_{i,j}\ol{X_i}B_{ij}X_j-\sum_is_i^2B_{ii}\Big\|_q&\le C_{\rm LD}\,(q+\log n)^{2\fa+\frac32}\Big(\sum_{i,j}s_i^2|B_{ij}|^2s_j^2\Big)^{1/2},\label{eq:LDm2}
\end{align}
where $\norm{F}_q:=(\EE|F|^q)^{1/q}$ and $C_{\rm LD}:=150\,(4\fa+2+\kappa_\vth)^{2\fa+1}$ depends only on $\vth$. The same holds with $X_i$ replaced by $\ol{X_i}$.
\end{lemma}
\begin{proof}
\emph{Step 1: from a tail bound to a moment bound.} We claim: let $F\ge0$ be a random variable, and let $\mathsf c>0$, $\rho>0$, $A\ge0$, $\kappa_0\ge1$ be constants such that
\begin{equation}\tag{M1}\label{eq:M1}
 \PP\big(F>\mathsf c\kappa^{\rho}\big)\le A\,e^{-\kappa}\qquad\text{for all }\kappa\ge\kappa_0 .
\end{equation}
Then for every real $q\ge1$ and every $\kappa_1\ge\max(\kappa_0,2\rho q)$:
\begin{equation}\tag{M2}\label{eq:M2}
 \EE F^q\le\big(\mathsf c\kappa_1^{\rho}\big)^q\big(1+5Ae^{-\kappa_1}\big).
\end{equation}
Indeed, put $t_k:=\mathsf c(\kappa_1+k)^\rho$ for $k=0,1,2,\dots$; this sequence is increasing and unbounded, so at every point
\[
 F^q\le t_0^q\,\bone(F\le t_0)+\sum_{k\ge0}t_{k+1}^q\,\bone(t_k<F\le t_{k+1})\le t_0^q+\sum_{k\ge0}t_{k+1}^q\,\bone(F>t_k).
\]
Taking expectations (monotone convergence) and using \eqref{eq:M1} with $\kappa=\kappa_1+k\ge\kappa_0$,
\[
 \EE F^q\le t_0^q+\sum_{k\ge0}t_{k+1}^q\,Ae^{-\kappa_1-k}.
\]
Since $1+x\le e^x$ and $\rho q/\kappa_1\le\frac12$: $\ t_{k+1}^q=t_0^q\big(1+\frac{k+1}{\kappa_1}\big)^{\rho q}\le t_0^q\exp\big(\frac{\rho q(k+1)}{\kappa_1}\big)\le t_0^qe^{(k+1)/2}$. Hence
\[
 \EE F^q\le t_0^q\Big(1+Ae^{-\kappa_1}e^{1/2}\sum_{k\ge0}e^{-k/2}\Big)=t_0^q\Big(1+\frac{e^{1/2}}{1-e^{-1/2}}Ae^{-\kappa_1}\Big)\le t_0^q\big(1+5Ae^{-\kappa_1}\big),
\]
because $e^{1/2}/(1-e^{-1/2})<4.2$. (This discrete summation replaces the layer-cake formula and needs no change of variables.)

\emph{Step 2: proof of \eqref{eq:LDm1}.} Let $S:=(\sum_i|a_i|^2s_i^2)^{1/2}$ and $r:=\fa+\frac12$. If $S=0$, then for every $i$ either $a_i=0$ or $s_i=0$, and in the second case $X_i=0$ almost surely; so $\sum_ia_iX_i=0$ almost surely and \eqref{eq:LDm1} holds.
Let $S>0$ and $F:=|\sum_ia_iX_i|/S$. By \eqref{eq:LD1}, \eqref{eq:M1} holds with $\mathsf c=5$, $\rho=r$, $\kappa_0=\kappa_\vth+\log n$ and $A=4+n\vth^{-1}\le5n\fa$.
Take $\kappa_1:=2rq+\log n+\kappa_\vth$. Then $Ae^{-\kappa_1}\le5n\fa\,e^{-\log n-\kappa_\vth}\le5e^{-23}$, since $\kappa_\vth\ge\log\fa+23$. By \eqref{eq:M2},
$\EE F^q\le2\,(5\kappa_1^r)^q$, and since $q\ge2$ and $2r+\kappa_\vth\ge1$,
\[
 \kappa_1\le(2r+\kappa_\vth)(q+\log n),\qquad\norm F_q\le2^{1/q}\cdot5\,\kappa_1^r\le10\,(2\fa+1+\kappa_\vth)^{\fa+\frac12}(q+\log n)^{\fa+\frac12}.
\]
The constant is at most $C_{\rm LD}$.

\emph{Step 3: proof of \eqref{eq:LDm2}.} Let $S_B:=(\sum_{i,j}s_i^2|B_{ij}|^2s_j^2)^{1/2}$, $D:=\lceil\log_2n\rceil$ and
\[
 \mathcal Q:=\sum_{i,j}\ol{X_i}B_{ij}X_j-\sum_is_i^2B_{ii}=\sum_iB_{ii}\big(|X_i|^2-s_i^2\big)+\sum_{i\ne j}\ol{X_i}B_{ij}X_j .
\]
Both $(\sum_i|B_{ii}|^2s_i^4)^{1/2}$ and $(\sum_{i\ne j}s_i^2|B_{ij}|^2s_j^2)^{1/2}$ are at most $S_B$. So by \eqref{eq:LD2} (with $b_i=B_{ii}$) and \eqref{eq:LD3}, for every $\kappa\ge\kappa_\vth+\log n$,
\[
 \PP\Big(|\mathcal Q|>\big(\ell(\kappa)^2+\tfrac12\ell_2(\kappa,n)\big)S_B\Big)\le\big(4+n\vth^{-1}+20n^2\lceil\log_22n\rceil\vth^{-1}\big)e^{-\kappa}\le25\,n^3\fa\,e^{-\kappa},
\]
where we used $\lceil\log_22n\rceil=1+\lceil\log_2n\rceil\le n$ (as $2^{n-1}\ge n$) and $4+n\fa+20n^3\fa\le25n^3\fa$. By \eqref{eq:loss},
\[
 \ell(\kappa)^2+\tfrac12\ell_2(\kappa,n)=c_n\,\kappa^{2\fa+1},\qquad c_n:=25\big(1+\sqrt{2D+2}\big).
\]
If $S_B=0$, the tail bound for $\kappa\to\infty$ shows $\mathcal Q=0$ almost surely. Let $S_B>0$ and $F:=|\mathcal Q|/S_B$. Then \eqref{eq:M1} holds with $\mathsf c=c_n$, $\rho=2r=2\fa+1$, $\kappa_0=\kappa_\vth+\log n$, $A=25n^3\fa$.
Because of the prefactor $n^3$ we take
\[
 \kappa_1:=4rq+3\log n+\kappa_\vth\qquad\big(\ge\max(\kappa_0,\,2\rho q)\big),
\]
so that $Ae^{-\kappa_1}\le25\fa e^{-\kappa_\vth}\le25e^{-23}$ and, by \eqref{eq:M2}, $\EE F^q\le2\,(c_n\kappa_1^{2\fa+1})^q$. It remains to bound $c_n$ and $\kappa_1$ by powers of $q+\log n$. First,
$D\le\log_2n+1$, so $2D+2\le\frac2{\log2}\log n+4\le3(q+\log n)$ (as $\frac2{\log 2}<3$ and $3q\ge6$), and $1\le\sqrt{(q+\log n)/2}$; hence
\[
 c_n\le25\big(\tfrac1{\sqrt2}+\sqrt3\big)\sqrt{q+\log n}\le75\sqrt{q+\log n}.
\]
Second, $4r=4\fa+2\ge3$ and $q\ge2$ give $\kappa_1\le(4\fa+2+\kappa_\vth)(q+\log n)$. Therefore
\[
 \norm{\mathcal Q}_q=S_B\norm F_q\le2^{1/q}\cdot75\,(4\fa+2+\kappa_\vth)^{2\fa+1}(q+\log n)^{2\fa+\frac32}S_B\le C_{\rm LD}(q+\log n)^{2\fa+\frac32}S_B .
\]
The statement for $\ol{X_i}$ follows by applying the lemma to the family $(\ol{X_i})$, which satisfies \eqref{eq:LDhyp} with the same $s_i$.
\end{proof}
\subsection{Row events at a fixed spectral parameter}
\paragraph{Coordinates, rows and partial integration.} Let $\nu\in\mathcal E_N(\vth,\dpm,C_0)$; recall $\Omega=\kk^{\mathcal P_N}$ and $\PP=\bigotimes_{(i,j)\in\mathcal P_N}\nu_{ij}$ (Definition~\ref{def:ens}).
A \emph{coordinate} is an element of $\mathcal P_N$; for $i,j\in\{1,\dots,N\}$ we write $\{i,j\}$ for the coordinate $(\min(i,j),\max(i,j))$. For $\mathcal Q\subset\mathcal P_N$ let $\mathcal Q^c:=\mathcal P_N\setminus\mathcal Q$,
$\PP_{\mathcal Q}:=\bigotimes_{(i,j)\in\mathcal Q}\nu_{ij}$ on $\kk^{\mathcal Q}$, and write $x=(x_{\mathcal Q^c},x_{\mathcal Q})$ for $x\in\Omega$.
A function $f$ on $\Omega$ \emph{does not depend on $\mathcal Q$} if $f(x)=f(x')$ whenever $x_{\mathcal Q^c}=x'_{\mathcal Q^c}$. For $T\subset\{1,\dots,N\}$,
\[
 \mathcal Q(T):=\big\{(i,j)\in\mathcal P_N:\ \{i,j\}\cap T\ne\emptyset\big\}\qquad(\text{the coordinates in the rows }T).
\]
By Lemma~\ref{lem:split} (whose statement and proof use nothing from Sections~\ref{sec:prob}--\ref{sec:weak}; it is placed at the end of this section only in order to keep the numbering of the other statements),
the law $\PP$ is the product $\PP_{\mathcal Q^c}\otimes\PP_{\mathcal Q}$, and the partial integral
\[
 (P_{\mathcal Q}f)(x):=\int_{\kk^{\mathcal Q}}f(x_{\mathcal Q^c},y)\,\PP_{\mathcal Q}(\rd y)
\]
is defined for every measurable $f:\Omega\to[0,\infty]$, and for complex $f$ at every point $x$ where $\int|f(x_{\mathcal Q^c},y)|\PP_{\mathcal Q}(\rd y)<\infty$. We write
$P_T:=P_{\mathcal Q(T)}$ (integration over the rows in $T$), $P_i:=P_{\{i\}}$ and $Q_i:=1-P_i$.

The entries of the minor $H^{(T)}$ are the $h_{kl}$ with $k,l\notin T$, and these are functions of the coordinates $\{k,l\}\notin\mathcal Q(T)$. Hence, for fixed $z$,
\begin{equation}\tag{R}\label{eq:rowdep}
 H^{(T)},\quad G^{(T)}(z),\quad m^{(T)}(z)\qquad\text{do not depend on }\mathcal Q(T).
\end{equation}
On the other hand, for $i\ne j$, the row variables $(h_{ik})_{k\ne i}$ are functions of the coordinates in $\mathcal Q_i:=\{\{i,k\}:k\ne i\}\subset\mathcal Q(\{i\})$, and the two families $(h_{ik})_{k\notin\{i,j\}}$, $(h_{lj})_{l\notin\{i,j\}}$ are functions of the coordinates in
\[
 \mathcal Q_{ij}:=\big\{\{i,k\}:k\notin\{i,j\}\big\}\ \cup\ \big\{\{l,j\}:l\notin\{i,j\}\big\}\subset\mathcal Q(\{i,j\}),
\]
where the two sets on the right are disjoint (if $\{i,k\}=\{l,j\}$ with $k\ne j$ and $i\ne j$, then $i=j$ or $k=j$, which is excluded).

\begin{definition}[Good event at $z$]\label{def:good}
For $z\in\C_+$ and $\kappa\ge1$ let $\cG_\kappa(z)\subset\Omega$ be the event on which
\begin{enumerate}[label=\textup{(G\arabic*)},leftmargin=2.8em,itemsep=0pt]
\item $|h_{kl}|\le\kappa^{\fa}\sigma_{kl}$ for all $k,l$;
\item $|Z_i|\le\ell_2(\kappa,N)\big(\sum^{(i)}_{k,l}\sigma_{ik}^2|G^{(i)}_{kl}|^2\sigma_{il}^2\big)^{1/2}$ for all $i$;
\item $\big|\sum^{(i)}_kh_{ik}G^{(i)}_{kj}\big|\le\ell(\kappa)\big(\sum^{(i)}_k\sigma_{ik}^2|G^{(i)}_{kj}|^2\big)^{1/2}$ and $\big|\sum^{(i)}_kG^{(i)}_{jk}h_{ki}\big|\le\ell(\kappa)\big(\sum^{(i)}_k|G^{(i)}_{jk}|^2\sigma_{ki}^2\big)^{1/2}$ for all $i\ne j$;
\item $\big|\sum^{(ij)}_{k,l}h_{ik}G^{(ij)}_{kl}h_{lj}\big|\le\ell(\kappa)^2\big(\sum^{(ij)}_{k,l}\sigma_{ik}^2|G^{(ij)}_{kl}|^2\sigma_{lj}^2\big)^{1/2}$ for all $i\ne j$.
\end{enumerate}
Here $G^{(T)}=G^{(T)}(z)$, $Z_i=Z_i(z)$ is as in Definition~\ref{def:XA}, and $\ell,\ell_2$ are the loss factors \eqref{eq:loss}.
\end{definition}

\begin{lemma}[Probability and consequences of the good event]\label{lem:good}
Let $\nu\in\mathcal E_N(\vth,\dpm,C_0)$, $z\in\C_+$, $\kappa\ge\kappa_\vth+\log N$. Then $\cG_\kappa(z)$ is a closed (in particular measurable) subset of $\Omega$ and
$\PP(\cG_\kappa(z)^c)\le46\,\vth^{-1}N^4e^{-\kappa}\le CN^5e^{-\kappa}$.
On $\cG_\kappa(z)$, with
\begin{equation}\label{eq:Psim}
 \Psi_m(z):=\sqrt{\frac{\im m(z)}{N\eta}}+\frac1{N\eta}+\frac1{\sqrt N},
\end{equation}
we have for all $i\ne j$ (no a priori assumption on $G$):
\begin{align}
 |X_i|&\le2C_0\,\ell_2(\kappa,N)\,\Psi_m,\label{eq:Xi}\\
 |A_i|&\le2C_0^2\,\ell(\kappa)^2\,|G_{ii}|\,\Psi_m^2,\label{eq:Ai}\\
 |G_{ij}|&\le3C_0\,\ell(\kappa)^2\,|G_{ii}|\,|G^{(i)}_{jj}|\,\Psi_m.\label{eq:Gij}
\end{align}
\end{lemma}
\begin{proof}
\emph{Measurability.} For fixed $z$ the functions $h_{kl}$, $G^{(T)}_{kl}(z)$ and $Z_i(z)$ are continuous on $\Omega$ (Lemma~\ref{lem:weyl}(e)). Each of the conditions (G1)--(G4) is a finite conjunction of inequalities $|f|\le g$ between continuous functions; so $\cG_\kappa(z)$ is closed.

\emph{Probability of \textup{(G1)}.} By Lemma~\ref{lem:moments} with $T=\kappa^{\fa}\ge1$, (G1) fails with probability at most $N^2\vth^{-1}e^{-\kappa}$.

\emph{Probability of \textup{(G2)}.} Fix $i$ and let $n:=N-1$. For $k\ne i$ put $X_k:=h_{ki}=\ol{h_{ik}}$. By Definition~\ref{def:XA},
\[
 Z_i=\sum^{(i)}_{k,l}\ol{X_k}\,G^{(i)}_{kl}X_l-\sum^{(i)}_k\sigma_{ik}^2G^{(i)}_{kk}=\sum^{(i)}_kG^{(i)}_{kk}\big(|X_k|^2-\sigma_{ik}^2\big)+\sum^{(i)}_{k\ne l}\ol{X_k}\,G^{(i)}_{kl}X_l .
\]
Fix a point $x'\in\kk^{\mathcal Q_i^c}$. By \eqref{eq:rowdep} the numbers $G^{(i)}_{kl}$ are determined by $x'$. By Lemma~\ref{lem:split}(e), on the probability space $(\kk^{\mathcal Q_i},\PP_{\mathcal Q_i})$ the coordinate functions $y\mapsto y_{\{i,k\}}$, $k\ne i$, are independent with laws $\nu_{\{i,k\}}$;
$X_k$ is such a coordinate or its conjugate, so by Remark~\ref{rem:laws} the family $(X_k)_{k\ne i}$ satisfies \eqref{eq:LDhyp} with $s_k=\sigma_{ik}$. Since $\kappa\ge\kappa_\vth+\log n$, \eqref{eq:LD2} (with $b_k=G^{(i)}_{kk}$) and \eqref{eq:LD3} (with $B_{kl}=G^{(i)}_{kl}$) show that, outside a set of $\PP_{\mathcal Q_i}$-probability at most
\[
 p_2:=\big(4+n\vth^{-1}+20n^2\lceil\log_22n\rceil\vth^{-1}\big)e^{-\kappa}\le25N^3\vth^{-1}e^{-\kappa}
\]
(we used $\lceil\log_22n\rceil\le n$), we have
\[
 |Z_i|\le\ell(\kappa)^2\Big(\sum^{(i)}_k\sigma_{ik}^4|G^{(i)}_{kk}|^2\Big)^{1/2}+\tfrac12\ell_2(\kappa,n)\Big(\sum^{(i)}_{k\ne l}\sigma_{ik}^2|G^{(i)}_{kl}|^2\sigma_{il}^2\Big)^{1/2}\le\ell_2(\kappa,N)\Big(\sum^{(i)}_{k,l}\sigma_{ik}^2|G^{(i)}_{kl}|^2\sigma_{il}^2\Big)^{1/2},
\]
by Remark~\ref{rem:LDn}(iii). Thus the set $\mathcal B_i\subset\Omega$ where (G2) fails for the index $i$ (an open set) satisfies $\PP_{\mathcal Q_i}(\{y:(x',y)\in\mathcal B_i\})\le p_2$ for every $x'$, and Lemma~\ref{lem:split}(b) gives $\PP(\mathcal B_i)\le p_2$.

\emph{Probability of \textup{(G3)}.} Fix $i\ne j$. As before, for fixed $x'\in\kk^{\mathcal Q_i^c}$ the coefficients $G^{(i)}_{kj}$, $G^{(i)}_{jk}$ ($k\ne i$) are constants and the families $(h_{ik})_{k\ne i}$, $(h_{ki})_{k\ne i}$ satisfy \eqref{eq:LDhyp} with $s_k=\sigma_{ik}$ under $\PP_{\mathcal Q_i}$. By \eqref{eq:LD1} and Lemma~\ref{lem:split}(b),
each of the two inequalities in (G3) fails with probability at most $(4+n\vth^{-1})e^{-\kappa}\le5N\vth^{-1}e^{-\kappa}$.

\emph{Probability of \textup{(G4)}.} Fix $i\ne j$ (recall $N\ge3$) and $x'\in\kk^{\mathcal Q_{ij}^c}$. By \eqref{eq:rowdep} the numbers $G^{(ij)}_{kl}$ are determined by $x'$. Under $\PP_{\mathcal Q_{ij}}$ the $2(N-2)$ coordinate functions are independent; hence the variables
$X_k:=h_{ik}$ $(k\notin\{i,j\})$ and $Y_l:=h_{lj}$ $(l\notin\{i,j\})$ are altogether independent, and the two families satisfy \eqref{eq:LDhyp} with $s_k=\sigma_{ik}$, $t_l=\sigma_{lj}$. By \eqref{eq:LD4} with $n=n'=N-2$ and Lemma~\ref{lem:split}(b), (G4) fails for $(i,j)$ with probability at most $10N^2\vth^{-1}e^{-\kappa}$.

\emph{Union bound.} There are $N$ indices $i$ and $N(N-1)$ ordered pairs $i\ne j$. So
\[
 \PP(\cG_\kappa(z)^c)\le\big(N^2+25N^4+2\cdot5N^3+10N^4\big)\vth^{-1}e^{-\kappa}\le46\,\vth^{-1}N^4e^{-\kappa}.
\]

\emph{Consequences.} From now on we work at a point of $\cG_\kappa(z)$ and write $\ell=\ell(\kappa)$, $\ell_2=\ell_2(\kappa,N)$. Let $T\subset\{1,\dots,N\}$ with $|T|\le2$ and $a,b\in T$. By (C) and the Ward identity for the minor $H^{(T)}$ (Lemma~\ref{lem:basic}(b)),
\begin{equation}\tag{W}\label{eq:WardT}
 \sum^{(T)}_{k}\sigma_{ak}^2|G^{(T)}_{kl}|^2\le\frac{C_0}N\cdot\frac{\im G^{(T)}_{ll}}{\eta},\qquad\sum^{(T)}_{k}|G^{(T)}_{lk}|^2\sigma_{ka}^2\le\frac{C_0}N\cdot\frac{\im G^{(T)}_{ll}}{\eta}\qquad(l\notin T),
\end{equation}
\begin{equation}\tag{W$'$}\label{eq:WardT2}
 \sum^{(T)}_{k,l}\sigma_{ak}^2|G^{(T)}_{kl}|^2\sigma_{lb}^2\le\frac{C_0}N\sum^{(T)}_l\sigma_{lb}^2\frac{\im G^{(T)}_{ll}}\eta\le\frac{C_0^2}{N^2}\sum^{(T)}_l\frac{\im G^{(T)}_{ll}}\eta=C_0^2\,\frac{\im m^{(T)}}{N\eta}\le C_0^2\Big(\frac{\im m}{N\eta}+\frac{|T|}{(N\eta)^2}\Big),
\end{equation}
where the last step is (I5) of Lemma~\ref{lem:schur} (and $\im G^{(T)}_{ll}>0$).

\emph{Proof of \eqref{eq:Xi}.} $X_i=h_{ii}-Z_i$. By (G1) and (C), $|h_{ii}|\le\kappa^{\fa}\sqrt{C_0/N}\le C_0\ell_2N^{-1/2}$ (as $\ell_2\ge\kappa^{\fa}$ and $C_0\ge1$). By (G2) and \eqref{eq:WardT2} with $T=\{i\}$, $a=b=i$,
$|Z_i|\le\ell_2C_0\big(\sqrt{\im m/(N\eta)}+1/(N\eta)\big)$. Hence $|X_i|\le C_0\ell_2\Psi_m$.

\emph{Proof of \eqref{eq:Ai}.} By the last formula of Lemma~\ref{lem:sce}, by (G3) and \eqref{eq:WardT} with $T=\{i\}$, $a=i$, $l=j$ (for each of the two factors), and then by (C) and (I5),
\[
 \frac{|A_i|}{|G_{ii}|}\le\sigma_{ii}^2+\sum_{j\ne i}\sigma_{ij}^2\,\ell^2\,\frac{C_0}N\frac{\im G^{(i)}_{jj}}\eta\le\frac{C_0}N+\ell^2C_0^2\,\frac{\im m^{(i)}}{N\eta}\le C_0^2\ell^2\Big(\frac1N+\frac{\im m}{N\eta}+\frac1{(N\eta)^2}\Big)\le C_0^2\ell^2\Psi_m^2 .
\]

\emph{Proof of \eqref{eq:Gij}.} By the second formula of (I2) in Lemma~\ref{lem:schur}, (G1), (C), (G4) and \eqref{eq:WardT2} with $T=\{i,j\}$, $a=i$, $b=j$,
\[
 \frac{|G_{ij}|}{|G_{ii}||G^{(i)}_{jj}|}\le|h_{ij}|+\Big|\sum^{(ij)}_{k,l}h_{ik}G^{(ij)}_{kl}h_{lj}\Big|\le\kappa^{\fa}\sqrt{\frac{C_0}N}+\ell^2C_0\Big(\sqrt{\frac{\im m}{N\eta}}+\frac{\sqrt2}{N\eta}\Big)\le\sqrt2\,C_0\ell^2\Psi_m .
\]
The three bounds are slightly better than \eqref{eq:Xi}--\eqref{eq:Gij}.
\end{proof}
\subsection{A crude bound on the norm}
\begin{lemma}[A priori norm bound]\label{lem:norm0}
Let $\nu\in\mathcal E_N(\vth,\dpm,C_0)$, $N\ge N_0:=\max\{13,\sqrt{C_0}\}$, $\kappa\ge\kappa_\vth+3\log N$. Then
\[
 \PP\Big(\norm H>25\,\kappa^{\fa}\sqrt{C_0\log N}\Big)\le N^2\vth^{-1}e^{-\kappa}+2e^{-N\log N}.
\]
\end{lemma}
\begin{proof}
Let $T:=\kappa^{\fa}\ge1$ and $\mathcal T:=\{x\in\Omega:\ |x_{ij}|\le T\sigma_{ij}\text{ for all }(i,j)\in\mathcal P_N\}$; by Lemma~\ref{lem:moments}, $\PP(\mathcal T^c)\le N^2\vth^{-1}e^{-\kappa}$.

\emph{Step 1: truncation and recentring.} For $(i,j)\in\mathcal P_N$ put
\[
 \mu_{ij}:=\EE\big[x_{ij}\bone(|x_{ij}|\le T\sigma_{ij})\big]\in\kk,\qquad\wh x_{ij}:=x_{ij}\bone(|x_{ij}|\le T\sigma_{ij})-\mu_{ij},
\]
and let $\wh H:=H(\wh x)$, $M:=H(\mu)$ be the Hermitian matrices that Definition~\ref{def:ens} associates with the points $\wh x=(\wh x_{ij})$ and $\mu=(\mu_{ij})$ of $\Omega$; both have entries in $\kk$.
The random variables $\wh x_{ij}$, $(i,j)\in\mathcal P_N$, are independent (each is a Borel function of the coordinate $x_{ij}$), centred, and
\[
 |\wh x_{ij}|\le2T\sigma_{ij},\qquad|\mu_{ij}|\le\sigma_{ij}\vth^{-1/2}e^{-\kappa/2}\le\sqrt{C_0/N}\,\vth^{-1/2}e^{-\kappa/2},
\]
by the centering condition of Definition~\ref{def:ass}, Lemma~\ref{lem:moments}(b) and (C), exactly as in Step~1 of the proof of Lemma~\ref{lem:LD}. Since $x\mapsto H(x)$ is $\R$-linear (Lemma~\ref{lem:weyl}(d)) and $x=\wh x+\mu$ on $\mathcal T$,
\[
 H=\wh H+M\quad\text{on }\mathcal T,\qquad\norm M\le\Big(\sum_{i,j}|M_{ij}|^2\Big)^{1/2}\le N\max_{(i,j)}|\mu_{ij}|\le\sqrt{NC_0\fa}\;e^{-\kappa/2}\le1 .
\]
The last inequality holds because $\kappa\ge\log\fa+3\log N$ and $N^2\ge C_0$.

\emph{Step 2: one vector.} Fix $y\in\kk^N$ with $|y|\le1$. Since $\wh H$ is Hermitian with $\wh H_{ii}=\re\wh x_{ii}$ and $\wh H_{ij}=\wh x_{ij}$, $\wh H_{ji}=\ol{\wh x_{ij}}$ for $i<j$,
\[
 \langle y,\wh Hy\rangle=\sum_i\xi_{ii}+\sum_{i<j}\xi_{ij},\qquad\xi_{ii}:=|y_i|^2\re\wh x_{ii},\qquad\xi_{ij}:=2\re\big(\ol{y_i}\,\wh x_{ij}\,y_j\big)\ \ (i<j).
\]
The $\xi_{ij}$, $(i,j)\in\mathcal P_N$, are independent, real, centred, with $|\xi_{ii}|\le2T\sigma_{ii}|y_i|^2$ and $|\xi_{ij}|\le4T\sigma_{ij}|y_i||y_j|$, and by (C)
\[
 V:=\sum_i4T^2\sigma_{ii}^2|y_i|^4+\sum_{i<j}16T^2\sigma_{ij}^2|y_i|^2|y_j|^2\le8T^2\sum_{i,j}\sigma_{ij}^2|y_i|^2|y_j|^2\le\frac{8T^2C_0}N|y|^4\le\frac{8T^2C_0}N .
\]
By (H) of Section~\ref{sec:prob} with $\kappa$ replaced by $4N\log N$: $\ \PP\big(|\langle y,\wh Hy\rangle|>8T\sqrt{C_0\log N}\big)\le2e^{-4N\log N}$, because $\sqrt{2\cdot4N\log N\cdot V}\le8T\sqrt{C_0\log N}$.

\emph{Step 3: the net.} Let $\mathcal N\subset\kk^N$ be the net of Lemma~\ref{lem:net}; $|\mathcal N|\le e^{3N\log N}$ as $N\ge13$, and every $y\in\mathcal N$ has $|y|\le1$. By Step~2 and a union bound, outside an event $\mathcal H_3$ with
$\PP(\mathcal H_3)\le2e^{3N\log N-4N\log N}=2e^{-N\log N}$ we have $\max_{y\in\mathcal N}|\langle y,\wh Hy\rangle|\le8T\sqrt{C_0\log N}$, hence $\norm{\wh H}\le16T\sqrt{C_0\log N}$ by Lemma~\ref{lem:net}
(applicable since $\wh H$ is Hermitian with entries in $\kk$). On $\mathcal T\setminus\mathcal H_3$, by Step~1,
\[
 \norm H\le\norm{\wh H}+\norm M\le16T\sqrt{C_0\log N}+1\le25\,\kappa^{\fa}\sqrt{C_0\log N},
\]
since $T\sqrt{C_0\log N}\ge1$. The set $\{\norm H>t\}$ is open in $\Omega$ (Lemma~\ref{lem:weyl}(a),(d)), hence measurable, and it is contained in $\mathcal T^c\cup\mathcal H_3$.
\end{proof}
\subsection{Levels}\label{sec:levels}
\begin{definition}[The three levels]\label{def:levels}
Given $K$ as in \eqref{eq:Krange} put
\[
 K_w:=K^2,\qquad p_1:=2\lceil K^{3/2}/2\rceil,\qquad p_2:=2\lceil K/2\rceil,\qquad\cG(z):=\cG_{K_w}(z),
\]
\[
 \fL:=3C_0^2\,\ell_2(K_w,N)\ \ (\le K^{4\fa+3}\text{ for }N\ge N_0),\qquad\Emax:=25K_w^{\fa}\sqrt{C_0\log N}\ \ (\le K^{2\fa+1}\text{ for }N\ge N_0).
\]
\end{definition}
\noindent So $\PP(\cG(z)^c)\le CN^5e^{-K^2}$ and on $\cG(z)$: $|X_i|\le\fL\Psi_m$, $|A_i|\le\fL|G_{ii}|\Psi_m^2$, $|G_{ij}|\le\fL|G_{ii}||G^{(i)}_{jj}|\Psi_m$ (see Lemma~\ref{lem:levels}).
The weak law is proved at level $K_w$; fluctuation averaging is used with $p=p_1$ (first pass) and $p=p_2$ (second pass). The relations used are
\begin{equation}\label{eq:levelrel}
 K_w\ge40\,p_1\log N,\qquad p_1\ge40\,p_2\log N,\qquad p_2\ge K\ge(\log N)^3\qquad(N\ge N_0),
\end{equation}
which hold because $K^{1/2}\ge(\log N)^{3/2}\ge 41\log N$ for $N\ge N_0$.

The next lemma collects, once and for all, the elementary inequalities between $N$, $K$, $\fL$, $\Emax$ that are used in Section~\ref{sec:weak}. We recall from \eqref{eq:ledger0} and \eqref{eq:Krange} that $(\log N)^3\le K\le N^{c_0}$; of the exponent ledger we use only
\begin{equation}\tag{L}\label{eq:ledgerA3}
 C_1-1\ \ge\ 12\fa+18,\qquad\qquad c_0\,(6\fa+6)\ \le\ \tfrac1{32}.
\end{equation}
Both hold for the values in \eqref{eq:ledger0}: $C_1-1=18\fa+19$ and $c_0=1/(192\fa+288)$.

\begin{lemma}[Standing inequalities]\label{lem:levels}
Assume \eqref{eq:ledgerA3}. For every constant $K_0\ge1$ depending only on the model parameters there is $N_0$, depending only on the model parameters and on $K_0$, such that for all $N\ge N_0$ and all $K$ with $(\log N)^3\le K\le N^{c_0}$:
\begin{enumerate}[label=\textup{(S\arabic*)},leftmargin=2.8em,itemsep=1pt]
\item $K\ge K_0$, $\ N\ge13$, $\ N^2\ge C_0$, and $K_w=K^2\ge\kappa_\vth+3\log N$;
\item $150\le\fL\le K^{4\fa+3}$ and $1\le\Emax\le K^{2\fa+1}$;
\item the relations \eqref{eq:levelrel} hold, and $K^{3/2}\le p_1\le2K^{3/2}$;
\item $K^{6\fa+6}\le N^{1/32}$ and $K_0\,N^{1/32}\le N^{1/8}$;
\item for every $z\in\C_+$: $\PP(\cG(z)^c)\le46\,\vth^{-1}N^4e^{-K^2}$, and on $\cG(z)$, for all $i\ne j$,
\[
 |X_i|\le\fL\Psi_m,\qquad|A_i|\le\fL|G_{ii}|\Psi_m^2,\qquad|G_{ij}|\le\fL|G_{ii}||G^{(i)}_{jj}|\Psi_m .
\]
\end{enumerate}
\end{lemma}
\begin{proof}
Let $\mathsf L:=\log N$. Since $K\ge\mathsf L^3$, every condition of the form ``$K\ge$ a constant'' or ``$K^{\alpha}\ge c\,\mathsf L^{\beta}$ with $3\alpha>\beta$'' holds for $N\ge N_0$.

(S1) $N\ge13$ and $N^2\ge C_0$ for $N\ge N_0$; $K\ge\mathsf L^3\ge K_0$ and $K^2\ge\mathsf L^6\ge\kappa_\vth+3\mathsf L$ for $N\ge N_0$.

(S2) By \eqref{eq:loss}, $\fL=150\,C_0^2\,K^{4\fa+2}\sqrt{2\lceil\log_2N\rceil+2}\ge150$, and $2\lceil\log_2N\rceil+2\le3\mathsf L+4\le4\mathsf L$ for $\mathsf L\ge4$. So $\fL\le K^{4\fa+3}$ as soon as $K\ge300\,C_0^2\mathsf L^{1/2}$,
which holds when $\mathsf L^{5/2}\ge300C_0^2$. Similarly $\Emax=25K^{2\fa}\sqrt{C_0\mathsf L}\le K^{2\fa+1}$ when $\mathsf L^{5/2}\ge25\sqrt{C_0}$, and $\Emax\ge1$.

(S3) By definition $K^{3/2}\le p_1\le K^{3/2}+2\le2K^{3/2}$ and $K\le p_2\le K+2$. If $\mathsf L\ge41^2$ then $K^{1/2}\ge\mathsf L^{3/2}\ge41\mathsf L$, hence
$40p_1\mathsf L\le40K^{3/2}\mathsf L+80\mathsf L\le41K^{3/2}\mathsf L\le K^2$ (using $K^{3/2}\ge80$) and $40p_2\mathsf L\le41K\mathsf L\le K^{3/2}\le p_1$.

(S4) $K^{6\fa+6}\le N^{c_0(6\fa+6)}\le N^{1/32}$ by \eqref{eq:ledgerA3}, and $K_0\le N^{3/32}$ for $N\ge N_0$.

(S5) By (S1), $\kappa:=K_w$ satisfies the hypothesis of Lemma~\ref{lem:good}. The three bounds follow from \eqref{eq:Xi}--\eqref{eq:Gij} because $C_0\ge1$, $2C_0\ell_2(K_w,N)\le\fL$ and $2C_0^2\ell(K_w)^2\le3C_0\ell(K_w)^2\cdot C_0\le3C_0^2\ell_2(K_w,N)=\fL$ (Remark~\ref{rem:LDn}(iii): $\ell^2\le\ell_2$).
\end{proof}
\noindent\emph{Convention.} In Section~\ref{sec:weak}, ``for $N\ge N_0$'' means: (S1)--(S5) hold with a constant $K_0$ that is the maximum of the finitely many explicit constants (depending only on the model parameters) appearing in the proofs;
each use is marked by ``(S1)'', \dots, or by the words ``for $N\ge N_0$''. $N_0$ is very large (for instance \eqref{eq:levelrel} is derived from $\log N\ge41^2$); this is harmless for the statements.
\subsection{Splitting of the product law}\label{sec:split}
The following lemma is pure measure theory: it does not use the assumptions (A)--(C), \eqref{eq:subexp}, or anything else from Sections~\ref{sec:prob}--\ref{sec:weak}.
It is used in the proof of Lemma~\ref{lem:good} and, systematically, in Section~\ref{sec:FA}.

\begin{lemma}[Splitting of the product law; partial integration]\label{lem:split}
Let $I$ be a finite set and $(\nu_\iota)_{\iota\in I}$ Borel probability measures on $\kk$. For $J\subset I$ let $\kk^J$ carry the product $\sigma$-algebra (which is its Borel $\sigma$-algebra) and the product measure $\PP_J:=\bigotimes_{\iota\in J}\nu_\iota$;
$\kk^\emptyset$ is a one-point space and $\PP_\emptyset$ its unit mass. Fix $J\subset I$, put $J^c:=I\setminus J$, and for $x\in\kk^I$ write $x_J:=(x_\iota)_{\iota\in J}$.
\begin{enumerate}[label=\textup{(\alph*)},leftmargin=2.2em,itemsep=1pt]
\item \emph{(Splitting)} The map $\pi_J:\kk^I\to\kk^{J^c}\times\kk^J$, $x\mapsto(x_{J^c},x_J)$, is a bijection, it is measurable and so is its inverse, and it maps $\PP_I$ to the product measure: $(\pi_J)_*\PP_I=\PP_{J^c}\otimes\PP_J$.
We identify $x$ with $(x_{J^c},x_J)$.
\item \emph{(Tonelli)} For every measurable $f:\kk^I\to[0,\infty]$ the function
\[
 (P_Jf)(x):=\int_{\kk^J}f(x_{J^c},y)\,\PP_J(\rd y)
\]
is measurable, depends only on $x_{J^c}$, and $\int f\,\rd\PP_I=\int P_Jf\,\rd\PP_I$. In particular, if $\mathcal B\subset\kk^I$ is measurable, $p\ge0$, and $\PP_J(\{y:(x',y)\in\mathcal B\})\le p$ for every $x'\in\kk^{J^c}$, then $\PP_I(\mathcal B)\le p$.
\item \emph{(Fubini)} Let $f:\kk^I\to\C$ be $\PP_I$-integrable. The set $\mathcal D_f$ of all $x$ with $\int|f(x_{J^c},y)|\PP_J(\rd y)<\infty$ is measurable, depends only on $x_{J^c}$, and has $\PP_I(\mathcal D_f)=1$.
Define $P_Jf$ by the formula in \textup{(b)} on $\mathcal D_f$ and by $0$ on the complement. Then $P_Jf$ is measurable and integrable, depends only on $x_{J^c}$, and $\int f\,\rd\PP_I=\int P_Jf\,\rd\PP_I$.
\item \emph{(Conditional expectation $=$ integration over one factor)} Let $f$ be as in \textup{(c)} and let $g:\kk^I\to\C$ be bounded, measurable and independent of the coordinates in $J$ (i.e.\ $g(x)=g(\wt x)$ whenever $x_{J^c}=\wt x_{J^c}$). Then
$P_J(gf)=g\,P_Jf$ and
\[
 \int g\,f\,\rd\PP_I=\int g\,P_Jf\,\rd\PP_I .
\]
Equivalently: $P_Jf$ is a version of the conditional expectation of $f$ given the $\sigma$-algebra generated by the coordinates $x_\iota$, $\iota\in J^c$.
\item \emph{(Independence; freezing)} Under $\PP_I$ the coordinate functions $x\mapsto x_\iota$, $\iota\in I$, are independent and $x_\iota$ has law $\nu_\iota$; the random vectors $x_{J^c}$ and $x_J$ are independent with laws $\PP_{J^c}$ and $\PP_J$.
For every fixed $x'\in\kk^{J^c}$, the coordinate functions $y\mapsto y_\iota$, $\iota\in J$, on the probability space $(\kk^J,\PP_J)$ are independent with laws $\nu_\iota$, and $(P_Jf)(x',\cdot)$ is the expectation of the random variable $y\mapsto f(x',y)$ on this space.
\item \emph{(Tower property)} The statements \textup{(a)--(e)} hold for the index set $J$ and any $J_1\subset J$ in place of $I$ and $J$: $\PP_J=\PP_{J\setminus J_1}\otimes\PP_{J_1}$. Consequently, for measurable $f:\kk^I\to[0,\infty]$ and $J_1\subset J$:
$P_JP_{J_1}f=P_Jf$ at every point of $\kk^I$.
\end{enumerate}
For the ensemble of Definition~\ref{def:ens} the lemma applies with $I=\mathcal P_N$, $\kk^I=\Omega$, $\PP_I=\PP$ and $J=\mathcal Q\subset\mathcal P_N$; then $P_J$ is the operator $P_{\mathcal Q}$ of Section~\ref{sec:prob}, and $P_T=P_{\mathcal Q(T)}$.
\end{lemma}
\begin{proof}
(a) The inverse of $\pi_J$ is the map that glues $(x',y)\in\kk^{J^c}\times\kk^J$ to the point $x$ with $x_\iota=x'_\iota$ for $\iota\in J^c$ and $x_\iota=y_\iota$ for $\iota\in J$. A map into a product space is measurable if and only if its coordinates are;
every coordinate of $\pi_J$ and of $\pi_J^{-1}$ is a coordinate function, hence measurable. The measures $(\pi_J)_*\PP_I$ and $\PP_{J^c}\otimes\PP_J$ are probability measures on $\kk^{J^c}\times\kk^J$. On a box
$R=\prod_{\iota\in J^c}A_\iota\times\prod_{\iota\in J}A_\iota$ with Borel sets $A_\iota\subset\kk$ they agree:
\[
 \PP_I\big(\pi_J^{-1}(R)\big)=\PP_I\Big(\prod_{\iota\in I}A_\iota\Big)=\prod_{\iota\in I}\nu_\iota(A_\iota)=\PP_{J^c}\Big(\prod_{\iota\in J^c}A_\iota\Big)\PP_J\Big(\prod_{\iota\in J}A_\iota\Big).
\]
The boxes form a $\pi$-system that generates the product $\sigma$-algebra ($I$ is finite), so the two measures are equal by the uniqueness theorem for finite measures.
(Mathlib: \texttt{MeasureTheory.measurePreserving\_piEquivPiSubtypeProd}.)

(b) Let $F:=f\circ\pi_J^{-1}$, a measurable function on $\kk^{J^c}\times\kk^J$. By Tonelli's theorem, $\varphi(x'):=\int F(x',y)\PP_J(\rd y)$ is a measurable function of $x'\in\kk^{J^c}$ and $\int F\,\rd(\PP_{J^c}\otimes\PP_J)=\int\varphi\,\rd\PP_{J^c}$.
By (a) and the change of variables formula for image measures, $\int f\,\rd\PP_I=\int F\,\rd(\PP_{J^c}\otimes\PP_J)$. Now $P_Jf=\varphi\circ\mathrm{pr}$, where $\mathrm{pr}(x):=x_{J^c}$; this is measurable and depends only on $x_{J^c}$. By (a), the image of $\PP_I$ under $\mathrm{pr}$ is the first marginal of
$\PP_{J^c}\otimes\PP_J$, i.e.\ $\PP_{J^c}$; hence $\int P_Jf\,\rd\PP_I=\int\varphi\,\rd\PP_{J^c}=\int f\,\rd\PP_I$. For the last claim take $f=\bone_{\mathcal B}$: then $(P_Jf)(x)=\PP_J(\{y:(x_{J^c},y)\in\mathcal B\})\le p$ at every point, so $\PP_I(\mathcal B)=\int P_Jf\,\rd\PP_I\le p$.

(c) By (b) applied to $|f|$, the function $P_J|f|$ is measurable, depends only on $x_{J^c}$, and $\int P_J|f|\rd\PP_I=\int|f|\rd\PP_I<\infty$; so $\mathcal D_f=\{P_J|f|<\infty\}$ is measurable, depends only on $x_{J^c}$, and has full measure.
Write $f=f_1-f_2+\ii(f_3-f_4)$ with $f_k\ge0$ measurable and $f_k\le|f|$. On $\mathcal D_f$ all $P_Jf_k$ are finite and $P_Jf=P_Jf_1-P_Jf_2+\ii(P_Jf_3-P_Jf_4)$ by linearity of the integral; so $P_Jf$ is measurable, depends only on $x_{J^c}$, $|P_Jf|\le P_J|f|$,
and (b) for each $f_k$ gives $\int P_Jf\,\rd\PP_I=\int f\,\rd\PP_I$ (the complement of $\mathcal D_f$ is a null set).

(d) Fix $x$. Since $g(x_{J^c},y)=g(x)$ for all $y$: if $x\in\mathcal D_f$ then $x\in\mathcal D_{gf}$ and $(P_J(gf))(x)=g(x)(P_Jf)(x)$. If $x\notin\mathcal D_f$ then $(P_Jf)(x)=0$ by convention, and also $(P_J(gf))(x)=0$:
either $g(x)\ne0$, and then $x\notin\mathcal D_{gf}$; or $g(x)=0$, and then $y\mapsto(gf)(x_{J^c},y)$ vanishes identically. So $P_J(gf)=g\,P_Jf$ at every point. Now (c) for the integrable function $gf$ gives $\int gf\,\rd\PP_I=\int P_J(gf)\,\rd\PP_I=\int g\,P_Jf\,\rd\PP_I$.
For the reformulation: $P_Jf$ is integrable and depends only on $x_{J^c}$, so it is measurable with respect to $\Sigma:=\sigma(x_\iota:\iota\in J^c)$ (it equals $\varphi\circ\mathrm{pr}$ with $\varphi$ measurable);
and every bounded $\Sigma$-measurable $g$ is of the form $\wt g\circ\mathrm{pr}$, hence independent of the coordinates in $J$. The displayed identity for all such $g$ is the defining property of the conditional expectation $\EE[f\,|\,\Sigma]$.

(e) By definition of the product measure the joint law of the coordinate functions under $\PP_I$ is $\bigotimes_\iota\nu_\iota$, which is the definition of independence with marginal laws $\nu_\iota$; the same for $\PP_J$.
By (a), the joint law of $(x_{J^c},x_J)$ is $\PP_{J^c}\otimes\PP_J$. The last statement is the definition of $P_Jf$.

(f) The first sentence is (a)--(e) for the finite index set $J$, the measures $(\nu_\iota)_{\iota\in J}$ and the subset $J_1$. For the consequence fix $x\in\kk^I$ and apply (b), in this setting, to the nonnegative measurable function $y\mapsto f(x_{J^c},y)$ on $\kk^J$:
\[
 (P_Jf)(x)=\int_{\kk^{J}}\Big[\int_{\kk^{J_1}}f\big(x_{J^c},y_{J\setminus J_1},y''\big)\PP_{J_1}(\rd y'')\Big]\PP_J(\rd y)=\int_{\kk^J}(P_{J_1}f)(x_{J^c},y)\,\PP_J(\rd y)=(P_JP_{J_1}f)(x).\qedhere
\]
\end{proof}
\begin{remark}\label{rem:split}
\textup{(i)} In Section~\ref{sec:FA} the lemma is used through the calculus (F0)--(F4) of Lemma~\ref{lem:FAfub}, for functions of polynomial growth; for those, $\mathcal D_f=\Omega$ and no exceptional null set occurs.
\textup{(ii)} A typical use of \textup{(b)} is the following ``conditional'' large deviation principle: if an event $\mathcal B$ is such that, for every fixed value $x'$ of the coordinates outside $J$, a large deviation bound for independent variables indexed by $J$
shows $\PP_J(\{y:(x',y)\in\mathcal B\})\le p$, then $\PP(\mathcal B)\le p$. No regular conditional probability is needed.
\end{remark}
\section{Entrywise bounds and the weak local law}\label{sec:weak}
From now on $N\ge N_0$, $\nu\in\mathcal E_N(\vth,\dpm,C_0)$ and $K$ in the range \eqref{eq:Krange} are fixed; $K_w,p_1,p_2,\fL,\Emax,\cG(z)$ are as in Definition~\ref{def:levels}, and (S1)--(S5) of Lemma~\ref{lem:levels} are in force
(in particular the two ledger inequalities \eqref{eq:ledgerA3}). Recall $z=E+\ii\eta$, $\kappa=||E|-2|$, $\varkappa=\sqrt{\kappa+\eta}$, the constants $c_*,C_*$ of Lemma~\ref{lem:mscsize}, $C_s=1+\sqrt{C_0}/\delta_{\min}$ of Lemma~\ref{lem:vecstab},
$\lambda_{\rm st}=1/(576C_s)$ of \eqref{eq:lambda1}, and $w=\avg v$, $w_0=\msc(1-\msc^2)$, $D=\max_i|v_i-w|$, $\Upsilon_{\max}=\max_i|\Upsilon_i|$ of Lemma~\ref{lem:stabnear}. We use the row form of (A): $\sum_j\sigma_{ij}^2=1$ for every $i$ (since $\sigma_{ij}=\sigma_{ji}$).

\subsection{Domains and the lattice}
\begin{definition}\label{def:domains}
\begin{align*}
 \bD^{\rm near}&:=\{z:\ |E|\le5,\ K^{C_1-1}N^{-1}\le\eta\le10\},\\
 \bD^{\rm far}&:=\{z:\ 4\le|E|\le\Emax+1,\ \tfrac12N^{-1/2}\le\eta\le10\},\\
 \mathbb L&:=(\bD^{\rm near}\cup\bD^{\rm far})\cap\big(N^{-4}\Z+\ii(10-N^{-4}\Z_{\ge 0})\big),\\
 |\mathbb L|&\le N^{10}.
\end{align*}
For $z=E+\ii\eta\in\mathbb L$ the \emph{column} above $z$ is ${\rm col}(z):=\{E+\ii\eta'\in\mathbb L:\ \eta'\ge\eta\}$ (the lattice points with the same $E$ between $z$ and $E+10\ii$), and
$\cG_{\rm col}(z):=\bigcap_{z'\in{\rm col}(z)}\cG(z')$. Thus $\PP(\cG_{\rm col}(z)^c)\le CN^{9}e^{-K^2}$, and $\cG_*:=\bigcap_{z\in\mathbb L}\cG(z)$ has $\PP(\cG_*^c)\le CN^{15}e^{-K^2}$.
\end{definition}
\noindent\emph{Justification of the claims in Definition~\ref{def:domains}.} A point of $\mathbb L$ is $E+\ii\eta$ with $E\in N^{-4}\Z$, $|E|\le\Emax+1$, and $\eta=10-kN^{-4}>0$ with an integer $k\ge0$.
There are at most $2(\Emax+1)N^4+1\le5\Emax N^4\le5N^{4+1/32}$ such $E$ (by (S2), (S4): $1\le\Emax\le K^{2\fa+1}\le N^{1/32}$) and at most $10N^4$ such $\eta$; so $|\mathbb L|\le50N^{8+1/32}\le N^{10}$.
\begin{itemize}[leftmargin=2em,itemsep=1pt]
\item \emph{Structure of a column.} Let $z=E+\ii\eta\in\mathbb L$, $\eta=10-k_*N^{-4}$, and $z_k:=E+\ii(10-kN^{-4})$ for $k=0,\dots,k_*$. The conditions defining $\bD^{\rm near}$ and $\bD^{\rm far}$ are of the form ``$E\in$ a set, $\eta_{\min}\le\eta\le10$''.
Hence if $z\in\bD^{\rm near}$ then $z_k\in\bD^{\rm near}$ for all $k\le k_*$, and likewise for $\bD^{\rm far}$. So ${\rm col}(z)=\{z_0,\dots,z_{k_*}\}$, $z_0=E+10\ii$, $z_{k_*}=z$, $|{\rm col}(z)|\le10N^4$, and all points of the column lie in each of the two domains that contains $z$.
\item \emph{Probabilities.} $\cG(z')$ is closed (Lemma~\ref{lem:good}), so $\cG_{\rm col}(z)$ and $\cG_*$ are closed. By (S5) and the union bound, $\PP(\cG_{\rm col}(z)^c)\le460\,\vth^{-1}N^8e^{-K^2}$ and $\PP(\cG_*^c)\le46\,\vth^{-1}N^{14}e^{-K^2}$.
\item \emph{Sizes.} For $z\in\bD^{\rm near}$: $N\eta\ge K^{C_1-1}\ge K^{12\fa+18}$ by \eqref{eq:ledgerA3}, $\eta\ge N^{-1}$, and $|z|\le\sqrt{125}$. For $z\in\bD^{\rm far}$: $N\eta\ge\frac12N^{1/2}$ and $|z|\le\Emax+11$.
\end{itemize}
\noindent For $z\in\bD^{\rm near}$ we use the control parameter
\begin{equation}\label{eq:Psi}
 \Psi(z):=\sqrt{\frac{\im\msc(z)+\Lambda(z)}{N\eta}}+\frac1{N\eta}.
\end{equation}
\subsection{One-point estimates}
The statements of this subsection are deterministic: they hold at every point of the indicated event, for one fixed $z$.
\begin{lemma}[Bounds under the a priori assumption, near regime]\label{lem:apriori}
Let $z\in\bD^{\rm near}$. On $\cG(z)\cap\{\Ld(z)+\Lo(z)\le\lambda_{\rm st}\}$, for all $i\ne j$:
\begin{align}
 &\tfrac3{52}\le|G_{ii}|\le\tfrac54,\qquad|G^{(i)}_{jj}|\le2,\qquad\Psi_m\le C\Psi,\qquad\Psi\le1,\label{eq:ap1}\\
 &\Lo\le C\fL\Psi,\qquad\max_i|X_i|\le C\fL\Psi,\qquad\max_i|A_i|\le C\fL\Psi^2,\qquad\Upsilon_{\max}\le C\fL\Psi,\label{eq:ap2}\\
 &|w^2-w_0w|\le|\avg X|+\mathfrak E_0\Big(\frac{\im\msc+\Lambda}{N\eta}+\frac1{(N\eta)^2}\Big)+\tfrac34\Lambda^2,\qquad\mathfrak E_0:=C\fL^2\ (\le K^{8\fa+7}),\label{eq:ap3}\\
 &D\le C\fL\Psi+C\Lambda^2,\qquad\Ld\le\Lambda+D.\label{eq:ap4}
\end{align}
The constants are explicit: \eqref{eq:ap1} holds with $C=14$; \eqref{eq:ap2} with the constants $35$, $14$, $245$, $259$ (in this order); $\mathfrak E_0:=6\cdot10^6\,C_s^2\,\fL^2$; \eqref{eq:ap4} reads $D\le1036\,C_s\fL\Psi+144\,C_s\Lambda^2$.
Moreover, on the same event,
\begin{equation}\tag{5.2$'$}\label{eq:ap1b}
 \Psi\le\frac2{\sqrt{N\eta}},\qquad\fL\Psi\le1,\qquad|w^2-w_0w|\le C_a\,\fL\Psi+\tfrac3{10}\Lambda^2,\qquad C_a:=3\cdot10^6\,C_s^2 .
\end{equation}
\end{lemma}
\begin{proof}
We work at a point of the event and write $u:=N\eta$. By Definition~\ref{def:domains}, \eqref{eq:ledgerA3} and (S2), $u\ge K^{12\fa+18}\ge4K^{8\fa+6}\ge4\fL^2\ge16$.

\emph{\eqref{eq:ap1} and the first two claims of \eqref{eq:ap1b}.} Since $\Ld\le\lambda_{\rm st}$, Lemma~\ref{lem:stabnear} gives $\frac3{52}\le|G_{ii}|\le\frac54$ and $\Lambda\le\Ld\le\lambda_{\rm st}$. By (I3) of Lemma~\ref{lem:schur} with $A=\{i\}$,
$|G^{(i)}_{jj}|\le|G_{jj}|+|G_{ji}||G_{ij}|/|G_{ii}|\le\frac54+\frac{52}3\lambda_{\rm st}^2\le2$. Next, $\im m\le\im\msc+|m-\msc|=\im\msc+\Lambda$, and by Lemma~\ref{lem:msc}(c) with $|z|\le\sqrt{125}$, $\im\msc\ge\eta/(|z|+1)^2\ge\eta/169$; hence
$N^{-1/2}\le13\sqrt{\im\msc/(N\eta)}$ and
\[
 \Psi_m=\sqrt{\frac{\im m}{N\eta}}+\frac1{N\eta}+\frac1{\sqrt N}\le14\Big(\sqrt{\frac{\im\msc+\Lambda}{N\eta}}+\frac1{N\eta}\Big)=14\,\Psi .
\]
Since $\im\msc\le|\msc|\le1$ and $\Lambda\le1$: $\Psi\le\sqrt{2/u}+1/u\le(\sqrt2+\frac14)u^{-1/2}\le2u^{-1/2}$ (as $u^{-1/2}\le\frac14$). So $\Psi\le\frac12\le1$ and $\fL\Psi\le2\fL u^{-1/2}\le1$.

\emph{\eqref{eq:ap2}.} By (S5) and \eqref{eq:ap1}: $|X_i|\le\fL\Psi_m\le14\fL\Psi$; $|A_i|\le\fL|G_{ii}|\Psi_m^2\le\frac54\cdot196\,\fL\Psi^2=245\,\fL\Psi^2$; $|G_{ij}|\le\fL|G_{ii}||G^{(i)}_{jj}|\Psi_m\le\frac54\cdot2\cdot14\,\fL\Psi=35\,\fL\Psi$;
and $\Upsilon_{\max}\le\max_i(|X_i|+|A_i|)\le(14+245\Psi)\fL\Psi\le259\,\fL\Psi$.

\emph{\eqref{eq:ap4}.} This is \eqref{eq:Dbound}: $D\le4C_s\Upsilon_{\max}+144C_s\Lambda^2\le1036\,C_s\fL\Psi+144C_s\Lambda^2$, and $\Ld\le\Lambda+D$.

\emph{\eqref{eq:ap3} and the last claim of \eqref{eq:ap1b}.} By \eqref{eq:quad} (in the explicit form stated in Lemma~\ref{lem:stabnear}) and $\avg\Upsilon=\avg A+\avg X$,
\[
 |w^2-w_0w|\le|\avg X|+|\avg A|+40C_s^2\Upsilon_{\max}^2+(90+72C_s)\Lambda^3 .
\]
Here $(90+72C_s)\Lambda^3\le162\,C_s\lambda_{\rm st}\Lambda^2\le\frac3{10}\Lambda^2$ because $\Lambda\le\lambda_{\rm st}=1/(576C_s)$ and $C_s\ge1$; and by \eqref{eq:ap2}, $|\avg A|\le245\fL\Psi^2$, $\Upsilon_{\max}^2\le259^2\fL^2\Psi^2$. Since $\fL\ge1$, $C_s\ge1$ and $245+40\cdot259^2\le3\cdot10^6$,
\[
 |w^2-w_0w|\le|\avg X|+3\cdot10^6\,C_s^2\fL^2\,\Psi^2+\tfrac3{10}\Lambda^2 .
\]
With $\Psi^2\le2\frac{\im\msc+\Lambda}{N\eta}+\frac2{(N\eta)^2}$ this is \eqref{eq:ap3} with $\mathfrak E_0=6\cdot10^6C_s^2\fL^2$ (and $\frac3{10}$ instead of $\frac34$).
For the last claim of \eqref{eq:ap1b} we use instead $|\avg X|\le14\fL\Psi$, $|\avg A|\le245\fL\Psi$ (as $\Psi\le1$) and $\fL^2\Psi^2\le\fL\Psi$ (as $\fL\Psi\le1$):
$|w^2-w_0w|\le(14+245+40\cdot259^2C_s^2)\fL\Psi+\frac3{10}\Lambda^2\le3\cdot10^6C_s^2\,\fL\Psi+\frac3{10}\Lambda^2$.
Finally, by (S2), $\mathfrak E_0\le6\cdot10^6C_s^2\cdot(300C_0^2)^2\log N\cdot K^{8\fa+4}\le K^{8\fa+7}$ for $N\ge N_0$, since $K^3\ge(\log N)^9$ (see the proof of (S2) for the bound on $\fL$).
\end{proof}

\begin{lemma}[Weak dichotomy, near regime]\label{lem:weakdich}
There are constants $U_w\ge1$ and $C$ (depending only on $C_s$, i.e.\ on the model parameters) such that for $z\in\bD^{\rm near}$, on $\cG(z)\cap\{\Ld+\Lo\le\lambda_{\rm st}\}$, with $\varkappa_w(z):=U_w\fL^{2/3}(N\eta)^{-1/3}$:
\begin{enumerate}[label=\textup{(\alph*)},leftmargin=2.2em,itemsep=0pt]
\item $\Lambda\le C\big(\varkappa+\fL^{2/3}(N\eta)^{-1/3}\big)$;
\item if $\varkappa\ge\varkappa_w$ and $\Lambda\le|w_0|/2$, then $\Lambda\le C\fL(\varkappa N\eta)^{-1/2}\le|w_0|/4$;
\item if $\varkappa<\varkappa_w$, or if $\varkappa\ge\varkappa_w$ and $\Lambda\le|w_0|/2$, then
$\Ld+\Lo\le K\fL(N\eta)^{-1/3}\le\lambda_{\rm st}/2$.
\end{enumerate}
\end{lemma}
\begin{proof}
We work at a point of the event, write $u:=N\eta$, and use Lemma~\ref{lem:apriori}. Two elementary facts will be used.

\smallskip\noindent\textbf{(E1)} \emph{If $x,b,c,e\ge0$ are reals with $x^2\le bx+c+e\sqrt x$, then $x\le\frac43b+\sqrt{4c/3}+e^{2/3}$.} Indeed, by Young's inequality ($\alpha\beta\le\frac14\alpha^4+\frac34\beta^{4/3}$ with $\alpha=\sqrt x$, $\beta=e$), $e\sqrt x\le\frac14x^2+\frac34e^{4/3}$; so
$x^2\le\frac43bx+\frac43c+e^{4/3}$, and the last statement of Lemma~\ref{lem:scalar} gives $x\le\frac43b+(\frac43c+e^{4/3})^{1/2}\le\frac43b+\sqrt{4c/3}+e^{2/3}$.

\smallskip\noindent\textbf{(E2)} \emph{$\im\msc\le C_*\varkappa$ and $\frac{c_*}{13}\varkappa\le|w_0|\le C_*\varkappa$ on $\bD^{\rm near}$.} The second statement is part of Lemma~\ref{lem:stabnear}. For the first: by \eqref{eq:imm}, either $\im\msc\le C_*\varkappa$, or $\im\msc\le C_*\eta/\varkappa\le C_*\varkappa$ because $\eta\le\kappa+\eta=\varkappa^2$.

\smallskip By \eqref{eq:ap1b}, (E2) and $\sqrt{\alpha+\beta}\le\sqrt\alpha+\sqrt\beta$,
\begin{equation}\tag{5.a}\label{eq:wd1}
 |w^2-w_0w|\le C_a\fL\Psi+\tfrac3{10}\Lambda^2,\qquad\Psi\le\sqrt{\frac{C_*\varkappa}u}+\sqrt{\frac\Lambda u}+\frac1u .
\end{equation}

\emph{(a).} Since $\Lambda=|w|$: $\Lambda^2\le|w_0|\Lambda+|w^2-w_0w|$, so by \eqref{eq:wd1}, $\frac7{10}\Lambda^2\le|w_0|\Lambda+C_a\fL\Psi$, i.e.\ $\Lambda^2\le b\Lambda+c+e\sqrt\Lambda$ with
\[
 b:=\tfrac{10}7|w_0|\le\tfrac{10}7C_*\varkappa,\qquad c:=\tfrac{10}7C_a\fL\Big(\sqrt{\frac{C_*\varkappa}u}+\frac1u\Big),\qquad e:=\tfrac{10}7C_a\fL u^{-1/2}.
\]
By (E1), using $\frac43\cdot\frac{10}7\le2$:
\[
 \Lambda\le2C_*\varkappa+\sqrt{2C_a}\,C_*^{1/4}\,\fL^{1/2}\varkappa^{1/4}u^{-1/4}+\sqrt{2C_a}\,\fL^{1/2}u^{-1/2}+(2C_a)^{2/3}\fL^{2/3}u^{-1/3}.
\]
Now $\fL^{1/2}\varkappa^{1/4}u^{-1/4}=\varkappa^{1/4}\big(\fL^{2/3}u^{-1/3}\big)^{3/4}\le\varkappa+\fL^{2/3}u^{-1/3}$ (weighted arithmetic--geometric mean inequality) and $\fL^{1/2}u^{-1/2}\le\fL^{2/3}u^{-1/3}$ (as $\fL u\ge1$).
This proves (a) with $C:=C_{(a)}:=2C_*+4C_aC_*$ (a crude upper bound for the sum of the four coefficients; recall $C_*,C_a\ge1$).

\emph{(b).} Assume $\varkappa\ge\varkappa_w$ and $\Lambda\le|w_0|/2$. Then $|w-w_0|\ge|w_0|/2$, so $\frac12|w_0|\Lambda\le|w^2-w_0w|$; by \eqref{eq:wd1} and $\frac3{10}\Lambda^2\le\frac3{20}|w_0|\Lambda$,
\[
 \tfrac7{20}|w_0|\Lambda\le C_a\fL\Psi .
\]
In $\Psi$ we use $\Lambda\le\frac12|w_0|\le\frac12C_*\varkappa$, so $\im\msc+\Lambda\le\frac32C_*\varkappa$ and $\Psi\le\sqrt{3C_*\varkappa/(2u)}+1/u$. With $|w_0|\ge\frac{c_*}{13}\varkappa$ (E2),
\[
 \Lambda\le\frac{20\cdot13\,C_a\fL}{7c_*\varkappa}\Big(\sqrt{\frac{3C_*\varkappa}{2u}}+\frac1u\Big)\le C_b\,\fL\Big(\frac1{\sqrt{\varkappa u}}+\frac1{\varkappa u}\Big),\qquad C_b:=\frac{38\,C_a}{c_*}\sqrt{\tfrac32C_*}\ .
\]
Since $\varkappa\ge\varkappa_w\ge u^{-1/3}$ (as $U_w,\fL\ge1$) we have $\varkappa u\ge u^{2/3}\ge1$, hence $(\varkappa u)^{-1}\le(\varkappa u)^{-1/2}$ and
\[
 \Lambda\le2C_b\,\fL\,(\varkappa u)^{-1/2}.
\]
This is the first inequality in (b), with $C=2C_b$. The second one, $2C_b\fL(\varkappa u)^{-1/2}\le\frac{c_*}{52}\varkappa\ (\le\frac14|w_0|)$, is equivalent to $\varkappa^{3/2}\ge\frac{104C_b}{c_*}\fL u^{-1/2}$, i.e.\ to
$\varkappa\ge(104C_b/c_*)^{2/3}\fL^{2/3}u^{-1/3}$. It holds because $\varkappa\ge\varkappa_w$, if we define
\[
 U_w:=\max\Big\{1,\ \Big(\frac{104\,C_b}{c_*}\Big)^{2/3}\Big\}.
\]

\emph{(c).} \emph{First case: $\varkappa<\varkappa_w$.} By (a), $\Lambda\le C_{(a)}(U_w+1)\fL^{2/3}u^{-1/3}$.
\emph{Second case: $\varkappa\ge\varkappa_w$ and $\Lambda\le|w_0|/2$.} By (b) and $\varkappa\ge\varkappa_w\ge\fL^{2/3}u^{-1/3}$: $\ \Lambda\le2C_b\fL\,\varkappa^{-1/2}u^{-1/2}\le2C_b\,\fL^{2/3}u^{-1/3}$.
So in both cases, with $C_c:=\max\{C_{(a)}(U_w+1),2C_b\}$ and using $\fL\ge1$,
\[
 \Lambda\le C_c\,\fL^{2/3}u^{-1/3}\le C_c\,\fL\,u^{-1/3}.
\]
By \eqref{eq:ap2}, \eqref{eq:ap4}, \eqref{eq:ap1b}, and $144C_s\Lambda\le144C_s\lambda_{\rm st}=\frac14$,
\[
 \Lo\le35\fL\Psi\le70\,\fL u^{-1/2},\qquad\Ld\le\Lambda+D\le\tfrac54\Lambda+1036\,C_s\fL\Psi\le\tfrac54\Lambda+2072\,C_s\,\fL u^{-1/2}.
\]
Since $u^{-1/2}\le u^{-1/3}$,
\[
 \Ld+\Lo\le\big(\tfrac54C_c+2072\,C_s+70\big)\fL u^{-1/3}\le K\fL u^{-1/3}
\]
for $N\ge N_0$ (by (S1), with $K_0\ge\frac54C_c+2072C_s+70$). Finally $K\fL u^{-1/3}\le\lambda_{\rm st}/2$ is equivalent to $u\ge(2K\fL/\lambda_{\rm st})^3=(1152C_s)^3K^3\fL^3$, and by (S2), \eqref{eq:ledgerA3} and (S1) (with $K_0\ge(1152C_s)^{1/2}$)
\[
 (1152C_s)^3K^3\fL^3\le(1152C_s)^3K^{12\fa+12}\le K^{12\fa+18}\le K^{C_1-1}\le u .\qedhere
\]
\end{proof}

\begin{lemma}[Initial estimate]\label{lem:initial}
For every $z$ with $\eta=10$ (any $E\in\R$): on $\cG(z)$, $\ \max_i|G_{ii}|\le\frac1{10}$ and $\Ld+\Lo\le4\fL N^{-1/2}$.
\end{lemma}
\begin{proof}
By Lemma~\ref{lem:basic}(a) (for $H$ and for its minors), $|G_{ii}|\le\eta^{-1}=\frac1{10}$, $|G^{(i)}_{jj}|\le\frac1{10}$, and $\im m\le\max_i|G_{ii}|\le\frac1{10}$. Hence
\[
 \Psi_m=\sqrt{\frac{\im m}{10N}}+\frac1{10N}+\frac1{\sqrt N}\le\Big(\frac1{10}+\frac1{10}+1\Big)N^{-1/2}=\tfrac65N^{-1/2}.
\]
By (S5): $|X_i|\le\frac65\fL N^{-1/2}$; $|A_i|\le\fL\cdot\frac1{10}\cdot\frac{36}{25}N^{-1}\le\frac1{10}\fL N^{-1/2}$ (as $N\ge3$); $\Lo\le\fL\cdot\frac1{100}\cdot\frac65N^{-1/2}$. So $\Upsilon_{\max}\le\frac{13}{10}\fL N^{-1/2}$.
By the second identity in \eqref{eq:sce1}, $|\msc|\le1$ (Lemma~\ref{lem:msc}(b)) and $\sum_j\sigma_{ij}^2=1$,
\[
 |v_i|=|\msc|\,|G_{ii}|\,\Big|\sum_j\sigma_{ij}^2v_j-\Upsilon_i\Big|\le\tfrac1{10}\big(\Ld+\Upsilon_{\max}\big)\quad\text{for all }i,\qquad\text{so}\quad\Ld\le\tfrac19\Upsilon_{\max}\le\tfrac3{20}\fL N^{-1/2}.
\]
Altogether $\Ld+\Lo\le\frac15\fL N^{-1/2}\le4\fL N^{-1/2}$. No a priori assumption and no stability analysis is needed.
\end{proof}

\begin{lemma}[One step, far regime]\label{lem:farstep}
Let $z\in\bD^{\rm far}$ and $\lambda_f(z):=\frac1{2(|z|+3)}$. On $\cG(z)\cap\{\max_i|G_{ii}|\le1,\ \Lo\le\lambda_f(z)\}$:
\begin{gather*}
 \max_i|G_{ii}|\le\tfrac12,\qquad\frac1{|G_{ii}|}\le|z|+2,\\
 \Ld\le\fL\Psi_m\le6\fL(N\eta)^{-1/2},\qquad\Lo\le\tfrac12\fL\Psi_m\le3\fL(N\eta)^{-1/2}\le\tfrac12\lambda_f(z).
\end{gather*}
\end{lemma}
\begin{proof}
Let $u:=N\eta\ge\frac12N^{1/2}\ge1$. Since $\im m\le\max_i|G_{ii}|\le1$, $u^{-1}\le u^{-1/2}$ and $N^{-1/2}=\eta^{1/2}u^{-1/2}\le\sqrt{10}\,u^{-1/2}$,
\[
 \Psi_m\le\big(1+1+\sqrt{10}\big)u^{-1/2}\le6u^{-1/2}\le6\sqrt2\,N^{-1/4}.
\]
By (S2), (S4): $\fL\le K^{4\fa+3}\le N^{1/32}$, so $2\fL\Psi_m\le17N^{1/32-1/4}\le\frac14$ for $N\ge N_0$ ((S4) with $K_0\ge68$); in particular $\Psi_m\le1$. By (S5),
\[
 \Upsilon_{\max}\le\max_i\big(|X_i|+|A_i|\big)\le\fL\Psi_m+\fL\Psi_m^2\le2\fL\Psi_m\le\tfrac14 .
\]
Since $|E|\ge4$ and $\max_i|G_{ii}|\le1$, Lemma~\ref{lem:stabfar} applies: $\Ld\le\frac12\Upsilon_{\max}\le\fL\Psi_m\le6\fL u^{-1/2}$ and $\Ld\le\frac18$; with $|\msc|\le\frac13$ (Lemma~\ref{lem:msc}(g)), $|G_{ii}|\le\frac13+\frac18\le\frac12$; and
$1/|G_{ii}|\le|z|+1+\Upsilon_{\max}\le|z|+2$. By (I3) of Lemma~\ref{lem:schur} with $A=\{i\}$, for $j\ne i$,
\[
 |G^{(i)}_{jj}|\le|G_{jj}|+\frac{|G_{ji}||G_{ij}|}{|G_{ii}|}\le\frac12+\lambda_f(z)^2(|z|+2)\le\frac12+\frac1{4(|z|+3)}\le1 .
\]
Then (S5) gives $|G_{ij}|\le\fL\cdot\frac12\cdot1\cdot\Psi_m\le3\fL u^{-1/2}$ for all $i\ne j$. Finally, $|z|+3\le\Emax+14\le15\Emax\le15K^{2\fa+1}$ by (S2), so
\[
 3\fL u^{-1/2}\cdot\frac2{\lambda_f(z)}=12(|z|+3)\fL u^{-1/2}\le12\cdot15\cdot\sqrt2\;K^{6\fa+4}N^{-1/4}\le255\,N^{1/32-1/4}\le1
\]
for $N\ge N_0$, by (S4) with $K_0\ge255$. This is $\Lo\le\frac12\lambda_f(z)$.
\end{proof}
\subsection{The continuity argument}
\begin{theorem}[Weak local law on the lattice; deterministic form]\label{thm:weak}
For every $z\in\mathbb L$, on the event $\cG_{\rm col}(z)$:
\begin{enumerate}[label=\textup{(\alph*)},leftmargin=2.2em,itemsep=0pt]
\item if $z\in\bD^{\rm near}$: $\ \Ld(z)+\Lo(z)\le K\fL(N\eta)^{-1/3}\le\lambda_{\rm st}/2$; and if moreover $\varkappa(z)\ge\varkappa_w(z)$, then $\Lambda(z)\le|w_0(z)|/4$;
\item if $z\in\bD^{\rm far}$: $\ \max_i|G_{ii}(z)|\le\frac12$ and $\Ld(z)+\Lo(z)\le9\fL(N\eta)^{-1/2}$;
\item in both cases the event
\begin{equation}\label{eq:Xiev}
 \Xi(z):=\Big\{\Ld(z)\le\tfrac12|\msc(z)|,\ \ \Lo(z)\le\frac{|\msc(z)|}{16\,p_1}\Big\}
\end{equation}
holds. Hence $\PP(\Xi(z)^c)\le CN^9e^{-K^2}$ for every $z\in\mathbb L$.
\end{enumerate}
\end{theorem}
\begin{proof}
Fix $z=E+\ii\eta\in\mathbb L$ and let ${\rm col}(z)=\{z_0,\dots,z_{k_*}\}$, $z_k=E+\ii\eta_k$, $\eta_k=10-kN^{-4}$, $z_{k_*}=z$, as in the justification of Definition~\ref{def:domains}. We fix a point of
$\cG_{\rm col}(z)=\bigcap_{k\le k_*}\cG(z_k)$; all quantities below are evaluated at this point, and the argument is deterministic. We shall use the following Lipschitz bounds for consecutive points of the column ($|z_k-z_{k+1}|=N^{-4}$).

\smallskip\noindent\textbf{(Lip)} \emph{Let $0\le k<k_*$ and $\eta_0:=\eta_{k+1}$ (the smaller of the two imaginary parts). Then
\[
 |F(z_k)-F(z_{k+1})|\le\big(\eta_0^{-2}+\eta_0^{-1}+1\big)N^{-4}\ \ \text{for }F\in\{\Lambda,\Ld,\Lo\},\qquad|G_{ii}(z_k)-G_{ii}(z_{k+1})|\le\eta_0^{-2}N^{-4},
\]
and $|w_0(z_k)-w_0(z_{k+1})|\le4(1+\eta_0^{-1})N^{-4}$.} The first two bounds are Lemma~\ref{lem:basic}(c). For the third, write $m=\msc(z_k)$, $m'=\msc(z_{k+1})$; then
$w_0(z_k)-w_0(z_{k+1})=(m-m')\big(1-m^2-mm'-m'^2\big)$, $|m|,|m'|\le1$, and $|m-m'|\le(1+\eta_0^{-1})N^{-4}$ by Lemma~\ref{lem:msc}(f).

\medskip\emph{Near regime: proof of \textup{(a)}.} Let $z\in\bD^{\rm near}$; then $z_k\in\bD^{\rm near}$ for all $k\le k_*$, in particular $\eta_k\ge N^{-1}$, and (Lip) gives, for $N\ge13$ (see (S1)),
\begin{equation}\tag{5.b}\label{eq:lipnear}
\begin{aligned}
 |F(z_k)-F(z_{k+1})|&\le(N^2+N+1)N^{-4}\le2N^{-2}\qquad(F\in\{\Lambda,\Ld,\Lo\}),\\
 |w_0(z_k)-w_0(z_{k+1})|&\le4(N+1)N^{-4}\le N^{-2}.
\end{aligned}
\end{equation}
We prove by induction on $k\in\{0,\dots,k_*\}$ the statement
\[
 \mathcal H_k:\qquad\Ld(z_k)+\Lo(z_k)\le K\fL(N\eta_k)^{-1/3}\le\frac{\lambda_{\rm st}}2,\qquad\text{and}\quad\Lambda(z_k)\le\frac{|w_0(z_k)|}4\ \text{ if }\ \varkappa(z_k)\ge\varkappa_w(z_k).
\]
Recall that $K\fL(N\eta_k)^{-1/3}\le\lambda_{\rm st}/2$ holds for every point of $\bD^{\rm near}$ (end of the proof of Lemma~\ref{lem:weakdich}), and that $|w_0(z_k)|\ge\frac{c_*}{13}\varkappa(z_k)\ge\frac{c_*}{13}\eta_k^{1/2}$.

\emph{$\mathcal H_0$.} Here $\eta_0=10$. By Lemma~\ref{lem:initial} (on $\cG(z_0)$), $\Ld+\Lo\le4\fL N^{-1/2}\le K\fL(10N)^{-1/3}$ (as $K\ge9$), and $\Lambda\le\Ld\le4\fL N^{-1/2}\le4N^{-15/32}\le\frac{c_*}{52}\sqrt{10}\le\frac14|w_0(z_0)|$ for $N\ge N_0$
(by (S2), (S4), $\fL\le N^{1/32}$).

\emph{$\mathcal H_k\Rightarrow\mathcal H_{k+1}$.} Assume $\mathcal H_k$, $k<k_*$. By \eqref{eq:lipnear}, at $z_{k+1}$,
\[
 \Ld(z_{k+1})+\Lo(z_{k+1})\le\frac{\lambda_{\rm st}}2+4N^{-2}\le\lambda_{\rm st}
\]
for $N\ge N_0$ (namely $N^2\ge8/\lambda_{\rm st}$; this follows from (S4) with $K_0\ge8/\lambda_{\rm st}$). Since our point lies in $\cG(z_{k+1})$, we are on the event $\cG(z_{k+1})\cap\{\Ld+\Lo\le\lambda_{\rm st}\}$ of Lemma~\ref{lem:weakdich} at $z_{k+1}\in\bD^{\rm near}$.

\emph{Case 1: $\varkappa(z_{k+1})<\varkappa_w(z_{k+1})$.} Lemma~\ref{lem:weakdich}(c) gives $\Ld+\Lo\le K\fL(N\eta_{k+1})^{-1/3}$ at $z_{k+1}$; the second part of $\mathcal H_{k+1}$ is void.

\emph{Case 2: $\varkappa(z_{k+1})\ge\varkappa_w(z_{k+1})$.} The function $\eta\mapsto\varkappa(E+\ii\eta)\,(N\eta)^{1/3}=\sqrt{\kappa+\eta}\,(N\eta)^{1/3}$ is increasing ($\kappa$ depends only on $E$), and $\varkappa\ge\varkappa_w$ means $\varkappa\,(N\eta)^{1/3}\ge U_w\fL^{2/3}$. Since $\eta_k>\eta_{k+1}$,
also $\varkappa(z_k)\ge\varkappa_w(z_k)$, so $\mathcal H_k$ gives $\Lambda(z_k)\le|w_0(z_k)|/4$. By \eqref{eq:lipnear},
\[
 \Lambda(z_{k+1})\le\Lambda(z_k)+2N^{-2}\le\tfrac14|w_0(z_k)|+2N^{-2}\le\tfrac14|w_0(z_{k+1})|+3N^{-2}\le\tfrac12|w_0(z_{k+1})|,
\]
where the last step uses $|w_0(z_{k+1})|\ge\frac{c_*}{13}\eta_{k+1}^{1/2}\ge\frac{c_*}{13}N^{-1/2}$ and $N^{3/2}\ge156/c_*$ ($N\ge N_0$). So the hypotheses of Lemma~\ref{lem:weakdich}(b) and of the second alternative of (c) hold at $z_{k+1}$, and we get
$\Lambda(z_{k+1})\le|w_0(z_{k+1})|/4$ and $\Ld+\Lo\le K\fL(N\eta_{k+1})^{-1/3}$ at $z_{k+1}$.

In both cases $\mathcal H_{k+1}$ holds. This proves $\mathcal H_{k_*}$, which is (a).

\medskip\emph{Far regime: proof of \textup{(b)}.} Let $z\in\bD^{\rm far}$; then $z_k\in\bD^{\rm far}$ for all $k\le k_*$, in particular $\eta_k\ge\frac12N^{-1/2}$, so $\eta_k^{-2}\le4N$ and (Lip) gives
\begin{equation}\tag{5.c}\label{eq:lipfar}
 |G_{ii}(z_k)-G_{ii}(z_{k+1})|\le4N^{-3},\qquad|\Lo(z_k)-\Lo(z_{k+1})|\le(4N+2N^{1/2}+1)N^{-4}\le5N^{-3}.
\end{equation}
We prove by induction on $k$ the statement
\[
 \mathcal F_k:\qquad\max_i|G_{ii}(z_k)|\le1\qquad\text{and}\qquad\Lo(z_k)\le\lambda_f(z_k).
\]
\emph{$\mathcal F_0$.} By Lemma~\ref{lem:initial}, $\max_i|G_{ii}(z_0)|\le\frac1{10}$ and $\Lo(z_0)\le4\fL N^{-1/2}\le4N^{-15/32}$, while $\lambda_f(z_0)=\frac1{2(|z_0|+3)}\ge\frac1{30K^{2\fa+1}}\ge\frac1{30}N^{-1/32}$ (as in the proof of Lemma~\ref{lem:farstep}); so $\Lo(z_0)\le\lambda_f(z_0)$ for $N\ge N_0$.

\emph{$\mathcal F_k\Rightarrow\mathcal F_{k+1}$.} Assume $\mathcal F_k$, $k<k_*$. Our point lies in $\cG(z_k)$, so Lemma~\ref{lem:farstep} at $z_k$ gives $\max_i|G_{ii}(z_k)|\le\frac12$ and $\Lo(z_k)\le\frac12\lambda_f(z_k)$. Since $|z_{k+1}|\le|z_k|$ (same $E$, smaller $\eta$), $\lambda_f(z_{k+1})\ge\lambda_f(z_k)$.
By \eqref{eq:lipfar}, $\max_i|G_{ii}(z_{k+1})|\le\frac12+4N^{-3}\le1$ and
\[
 \Lo(z_{k+1})\le\tfrac12\lambda_f(z_k)+5N^{-3}\le\tfrac12\lambda_f(z_{k+1})+5N^{-3}\le\lambda_f(z_{k+1}),
\]
because $\lambda_f(z_{k+1})\ge\frac1{30}N^{-1/32}\ge10N^{-3}$. This is $\mathcal F_{k+1}$.

So $\mathcal F_{k_*}$ holds, and Lemma~\ref{lem:farstep} at $z=z_{k_*}$ (on $\cG(z)$) gives $\max_i|G_{ii}(z)|\le\frac12$, $\Ld(z)\le6\fL(N\eta)^{-1/2}$ and $\Lo(z)\le3\fL(N\eta)^{-1/2}$. This is (b).

\medskip\emph{Proof of \textup{(c)}.} \emph{Near regime.} By Lemma~\ref{lem:mscsize}, $|\msc(z)|\ge\frac1{13}$. By (a), $\Ld\le\lambda_{\rm st}/2\le\frac1{1152}\le\frac12|\msc|$. By (a), (S2), (S3) and \eqref{eq:ledgerA3},
\[
 \Lo\cdot\frac{16p_1}{|\msc|}\le K\fL(N\eta)^{-1/3}\cdot16\cdot2K^{3/2}\cdot13\le416\,K^{4\fa+\frac{11}2}\big(K^{12\fa+18}\big)^{-1/3}=416\,K^{-1/2}\le1
\]
for $N\ge N_0$ ((S1) with $K_0\ge416^2$).

\emph{Far regime.} By Lemma~\ref{lem:msc}(c), $|\msc(z)|\ge\frac1{|z|+1}\ge\frac1{15K^{2\fa+1}}$. By (b) (more precisely by the two bounds at the end of its proof), $N\eta\ge\frac12N^{1/2}$, (S2), (S3), (S4),
\begin{align*}
 \frac{2\Ld}{|\msc|}&\le12\sqrt2\,\fL N^{-1/4}\cdot15K^{2\fa+1}\le255\,K^{6\fa+4}N^{-1/4}\le1,\\
 \Lo\cdot\frac{16p_1}{|\msc|}&\le3\sqrt2\,\fL N^{-1/4}\cdot32K^{3/2}\cdot15K^{2\fa+1}\le2037\,K^{6\fa+\frac{11}2}N^{-1/4}\le1,
\end{align*}
for $N\ge N_0$ ((S4) with $K_0\ge2037$: $K^{6\fa+6}\le N^{1/32}$ and $2037N^{1/32}\le N^{1/8}$).

\emph{Probability.} $\Xi(z)$ is closed, hence measurable (the functions $\Ld(z)$, $\Lo(z)$ are continuous on $\Omega$, Lemma~\ref{lem:weyl}(e)). We have shown $\cG_{\rm col}(z)\subset\Xi(z)$, so by Definition~\ref{def:domains},
$\PP(\Xi(z)^c)\le\PP(\cG_{\rm col}(z)^c)\le460\,\vth^{-1}N^8e^{-K^2}\le CN^9e^{-K^2}$.
\end{proof}

\begin{remark}
Theorem~\ref{thm:weak} corresponds to Theorem~3.1 of the original (errors $(N\eta)^{-1/3}$, $(N\eta)^{-1/2}$ up to powers of $K$). It is used below only through the a priori event $\Xi(z)$ and the bounds of Lemma~\ref{lem:apriori}.
The induction is deterministic: probability enters only through $\PP(\cG(z)^c)\le CN^5e^{-K^2}$ and a union bound over ${\rm col}(z)$.
\end{remark}
\begin{remark}[What Section~\ref{sec:weak} uses from the exponent ledger]\label{rem:ledgerA3}
Only \eqref{eq:ledgerA3}: $C_1-1\ge12\fa+18$ (used for $N\eta\ge4\fL^2$ in Lemma~\ref{lem:apriori}, for $K\fL(N\eta)^{-1/3}\le\lambda_{\rm st}/2$ in Lemma~\ref{lem:weakdich}(c), and for $\Lo\le|\msc|/(16p_1)$ in Theorem~\ref{thm:weak}(c)), and $c_0(6\fa+6)\le\frac1{32}$
(used for $|\mathbb L|\le N^{10}$ and in the far regime). Both hold for the values \eqref{eq:ledger0} (Appendix~\ref{app:ledger}).
\end{remark}

\section{Fluctuation averaging}\label{sec:FA}
This section replaces Section~7 of the original (Lemma~4.1 there). It is self-contained given the following statements of
Sections~\ref{sec:resid} and~\ref{sec:prob}: Lemma~\ref{lem:msc}(a),(c), Lemma~\ref{lem:basic}(a),(b), Lemma~\ref{lem:schur}\,(I1),(I4),(I5), Lemma~\ref{lem:minor},
Lemma~\ref{lem:moments}, Lemma~\ref{lem:LDmom} and Lemma~\ref{lem:split} (splitting of the product law); Corollary~\ref{cor:FA} also uses
Theorem~\ref{thm:weak}(c) and the relations \eqref{eq:levelrel}.

Throughout the section, $N\ge3$, $\nu\in\mathcal E_N(\vth,\dpm,C_0)$ (normalised as in Convention~\ref{conv:const}: $\vth\le1$, $C_0\ge1$, $\fa=1/\vth$), and
$z=E+\ii\eta$ with $\eta>0$ is fixed; $G$, $G^{(T)}$, $m$, $m^{(T)}$, $\Ld$, $\Lo$ are as in \eqref{eq:green}, \eqref{eq:defLambda} and Section~\ref{sec:resid},
and (Definition~\ref{def:XA})
\[
 Z_k=\sum^{(k)}_{i,j}h_{ki}G^{(k)}_{ij}h_{jk}-\sum^{(k)}_i\sigma_{ki}^2G^{(k)}_{ii},\qquad X_k=h_{kk}-Z_k .
\]
We shall see (Lemma~\ref{lem:FAstruct}(d)) that $X_k=Q_k(1/G_{kk})$, where $Q_k=1-P_k$ and $P_k$ integrates out the $k$-th row.

\subsection{Statements}

\begin{theorem}[Fluctuation averaging]\label{thm:FA}
Let $c_F:=8\fa+7$. There are constants $C_F$ and $N_0$ (depending only on the model parameters) such that the following holds for $N\ge N_0$, $\nu\in\mathcal E_N(\vth,\dpm,C_0)$.
Let $z=E+\ii\eta$ with $|z|\le N$, let $p$ be an even integer with
\begin{equation}\label{eq:FAhyp}
 \log N\le p\le N^{1/(40\fa+40)},\qquad p\ge p_0,\qquad N^{-1}\le\eta\le10,\qquad N\eta\ge p^{12\fa+11},
\end{equation}
where $p_0:=25\,e\,C_0\,C_m\,13^{2\fa+\frac32}$ with the constant $C_m$ of \eqref{eq:Cm} (it depends only on the model parameters),
and let $t_1,\dots,t_N\in\C$ be deterministic with $|t_k|\le1/N$. Let
$\Xi_p(z):=\{\Ld(z)\le\frac12|\msc(z)|,\ \Lo(z)\le\frac{|\msc(z)|}{16p}\}$ and
\begin{equation}\label{eq:Phistar}
 \Phi_*(z)^2:=\frac1N+\frac1{N\eta}\Big(\im m(z)+\frac{p}{N\eta}\Big)\qquad(\text{a random variable; it depends on }p).
\end{equation}
Then
\begin{equation}\label{eq:FA}
 \EE\Big|\sum_{k=1}^Nt_kX_k(z)\Big|^p\le\big(C_Fp^{c_F}\big)^p\ \EE\big[\Phi_*(z)^{2p}\big]+N^{2p}\,\PP\big(\Xi_p(z)^c\big)^{1/2}.
\end{equation}
\end{theorem}
\noindent The right-hand side is \emph{random-parameter} (it contains $\im m=\im\msc+O(\Lambda)$ under the expectation); this is what makes the one-step argument of Section~\ref{sec:strong} possible.
The statement is uniform in $E$ (near and far regime) and the $p$-dependence is explicit. The constants are given in \eqref{eq:CF} and \eqref{eq:N0FA}.

\begin{corollary}[Form used below]\label{cor:FA}
Let $N\ge N_0$, let $K$ be as in \eqref{eq:Krange}, $z\in\mathbb L$ and $p\in\{p_1,p_2\}$ (Definitions~\ref{def:levels}, \ref{def:domains}). Then, with $\Phi_*$ defined by \eqref{eq:Phistar} for this $p$,
\[
 \norm{\avg{X}(z)}_p\le\mathfrak F\,\norm{\Phi_*(z)^2}_p+N^{-8},\qquad\mathfrak F:=C_Fp_1^{c_F}\ \ (\le K^{12\fa+11}).
\]
\end{corollary}

\paragraph{Organisation.} Section~\ref{sec:FAfub} fixes the notation for integration over rows (Lemma~\ref{lem:split}) and the algebraic structure obtained by conditioning on a minor (Lemma~\ref{lem:FAstruct}).
Section~\ref{sec:FAmom} contains the conditional moment bounds (Lemma~\ref{lem:FAmom}), the deterministic comparison of the a priori event with its minor versions (Lemma~\ref{lem:FAevent}) and a crude bound.
Section~\ref{sec:FAtwo} proves the case $p=2$ (Proposition~\ref{prop:FAtwo}): it contains every idea of the general proof and no combinatorics.
Sections~\ref{sec:FAexp}--\ref{sec:FAproof} prove Theorem~\ref{thm:FA}: the finite expansion (Lemma~\ref{lem:FAexp}), the vanishing of terms with an unvisited lone index (Lemma~\ref{lem:FAvanish}),
the size of the terms (Lemma~\ref{lem:FAsize}), the bound for one index tuple (Lemma~\ref{lem:FAtuple}), the counting of tuples (Lemma~\ref{lem:FAcount}). Corollary~\ref{cor:FA} is proved in Section~\ref{sec:FAcor}.

\subsection{Integration over rows; the block matrix given a minor}\label{sec:FAfub}
\paragraph{Partial integration.} Recall $\Omega=\kk^{\mathcal P_N}$, $\PP=\bigotimes_{(i,j)\in\mathcal P_N}\nu_{ij}$ (Definition~\ref{def:ens}). A \emph{coordinate} is an element of $\mathcal P_N$; we also write $\{i,j\}$ for the coordinate $(\min(i,j),\max(i,j))$.
For $\mathcal Q\subset\mathcal P_N$ let $\mathcal Q^c:=\mathcal P_N\setminus\mathcal Q$, $\PP_{\mathcal Q}:=\bigotimes_{(i,j)\in\mathcal Q}\nu_{ij}$ on $\kk^{\mathcal Q}$, and write $x=(x_{\mathcal Q^c},x_{\mathcal Q})$ for $x\in\Omega$.
By Lemma~\ref{lem:split}, $x\mapsto(x_{\mathcal Q^c},x_{\mathcal Q})$ is a measurable equivalence $\Omega\to\kk^{\mathcal Q^c}\times\kk^{\mathcal Q}$ which maps $\PP$ to $\PP_{\mathcal Q^c}\otimes\PP_{\mathcal Q}$; the same holds
for $\PP_{\mathcal Q}$ and a subset $\mathcal Q_1\subset\mathcal Q$ (it is the same statement for the index set $\mathcal Q$).
A function $f$ on $\Omega$ \emph{does not depend on $\mathcal Q$} if $f(x)=f(x')$ whenever $x_{\mathcal Q^c}=x'_{\mathcal Q^c}$. For $T\subset\{1,\dots,N\}$,
\[
 \mathcal Q(T):=\big\{(i,j)\in\mathcal P_N:\ \{i,j\}\cap T\ne\emptyset\big\}\qquad(\text{the coordinates in the rows }T).
\]
Let $|x|_1:=\sum_{(i,j)\in\mathcal P_N}|x_{ij}|$, and let $\mathcal A$ be the set of all measurable $f:\Omega\to\C$ for which there are $c,m\in[0,\infty)$ with
\begin{equation}\label{eq:classA}
 |f(x)|\le c\,(1+|x|_1)^m\qquad\text{for all }x\in\Omega .
\end{equation}
For $f\in\mathcal A$, or $f:\Omega\to[0,\infty]$ measurable, and for real $q\ge1$ put
\begin{equation}\label{eq:PQ}
 (P_{\mathcal Q}f)(x):=\int_{\kk^{\mathcal Q}}f(x_{\mathcal Q^c},y)\,\PP_{\mathcal Q}(\rd y),\qquad\norm{f}_{q,\mathcal Q}(x):=\Big(\big(P_{\mathcal Q}|f|^q\big)(x)\Big)^{1/q},
\end{equation}
and $P_T:=P_{\mathcal Q(T)}$, $\norm{\cdot}_{q,T}:=\norm{\cdot}_{q,\mathcal Q(T)}$, $P_k:=P_{\{k\}}$, $Q_k:=1-P_k$. This is the operator $P_T$ of Section~\ref{sec:prob}; $P_{\mathcal Q}f$ is a version of the conditional expectation of $f$
given the coordinates in $\mathcal Q^c$, but we shall only use the properties (F0)--(F5) below, which are statements about iterated integrals.

\begin{lemma}[Calculus of partial integration]\label{lem:FAfub}
Let $\mathcal Q_1,\mathcal Q\subset\mathcal P_N$, $f,g,f_1,\dots,f_n\in\mathcal A$.
\begin{enumerate}[label=\textup{(F\arabic*)},leftmargin=2.6em,itemsep=1pt,start=0]
\item $\mathcal A$ is an algebra, closed under complex conjugation and under $f\mapsto|f|^q$ ($q>0$); it contains the functions $h_{ij}$; $\EE|f|^q<\infty$ for all $q<\infty$.
For every $x\in\Omega$ the integral $(P_{\mathcal Q}f)(x)$ converges absolutely; $P_{\mathcal Q}f\in\mathcal A$; $P_{\mathcal Q}f$ does not depend on $\mathcal Q$; and if $f$ does not depend on $\mathcal Q_1$, then $P_{\mathcal Q}f$ does not depend on $\mathcal Q_1$.
Moreover $P_{\mathcal Q}\ol f=\ol{P_{\mathcal Q}f}$.
\item \emph{(Fubini)} $\EE f=\EE[P_{\mathcal Q}f]$, and if $\mathcal Q_1\subset\mathcal Q$ then $P_{\mathcal Q}P_{\mathcal Q_1}f=P_{\mathcal Q}f$. In particular $P_{\mathcal Q}(f-P_{\mathcal Q}f)=0$, i.e.\ $P_kQ_kf=0$.
\item If $g$ does not depend on $\mathcal Q$, then $P_{\mathcal Q}(gf)=g\,P_{\mathcal Q}f$ and $P_{\mathcal Q}g=g$. If $\Xi\subset\Omega$ is measurable and $\bone_\Xi$ does not depend on $\mathcal Q$, then $\EE[\bone_\Xi f]=\EE[\bone_\Xi P_{\mathcal Q}f]$.
\item \emph{(H\"older, Minkowski, Jensen)} If $q,q_1,\dots,q_n\in[1,\infty)$ and $\sum_i1/q_i=1/q$, then $\norm{f_1\cdots f_n}_{q,\mathcal Q}\le\prod_{i}\norm{f_i}_{q_i,\mathcal Q}$ pointwise on $\Omega$; in particular ($q=1$) $|P_{\mathcal Q}(f_1\cdots f_n)|\le\prod_{i}\norm{f_i}_{q_i,\mathcal Q}$.
For $q\ge1$: $\norm{f+g}_{q,\mathcal Q}\le\norm f_{q,\mathcal Q}+\norm g_{q,\mathcal Q}$; $\norm{f}_{q,\mathcal Q}\le\norm{f}_{q',\mathcal Q}$ if $q\le q'$; and if $\mathcal Q_1\subset\mathcal Q$ then $\norm{P_{\mathcal Q_1}f}_{q,\mathcal Q}\le\norm{f}_{q,\mathcal Q}$, hence $\norm{Q_kf}_{q,\mathcal Q}\le2\norm f_{q,\mathcal Q}$ if $\mathcal Q(\{k\})\subset\mathcal Q$.
\item \emph{(Freezing)} For fixed $x_{\mathcal Q^c}$, the coordinate functions $y\mapsto y_{ij}$, $(i,j)\in\mathcal Q$, are independent random variables on the probability space $(\kk^{\mathcal Q},\PP_{\mathcal Q})$ with laws $\nu_{ij}$,
and $\norm{f}_{q,\mathcal Q}(x)$ is the $L^q(\PP_{\mathcal Q})$-norm of $y\mapsto f(x_{\mathcal Q^c},y)$.
\item \emph{(Irrelevant coordinates)} If $f$ does not depend on $\mathcal Q_1$ and does not depend on $\mathcal Q_2$ ($\mathcal Q_2\subset\mathcal P_N$), then $f$ does not depend on $\mathcal Q_1\cup\mathcal Q_2$.
If $\mathcal Q_1\subset\mathcal Q$ and $f$ does not depend on $\mathcal Q\setminus\mathcal Q_1$, then $P_{\mathcal Q}f=P_{\mathcal Q_1}f$ and $\norm f_{q,\mathcal Q}=\norm f_{q,\mathcal Q_1}$ for every real $q\ge1$.
\end{enumerate}
\end{lemma}
\begin{proof}
(F0) Sums, products, conjugates and powers of functions satisfying \eqref{eq:classA} satisfy \eqref{eq:classA}; $|h_{ij}|\le|x_{ij}|\le|x|_1$ (recall $h_{ii}=\re x_{ii}$). By Minkowski's inequality and Lemma~\ref{lem:moments},
$\norm{1+|x|_1}_{L^{s}(\PP)}\le1+\sum_{(i,j)}\big(\int|y|^{s}\nu_{ij}(\rd y)\big)^{1/s}<\infty$ for every $s\ge1$; this gives $\EE|f|^q<\infty$.
For $x\in\Omega$ and $y\in\kk^{\mathcal Q}$ we have $|(x_{\mathcal Q^c},y)|_1\le|x|_1+|y|_1$ and $(a+b)^m\le2^m(a^m+b^m)$ for $a,b\ge0$, so
\[
 \int|f(x_{\mathcal Q^c},y)|\,\PP_{\mathcal Q}(\rd y)\le c\,2^m\Big((1+|x|_1)^m+\int|y|_1^m\,\PP_{\mathcal Q}(\rd y)\Big)\le c'(1+|x|_1)^m,
\]
with $c':=c\,2^m(1+\int|y|_1^m\PP_{\mathcal Q}(\rd y))<\infty$ (same argument as before). Thus the integral converges absolutely for every $x$ and $P_{\mathcal Q}f$ satisfies \eqref{eq:classA}. Measurability of $P_{\mathcal Q}f$ is the
measurability of partial integrals of jointly measurable functions (Fubini--Tonelli theorem). By \eqref{eq:PQ}, $(P_{\mathcal Q}f)(x)$ is a function of $x_{\mathcal Q^c}$ alone. If $f$ does not depend on $\mathcal Q_1$ and $x_{\mathcal Q_1^c}=x'_{\mathcal Q_1^c}$,
then the points $(x_{\mathcal Q^c},y)$ and $(x'_{\mathcal Q^c},y)$ agree outside $\mathcal Q_1$ for every $y$, so the integrands in \eqref{eq:PQ} coincide. The last claim is the linearity of the integral.

(F1) The first identity is Fubini's theorem for $\PP=\PP_{\mathcal Q^c}\otimes\PP_{\mathcal Q}$ (Lemma~\ref{lem:split}) and the integrable function $f$. For $\mathcal Q_1\subset\mathcal Q$ apply Fubini's theorem to
$\PP_{\mathcal Q}=\PP_{\mathcal Q\setminus\mathcal Q_1}\otimes\PP_{\mathcal Q_1}$ and the function $y\mapsto f(x_{\mathcal Q^c},y)$, which is $\PP_{\mathcal Q}$-integrable by (F0):
$(P_{\mathcal Q}f)(x)=\int\!\big[\int f(x_{\mathcal Q^c},y',y'')\PP_{\mathcal Q_1}(\rd y'')\big]\PP_{\mathcal Q\setminus\mathcal Q_1}(\rd y')=(P_{\mathcal Q}P_{\mathcal Q_1}f)(x)$. With $\mathcal Q_1=\mathcal Q$: $P_{\mathcal Q}P_{\mathcal Q}f=P_{\mathcal Q}f$.

(F2) In \eqref{eq:PQ} for $gf$, the factor $g(x_{\mathcal Q^c},y)=g(x)$ is constant in $y$. For the last claim apply this to $g=\bone_\Xi$ and use (F1).

(F3) By (F0) all functions $y\mapsto f_i(x_{\mathcal Q^c},y)$ lie in every $L^q(\PP_{\mathcal Q})$. The first three statements are H\"older's inequality (for $n$ factors, applied to the functions $|f_i|^q$ with exponents $q_i/q$), Minkowski's inequality and the monotonicity of $L^q$-norms on the probability space $(\kk^{\mathcal Q},\PP_{\mathcal Q})$.
For the fourth: by Jensen's inequality on $(\kk^{\mathcal Q_1},\PP_{\mathcal Q_1})$, $|P_{\mathcal Q_1}f|^q\le P_{\mathcal Q_1}|f|^q$ pointwise; apply $P_{\mathcal Q}$ and (F1) (for the function $|f|^q\in\mathcal A$): $P_{\mathcal Q}|P_{\mathcal Q_1}f|^q\le P_{\mathcal Q}P_{\mathcal Q_1}|f|^q=P_{\mathcal Q}|f|^q$.

(F4) $\PP_{\mathcal Q}$ is a product measure, and the second statement is the definition \eqref{eq:PQ}.

(F5) Let $x,x'\in\Omega$ agree outside $\mathcal Q_1\cup\mathcal Q_2$, and let $x''\in\Omega$ coincide with $x'$ on $\mathcal Q_1$ and with $x$ outside $\mathcal Q_1$. Then $x$ and $x''$ agree outside $\mathcal Q_1$,
so $f(x)=f(x'')$; and $x''$ and $x'$ agree outside $\mathcal Q_2$ (on $\mathcal Q_1$ by construction, and outside $\mathcal Q_1\cup\mathcal Q_2$ because there $x''=x=x'$), so $f(x'')=f(x')$.
For the second statement write $y=(y',y'')$ with $y'\in\kk^{\mathcal Q\setminus\mathcal Q_1}$, $y''\in\kk^{\mathcal Q_1}$, so that $\PP_{\mathcal Q}=\PP_{\mathcal Q\setminus\mathcal Q_1}\otimes\PP_{\mathcal Q_1}$ (Lemma~\ref{lem:split}).
Since $f$ does not depend on $\mathcal Q\setminus\mathcal Q_1$, $f(x_{\mathcal Q^c},y',y'')=f(x_{\mathcal Q^c},x_{\mathcal Q\setminus\mathcal Q_1},y'')=f(x_{\mathcal Q_1^c},y'')$ for all $y',y''$. By Fubini's theorem (the integrand is $\PP_{\mathcal Q}$-integrable by (F0))
and because $\PP_{\mathcal Q\setminus\mathcal Q_1}$ is a probability measure,
$(P_{\mathcal Q}f)(x)=\int\!\big[\int f(x_{\mathcal Q_1^c},y'')\,\PP_{\mathcal Q_1}(\rd y'')\big]\PP_{\mathcal Q\setminus\mathcal Q_1}(\rd y')=(P_{\mathcal Q_1}f)(x)$.
Applied to $|f|^q$, which does not depend on $\mathcal Q\setminus\mathcal Q_1$ either, this gives $\norm f_{q,\mathcal Q}=\norm f_{q,\mathcal Q_1}$.
\end{proof}

\paragraph{The block matrix given a minor.} Let $L\subset\{1,\dots,N\}$ with $1\le r:=|L|<N$. We write $M:=M^L$ for the $L\times L$ matrix \eqref{eq:MA},
\begin{equation}\label{eq:ML}
 M_{ab}=h_{ab}-z\delta_{ab}-\sum_{k,l\notin L}h_{ak}G^{(L)}_{kl}h_{lb}\qquad(a,b\in L),
\end{equation}
and define
\begin{equation}\label{eq:PhiL}
\begin{aligned}
 m^{(L)}_a&:=\sum_{l\notin L}\sigma_{al}^2G^{(L)}_{ll}\ \ (a\in L),\qquad\qquad\Phi_L:=\sqrt{\frac{C_0}N}+C_0\sqrt{\frac{\im m^{(L)}}{N\eta}},\\
 Z^{(L)}_a&:=\sum_{k,l\notin L}h_{ak}G^{(L)}_{kl}h_{la}-\sum_{k\notin L}\sigma_{ak}^2G^{(L)}_{kk}\ \ (a\in L),
\end{aligned}
\end{equation}
where $m^{(L)}=\frac1N\sum_{k\notin L}G^{(L)}_{kk}$, so that $\im m^{(L)}>0$ by Lemma~\ref{lem:basic}(a) for the minor $H^{(L)}$. (For $L=\{k\}$: $Z^{(L)}_k=Z_k$.)

\begin{lemma}[Structure given the minor $H^{(L)}$]\label{lem:FAstruct}
Let $L$, $M$ be as above. Then, at every point of $\Omega$:
\begin{enumerate}[label=\textup{(\alph*)},leftmargin=2.2em,itemsep=1pt]
\item \emph{(Algebra)} $(G_{ab})_{a,b\in L}=M^{-1}$. For every $a\in L$: $M_{aa}=1/G^{(L\setminus\{a\})}_{aa}\ne0$ and $|1/M_{aa}|\le\eta^{-1}$. If $L=\{k\}$, then $1/G_{kk}=M_{kk}$. If $r\ge2$, $k\in L$ and $L':=L\setminus\{k\}$, then
the matrix $M_{L'L'}:=(M_{ab})_{a,b\in L'}$ is invertible, $(M_{L'L'})^{-1}=(G^{(k)}_{ab})_{a,b\in L'}$, $|G^{(k)}_{ab}|\le\eta^{-1}$, and
\begin{equation}\label{eq:1overGL}
 \frac1{G_{kk}}=M_{kk}-\sum_{a,b\in L'}M_{ka}\,G^{(k)}_{ab}\,M_{bk}.
\end{equation}
\item \emph{(Row dependence)} The functions $G^{(L)}_{kl}$ ($k,l\notin L$), $m^{(L)}$, $m^{(L)}_a$, $\Phi_L$ do not depend on $\mathcal Q(L)$. For $a,b\in L$ and every $c\in L\setminus\{a,b\}$, the function $M_{ab}$ does not depend on $\mathcal Q(\{c\})$.
\item \emph{(Regularity)} The functions $M_{ab}$, $1/M_{aa}$ ($a,b\in L$), $G^{(k)}_{ab}$ ($k\in L$, $a,b\in L\setminus\{k\}$), $1/G_{kk}$, $X_k$, $Z_a^{(L)}$ are continuous on $\Omega$ and belong to $\mathcal A$; $\Phi_L$ and $m_a^{(L)}$ are continuous and bounded.
\item \emph{(Expectation in one row)} For $a\in L$: $\ P_aM_{aa}=-z-m^{(L)}_a$ and $Q_aM_{aa}=h_{aa}-Z^{(L)}_a$. In particular, for every $k\in\{1,\dots,N\}$,
\[
 P_k\Big(\frac1{G_{kk}}\Big)=-z-\sum^{(k)}_l\sigma_{kl}^2G^{(k)}_{ll},\qquad X_k=Q_k\Big(\frac1{G_{kk}}\Big).
\]
\end{enumerate}
\end{lemma}
\begin{proof}
(a) The first statement is Lemma~\ref{lem:schur} with $A=L$. The identity $M_{aa}=1/G^{(L\setminus\{a\})}_{aa}$ is the last sentence of (I4) of Lemma~\ref{lem:schur} with $A=L$
(see also Remark~\ref{rem:rows}); explicitly, $1/G^{(L\setminus\{a\})}_{aa}=h_{aa}-z-\sum_{k,l\notin L}h_{ak}G^{(L)}_{kl}h_{la}=M_{aa}$. Further $|G^{(L\setminus\{a\})}_{aa}|\le\eta^{-1}$ by Lemma~\ref{lem:basic}(a) for this minor. The case $L=\{k\}$ is the same identity.
The statements for $r\ge2$ are (I4) of Lemma~\ref{lem:schur} with $A=L$, and Lemma~\ref{lem:basic}(a) for the minor $H^{(k)}$.

(b) $H^{(L)}$ is a function of the coordinates $\{i,j\}$ with $i,j\notin L$, i.e.\ of $x_{\mathcal Q(L)^c}$; so are $G^{(L)}$, $m^{(L)}$, $m^{(L)}_a$, $\Phi_L$. Let $a,b\in L$ and $c\in L\setminus\{a,b\}$. In \eqref{eq:ML}, $h_{ab}$ is a function of the coordinate $\{a,b\}$,
$h_{ak}$ of $\{a,k\}$, $h_{lb}$ of $\{l,b\}$ (with $k,l\notin L$), and $G^{(L)}_{kl}$ of coordinates $\{i,j\}$ with $i,j\notin L$. None of these coordinates contains $c$, i.e.\ none belongs to $\mathcal Q(\{c\})$.

(c) Entries of resolvents of $H$ and of all its minors $H^{(T)}$, $T\subsetneq\{1,\dots,N\}$, are continuous functions of $x\in\Omega$ (Lemma~\ref{lem:weyl}(e)) and bounded by $\eta^{-1}$ (Lemma~\ref{lem:basic}(a)). Hence $M_{ab}$ is continuous and
$|M_{ab}|\le|x|_1+|z|+\eta^{-1}|x|_1^2$; $1/M_{aa}=G^{(L\setminus\{a\})}_{aa}$ and $G^{(k)}_{ab}$ are continuous and bounded; $1/G_{kk}=M^{\{k\}}_{kk}$, $Z^{(L)}_a$ and $X_k=h_{kk}-Z^{(\{k\})}_k$ are continuous and satisfy \eqref{eq:classA} with $m=2$
(use $\sum_k\sigma^2_{ak}=1$). $\Phi_L$ and $m_a^{(L)}$ are continuous, $|m^{(L)}_a|\le\eta^{-1}$ and $\im m^{(L)}\le\eta^{-1}$.

(d) $\mathcal Q(\{a\})$ consists of the coordinates $\{a,j\}$, $1\le j\le N$. By (b) and (F2), $P_a(h_{ak}G^{(L)}_{kl}h_{la})=G^{(L)}_{kl}\,P_a(h_{ak}h_{la})$ for $k,l\notin L$. Now $h_{la}=\ol{h_{al}}$; for $k\ne l$ the variables $h_{ak}$, $h_{al}$ are functions of different coordinates of $\mathcal Q(\{a\})$, hence
independent under $\PP_{\mathcal Q(\{a\})}$ (F4), and centred by the centering assumption of Definition~\ref{def:ass}; so $P_a(h_{ak}h_{la})=0$. For $k=l$, $P_a(|h_{ak}|^2)=\sigma_{ak}^2$ by \eqref{eq:aver}. Finally $P_ah_{aa}=0$ and $P_az=z$. This gives $P_aM_{aa}=-z-m_a^{(L)}$, and
$Q_aM_{aa}=M_{aa}-P_aM_{aa}=h_{aa}-Z_a^{(L)}$ by \eqref{eq:ML}, \eqref{eq:PhiL}. For $L=\{k\}$, (a) gives $1/G_{kk}=M_{kk}$, hence $Q_k(1/G_{kk})=h_{kk}-Z_k=X_k$.
\end{proof}

\subsection{Conditional moment bounds and the a priori event}\label{sec:FAmom}
Let $C_\vth\ge1$ be the constant of Lemma~\ref{lem:moments} and $C_{\rm LD}\ge1$ the constant $C$ of Lemma~\ref{lem:LDmom} (both depend only on $\vth$), and put
\begin{equation}\label{eq:Cm}
 C_m:=\max\{C_\vth^{\fa},\,C_{\rm LD}^2\}\ \ (\ge1).
\end{equation}

\begin{lemma}[Conditional moments]\label{lem:FAmom}
Let $L\subset\{1,\dots,N\}$, $1\le|L|<N$, and $M=M^L$. For every real $q\ge2$, at every point of $\Omega$:
\begin{enumerate}[label=\textup{(m\arabic*)},leftmargin=2.8em,itemsep=1pt]
\item for $a\ne b\in L$:\quad $\norm{M_{ab}}_{q,L}\le C_m\,(q+\log N)^{2\fa+1}\,\Phi_L$;
\item for $a\in L$:\quad $\norm{Q_aM_{aa}}_{q,L}\le C_m\,(q+\log N)^{2\fa+\frac32}\,\Phi_L$.
\end{enumerate}
Moreover, let $\eta\ge N^{-1}$, $q_*\ge2$ real and $a\in L$. Then:
\begin{enumerate}[label=\textup{(m\arabic*)},leftmargin=2.8em,itemsep=1pt,start=3]
\item at every point of $\Omega$ at which
\begin{equation}\label{eq:m3hyp}
 |z+m_a^{(L)}|\ge\tfrac25\qquad\text{and}\qquad5e\,C_m\big((q_*+1)\lceil\log N\rceil\big)^{2\fa+\frac32}\,\Phi_L\le1,
\end{equation}
we have $\norm{1/M_{aa}}_{q,L}\le6$ for all real $q\in[1,q_*]$.
\end{enumerate}
\end{lemma}
\noindent Note that (m1), (m2) need no assumption on $G$, and that the hypothesis \eqref{eq:m3hyp} concerns functions that do not depend on $\mathcal Q(L)$ (Lemma~\ref{lem:FAstruct}(b)): it is a condition on the minor $H^{(L)}$ only.
\begin{proof}
Fix $L$ and put $\mathcal R_a:=\{\{a,k\}:k\notin L\}\subset\mathcal Q(L)$ for $a\in L$; these sets are pairwise disjoint. Let $n:=N-|L|\in[1,N]$. By (F4), for fixed $x_{\mathcal R_a^c}$ the variables $(h_{ak})_{k\notin L}$ are independent on $(\kk^{\mathcal R_a},\PP_{\mathcal R_a})$; $h_{ak}$ is the coordinate $x_{ak}$ or its complex conjugate, so by
Definition~\ref{def:ass} it is centred, has variance $\sigma_{ak}^2$ and satisfies the tail bound in \eqref{eq:LDhyp} with $s_k=\sigma_{ak}$. The same holds for the family $(h_{ka})_{k\notin L}=(\ol{h_{ak}})_{k\notin L}$. So Lemma~\ref{lem:LDmom} applies to these families, with coefficients that are constant on $\kk^{\mathcal R_a}$,
and $\log n\le\log N$. We shall also use that, by (C) and the Ward identity (Lemma~\ref{lem:basic}(b)) for the minor $H^{(L)}$, for all $a,b\in L$,
\begin{equation}\label{eq:WardL}
 \sum_{k,l\notin L}\sigma_{ak}^2\,|G^{(L)}_{kl}|^2\,\sigma_{lb}^2\le\frac{C_0}N\sum_{l\notin L}\sigma_{lb}^2\sum_{k\notin L}|G^{(L)}_{kl}|^2=\frac{C_0}N\sum_{l\notin L}\sigma_{lb}^2\frac{\im G^{(L)}_{ll}}{\eta}\le C_0^2\,\frac{\im m^{(L)}}{N\eta}.
\end{equation}

\emph{(m1).} Let $a\ne b\in L$ and $B_{ab}:=\sum_{k,l\notin L}h_{ak}G^{(L)}_{kl}h_{lb}=\sum_{k\notin L}h_{ak}Y_k$ with $Y_k:=\sum_{l\notin L}G^{(L)}_{kl}h_{lb}$. The functions $Y_k$ do not depend on $\mathcal R_a$ (they depend on $x_{\mathcal Q(L)^c}$ and on $\mathcal R_b$).
By \eqref{eq:LDm1} for the family $(h_{ak})_{k\notin L}$, with $x_{\mathcal R_a^c}$ fixed,
\[
 \norm{B_{ab}}_{q,\mathcal R_a}\le C_{\rm LD}(q+\log N)^{\fa+\frac12}\Big(\sum_{k\notin L}\sigma_{ak}^2|Y_k|^2\Big)^{1/2}.
\]
Raise this to the power $q$, apply $P_L$, and use $P_LP_{\mathcal R_a}=P_L$ (F1), then Minkowski's inequality in $L^{q/2}(\PP_{\mathcal Q(L)})$ ($q/2\ge1$; (F3)):
\[
 \norm{B_{ab}}_{q,L}\le C_{\rm LD}(q+\log N)^{\fa+\frac12}\Big\|\sum_{k\notin L}\sigma_{ak}^2|Y_k|^2\Big\|_{q/2,L}^{1/2}\le C_{\rm LD}(q+\log N)^{\fa+\frac12}\Big(\sum_{k\notin L}\sigma_{ak}^2\norm{Y_k}_{q,L}^2\Big)^{1/2}.
\]
Since $Y_k$ depends, among the coordinates of $\mathcal Q(L)$, only on those of $\mathcal R_b$, $\norm{Y_k}_{q,L}=\norm{Y_k}_{q,\mathcal R_b}$ (by (F5), since $Y_k$ does not depend on $\mathcal Q(L)\setminus\mathcal R_b$), and by \eqref{eq:LDm1} for the family $(h_{lb})_{l\notin L}$ (with $x_{\mathcal R_b^c}$ fixed),
$\norm{Y_k}_{q,\mathcal R_b}\le C_{\rm LD}(q+\log N)^{\fa+\frac12}(\sum_{l\notin L}|G^{(L)}_{kl}|^2\sigma_{lb}^2)^{1/2}$. With \eqref{eq:WardL},
\[
 \norm{B_{ab}}_{q,L}\le C_{\rm LD}^2(q+\log N)^{2\fa+1}\,C_0\sqrt{\frac{\im m^{(L)}}{N\eta}}.
\]
The coordinate $\{a,b\}$ belongs to $\mathcal Q(L)$ and $h_{ab}$ depends on no other coordinate, so by (F5) and (F4) $\norm{h_{ab}}_{q,L}=\norm{h_{ab}}_{q,\{\{a,b\}\}}=(\EE|h_{ab}|^q)^{1/q}$, and by Lemma~\ref{lem:moments} and (C), $(\EE|h_{ab}|^q)^{1/q}\le(C_\vth q)^{\fa}\sigma_{ab}\le C_\vth^{\fa}(q+\log N)^{2\fa+1}\sqrt{C_0/N}$. Since $M_{ab}=h_{ab}-B_{ab}$, (m1) follows from Minkowski's inequality and \eqref{eq:Cm}, \eqref{eq:PhiL}.

\emph{(m2).} By Lemma~\ref{lem:FAstruct}(d), $Q_aM_{aa}=h_{aa}-Z_a^{(L)}$. With $\xi_k:=h_{ka}=\ol{h_{ak}}$ and $\mathsf B_{kl}:=G^{(L)}_{kl}$ ($k,l\notin L$) we have $Z_a^{(L)}=\sum_{k,l}\ol{\xi_k}\mathsf B_{kl}\xi_l-\sum_k\sigma_{ak}^2\mathsf B_{kk}$, and $Z_a^{(L)}$ depends, among the coordinates of $\mathcal Q(L)$, only on those of $\mathcal R_a$.
By (F5), \eqref{eq:LDm2} (with $x_{\mathcal R_a^c}$ fixed) and \eqref{eq:WardL} with $b=a$,
\[
 \norm{Z_a^{(L)}}_{q,L}=\norm{Z_a^{(L)}}_{q,\mathcal R_a}\le C_{\rm LD}(q+\log N)^{2\fa+\frac32}\,C_0\sqrt{\frac{\im m^{(L)}}{N\eta}},
\]
and $\norm{h_{aa}}_{q,L}\le(C_\vth q)^{\fa}\sqrt{C_0/N}$ as before. Minkowski's inequality gives (m2).

\emph{(m3).} Fix a point $x$ at which \eqref{eq:m3hyp} holds; the functions in \eqref{eq:m3hyp} are constant on the fibre $\{x_{\mathcal Q(L)^c}\}\times\kk^{\mathcal Q(L)}$, on which we now work. By Lemma~\ref{lem:FAstruct}(d), $M_{aa}=-(z+m_a^{(L)})+Q_aM_{aa}$.
Hence on the set $\{|Q_aM_{aa}|\le\frac15\}$ we have $|M_{aa}|\ge\frac25-\frac15$, i.e.\ $|1/M_{aa}|\le5$; and always $|1/M_{aa}|\le\eta^{-1}\le N$ (Lemma~\ref{lem:FAstruct}(a)).
Let $q':=q_*\lceil\log N\rceil\ge2$; then $q'+\log N\le(q_*+1)\lceil\log N\rceil$. By Markov's inequality on $(\kk^{\mathcal Q(L)},\PP_{\mathcal Q(L)})$, (m2) with exponent $q'$, and \eqref{eq:m3hyp},
\[
 \PP_{\mathcal Q(L)}\big(|Q_aM_{aa}|>\tfrac15\big)\le\big(5\norm{Q_aM_{aa}}_{q',L}\big)^{q'}\le\Big(5C_m\big((q_*+1)\lceil\log N\rceil\big)^{2\fa+\frac32}\Phi_L\Big)^{q'}\le e^{-q'}\le N^{-q_*}.
\]
Therefore, for $1\le q\le q_*$: $\ P_L|1/M_{aa}|^q\le5^q+N^q\,N^{-q_*}\le5^q+1\le6^q$.
\end{proof}

\begin{lemma}[The a priori event controls all small minors]\label{lem:FAevent}
Let $p\ge1$ be an integer with $12\,C_0\,p\le N$, let $L\subset\{1,\dots,N\}$ with $1\le|L|\le p$ (so that $|L|\le p<N$, since $C_0\ge1$), and let $\Phi_*$ be defined by \eqref{eq:Phistar} (with this $p$). Then:
\begin{enumerate}[label=\textup{(\alph*)},leftmargin=2.2em,itemsep=1pt]
\item at every point of $\Omega$:\quad $N^{-1}\le\Phi_*^2$\quad and\quad $\Phi_L^2\le2C_0^2\,\Phi_*^2$;
\item on the event $\Xi_p(z)=\{\Ld\le\frac12|\msc|,\ \Lo\le\frac{|\msc|}{16p}\}$:\quad $\im m\le\frac32$\quad and\quad $|z+m_a^{(L)}|\ge\frac25$ for all $a\in L$;
\item if moreover $\eta\le10$ and $N\eta\ge p$, then on $\Xi_p(z)$:\quad $\Phi_L^2\le25\,C_0^2/(N\eta)$.
\end{enumerate}
\end{lemma}
\begin{proof}
(a) The first inequality is clear. By (I5) of Lemma~\ref{lem:schur}, $\im m^{(L)}\le\im m+|L|/(N\eta)\le\im m+p/(N\eta)$. Using $(u+v)^2\le2u^2+2v^2$ and $C_0\ge1$,
$\Phi_L^2\le2C_0/N+2C_0^2\im m^{(L)}/(N\eta)\le2C_0^2\Phi_*^2$.

(b) Recall $|\msc|\le1$ and $|z+\msc|=1/|\msc|\ge1$ (Lemma~\ref{lem:msc}(a),(c)). On $\Xi_p(z)$: $|G_{kk}|\le|\msc|+\Ld\le\frac32$ for all $k$, so $\im m\le|m|\le\frac32$. Next, let $g:=\frac12|\msc|$ and $r:=|L|\le p$. On $\Xi_p(z)$, $|G_{aa}|\ge|\msc|-\Ld\ge g$ for all $a$,
and $\max_{a\ne b}|G_{ab}|=\Lo\le|\msc|/(16p)\le g/(2r)$. By Lemma~\ref{lem:minor} (with $A=L$), for all $l\notin L$,
\[
 |G^{(L)}_{ll}-G_{ll}|\le\frac{2r^2}{g}\Lo^2\le\frac{4p^2}{|\msc|}\cdot\frac{|\msc|^2}{256\,p^2}=\frac{|\msc|}{64},\qquad\text{hence}\quad|G^{(L)}_{ll}-\msc|\le\Big(\frac12+\frac1{64}\Big)|\msc|.
\]
By (A) and (C), $\sum_{l\notin L}\sigma_{al}^2=1-\sum_{l\in L}\sigma_{al}^2\in[1-pC_0/N,1]$, so for $a\in L$
\[
 |m_a^{(L)}-\msc|\le\sum_{l\notin L}\sigma_{al}^2|G^{(L)}_{ll}-\msc|+|\msc|\sum_{l\in L}\sigma_{al}^2\le\Big(\frac{33}{64}+\frac{pC_0}N\Big)|\msc|\le\Big(\frac{33}{64}+\frac1{12}\Big)|\msc|\le\frac35,
\]
and $|z+m_a^{(L)}|\ge|z+\msc|-|m_a^{(L)}-\msc|\ge1-\frac35=\frac25$.

(c) By (a), (b), $N^{-1}\le10/(N\eta)$ and $p/(N\eta)\le1$: $\ \Phi_L^2\le2C_0^2\big(\frac{10}{N\eta}+\frac{3/2+1}{N\eta}\big)=25C_0^2/(N\eta)$.
\end{proof}

\begin{lemma}[Crude moment bound]\label{lem:FAcrude}
Let $\eta\ge N^{-1}$. For every $k$ and every real $q\ge2$: $\ \norm{X_k}_{L^q(\PP)}\le3N(2C_\vth q)^{2\fa}$. In particular $\norm{X_k}_{L^{2p}(\PP)}\le N^2$ whenever $p\ge1$ and $3(4C_\vth p)^{2\fa}\le N$.
\end{lemma}
\begin{proof}
Since $\norm{G^{(k)}}\le\eta^{-1}\le N$ and $|G^{(k)}_{ii}|\le N$ (Lemma~\ref{lem:basic}(a)) and $\sum_i\sigma_{ki}^2=1$: $|Z_k|\le N\sum_{i\ne k}|h_{ki}|^2+N$. By Minkowski's inequality and Lemma~\ref{lem:moments} (with exponents $q$ and $2q$),
\[
 \norm{X_k}_{L^q}\le\norm{h_{kk}}_{L^q}+N\sum_{i\ne k}\norm{h_{ki}}_{L^{2q}}^2+N\le(C_\vth q)^{\fa}\sigma_{kk}+N(2C_\vth q)^{2\fa}\sum_i\sigma_{ki}^2+N\le3N(2C_\vth q)^{2\fa}.\qedhere
\]
\end{proof}

\subsection{The case $p=2$}\label{sec:FAtwo}
The second moment is treated separately. It is not literally a case of Theorem~\ref{thm:FA} (which assumes $p\ge\log N$), and the constant carries a power of $\log N$ which comes from the term $\log n$ in Lemma~\ref{lem:LDmom}.
Its proof is the proof of Theorem~\ref{thm:FA} in the situation where the expansion of Section~\ref{sec:FAexp} terminates after one step and has no remainder.

\begin{proposition}[Variance bound]\label{prop:FAtwo}
There are constants $C_V$ and $N_0$, depending only on the model parameters, such that for $N\ge N_0$, $\nu\in\mathcal E_N(\vth,\dpm,C_0)$, every $z=E+\ii\eta$ with
\[
 N^{-1}\le\eta\le10,\qquad N\eta\ge(\log N)^{4\fa+4},
\]
and all deterministic $t_1,\dots,t_N\in\C$ with $|t_k|\le1/N$:
\begin{equation}\label{eq:FAtwo}
 \EE\Big|\sum_{k=1}^Nt_kX_k(z)\Big|^2\le C_V(\log N)^{8\fa+4}\ \EE\big[\Phi_*(z)^4\big]+N^4\,\PP\big(\Xi_2(z)^c\big)^{1/2},
\end{equation}
where $\Xi_2(z)=\{\Ld\le\frac12|\msc|,\ \Lo\le\frac1{32}|\msc|\}$ and $\Phi_*^2=\frac1N+\frac1{N\eta}\big(\im m+\frac2{N\eta}\big)$ (i.e.\ \eqref{eq:Phistar} with $p=2$). One can take
\begin{equation}\label{eq:CV}
 C_V:=578\,C_0^4\,C_m^4\,7^{8\fa+4},\qquad N_0:=\max\Big\{24C_0,\ 3(8C_\vth)^{2\fa},\ \exp\big(625\,e^2C_0^2C_m^2\,14^{4\fa+3}\big)\Big\}.
\end{equation}
\end{proposition}
\begin{proof}
All functions below belong to $\mathcal A$ (Lemma~\ref{lem:FAstruct}(c), (F0)), so every expectation and every partial integral is well defined. Expanding the square,
\begin{equation}\label{eq:two0}
 \EE\Big|\sum_kt_kX_k\Big|^2=\sum_k|t_k|^2\,\EE|X_k|^2+\sum_{k\ne l}t_k\ol{t_l}\ \EE\big[X_k\ol{X_l}\big].
\end{equation}

\emph{Step 1: the diagonal terms.} Fix $k$ and let $L=\{k\}$, $M=M^{\{k\}}$. By Lemma~\ref{lem:FAstruct}(a),(d), $X_k=Q_kM_{kk}$, so by (m2) of Lemma~\ref{lem:FAmom} with $q=2$ and Lemma~\ref{lem:FAevent}(a) with $p=2$ (note $24C_0\le N$), at every point,
$P_k|X_k|^2\le C_m^2(2+\log N)^{4\fa+3}\Phi_{\{k\}}^2\le2C_0^2C_m^2(2+\log N)^{4\fa+3}\Phi_*^2$. By (F1), $\EE|X_k|^2=\EE[P_k|X_k|^2]$. Since $|t_k|^2\le N^{-2}$ and $N^{-1}\le\Phi_*^2$,
\begin{equation}\label{eq:two1}
 \sum_k|t_k|^2\,\EE|X_k|^2\le2C_0^2C_m^2(2+\log N)^{4\fa+3}\,\EE\big[N^{-1}\Phi_*^2\big]\le2C_0^2C_m^2(2+\log N)^{4\fa+3}\,\EE\big[\Phi_*^4\big].
\end{equation}

\emph{Step 2: off-diagonal terms, algebra.} Fix $k\ne l$ and let $L:=\{k,l\}$, $M:=M^L$ (a $2\times2$ matrix). By Lemma~\ref{lem:FAstruct}(a) with $L'=\{l\}$, $G^{(k)}_{ll}=(M_{ll})^{-1}$ and \eqref{eq:1overGL} reads
$1/G_{kk}=M_{kk}-M_{kl}M_{lk}/M_{ll}$; the same holds with $k$ and $l$ exchanged. Applying $Q_k$, resp.\ $Q_l$ (Lemma~\ref{lem:FAstruct}(d)),
\begin{equation}\label{eq:two2}
 X_k=Q_kM_{kk}-Q_kW_k,\qquad X_l=Q_lM_{ll}-Q_lW_l,\qquad W_k:=\frac{M_{kl}M_{lk}}{M_{ll}},\quad W_l:=\frac{M_{lk}M_{kl}}{M_{kk}}.
\end{equation}

\emph{Step 3: off-diagonal terms, vanishing.} We claim that
\begin{equation}\label{eq:two3}
 P_L\big(X_k\ol{X_l}\big)=P_L\big(Q_kW_k\cdot\ol{Q_lW_l}\big).
\end{equation}
By Lemma~\ref{lem:FAstruct}(b) (with $a=b=k$, $c=l$), $M_{kk}$ does not depend on $\mathcal Q(\{l\})$; by (F0) neither does $Q_kM_{kk}=M_{kk}-P_kM_{kk}$. Hence, by (F2) and (F1) ($P_lQ_l=0$), for $f\in\{M_{ll},W_l\}$:
$P_l\big(Q_kM_{kk}\cdot\ol{Q_lf}\big)=Q_kM_{kk}\cdot\ol{P_lQ_lf}=0$. Symmetrically, $Q_lM_{ll}$ does not depend on $\mathcal Q(\{k\})$ and $P_k\big(Q_kW_k\cdot\ol{Q_lM_{ll}}\big)=\ol{Q_lM_{ll}}\cdot P_kQ_kW_k=0$.
Since $P_L=P_LP_l=P_LP_k$ (F1), three of the four terms of $P_L(X_k\ol{X_l})$ obtained from \eqref{eq:two2} vanish, which is \eqref{eq:two3}.
(In the language of Section~\ref{sec:FAexp}: both indices are lone; the term $Q_kM_{kk}$ visits no index and the term $Q_kW_k$ visits $l$; a product survives only if $k$ and $l$ are both visited.)

\emph{Step 4: off-diagonal terms, size.} Let
\begin{align*}
 \tau_2&:=C_m(6+\log N)^{2\fa+1},\qquad\varphi_2:=\big(5e\,C_m(7\lceil\log N\rceil)^{2\fa+\frac32}\big)^{-1},\\
 \Xi_L&:=\{\Phi_L\le\varphi_2\}\cap\bigcap_{a\in L}\big\{|z+m_a^{(L)}|\ge\tfrac25\big\}.
\end{align*}
By Lemma~\ref{lem:FAstruct}(b),(c), $\Xi_L$ is a closed subset of $\Omega$ and $\bone_{\Xi_L}$ does not depend on $\mathcal Q(L)$. At every point of $\Xi_L$, \eqref{eq:m3hyp} holds with $q_*=6$, so by (F3) (H\"older with three factors in $L^6$), (m1) and (m3):
\[
 \norm{W_k}_{2,L}\le\norm{M_{kl}}_{6,L}\,\norm{M_{lk}}_{6,L}\,\norm{1/M_{ll}}_{6,L}\le6\,\tau_2^2\,\Phi_L^2,
\]
and the same for $W_l$. By \eqref{eq:two3}, (F3) (Cauchy--Schwarz, and $\norm{Q_kf}_{2,L}\le2\norm f_{2,L}$) and Lemma~\ref{lem:FAevent}(a), on $\Xi_L$:
\begin{equation}\label{eq:two4}
 \big|P_L\big(X_k\ol{X_l}\big)\big|\le4\,\norm{W_k}_{2,L}\norm{W_l}_{2,L}\le144\,\tau_2^4\,\Phi_L^4\le576\,C_0^4\,\tau_2^4\,\Phi_*^4 .
\end{equation}

\emph{Step 5: the complement of $\Xi_L$.} By (F2), the Cauchy--Schwarz and H\"older inequalities, and Lemma~\ref{lem:FAcrude} (with $p=2$; note $3(8C_\vth)^{2\fa}\le N$),
\begin{align*}
 \EE\big[X_k\ol{X_l}\big]&=\EE\big[\bone_{\Xi_L}P_L(X_k\ol{X_l})\big]+\EE\big[\bone_{\Xi_L^c}X_k\ol{X_l}\big],\\
 \big|\EE\big[\bone_{\Xi_L^c}X_k\ol{X_l}\big]\big|&\le\PP(\Xi_L^c)^{1/2}\norm{X_k}_{L^4(\PP)}\norm{X_l}_{L^4(\PP)}\le N^4\,\PP(\Xi_L^c)^{1/2}.
\end{align*}
We claim $\Xi_2(z)\subset\Xi_L$. By Lemma~\ref{lem:FAevent}(b),(c) with $p=2$ (note $N\eta\ge2$), on $\Xi_2(z)$ we have $|z+m_a^{(L)}|\ge\frac25$ for $a\in L$ and $\Phi_L^2\le25C_0^2/(N\eta)$; and
$25C_0^2/(N\eta)\le\varphi_2^2$ because $N\eta\ge(\log N)^{4\fa+4}$, $\lceil\log N\rceil\le2\log N$ and $\log N\ge625e^2C_0^2C_m^214^{4\fa+3}$ for $N\ge N_0$. Hence $\PP(\Xi_L^c)\le\PP(\Xi_2(z)^c)$.

\emph{Step 6: conclusion.} By Steps 4 and 5, $|\EE[X_k\ol{X_l}]|\le576C_0^4\tau_2^4\,\EE[\Phi_*^4]+N^4\PP(\Xi_2(z)^c)^{1/2}$ for all $k\ne l$, and $\sum_{k\ne l}|t_k||t_l|\le1$. With \eqref{eq:two0}, \eqref{eq:two1} and $2+\log N\le6+\log N\le7\log N$:
\[
 \EE\Big|\sum_kt_kX_k\Big|^2\le\big(2C_0^2C_m^2+576C_0^4C_m^4\big)(7\log N)^{8\fa+4}\,\EE[\Phi_*^4]+N^4\,\PP(\Xi_2(z)^c)^{1/2},
\]
which is \eqref{eq:FAtwo} with the constant \eqref{eq:CV}.
\end{proof}

\begin{remark}[Reading of the case $p=2$]\label{rem:FAtwo}
(1) The gain is in \eqref{eq:two3}--\eqref{eq:two4}: each $X_k$ is of size $\Phi$ (Step~1), but the covariance of $X_k$ and $X_l$, $k\ne l$, is of size $\Phi^4$, not $\Phi^2$; there are $N^2$ covariances with weight $N^{-2}$ and $N$ variances with weight $N^{-2}$, and $N^{-1}\Phi^2\le\Phi^4$. Thus $\sum_kt_kX_k$ is of size $\Phi^2$.
(2) The three ingredients are: the exact identity \eqref{eq:two2}, in which every term is a function of the minor $H^{(L)}$ and of the rows it ``visits''; the vanishing \eqref{eq:two3}, which is Fubini's theorem; and conditional moment bounds given $H^{(L)}$, valid on an event $\Xi_L$ that is measurable with respect to $H^{(L)}$ and contains the a priori event $\Xi_2(z)$ deterministically.
No bound on the probability of $\Xi_2(z)$ is assumed: it enters \eqref{eq:FAtwo} additively.
(3) For general $p$ the identity \eqref{eq:two2} is replaced by a finite expansion with a remainder (Lemma~\ref{lem:FAexp}), and the powers of $\log N$ by powers of $p$, because moments of order up to $8p^2$ are needed.
\end{remark}

\subsection{The finite expansion and the vanishing of terms}\label{sec:FAexp}
Let $L\subset\{1,\dots,N\}$, $1\le r=|L|<N$, $M=M^L$, $k\in L$, $L':=L\setminus\{k\}$, and let $n_0\ge1$ be an integer.
\begin{definition}[Paths]\label{def:FApaths}
For $n\ge0$, a \emph{path of length $n$ in $L'$} is a tuple $\pi=(a_0,\dots,a_n)\in(L')^{n+1}$ with $a_t\ne a_{t+1}$ for $0\le t<n$; $\Pi_n(k)$ is the set of these paths (it is empty if $L'=\emptyset$).
For $\pi\in\Pi_n(k)$ put $V(\pi):=\{a_0,\dots,a_n\}\subset L'$ (the indices \emph{visited} by $\pi$), $o(\pi):=n+2$ (the \emph{order}), and
\begin{equation}\label{eq:Wk}
 W_k(\pi):=\frac{M_{ka_0}\,M_{a_0a_1}\cdots M_{a_{n-1}a_n}\,M_{a_nk}}{M_{a_0a_0}\,M_{a_1a_1}\cdots M_{a_na_n}}
\end{equation}
($n+2$ off-diagonal entries of $M$ in the numerator, $n+1$ diagonal entries in the denominator). The \emph{empty path} $\emptyset$ has $V(\emptyset):=\emptyset$, $o(\emptyset):=1$, $W_k(\emptyset):=M_{kk}$. Let
\[
 \mathcal T(k):=\{\emptyset\}\cup\bigcup_{n=0}^{n_0-1}\Pi_n(k),\qquad\epsilon(\emptyset):=1,\qquad\epsilon(\pi):=(-1)^{n+1}\ \ (\pi\in\Pi_n(k)),
\]
and define the \emph{remainder}
\begin{equation}\label{eq:Rk}
 R_k:=(-1)^{n_0}\sum_{(a_0,\dots,a_{n_0})\in\Pi_{n_0}(k)}\ \sum_{b\in L'}M_{ka_0}\Big(\prod_{t=0}^{n_0-1}\frac{M_{a_ta_{t+1}}}{M_{a_ta_t}}\Big)G^{(k)}_{a_{n_0}b}\,M_{bk}\qquad(R_k:=0\text{ if }L'=\emptyset).
\end{equation}
\end{definition}

\begin{lemma}[Finite expansion]\label{lem:FAexp}
With the notation above, at every point of $\Omega$,
\begin{equation}\label{eq:expansion}
 \frac1{G_{kk}}=\sum_{\pi\in\mathcal T(k)}\epsilon(\pi)\,W_k(\pi)-R_k=M_{kk}-\sum_{n=0}^{n_0-1}(-1)^n\!\!\sum_{\pi\in\Pi_n(k)}W_k(\pi)-R_k .
\end{equation}
All $W_k(\pi)$ and $R_k$ belong to $\mathcal A$, and
\begin{equation}\label{eq:FAX}
 X_k=\sum_{\pi\in\mathcal T(k)}\epsilon(\pi)\,Q_kW_k(\pi)-Q_kR_k .
\end{equation}
Moreover, for every $\pi\in\mathcal T(k)$ and every integer $o\ge1$:
\begin{equation}\label{eq:FAcountpaths}
 o(\pi)\ge|V(\pi)|+1,\qquad\#\{\pi\in\mathcal T(k):\ o(\pi)=o\}\le r^{\,o-1},\qquad|\Pi_{n_0}(k)|\cdot|L'|\le r^{\,n_0+2}.
\end{equation}
\end{lemma}
\begin{proof}
If $L'=\emptyset$ then $\mathcal T(k)=\{\emptyset\}$, $R_k=0$, and \eqref{eq:expansion} is $1/G_{kk}=M_{kk}$ (Lemma~\ref{lem:FAstruct}(a)). Let $r\ge2$. Write $A:=M_{L'L'}=D+O$ with $D$ the diagonal part and $O$ the off-diagonal part of $A$. By Lemma~\ref{lem:FAstruct}(a), $A$ is invertible with
$A^{-1}=(G^{(k)}_{ab})_{a,b\in L'}$, and $D$ is invertible because $D_{aa}=M_{aa}\ne0$. Let $T:=-D^{-1}O$, so $A=D(1-T)$; thus $1-T=D^{-1}A$ is invertible and $A^{-1}=(1-T)^{-1}D^{-1}$. From the telescoping identity $(1-T)\sum_{n=0}^{n_0-1}T^n=1-T^{n_0}$ we get
$(1-T)^{-1}=\sum_{n=0}^{n_0-1}T^n+T^{n_0}(1-T)^{-1}$, hence the purely algebraic identity
\begin{equation}\label{eq:Neumann}
 A^{-1}=\sum_{n=0}^{n_0-1}T^nD^{-1}+T^{n_0}A^{-1}.
\end{equation}
(No convergence is involved.) The entries of $T$ are $T_{aa}=0$ and $T_{ab}=-M_{ab}/M_{aa}$ for $a\ne b\in L'$. Therefore, for $a,b\in L'$ and $n\ge0$,
\[
 (T^n)_{ab}=(-1)^n\sum_{\substack{(a_0,\dots,a_n)\in\Pi_n(k)\\ a_0=a,\ a_n=b}}\ \prod_{t=0}^{n-1}\frac{M_{a_ta_{t+1}}}{M_{a_ta_t}},\qquad(T^nD^{-1})_{ab}=(-1)^n\sum_{\substack{(a_0,\dots,a_n)\in\Pi_n(k)\\ a_0=a,\ a_n=b}}\ \frac{\prod_{t=0}^{n-1}M_{a_ta_{t+1}}}{\prod_{t=0}^{n}M_{a_ta_t}} .
\]
Insert \eqref{eq:Neumann} into \eqref{eq:1overGL}: $\sum_{a,b\in L'}M_{ka}(T^nD^{-1})_{ab}M_{bk}=(-1)^n\sum_{\pi\in\Pi_n(k)}W_k(\pi)$ by \eqref{eq:Wk}, and
$\sum_{a,b\in L'}M_{ka}(T^{n_0}A^{-1})_{ab}M_{bk}=\sum_{a,c,b\in L'}M_{ka}(T^{n_0})_{ac}G^{(k)}_{cb}M_{bk}=R_k$ by \eqref{eq:Rk}. This proves \eqref{eq:expansion}.

By Lemma~\ref{lem:FAstruct}(c) and (F0), all $W_k(\pi)$ and $R_k$ are in $\mathcal A$ (finite sums of products of functions of $\mathcal A$); so $Q_k$ may be applied term by term to \eqref{eq:expansion}, and $Q_k(1/G_{kk})=X_k$ by Lemma~\ref{lem:FAstruct}(d). This is \eqref{eq:FAX}.

Finally, $o(\emptyset)=1=|V(\emptyset)|+1$, and for $\pi\in\Pi_n(k)$: $|V(\pi)|\le n+1=o(\pi)-1$. The only $\pi\in\mathcal T(k)$ of order $1$ is $\emptyset$; those of order $o\ge2$ are the elements of $\Pi_{o-2}(k)$ (if $o-2<n_0$), and $|\Pi_{n}(k)|\le|L'|^{n+1}\le r^{n+1}$.
The last bound is $|L'|^{n_0+1}\cdot|L'|\le r^{n_0+2}$.
\end{proof}

\begin{lemma}[Row dependence and vanishing]\label{lem:FAvanish}
Let $L$, $M$, $n_0$ be as above.
\begin{enumerate}[label=\textup{(\alph*)},leftmargin=2.2em,itemsep=1pt]
\item Let $k\in L$, $\pi\in\mathcal T(k)$ and $c\in L\setminus(\{k\}\cup V(\pi))$. Then $W_k(\pi)$ and $Q_kW_k(\pi)$ do not depend on $\mathcal Q(\{c\})$.
\item Let $j\ge1$, $k_1,\dots,k_j\in L$, $\pi_s\in\mathcal T(k_s)$ for $1\le s\le j$, and let $f^{\#_s}$ denote either $f$ or $\ol f$, for each $s$. Suppose that there is an $s$ such that
\begin{equation}\label{eq:lonecond}
 k_s\notin\{k_u\}\cup V(\pi_u)\qquad\text{for all }u\ne s .
\end{equation}
Then
\[
 P_L\Big(\prod_{u=1}^j\big(Q_{k_u}W_{k_u}(\pi_u)\big)^{\#_u}\Big)=0\qquad\text{at every point of }\Omega .
\]
\end{enumerate}
\end{lemma}
\begin{proof}
(a) Every factor of $W_k(\pi)$ in \eqref{eq:Wk} (resp.\ $W_k(\emptyset)=M_{kk}$) is of the form $M_{ab}$ or $1/M_{aa}$ with $a,b\in\{k\}\cup V(\pi)$. By Lemma~\ref{lem:FAstruct}(b), none of them depends on $\mathcal Q(\{c\})$, since $c\in L\setminus\{a,b\}$.
So $W_k(\pi)$ does not depend on $\mathcal Q(\{c\})$, and by (F0) neither does $P_kW_k(\pi)$, nor $Q_kW_k(\pi)=W_k(\pi)-P_kW_k(\pi)$.

(b) Put $Y_u:=(Q_{k_u}W_{k_u}(\pi_u))^{\#_u}\in\mathcal A$ and $c:=k_s$. For $u\ne s$, \eqref{eq:lonecond} says $c\in L\setminus(\{k_u\}\cup V(\pi_u))$, so by (a) (and since conjugation does not change the dependence) $Y_u$ does not depend on $\mathcal Q(\{c\})$; hence neither does $g:=\prod_{u\ne s}Y_u\in\mathcal A$.
By (F2), $P_c(g\,Y_s)=g\,P_cY_s$, and by (F0), (F1), $P_cY_s=(P_cQ_cW_c(\pi_s))^{\#_s}=0$. Since $\mathcal Q(\{c\})\subset\mathcal Q(L)$, (F1) gives $P_L(gY_s)=P_LP_c(gY_s)=0$.
\end{proof}
\noindent In words: given the minor $H^{(L)}$, a term $Q_kW_k(\pi)$ is a function of the rows $k$ and $V(\pi)$ only; if some index $k_s$ occurs in exactly one factor, as the index on which $Q_{k_s}$ acts, then integrating the row $k_s$ first gives zero. Remainder terms $Q_kR_k$ are not claimed to vanish.

\subsection{Size of the terms; one index tuple}\label{sec:FAsize}
From now on $p\ge2$ is an even integer, $N\ge8$ and $\log N\le p$, $\eta\ge N^{-1}$. We fix the parameters
\begin{equation}\label{eq:n0}
 n_0:=2p+\lceil\log N\rceil,\quad\ q_{\max}:=p\,(2n_0+2),\quad\ \tau:=C_m\,(9p^2)^{2\fa+1},\quad\ \varphi_0:=\big(5e\,C_m\,(13p^3)^{2\fa+\frac32}\big)^{-1},
\end{equation}
and for $L\subset\{1,\dots,N\}$ with $1\le|L|\le p$, $|L|<N$, the event and the random variable
\begin{equation}\label{eq:XiL}
 \Xi_L:=\{\Phi_L\le\varphi_0\}\cap\bigcap_{a\in L}\big\{|z+m_a^{(L)}|\ge\tfrac25\big\},\qquad\rho:=6\,\tau\,\Phi_L .
\end{equation}
By Lemma~\ref{lem:FAstruct}(b),(c), $\Xi_L$ is closed and $\bone_{\Xi_L}$, $\rho$ do not depend on $\mathcal Q(L)$. The expansion of Section~\ref{sec:FAexp} is used with $n_0$ from \eqref{eq:n0}. Elementary consequences of $2\le\log N\le p$, used below:
\begin{equation}\label{eq:elem}
 \lceil\log N\rceil\le p+1,\quad\ n_0\le3p+1,\quad\ q_{\max}\le8p^2,\quad\ q_{\max}+\log N\le9p^2,\quad\ (q_{\max}+1)\lceil\log N\rceil\le13p^3 .
\end{equation}
(Indeed $2n_0+2\le6p+4\le8p$, and $(8p^2+1)(p+1)=8p^3+8p^2+p+1\le(8+4+\frac14+\frac18)p^3$ for $p\ge2$.)

\begin{lemma}[Size of the terms]\label{lem:FAsize}
Let $L$ be as above, $r=|L|$, $k\in L$. At every point of $\Xi_L$:
\begin{enumerate}[label=\textup{(\alph*)},leftmargin=2.2em,itemsep=1pt]
\item $p\,\rho\le\frac18$;
\item $\norm{Q_kW_k(\pi)}_{p,L}\le2\,\rho^{\,o(\pi)}$ for every $\pi\in\mathcal T(k)$;
\item $\norm{Q_kR_k}_{p,L}\le2\,(r\rho)^{2p+2}$.
\end{enumerate}
\end{lemma}
\begin{proof}
(a) On $\Xi_L$, by \eqref{eq:n0},
$p\rho\le6p\,\tau\,\varphi_0=\frac{6}{5e}\cdot\frac{9^{2\fa+1}}{13^{2\fa+3/2}}\cdot p^{\,1+(4\fa+2)-(6\fa+\frac92)}\le\frac12\,p^{-(2\fa+\frac32)}\le\frac12\,2^{-7/2}\le\frac18$.

Preparations for (b), (c). Let $2\le q\le q_{\max}$. By (m1) of Lemma~\ref{lem:FAmom} and \eqref{eq:elem}, for $a\ne b\in L$,
\begin{equation}\label{eq:size1}
 \norm{M_{ab}}_{q,L}\le C_m(q_{\max}+\log N)^{2\fa+1}\Phi_L\le\tau\,\Phi_L .
\end{equation}
At a point of $\Xi_L$ the hypothesis \eqref{eq:m3hyp} holds for every $a\in L$ with $q_*:=q_{\max}$, because $5eC_m((q_{\max}+1)\lceil\log N\rceil)^{2\fa+3/2}\Phi_L\le5eC_m(13p^3)^{2\fa+3/2}\varphi_0=1$ by \eqref{eq:elem}. Hence, by (m3),
\begin{equation}\label{eq:size2}
 \norm{1/M_{aa}}_{q,L}\le6\qquad(a\in L,\ 1\le q\le q_{\max}).
\end{equation}

(b) For $\pi=\emptyset$: $Q_kW_k(\emptyset)=Q_kM_{kk}$, and by (m2) with $q=p$: $\norm{Q_kM_{kk}}_{p,L}\le C_m(p+\log N)^{2\fa+3/2}\Phi_L\le C_m(2p)^{2\fa+3/2}\Phi_L\le\tau\Phi_L\le\rho$ (as $(2p)^{2\fa+3/2}\le(9p^2)^{2\fa+1}$).
Let $\pi\in\Pi_n(k)$ with $0\le n\le n_0-1$. $W_k(\pi)$ is a product of $2n+3$ factors: $n+2$ factors $M_{ab}$ with $a\ne b$ (namely $M_{ka_0}$, $M_{a_ta_{t+1}}$, $M_{a_nk}$; recall $k\notin L'$ and $a_t\ne a_{t+1}$) and $n+1$ factors $1/M_{aa}$.
By H\"older's inequality (F3) with $2n+3$ factors, each in $L^{q}$ with $q:=p(2n+3)\in[2,q_{\max}]$, and \eqref{eq:size1}, \eqref{eq:size2}:
\[
 \norm{W_k(\pi)}_{p,L}\le(\tau\Phi_L)^{n+2}\,6^{\,n+1}\le\rho^{\,n+2}=\rho^{\,o(\pi)} .
\]
By (F3), $\norm{Q_kW_k(\pi)}_{p,L}\le2\norm{W_k(\pi)}_{p,L}$.

(c) If $r=1$, $R_k=0$. Let $r\ge2$. Each of the at most $r^{n_0+2}$ terms of \eqref{eq:Rk} (see \eqref{eq:FAcountpaths}) is a product of the deterministic-size factor $G^{(k)}_{a_{n_0}b}$, with $|G^{(k)}_{a_{n_0}b}|\le\eta^{-1}\le N$ (Lemma~\ref{lem:FAstruct}(a)), of $n_0+2$ factors $M_{ab}$ with $a\ne b$
(namely $M_{ka_0}$, $M_{a_ta_{t+1}}$ for $t<n_0$, and $M_{bk}$), and of $n_0$ factors $1/M_{aa}$. By H\"older's inequality (F3) with $2n_0+2$ factors, each in $L^{q_{\max}}$, and \eqref{eq:size1}, \eqref{eq:size2}, each term has $\norm{\cdot}_{p,L}\le N(\tau\Phi_L)^{n_0+2}6^{n_0}\le N\rho^{n_0+2}$.
By Minkowski's inequality (F3), \eqref{eq:n0} and (a) (which gives $r\rho\le p\rho\le\frac18\le e^{-1}$, so $N(r\rho)^{\lceil\log N\rceil}\le1$):
\[
 \norm{R_k}_{p,L}\le N\,(r\rho)^{n_0+2}=N(r\rho)^{\lceil\log N\rceil}\,(r\rho)^{2p+2}\le(r\rho)^{2p+2},
\]
and $\norm{Q_kR_k}_{p,L}\le2\norm{R_k}_{p,L}$.
\end{proof}

\paragraph{Index tuples.} For $\mathbf k=(k_1,\dots,k_p)\in\{1,\dots,N\}^p$ let $L(\mathbf k):=\{k_1,\dots,k_p\}$ and $r(\mathbf k):=|L(\mathbf k)|\le p$. A position $s\in\{1,\dots,p\}$ is \emph{lone} if $k_s\ne k_{u}$ for all $u\ne s$; $\ell(\mathbf k)$ is the number of lone positions.
For a function $f$ we write $f^{\#_s}:=f$ if $s\le p/2$ and $f^{\#_s}:=\ol f$ if $s>p/2$.

\begin{lemma}[One index tuple]\label{lem:FAtuple}
Let $\mathbf k\in\{1,\dots,N\}^p$, $L:=L(\mathbf k)$, $\ell:=\ell(\mathbf k)$, and assume $p<N$. At every point of $\Xi_L$:
\begin{equation}\label{eq:FAtuple}
 \Big|P_L\Big(\prod_{s=1}^pX_{k_s}^{\#_s}\Big)\Big|\le\big(8\,p\,\rho\big)^{p+\ell}=\big(48\,p\,\tau\,\Phi_L\big)^{p+\ell}.
\end{equation}
\end{lemma}
\begin{proof}
Let $r:=|L|$ and $x:=r\rho$; by Lemma~\ref{lem:FAsize}(a), $x\le p\rho\le\frac18$ on $\Xi_L$. Let $\mathfrak r$ be a symbol (``remainder'') and, for $k\in L$ and $\xi\in\mathcal T(k)\cup\{\mathfrak r\}$,
\[
 Y_k(\xi):=\epsilon(\xi)\,Q_kW_k(\xi)\ \ (\xi\in\mathcal T(k)),\qquad Y_k(\mathfrak r):=-Q_kR_k .
\]
By \eqref{eq:FAX}, $X_k=\sum_{\xi\in\mathcal T(k)\cup\{\mathfrak r\}}Y_k(\xi)$ (a finite sum of functions of $\mathcal A$). Multiplying out,
\begin{equation}\label{eq:tuple1}
 P_L\Big(\prod_{s=1}^pX_{k_s}^{\#_s}\Big)=\sum_{\boldsymbol\xi\in\mathcal X}P_L\Big(\prod_{s=1}^pY_{k_s}(\xi_s)^{\#_s}\Big),\qquad\mathcal X:=\prod_{s=1}^p\big(\mathcal T(k_s)\cup\{\mathfrak r\}\big).
\end{equation}
We split $\mathcal X=\mathcal X_0\cup\mathcal X_1\cup\mathcal X_2$ (disjoint), where $\mathcal X_2$ is the set of $\boldsymbol\xi$ with $\xi_s=\mathfrak r$ for at least one $s$; $\mathcal X_0$ is the set of $\boldsymbol\xi\in\prod_s\mathcal T(k_s)$ for which there is a lone position $s$ with
$k_s\notin\bigcup_{u\ne s}V(\xi_u)$; and $\mathcal X_1:=\prod_s\mathcal T(k_s)\setminus\mathcal X_0$.

\emph{Terms in $\mathcal X_0$ vanish.} If $s$ is lone, then $k_s\ne k_u$ for all $u\ne s$; so for $\boldsymbol\xi\in\mathcal X_0$ the condition \eqref{eq:lonecond} holds, and Lemma~\ref{lem:FAvanish}(b) (with $j=p$) gives $P_L(\prod_sY_{k_s}(\xi_s)^{\#_s})=0$.

\emph{Terms in $\mathcal X_1$ have total order at least $p+\ell$.} Let $\boldsymbol\xi\in\mathcal X_1$. For every lone position $s$ there is $u\ne s$ with $k_s\in V(\xi_u)$. The values $k_s$ at the $\ell$ lone positions are pairwise distinct, so $|\bigcup_uV(\xi_u)|\ge\ell$ and, by \eqref{eq:FAcountpaths},
\begin{equation}\label{eq:tuple2}
 \sum_{s=1}^po(\xi_s)\ \ge\ \sum_{s=1}^p\big(|V(\xi_s)|+1\big)\ \ge\ p+\ell .
\end{equation}
By H\"older's inequality (F3) with $p$ factors in $L^p$ and Lemma~\ref{lem:FAsize}(b), $|P_L(\prod_sY_{k_s}(\xi_s)^{\#_s})|\le\prod_s\norm{Y_{k_s}(\xi_s)}_{p,L}\le\prod_s2\rho^{o(\xi_s)}$. We sum over $\boldsymbol\xi\in\mathcal X_1$ by first fixing the orders $o_s=o(\xi_s)=:1+e_s$, $e_s\ge0$;
by \eqref{eq:FAcountpaths} there are at most $\prod_sr^{e_s}$ tuples $\boldsymbol\xi$ with given $(e_s)$. Hence, using \eqref{eq:tuple2},
\begin{align}
 \sum_{\boldsymbol\xi\in\mathcal X_1}\prod_{s}\norm{Y_{k_s}(\xi_s)}_{p,L}&\le2^p\sum_{\substack{(e_s)\in\Z_{\ge0}^p\\ \sum_se_s\ge\ell}}\ \prod_{s=1}^p\rho\,(r\rho)^{e_s}=2^p\rho^p\sum_{j\ge\ell}\binom{j+p-1}{p-1}x^j\nonumber\\
 &\le2^p\rho^p\sum_{j\ge\ell}2^{j+p-1}x^j=2^p\rho^p\,\frac{2^{p-1}(2x)^\ell}{1-2x}\le(4\rho)^p(2x)^\ell\le(4p\rho)^{p+\ell}.\label{eq:tuple3}
\end{align}
Here $\binom{j+p-1}{p-1}$ is the number of $(e_s)\in\Z_{\ge0}^p$ with $\sum e_s=j$; it is at most $2^{j+p-1}$; and $2x\le\frac14$, $x\le p\rho$, $p\ge1$.

\emph{Terms in $\mathcal X_2$.} For every $s$, by Lemma~\ref{lem:FAsize}(b),(c) and \eqref{eq:FAcountpaths},
\[
 \sum_{\xi\in\mathcal T(k_s)}\norm{Y_{k_s}(\xi)}_{p,L}\le\sum_{o\ge1}r^{o-1}\,2\rho^{\,o}=\frac{2\rho}{1-x}\le3\rho=:b,\qquad\norm{Y_{k_s}(\mathfrak r)}_{p,L}\le2x^{2p+2}=:a .
\]
By H\"older's inequality as before, and since $\boldsymbol\xi\in\mathcal X_2$ means that the set $J:=\{s:\xi_s=\mathfrak r\}$ is non-empty,
\[
 \sum_{\boldsymbol\xi\in\mathcal X_2}\prod_s\norm{Y_{k_s}(\xi_s)}_{p,L}\le\sum_{\emptyset\ne J\subset\{1..p\}}a^{|J|}b^{p-|J|}=(a+b)^p-b^p\le p\,a\,(a+b)^{p-1}\le p\,2^{1-p}a\le a,
\]
using $a+b\le2\cdot8^{-6}+\frac38\le\frac12$ and $p\,2^{1-p}\le1$. Finally $a=2x^2\,x^{2p}\le x^{2p}\le(p\rho)^{p+\ell}$, because $x\le p\rho\le1$ and $\ell\le p$.

\emph{Conclusion.} By \eqref{eq:tuple1} and the three cases, $|P_L(\prod_sX_{k_s}^{\#_s})|\le(4p\rho)^{p+\ell}+(p\rho)^{p+\ell}\le2(4p\rho)^{p+\ell}\le(8p\rho)^{p+\ell}$.
\end{proof}

\begin{lemma}[Counting tuples]\label{lem:FAcount}
For $0\le\ell\le p$: $\ \#\{\mathbf k\in\{1,\dots,N\}^p:\ \ell(\mathbf k)=\ell\}\le p^p\,N^{(p+\ell)/2}$. Moreover $r(\mathbf k)\le\ell(\mathbf k)+\frac12(p-\ell(\mathbf k))$.
\end{lemma}
\begin{proof}
To $\mathbf k$ associate the partition $\mathscr P(\mathbf k)$ of $\{1,\dots,p\}$ into the non-empty level sets of $s\mapsto k_s$. It has $r(\mathbf k)$ blocks; the lone positions are its blocks of size one, and all other blocks have size at least two, so $r(\mathbf k)\le\ell+\frac12(p-\ell)$ if $\ell(\mathbf k)=\ell$.
The number of tuples $\mathbf k$ with a given partition with $r$ blocks is at most $N^r\le N^{(p+\ell)/2}$ (choose the value of $k$ on each block). The number of partitions of $\{1,\dots,p\}$ is at most $p^p$, since a partition is determined by the map $\{1..p\}\to\{1..p\}$ sending $s$ to the smallest element of its block.
\end{proof}

\subsection{Proof of Theorem \ref{thm:FA}}\label{sec:FAproof}
\paragraph{Constants.} Recall $C_m$ from \eqref{eq:Cm}. We prove Theorem~\ref{thm:FA} with
\begin{equation}\label{eq:CF}
 C_F:=2\,\big(48\sqrt2\,C_0\,C_m\,9^{2\fa+1}\big)^2,\qquad c_F=8\fa+7,
\end{equation}
and with $N_0$ the smallest integer such that for all $N\ge N_0$:
\begin{equation}\label{eq:N0FA}
 N\ge8,\qquad12\,C_0\,N^{1/80}\le N,\qquad3\,(4C_\vth)^{2\fa}N^{1/20}\le N.
\end{equation}
(The condition $p\ge p_0$ of \eqref{eq:FAhyp} is a hypothesis on $p$ rather than a condition on $N$; together with $p\ge\log N$ it would follow from $\log N\ge p_0$, but
in the application, Corollary~\ref{cor:FA}, it follows from the much weaker condition $(\log N)^3\ge p_0$, since there $p\ge(\log N)^3$.)

\begin{proof}[Proof of Theorem~\ref{thm:FA}]
Let $N\ge N_0$, and let $z,p,(t_k)$ satisfy the hypotheses. Then $p\ge\log N\ge2$, $p\le N^{1/(40\fa+40)}\le N^{1/80}$, so $12C_0p\le N$ (in particular $p<N$), and $3(4C_\vth p)^{2\fa}\le3(4C_\vth)^{2\fa}N^{2\fa/(40\fa+40)}\le N$; moreover $N\eta\ge p$.
Thus the standing assumptions of Section~\ref{sec:FAsize} and the hypotheses of Lemmas~\ref{lem:FAevent}, \ref{lem:FAcrude}, \ref{lem:FAtuple} hold for every $L$ with $1\le|L|\le p$.

\emph{Step 1: expansion of the $p$-th moment.} Since $p$ is even, $|w|^p=w^{p/2}\,\ol w^{\,p/2}$ for $w\in\C$; with $w=\sum_kt_kX_k$ and $X_k\in\mathcal A$,
\begin{equation}\label{eq:pf1}
 \EE\Big|\sum_kt_kX_k\Big|^p=\sum_{\mathbf k\in\{1..N\}^p}\ \prod_{s=1}^pt_{k_s}^{\#_s}\ \cdot\ \EE\Big[\prod_{s=1}^pX_{k_s}^{\#_s}\Big].
\end{equation}

\emph{Step 2: one tuple.} Fix $\mathbf k$ and let $L:=L(\mathbf k)$, $\ell:=\ell(\mathbf k)$. Since $\bone_{\Xi_L}$ does not depend on $\mathcal Q(L)$, (F2) gives
\begin{equation}\label{eq:pf2}
 \EE\Big[\prod_sX_{k_s}^{\#_s}\Big]=\EE\Big[\bone_{\Xi_L}\,P_L\Big(\prod_sX_{k_s}^{\#_s}\Big)\Big]+\EE\Big[\bone_{\Xi_L^c}\prod_sX_{k_s}^{\#_s}\Big].
\end{equation}
For the first term we use Lemma~\ref{lem:FAtuple} and Lemma~\ref{lem:FAevent}(a) ($\Phi_L\le\sqrt2\,C_0\Phi_*$): with
\[
 \mathfrak p:=48\sqrt2\,C_0\,p\,\tau=48\sqrt2\,C_0C_m9^{2\fa+1}\,p^{4\fa+3}\ \ (\ge1),
\]
\begin{equation}\label{eq:pf3}
 \Big|\EE\Big[\bone_{\Xi_L}\,P_L\Big(\prod_sX_{k_s}^{\#_s}\Big)\Big]\Big|\le\EE\big[\bone_{\Xi_L}(48\,p\,\tau\,\Phi_L)^{p+\ell}\big]\le\EE\big[(\mathfrak p\,\Phi_*)^{p+\ell}\big].
\end{equation}
For the second term, by the Cauchy--Schwarz inequality, H\"older's inequality with $p$ factors and Lemma~\ref{lem:FAcrude},
\begin{equation}\label{eq:pf4}
 \Big|\EE\Big[\bone_{\Xi_L^c}\prod_sX_{k_s}^{\#_s}\Big]\Big|\le\PP(\Xi_L^c)^{1/2}\Big(\EE\prod_s|X_{k_s}|^2\Big)^{1/2}\le\PP(\Xi_L^c)^{1/2}\prod_s\norm{X_{k_s}}_{L^{2p}(\PP)}\le\PP(\Xi_L^c)^{1/2}\,N^{2p}.
\end{equation}

\emph{Step 3: $\Xi_p(z)\subset\Xi_L$.} By Lemma~\ref{lem:FAevent}(b),(c), on $\Xi_p(z)$ we have $|z+m_a^{(L)}|\ge\frac25$ for all $a\in L$, and $\Phi_L^2\le25C_0^2/(N\eta)\le25C_0^2\,p^{-(12\fa+11)}$ by \eqref{eq:FAhyp}. By \eqref{eq:n0}, $\varphi_0^2=(5eC_m13^{2\fa+3/2})^{-2}p^{-(12\fa+9)}$, so
$\Phi_L\le\varphi_0$ as soon as $p\ge5C_0\cdot5eC_m13^{2\fa+3/2}=p_0$, which holds by \eqref{eq:FAhyp}. Hence $\Xi_p(z)\subset\Xi_L$ and $\PP(\Xi_L^c)\le\PP(\Xi_p(z)^c)$, for every $\mathbf k$.

\emph{Step 4: summation over tuples.} Since $|t_k|\le N^{-1}$, $\sum_{\mathbf k}\prod_s|t_{k_s}|\le1$, so by \eqref{eq:pf4} and Step~3 the second terms of \eqref{eq:pf2} contribute to \eqref{eq:pf1} at most $N^{2p}\PP(\Xi_p(z)^c)^{1/2}$.
For the first terms we group the tuples according to $\ell=\ell(\mathbf k)$ and use Lemma~\ref{lem:FAcount}, \eqref{eq:pf3}, and $N^{-1/2}\le\Phi_*\le\mathfrak p\,\Phi_*$ (Lemma~\ref{lem:FAevent}(a)):
\begin{align*}
 \sum_{\mathbf k:\,\ell(\mathbf k)=\ell}\ \prod_s|t_{k_s}|\ \EE\big[(\mathfrak p\Phi_*)^{p+\ell}\big]&\le p^p\,N^{\frac{p+\ell}2-p}\,\EE\big[(\mathfrak p\Phi_*)^{p+\ell}\big]\\
 &=p^p\,\EE\big[\big(N^{-1/2}\big)^{p-\ell}(\mathfrak p\Phi_*)^{p+\ell}\big]\le p^p\,\EE\big[(\mathfrak p\Phi_*)^{2p}\big].
\end{align*}
Summing over $\ell\in\{0,\dots,p\}$ and using $p+1\le2^p$,
\[
 \EE\Big|\sum_kt_kX_k\Big|^p\le(p+1)\,p^p\,\mathfrak p^{2p}\,\EE\big[\Phi_*^{2p}\big]+N^{2p}\,\PP\big(\Xi_p(z)^c\big)^{1/2}\le\big(2\,p\,\mathfrak p^2\big)^p\,\EE\big[\Phi_*^{2p}\big]+N^{2p}\,\PP\big(\Xi_p(z)^c\big)^{1/2},
\]
and $2p\,\mathfrak p^2=C_F\,p^{8\fa+7}$ by \eqref{eq:CF}. This is \eqref{eq:FA}.
\end{proof}

\begin{remark}[Comments on the proof]\label{rem:FAcheck}
(1) \emph{Where the powers of $p$ come from.} The exponent $c_F=8\fa+7=2(4\fa+3)+1$: a factor $p^{4\fa+2}$ in $\tau$ (moments of order up to $q_{\max}\le8p^2$ of bilinear forms, each costing $(q+\log N)^{2\fa+1}$; this is the only place where $p^2$ enters),
a factor $p$ from the number of paths ($r\le p$ choices per step), squared because the bound is $\Phi_*^{2p}$; and a factor $p^p$ from the number of partitions. The exponent $12\fa+11$ in \eqref{eq:FAhyp} is used only in Step~3 (it could be replaced by $12\fa+10$, at the price of the stronger hypothesis $p\ge p_0^2$ instead of $p\ge p_0$); it comes from
$\varphi_0\asymp p^{-(6\fa+9/2)}$, i.e.\ from the order $q'\le13p^3$ of the moment used for the tail of $Q_aM_{aa}$ in (m3).
(2) \emph{The inverse diagonal entries.} The conditional moments of $1/M_{aa}$ are controlled by the deterministic bound $|1/M_{aa}|\le N$ together with a conditional tail bound for $Q_aM_{aa}$ obtained from Markov's inequality at order $q'=q_{\max}\lceil\log N\rceil$; no large-deviation event at a fixed level is involved, so there is no competition between $p$ and the level of such an event.
(3) \emph{The bad set} enters only through the additive term $N^{2p}\PP(\Xi_p(z)^c)^{1/2}$; the event $\Xi_L$ on which the conditional bounds are used is measurable with respect to the minor $H^{(L)}$ and contains $\Xi_p(z)$ deterministically. This is why $\Phi_*$ contains $p/(N\eta)$ (from $\im m^{(L)}\le\im m+|L|/(N\eta)$).
(4) \emph{Unused hypotheses.} The hypothesis $|z|\le N$ of Theorem~\ref{thm:FA} is not used in the proof (the crude bound of Lemma~\ref{lem:FAcrude} is for $X_k=h_{kk}-Z_k$ and does not involve $z$); the upper bound on $p$ is used only through $12C_0p\le N$ and $3(4C_\vth p)^{2\fa}\le N$; the lower bound $p\ge\log N$ is used through \eqref{eq:elem} (the remainder needs $n_0-2p\ge\log N$ terms, and $q+\log N\le Cp^2$), and the hypothesis $p\ge p_0$ only in Step~3.
(5) \emph{Sanity checks.} For $r=1$ (all $k_s$ equal) Lemma~\ref{lem:FAtuple} reduces to $\norm{Q_kM_{kk}}_{p,L}^p\le\rho^p$. For two lone indices the expansion is \eqref{eq:two2} and the surviving term is \eqref{eq:two3}.
The algebraic identities \eqref{eq:1overGL}, \eqref{eq:expansion}, Lemma~\ref{lem:FAstruct}(d) and the vanishing of Lemma~\ref{lem:FAvanish}(b) were also checked numerically (the last two exactly, by enumeration over a finite product law), and a Monte Carlo experiment illustrates the sizes $\EE|\avg X|^2\lesssim\Phi_*^4$, $\EE|X_k|^2\approx\Phi_*^2$.
\end{remark}

\subsection{Proof of Corollary \ref{cor:FA}}\label{sec:FAcor}
\begin{proof}[Proof of Corollary~\ref{cor:FA}]
Let $K$ be as in \eqref{eq:Krange}, $z=E+\ii\eta\in\mathbb L$ and $p\in\{p_1,p_2\}$, where $p_1=2\lceil K^{3/2}/2\rceil$, $p_2=2\lceil K/2\rceil$. We apply Theorem~\ref{thm:FA} with $t_k:=1/N$, so that $\sum_kt_kX_k=\avg X$.
In the following, ``for $N\ge N_0$'' refers to finitely many conditions depending only on the model parameters, which we add to \eqref{eq:N0FA}.

\emph{Hypotheses of Theorem~\ref{thm:FA}.} (i) $p$ is an even integer and $K\le p_2\le p\le p_1\le K^{3/2}+2\le2K^{3/2}$. (ii) $p\ge K\ge(\log N)^3\ge\log N$, and $(\log N)^3\ge p_0$ for $N\ge N_0$, so $p\ge p_0$. (iii) By \eqref{eq:Krange}, $p\le2N^{3c_0/2}$, and $\frac32c_0<\frac1{40\fa+40}$ for the value of $c_0$ in \eqref{eq:ledger0}
(indeed $\frac32c_0=\frac{3}{16C_3}=\frac1{128\fa+192}$); hence $p\le N^{1/(40\fa+40)}$ for $N\ge N_0$. (iv) On $\mathbb L\subset\bD^{\rm near}\cup\bD^{\rm far}$ we have $\eta\le10$ and $\eta\ge\min\{K^{C_1-1}N^{-1},\frac12N^{-1/2}\}\ge N^{-1}$.
(v) If $z\in\bD^{\rm near}$: $N\eta\ge K^{C_1-1}=K^{18\fa+19}$, while $p^{12\fa+11}\le2^{12\fa+11}K^{18\fa+\frac{33}2}\le K^{18\fa+17}$ for $N\ge N_0$ (as $K^{1/2}\ge(\log N)^{3/2}\ge2^{12\fa+11}$). If $z\in\bD^{\rm far}$: $N\eta\ge\frac12N^{1/2}$, while
$p^{12\fa+11}\le N^{(12\fa+11)/(40\fa+40)}\le N^{3/10}\le\frac12N^{1/2}$. So $N\eta\ge p^{12\fa+11}$ in both cases. (vi) $|z|\le\Emax+11\le K^{2\fa+1}+11\le N$ (Definition~\ref{def:levels} and \eqref{eq:Krange}).
Hence \eqref{eq:FAhyp} holds.

\emph{Conclusion.} Taking $p$-th roots in \eqref{eq:FA} and using $(u+v)^{1/p}\le u^{1/p}+v^{1/p}$ for $u,v\ge0$,
\[
 \norm{\avg X(z)}_p\le C_Fp^{c_F}\,\norm{\Phi_*(z)^2}_p+N^2\,\PP\big(\Xi_p(z)^c\big)^{1/(2p)} .
\]
Here $C_Fp^{c_F}\le C_Fp_1^{c_F}=\mathfrak F$. Since $p\le p_1$, the event $\Xi(z)$ of \eqref{eq:Xiev} is contained in $\Xi_p(z)$, so by Theorem~\ref{thm:weak}(c), $\PP(\Xi_p(z)^c)\le\PP(\Xi(z)^c)\le C_wN^9e^{-K^2}$ with a constant $C_w\ge1$ depending only on the model parameters. By $p\ge\log N$, $p\le p_1$ and
$K^2=K_w\ge40p_1\log N$ \eqref{eq:levelrel},
\[
 N^2\,\PP\big(\Xi_p(z)^c\big)^{1/(2p)}\le N^2\,C_w\,e^{\frac{9\log N}{2p}}\,e^{-\frac{K^2}{2p_1}}\le C_w\,e^{9/2}\,N^{2-20}\le N^{-8}\qquad(N\ge N_0).
\]
Finally $\mathfrak F=C_Fp_1^{8\fa+7}\le C_F2^{8\fa+7}K^{12\fa+\frac{21}2}\le K^{12\fa+11}$ for $N\ge N_0$, because $K^{1/2}\ge(\log N)^{3/2}\ge C_F2^{8\fa+7}$.
\end{proof}

\begin{remark}[On the statements used from other sections]\label{rem:FAext}
(1) Lemma~\ref{lem:LDmom} is used for all real $q$ in $[2,13p^3]$ and for families of $n=N-|L|\le N$ variables, only through the constant $C_{\rm LD}$ in \eqref{eq:Cm}. A different value of this constant
changes $C_m$, $C_F$, $C_V$, $p_0$ and $N_0$, and nothing else.
(2) Lemma~\ref{lem:split} is used only through (F1), (F2), (F4) of Lemma~\ref{lem:FAfub}, i.e.\ as Fubini's theorem for the splittings $\mathcal P_N=\mathcal Q^c\cup\mathcal Q$ and $\mathcal Q=(\mathcal Q\setminus\mathcal Q_1)\cup\mathcal Q_1$.
(3) Corollary~\ref{cor:FA} uses the numerical values $C_1-1\ge18\fa+17$ and $c_0<\frac{2}{3(40\fa+40)}$ of the ledger \eqref{eq:ledger0} (see Appendix~\ref{app:ledger}).
\end{remark}

\section{The strong local law: proof of Theorem \ref{thm:K}\,(i),(ii)}\label{sec:strong}
In this section and the next, $N\ge N_0$, $\kk\in\{\R,\C\}$, $\nu\in\mathcal E_N(\vth,\dpm,C_0)$ (normalised as in Convention~\ref{conv:const}) and $K$ in the range \eqref{eq:Krange} are fixed, as in Section~\ref{sec:weak};
$K_w,p_1,p_2,\fL,\Emax,\cG(z)$ are as in Definition~\ref{def:levels}; $\bD^{\rm near}$, $\bD^{\rm far}$, $\mathbb L$, ${\rm col}(z)$, $\cG_{\rm col}(z)$ as in Definition~\ref{def:domains};
$\lambda_{\rm st}=1/(576C_s)$ as in \eqref{eq:lambda1}; $w=\avg v$, $w_0=\msc(1-\msc^2)$ as in Lemma~\ref{lem:stabnear}; and $\avg X=\frac1N\sum_kX_k$ with $X_k$ from Definition~\ref{def:XA}.
For a random variable $f$ on $(\Omega,\PP)$ and real $p\ge1$, $\norm f_p:=(\EE|f|^p)^{1/p}$. All quantities depend on the spectral parameter $z=E+\ii\eta$; recall $\kappa=||E|-2|$ and $\varkappa=\sqrt{\kappa+\eta}$ (Convention~\ref{conv:z}).

\paragraph{Exponents of this section.} Throughout,
\[
 b:=12\fa+12,\qquad\mathfrak E:=K^{b},\qquad\mathfrak E_3:=K^{b+3}.
\]
From the previous sections we use the two quantities $\mathfrak E_0=6\cdot10^6C_s^2\fL^2$ of Lemma~\ref{lem:apriori} and $\mathfrak F=C_Fp_1^{c_F}$ of Corollary~\ref{cor:FA}, with the bounds stated there:
\begin{equation}\tag{T1}\label{eq:a5-T1}
 1\le\mathfrak E_0\le K^{8\fa+7},\qquad1\le\mathfrak F\le K^{12\fa+11},\qquad K^{3/2}\le p_1\le2K^{3/2},\qquad K\le p_2\le K+2\le p_1 .
\end{equation}
(The lower bounds hold because $C_s,\fL,C_F,p_1\ge1$; the bounds on $p_1,p_2$ are Definition~\ref{def:levels} and Lemma~\ref{lem:levels}(S3), and $K+2\le K^{3/2}$ as $K\ge4$.)

\paragraph{Standing conditions.} ``For $N\ge N_0$'' has the same meaning as in Section~\ref{sec:weak}: (S1)--(S5) of Lemma~\ref{lem:levels} hold with a constant $K_0$ which is at least the maximum of the explicit constants of Sections~\ref{sec:weak}--\ref{sec:norm};
for the present two sections it suffices that
\begin{equation}\tag{T2}\label{eq:a5-T2}
 K\ \ge\ K_0\ \ge\ \max\{124800,\ 2142\,C_s\},\qquad N^{31}\ \ge\ 920\,\fa\,e^{8},
\end{equation}
that $N_0$ is at least the thresholds of Lemma~\ref{lem:norm0}, Theorem~\ref{thm:FA} and Corollary~\ref{cor:FA}, and that the finitely many further inequalities marked ``$N\ge N_0$'' below hold; each of them involves only $N$ and absolute constants.
The complete list of thresholds is collected in Appendix~\ref{app:ledger}. Of the exponent ledger \eqref{eq:ledger0} we use, besides \eqref{eq:ledgerA3} and the two inequalities of Remark~\ref{rem:FAext}(3) (which are hypotheses of the statements we quote), only
\begin{equation}\tag{T3}\label{eq:a5-T3}
 C_2\ge b+7,\qquad C_3\ge2b+12,\qquad c_0\,(C_1-1)\le\tfrac13 ;
\end{equation}
all of these hold for the values \eqref{eq:ledger0} (Appendix~\ref{app:ledger}). We record three consequences which are used repeatedly.

\smallskip\noindent\textbf{(T4)} \emph{Sizes on $\bD^{\rm near}$.} For $z\in\bD^{\rm near}$: $\eta\ge N^{-1}$, $\ N\eta\ge K^{C_1-1}\ge K^{12\fa+18}\ge\mathfrak F$, $\ \varkappa\ge\sqrt\eta\ge N^{-1/2}$, and, with $c_*=\frac1{150}$, $C_*=4$ of Lemma~\ref{lem:mscsize},
\[
 \frac{\eta}{150}\le\im\msc\le\varkappa,\qquad\frac{\varkappa}{1950}\le|w_0|\le4\varkappa,\qquad\Lambda\le\frac1\eta+1\le2N .
\]
Indeed, the first three statements are Definition~\ref{def:domains} with \eqref{eq:ledgerA3}, \eqref{eq:a5-T1} and $\kappa\ge0$; the bounds on $\im\msc$ are the ``in particular'' line of Lemma~\ref{lem:mscsize}; those on $w_0=\msc(1-\msc^2)$ follow from \eqref{eq:alpha}: $\frac1{13}\le|\msc|\le1$ and $c_*\varkappa\le|1-\msc^2|\le C_*\varkappa$, with $\frac{c_*}{13}=\frac1{1950}$ (this is also the last statement of Lemma~\ref{lem:stabnear}; it does not involve $H$);
and $\Lambda=|m-\msc|\le|m|+|\msc|\le\eta^{-1}+1$ by Lemma~\ref{lem:basic}(a) and Lemma~\ref{lem:msc}(b).

\smallskip\noindent\textbf{(T5)} \emph{Powers of $K$ against powers of $N$.} By Lemma~\ref{lem:levels}(S4), $K^{6\fa+6}\le N^{1/32}$; hence
\[
 \max\{\mathfrak F,\ \mathfrak E,\ K^{10\fa+8}\}\le K^{12\fa+12}\le N^{1/16},\qquad K^{b+2}=K^{12\fa+14}\le K^{18\fa+18}\le N^{3/32},
\]
and, by the third relation in \eqref{eq:a5-T3}, $K^{C_1-1}\le N^{c_0(C_1-1)}\le N^{1/3}$.

\smallskip\noindent\textbf{(T6)} \emph{Bad sets in $L^p$.} Let $z\in\mathbb L$ and $p\in\{p_1,p_2\}$. Then
\[
 2N\,\PP\big(\cG_{\rm col}(z)^c\big)^{1/p}\le N^{-8},\qquad\text{and}\qquad2N\,\big(N^{10}e^{-p_1}\big)^{1/p_2}\le N^{-8}.
\]
Indeed, by Definition~\ref{def:domains}, $\PP(\cG_{\rm col}(z)^c)\le460\,\fa N^8e^{-K^2}$. Since $\log N\le K\le p\le p_1$ and $K^2\ge40p_1\log N$ \eqref{eq:levelrel}: $(460\fa)^{1/p}\le460\fa$, $N^{8/p}=e^{8\log N/p}\le e^8$ and $e^{-K^2/p}\le e^{-K^2/p_1}\le N^{-40}$.
So the first quantity is at most $920\,\fa e^8N^{-39}\le N^{-8}$ by \eqref{eq:a5-T2}. Similarly, $N^{10/p_2}\le e^{10}$ and $e^{-p_1/p_2}\le N^{-40}$ because $p_1\ge40p_2\log N$ \eqref{eq:levelrel}; so the second quantity is at most $2e^{10}N^{-39}\le N^{-8}$ for $N\ge N_0$.

\smallskip We shall also use, without further comment, the elementary inequalities $\sqrt{\alpha+\beta}\le\sqrt\alpha+\sqrt\beta$ and $2\sqrt{\alpha\beta}\le t\alpha+t^{-1}\beta$ ($\alpha,\beta\ge0$, $t>0$), Minkowski's inequality $\norm{f+g}_p\le\norm f_p+\norm g_p$,
and the following two facts about a random variable $f\ge0$ and $p\ge1$:
\begin{equation}\tag{M}\label{eq:a5-M}
 \norm{f^{1/2}}_p\le\norm f_p^{1/2},\qquad\qquad\PP(f>t)\le\big(\norm f_p/t\big)^p\quad(t>0).
\end{equation}
(The first is $\EE f^{p/2}\le(\EE f^p)^{1/2}$, the Cauchy--Schwarz inequality; the second is Markov's inequality for $f^p$.)

\subsection{First pass: a moment inequality that improves itself}
\begin{proposition}[Upper-branch moment bound, near regime]\label{prop:S1}
Let $z\in\mathbb L\cap\bD^{\rm near}$ and $p=p_1$. Then
\begin{align}
 Y(z):=\norm{\Lambda(z)\bone_{\cG_{\rm col}(z)}}_{p}&\le C\Big(\varkappa+\frac{\mathfrak E}{N\eta}\Big),\label{eq:Y}\\
 \norm{\avg X(z)}_{p}&\le\frac{C\mathfrak F}{N\eta}\Big(\varkappa+\frac{\mathfrak E}{N\eta}\Big),\label{eq:Xmom}
\end{align}
and the event $\cG^1(z):=\cG_{\rm col}(z)\cap\big\{|\avg X(z)|\le\frac{K\mathfrak E}{N\eta}\big(\varkappa+\frac{\mathfrak E}{N\eta}\big)\big\}$ satisfies $\PP(\cG_{\rm col}(z)\setminus\cG^1(z))\le e^{-p_1}$.
The constants are absolute: \eqref{eq:Y} holds with $C=42$ and \eqref{eq:Xmom} with $C=194$. The set $\cG^1(z)$ is closed.
\end{proposition}
\begin{proof}
Write $u:=N\eta$, $p:=p_1$, $\cG_{\rm c}:=\cG_{\rm col}(z)$. Since $z\in{\rm col}(z)$, $\cG_{\rm c}\subset\cG(z)$, and by Theorem~\ref{thm:weak}(a), $\Ld+\Lo\le\lambda_{\rm st}/2$ on $\cG_{\rm c}$. So $\cG_{\rm c}$ is contained in the event $\cG(z)\cap\{\Ld+\Lo\le\lambda_{\rm st}\}$ of Lemma~\ref{lem:apriori}.
No information on the branch of the quadratic equation (i.e.\ on the size of $\Lambda$ relative to $|w_0|$) is used in this proof.

\emph{Step 1: a pointwise bound on $\cG_{\rm c}$.} Since $\Lambda=|w|$, we have $\Lambda^2=|w^2|\le|w_0|\Lambda+|w^2-w_0w|$, and \eqref{eq:ap3} gives
\[
 \tfrac14\Lambda^2\le\Big(|w_0|+\frac{\mathfrak E_0}{u}\Big)\Lambda+|\avg X|+\mathfrak E_0\Big(\frac{\im\msc}{u}+\frac1{u^2}\Big).
\]
By the last statement of Lemma~\ref{lem:scalar} (applied to $x=\Lambda$ after multiplication by $4$) and $\sqrt{\alpha+\beta+\gamma}\le\sqrt\alpha+\sqrt\beta+\sqrt\gamma$,
\[
 \Lambda\le4|w_0|+\frac{4\mathfrak E_0}{u}+2|\avg X|^{1/2}+2\Big(\frac{\mathfrak E_0\im\msc}{u}\Big)^{1/2}+\frac{2\mathfrak E_0^{1/2}}{u}.
\]
By (T4), $|w_0|\le4\varkappa$ and $\im\msc\le\varkappa$, so $2(\mathfrak E_0\im\msc/u)^{1/2}\le\varkappa+\mathfrak E_0/u$; and $\mathfrak E_0^{1/2}\le\mathfrak E_0$. Hence, on $\cG_{\rm c}$,
\begin{equation}\tag{7.a}\label{eq:a5-7a}
 \Lambda\le17\varkappa+\frac{7\mathfrak E_0}{u}+2|\avg X|^{1/2}.
\end{equation}

\emph{Step 2: the moment of $\avg X$ in terms of $Y$.} By Lemma~\ref{lem:weyl}(e) the functions $\Lambda(z)$, $\avg X(z)$, $\im m(z)$ are continuous on $\Omega$; $\cG_{\rm c}$ is closed (Definition~\ref{def:domains}). By (T4), $0\le\Lambda\le2N$, so $Y:=Y(z)\le2N<\infty$, and by (T6)
\begin{equation}\tag{7.b}\label{eq:a5-7b}
 \norm\Lambda_p\le\norm{\Lambda\bone_{\cG_{\rm c}}}_p+\norm{\Lambda\bone_{\cG_{\rm c}^c}}_p\le Y+2N\,\PP(\cG_{\rm c}^c)^{1/p}\le Y+N^{-8}.
\end{equation}
Next, $\im m\le\im\msc+|m-\msc|=\im\msc+\Lambda$, and $\frac1N\le\frac{150\,\im\msc}{u}$ by (T4). So the random variable \eqref{eq:Phistar} (with $p=p_1$) satisfies, at every point of $\Omega$,
\[
 0\le\Phi_*^2=\frac1N+\frac{\im m}{u}+\frac{p_1}{u^2}\le\frac{151\,\im\msc}{u}+\frac{\Lambda}{u}+\frac{p_1}{u^2},
\]
and by Minkowski's inequality, \eqref{eq:a5-7b} and $\im\msc\le\varkappa$: $\norm{\Phi_*^2}_p\le\frac{151\varkappa}u+\frac{Y+N^{-8}}u+\frac{p_1}{u^2}$. Corollary~\ref{cor:FA} (with $p=p_1$) and $\mathfrak F\le u$ (T4) now give
\begin{equation}\tag{7.c}\label{eq:a5-7c}
 \norm{\avg X}_p\le\mathfrak F\,\norm{\Phi_*^2}_p+N^{-8}\le\frac{\mathfrak F}{u}\Big(151\varkappa+Y+\frac{p_1}u\Big)+2N^{-8}.
\end{equation}

\emph{Step 3: closing the inequality.} Multiply \eqref{eq:a5-7a} by $\bone_{\cG_{\rm c}}$, take the $L^p$-norm, and use Minkowski's inequality and \eqref{eq:a5-M}:
\[
 Y\le17\varkappa+\frac{7\mathfrak E_0}{u}+2\norm{\avg X}_p^{1/2}.
\]
By \eqref{eq:a5-7c},
\begin{align*}
 2\norm{\avg X}_p^{1/2}&\le2\Big(\frac{151\mathfrak F\varkappa}{u}\Big)^{1/2}+2\Big(\frac{\mathfrak FY}{u}\Big)^{1/2}+\frac{2(\mathfrak Fp_1)^{1/2}}{u}+2\sqrt2\,N^{-4}\\
 &\le\varkappa+\frac{151\mathfrak F}{u}+\frac Y2+\frac{2\mathfrak F}{u}+\frac{2(\mathfrak Fp_1)^{1/2}}{u}+3N^{-4}.
\end{align*}
Therefore
\[
 Y\le\frac Y2+18\varkappa+\frac{7\mathfrak E_0+153\mathfrak F+2(\mathfrak Fp_1)^{1/2}}{u}+3N^{-4}.
\]
\emph{Here the right-hand side contains the unknown $Y$ itself with the factor $\frac12$; since $Y$ is finite, it can be absorbed.} This gives $Y\le36\varkappa+2\big(7\mathfrak E_0+153\mathfrak F+2(\mathfrak Fp_1)^{1/2}\big)/u+6N^{-4}$.
By \eqref{eq:a5-T1}, $(\mathfrak Fp_1)^{1/2}\le(2K^{12\fa+12.5})^{1/2}\le2K^{12\fa+11}$ and
\[
 14\mathfrak E_0+306\mathfrak F+4(\mathfrak Fp_1)^{1/2}\le328\,K^{12\fa+11}\le K^{12\fa+12}=\mathfrak E
\]
by \eqref{eq:a5-T2}; and $6N^{-4}\le6\varkappa$ by (T4). Hence $Y\le42\varkappa+\mathfrak E/u$, which proves \eqref{eq:Y} with $C=42$.

Inserting \eqref{eq:Y} into \eqref{eq:a5-7c}, and using $p_1\le2K^{3/2}\le\mathfrak E$ and $2N^{-8}\le\frac{1}{10}N^{-3/2}\le\frac{\mathfrak F\varkappa}{u}$ (as $\mathfrak F\ge1$, $\varkappa\ge N^{-1/2}$, $u\le10N$, and $N^{13/2}\ge20$),
\[
 \norm{\avg X}_p\le\frac{\mathfrak F}{u}\Big(151\varkappa+42\varkappa+\frac{42\mathfrak E}{u}+\frac{\mathfrak E}{u}\Big)+\frac{\mathfrak F\varkappa}{u}\le\frac{194\,\mathfrak F}{u}\Big(\varkappa+\frac{\mathfrak E}u\Big),
\]
which is \eqref{eq:Xmom} with $C=194$.

\emph{Step 4: Markov's inequality.} Let $t:=\frac{K\mathfrak E}{u}\big(\varkappa+\frac{\mathfrak E}u\big)>0$. By \eqref{eq:a5-M}, \eqref{eq:Xmom}, \eqref{eq:a5-T1} and \eqref{eq:a5-T2},
\[
 \PP\big(\cG_{\rm col}(z)\setminus\cG^1(z)\big)\le\PP\big(|\avg X(z)|>t\big)\le\Big(\frac{194\,\mathfrak F}{K\mathfrak E}\Big)^{p_1}\le\Big(\frac{194}{K^2}\Big)^{p_1}\le e^{-p_1}.
\]
Finally $\cG^1(z)$ is closed, as the intersection of the closed set $\cG_{\rm col}(z)$ with a set $\{|f|\le t\}$, $f=\avg X(z)$ continuous.
\end{proof}

\begin{remark}
This is the point where NP departs most from the original: there, the bound $\Lambda\lesssim(N\eta)^{-\tau}$ is improved to exponent $\frac{1+\tau}2$ and the step is iterated $2\log\log N$ times, each time losing a factor $(\log N)^{3}$ in the exponent of the probability
(Lemmas 4.3--4.4 of the original; this is the origin of $L\ge A_0\log\log N$). Here the inequality of Step 3 closes in one step because the right-hand side of \eqref{eq:FA} contains $\Lambda$ itself under the expectation (through $\im m\le\im\msc+\Lambda$ in $\Phi_*^2$, linearly),
and because the a priori bound $Y\le2N$ makes the absorption $Y\le\frac12Y+\dots$ legitimate. The square root in \eqref{eq:a5-7a} comes from solving the quadratic inequality on the \emph{upper} branch, $|w|\le|w_0|+|\,\cdot\,|^{1/2}$ (Lemma~\ref{lem:scalar}), which needs no dichotomy.
\end{remark}

\subsection{Dichotomy and continuity at the level of events}
\begin{proposition}[Averaged and entrywise law on the lattice]\label{prop:S2}
For $z\in\mathbb L\cap\bD^{\rm near}$ let $\cG^2(z):=\bigcap_{z'\in{\rm col}(z)}\cG^1(z')$; this set is closed, and $\PP(\cG^2(z)^c)\le N^{10}e^{-p_1}$. On $\cG^2(z)$:
\begin{enumerate}[label=\textup{(\alph*)},leftmargin=2.2em,itemsep=0pt]
\item if $\varkappa N\eta\ge\mathfrak E_3$:\quad $\Lambda(z)\le K^{b+2}/(N\eta)\le|w_0(z)|/4$;
\item in all cases:\quad $\Lambda(z)\le K^{b+4}/(N\eta)$;
\item $\Ld(z)+\Lo(z)\le K^{b+5}\big(\sqrt{\im\msc(z)/(N\eta)}+1/(N\eta)\big)$.
\end{enumerate}
\end{proposition}
\begin{proof}
Let ${\rm col}(z)=\{z_0,\dots,z_{k_*}\}$, $z_k=E+\ii\eta_k$, $\eta_k=10-kN^{-4}$, $z_{k_*}=z$, as in the justification of Definition~\ref{def:domains}; all $z_k$ lie in $\mathbb L\cap\bD^{\rm near}$, $k_*+1\le10N^4$, and ${\rm col}(z_k)=\{z_0,\dots,z_k\}\subset{\rm col}(z)$.

\emph{Probability.} Each $\cG^1(z_k)$ is closed (Proposition~\ref{prop:S1}), so $\cG^2(z)$ is closed. Since ${\rm col}(z_k)\subset{\rm col}(z)$ we have $\cG_{\rm col}(z)\subset\cG_{\rm col}(z_k)$, hence
\[
 \cG^2(z)^c\subset\cG_{\rm col}(z)^c\ \cup\ \bigcup_{k=0}^{k_*}\big(\cG_{\rm col}(z_k)\setminus\cG^1(z_k)\big),
\]
and by Definition~\ref{def:domains}, Proposition~\ref{prop:S1} and $K^2\ge p_1$ \eqref{eq:levelrel},
$\PP(\cG^2(z)^c)\le460\,\fa N^8e^{-K^2}+10N^4e^{-p_1}\le N^{10}e^{-p_1}$ for $N\ge N_0$ (namely $N^{2}\ge460\fa+1$).

\emph{The quadratic inequality on $\cG^1(z')$.} Let $z'=E+\ii\eta'\in{\rm col}(z)$; in the next display all quantities are evaluated at $z'$, $u:=N\eta'$ and
\[
 \beta:=\varkappa+\frac{\mathfrak E}{u}.
\]
We claim that at every point of $\cG^1(z')$
\begin{equation}\label{eq:Q1}
 |w^2-w_0w|\le\frac{2K\mathfrak E\beta}{u}+\frac{2\mathfrak E_0}{u}\Lambda+\tfrac34\Lambda^2 .
\end{equation}
Indeed $\cG^1(z')\subset\cG_{\rm col}(z')\subset\cG(z')\cap\{\Ld+\Lo\le\lambda_{\rm st}\}$ (Theorem~\ref{thm:weak}(a)), so \eqref{eq:ap3} holds; by the definition of $\cG^1(z')$, $|\avg X|\le K\mathfrak E\beta/u$; and, by $\im\msc\le\varkappa$ (T4) and $\mathfrak E_0\le\mathfrak E\le K\mathfrak E$ \eqref{eq:a5-T1},
\[
 \mathfrak E_0\Big(\frac{\im\msc}{u}+\frac1{u^2}\Big)\le\frac{\mathfrak E_0}{u}\Big(\varkappa+\frac{\mathfrak E}{u}\Big)\le\frac{K\mathfrak E\beta}{u}.
\]

\emph{Proof of \textup{(a)}.} Call a point $z'=E+\ii\eta'$ of the column \emph{of type \textup{(a)}} if $\varkappa(z')N\eta'\ge\mathfrak E_3$. Since $\kappa$ depends only on $E$, the function $\eta'\mapsto\varkappa(E+\ii\eta')\,N\eta'=\sqrt{\kappa+\eta'}\,N\eta'$ is increasing. Hence, if $z=z_{k_*}$ is of type (a), then all points $z_0,\dots,z_{k_*}$ of the column are of type (a).
Assume this, and fix a point of $\cG^2(z)$.

\smallskip\noindent\emph{Dichotomy.} Let $z'=z_k$ ($0\le k\le k_*$), write $u=N\eta_k$, and suppose that $\Lambda(z_k)\le\frac12|w_0(z_k)|$. We show that then
\begin{equation}\tag{7.d}\label{eq:a5-7d}
 \Lambda(z_k)\le\frac{K^{b+2}}{N\eta_k}\le\frac{|w_0(z_k)|}4 .
\end{equation}
All quantities are evaluated at $z_k$. Since $z_k$ is of type (a), $\mathfrak E/u\le\varkappa\mathfrak E/\mathfrak E_3=K^{-3}\varkappa\le\varkappa$, so $\beta\le2\varkappa$. From $|w|=\Lambda\le\frac12|w_0|$ we get $|w-w_0|\ge\frac12|w_0|$, hence $\frac12|w_0|\Lambda\le|w^2-w_0w|$.
In \eqref{eq:Q1} we have $\frac34\Lambda^2\le\frac38|w_0|\Lambda$, and, by (T4), \eqref{eq:a5-T1} and \eqref{eq:a5-T2},
\[
 \frac{2\mathfrak E_0}{u}\le\frac{2\mathfrak E_0\,\varkappa}{\mathfrak E_3}\le\frac{2\varkappa}{K^{4\fa+8}}\le\frac{\varkappa}{16\cdot1950}\le\frac{|w_0|}{16}.
\]
Therefore $\big(\frac12-\frac38-\frac1{16}\big)|w_0|\Lambda\le\frac{4K\mathfrak E\varkappa}{u}$, i.e., using $|w_0|\ge\varkappa/1950$ and $K\ge124800$,
\[
 \Lambda\le\frac{64\,K\mathfrak E\,\varkappa}{|w_0|\,u}\le\frac{64\cdot1950\,K^{b+1}}{u}=\frac{124800\,K^{b+1}}{u}\le\frac{K^{b+2}}{u}.
\]
Moreover $\frac{K^{b+2}}{u}\le\frac{\varkappa}{7800}\le\frac{|w_0|}4$, because $\varkappa u\ge\mathfrak E_3=K^{b+3}\ge7800\,K^{b+2}$. This proves \eqref{eq:a5-7d}.

\smallskip\noindent\emph{Continuity.} We prove by induction on $k\in\{0,\dots,k_*\}$ that \eqref{eq:a5-7d} holds at $z_k$. By the dichotomy it suffices to show $\Lambda(z_k)\le\frac12|w_0(z_k)|$ for each $k$.

$k=0$: here $\eta_0=10$. Our point lies in $\cG(z_0)$, so by Lemma~\ref{lem:initial}, (S2) and (S4), $\Lambda(z_0)\le\Ld(z_0)\le4\fL N^{-1/2}\le4N^{-15/32}$, while $\frac12|w_0(z_0)|\ge\frac{\varkappa(z_0)}{3900}\ge\frac{\sqrt{10}}{3900}$; so $\Lambda(z_0)\le\frac12|w_0(z_0)|$ for $N\ge N_0$.

$k\to k+1$ ($k<k_*$): the points $z_k,z_{k+1}\in\bD^{\rm near}$ have $|z_k-z_{k+1}|=N^{-4}$ and imaginary parts $\ge N^{-1}$, so by \eqref{eq:lipnear} (which is a consequence of Lemma~\ref{lem:basic}(c) and Lemma~\ref{lem:msc}(f) and holds for every Hermitian matrix)
\[
 |\Lambda(z_k)-\Lambda(z_{k+1})|\le2N^{-2},\qquad|w_0(z_k)-w_0(z_{k+1})|\le N^{-2}.
\]
By the induction hypothesis $\Lambda(z_k)\le\frac14|w_0(z_k)|$. Hence
\[
 \Lambda(z_{k+1})\le\tfrac14|w_0(z_k)|+2N^{-2}\le\tfrac14|w_0(z_{k+1})|+3N^{-2}\le\tfrac12|w_0(z_{k+1})|,
\]
where the last step uses $|w_0(z_{k+1})|\ge\frac{\varkappa(z_{k+1})}{1950}\ge\frac{N^{-1/2}}{1950}$ (T4) and $N^{3/2}\ge23400$. This completes the induction; for $k=k_*$ we obtain (a).

\emph{Proof of \textup{(b)}.} Fix a point of $\cG^2(z)\subset\cG^1(z)$; all quantities are evaluated at $z$, $u=N\eta$. If $\varkappa u\ge\mathfrak E_3$, then (b) follows from (a). Let $\varkappa u<\mathfrak E_3$. From $\Lambda^2\le|w_0|\Lambda+|w^2-w_0w|$ and \eqref{eq:Q1},
\[
 \tfrac14\Lambda^2\le\Big(|w_0|+\frac{2\mathfrak E_0}{u}\Big)\Lambda+\frac{2K\mathfrak E\beta}{u},
\]
so by the last statement of Lemma~\ref{lem:scalar}, $\Lambda\le4|w_0|+\frac{8\mathfrak E_0}{u}+\big(\frac{8K\mathfrak E\beta}{u}\big)^{1/2}$. Now $|w_0|\le4\varkappa<4\mathfrak E_3/u$ and $\beta<(\mathfrak E_3+\mathfrak E)/u\le2\mathfrak E_3/u$, and $K\mathfrak E\mathfrak E_3=K^{2b+4}$. Hence
\[
 \Lambda\le\frac{16\,\mathfrak E_3+8\,\mathfrak E_0+4\,(K\mathfrak E\mathfrak E_3)^{1/2}}{u}\le\frac{(16+8+4)K^{b+3}}{u}\le\frac{K^{b+4}}{u}.
\]

\emph{Proof of \textup{(c)}.} Again at a point of $\cG^2(z)\subset\cG(z)\cap\{\Ld+\Lo\le\lambda_{\rm st}\}$, with $u=N\eta$. By \eqref{eq:ap2} and \eqref{eq:ap4} in their explicit form (Lemma~\ref{lem:apriori}), and by $144C_s\Lambda\le144C_s\lambda_{\rm st}=\frac14$ (Lemma~\ref{lem:stabnear}: $\Lambda\le\Ld\le\lambda_{\rm st}$),
\[
 \Lo\le35\,\fL\Psi,\qquad\Ld\le\Lambda+D\le\tfrac54\Lambda+1036\,C_s\fL\Psi,\qquad\text{so}\qquad\Ld+\Lo\le\tfrac54\Lambda+1071\,C_s\fL\Psi .
\]
By \eqref{eq:Psi} and (b),
\[
 \Psi\le\sqrt{\frac{\im\msc}{u}}+\sqrt{\frac{\Lambda}{u}}+\frac1u\le\sqrt{\frac{\im\msc}{u}}+\frac{K^{(b+4)/2}+1}{u}\le\sqrt{\frac{\im\msc}{u}}+\frac{2K^{6\fa+8}}{u}.
\]
By (S2) and \eqref{eq:a5-T2} ($K\ge K_0\ge2142\,C_s$): $1071C_s\fL\le K^{4\fa+4}\le K^{b+5}$ and $2142\,C_s\fL K^{6\fa+8}\le K^{10\fa+12}\le K^{b}$. Together with (b),
\[
 \Ld+\Lo\le K^{b+5}\sqrt{\frac{\im\msc}{u}}+\frac{\frac54K^{b+4}+K^{b}}{u}\le K^{b+5}\Big(\sqrt{\frac{\im\msc}{u}}+\frac1u\Big).\qedhere
\]
\end{proof}
\noindent\emph{Remark on the two regimes.} The set of type-(a) points is much larger than the set $\{\varkappa\ge\varkappa_w\}$ on which the weak law selects the small branch (Theorem~\ref{thm:weak}(a)): at $E=\pm2$ the former is $\eta\ge K^{2(b+3)/3}N^{-2/3}$, the latter
$\eta\ge(U_w\fL^{2/3})^{6/5}N^{-2/5}$. This is why the continuity argument has to be repeated here with the stronger input $\cG^1$; it is again deterministic on the intersection of the events along the column.

\subsection{Second pass: the improved bound away from the spectrum}
\begin{proposition}[Improved averaged law, near regime]\label{prop:S3}
Let $z\in\mathbb L\cap\bD^{\rm near}$ with $\varkappa N\eta\ge\mathfrak E_3$. There is a closed event $\cG^3(z)\subset\cG^2(z)$ with $\PP(\cG^2(z)\setminus\cG^3(z))\le e^{-p_2}$ on which
\begin{equation}\label{eq:improved}
 \Lambda(z)\le K^{b+2}\Big(\frac{\im\msc(z)}{\varkappa N\eta}+\frac1{\varkappa(N\eta)^2}\Big).
\end{equation}
Namely, $\cG^3(z)$ is the set of points of $\cG^2(z)$ at which \eqref{eq:improved} holds.
\end{proposition}
\begin{proof}
All quantities are evaluated at $z$; $u:=N\eta$, $p:=p_2$, and $\cG^2:=\cG^2(z)$.

\emph{Step 1: a pointwise bound on $\cG^2$, on the small branch.} At a point of $\cG^2$, Proposition~\ref{prop:S2}(a) gives $\Lambda\le\frac14|w_0|$. As in the dichotomy step of that proof, $|w-w_0|\ge\frac34|w_0|\ge\frac12|w_0|$, so by \eqref{eq:ap3}
(which holds on $\cG^2\subset\cG(z)\cap\{\Ld+\Lo\le\lambda_{\rm st}\}$)
\[
 \tfrac12|w_0|\Lambda\le|w^2-w_0w|\le|\avg X|+\mathfrak E_0\Big(\frac{\im\msc}{u}+\frac1{u^2}\Big)+\frac{\mathfrak E_0}{u}\Lambda+\tfrac34\Lambda^2 .
\]
Here $\frac34\Lambda^2\le\frac3{16}|w_0|\Lambda$ and $\frac{\mathfrak E_0}{u}\le\frac1{32}|w_0|$ (shown in the proof of Proposition~\ref{prop:S2}(a)). Hence $\frac9{32}|w_0|\Lambda\le|\avg X|+\mathfrak E_0\big(\frac{\im\msc}{u}+\frac1{u^2}\big)$, and with $|w_0|\ge\varkappa/1950$ (T4) and $\frac{32\cdot1950}{9}\le7000$:
\begin{equation}\tag{7.e}\label{eq:a5-7e}
 \Lambda\le\frac{7000}{\varkappa}\Big(|\avg X|+\mathfrak E_0\frac{\im\msc}{u}+\frac{\mathfrak E_0}{u^2}\Big)\qquad\text{on }\cG^2 .
\end{equation}

\emph{Step 2: moments.} Let $Y_2:=\norm{\Lambda\bone_{\cG^2}}_p\le2N$. By (T4), Proposition~\ref{prop:S2} and (T6),
\[
 \norm\Lambda_p\le Y_2+2N\,\PP\big((\cG^2)^c\big)^{1/p_2}\le Y_2+2N\big(N^{10}e^{-p_1}\big)^{1/p_2}\le Y_2+N^{-8}.
\]
As in Step 2 of the proof of Proposition~\ref{prop:S1} (now with $p=p_2$ in \eqref{eq:Phistar}), $\Phi_*^2\le\frac{151\im\msc}{u}+\frac\Lambda u+\frac{p_2}{u^2}$, and Corollary~\ref{cor:FA} with $p=p_2$ gives, using $\mathfrak F\le u$,
\[
 \norm{\avg X}_p\le\frac{\mathfrak F}{u}\Big(151\,\im\msc+Y_2+\frac{p_2}{u}\Big)+2N^{-8}.
\]
Multiplying \eqref{eq:a5-7e} by $\bone_{\cG^2}$ and taking $L^p$-norms,
\[
 Y_2\le\frac{7000\,\mathfrak F}{\varkappa u}\,Y_2+\frac{7000}{\varkappa}\Big[(151\mathfrak F+\mathfrak E_0)\frac{\im\msc}{u}+\frac{\mathfrak Fp_2+\mathfrak E_0}{u^2}+2N^{-8}\Big].
\]
Since $\varkappa u\ge\mathfrak E_3$, the first coefficient is at most $7000\,\mathfrak F/\mathfrak E_3\le7000\,K^{-4}\le\frac12$ \eqref{eq:a5-T1}, \eqref{eq:a5-T2}; as $Y_2<\infty$, the term is absorbed.
Further $2N^{-8}\le\frac1{150N}\le\frac{\im\msc}{u}$ by (T4) (and $N^7\ge300$); $151\mathfrak F+\mathfrak E_0+1\le153\,K^{12\fa+11}\le K^b$; and $\mathfrak Fp_2+\mathfrak E_0\le K^{12\fa+11}(K+3)\le2K^b$. Therefore
\[
 Y_2\le B:=\frac{28000\,K^b}{\varkappa}\Big(\frac{\im\msc}{u}+\frac1{u^2}\Big).
\]

\emph{Step 3: Markov's inequality.} $B>0$, and $eB\le K^{b+2}\big(\frac{\im\msc}{\varkappa u}+\frac1{\varkappa u^2}\big)$ because $28000\,e\le K^2$. So $\cG^2\setminus\cG^3(z)\subset\{\Lambda\bone_{\cG^2}>eB\}$ and, by \eqref{eq:a5-M},
$\PP(\cG^2\setminus\cG^3(z))\le(Y_2/(eB))^{p_2}\le e^{-p_2}$. The set $\cG^3(z)$ is closed because $\cG^2$ is closed and $\Lambda(z)$ is continuous on $\Omega$.
\end{proof}
\noindent In contrast to the first pass, the factor in front of $Y_2$ is small because of the hypothesis $\varkappa N\eta\ge\mathfrak E_3$ (and not because of a square root); the gain is the factor $1/\varkappa$ times $\im\msc$, which is what excludes eigenvalues outside the spectrum in Section~\ref{sec:norm}.

\begin{proposition}[Averaged law, far regime]\label{prop:S4}
Let $z\in\mathbb L\cap\bD^{\rm far}$. There is a closed event $\cG^{\rm f}(z)\subset\cG_{\rm col}(z)$ with $\PP(\cG_{\rm col}(z)\setminus\cG^{\rm f}(z))\le e^{-p_2}$ on which
\[
 \Lambda(z)\le K^{b+2}\Big(\frac{\im\msc(z)}{N\eta}+\frac1{(N\eta)^2}+\frac1N\Big).
\]
Namely, $\cG^{\rm f}(z)$ is the set of points of $\cG_{\rm col}(z)$ at which this inequality holds.
\end{proposition}
\begin{proof}
All quantities are evaluated at $z$; $u:=N\eta\ge\frac12N^{1/2}\ge1$, $p:=p_2$, $\cG_{\rm c}:=\cG_{\rm col}(z)\subset\cG(z)$. Recall $4\le|E|\le\Emax+1$ and $\eta\le10$, so $|z|\le\Emax+11$.

\emph{Step 1: a pointwise bound on $\cG_{\rm c}$.} Fix a point of $\cG_{\rm c}$. By Theorem~\ref{thm:weak}(b), $\max_i|G_{ii}|\le\frac12$; hence $\im m\le\frac12$ and, as in the proof of Lemma~\ref{lem:farstep} ($u^{-1}\le u^{-1/2}$, $N^{-1/2}=\eta^{1/2}u^{-1/2}\le\sqrt{10}\,u^{-1/2}$),
\[
 \Psi_m=\sqrt{\frac{\im m}{u}}+\frac1u+\frac1{\sqrt N}\le6u^{-1/2}\le6\sqrt2\,N^{-1/4},\qquad\fL\Psi_m\le9N^{\frac1{32}-\frac14}\le\tfrac18\quad(N\ge N_0),
\]
by (S2), (S4). By Lemma~\ref{lem:levels}(S5), $|X_i|\le\fL\Psi_m$ and $|A_i|\le\fL|G_{ii}|\Psi_m^2\le\frac12\fL\Psi_m^2$ for all $i$. Hence, with $\Upsilon_i=A_i+X_i$ (Definition~\ref{def:XA}),
\[
 \Upsilon_{\max}\le\tfrac32\fL\Psi_m\le\tfrac3{16},\qquad|\avg\Upsilon|\le|\avg X|+\tfrac12\fL\Psi_m^2 .
\]
Since $|E|\ge4$ and $\max_i|G_{ii}|\le1$, Lemma~\ref{lem:stabfar} applies: $\Ld\le\frac12\Upsilon_{\max}\le\frac34\fL\Psi_m$ and
\[
 \Lambda\le\tfrac18|\avg\Upsilon|+\tfrac98\Ld^2\big(|z|+1+\Upsilon_{\max}\big)\le\tfrac18|\avg X|+\Big[\tfrac1{16}\fL+\tfrac{81}{128}\big(\Emax+12+\tfrac3{16}\big)\fL^2\Big]\Psi_m^2\le\tfrac18|\avg X|+\tfrac13\mathfrak E_0'\,\Psi_m^2,
\]
where $\mathfrak E_0':=3(\Emax+14)\fL^2$. By $(\alpha+\beta+\gamma)^2\le3(\alpha^2+\beta^2+\gamma^2)$ and $\im m\le\im\msc+\Lambda$: $\Psi_m^2\le3\big(\frac{\im\msc+\Lambda}{u}+\frac1{u^2}+\frac1N\big)$. By (S2) and (T5),
$\mathfrak E_0'\le45\,\Emax\fL^2\le45K^{10\fa+7}\le K^{10\fa+8}\le N^{1/16}$, so $\mathfrak E_0'/u\le2N^{\frac1{16}-\frac12}\le\frac12$. Absorbing the term $\frac{\mathfrak E_0'}{u}\Lambda$ we get
\begin{equation}\tag{7.f}\label{eq:a5-7f}
 \Lambda\le\tfrac14|\avg X|+2\mathfrak E_0'\Big(\frac{\im\msc}{u}+\frac1{u^2}+\frac1N\Big)\qquad\text{on }\cG_{\rm c}.
\end{equation}
(No square root and no branch: in the far regime the equation for $\Lambda$ is linear up to a small quadratic term.)

\emph{Step 2: moments.} Let $Y_{\rm f}:=\norm{\Lambda\bone_{\cG_{\rm c}}}_p$. Since $\Lambda\le\eta^{-1}+1\le2N^{1/2}+1\le2N$, $Y_{\rm f}\le2N$ and, by (T6), $\norm\Lambda_p\le Y_{\rm f}+N^{-8}$. From \eqref{eq:Phistar} with $p=p_2$ and $\im m\le\im\msc+\Lambda$,
$\Phi_*^2\le\frac1N+\frac{\im\msc}{u}+\frac{\Lambda}{u}+\frac{p_2}{u^2}$; by (T5), $\mathfrak F/u\le2N^{\frac1{16}-\frac12}\le1$. So Corollary~\ref{cor:FA} gives
\[
 \norm{\avg X}_p\le\mathfrak F\Big(\frac1N+\frac{\im\msc}{u}+\frac{p_2}{u^2}\Big)+\frac{\mathfrak F}{u}Y_{\rm f}+2N^{-8}.
\]
Multiply \eqref{eq:a5-7f} by $\bone_{\cG_{\rm c}}$ and take $L^p$-norms; since $\frac{\mathfrak F}{4u}\le\frac12$ and $Y_{\rm f}<\infty$, the term with $Y_{\rm f}$ is absorbed, and with $N^{-8}\le\frac1N$:
\[
 Y_{\rm f}\le\big(\tfrac12\mathfrak F+4\mathfrak E_0'+1\big)\Big(\frac{\im\msc}{u}+\frac1N\Big)+\frac{\frac12\mathfrak Fp_2+4\mathfrak E_0'}{u^2}\le B_{\rm f}:=K^{b}\Big(\frac{\im\msc}{u}+\frac1{u^2}+\frac1N\Big),
\]
because $\mathfrak E_0'\le K^{10\fa+8}\le K^{12\fa+11}$, $\frac12\mathfrak F+4\mathfrak E_0'+1\le6K^{12\fa+11}\le K^b$ and $\frac12\mathfrak Fp_2+4\mathfrak E_0'\le K^{12\fa+11}\big(\frac{K+2}2+4\big)\le K^b$.

\emph{Step 3: Markov's inequality.} As in the previous proof: $B_{\rm f}>0$, $eB_{\rm f}\le K^{b+2}\big(\frac{\im\msc}{u}+\frac1{u^2}+\frac1N\big)$, and $\PP(\cG_{\rm c}\setminus\cG^{\rm f}(z))\le\PP(\Lambda\bone_{\cG_{\rm c}}>eB_{\rm f})\le e^{-p_2}$; $\cG^{\rm f}(z)$ is closed.
\end{proof}

\subsection{From the lattice to the domain; conclusion}
\begin{lemma}[Lipschitz extension]\label{lem:lip}
Let $N\ge13$ and $K\ge1$. For every $z\in\bD_K$ there is $z'\in\mathbb L\cap\bD^{\rm near}$ with $|z-z'|\le2N^{-4}$, $\im z'\ge\im z-N^{-4}$, and then, for every Hermitian $N\times N$ matrix $H$:
$|\Lambda(z)-\Lambda(z')|+|\Ld(z)-\Ld(z')|+|\Lo(z)-\Lo(z')|\le8N^{-2}$, $|\im\msc(z)-\im\msc(z')|\le4N^{-3}$, $|\frac1{N\im z}-\frac1{N\im z'}|\le N^{-3}$.
More precisely one can take $z'=E'+\ii\eta'$ with
\[
 E':=\operatorname{sgn}(E)\,N^{-4}\lfloor|E|N^4\rfloor,\qquad\eta':=10-N^{-4}\lfloor(10-\eta)N^4\rfloor ;
\]
then $z'\in\bD_K$, $|E'|\le|E|$, $|E-E'|<N^{-4}$, $\eta\le\eta'<\eta+N^{-4}$, and each of $|\Lambda(z)-\Lambda(z')|$, $|\Ld(z)-\Ld(z')|$, $|\Lo(z)-\Lo(z')|$ is at most $2N^{-2}$.
\end{lemma}
\begin{proof}
Let $z=E+\ii\eta\in\bD_K$, i.e.\ $|E|\le5$ and $K^{C_1}N^{-1}\le\eta\le10$ \eqref{eq:DK}, and define $E',\eta'$ as stated. Then $E'\in N^{-4}\Z$, $|E'|\le|E|\le5$ and $|E-E'|<N^{-4}$.
The integer $k:=\lfloor(10-\eta)N^4\rfloor$ satisfies $k\ge0$ and $(10-\eta)N^4-1<k\le(10-\eta)N^4$, so $\eta'=10-kN^{-4}\in10-N^{-4}\Z_{\ge0}$ and $\eta\le\eta'<\eta+N^{-4}$; in particular $K^{C_1}N^{-1}\le\eta'\le10$.
Hence $z'\in\bD_K\subset\bD^{\rm near}$ (as $K^{C_1-1}\le K^{C_1}$) and $z'$ belongs to the lattice of Definition~\ref{def:domains}, so $z'\in\mathbb L\cap\bD^{\rm near}$. Moreover $|z-z'|<\sqrt2\,N^{-4}$ and $\im z'\ge\im z$.

Both imaginary parts are at least $\eta_0:=\eta\ge K^{C_1}N^{-1}\ge N^{-1}$. By Lemma~\ref{lem:basic}(c), for $F\in\{\Lambda,\Ld,\Lo\}$,
\[
 |F(z)-F(z')|\le\big(\eta_0^{-2}+\eta_0^{-1}+1\big)|z-z'|\le(N^2+N+1)\sqrt2\,N^{-4}\le2N^{-2}\qquad(N\ge13);
\]
the sum of the three differences is therefore at most $6N^{-2}\le8N^{-2}$. By Lemma~\ref{lem:msc}(f), $|\im\msc(z)-\im\msc(z')|\le|\msc(z)-\msc(z')|\le(1+\eta_0^{-1})|z-z'|\le(1+N)\sqrt2\,N^{-4}\le4N^{-3}$.
Finally $\big|\frac1{N\eta}-\frac1{N\eta'}\big|=\frac{\eta'-\eta}{N\eta\eta'}\le\frac{N^{-4}}{N\cdot N^{-2}}=N^{-3}$.
\end{proof}

\medskip\noindent\textbf{Proof of Theorem~\ref{thm:K}\,(i),(ii).} Let
\[
 \Xi_K^{(1)}:=\bigcap_{z\in\mathbb L\cap\bD^{\rm near}}\cG^2(z).
\]
This is a closed subset of $\Omega$, and by Proposition~\ref{prop:S2} and $|\mathbb L|\le N^{10}$ (Definition~\ref{def:domains}), $\PP\big((\Xi_K^{(1)})^c\big)\le N^{20}e^{-p_1}$. Fix a point of $\Xi_K^{(1)}$. Then Proposition~\ref{prop:S2}(b),(c) hold at all points of $\mathbb L\cap\bD^{\rm near}$.
Let $z=E+\ii\eta\in\bD_K$ and let $z'=E'+\ii\eta'$ be the lattice point of Lemma~\ref{lem:lip} (recall $N\ge13$ by (S1)); thus $\eta\le\eta'$, and $N\eta\le10N$.

\emph{\textup{(i)}.} By Proposition~\ref{prop:S2}(b) at $z'$ and Lemma~\ref{lem:lip},
\[
 \Lambda(z)\le\Lambda(z')+2N^{-2}\le\frac{K^{b+4}}{N\eta'}+2N^{-2}\le\frac{K^{b+4}}{N\eta}+\frac{20}{N}\cdot\frac{1}{N\eta}\le\frac{2K^{b+4}}{N\eta}\le\frac{K^{b+5}}{N\eta},
\]
where we used $2N^{-2}=\frac{2\eta}{N}\cdot\frac1{N\eta}\le\frac{20}N\cdot\frac1{N\eta}$.

\emph{\textup{(ii)}.} By Proposition~\ref{prop:S2}(c) at $z'$, Lemma~\ref{lem:lip} and $\eta'\ge\eta$,
\[
 \Ld(z)+\Lo(z)\le K^{b+5}\Big(\sqrt{\frac{\im\msc(z)+4N^{-3}}{N\eta}}+\frac1{N\eta}\Big)+4N^{-2}.
\]
Since $1\le N\eta\le10N$, we have $\sqrt{4N^{-3}/(N\eta)}\le2N^{-3/2}\le\frac1{10N}\le\frac1{N\eta}$ (as $N\ge400$), and $4N^{-2}\le\frac{40}{N}\cdot\frac1{N\eta}\le\frac{K^{b+5}}{N\eta}$. Hence, by $\sqrt{\alpha+\beta}\le\sqrt\alpha+\sqrt\beta$,
\[
 \Ld(z)+\Lo(z)\le K^{b+5}\sqrt{\frac{\im\msc(z)}{N\eta}}+\frac{3K^{b+5}}{N\eta}\le K^{b+6}\Big(\sqrt{\frac{\im\msc(z)}{N\eta}}+\frac1{N\eta}\Big).
\]
Since $C_2\ge b+7$ by \eqref{eq:a5-T3} and $K\ge1$, parts (i) and (ii) of Theorem~\ref{thm:K} hold at every point of $\Xi_K^{(1)}$, for all $z\in\bD_K$ simultaneously (with one spare power of $K$). \qed
\section{Outside the spectrum: the norm bound \eqref{eq:443}}\label{sec:norm}
The standing assumptions and conditions of Section~\ref{sec:strong} remain in force. Recall that $\norm H=\max_j|\lambda_j|$ (Lemma~\ref{lem:weyl}(a)) and that, by Lemma~\ref{lem:basic}(e), for every $E\in\R$ and $\eta>0$:
\begin{equation}\tag{8.a}\label{eq:a5-8a}
 N\eta\,\im m(E+\ii\eta)<\tfrac12\quad\Longrightarrow\quad H\text{ has no eigenvalue in }[E-\eta,E+\eta].
\end{equation}
We also use $\im m\le\im\msc+\Lambda$. Both edges are treated directly (the lattice $\mathbb L$, the domains and all statements of Sections~\ref{sec:weak}--\ref{sec:strong} are symmetric under $E\mapsto-E$); no reflection $H\mapsto-H$ is needed.

\subsection{Exclusion of eigenvalues in three ranges}
\begin{proposition}[No eigenvalues outside $2+K^{C_3}N^{-2/3}$]\label{prop:norm}
There is a closed event $\Xi_K^{(2)}$ with $\PP((\Xi_K^{(2)})^c)\le N^{21}e^{-p_1}+N^{6}e^{-p_2}$ on which $\max_j|\lambda_j|\le2+K^{C_3}N^{-2/3}$, where $C_3=2b+12=24\fa+36$.
\end{proposition}
\begin{proof}
We define two heights of the lattice,
\[
 \eta_f:=10-N^{-4}\big\lfloor(10-\tfrac12N^{-1/2})N^4\big\rfloor,\qquad\eta_e:=10-N^{-4}\big\lfloor(10-N^{-2/3})N^4\big\rfloor ,
\]
so that $\eta_f,\eta_e\in10-N^{-4}\Z_{\ge0}$ and (as in the proof of Lemma~\ref{lem:lip})
\[
 \tfrac12N^{-1/2}\le\eta_f<\tfrac12N^{-1/2}+N^{-4}\le N^{-1/2},\qquad N^{-2/3}\le\eta_e<N^{-2/3}+N^{-4}\le2N^{-2/3},
\]
and two finite sets of energies,
\[
 \mathbb E_f:=\big\{E\in N^{-4}\Z:\ 4.5\le|E|\le\Emax+1\big\},\qquad\mathbb E_e:=\big\{E\in N^{-4}\Z:\ 2+\tfrac12K^{C_3}N^{-2/3}\le|E|\le5\big\}.
\]
As in the justification of Definition~\ref{def:domains}, $|\mathbb E_f|\le5N^{4+1/32}\le N^5$ and $|\mathbb E_e|\le10N^4+1\le N^5$.

\emph{Step 1 (crude bound: no eigenvalue with $|\lambda|>\Emax$).} Let $\mathcal N_0:=\{\norm H\le\Emax\}$, a closed set (Lemma~\ref{lem:weyl}(a),(d)). By (S1) the hypotheses of Lemma~\ref{lem:norm0} hold for $\kappa:=K_w=K^2$, and $25K_w^{\fa}\sqrt{C_0\log N}=\Emax$ (Definition~\ref{def:levels}). Hence
\[
 \PP(\mathcal N_0^c)\le N^2\fa\,e^{-K^2}+2e^{-N\log N}\le(N^2\fa+1)\,e^{-K^2},
\]
because $K^2\le K^{6\fa+6}\le N^{1/32}$ by (S4), so that $K^2+\log2\le N\log N$.

\emph{Step 2 (far regime: no eigenvalue with $4.5\le|\lambda|\le\Emax$).} Let $E\in\mathbb E_f$ and $z:=E+\ii\eta_f$; then $z\in\mathbb L\cap\bD^{\rm far}$. Fix a point of $\cG^{\rm f}(z)$. By Lemma~\ref{lem:msc}(g), $\im\msc(z)\le\eta_f/8$, and by Proposition~\ref{prop:S4} and $N\eta_f\ge\frac12N^{1/2}$,
\[
 \im m(z)\le\frac{\eta_f}8+K^{b+2}\Big(\frac1{8N}+\frac1{(N\eta_f)^2}+\frac1N\Big)\le\frac{\eta_f}8+\frac{41}{8}\cdot\frac{K^{b+2}}{N}.
\]
Since $\eta_f\le N^{-1/2}$ and $K^{b+2}\le N^{3/32}$ (T5),
\[
 N\eta_f\,\im m(z)\le\frac{N\eta_f^2}8+\frac{41}{8}K^{b+2}\eta_f\le\frac18+6N^{\frac3{32}-\frac12}\le\frac38<\frac12\qquad(N\ge N_0).
\]
By \eqref{eq:a5-8a}, $H$ has no eigenvalue in $[E-\eta_f,E+\eta_f]$. Now let $x\in\R$ with $4.5\le|x|\le\Emax$ and put $E:=\operatorname{sgn}(x)N^{-4}\lceil|x|N^4\rceil$. Then $|x|\le|E|<|x|+N^{-4}\le\Emax+1$, so $E\in\mathbb E_f$, and $|x-E|<N^{-4}\le\eta_f$.
Consequently, on $\bigcap_{E\in\mathbb E_f}\cG^{\rm f}(E+\ii\eta_f)$ there is no eigenvalue $\lambda$ with $4.5\le|\lambda|\le\Emax$.

\emph{Step 3 (near the edges: no eigenvalue with $2+K^{C_3}N^{-2/3}\le|\lambda|\le5$).} Let $E\in\mathbb E_e$ and $z:=E+\ii\eta_e$. By (T5), $K^{C_1-1}N^{-1}\le N^{-2/3}\le\eta_e$, so $z\in\mathbb L\cap\bD^{\rm near}$. Since $|E|>2$, $\kappa=|E|-2$, and by \eqref{eq:a5-T3},
\begin{equation}\tag{8.b}\label{eq:a5-8b}
 \kappa\ge\tfrac12K^{C_3}N^{-2/3}\ge\tfrac12K^{2b+12}N^{-2/3},\qquad\sqrt\kappa\ge\tfrac1{\sqrt2}K^{b+6}N^{-1/3},\qquad\eta_e\le2N^{-2/3}\le\kappa .
\end{equation}
Hence $\varkappa N\eta_e\ge\sqrt\kappa\,N^{1/3}\ge\frac1{\sqrt2}K^{b+6}\ge K^{b+3}=\mathfrak E_3$: the point $z$ is of type (a), and Proposition~\ref{prop:S3} applies. Fix a point of $\cG^3(z)$.
Since $|E|>2$ and $\kappa\ge\eta_e$, \eqref{eq:imm} gives $\im\msc(z)\le4\eta_e/\varkappa$, and $\varkappa\ge\sqrt\kappa$. By \eqref{eq:improved},
\[
 \Lambda(z)\le K^{b+2}\Big(\frac{4}{N\varkappa^2}+\frac1{\varkappa(N\eta_e)^2}\Big)\le K^{b+2}\Big(\frac{4}{N\kappa}+\frac{1}{\sqrt\kappa\,(N\eta_e)^2}\Big).
\]
We multiply by $N\eta_e$ and use \eqref{eq:a5-8b} and $N^{1/3}\le N\eta_e\le2N^{1/3}$:
\begin{align*}
 N\eta_e\,\im\msc(z)&\le\frac{4N\eta_e^2}{\sqrt\kappa}\le\frac{16\sqrt2\,N^{-1/3}}{K^{b+6}N^{-1/3}}=\frac{16\sqrt2}{K^{b+6}}\le\frac18,\\
 N\eta_e\cdot\frac{4K^{b+2}}{N\kappa}&=\frac{4K^{b+2}\eta_e}{\kappa}\le\frac{16\,K^{b+2}}{K^{2b+12}}\le\frac1{16},\\
 N\eta_e\cdot\frac{K^{b+2}}{\sqrt\kappa\,(N\eta_e)^2}&=\frac{K^{b+2}}{\sqrt\kappa\,N\eta_e}\le\frac{\sqrt2\,K^{b+2}}{K^{b+6}}\le\frac1{16}.
\end{align*}
(These are the three places where the exponent of the norm bound is fixed: they need $\frac12C_3>b+2$, and $C_3=2b+12$ leaves room.) So $N\eta_e\,\im m(z)\le\frac18+\frac1{16}+\frac1{16}=\frac14<\frac12$, and by \eqref{eq:a5-8a} there is no eigenvalue in $[E-\eta_e,E+\eta_e]$.
Now let $x\in\R$ with $2+K^{C_3}N^{-2/3}\le|x|\le5$ and put $E:=\operatorname{sgn}(x)N^{-4}\lfloor|x|N^4\rfloor$. Then $|x|-N^{-4}<|E|\le|x|\le5$ and $|E|\ge2+K^{C_3}N^{-2/3}-N^{-4}\ge2+\frac12K^{C_3}N^{-2/3}$ (as $N^{-4}\le\frac12N^{-2/3}$), so
$E\in\mathbb E_e$, and $|x-E|<N^{-4}\le\eta_e$. Consequently, on $\bigcap_{E\in\mathbb E_e}\cG^3(E+\ii\eta_e)$ there is no eigenvalue $\lambda$ with $2+K^{C_3}N^{-2/3}\le|\lambda|\le5$.

\emph{Conclusion.} Let
\[
 \Xi_K^{(2)}:=\mathcal N_0\ \cap\bigcap_{E\in\mathbb E_f}\cG^{\rm f}(E+\ii\eta_f)\ \cap\bigcap_{E\in\mathbb E_e}\cG^3(E+\ii\eta_e).
\]
It is closed (Propositions~\ref{prop:S3}, \ref{prop:S4}). At a point of $\Xi_K^{(2)}$ every eigenvalue satisfies $|\lambda|\le\Emax$ (Step 1), not $4.5\le|\lambda|\le\Emax$ (Step 2), and not $2+K^{C_3}N^{-2/3}\le|\lambda|\le5$ (Step 3); the three ranges cover $[2+K^{C_3}N^{-2/3},\infty)$, so $|\lambda|<2+K^{C_3}N^{-2/3}$.

\emph{Probability.} For $z\in\mathbb L\cap\bD^{\rm far}$: $\PP(\cG^{\rm f}(z)^c)\le\PP(\cG_{\rm col}(z)^c)+\PP(\cG_{\rm col}(z)\setminus\cG^{\rm f}(z))\le460\,\fa N^8e^{-K^2}+e^{-p_2}$ (Definition~\ref{def:domains}, Proposition~\ref{prop:S4}).
For $z\in\mathbb L\cap\bD^{\rm near}$ of type (a): $\PP(\cG^3(z)^c)\le\PP(\cG^2(z)^c)+\PP(\cG^2(z)\setminus\cG^3(z))\le N^{10}e^{-p_1}+e^{-p_2}$ (Propositions~\ref{prop:S2}, \ref{prop:S3}). With Step~1, the union bound and $K^2\ge p_1$,
\[
 \PP\big((\Xi_K^{(2)})^c\big)\le\big(N^2\fa+1+460\,\fa N^{13}+N^{15}\big)e^{-p_1}+2N^5e^{-p_2}\le N^{21}e^{-p_1}+N^6e^{-p_2}\qquad(N\ge N_0).\qedhere
\]
\end{proof}

\begin{remark}
The original excludes eigenvalues in $[2+N^{-2/3}(\log N)^{9L},5]$ by Lemma~4.5 and cites a moment bound for $\PP(\lambda_N\ge3)$ (Lemma~4.6, from [EYY]). NP has no such external bound; this is why the far regime $4\le|E|\le\Emax+1$ is carried through Sections~\ref{sec:weak}--\ref{sec:strong}.
The cost is small because stability there is a contraction (Lemmas~\ref{lem:stabfar}, \ref{lem:farstep}) and Theorem~\ref{thm:FA} is uniform in $E$. The three ranges overlap ($[2+K^{C_3}N^{-2/3},5]$, $[4.5,\Emax]$, $(\Emax,\infty)$), so no gap can occur at their junctions.
The far regime is used at the single height $\eta_f\approx\frac12N^{-1/2}$, where $\im\msc\le\eta_f/8$ contributes $\frac18$ to $N\eta_f\im m$ and all other terms are $O(K^{b+2}N^{-1/2})$.
\end{remark}

\subsection{The event \texorpdfstring{$\Xi_K$}{Xi\_K}; proof of Theorem \ref{thm:K}\,(i)--(iii)}\label{sec:XiK}
\begin{definition}[The event of Theorem~\ref{thm:K}]\label{def:XiK}
For $N\ge N_0$, $\nu\in\mathcal E_N(\vth,\dpm,C_0)$ and $K$ as in \eqref{eq:Krange}, let
\[
 \Xi_K:=\Xi_K^{(1)}\cap\Xi_K^{(2)},
\]
with $\Xi_K^{(1)}=\bigcap_{z\in\mathbb L\cap\bD^{\rm near}}\cG^2(z)$ from the end of Section~\ref{sec:strong} and $\Xi_K^{(2)}$ from Proposition~\ref{prop:norm}. Explicitly, $\Xi_K$ is the set of $x\in\Omega$ such that
\begin{enumerate}[label=\textup{(\roman*)},leftmargin=2.2em,itemsep=0pt]
\item $x\in\cG(z)$ for all $z\in\mathbb L\cap\bD^{\rm near}$ and for all $z\in{\rm col}(E+\ii\eta_f)$, $E\in\mathbb E_f$ \quad(large deviation events, level $K_w=K^2$);
\item $|\avg X(z)|\le\frac{K\mathfrak E}{N\eta}\big(\varkappa+\frac{\mathfrak E}{N\eta}\big)$ for all $z=E+\ii\eta\in\mathbb L\cap\bD^{\rm near}$ \quad(first pass, level $p_1$);
\item $\Lambda(z)\le K^{b+2}\big(\frac{\im\msc}{\varkappa N\eta}+\frac1{\varkappa(N\eta)^2}\big)$ for $z=E+\ii\eta_e$, $E\in\mathbb E_e$, and $\Lambda(z)\le K^{b+2}\big(\frac{\im\msc}{N\eta}+\frac1{(N\eta)^2}+\frac1N\big)$ for $z=E+\ii\eta_f$, $E\in\mathbb E_f$ \quad(second pass, level $p_2$);
\item $\norm H\le\Emax$.
\end{enumerate}
\end{definition}
\noindent $\Xi_K$ is a closed, hence Borel, subset of $\Omega$: it is defined by finitely many inequalities ``$\le$'' between continuous functions of $x$ (Lemmas~\ref{lem:weyl}(e), \ref{lem:good}). The description (i)--(iv) is the definition unfolded:
$\cG^2(z)=\bigcap_{z'\in{\rm col}(z)}\big(\cG_{\rm col}(z')\cap\{|\avg X(z')|\le\dots\}\big)$, and the columns of the points of $\mathbb L\cap\bD^{\rm near}$ lie in $\mathbb L\cap\bD^{\rm near}$.

\begin{proposition}[Probability of $\Xi_K$]\label{prop:XiK}
For $N\ge N_0$: $\ \PP(\Xi_K^c)\le2N^{21}e^{-p_1}+N^6e^{-p_2}\le e^{-K/2}$.
\end{proposition}
\begin{proof}
By the end of Section~\ref{sec:strong} and Proposition~\ref{prop:norm}, $\PP(\Xi_K^c)\le N^{20}e^{-p_1}+N^{21}e^{-p_1}+N^6e^{-p_2}\le2N^{21}e^{-K^{3/2}}+N^6e^{-K}$, since $p_1\ge K^{3/2}$ and $p_2\ge K$ \eqref{eq:a5-T1}.
Let $\mathsf L:=\log N$ and assume $\mathsf L^2\ge46$ (i.e.\ $N\ge e^{7}$). Since $K\ge\mathsf L^3$ \eqref{eq:Krange} and $K\ge1$: $\ \frac12K\ge\frac12\mathsf L^3\ge23\,\mathsf L\ge22\,\mathsf L+\log4$. Hence
\[
 2N^{21}e^{-K^{3/2}}\le2N^{21}e^{-K}\le\tfrac12e^{-K/2},\qquad N^6e^{-K}\le\tfrac12e^{-K/2},
\]
and the sum is at most $e^{-K/2}$.
\end{proof}

\medskip\noindent\textbf{Proof of Theorem~\ref{thm:K}.} Let the model parameters be normalised as in Convention~\ref{conv:const}, and let $N_0$ be the threshold described in Appendix~\ref{app:ledger} (paragraph ``Thresholds''); in particular $N_0\ge e^{7}$.
Let $N\ge N_0$, $\kk\in\{\R,\C\}$, $\nu\in\mathcal E_N(\vth,\dpm,C_0)$ and $K$ as in \eqref{eq:Krange}. The event $\Xi_K$ of Definition~\ref{def:XiK} is measurable and $\PP(\Xi_K^c)\le e^{-K/2}$ (Proposition~\ref{prop:XiK}).
At every point of $\Xi_K\subset\Xi_K^{(1)}$, parts (i) and (ii) hold for all $z\in\bD_K$ (end of Section~\ref{sec:strong}). At every point of $\Xi_K\subset\Xi_K^{(2)}$, part (iii) holds (Proposition~\ref{prop:norm}; $C_3=2b+12$ is the value in \eqref{eq:ledger0}).
Finally, since $N\ge e^6$ and the relations \eqref{eq:ledgerrel} hold for the values \eqref{eq:ledger0} (Appendix~\ref{app:ledger}), Proposition~\ref{prop:rig} shows that parts (iv) and (v) hold at every point at which (i) and (iii) hold, in particular on $\Xi_K$. \qed

\begin{remark}[What the event contains]\label{rem:XiK}
On $\Xi_K$ more is true than Theorem~\ref{thm:K} states: the bounds of Proposition~\ref{prop:S2} with the exponents $b+4$, $b+5$ at all points of $\mathbb L\cap\bD^{\rm near}$, and the improved bounds \eqref{eq:improved} at the heights $\eta_e$, $\eta_f$ on the sets $\mathbb E_e$, $\mathbb E_f$.
The improved bound \eqref{eq:improved} at \emph{all} lattice points of type (a), and the far-regime bound of Proposition~\ref{prop:S4} at all points of $\mathbb L\cap\bD^{\rm far}$, hold on a smaller closed event $\Xi_K^{(4)}\subset\Xi_K$ whose complement has probability at most $2N^{21}e^{-p_1}+2N^{10}e^{-p_2}\le e^{-K/2}$ as well (see (X4) in Section~\ref{sec:export}).
The three levels are visible in Definition~\ref{def:XiK}: the set where (i) or (iv) fails has probability at most $47\fa N^{14}e^{-K^2}$; the set where (i) holds and (ii) fails has probability at most $N^{10}e^{-p_1}$;
the set where (i), (ii) hold and (iii) fails has probability at most $2N^{5}e^{-p_2}$. The passage from one level to the next costs the factor $40\log N$ of \eqref{eq:levelrel}, which is spent in (T6).
\end{remark}

\section{Rigidity: proof of Theorem \ref{thm:K}\,(iv),(v)}\label{sec:rig}
This section is deterministic. It converts the averaged local law (Theorem~\ref{thm:K}(i)) and the norm bound (Theorem~\ref{thm:K}(iii)) into the bound on the counting function and into rigidity.
The exponents of \eqref{eq:ledger0} enter only through the following relations, which hold for the values in \eqref{eq:ledger0} (recall $\fa\ge1$):
\begin{equation}\label{eq:ledgerrel}
 C_5\ge\max\big(C_1,\,C_2+\tfrac13\big)+1,\qquad C_4\ge\max\big(C_5,\,C_3+\tfrac13C_5\big)+1,\qquad c_0\max(C_1,C_3,C_5)\le\tfrac18 .
\end{equation}
Indeed $C_5=C_1+1$ and $C_2+\frac13<C_1$; $C_3+\frac13C_5+1=24\fa+36+6\fa+7+1=C_4$ and $C_5<C_3$; and $C_1<C_5<C_3$, $c_0C_3=\frac18$.

\begin{lemma}[From counting to rigidity; deterministic]\label{lem:conv}
Let $C_7:=9$. Let $\lambda_1\le\dots\le\lambda_N$ be real, $\fn(E):=\frac1N\#\{j:\lambda_j\le E\}$, $\Delta\ge1$, $0\le\omega\le1$, $N\ge 100\Delta$, and assume
\[
 \sup_{E\in\R}|\fn(E)-\nsc(E)|\le\frac\Delta N,\qquad-2-\omega\le\lambda_1,\quad\lambda_N\le2+\omega .
\]
Then for every $j\in\{1,\dots,N\}$, with $\hat\jmath:=\min(j,N+1-j)$:
\[
 |\lambda_j-\gamma_j|\le C_7\,\Delta\,\hat\jmath^{-1/3}N^{-2/3}+\bone(\hat\jmath\le3\Delta)\big(\omega+C_7\Delta^{2/3}N^{-2/3}\big).
\]
More precisely: if $\hat\jmath>3\Delta$ then $|\lambda_j-\gamma_j|\le3.5\,\Delta\,\hat\jmath^{-1/3}N^{-2/3}$, and if $\hat\jmath\le3\Delta$ then $|\lambda_j-\gamma_j|\le\omega+9\,\Delta^{2/3}N^{-2/3}$.
\end{lemma}
\begin{proof}
\emph{Step 1: for every $j$,}
\begin{equation}\tag{R1}\label{eq:a1-R1}
 \frac{j-\Delta}N\le\nsc(\lambda_j)\le\frac{j-1+\Delta}N .
\end{equation}
At least $j$ of the numbers $\lambda_k$ are $\le\lambda_j$, so $\fn(\lambda_j)\ge j/N$ and $\nsc(\lambda_j)\ge\fn(\lambda_j)-\Delta/N\ge(j-\Delta)/N$.
If $E<\lambda_j$, then $\lambda_k\le E$ implies $k\le j-1$, so $\fn(E)\le(j-1)/N$ and $\nsc(E)\le(j-1+\Delta)/N$; letting $E\uparrow\lambda_j$ and using the continuity of $\nsc$ (Lemma~\ref{lem:nsc}(a)) gives the upper bound.
Recall also $\nsc(\gamma_j)=j/N$ (Lemma~\ref{lem:nsc}(c)); since $\Delta\ge1$, \eqref{eq:a1-R1} holds with $\gamma_j$ in place of $\lambda_j$ as well.

\emph{Step 2: the case $\hat\jmath>3\Delta$.} Let $a:=\min(\lambda_j,\gamma_j)$, $b:=\max(\lambda_j,\gamma_j)$. By Step 1 and the monotonicity of $\nsc$,
\[
 \nsc(a)\ge\frac{j-\Delta}N,\qquad1-\nsc(b)\ge\frac{N+1-j-\Delta}N,\qquad0\le\nsc(b)-\nsc(a)=|\nsc(\lambda_j)-\nsc(\gamma_j)|\le\frac\Delta N ,
\]
where the last bound follows from \eqref{eq:a1-R1} and $\nsc(\gamma_j)=j/N$: $-\Delta\le N\nsc(\lambda_j)-j\le\Delta-1$. Hence $\min(\nsc(a),1-\nsc(b))\ge(\hat\jmath-\Delta)/N\ge\frac{2\hat\jmath}{3N}>0$, and Lemma~\ref{lem:nsc}(f) gives
\[
 |\lambda_j-\gamma_j|=b-a\le3\cdot\frac\Delta N\cdot\Big(\frac{2\hat\jmath}{3N}\Big)^{-1/3}=3\big(\tfrac32\big)^{1/3}\Delta\,\hat\jmath^{-1/3}N^{-2/3}\le3.5\,\Delta\,\hat\jmath^{-1/3}N^{-2/3}.
\]

\emph{Step 3: the case $j\le3\Delta$ (left edge).} Since $N\ge100\Delta$ we have $j\le N/2$. Put $\theta:=(\Delta/N)^{2/3}$.

By Lemma~\ref{lem:nsc}(c), $0<\gamma_j+2\le3.6\,(j/N)^{2/3}\le3.6\cdot3^{2/3}\theta\le7.5\,\theta$.

By \eqref{eq:a1-R1}, $\nsc(\lambda_j)\le(j-1+\Delta)/N\le4\Delta/N\le\frac1{25}<\frac12=\nsc(0)$, so $\lambda_j<0$. If $\lambda_j\le-2$, then $-\omega\le\lambda_1+2\le\lambda_j+2\le0$. If $\lambda_j>-2$, then $u:=\lambda_j+2\in(0,2)$ and Lemma~\ref{lem:nsc}(b) gives
$\frac{\sqrt2}{3\pi}u^{3/2}\le\nsc(\lambda_j)\le4\Delta/N$, i.e.\ $u\le(6\sqrt2\pi)^{2/3}\theta\le9\,\theta$. In both cases $-\omega\le\lambda_j+2\le9\theta$. Therefore
\[
 |\lambda_j-\gamma_j|=|(\lambda_j+2)-(\gamma_j+2)|\le\max\big(\omega+7.5\,\theta,\ 9\,\theta\big)\le\omega+9\,\Delta^{2/3}N^{-2/3}.
\]

\emph{Step 4: the case $N+1-j\le3\Delta$ (right edge).} Now $j\ge N/2$, and by Lemma~\ref{lem:nsc}(c), $0\le2-\gamma_j\le3.6\,((N-j)/N)^{2/3}\le7.5\,\theta$.
By \eqref{eq:a1-R1}, $1-\nsc(\lambda_j)\le(N-j+\Delta)/N\le4\Delta/N<\frac12$, so $\lambda_j>0$. If $\lambda_j\ge2$, then $0\le\lambda_j-2\le\lambda_N-2\le\omega$. If $\lambda_j<2$, then $u:=2-\lambda_j\in(0,2)$ and, by the symmetry $1-\nsc(2-u)=\nsc(-2+u)$ and Lemma~\ref{lem:nsc}(b),
$\frac{\sqrt2}{3\pi}u^{3/2}\le1-\nsc(\lambda_j)\le4\Delta/N$, so $u\le9\theta$. Hence $-\omega\le2-\lambda_j\le9\theta$, and as in Step 3, $|\lambda_j-\gamma_j|\le\omega+9\,\Delta^{2/3}N^{-2/3}$.

Steps 2--4 cover all $j$ and give the two ``more precise'' statements; these imply the displayed bound with $C_7=9$.
\end{proof}
\noindent The lemma treats the two edges directly; no reflection $H\mapsto-H$ and no shift of the index are needed.

\begin{proposition}[Counting function and rigidity from the local law; deterministic]\label{prop:rig}
Let $\vth\in(0,1]$, let $C_1,\dots,C_5,c_0$ be as in \eqref{eq:ledger0}, and let $C_6=10$, $C_7=9$ be the constants of Lemmas~\ref{lem:count} and~\ref{lem:conv}. Let $N\ge e^{6}$ and let $K$ be a real number with $(\log N)^3\le K\le N^{c_0}$.
Let $H$ be a Hermitian $N\times N$ matrix with eigenvalues $\lambda_1\le\dots\le\lambda_N$ such that
\begin{enumerate}[label=\textup{(\alph*)},leftmargin=2.2em,itemsep=1pt]
\item $|m(z)-\msc(z)|\le K^{C_2}/(N\eta)$ for all $z=E+\ii\eta$ with $|E|\le\frac72$ and $K^{C_1}N^{-1}\le\eta\le2$;
\item $\max_j|\lambda_j|\le2+K^{C_3}N^{-2/3}$.
\end{enumerate}
Then
\[
 \sup_{E\in\R}|\fn(E)-\nsc(E)|\le\frac{K^{C_5}}N\qquad\text{and}\qquad|\lambda_j-\gamma_j|\le K^{C_4}\,[\min(j,N+1-j)]^{-1/3}N^{-2/3}\quad\text{for all }j .
\]
In particular, parts \textup{(iv)} and \textup{(v)} of Theorem~\ref{thm:K} hold at every point of any event on which parts \textup{(i)} and \textup{(iii)} hold, provided $N\ge e^6$.
\end{proposition}
\begin{proof}
\emph{Preliminaries.} Since $N\ge e^6$: $\log N\ge6$, $K\ge(\log N)^3\ge216$, and $N^{1/8}\le N^{2/3}$. By the third relation in \eqref{eq:ledgerrel}, for $i\in\{1,3,5\}$,
\begin{equation}\tag{R2}\label{eq:a1-R2}
 K^{C_i}\le N^{c_0C_i}\le N^{1/8} .
\end{equation}
Moreover $\log(2N)\le2\log N\le2K^{1/3}$, and $N^{7/8}\ge e^{5.25}>100$.

\emph{Part \textup{(iv)}.} By (b) and \eqref{eq:a1-R2}, $|\lambda_j|\le2+N^{1/8-2/3}<3$ for all $j$. Hence for $E\ge3$, $\fn(E)=1=\nsc(E)$, and for $E\le-3$, $\fn(E)=0=\nsc(E)$ (Lemma~\ref{lem:nsc}(a)).
Let $E\in(-3,3)$. We apply Lemma~\ref{lem:count} with
\[
 E_1:=-3,\qquad E_2:=E,\qquad\tilde\eta:=K^{C_1}N^{-1},\qquad U:=K^{C_2}.
\]
Its assumptions hold: $E_1<E_2$ and $E_2-E_1\le6\le8$; $U\ge1$; by \eqref{eq:a1-R2}, $0<\tilde\eta\le N^{-7/8}\le\frac12$; and \eqref{eq:HShyp} is hypothesis (a), because $m(z)=\frac1N\tr G(z)=\frac1N\sum_j(\lambda_j-z)^{-1}$ (Lemma~\ref{lem:basic}(d)),
$[E_1-\tilde\eta,E_2+\tilde\eta]\subset[-\frac72,\frac72]$ and $[\tilde\eta,2]=[K^{C_1}N^{-1},2]$. Since $\fn(-3)=\nsc(-3)=0$ and $\log(2/\tilde\eta)\le\log(2N)\le2K^{1/3}$, Lemma~\ref{lem:count} gives
\[
 |\fn(E)-\nsc(E)|\le\frac{C_6}N\Big(K^{C_2}\log(2/\tilde\eta)+K^{C_1}\Big)\le\frac{C_6}N\Big(2K^{C_2+1/3}+K^{C_1}\Big)\le\frac{3C_6}{N}K^{C_5-1}\le\frac{K^{C_5}}N,
\]
where we used the first relation in \eqref{eq:ledgerrel} and $K\ge216\ge3C_6$.

\emph{Part \textup{(v)}.} We apply Lemma~\ref{lem:conv} with $\Delta:=K^{C_5}\ge1$ and $\omega:=K^{C_3}N^{-2/3}$. Its assumptions hold: $\omega\le N^{1/8-2/3}\le1$ and $100\Delta\le100N^{1/8}\le N$ by \eqref{eq:a1-R2} and $N^{7/8}\ge100$;
the bound on $\fn-\nsc$ is part (iv); and $-2-\omega\le\lambda_1$, $\lambda_N\le2+\omega$ by (b). Let $j\in\{1,\dots,N\}$ and $\hat\jmath=\min(j,N+1-j)$. If $\hat\jmath>3\Delta$, the lemma gives
$|\lambda_j-\gamma_j|\le3.5K^{C_5}\hat\jmath^{-1/3}N^{-2/3}$. If $\hat\jmath\le3\Delta$, then $1\le(3K^{C_5})^{1/3}\hat\jmath^{-1/3}$ and the lemma gives
\[
 |\lambda_j-\gamma_j|\le\big(K^{C_3}+9K^{2C_5/3}\big)N^{-2/3}\le3^{1/3}\big(K^{C_3+C_5/3}+9K^{C_5}\big)\hat\jmath^{-1/3}N^{-2/3}.
\]
In both cases, by the second relation in \eqref{eq:ledgerrel} and $K\ge216$,
\[
 |\lambda_j-\gamma_j|\le3^{1/3}\cdot10\,K^{\max(C_5,\,C_3+C_5/3)}\,\hat\jmath^{-1/3}N^{-2/3}\le15\,K^{C_4-1}\hat\jmath^{-1/3}N^{-2/3}\le K^{C_4}\hat\jmath^{-1/3}N^{-2/3}.
\]
The last sentence of the proposition follows because (a) is contained in Theorem~\ref{thm:K}(i) ($\{|E|\le\frac72,\ K^{C_1}N^{-1}\le\eta\le2\}\subset\bD_K$) and (b) is Theorem~\ref{thm:K}(iii).
\end{proof}

\medskip\noindent\textbf{Proof of Theorem~\ref{thm:K}\,(iv),(v).} Let $N_0\ge e^6$ be such that parts (i)--(iii) of Theorem~\ref{thm:K} hold for $N\ge N_0$ on the event $\Xi_K$ constructed in Sections~\ref{sec:strong} and~\ref{sec:norm}.
By Proposition~\ref{prop:rig}, (iv) and (v) hold at every point of the same event $\Xi_K$; no further exceptional set is needed, so $\PP(\Xi_K^c)$ is unchanged. \qed

\begin{remark}\label{rem:rigexp}
If the ledger is revised, Proposition~\ref{prop:rig} remains valid for any exponents satisfying \eqref{eq:ledgerrel} (and $C_1,\dots,C_5\ge1$); the proof uses nothing else.
The supremum in \textup{(iv)} is over all real $E$; the measurability of $\sup_E|\fn(E)-\nsc(E)|$ (Lemma~\ref{lem:weyl}(f)) is needed only to state \eqref{eq:nn}, not in the proof.
\end{remark}

\section{Exported results (for Parts~II and~III)}\label{sec:export}
All constants depend only on $(\vth,\delta_-,\delta_+,C_0)$ (normalised as in Convention~\ref{conv:const}); $C_1,\dots,C_5,c_0$ are those of \eqref{eq:ledger0}, $b=12\fa+12$ is the exponent of Section~\ref{sec:strong}, and $N_0$ is the threshold of Theorem~\ref{thm:K}.
All statements hold for every $N\ge N_0$, every $\kk\in\{\R,\C\}$ and every $\nu\in\mathcal E_N(\vth,\dpm,C_0)$, with the \emph{same} constants. Items (X1), (X8)--(X13) are pointers to numbered statements; the items (X2)--(X7) and (X14) are short consequences, proved in Section~\ref{sec:exportproofs}.
Section~\ref{sec:exportuse} lists which of these items Parts~II and~III actually use.

\subsection{The list}\label{sec:exportlist}
\begin{enumerate}[label=\textbf{(X\arabic*)},leftmargin=3em,itemsep=3pt]
\item \textbf{Theorems \ref{thm:21}, \ref{thm:22}} as stated (with one measurable event for all $z\in\bS_L$, and the same constants $A_0,C,c,\phi,N_0$ in both theorems), and their level form \textbf{Theorem~\ref{thm:K}}: for every $K\in[(\log N)^3,N^{c_0}]$ a closed event $\Xi_K$ (Definition~\ref{def:XiK}), $\PP(\Xi_K^c)\le e^{-K/2}$, on which (i)--(v) hold \emph{simultaneously} for all $z\in\bD_K$, all $E\in\R$ and all $j$.
The constants of Theorems~\ref{thm:21}, \ref{thm:22} are $\phi=1/(60\fa+88)$, $A_0=180\fa+264$, $C=1$, $c=\frac12$, with $\fa=1/\min(\vth,1)$ (Proposition~\ref{prop:reduction} and Appendix~\ref{app:ledger}).
\item \textbf{$N^{\eps}$-form.} Let $0<\eps\le c_0$ and $N^\eps\ge(\log N)^3$. With probability at least $1-\exp(-\frac12N^\eps)$, for all $|E|\le5$, $N^{-1+C_1\eps}\le\eta\le10$:
$\Lambda\le N^{C_2\eps}/(N\eta)$, $\ \Ld+\Lo\le N^{C_2\eps}\big(\sqrt{\im\msc/(N\eta)}+1/(N\eta)\big)$; and $\norm H\le2+N^{-2/3+C_3\eps}$; $|\lambda_j-\gamma_j|\le N^{C_4\eps}\,[\min(j,N+1-j)]^{-1/3}N^{-2/3}$ for all $j$; $\sup_{E\in\R}|\fn(E)-\nsc(E)|\le N^{-1+C_5\eps}$.
\item \textbf{Local law at one $z$, without the lattice.} For a fixed $z=E+\ii\eta\in\bD_K$ there is a closed event of probability at least $1-N^{10}e^{-p_1}\ge1-N^{10}\exp(-K^{3/2})$ on which
$\Lambda(z)\le K^{b+5}/(N\eta)$ and $\Ld(z)+\Lo(z)\le K^{b+6}\big(\sqrt{\im\msc(z)/(N\eta)}+1/(N\eta)\big)$; in particular the bounds of Theorem~\ref{thm:K}(i),(ii) hold at $z$.
\item \textbf{Improved averaged law outside the spectrum} (Propositions~\ref{prop:S3}, \ref{prop:S4}). With probability at least $1-N^{20}e^{-p_1}-N^{10}e^{-p_2}\ \ge1-N^{21}\exp(-K^{3/2})-N^{10}e^{-K}$, for all $z=E+\ii\eta\in\mathbb L\cap\bD^{\rm near}$ with $\varkappa N\eta\ge K^{b+3}$:
$\ \Lambda(z)\le K^{b+2}\big(\frac{\im\msc}{\varkappa N\eta}+\frac1{\varkappa(N\eta)^2}\big)$; in particular, if moreover $|E|>2$ and $\kappa=|E|-2\ge\eta$: $\ \Lambda(z)\le K^{b+2}\big(\frac4{N\kappa}+\frac1{\sqrt\kappa(N\eta)^2}\big)$.
Moreover there is a closed event $\Xi_K^{(4)}\subset\Xi_K$ with $\PP((\Xi_K^{(4)})^c)\le2N^{21}e^{-p_1}+2N^{10}e^{-p_2}\le e^{-K/2}$ on which, in addition to these bounds and to (i)--(v) of Theorem~\ref{thm:K}, for all $z\in\mathbb L\cap\bD^{\rm far}$ (i.e.\ $4\le|E|\le\Emax+1$, $\frac12N^{-1/2}\le\eta\le10$, on the lattice):
$\ \Lambda(z)\le K^{b+2}\big(\frac{\im\msc}{N\eta}+\frac1{(N\eta)^2}+\frac1N\big)$.
\item \textbf{Norm.} $\PP\big(\norm H>25K^{2\fa}\sqrt{C_0\log N}\big)\le N^2\vth^{-1}e^{-K^2}+2e^{-N\log N}\le2N^2\vth^{-1}e^{-K^2}$ (Lemma~\ref{lem:norm0} with $\kappa=K^2$), and $\norm H\le2+K^{C_3}N^{-2/3}$ on $\Xi_K$ (Proposition~\ref{prop:norm}).
\item \textbf{Local density.} On $\Xi_K$, for all $z=E+\ii\eta\in\bD_K$: $\ \im m(z)\le\im\msc(z)+K^{C_2}/(N\eta)$ and $\#\{j:|\lambda_j-E|\le\eta\}\le2N\eta\,\im\msc(z)+2K^{C_2}$.
\item \textbf{Below the domain (deterministic).} For every Hermitian matrix $H$, $E\in\R$, $0<\eta\le\eta'$ and all $i,j$:
$\im G_{ii}(E+\ii\eta)\le\frac{\eta'}\eta\im G_{ii}(E+\ii\eta')$ and
$|G_{ij}(E+\ii\eta)|\le|G_{ij}(E+\ii\eta')|+\frac{\eta'}{\eta}\max_k\im G_{kk}(E+\ii\eta')$.
\item \textbf{Large deviation bounds} for weighted linear, quadratic and bilinear forms with unequal, possibly vanishing variances: moments and truncation of a single entry (Lemma~\ref{lem:moments}); tail form (Lemma~\ref{lem:LD}) and moment form (Lemma~\ref{lem:LDmom});
the good event and its consequences \eqref{eq:Xi}--\eqref{eq:Gij} at any $z\in\C_+$ (Definition~\ref{def:good}, Lemma~\ref{lem:good}); the splitting of the product law and partial integration (Lemma~\ref{lem:split}).
\item \textbf{Fluctuation averaging} with arbitrary deterministic weights $|t_k|\le1/N$, any $z$ with $|z|\le N$, explicit $p$-dependence and random control parameter (Theorem~\ref{thm:FA}; the case $p=2$ is Proposition~\ref{prop:FAtwo}, which holds for $N$ above its own threshold $N_0$ of \eqref{eq:CV}, not necessarily the $N_0$ of Theorem~\ref{thm:K}).
\item \textbf{Resolvent algebra}: Lemma~\ref{lem:basic} (including the spectral form of $\tr G$ and the two-matrix resolvent identity), Lemma~\ref{lem:schur} (including the block form $G_{AA}=(M^A)^{-1}$, \eqref{eq:1overG}, the transposed form of (I2) and the determinant identity (I6)),
Lemmas~\ref{lem:minor}, \ref{lem:sce}; stability: Lemmas~\ref{lem:vecstab}--\ref{lem:stabfar}.
\item \textbf{Semicircle facts}: Lemmas~\ref{lem:msc}--\ref{lem:nsc} (algebraic definition, sizes with explicit constants $c_*=\frac1{150}$, $C_*=4$, Lipschitz bound, Stieltjes representation, $\rhosc\le1/\pi$, $\nsc$, $\gamma_j$, the bulk lower bound on $\rhosc$ and the count of classical locations).
\item \textbf{Helffer--Sj\"ostrand} formula (Lemma~\ref{lem:HS}), the counting lemma (Lemma~\ref{lem:count}), the conversion lemma (Lemma~\ref{lem:conv}) and the deterministic implication ``(i), (iii) $\Rightarrow$ (iv), (v)'' (Proposition~\ref{prop:rig}).
\item \textbf{Weyl's inequality}, reflection $\lambda_k(-A)=-\lambda_{n+1-k}(A)$ (Lemma~\ref{lem:weyl}(c)), Lipschitz continuity and measurability of eigenvalues, resolvent entries and the counting function (Lemma~\ref{lem:weyl}); grid net for the norm (Lemma~\ref{lem:net}).
\item \textbf{Families of ensembles.} If $(\nu^{(\theta)})_{\theta\in\Theta}$ is a finite family in $\mathcal E_N(\vth,\dpm,C_0)$, then for every $\theta$ the statements above hold for $\nu^{(\theta)}$ with constants and thresholds that do not depend on $\theta$;
if the matrices $H^{(\theta)}$ are realised on one probability space, all of them satisfy a given statement simultaneously off an event of probability at most $|\Theta|$ times the stated bound.
The class is closed under $H\mapsto-H$, monotone in the parameters, and \eqref{eq:subexp} is invariant under scaling with the same $\vth$ (Lemma~\ref{lem:mono}(ii),(iv),(v)).
\end{enumerate}
\noindent\emph{Not provided by Part~I:} bounds for $\eta<K^{C_1}/N$ other than (X7); isotropic bounds; anything for ensembles violating (A) (e.g.\ $H_t$ in the real case before rescaling, cf.\ Part~III; or a matrix with one entry or one row removed, cf.\ Part~II).

\subsection{Proofs of \textup{(X2)--(X7)} and \textup{(X14)}}\label{sec:exportproofs}
\emph{\textup{(X2)}.} Put $K:=N^\eps$. Then $(\log N)^3\le K\le N^{c_0}$, so Theorem~\ref{thm:K} applies; $K^{C_i}=N^{C_i\eps}$, $\bD_K=\{|E|\le5,\ N^{-1+C_1\eps}\le\eta\le10\}$ and $\PP(\Xi_K^c)\le\exp(-\frac12N^\eps)$. (For a fixed $\eps\in(0,c_0]$ the condition $N^\eps\ge(\log N)^3$ holds for all large $N$.)

\emph{\textup{(X3)}.} Let $z'\in\mathbb L\cap\bD^{\rm near}$ be the lattice point of Lemma~\ref{lem:lip}. The proof of Theorem~\ref{thm:K}(i),(ii) at the end of Section~\ref{sec:strong} uses, for the given $z$, only that the point of $\Omega$ lies in $\cG^2(z')$;
it gives $\Lambda(z)\le2K^{b+4}/(N\eta)$ and $\Ld(z)+\Lo(z)\le K^{b+5}\sqrt{\im\msc/(N\eta)}+3K^{b+5}/(N\eta)$ there. The event $\cG^2(z')$ is closed and $\PP(\cG^2(z')^c)\le N^{10}e^{-p_1}$ (Proposition~\ref{prop:S2}); $p_1\ge K^{3/2}$.

\emph{\textup{(X4)}.} Let $\mathbb L_a$ be the set of points $z=E+\ii\eta\in\mathbb L\cap\bD^{\rm near}$ with $\varkappa N\eta\ge\mathfrak E_3=K^{b+3}$, and $\Xi_K^{(3)}:=\bigcap_{z\in\mathbb L_a}\cG^3(z)$, a closed set. On $\Xi_K^{(3)}$ the bound \eqref{eq:improved} holds at all $z\in\mathbb L_a$, and by Propositions~\ref{prop:S2}, \ref{prop:S3} and $|\mathbb L|\le N^{10}$,
\[
 \PP\big((\Xi_K^{(3)})^c\big)\le\sum_{z\in\mathbb L_a}\Big(\PP\big(\cG^2(z)^c\big)+\PP\big(\cG^2(z)\setminus\cG^3(z)\big)\Big)\le N^{20}e^{-p_1}+N^{10}e^{-p_2},
\]
and $p_1\ge K^{3/2}$, $p_2\ge K$. For the particular case: if $|E|>2$ and $\kappa\ge\eta$ then $\im\msc\le4\eta/\varkappa$ by \eqref{eq:imm} and $\varkappa\ge\sqrt\kappa$, so $\frac{\im\msc}{\varkappa N\eta}\le\frac4{N\kappa}$ and $\frac1{\varkappa(N\eta)^2}\le\frac1{\sqrt\kappa(N\eta)^2}$.

Next let $\Xi_K^{(4)}:=\Xi_K\cap\Xi_K^{(3)}\cap\bigcap_{z'\in\mathbb L\cap\bD^{\rm far}}\cG^{\rm f}(z')$, a closed set; by Proposition~\ref{prop:S4} the far-regime bound holds on it. Since $\Xi_K\subset\cG^2(z)$ for all $z\in\mathbb L\cap\bD^{\rm near}$,
\[
 \big(\Xi_K^{(4)}\big)^c\subset\Xi_K^c\cup\bigcup_{z\in\mathbb L_a}\big(\cG^2(z)\setminus\cG^3(z)\big)\cup\bigcup_{z'}\cG_{\rm col}(z')^c\cup\bigcup_{z'}\big(\cG_{\rm col}(z')\setminus\cG^{\rm f}(z')\big).
\]
By the bound $\PP(\Xi_K^c)\le(N^{20}+N^{21})e^{-p_1}+N^6e^{-p_2}$ (proof of Proposition~\ref{prop:XiK}), Propositions~\ref{prop:S3}, \ref{prop:S4}, Definition~\ref{def:domains} ($\bigcup_{z'}\cG_{\rm col}(z')^c\subset\cG_*^c$, $\PP(\cG_*^c)\le46\,\fa N^{14}e^{-K^2}$), $|\mathbb L|\le N^{10}$ and $K^2\ge p_1$,
\[
 \PP\big((\Xi_K^{(4)})^c\big)\le\big(N^{20}+N^{21}+46\,\fa N^{14}\big)e^{-p_1}+\big(N^6+N^{10}\big)e^{-p_2}\le2N^{21}e^{-p_1}+2N^{10}e^{-p_2}\le e^{-K/2},
\]
the last step as in the proof of Proposition~\ref{prop:XiK} ($21\log N+\log4\le\frac12K$).

\emph{\textup{(X5)}.} This is Lemma~\ref{lem:norm0} (see Step~1 of the proof of Proposition~\ref{prop:norm}, where $2e^{-N\log N}\le e^{-K^2}\le N^2\vth^{-1}e^{-K^2}$ is shown) and Theorem~\ref{thm:K}(iii).

\emph{\textup{(X6)}.} On $\Xi_K$, by Theorem~\ref{thm:K}(i), $\im m\le\im\msc+|m-\msc|\le\im\msc+K^{C_2}/(N\eta)$, and by Lemma~\ref{lem:basic}(e), $\#\{j:|\lambda_j-E|\le\eta\}\le2N\eta\,\im m(z)$.

\emph{\textup{(X7)}.} Write $z=E+\ii\eta$, $z'=E+\ii\eta'$. The first inequality is Lemma~\ref{lem:basic}(f) ($\eta\mapsto\eta\,\im G_{ii}(E+\ii\eta)$ is nondecreasing and $\eta\le\eta'$: $\eta\im G_{ii}(z)\le\eta'\im G_{ii}(z')$).
For the second, the resolvent identity (Lemma~\ref{lem:basic}(c)) gives $G_{ij}(z)-G_{ij}(z')=\ii(\eta-\eta')\sum_kG_{ik}(z)G_{kj}(z')$, and by the Cauchy--Schwarz inequality, the Ward identity (Lemma~\ref{lem:basic}(b)) and the first inequality,
\begin{align*}
 \Big|\sum_kG_{ik}(z)G_{kj}(z')\Big|&\le\Big(\frac{\im G_{ii}(z)}{\eta}\Big)^{1/2}\Big(\frac{\im G_{jj}(z')}{\eta'}\Big)^{1/2}\le\Big(\frac{\eta'\,\im G_{ii}(z')}{\eta^2}\Big)^{1/2}\Big(\frac{\im G_{jj}(z')}{\eta'}\Big)^{1/2}\\
 &=\frac1\eta\big(\im G_{ii}(z')\,\im G_{jj}(z')\big)^{1/2}\le\frac1\eta\max_k\im G_{kk}(z').
\end{align*}
Hence $|G_{ij}(z)-G_{ij}(z')|\le\frac{\eta'-\eta}{\eta}\max_k\im G_{kk}(z')\le\frac{\eta'}\eta\max_k\im G_{kk}(z')$.

\emph{\textup{(X14)}.} All constants and thresholds of Part~I are functions of $(\vth,\delta_-,\delta_+,C_0)$ only (Convention~\ref{conv:const}); nothing depends on the individual laws $\nu_{ij}$ or on the variances beyond (A)--(C). The simultaneous statement is the union bound.
The closure properties are Lemma~\ref{lem:mono}(ii),(iv),(v).

\subsection{What Parts~II and~III use}\label{sec:exportuse}
For the convenience of the reader we list the statements of Part~I that are cited in Parts~II and~III, with the names used there.
\begin{center}
\renewcommand{\arraystretch}{1.12}\small
\begin{tabular}{|p{0.11\textwidth}|p{0.50\textwidth}|p{0.29\textwidth}|}
\hline
name in B, C & statement of Part~I & item above\\\hline
(BB1) & Theorem~\ref{thm:21}(ii), i.e.\ \eqref{eq:Lambdaodfinal}, with outer probability over all $z\in\bS_L$ and one measurable event & (X1)\\
(BB2), (RA1) & Theorem~\ref{thm:22}: \eqref{eq:rigidity} and \eqref{eq:nn} & (X1)\\
(S1) & Lemma~\ref{lem:msc}(a)--(c) & (X11)\\
(S2) & Lemma~\ref{lem:basic}(a),(b),(c),(d),(h) & (X10)\\
(S3) & Lemma~\ref{lem:schur}\,(I2), (I6) & (X10)\\
(S4) & Lemma~\ref{lem:weyl}(b)--(e); Lemma~\ref{lem:mono}(iv) & (X13), (X14)\\
(S5), (RA6) & Lemma~\ref{lem:moments} & (X8)\\
(S6) & Lemma~\ref{lem:split} & (X8)\\
(RA2), (RA3) & Lemma~\ref{lem:mono}(ii), (v) & (X14)\\
(RA4) & Lemma~\ref{lem:weyl}(a),(b) & (X13)\\
(RA5) & Lemma~\ref{lem:nsc}(a),(c),(d),(e) (Part~II also uses (b)) & (X11)\\
--- & Definitions~\ref{def:ens}, \ref{def:ass}; \eqref{eq:green}--\eqref{eq:defS}; statements of Theorems~\ref{thm:23}, \ref{thm:24} & Section~\ref{sec:main}\\\hline
\end{tabular}
\end{center}
\noindent Parts~II and~III do \emph{not} use: the level form (Theorem~\ref{thm:K}) and the items (X2)--(X7), (X9), (X12); Theorem~\ref{thm:21}(i),(iii); any bound for a matrix that is not in the class $\mathcal E_N$.
Part~II uses the simultaneity in $z$ of Theorem~\ref{thm:21}(ii) (on a segment of energies at fixed $\eta$), which is provided by the measurable event $\Xi_K$; Part~III uses Theorem~\ref{thm:22} at each fixed time for the rescaled matrix $\widetilde H_t$, which lies in one class $\mathcal E_N$ for all $t$.
The items (X2)--(X7), (X9), (X12) are recorded because they are the natural inputs for extensions (other observables at the edge, the removal of the moment matching conditions, eigenvector bounds).


% ================= PART II =================
\clearpage
\part{Universality of the extreme eigenvalues (Theorem 2.4)}
\useSharedCounters
\section{Edge universality}\label{sec:edge}
\subsection{Inputs from Part~I, conventions, statement}\label{sec:inputs}

\paragraph{Setting.} We use the model of \S\ref{sec:main}. $\K\in\{\R,\C\}$, $N\ge3$, $\mathcal P_N=\{(i,j):1\le i\le j\le N\}$. An
\emph{ensemble} $M=(\nu_p)_{p\in\mathcal P_N}$ (of size $N$ over $\K$) is a family of Borel probability measures on $\K$
(Definition~\ref{def:ens}); the sample space is $\Omega=\K^{\mathcal P_N}$, the law is the product measure $\P_M=\bigotimes_p\nu_p$, $\E_M$ is the
expectation, and $H=H(x)$ is the Hermitian matrix with $h_{ij}=x_{ij}$, $h_{ji}=\overline{x_{ij}}$ $(i<j)$, $h_{ii}=\re x_{ii}$. Further
$\sigma_p^2=\sigma_{ij}^2=\int|y|^2\nu_{ij}(dy)$ and $B=(\sigma_{ij}^2)$. For $\Par=(\vartheta,\delta_-,\delta_+,C_0)\in(0,\infty)^4$ we say that
$M$ \emph{satisfies} $\fA(\Par)$ if $M\in\mathcal E_N(\vartheta,\delta_-,\delta_+,C_0)$ (Definition~\ref{def:ass}): real diagonal (R), the
subexponential decay (D) \eqref{eq:subexp} (note the strict inequality $|y|>x\sigma_p$ and
that nothing is divided by $\sigma_p$), centring (Z), and \textup{(A)}, \textup{(B)}, \textup{(C)}. Two ensembles $M^v=(\nu^v_p)$, $M^w=(\nu^w_p)$
of the same size over the same $\K$ \emph{have matching second moments} if \eqref{eq:2m} holds (this is (2.40) of the original, stated on the
laws of the entries); we write $\P^v=\P_{M^v}$, $\E^v=\E_{M^v}$, and similarly for $w$.

\paragraph{Conventions.} For a Hermitian $A\in M_n(\C)$, $n\ge1$, and $z\in\C_+$: $G_A(z)=(A-z)^{-1}$, and $\lambda_1(A)\le\dots\le\lambda_n(A)$
are the eigenvalues counted with multiplicity; $\lambda_j=\lambda_j(H)$. $\msc$ is the function \eqref{eq:defmsc}, i.e.\ (Lemma~\ref{lem:msc}(a))
the root of $m^2+zm+1=0$ with $\im m>0$. $\varrho_{sc}$, $n_{sc}$ and the classical locations $\gamma_j$ are those of \eqref{eq:sc}; by
Lemma~\ref{lem:nsc}(c), $Nn_{sc}(\gamma_j)=j$ and $\gamma_N=2$. $\Lambda_d,\Lambda_o$ and $S_L$ are defined in \eqref{eq:defLambda}, \eqref{eq:defS}.
$\Log$ is the principal branch of the logarithm on $\C\setminus(-\infty,0]$. For a function $g$ of the spectral parameter and fixed
$z_1,z_2$ we write $\br{g}:=g(z_2)-g(z_1)$. $\P^*$ denotes outer probability (\S\ref{sec:green}). All logarithms are natural; $\log_+x:=\max(\log x,0)$ for $x>0$ and $\log_+0:=0$.
\emph{Local constants.} The absolute constants $c_1,c_2,\dots$ of \S\ref{sec:step}, the named constants $C_2,C_3,C_4$, $C_F$, $C_q$, $C_\Phi$, the moment constants $\mathfrak m_k$, the control parameter $\Psi$ and the thresholds $N_1(\Par,\eg)$, $N_0^{\rm edge}(\Par)$
are defined in this section and are used only here; they are unrelated to the constants and quantities with the same names in Parts~I and~III (for instance the exponents $C_1,\dots,C_5$ of Appendix~\ref{app:ledger},
the constant $C_F$ of Theorem~\ref{thm:FA}, or the control parameter $\Psi$ of Section~\ref{sec:weak}).
Likewise $\delta$ denotes only the exponent of Theorem~\ref{thm:24}.

\paragraph{The probabilistic input from Part~I.} Part~II uses Theorems~\ref{thm:21}(ii) and \ref{thm:22} as black boxes, in the following form.
For every $\Par$ there are constants $A_0>1$, $C>0$, $c>0$, $\phi\in(0,1)$ and $N_0\ge3$, depending only on $\Par$, the same in both
theorems (see the remark after Theorem~\ref{thm:22}), such that for every $N\ge N_0$, $\K\in\{\R,\C\}$, every ensemble $M$ of size $N$ over $\K$
with $\fA(\Par)$ and every real $L$ with $A_0\log\log N\le L\le\log(10N)/(10\log\log N)$ \eqref{eq:Lbound}:
\begin{align*}
U_L&:=\bigcup_{z\in S_L}\Big\{x\in\Omega:\ \Lambda_d(z)+\Lambda_o(z)\ge(\log N)^{4L}\sqrt{\frac{\im\msc(z)}{N\eta}}+\frac{(\log N)^{4L}}{N\eta}\Big\}\qquad(\eta=\im z),\\
\textbf{(BB1)}\quad&\P^*_M(U_L)\le Ce^{-c(\log N)^{\phi L}},\\
\textbf{(BB2)}\quad&\P_M\Big(\exists j:\ |\lambda_j-\gamma_j|\ge(\log N)^{L}\big[\min(j,N-j+1)\big]^{-1/3}N^{-2/3}\Big)\le Ce^{-c(\log N)^{\phi L}}.
\end{align*}
(BB1) is \eqref{eq:Lambdaodfinal} and (BB2) is \eqref{eq:rigidity}. The union in (BB1) is over \emph{all} $z\in S_L$, so the bound holds
simultaneously for all $z\in S_L$; since the union need not be measurable a priori, (BB1) is stated with the outer probability, and
Theorem~\ref{thm:21} provides in addition one measurable event of probability $\ge1-Ce^{-c(\log N)^{\phi L}}$ on which none of these events occurs.
We use (BB1) only through the following consequence of the definition of $\P^*$:
\emph{if $\Omega'\subset\Omega$ is a Borel set with $\Omega'\subset U_L$, then $\P_M(\Omega')\le Ce^{-c(\log N)^{\phi L}}$.}
The event $\Omega_{\rm edge}$ of Lemma~\ref{lem:edge} is of this kind: it is \emph{open}, in particular Borel (Step~0 of the proof of Lemma~\ref{lem:edge}). The event of (BB2) is closed, because each $x\mapsto\lambda_j(H(x))$ is continuous
(Lemma~\ref{lem:weyl}(d)).
\emph{Nothing else from Theorems~\ref{thm:21}--\ref{thm:23} is used:} not \eqref{eq:Lambdafinal}, \eqref{eq:443}, \eqref{eq:nn}; no large deviation
or fluctuation averaging lemma, no norm bound, no resolvent bound below the scale of $S_L$, and no bound for matrices with an entry or a row removed.

\paragraph{Deterministic and elementary facts from Part~I.} We label them (S1)--(S6) (labels local to this section) and cite them by their numbers in Part~I.
\begin{itemize}\itemsep0pt
\item[(S1)] Lemma~\ref{lem:msc}(a)--(c): $\msc(z)^2+z\msc(z)+1=0$, $\im\msc(z)>0$ and $|\msc(z)|\le1$ for $z\in\C_+$.
\item[(S2)] Lemma~\ref{lem:basic}(a),(c),(d),(h): $G_A(z)$ exists, $\|G_A(z)\|\le(\im z)^{-1}$, $\tr G_A(z)=\sum_j(\lambda_j(A)-z)^{-1}$,
$G_A-G_B=G_A(B-A)G_B$ for Hermitian $A,B$ of equal size.
\item[(S3)] Lemma~\ref{lem:schur}: the Schur complement formulas (I2) and the determinant identity (I6), $\det(A-z)=\det(A^{(a)}-z)/G_{aa}$.
\item[(S4)] Lemma~\ref{lem:weyl}(c),(d): $\lambda_1(A)=-\lambda_n(-A)$ for Hermitian $A\in M_n(\C)$; $x\mapsto H(x)$ is $\R$-linear and
$x\mapsto\lambda_j(H(x))$ is ($\sqrt2$-Lipschitz) continuous on $\Omega$. Lemma~\ref{lem:mono}(iv): reflection of an ensemble.
\item[(S5)] Lemma~\ref{lem:moments}: moments under \eqref{eq:subexp}; if $\sigma_p=0$ then $\nu_p=\delta_0$.
\item[(S6)] Lemma~\ref{lem:split}: splitting of the product law and partial integration (Tonelli, Fubini).
\end{itemize}
In addition Lemma~\ref{lem:nsc} (the semicircle distribution function and the classical locations) is used in \S\ref{sec:sandwich}.

\begin{theorem}[Theorem~\ref{thm:24} with explicit exponents]\label{thm:24B}
Let $\e:=\frac1{100}$ and $\delta:=\frac1{20}$. For every $\Par=(\vartheta,\delta_-,\delta_+,C_0)\in(0,\infty)^4$ there is $N_0^{\rm edge}(\Par)$ such that for
all $N\ge N_0^{\rm edge}(\Par)$, $\K\in\{\R,\C\}$, all ensembles $M^v,M^w$ of size $N$ over $\K$ which satisfy $\fA(\Par)$ and have matching second moments
\eqref{eq:2m}, and all $s\in\R$, the inequalities \eqref{eq:tw} (for $\lambda_N$) and \eqref{eq:tw1} (for $\lambda_1$) hold.
\end{theorem}
\noindent Theorem~\ref{thm:24B} implies Theorem~\ref{thm:24}: given $\vartheta$, take $\e=\frac1{100}$, $\delta=\frac1{20}$ (absolute constants, the
same for every $\vartheta$); given $\delta_\pm,C_0$, take $N_0:=N_0^{\rm edge}(\vartheta,\delta_-,\delta_+,C_0)$. The quantifier order is
$\forall\vartheta\ \exists\e,\delta\ \forall\delta_\pm,C_0\ \exists N_0\ \forall N\ \forall\K,M^v,M^w\ \forall s$.
In particular $N_0$ depends only on $\Par$, and not on $s$, $\K$, $M^v$, $M^w$. The proof is in \S\ref{sec:proof24}.

\paragraph{Parameters (fixed for the whole section).} Let $A_0,C,c,\phi,N_0$ be the constants of (BB1), (BB2) for the given $\Par$. For
$N\ge16$ (so that $\log\log N>1$) and a parameter $\eg\in(0,\frac1{10}]$ put
\begin{equation}\label{eq:B2par}
L_0:=A_0\log\log N,\qquad W:=(\log N)^{L_0}N^{-2/3},\qquad E_L:=2+2W,
\end{equation}
\begin{equation}\label{eq:B2par2}
\eta:=N^{-2/3-\eg},\qquad \e:=\tfrac12\eg,\qquad \ell:=\tfrac13N^{-2/3-\e},\qquad \Psi:=4N^{-1/3+2\eg},\qquad t:=N^{-1/2+2\eg}.
\end{equation}
The comparison branch (\S\ref{sec:step}--\S\ref{sec:gfc}) works with any fixed $\eg\in(0,\frac1{10}]$; the rigidity branch (\S\ref{sec:sandwich})
and the proof of Theorem~\ref{thm:24B} use $\eg=\frac1{50}$, i.e.\ $\e=\frac1{100}$ (the $\e$ of Theorem~\ref{thm:24B}). Note
\begin{equation}\label{eq:B2etaell}
\frac{\eta}{\ell}=3N^{-\eg+\e}=3N^{-\e},\qquad 3\ell=N^{-2/3-\e}.
\end{equation}

\begin{lemma}[largeness of $N$; elementary]\label{lem:N0}
Let $\Par$ be given and let $A_0,C,c,\phi,N_0$ be as above. For $\eg\in(0,\frac1{10}]$ put
\begin{equation}\label{eq:B2ustar}
u_*:=\max\Big\{\sqrt{\tfrac{48(4A_0+5)}{\eg}},\ \tfrac2{\phi A_0},\ \sqrt{\tfrac{2\,b_2}{\phi A_0}},\ \sqrt{\tfrac{24(1+b_3)}{\eg\vartheta}},\ \log_+\tfrac{\log_+C_0}{2\eg}\Big\},
\end{equation}
where $b_2:=\log_+\big((\log_+C+20)/c\big)$ and $b_3:=\log\big(20+\log_+(1/\vartheta)\big)$, and
\[
N_1(\Par,\eg):=\max\big\{N_0,\ e^{27},\ \exp(e^{u_*})\big\}.
\]
Then $N_1$ depends only on $(\Par,\eg)$, and for every $N\ge N_1$, with the parameters \eqref{eq:B2par}, \eqref{eq:B2par2}:
\begin{itemize}\itemsep1pt
\item[\textup{(N1)}] $L_0$ satisfies \eqref{eq:Lbound}, and $\eta>N^{-1}(\log N)^{10L_0}$; hence $y+i\eta\in S_{L_0}$ for every $y\in[-5,5]$, and \textup{(BB1), (BB2)} hold with $L=L_0$.
\item[\textup{(N2)}] $3(\log N)^{4L_0}\le N^{\eg/2}$ and $2(\log N)^{L_0}+1\le N^{\eg}$; in particular $2W+N^{-2/3}\le N^{-2/3+\eg}$.
\item[\textup{(N3)}] $Ce^{-c(\log N)^{\phi L_0}}\le N^{-20}$ and $\vartheta^{-1}\exp(-N^{\eg\vartheta})\le N^{-20}$.
\item[\textup{(N4)}] $C_0^{1/2}\le N^{\eg}$, $\Psi\le\frac18$ and $40\,t\le\Psi$.
\item[\textup{(N5)}] $\frac3\pi\big(4(\log N)^{3L_0/2}+1\big)N^{-\e}+4N^{-2\e}(\log N)^{L_0}\le\frac19$, $\ \frac6\pi N^{-\e}\le\frac19$, and $2\ell\le W$.
\end{itemize}
\end{lemma}
\noindent This lemma is the only place where the threshold for $N$ is enlarged in the rigidity branch; later statements refer to (N1)--(N5) by number.
The threshold of the comparison theorem (Theorem~\ref{thm:gfc}) and the final condition (N6) are added in \S\ref{sec:proof24}.

\begin{proof}
Write $u:=\log\log N$, so that $\log N=e^u$ and $(\log N)^{aL_0}=\exp(aA_0u^2)$ for every $a>0$. Since $N\ge e^{27}$, $\log N\ge27$ and $u\ge\log27>1$.
Since $N\ge\exp(e^{u_*})$, $u\ge u_*$.

\emph{An elementary inequality.} For $u\ge0$, $e^u=\sum_{k\ge0}u^k/k!\ge u^4/24$. Hence, for $\alpha,\beta\ge0$ and $\gamma>0$:
\begin{equation}\tag{E}\label{eq:B2E}
u\ge1\ \text{ and }\ u^2\ge\frac{24(\alpha+\beta)}{\gamma}\quad\Longrightarrow\quad\alpha u^2+\beta\le(\alpha+\beta)u^2\le\frac{\gamma u^4}{24}\le\gamma e^u .
\end{equation}
By the choice of $u_*$ we obtain from \eqref{eq:B2E} and directly:
\begin{align}
&4A_0u^2+5\le\tfrac{\eg}2\,e^u\qquad&&(\alpha=4A_0,\ \beta=5,\ \gamma=\eg/2),\label{eq:B2U1}\\
&\phi A_0u^2\ge2u\ \text{ and }\ \phi A_0u^2\ge2b_2,\ \text{ hence }\ \phi A_0u^2\ge u+b_2,\label{eq:B2U2}\\
&u+b_3\le u^2+b_3\le\eg\vartheta\,e^u\qquad&&(\alpha=1,\ \beta=b_3,\ \gamma=\eg\vartheta),\label{eq:B2U3}\\
&e^u\ge\frac{\log_+C_0}{2\eg}\qquad&&(u\ge\log_+\tfrac{\log_+C_0}{2\eg}).\label{eq:B2U4}
\end{align}
(For \eqref{eq:B2U4}: if $\log_+C_0=0$ it is trivial; otherwise $u\ge\log\frac{\log C_0}{2\eg}$.)

\emph{\textup{(N1)}.} By \eqref{eq:B2U1} and $\eg\le\frac1{10}$: $4A_0u^2\le\frac1{20}e^u$, i.e.\ $80A_0u^2\le\log N$. Hence $10L_0\log\log N=10A_0u^2\le\log N<\log(10N)$,
which is the upper bound in \eqref{eq:Lbound}; the lower bound holds with equality. Moreover $(\log N)^{10L_0}=\exp(10A_0u^2)\le\exp(\frac18\log N)=N^{1/8}$, so
$N^{-1}(\log N)^{10L_0}\le N^{-7/8}<N^{-23/30}\le N^{-2/3-\eg}=\eta$, since $\frac23+\eg\le\frac{23}{30}<\frac78$ and $N>1$. For $|y|\le5$ the point $y+i\eta$
has $|\re|\le5$ and $N^{-1}(\log N)^{10L_0}<\eta\le1\le10$, so it lies in $S_{L_0}$ \eqref{eq:defS}. Finally $N\ge N_0$, so (BB1), (BB2) apply with $L=L_0$.

\emph{\textup{(N2)}.} By \eqref{eq:B2U1}, $\log3+4A_0u^2\le5+4A_0u^2\le\frac{\eg}2\log N$, i.e.\ $3(\log N)^{4L_0}\le N^{\eg/2}$. Since $\log N\ge1$ and $L_0>0$,
$(\log N)^{L_0}\ge1$, so $2(\log N)^{L_0}+1\le3(\log N)^{L_0}\le3(\log N)^{4L_0}\le N^{\eg/2}\le N^{\eg}$. Multiplying by $N^{-2/3}$ gives $2W+N^{-2/3}\le N^{-2/3+\eg}$.

\emph{\textup{(N3)}.} By \eqref{eq:B2U2}, $(\log N)^{\phi L_0}=\exp(\phi A_0u^2)\ge e^{b_2}e^u\ge\frac{\log_+C+20}c\,\log N$, so
$c(\log N)^{\phi L_0}\ge(\log_+C+20)\log N\ge\log C+20\log N$ (as $\log N\ge1$), i.e.\ $Ce^{-c(\log N)^{\phi L_0}}\le N^{-20}$.
By \eqref{eq:B2U3}, $N^{\eg\vartheta}=\exp(\eg\vartheta\log N)=\exp(\eg\vartheta e^u)\ge e^{b_3}e^u=(20+\log_+(1/\vartheta))\log N\ge20\log N+\log(1/\vartheta)$,
hence $\vartheta^{-1}\exp(-N^{\eg\vartheta})\le\vartheta^{-1}\cdot\vartheta N^{-20}=N^{-20}$.

\emph{\textup{(N4)}.} By \eqref{eq:B2U4}, $\eg\log N\ge\frac12\log C_0$, i.e.\ $C_0^{1/2}\le N^{\eg}$. Since $\eg\le\frac1{10}$, $\frac13-2\eg\ge\frac2{15}$ and
$N^{1/3-2\eg}\ge N^{2/15}\ge e^{(2/15)\cdot27}=e^{3.6}>32$, i.e.\ $\Psi=4N^{-(1/3-2\eg)}\le\frac18$. Next $\Psi/t=4N^{1/6}\ge4e^{27/6}>40$.

\emph{\textup{(N5)}.} Put $Y:=(\log N)^{3L_0/2}=\exp(\frac32A_0u^2)\ge1$. Then $(\log N)^{L_0}\le Y$, $N^{-2\e}\le N^{-\e}$, $\frac3\pi<1$, so the left side of the first
inequality is at most $5YN^{-\e}+4YN^{-\e}=9YN^{-\e}$. By \eqref{eq:B2U1} and $\e=\eg/2$: $\log81+\frac32A_0u^2\le5+4A_0u^2\le\e\log N$, i.e.\ $81Y\le N^{\e}$, so
$9YN^{-\e}\le\frac19$. Also $\frac6\pi N^{-\e}\le2N^{-\e}\le9YN^{-\e}\le\frac19$. Finally $2\ell=\frac23N^{-2/3-\e}\le N^{-2/3}\le(\log N)^{L_0}N^{-2/3}=W$.
\end{proof}

\noindent\emph{Architecture of the proof of Theorem~\ref{thm:24B}.}
(1) \S\ref{sec:prelim}: hybrids lie in one class (Lemma~\ref{lem:hyb}); $\msc$ near the edge (Lemma~\ref{lem:mscB}); edge resolvent bounds from (BB1) (Lemma~\ref{lem:edge}).
(2) \S\ref{sec:X}: the functional $X$ and exact determinant-ratio formulas (deterministic).
(3) \S\ref{sec:sandwich}: $\E q(X^-_E(H))-N^{-20}\le\P(\lambda_N<E)\le\E q(X^+_E(H))+N^{-20}$, from (BB2) alone.
(4) \S\ref{sec:step}--\S\ref{sec:decor}: one Lindeberg step. (5) \S\ref{sec:gfc}: telescoping over the $N(N+1)/2$ entries (Theorem~\ref{thm:gfc}).
(6) \S\ref{sec:proof24}: Proposition~\ref{prop:oneside} combines (3) and (5); Theorem~\ref{thm:24B} follows, for $\lambda_1$ via $H\mapsto-H$.
The \emph{rigidity branch} (3) uses (BB2) only; the \emph{comparison branch} (4)--(5) uses (BB1) only; they meet in Proposition~\ref{prop:oneside}.

\subsection{Preliminaries: hybrid ensembles and the edge regime}\label{sec:prelim}

\begin{definition}[hybrids, removed entry, negation]\label{def:hyb}
Let $M^v=(\nu^v_p)$, $M^w=(\nu^w_p)$ be ensembles of size $N$ over $\K$, $\Gamma:=N(N+1)/2=\#\mathcal P_N$, and fix any bijection
$\varphi:\mathcal P_N\to\{1,\dots,\Gamma\}$; $p_k:=\varphi^{-1}(k)$. For $0\le k\le\Gamma$ the \emph{hybrid} $M_k=(\nu^{(k)}_p)_p$ has
$\nu^{(k)}_p:=\nu^w_p$ if $\varphi(p)\le k$ and $\nu^{(k)}_p:=\nu^v_p$ otherwise; thus $M_0=M^v$ and $M_\Gamma=M^w$.
For a pair $p=(a,b)\in\mathcal P_N$ and $v\in\K$ put
\[
V_p(v):=vE^{(ab)}+\bar vE^{(ba)}\quad(a<b),\qquad V_p(v):=(\re v)\,E^{(aa)}\quad(a=b),
\]
where $E^{(ab)}$ is the matrix unit (so $V_{(a,a)}(v)=vE^{(aa)}$ for real $v$). For a Hermitian $H$ let $H^{[p]}:=H-V_p(h_{ab})$, i.e.\ $H$ with the entries
$(a,b)$ and $(b,a)$ set to $0$. For an ensemble $M=(\nu_p)$, $-M:=(\nu^-_p)$ denotes the ensemble of the image laws $\nu^-_p(S):=\nu_p(-S)$.
\end{definition}

\begin{lemma}[stability of the assumptions; structure of one step]\label{lem:hyb}
Let $M^v,M^w$ be ensembles of size $N$ over $\K$ which satisfy $\fA(\Par)$ and have matching second moments \eqref{eq:2m}.
\begin{itemize}\itemsep1pt
\item[\textup{(i)}] $\sigma^v_p=\sigma^w_p$ for all $p\in\mathcal P_N$; hence $B^v=B^w$.
\item[\textup{(ii)}] Every ensemble $M=(\nu_p)$ with $\nu_p\in\{\nu^v_p,\nu^w_p\}$ for each $p$ satisfies $\fA(\Par)$ \emph{with the same} $\Par$, and its variance
matrix is $B^v$. In particular every hybrid $M_k$, $0\le k\le\Gamma$, satisfies $\fA(\Par)$.
\emph{Structure of one step:} let $1\le k\le\Gamma$, $p:=p_k=(a,b)$, $J:=\{p\}$, $J^c:=\mathcal P_N\setminus\{p\}$, and write $x=(x',x_p)$ with $x'=x_{J^c}$. Then
$\P_{M_{k-1}}$ and $\P_{M_k}$ have the same marginal $\P':=\bigotimes_{q\ne p}\nu^{(k)}_q$ on $\K^{J^c}$, and under each of them $x'$ and $x_p$ are independent,
$x_p$ having the law $\nu^v_p$ under $\P_{M_{k-1}}$ and $\nu^w_p$ under $\P_{M_k}$. The matrix $Q(x):=H(x)^{[p]}$ depends only on $x'$, and
$H(x)=Q(x)+V_p(x_p)$ for every $x\in\Omega$. Consequently, for every bounded Borel function $\Phi$ on the Hermitian $N\times N$ matrices,
\begin{equation}\label{eq:B2step}
\begin{aligned}
\E_{M_{k-1}}\Phi(H)&=\int_{\K^{J^c}}\Big[\int_\K\Phi\big(Q(x')+V_p(v)\big)\,\nu^v_p(dv)\Big]\P'(dx'),\\
\E_{M_{k}}\Phi(H)&=\int_{\K^{J^c}}\Big[\int_\K\Phi\big(Q(x')+V_p(v)\big)\,\nu^w_p(dv)\Big]\P'(dx'),
\end{aligned}
\end{equation}
the inner integrals being Borel functions of $x'$; and $\E_{M_{k-1}}\Theta(Q)=\E_{M_k}\Theta(Q)=\int\Theta(Q(x'))\P'(dx')$ for every bounded Borel $\Theta$.
\item[\textup{(iii)}] $-M^v$ and $-M^w$ satisfy $\fA(\Par)$ and have matching second moments \eqref{eq:2m}. Moreover $H(-x)=-H(x)$, the law of $H$ under $\P_{-M}$ is the
law of $-H$ under $\P_M$ (for $M\in\{M^v,M^w\}$), and $\lambda_1(A)=-\lambda_N(-A)$ for every Hermitian $A\in M_N(\C)$.
\end{itemize}
\end{lemma}
\begin{proof}
(i) For $(l,u)=(1,1)$, \eqref{eq:2m} reads $\int|y|^2\nu^v_p(dy)=\int|y|^2\nu^w_p(dy)$, i.e.\ $(\sigma^v_p)^2=(\sigma^w_p)^2$; both are nonnegative.

(ii) Let $M=(\nu_p)$ with $\nu_p\in\{\nu^v_p,\nu^w_p\}$. The second moment of $\nu_p$ is $\sigma_p^2$ with $\sigma_p:=\sigma^v_p=\sigma^w_p$ by (i), so the variance matrix of $M$ is $B^v$.
We check the conditions of Definition~\ref{def:ass} one by one.
(R) and (Z) are properties of each single law $\nu_p$, and they hold for $\nu^v_p$ and for $\nu^w_p$.
\eqref{eq:subexp} for $M$ asks, for each $p$ and $x\ge1$, that $\nu_p(\{|y|>x\sigma_p\})\le\vartheta^{-1}e^{-x^\vartheta}$ with the standard deviation $\sigma_p$ of $\nu_p$;
since $\sigma_p=\sigma^v_p=\sigma^w_p$, this is \eqref{eq:subexp} for $M^v$ (if $\nu_p=\nu^v_p$) or for $M^w$ (if $\nu_p=\nu^w_p$) at the pair $p$.
(A), (B), (C) are conditions on the variance matrix alone, which is $B^v$, and $M^v$ satisfies them with the parameters $\delta_\pm$, $C_0$.
For the hybrids, $\nu^{(k)}_p\in\{\nu^v_p,\nu^w_p\}$ by Definition~\ref{def:hyb}.

\emph{Structure of one step.} For $q\ne p_k$ we have $\varphi(q)\ne k$, so $\varphi(q)\le k-1$ if and only if $\varphi(q)\le k$; hence $\nu^{(k-1)}_q=\nu^{(k)}_q$. For $q=p_k$:
$\nu^{(k-1)}_{p}=\nu^v_{p}$ (as $\varphi(p)=k>k-1$) and $\nu^{(k)}_{p}=\nu^w_{p}$. By Lemma~\ref{lem:split}(a), applied with $I=\mathcal P_N$ and $J=\{p\}$,
$\P_{M_{k-1}}$ is the image of $\P'\otimes\nu^v_p$ and $\P_{M_k}$ the image of $\P'\otimes\nu^w_p$ under $(x',x_p)\mapsto x$; this is the statement on the
marginals and on independence. By Definition~\ref{def:ens}, the entries of $H(x)$ outside $\{(a,b),(b,a)\}$ are functions of $x'$, and the entries
$(a,b),(b,a)$ of $H(x)$ are $x_p,\overline{x_p}$ if $a<b$, resp.\ the entry $(a,a)$ is $\re x_p$ if $a=b$. Hence $Q(x)=H(x)-V_p(h_{ab}(x))$ has the entries of $H(x)$
outside $\{(a,b),(b,a)\}$ and zeros there, so it depends only on $x'$, and $V_p(x_p)=V_p(h_{ab}(x))$ (Definition~\ref{def:hyb}: for $a=b$ only $\re x_p=h_{aa}$ enters)
gives $H(x)=Q(x)+V_p(x_p)$. The map $(x',v)\mapsto Q(x')+V_p(v)$ is $\R$-linear, hence continuous, so $(x',v)\mapsto\Phi(Q(x')+V_p(v))$ is a bounded Borel function on
$\K^{J^c}\times\K$. Lemma~\ref{lem:split}(b),(c) (partial integration over the coordinate $x_p$) gives \eqref{eq:B2step}, with inner integrals Borel in $x'$.
The last statement is \eqref{eq:B2step} for the function $\Phi$ with $\Phi(Q(x')+V_p(v)):=\Theta(Q(x'))$, i.e.\ it is the equality of the marginals on $\K^{J^c}$.

(iii) By Lemma~\ref{lem:mono}(iv), $-M^v$ and $-M^w$ belong to $\mathcal E_N(\vartheta,\delta_-,\delta_+,C_0)$, $H(-x)=-H(x)$, and the law of $H$ under $\P_{-M}$ is the
law of $-H$ under $\P_M$. For integers $l,u\ge0$ with $l+u\le2$, the function $\bar y^{\,l}y^u$ is integrable for $\nu^v_p$ and $\nu^w_p$ (it is bounded by $1+|y|^2$),
and $\int\bar y^{\,l}y^u\,\nu^-_p(dy)=\int\overline{(-y)}^{\,l}(-y)^u\,\nu_p(dy)=(-1)^{l+u}\int\bar y^{\,l}y^u\,\nu_p(dy)$; so \eqref{eq:2m} for $(M^v,M^w)$ implies
\eqref{eq:2m} for $(-M^v,-M^w)$. The last identity is Lemma~\ref{lem:weyl}(c) with $n=N$.
\end{proof}

\noindent\emph{Remarks.} (1) The matrix $Q=H^{[p]}$ does \emph{not} satisfy (A) (one variance is removed), so (BB1) is never applied to $Q$; bounds on
$G_Q$ are derived deterministically from those on $G_H$ (Lemma~\ref{lem:RfromS}). (2) Vanishing variances are allowed throughout: if $\sigma_{p_k}=0$ then
$\nu^v_{p_k}=\nu^w_{p_k}=\delta_0$ by Lemma~\ref{lem:moments}(a), so $M_{k-1}=M_k$ and the step $k$ is trivial; no statement of this section divides by a variance.
(3) No union bound over the hybrids is ever taken: each Lindeberg step uses events of its own hybrid only, and the uniformity of all bounds in $k$ is the uniformity
of (BB1), (BB2) over the class $\fA(\Par)$, which contains every hybrid by (ii).

\begin{lemma}[$\msc$ near the edge]\label{lem:mscB}
\textup{(a)} For every $z\in\C_+$:
\[
(\msc(z)+1)^2=(2-z)\,\msc(z),\qquad\text{hence}\qquad|\msc(z)+1|\le|z-2|^{1/2}\quad\text{and}\quad\im\msc(z)\le|z-2|^{1/2}.
\]
\textup{(b)} Let $N\ge1$, $\eg\in(0,\frac1{10}]$, $\eta=N^{-2/3-\eg}$, $y\in\R$ with $|y-2|\le N^{-2/3+\eg}$, and $z=y+i\eta$. Then
\[
|z-2|\le2N^{-2/3+\eg},\qquad|\msc(z)+1|\le\sqrt2\,N^{-1/3+\eg/2},\qquad\im\msc(z)\le\sqrt2\,N^{-1/3+\eg/2}.
\]
\end{lemma}
\begin{proof}
(a) Write $m=\msc(z)$. Then $(m+1)^2=m^2+2m+1=(m^2+zm+1)+(2-z)m=(2-z)m$ by (S1). Since $|m|\le1$ (S1), $|m+1|^2=|2-z||m|\le|z-2|$.
Finally $\im m=\im(m+1)\le|m+1|$.
(b) $|z-2|\le|y-2|+\eta\le N^{-2/3+\eg}+N^{-2/3-\eg}\le2N^{-2/3+\eg}$, and by (a) $|\msc(z)+1|$ and $\im\msc(z)$ are at most $(2N^{-2/3+\eg})^{1/2}=\sqrt2N^{-1/3+\eg/2}$.
\end{proof}
\noindent This exact identity replaces the asymptotics ``$\im\msc\sim\sqrt{\kappa+\eta}$'' of the original, (6.34). The edge resolvent bounds
(Lemma~\ref{lem:edge}) follow next; they are derived from (BB1) and Lemma~\ref{lem:mscB}(b).

\begin{lemma}[edge resolvent bounds, uniformly in the ensemble]\label{lem:edge}
Let $\eg\in(0,\frac1{10}]$, $N\ge N_1(\Par,\eg)$ \textup{(Lemma~\ref{lem:N0})}, $\eta=N^{-2/3-\eg}$, $\Psi=4N^{-1/3+2\eg}$, and
\[
I_{\rm edge}:=\big[\,2-N^{-2/3+\eg},\ 2+N^{-2/3+\eg}\,\big].
\]
Define the subset of $\Omega=\K^{\mathcal P_N}$
\begin{align*}
\Omega^c_{\rm edge}&:=\Big\{x\in\Omega:\ \big|G_{H(x)}(y+i\eta)_{kl}+\delta_{kl}\big|\le\tfrac12\Psi\ \ \text{for all }k,l\in\{1,\dots,N\}\text{ and all }y\in I_{\rm edge}\Big\},\\
\Omega_{\rm edge}&:=\Omega\setminus\Omega^c_{\rm edge}.
\end{align*}
Then $\Omega^c_{\rm edge}$ is closed, $\Omega_{\rm edge}$ is open (in particular both are Borel sets; they do not depend on the ensemble), and for every $\K\in\{\R,\C\}$ and
\emph{every} ensemble $M$ of size $N$ over $\K$ with $\fA(\Par)$:
\[
\P_M(\Omega_{\rm edge})\le N^{-20}.
\]
\end{lemma}
\begin{proof}
\emph{Step 0: topology.} Fix $k,l$ and $y\in I_{\rm edge}$, and put $z:=y+i\eta$. By Lemma~\ref{lem:basic}(h) and Lemma~\ref{lem:weyl}(d), for $x,x'\in\Omega$,
$|G_{H(x)}(z)_{kl}-G_{H(x')}(z)_{kl}|\le\eta^{-2}\|H(x)-H(x')\|\le\sqrt2\,\eta^{-2}|x-x'|$. So $x\mapsto G_{H(x)}(z)_{kl}$ is continuous and
$\{x:|G_{H(x)}(z)_{kl}+\delta_{kl}|\le\frac12\Psi\}$ is closed. $\Omega^c_{\rm edge}$ is the intersection of these sets over all $k,l$ and all $y\in I_{\rm edge}$, hence
closed; its complement $\Omega_{\rm edge}$ is open, hence Borel.

\emph{Step 1: the window lies in $S_{L_0}$.} Since $N\ge1$ and $-\frac23+\eg<0$, $N^{-2/3+\eg}\le1$, so $I_{\rm edge}\subset[1,3]\subset[-5,5]$. By \textup{(N1)},
$y+i\eta\in S_{L_0}$ for every $y\in I_{\rm edge}$, and \textup{(BB1)} holds with $L=L_0=A_0\log\log N$.

\emph{Step 2: the size of the control parameter in the window.} Let $y\in I_{\rm edge}$, $z=y+i\eta$. By Lemma~\ref{lem:mscB}(b),
\begin{equation}\label{eq:B3msc}
\im\msc(z)\le\sqrt2\,N^{-1/3+\eg/2},\qquad|\msc(z)+1|\le\sqrt2\,N^{-1/3+\eg/2}.
\end{equation}
Moreover $N\eta=N^{1/3-\eg}$. Hence
\[
\sqrt{\frac{\im\msc(z)}{N\eta}}\le\Big(\sqrt2\,N^{-\frac13+\frac{\eg}2}\cdot N^{-\frac13+\eg}\Big)^{1/2}=2^{1/4}N^{-\frac13+\frac34\eg}\le2N^{-\frac13+\eg},\qquad\frac1{N\eta}=N^{-\frac13+\eg},
\]
and therefore, by \textup{(N2)} (first inequality),
\begin{equation}\label{eq:B3ctrl}
(\log N)^{4L_0}\sqrt{\frac{\im\msc(z)}{N\eta}}+\frac{(\log N)^{4L_0}}{N\eta}\le3(\log N)^{4L_0}N^{-\frac13+\eg}\le N^{\eg/2}\cdot N^{-\frac13+\eg}=N^{-\frac13+\frac32\eg}.
\end{equation}

\emph{Step 3: $\Omega_{\rm edge}\subset U_{L_0}$.} Recall from \S\ref{sec:inputs} that $U_{L_0}$ is the union over $z\in S_{L_0}$ of the events
$\{\Lambda_d(z)+\Lambda_o(z)\ge(\log N)^{4L_0}\sqrt{\im\msc(z)/(N\eta)}+(\log N)^{4L_0}/(N\eta)\}$, where $\Lambda_d=\max_k|G_{kk}-\msc|$ and $\Lambda_o=\max_{k\ne l}|G_{kl}|$ \eqref{eq:defLambda},
$G=G_{H(x)}$. Let $x\in\Omega\setminus U_{L_0}$ and $y\in I_{\rm edge}$, $z=y+i\eta$. Since $z\in S_{L_0}$ (Step 1) and $x\notin U_{L_0}$, \eqref{eq:B3ctrl} gives
\[
\Lambda_d(z)+\Lambda_o(z)<N^{-\frac13+\frac32\eg}.
\]
Both $\Lambda_d$ and $\Lambda_o$ are nonnegative, so for all $k,l$: $|G_{kl}(z)-\delta_{kl}\msc(z)|\le\max(\Lambda_d(z),\Lambda_o(z))<N^{-1/3+3\eg/2}$. With \eqref{eq:B3msc},
\begin{align*}
\big|G_{kl}(z)+\delta_{kl}\big|&\le\big|G_{kl}(z)-\delta_{kl}\msc(z)\big|+\delta_{kl}\,|\msc(z)+1|\\
&<N^{-\frac13+\frac32\eg}+\sqrt2\,N^{-\frac13+\frac12\eg}=N^{-\frac13+2\eg}\Big(N^{-\frac{\eg}2}+\sqrt2\,N^{-\frac32\eg}\Big).
\end{align*}
By \textup{(N2)}, $N^{\eg/2}\ge3(\log N)^{4L_0}\ge3$ (recall $\log N\ge1$ and $L_0>0$), so $N^{-\eg/2}+\sqrt2N^{-3\eg/2}\le\frac13+\frac{\sqrt2}{27}<2$. Thus
$|G_{kl}(z)+\delta_{kl}|<2N^{-1/3+2\eg}=\frac12\Psi$ for all $k,l$ and all $y\in I_{\rm edge}$, i.e.\ $x\in\Omega^c_{\rm edge}$. We have shown
$\Omega\setminus U_{L_0}\subset\Omega^c_{\rm edge}$, that is, $\Omega_{\rm edge}\subset U_{L_0}$.

\emph{Step 4: probability.} $\Omega_{\rm edge}$ is a Borel subset of $U_{L_0}$ (Steps 0 and 3). Let $M$ be an ensemble of size $N$ with $\fA(\Par)$. By \textup{(BB1)} with $L=L_0$,
in the form stated in \S\ref{sec:inputs} (monotonicity of the outer probability), and by \textup{(N3)},
\[
\P_M(\Omega_{\rm edge})\le\P_M^*(U_{L_0})\le Ce^{-c(\log N)^{\phi L_0}}\le N^{-20}.\qedhere
\]
\end{proof}

\noindent\emph{Remarks.} (1) \emph{Exponent check.} The three contributions to $|G_{kl}+\delta_{kl}|$ are $(\log N)^{4L_0}(N\eta)^{-1}\approx N^{-1/3+\eg}$ (the dominant one),
$(\log N)^{4L_0}(\im\msc/(N\eta))^{1/2}\approx N^{-1/3+3\eg/4}$ and $|\msc+1|\approx N^{-1/3+\eg/2}$; the factor $(\log N)^{4L_0}\le\frac13N^{\eg/2}$ is absorbed by (N2), and the
remaining room $N^{\eg/2}$ up to $\frac12\Psi=2N^{-1/3+2\eg}$ is not used. The value $-1$ on the diagonal is the edge value $\msc(2)=-1$.

(2) \emph{Uniformity over the hybrids.} The constants $A_0,C,c,\phi,N_0$ of \textup{(BB1)}, hence $N_1(\Par,\eg)$, depend only on $\Par$ (and $\eg$), and every hybrid $M_k$,
$0\le k\le\Gamma$, satisfies $\fA(\Par)$ (Lemma~\ref{lem:hyb}(ii)). So the lemma holds for each $M_k$ with the same threshold and the same bound $N^{-20}$. The set $\Omega_{\rm edge}$
is the same for all ensembles; only its probability depends on $M$. Each Lindeberg step below uses the bound for one hybrid only; no union bound over $k$ is taken.

(3) \emph{Simultaneity in $y$.} The lemma controls all $y\in I_{\rm edge}$ at once; this is the simultaneity in $z\in S_L$ of Theorem~\ref{thm:21}(ii). By the remark~(4)
after Lemma~\ref{lem:expand}, the comparison branch would also work with the bound at two fixed points $y\in\{E_1,E_2\}$; we do not use this simplification.
(4) \emph{What is not used.} No bound for a matrix with an entry or a row removed is taken from Part~I: the matrix $Q=H^{[p]}$ of the next subsections does not satisfy (A),
and its resolvent is controlled deterministically by Lemma~\ref{lem:RfromS}(a).

\subsection{The smoothed counting functional and determinant ratios}\label{sec:X}

Throughout this subsection $\eta>0$ and $E_1\le E_2$ are arbitrary real numbers, $z_i:=E_i+i\eta$ $(i=1,2)$, and everything is
deterministic: no property of $N$, of the ensemble, or of the parameters of \S\ref{sec:inputs} is used. For a complex valued function
$g$ of the spectral parameter we write, as in \S\ref{sec:inputs},
\[
\br{g}:=g(z_2)-g(z_1).
\]

\paragraph{Matrices.} For $n\ge0$ and a Hermitian matrix $A\in M_n(\C)$ (a real symmetric matrix is regarded as a Hermitian one),
$G_A(z):=(A-z)^{-1}$ exists for all $z\in\C_+$ by Lemma~\ref{lem:basic}(a), and $\lambda_1(A)\le\dots\le\lambda_n(A)$ are the
eigenvalues of $A$. The case $n=0$ (the empty matrix, with the conventions of \S\ref{sec:resid}: its determinant is $1$, its trace is $0$, and an empty
sum is $0$) is admitted only because the minor of a $1\times1$ matrix is empty; it never needs a separate argument.

\paragraph{The principal logarithm.} $\Log:\C\setminus(-\infty,0]\to\C$ is the principal branch of the logarithm,
$\Log w=\log|w|+i\operatorname{Arg}w$ with $\operatorname{Arg}w\in(-\pi,\pi)$. It is holomorphic on $\C\setminus(-\infty,0]$, with
$\frac{d}{dw}\Log w=\frac1w$ and $\exp(\Log w)=w$. The identities of Lemma~\ref{lem:ratio} and Corollary~\ref{cor:Xchange} use $\Log$ on all of $\C\setminus(-\infty,0]$
(by \eqref{eq:b1:nocut} the relevant numbers never meet the cut). The bounds \eqref{eq:b1:L1}--\eqref{eq:b1:L3} below are used only for numbers $1+\zeta$ with $|\zeta|\le\frac12$;
they lie in the half plane $\{\re w\ge\frac12\}$, which does not meet the cut $(-\infty,0]$. We shall use three bounds: for all
$\zeta\in\C$ with $|\zeta|\le\frac12$,
\begin{align}
|\Log(1+\zeta)|&\le2|\zeta|,\tag{L1}\label{eq:b1:L1}\\
\big|\Log(1+\zeta)-\zeta+\tfrac12\zeta^2\big|&\le\tfrac23|\zeta|^3,\tag{L2}\label{eq:b1:L2}\\
\big|\Log(1+\zeta)-\zeta+\tfrac12\zeta^2-\tfrac13\zeta^3\big|&\le\tfrac12|\zeta|^4.\tag{L3}\label{eq:b1:L3}
\end{align}

\begin{proof}[Proof of \eqref{eq:b1:L1}--\eqref{eq:b1:L3}]
Fix $\zeta$ with $|\zeta|\le\frac12$. For $s\in[0,1]$ we have $|1+s\zeta|\ge1-s|\zeta|\ge\frac12$ and $\re(1+s\zeta)\ge\frac12$, so
$\varphi(s):=\Log(1+s\zeta)$ is continuously differentiable on $[0,1]$ with $\varphi'(s)=\zeta/(1+s\zeta)$ (chain rule) and $\varphi(0)=0$.
By the fundamental theorem of calculus, $\Log(1+\zeta)=\int_0^1\frac{\zeta}{1+s\zeta}\,ds$. For every integer $k\ge0$ and every $w\ne-1$,
\[
\frac1{1+w}=\sum_{j=0}^{k-1}(-w)^j+\frac{(-w)^k}{1+w}
\]
(multiply by $1+w$; the sum telescopes). Put $w=s\zeta$, multiply by $\zeta$ and integrate over $s\in[0,1]$, using $\int_0^1s^j\,ds=\frac1{j+1}$:
\[
\Log(1+\zeta)-\sum_{j=1}^{k}\frac{(-1)^{j-1}\zeta^j}{j}=\int_0^1\frac{\zeta\,(-s\zeta)^k}{1+s\zeta}\,ds,\qquad
\Big|\int_0^1\frac{\zeta\,(-s\zeta)^k}{1+s\zeta}\,ds\Big|\le\int_0^12|\zeta|^{k+1}s^k\,ds=\frac{2|\zeta|^{k+1}}{k+1}.
\]
The cases $k=0$, $k=2$, $k=3$ are \eqref{eq:b1:L1}, \eqref{eq:b1:L2}, \eqref{eq:b1:L3}.
\end{proof}

Finally, we use $\pi>3$ in the form $\frac1\pi<\frac13$, and the following trivial remark: for a function $\xi$ of the spectral parameter,
\begin{equation}\label{eq:b1:br}
\big|\im\br{\xi}\big|\le|\im\xi(z_1)|+|\im\xi(z_2)|\le|\xi(z_1)|+|\xi(z_2)|.
\end{equation}

\begin{definition}\label{def:X}
For a Hermitian matrix $A\in M_n(\C)$ (any $n\ge1$),
\[
X(A)=X_{E_1,E_2,\eta}(A):=\frac1\pi\int_{E_1}^{E_2}\im\tr G_A(y+i\eta)\,dy ,
\]
and $X(A):=0$ for the empty matrix $(n=0)$. The integrand is a continuous function of $y$ (Lemma~\ref{lem:basic}(c)), so the integral is
an ordinary Riemann integral. For the random matrix, $X(H)=\frac N\pi\int_{E_1}^{E_2}\im m(y+i\eta)\,dy=\tr(\chi_{[E_1,E_2]}*\theta_\eta)(H)$
with $\theta_\eta(x)=\frac1\pi\frac{\eta}{x^2+\eta^2}$, in the notation of the original (this identification is not used).
\end{definition}

\begin{lemma}[spectral form]\label{lem:Xspec}\infra{}
Let $A\in M_n(\C)$ be Hermitian, $n\ge0$. Then
\[
X(A)=\sum_{j=1}^nf(\lambda_j(A)),\qquad f(\lambda):=\frac1\pi\Big[\arctan\frac{E_2-\lambda}\eta-\arctan\frac{E_1-\lambda}\eta\Big]\in[0,1)\quad(\lambda\in\R).
\]
Consequently $0\le X(A)\le n$. Moreover,
\begin{align*}
\textup{(a)}\quad&f(\lambda)\ge1-\frac\eta\pi\Big(\frac1{E_2-\lambda}+\frac1{\lambda-E_1}\Big)&&\text{if }E_1<\lambda<E_2,\\
\textup{(b)}\quad&f(\lambda)\le\frac\eta\pi\min\Big(\frac1{E_1-\lambda},\ \frac{E_2-E_1}{(E_1-\lambda)^2}\Big)&&\text{if }\lambda<E_1,
\end{align*}
and \textup{(c)} for Hermitian $A,A'\in M_n(\C)$: $\ |X(A)-X(A')|\le\dfrac{n\,(E_2-E_1)}{\pi\eta^2}\,\|A-A'\|$. In particular $A\mapsto X(A)$ is continuous
on the Hermitian $n\times n$ matrices.
\end{lemma}
\begin{proof}
\emph{Spectral form.} For $n=0$ both sides are $0$. Let $n\ge1$ and write $\lambda_j:=\lambda_j(A)$. By Lemma~\ref{lem:basic}(d) applied to
$A$ (the lemma holds for Hermitian matrices of any size), for every $y\in\R$,
\[
\im\tr G_A(y+i\eta)=\sum_{j=1}^n\im\frac1{\lambda_j-y-i\eta}=\sum_{j=1}^n\frac{\eta}{(\lambda_j-y)^2+\eta^2}.
\]
For fixed $\lambda\in\R$ the function $y\mapsto\arctan\frac{y-\lambda}\eta$ is differentiable on $\R$ with derivative
$\frac1\eta\big(1+\frac{(y-\lambda)^2}{\eta^2}\big)^{-1}=\frac{\eta}{(\lambda-y)^2+\eta^2}$, which is continuous. By the fundamental theorem of calculus,
\[
\int_{E_1}^{E_2}\frac{\eta\,dy}{(\lambda-y)^2+\eta^2}=\arctan\frac{E_2-\lambda}\eta-\arctan\frac{E_1-\lambda}\eta=\pi f(\lambda).
\]
Summing over $\lambda=\lambda_j$ and dividing by $\pi$ gives $X(A)=\sum_jf(\lambda_j)$. Since $\arctan$ is nondecreasing with values in
$(-\frac\pi2,\frac\pi2)$ and $E_1\le E_2$, we get $0\le\pi f(\lambda)<\pi$, i.e.\ $f(\lambda)\in[0,1)$, and therefore $0\le X(A)\le n$.

\emph{Two inequalities for $\arctan$.} For $x>0$,
\begin{equation}\label{eq:b1:atan1}
\frac\pi2-\arctan x=\arctan\frac1x\le\frac1x ,
\end{equation}
because $\arctan x+\arctan\frac1x=\frac\pi2$ for $x>0$ and $\arctan s\le s$ for $s\ge0$ (the function $s-\arctan s$ vanishes at $0$ and has the
derivative $1-\frac1{1+s^2}\ge0$). For $0<b\le a$, the mean value theorem gives a $\xi\in[b,a]$ with
\begin{equation}\label{eq:b1:atan2}
\arctan a-\arctan b=\frac{a-b}{1+\xi^2}\le\frac{a-b}{1+b^2}\le\frac{a-b}{b^2}.
\end{equation}

\emph{\textup{(a)}.} Let $E_1<\lambda<E_2$ and $x_2:=\frac{E_2-\lambda}\eta>0$, $x_1:=\frac{\lambda-E_1}\eta>0$. Since $\arctan$ is odd,
$\pi f(\lambda)=\arctan x_2+\arctan x_1\ge\pi-\frac1{x_2}-\frac1{x_1}$ by \eqref{eq:b1:atan1}; this is (a).

\emph{\textup{(b)}.} Let $\lambda<E_1$ and $b:=\frac{E_1-\lambda}\eta>0$, $a:=\frac{E_2-\lambda}\eta\ge b$. Then
$\pi f(\lambda)=\arctan a-\arctan b\le\frac\pi2-\arctan b\le\frac1b=\frac{\eta}{E_1-\lambda}$ by \eqref{eq:b1:atan1}, and
$\pi f(\lambda)\le\frac{a-b}{b^2}=\frac{\eta\,(E_2-E_1)}{(E_1-\lambda)^2}$ by \eqref{eq:b1:atan2}.

\emph{\textup{(c)}.} For $n=0$ there is nothing to prove. For $n\ge1$ and every $y\in\R$, Lemma~\ref{lem:basic}(h) (with $n$ in place of $N$) gives
$|\tr G_A(y+i\eta)-\tr G_{A'}(y+i\eta)|\le n\eta^{-2}\|A-A'\|$. Since $|\im\xi|\le|\xi|$, integrating over $y\in[E_1,E_2]$ and dividing by $\pi$ gives (c).
\end{proof}

\noindent\emph{Remark (measurability).} The map $x\mapsto H(x)$ is linear, hence continuous (Lemma~\ref{lem:weyl}(d)), and so are the maps
$x\mapsto H(x)^{[p]}$ (an entry set to zero, Definition~\ref{def:hyb}) and $x\mapsto H(x)^{(a)}$ (a row and column removed).
By (c), the functions $x\mapsto X(H(x))$, $X(H(x)^{[p]})$, $X(H(x)^{(a)})$, $X\big((H(x)^{[p]})^{(a)}\big)$ are continuous on $\Omega$, in particular Borel.

\begin{lemma}[ratio lemma]\label{lem:ratio}\infra{}
Let $A\in M_n(\C)$, $B\in M_m(\C)$ be Hermitian $(n,m\ge0)$ and $p_A(y):=\det(A-y-i\eta)$ for $y\in\R$ (so $p_A\equiv1$ if $n=0$).
\begin{itemize}\itemsep0pt
\item[\textup{(i)}] $p_A$ is a polynomial function of $y$ with complex coefficients, $p_A(y)\ne0$ for all $y\in\R$, and
\[
p_A'(y)=-\,p_A(y)\,\tr G_A(y+i\eta)\qquad(y\in\R).
\]
\item[\textup{(ii)}] Let $D:[E_1,E_2]\to\C$ satisfy $D(y)\,p_B(y)=\kappa\,p_A(y)$ for all $y\in[E_1,E_2]$, with a constant $\kappa\in\C\setminus\{0\}$, and
\begin{equation}\label{eq:b1:cut}
D(y)\notin(-\infty,0]\qquad\text{for \emph{all} }y\in[E_1,E_2].
\end{equation}
Then
\begin{equation}\label{eq:b1:ratio}
X(A)-X(B)=-\frac1\pi\im\big[\Log D(E_2)-\Log D(E_1)\big].
\end{equation}
If moreover $|D(E_i)-1|\le\frac12$ for $i=1,2$, then
\begin{equation}\label{eq:b1:ratiobd}
|X(A)-X(B)|\le\tfrac2\pi\big(|D(E_1)-1|+|D(E_2)-1|\big).
\end{equation}
Condition \eqref{eq:b1:cut} holds in particular if $|D(y)-1|\le\frac12$ for all $y\in[E_1,E_2]$.
\end{itemize}
\end{lemma}
\begin{proof}
(i) For $n=0$ the claim reads $0=0$. Let $n\ge1$. By the spectral theorem there are a unitary matrix $U$ and the real diagonal matrix
$\Lambda=\operatorname{diag}(\lambda_1(A),\dots,\lambda_n(A))$ with $A=U\Lambda U^*$ (as in the proof of Lemma~\ref{lem:basic}(d)). For $z\in\C$,
$A-z=U(\Lambda-z)U^*$ and $\det U\det U^*=\det(UU^*)=1$, so $\det(A-z)=\prod_{j=1}^n(\lambda_j(A)-z)$. Hence, with $\lambda_j:=\lambda_j(A)$,
\[
p_A(y)=\prod_{j=1}^n(\lambda_j-i\eta-y),
\]
a polynomial in $y$. Each factor has imaginary part $-\eta\ne0$, so $p_A(y)\ne0$. By the product rule (each factor has derivative $-1$),
\[
p_A'(y)=-\sum_{j=1}^n\prod_{k\ne j}(\lambda_k-i\eta-y)=-\,p_A(y)\sum_{j=1}^n\frac1{\lambda_j-y-i\eta}=-\,p_A(y)\,\tr G_A(y+i\eta),
\]
the last step by Lemma~\ref{lem:basic}(d).

(ii) Define $\widetilde D(y):=\kappa\,p_A(y)/p_B(y)$ for all $y\in\R$. By (i) this is a quotient of two polynomial functions whose denominator
has no real zero; hence $\widetilde D:\R\to\C$ is continuously differentiable, $\widetilde D(y)\ne0$, and
\begin{equation}\label{eq:b1:logder}
\frac{\widetilde D'(y)}{\widetilde D(y)}=\frac{p_A'(y)}{p_A(y)}-\frac{p_B'(y)}{p_B(y)}=\tr G_B(y+i\eta)-\tr G_A(y+i\eta)\qquad(y\in\R).
\end{equation}
By hypothesis $\widetilde D=D$ on $[E_1,E_2]$. If $E_1=E_2$ both sides of the claimed identity vanish; let $E_1<E_2$.
For every $y\in[E_1,E_2]$ the number $\widetilde D(y)=D(y)$ lies in $\C\setminus(-\infty,0]$ by \eqref{eq:b1:cut}, and $\Log$ is
complex differentiable at each point of $\C\setminus(-\infty,0]$ with derivative $1/w$. By the chain rule the function
$\ell(y):=\Log\widetilde D(y)$ is differentiable at every $y\in[E_1,E_2]$, with
\[
\ell'(y)=\frac{\widetilde D'(y)}{\widetilde D(y)}=\tr G_B(y+i\eta)-\tr G_A(y+i\eta),
\]
and $\ell'$ is continuous on $[E_1,E_2]$. By the fundamental theorem of calculus (applied to the real and imaginary parts of $\ell$),
\[
\Log D(E_2)-\Log D(E_1)=\ell(E_2)-\ell(E_1)=\int_{E_1}^{E_2}\big(\tr G_B(y+i\eta)-\tr G_A(y+i\eta)\big)\,dy .
\]
Taking $\frac1\pi\im$ of both sides and using Definition~\ref{def:X} (also for $n=0$ or $m=0$, where the corresponding trace is $0$) gives
$\frac1\pi\im\big[\Log D(E_2)-\Log D(E_1)\big]=X(B)-X(A)$, which is \eqref{eq:b1:ratio}. For \eqref{eq:b1:ratiobd} apply \eqref{eq:b1:L1} with
$\zeta=D(E_i)-1$: $|\im\Log D(E_i)|\le|\Log D(E_i)|\le2|D(E_i)-1|$ for $i=1,2$. Finally, the closed disc $\{|w-1|\le\frac12\}$ lies in the half
plane $\{\re w\ge\frac12\}$, which gives the last assertion.
\end{proof}

\noindent\emph{Remarks.} (1) \emph{The branch; why a condition on the whole segment.} The identity \eqref{eq:b1:ratio} uses that the curve
$y\mapsto D(y)$, $y\in[E_1,E_2]$, does not meet the cut $(-\infty,0]$ of the principal branch; a condition at the endpoints alone is not
sufficient in general. If it is known only that $D(E_1),D(E_2)\notin(-\infty,0]$, then by \eqref{eq:b1:logder} the function
$y\mapsto\widetilde D(y)\exp\big(-\!\int_{E_1}^y(\tr G_B-\tr G_A)(s+i\eta)\,ds\big)$ has derivative $0$, hence
$\exp\big(\Log D(E_2)-\Log D(E_1)\big)=\exp\big(\int_{E_1}^{E_2}(\tr G_B-\tr G_A)(y+i\eta)\,dy\big)$, and one can only conclude that
\[
X(A)-X(B)+\tfrac1\pi\im\big[\Log D(E_2)-\Log D(E_1)\big]\in2\mathbb Z .
\]
The even integer can be nonzero (numerically it is $\pm2$ for about a third of random $6\times6$ pairs $A=B+{}$a positive matrix of rank $3$).
(2) The hypothesis $|D-1|\le\frac12$ on the whole segment is a special case of \eqref{eq:b1:cut}; see the last sentence of (ii).
In Lemma~\ref{lem:det} we show that for the two ratios that occur in Part~II (replacement of one entry, removal of one row and column)
condition \eqref{eq:b1:cut} holds \emph{automatically}, at every point of $\C_+$. Thus no smallness condition along the segment is needed for the
exact formulas of Corollary~\ref{cor:Xchange}; smallness is needed only at the two endpoints $z_1,z_2$, where the logarithm is expanded.
(3) The function $D$ is not assumed to be continuous: it is determined by $D=\kappa p_A/p_B$.

\begin{lemma}[one entry: $2\times2$ determinant]\label{lem:det}
Let $Q\in M_N(\C)$ be Hermitian, $R(z):=G_Q(z)$, $p=(a,b)$ with $1\le a\le b\le N$, and $v\in\K$ ($v\in\R$ if $a=b$). Let $V_p(v)$ be the matrix of
Definition~\ref{def:hyb}: $V_p(v)=vE^{(ab)}+\bar vE^{(ba)}$ if $a\ne b$ (its only nonzero entries are $v$ at $(a,b)$ and $\bar v$ at $(b,a)$), and
$V_p(v)=vE^{(aa)}$ if $a=b$. For all $z\in\C_+$,
\[
\det(Q+V_p(v)-z)=\det(Q-z)\cdot D_p(v;z)
\]
with
\[
D_p(v;z)=\begin{cases}1+u(v;z)-w(v;z),& a\ne b,\\ 1+vR_{aa}(z),& a=b,\end{cases}\qquad
\begin{aligned}u(v;z)&:=vR_{ba}(z)+\bar vR_{ab}(z),\\ w(v;z)&:=|v|^2g(z),\quad g:=R_{aa}R_{bb}-R_{ab}R_{ba}.\end{aligned}
\]
Moreover, the principal logarithm can be applied to $D_p$ and to $-R_{aa}$ at every point of $\C_+$:
\begin{equation}\label{eq:b1:nocut}
D_p(v;z)\in\C\setminus(-\infty,0]\qquad\text{and}\qquad-R_{aa}(z)\in\C\setminus(-\infty,0]\qquad\text{for all }z\in\C_+ .
\end{equation}
\textup{(The statement and its proof hold for every $v\in\C$ if $a\ne b$; for $\K=\R$ the matrix $Q$ is real symmetric and $v$ is real.)}
\end{lemma}
\begin{proof}
Write $V:=V_p(v)$; it is Hermitian, so $Q+V$ is Hermitian. Since $(Q-z)R=1$, we have $Q+V-z=(Q-z)(1+RV)$ and therefore
\begin{equation}\label{eq:b1:det1}
\det(Q+V-z)=\det(Q-z)\cdot\det(1_N+RV).
\end{equation}
\emph{Reduction to the block.} Let $m:=2$, $(i_1,i_2):=(a,b)$ if $a\ne b$, and $m:=1$, $i_1:=a$ if $a=b$. Let $U$ be the $N\times m$ matrix whose columns are the
standard basis vectors $e_{i_1},\dots,e_{i_m}$, and
\[
C:=\begin{pmatrix}0&v\\ \bar v&0\end{pmatrix}\ \ (a\ne b),\qquad C:=(v)\ \ (a=b).
\]
Then $UCU^*=\sum_{k,l=1}^mC_{kl}\,e_{i_k}e_{i_l}^*=V$, and $U^*RU=(R_{i_ki_l})_{k,l=1}^m=:R_P$ is the block of $R$ with rows and columns in $P:=\{a,b\}$.
We use the identity
\begin{equation}\label{eq:b1:sylv}
\det(1_N+\mathsf X\mathsf Y)=\det(1_m+\mathsf Y\mathsf X)\qquad\text{for all }\mathsf X\in M_{N\times m}(\C),\ \mathsf Y\in M_{m\times N}(\C)
\end{equation}
(Mathlib: \texttt{Matrix.det\_one\_add\_mul\_comm}). For completeness: apply the block determinant formula \textup{(BD)} from the proof of
Lemma~\ref{lem:schur} to the $(N+m)\times(N+m)$ matrix $\mathcal M:=\big(\begin{smallmatrix}1_N&-\mathsf X\\ \mathsf Y&1_m\end{smallmatrix}\big)$. With the
lower right block $1_m$ as pivot it gives $\det\mathcal M=\det(1_m)\det(1_N+\mathsf X\mathsf Y)$; with the upper left block $1_N$ as pivot (which is
\textup{(BD)} after the simultaneous relabelling of rows and columns that exchanges the two blocks) it gives $\det\mathcal M=\det(1_N)\det(1_m+\mathsf Y\mathsf X)$.

By \eqref{eq:b1:sylv} with $\mathsf X=RU$ and $\mathsf Y=CU^*$,
\begin{equation}\label{eq:b1:det2}
\det(1_N+RV)=\det\big(1_N+(RU)(CU^*)\big)=\det\big(1_m+CU^*RU\big)=\det(1_m+CR_P).
\end{equation}
\emph{The small determinant.} If $a=b$ then $\det(1_1+CR_P)=1+vR_{aa}$. If $a\ne b$ then
\[
CR_P=\begin{pmatrix}0&v\\ \bar v&0\end{pmatrix}\begin{pmatrix}R_{aa}&R_{ab}\\ R_{ba}&R_{bb}\end{pmatrix}
=\begin{pmatrix}vR_{ba}&vR_{bb}\\ \bar vR_{aa}&\bar vR_{ab}\end{pmatrix},
\]
\begin{align*}
\det(1_2+CR_P)&=(1+vR_{ba})(1+\bar vR_{ab})-|v|^2R_{aa}R_{bb}\\
&=1+vR_{ba}+\bar vR_{ab}-|v|^2\big(R_{aa}R_{bb}-R_{ab}R_{ba}\big)=1+u-w .
\end{align*}
Together with \eqref{eq:b1:det1} and \eqref{eq:b1:det2} this proves the determinant identity.

\emph{Proof of \eqref{eq:b1:nocut}.} Fix $z=y+i\eta'\in\C_+$ ($\eta'>0$) and write $R=R(z)$. By Lemma~\ref{lem:basic}(a), $\im R_{aa}>0$. Hence
$\im(-R_{aa})<0$, so $-R_{aa}$ is not real. If $a=b$, then $\im D_p(v;z)=v\,\im R_{aa}\ne0$ for $v\ne0$, and $D_p(0;z)=1$. Let now $a\ne b$.
If $v=0$ then $D_p=1$; let $v\ne0$. Since $C^2=|v|^2\,1_2$, the matrix $C$ is invertible with $C^{-1}=|v|^{-2}C$ and $\det C=-|v|^2$, so
\begin{equation}\label{eq:b1:DM}
D_p(v;z)=\det(1_2+CR_P)=\det C\cdot\det(C^{-1}+R_P)=-|v|^2\det\mathsf M,\qquad\mathsf M:=|v|^{-2}C+R_P .
\end{equation}
By the identity $R-R^*=2i\eta'RR^*$ (proof of Lemma~\ref{lem:basic}(b)) and $C=C^*$,
\[
\frac1{2i}(\mathsf M-\mathsf M^*)=\frac1{2i}\,U^*(R-R^*)U=\eta'\,U^*RR^*U=:\mathsf T,\qquad
\langle x,\mathsf Tx\rangle=\eta'\,|R^*Ux|^2>0\quad\text{for }x\in\C^2\setminus\{0\},
\]
because $Ux\ne0$ and $R^*$ is invertible. Let $\mu\in\C$ be an eigenvalue of $\mathsf M$ and $x$ a corresponding eigenvector with $|x|=1$. Then
$\mu=\langle x,\mathsf Mx\rangle$ and $\im\mu=\frac1{2i}\big(\langle x,\mathsf Mx\rangle-\overline{\langle x,\mathsf Mx\rangle}\big)=\langle x,\mathsf Tx\rangle>0$.
Thus both roots $\mu_1,\mu_2$ of the characteristic polynomial of the $2\times2$ matrix $\mathsf M$ lie in $\C_+$, and $\det\mathsf M=\mu_1\mu_2$.
Suppose that $\det\mathsf M=r$ with a real $r\ge0$. Then $r\ne0$ and $\mu_2=r/\mu_1=r\,\overline{\mu_1}/|\mu_1|^2$ has negative imaginary part, a contradiction.
Hence $\det\mathsf M\notin[0,\infty)$, and by \eqref{eq:b1:DM} $D_p(v;z)\notin(-\infty,0]$.
\end{proof}

\begin{corollary}[exact formulas for the change of $X$]\label{cor:Xchange}
Let $Q$, $p=(a,b)$, $v$ be as in Lemma~\ref{lem:det}.
\begin{itemize}\itemsep0pt
\item[\textup{(i)}] \emph{(one entry)} Without any further hypothesis,
\[
X(Q+V_p(v))-X(Q)=-\frac1\pi\im\br{\Log D_p(v;\cdot)}.
\]
\item[\textup{(ii)}] \emph{(one row and column)} Let $Q^{(a)}$ be the minor of $Q$ with row and column $a$ removed and $X(Q^{(a)})=X_{E_1,E_2,\eta}(Q^{(a)})$ the
functional of Definition~\ref{def:X} for this $(N-1)\times(N-1)$ matrix. Without any further hypothesis,
\[
X(Q)-X(Q^{(a)})=\frac1\pi\im\br{\Log(-R_{aa})} ,
\]
and if $|R_{aa}(z_i)+1|\le\frac12$ for $i=1,2$, then $\ |X(Q)-X(Q^{(a)})|\le\tfrac2\pi\big(|R_{aa}(z_1)+1|+|R_{aa}(z_2)+1|\big)$.
\end{itemize}
In particular \textup{(i)} holds if $|D_p(v;y+i\eta)-1|\le\frac12$ for all $y\in[E_1,E_2]$, and \textup{(ii)} holds, with the bound, if
$|R_{aa}(y+i\eta)+1|\le\frac12$ for all $y\in[E_1,E_2]$.
\end{corollary}
\begin{proof}
(i) Apply Lemma~\ref{lem:ratio}(ii) to $A:=Q+V_p(v)$, $B:=Q$ (both Hermitian, $n=m=N$), $D(y):=D_p(v;y+i\eta)$ and $\kappa:=1$: the relation
$D(y)\,p_Q(y)=p_{Q+V_p(v)}(y)$ is Lemma~\ref{lem:det} at $z=y+i\eta$, and condition \eqref{eq:b1:cut} is \eqref{eq:b1:nocut}. Since
$D(E_i)=D_p(v;z_i)$, \eqref{eq:b1:ratio} is the claim.

(ii) By Lemma~\ref{lem:schur}\,(I6) for the matrix $Q$ and the index $a$ (recall $R_{aa}\ne0$), for every $z\in\C_+$,
$\det(Q-z)=\det(Q^{(a)}-z)/R_{aa}(z)$, that is,
\[
\big(-R_{aa}(y+i\eta)\big)\,p_Q(y)=-\,p_{Q^{(a)}}(y)\qquad(y\in\R).
\]
Apply Lemma~\ref{lem:ratio}(ii) to $A:=Q^{(a)}$ $(n=N-1)$, $B:=Q$ $(m=N)$, $D(y):=-R_{aa}(y+i\eta)$ and $\kappa:=-1$; condition \eqref{eq:b1:cut} is
again \eqref{eq:b1:nocut}. Then \eqref{eq:b1:ratio} reads $X(Q^{(a)})-X(Q)=-\frac1\pi\im\br{\Log(-R_{aa})}$, and \eqref{eq:b1:ratiobd} with
$|D(E_i)-1|=|R_{aa}(z_i)+1|$ is the bound.
\end{proof}

\noindent\emph{Remarks.} (1) By \eqref{eq:b1:nocut} the exact formulas need no information on $R$ along the segment $[E_1,E_2]+i\eta$: the resolvent
enters only through the entries $R_{kl}(z_i)$, $k,l\in\{a,b\}$, at the two endpoints $z_1,z_2$. (Smallness conditions on the whole
segment are therefore not needed to stay on the principal branch; nevertheless the good sets of \S\ref{sec:step} are defined with conditions on the whole segment.
See the remark after Lemma~\ref{lem:expand} for what this means for the use of Theorem~\ref{thm:21}.)
(2) Both identities were checked numerically for real symmetric and complex Hermitian matrices, including large $v$, small $\eta$ and windows
inside the spectrum, where $|D_p-1|$ is not small.

\subsection{From the law of $\lambda_N$ to the functional, via rigidity}\label{sec:sandwich}
This subsection replaces Lemma~6.1 and Corollary~6.2 of the original. It uses (BB2) only (not (2.19), not (2.26) of the original), and no
averaging over the window edge is needed, because only the event $\{\lambda_N<E\}$ is considered.

\begin{lemma}[classical locations at the edge]\label{lem:gamma}
For $N\ge1$ and $1\le j\le N$:\quad $2-\gamma_j\ge\Big(\dfrac{3\pi(N-j)}{2N}\Big)^{2/3}$.
\end{lemma}
\begin{proof}
By Lemma~\ref{lem:nsc}(c), $\gamma_j\in(-2,2]$ and $n_{sc}(\gamma_j)=j/N$; by Lemma~\ref{lem:nsc}(a), $n_{sc}(2)=1$ and $n_{sc}'=\varrho_{sc}$. Put $r:=2-\gamma_j\in[0,4)$. Since
$\varrho_{sc}$ is even, Lemma~\ref{lem:nsc}(b) (upper bound, valid for all $u\ge0$) gives $\varrho_{sc}(2-u)=\varrho_{sc}(-2+u)\le\sqrt u/\pi$ for $u\ge0$. Hence,
substituting $x=2-u$,
\[
\frac{N-j}N=n_{sc}(2)-n_{sc}(\gamma_j)=\int_{\gamma_j}^2\varrho_{sc}(x)\,dx=\int_0^{r}\varrho_{sc}(2-u)\,du\le\frac1\pi\int_0^r\sqrt u\,du=\frac2{3\pi}\,r^{3/2}.
\]
Thus $r^{3/2}\ge\frac{3\pi(N-j)}{2N}\ge0$, and since $s\mapsto s^{2/3}$ is nondecreasing on $[0,\infty)$, $r\ge\big(\frac{3\pi(N-j)}{2N}\big)^{2/3}$.
\end{proof}

\begin{proposition}[sandwich]\label{prop:sandwich}
Let $\eg\in(0,\frac1{10}]$, $N\ge N_1(\Par,\eg)$ (Lemma~\ref{lem:N0}), and let $L_0,W,E_L,\eta,\e=\eg/2,\ell$ be as in \eqref{eq:B2par}, \eqref{eq:B2par2}.
For $E\in[2-2W,\,2+W]$ put
\[
X^\pm_E:=X_{E\pm\ell,\,E_L,\,\eta}\qquad\text{(Definition~\ref{def:X}; note $E-\ell<E+\ell<E_L$).}
\]
Let $q:\R\to[0,1]$ be a Borel function with $q=1$ on $(-\infty,\frac19]$ and $q=0$ on $[\frac29,\infty)$. Then for every $\K\in\{\R,\C\}$ and every ensemble $M$ of size
$N$ over $\K$ with $\fA(\Par)$, the functions $x\mapsto q(X^\pm_E(H(x)))$ are Borel, the set $\{x:\lambda_N(H(x))<E\}$ is open, and
\[
\E_M\,q\big(X^-_E(H)\big)-N^{-20}\ \le\ \P_M(\lambda_N<E)\ \le\ \E_M\,q\big(X^+_E(H)\big)+N^{-20}.
\]
\end{proposition}
\begin{proof}
\emph{Step 0: preliminaries and the rigidity event.} By (N5), $0<2\ell\le W$, so $\ell<W$, and for $E\in[2-2W,2+W]$:
$E-\ell<E+\ell\le2+W+\ell<2+2W=E_L$. Hence both windows $[E\pm\ell,E_L]$ are nondegenerate intervals and $X^\pm_E$ are defined. By Lemma~\ref{lem:Xspec}(c) and the
measurability remark after it, $x\mapsto X^\pm_E(H(x))$ is continuous on $\Omega$, so $x\mapsto q(X^\pm_E(H(x)))$ is Borel, with values in $[0,1]$.
By Lemma~\ref{lem:weyl}(d), $x\mapsto\lambda_j(H(x))$ is continuous for every $j$, so $\{\lambda_N<E\}$ is open. Let
\[
\cR:=\Big\{x\in\Omega:\ |\lambda_j-\gamma_j|<(\log N)^{L_0}\big[\min(j,N-j+1)\big]^{-1/3}N^{-2/3}\ \text{ for all }j\in\{1,\dots,N\}\Big\}.
\]
$\cR$ is open (a finite intersection of open sets), and $\Omega\setminus\cR$ is the event of (BB2) with $L=L_0$, which is admissible by (N1). By (BB2) and (N3),
\begin{equation}\label{eq:B2Rc}
\P_M(\Omega\setminus\cR)\le Ce^{-c(\log N)^{\phi L_0}}\le N^{-20}.
\end{equation}
On $\cR$ we have, since $\gamma_N=2$, $\min(N,1)=1$ and $\min(j,N-j+1)\ge1$:
\begin{equation}\label{eq:B2R12}
\text{(R1)}\quad|\lambda_N-2|<(\log N)^{L_0}N^{-2/3}=W,\qquad\qquad\text{(R2)}\quad\lambda_j<\gamma_j+W\quad\text{for all }j.
\end{equation}
Recall from Lemma~\ref{lem:Xspec} that $X_{E_1,E_2,\eta}(A)=\sum_jf(\lambda_j(A))$ with $f=f_{E_1,E_2,\eta}\ge0$.

\emph{Step 1: on $\cR\cap\{\lambda_N\ge E\}$ we have $X^-_E(H)>\frac29$.} Take $E_1=E-\ell$, $E_2=E_L$. Then $\lambda_N-E_1\ge E-E_1=\ell>0$, and by (R1)
$E_2-\lambda_N>2+2W-(2+W)=W\ge2\ell>\ell$. So $E_1<\lambda_N<E_2$, and Lemma~\ref{lem:Xspec}(a), \eqref{eq:B2etaell} and (N5) give
\[
f(\lambda_N)\ge1-\frac\eta\pi\Big(\frac1{E_2-\lambda_N}+\frac1{\lambda_N-E_1}\Big)\ge1-\frac{2\eta}{\pi\ell}=1-\frac6\pi N^{-\e}\ge\frac89 .
\]
Since all terms of $\sum_jf(\lambda_j)$ are nonnegative, $X^-_E(H)\ge f(\lambda_N)\ge\frac89>\frac29$.

\emph{Step 2: on $\cR\cap\{\lambda_N<E\}$ we have $X^+_E(H)\le\frac19$.} Take $E_1=E+\ell$, $E_2=E_L$. Then
$E_2-E_1=2+2W-E-\ell\le2+2W-(2-2W)=4W$. For every $j$, $\lambda_j\le\lambda_N<E$, so $d_j:=E_1-\lambda_j>\ell>0$, and Lemma~\ref{lem:Xspec}(b) applies:
\begin{equation}\label{eq:B2fb}
f(\lambda_j)\le\frac\eta\pi\min\Big(\frac1{d_j},\ \frac{E_2-E_1}{d_j^2}\Big)\le\frac\eta\pi\min\Big(\frac1{d_j},\ \frac{4W}{d_j^2}\Big).
\end{equation}
Index the eigenvalues from the top: $j=N-k+1$, $k\in\{1,\dots,N\}$, and let $K_1:=\lfloor4(\log N)^{3L_0/2}\rfloor+1$, so $4(\log N)^{3L_0/2}<K_1\le4(\log N)^{3L_0/2}+1$.

\emph{The top eigenvalues, $k\le K_1$.} By \eqref{eq:B2fb}, $d_j>\ell$ and \eqref{eq:B2etaell}: $f(\lambda_j)\le\frac{\eta}{\pi d_j}<\frac\eta{\pi\ell}=\frac3\pi N^{-\e}$.
There are at most $K_1$ such indices, so their total contribution is at most $\frac3\pi K_1N^{-\e}$.

\emph{The other eigenvalues, $k\ge K_1+1$.} Put $a_k:=\big(\frac{3\pi(k-1)}{2N}\big)^{2/3}$. By Lemma~\ref{lem:gamma} (note $N-j=k-1$), $2-\gamma_j\ge a_k$. By (R2) and
$E\ge2-2W$, dropping $\ell>0$:
\[
d_j=E+\ell-\lambda_j>E-\gamma_j-W\ge(2-\gamma_j)-3W\ge a_k-3W.
\]
Since $k-1\ge K_1>4(\log N)^{3L_0/2}$, we have $(k-1)^{2/3}>4^{2/3}(\log N)^{L_0}$, hence
$a_k>(3\pi/2)^{2/3}4^{2/3}(\log N)^{L_0}N^{-2/3}=(6\pi)^{2/3}W\ge6W$, because $(6\pi)^2\ge6^3$ (as $\pi\ge\sqrt6$). Therefore $3W<\frac12a_k$ and $d_j>\frac12a_k>0$.
By \eqref{eq:B2fb},
\[
f(\lambda_j)\le\frac\eta\pi\cdot\frac{4W}{d_j^2}\le\frac{16\,\eta W}{\pi\,a_k^2}=\kappa_0\,\eta WN^{4/3}(k-1)^{-4/3},\qquad
\kappa_0:=\frac{16}\pi\Big(\frac2{3\pi}\Big)^{4/3}.
\]
We have $\kappa_0\le1$: indeed $\kappa_0^3=\frac{16^3}{\pi^3}\cdot\frac{2^4}{3^4\pi^4}=\frac{65536}{81\,\pi^7}<\frac{65536}{81\cdot3^7}=\frac{65536}{177147}<1$, using $\pi>3$.
Moreover, for every integer $m_0\ge1$, $\sum_{m=1}^{m_0}m^{-4/3}\le1+\sum_{m=2}^{m_0}\int_{m-1}^mx^{-4/3}dx=1+3\big(1-m_0^{-1/3}\big)\le4$, since
$m^{-4/3}\le x^{-4/3}$ for $x\in[m-1,m]$. With $m=k-1$ ranging over a subset of $\{1,\dots,N-1\}$, the total contribution of the indices $k\ge K_1+1$ is at most
$4\eta WN^{4/3}$, and by \eqref{eq:B2par}, \eqref{eq:B2par2},
\[
\eta WN^{4/3}=N^{-2/3-\eg}\cdot(\log N)^{L_0}N^{-2/3}\cdot N^{4/3}=N^{-2\e}(\log N)^{L_0}.
\]

\emph{Conclusion of Step 2.} Adding the two groups and using (N5),
\[
X^+_E(H)\le\frac3\pi K_1N^{-\e}+4N^{-2\e}(\log N)^{L_0}\le\frac3\pi\big(4(\log N)^{3L_0/2}+1\big)N^{-\e}+4N^{-2\e}(\log N)^{L_0}\le\frac19 .
\]

\emph{Step 3: integration.} We claim that at every point of $\Omega$,
\begin{equation}\label{eq:B2pw}
q\big(X^-_E(H)\big)\le\mathbf 1(\lambda_N<E)+\mathbf 1_{\Omega\setminus\cR},\qquad
\mathbf 1(\lambda_N<E)\le q\big(X^+_E(H)\big)+\mathbf 1_{\Omega\setminus\cR}.
\end{equation}
For the first: if $x\in\cR$ and $\lambda_N\ge E$, then $X^-_E(H)>\frac29$ by Step 1, so $q(X^-_E(H))=0$; in all other cases the right side is $\ge1\ge q$.
For the second: if $x\in\cR$ and $\lambda_N<E$, then $X^+_E(H)\le\frac19$ by Step 2, so $q(X^+_E(H))=1$; if $\lambda_N\ge E$ the left side is $0\le$ the right side
(as $q\ge0$); if $x\notin\cR$ the right side is $\ge1$. All functions in \eqref{eq:B2pw} are bounded and Borel (Step 0). Taking $\E_M$ and using \eqref{eq:B2Rc}
gives the two inequalities of the proposition.
\end{proof}

\noindent\emph{Why rigidity at the scale $W\gg\ell$ suffices.} The input (BB2) locates $\lambda_N$ only up to $W=(\log N)^{L_0}N^{-2/3}$, which is much larger than
the shift $\ell=\frac13N^{-2/3-\e}$ of the window; the proposition never asks (BB2) to locate the top eigenvalues at the scale $\ell$. In Step~1 only the single
eigenvalue $\lambda_N$ is used, and its distance $\ge\ell$ to the left end $E-\ell$ of the window comes from the event $\{\lambda_N\ge E\}$ itself. In Step~2 the event
$\{\lambda_N<E\}$ puts \emph{every} eigenvalue at distance $>\ell$ to the left of $E+\ell$; hence each of the top $K_1\approx4(\log N)^{3L_0/2}$ eigenvalues contributes at most
$\eta/(\pi\ell)=\frac3\pi N^{-\e}$, wherever it is, and their number is only polylogarithmic. Only the remaining eigenvalues are placed by \eqref{eq:rigidity} ((2.25) of the original; bound (R2)); they lie
at distance $\ge\frac12a_k\gg W$ below $E$, where an error $W$ in their location is harmless, and they contribute at most $4N^{-2\e}(\log N)^{L_0}$ in total.
In particular no averaging over the window edge, no bound on the counting function ((2.19), (2.26) of the original) and no monotonicity of $q$ are needed.


\subsection{One-entry replacement: exact expansion of the functional}\label{sec:step}
\textbf{Standing assumptions for \S\ref{sec:step}--\ref{sec:gfc}.} $\eg\in(0,\frac1{10}]$, $N\ge N_1(\Par,\eg)$,
$\eta=N^{-2/3-\eg}$, $\Psi=4N^{-1/3+2\eg}$, $t=N^{-1/2+2\eg}$; $E_1\le E_2$ with $|E_i-2|\le N^{-2/3+\eg}$; $z_i=E_i+i\eta$;
$X=X_{E_1,E_2,\eta}$; $F\in C^3(\R)$ with $C_F:=\max_{0\le k\le3}\sup_\R|F^{(k)}|<\infty$. $c_1,c_2,\dots$ are absolute
numerical constants; $\mathfrak m_k=\mathfrak m_k(\vartheta)$ are the moment constants of (S5) in the explicit form \eqref{eq:b3:M}.
The probabilistic statements of \S\ref{sec:step}--\ref{sec:gfc} (Lemma~\ref{lem:RfromS}(b), Corollary~\ref{cor:entry}, Lemma~\ref{lem:decor}, Proposition~\ref{prop:step}, Theorem~\ref{thm:gfc}) hold
under these standing assumptions, which are repeated in each of them.

\medskip\noindent\textbf{What the deterministic statements use.} Definition~\ref{def:goodB}, Lemma~\ref{lem:RfromS}(a) and Lemma~\ref{lem:expand} are
deterministic, and they use the standing assumptions only through the following three consequences:
\begin{itemize}\itemsep0pt
\item[\textup{(D1)}] $\eta>0$; $E_1\le E_2$ are real; $z_i=E_i+i\eta$; $X=X_{E_1,E_2,\eta}$ (Definition~\ref{def:X});
\item[\textup{(D2)}] $0<\Psi\le\frac18$ and $0\le t\le\Psi/40$ \quad(this is \textup{(N4)} of Lemma~\ref{lem:N0});
\item[\textup{(D3)}] $F\in C^3(\R)$ and $\sup_\R|F'|,\ \sup_\R|F''|,\ \sup_\R|F'''|\le C_F$.
\end{itemize}
In particular neither $N$, nor the position of $E_1,E_2$ relative to the edge, nor the size of $\eta$ plays any role in these statements, and
they hold for every Hermitian matrix in $M_N(\C)$, $N\ge1$. All numerical constants below are derived using only $\frac1\pi<\frac13$ and
$\Psi\le\frac18$; they were also checked by exact rational arithmetic and by random sampling.

\begin{definition}[good set for the pair $p=(a,b)$]\label{def:goodB}
$\cG_p:=\big\{Q\in M_N(\C)\text{ Hermitian}:\ |G_Q(y+i\eta)_{kl}+\delta_{kl}|\le\Psi\ \text{for all }k,l\in\{a,b\},\ y\in[E_1,E_2]\big\}$ (a closed set).
\end{definition}
\noindent\emph{Closedness.} For fixed $y$ and $k,l$ the function $Q\mapsto G_Q(y+i\eta)_{kl}$ is continuous on the Hermitian matrices, because
$|G_Q(z)_{kl}-G_{Q'}(z)_{kl}|\le\eta^{-2}\|Q-Q'\|$ by Lemma~\ref{lem:basic}(h); so each of the conditions defines a closed set, and $\cG_p$ is their
intersection. Since $x\mapsto H(x)^{[p]}$ is linear, hence continuous, $\{x\in\Omega:\ H(x)^{[p]}\in\cG_p\}$ is a closed, in particular Borel, subset of $\Omega$.
Every $Q\in\cG_p$ satisfies in particular the \emph{endpoint condition}
\begin{equation}\label{eq:b1:end}
|G_Q(z_i)_{kl}+\delta_{kl}|\le\Psi\qquad\text{for all }k,l\in\{a,b\}\text{ and }i\in\{1,2\},
\end{equation}
and \eqref{eq:b1:end} is all that Lemma~\ref{lem:expand} uses.

\begin{lemma}[bounds for the matrix with one entry removed]\label{lem:RfromS}
\textup{(a) (deterministic)} Let $H\in M_N(\C)$ be Hermitian, $p=(a,b)$, $Q=H^{[p]}$. If $|h_{ab}|\le t$ and
$|G_H(y+i\eta)_{kl}+\delta_{kl}|\le\frac12\Psi$ for all $k,l\in\{a,b\}$, $y\in[E_1,E_2]$, then $Q\in\cG_p$.
More precisely, for every single $y\in\R$: if $|h_{ab}|\le t$ and $|G_H(y+i\eta)_{kl}+\delta_{kl}|\le\frac12\Psi$ for all $k,l\in\{a,b\}$, then for all $k,l\in\{a,b\}$
\begin{equation}\label{eq:b1:RS}
\big|G_Q(y+i\eta)_{kl}-G_H(y+i\eta)_{kl}\big|\le\tfrac43|h_{ab}|\le\tfrac1{30}\Psi\qquad\text{and}\qquad\big|G_Q(y+i\eta)_{kl}+\delta_{kl}\big|\le\Psi .
\end{equation}
\textup{(b)} Under the standing assumptions of \S\ref{sec:step} \textup{(}in particular $\eg\in(0,\frac1{10}]$, $N\ge N_1(\Par,\eg)$, $\eta=N^{-2/3-\eg}$, $|E_i-2|\le N^{-2/3+\eg}$\textup{)}:
for every ensemble $M$ of size $N$ with $\fA(\Par)$ and every pair $p$, $\ \P_M\big(H^{[p]}\notin\cG_p\big)\le2N^{-20}$.
\end{lemma}
\begin{proof}[Proof of \textup{(a)}]
It suffices to prove the statement for a single $y$. Fix $y\in\R$, $z:=y+i\eta$, and write $S:=G_H(z)$, $R:=G_Q(z)$, $h:=h_{ab}$ and
$V:=V_p(h)=H-Q$ (Definition~\ref{def:hyb}). By Lemma~\ref{lem:basic}(h) with $(A,B)=(H,Q)$, $S-R=S\,(Q-H)\,R=-SVR$, i.e.
\[
R=S+SVR\qquad\text{(exactly)}.
\]
Let $P:=\{a,b\}$ and $m:=|P|\in\{1,2\}$. Since $V_{ij}=0$ unless $i,j\in P$, we get for $k,l\in P$: $R_{kl}=S_{kl}+\sum_{i,j\in P}S_{ki}V_{ij}R_{jl}$.
In terms of the $P\times P$ blocks $R_P,S_P$ and $C:=V_P$ (so $C=\big(\begin{smallmatrix}0&h\\ \bar h&0\end{smallmatrix}\big)$ if $a\ne b$, and $C=(h)$ if $a=b$)
this is the \emph{closed} $m\times m$ identity
\begin{equation}\label{eq:b1:closed}
R_P=S_P+S_PC\,R_P .
\end{equation}
For an $m\times m$ matrix $\mathsf A$ let $\|\mathsf A\|_\infty:=\max_k\sum_l|\mathsf A_{kl}|$ (maximal row sum). Then $|\mathsf A_{kl}|\le\|\mathsf A\|_\infty$ and
$\|\mathsf A\mathsf B\|_\infty\le\|\mathsf A\|_\infty\|\mathsf B\|_\infty$, since $\sum_l|\sum_j\mathsf A_{kj}\mathsf B_{jl}|\le\sum_j|\mathsf A_{kj}|\sum_l|\mathsf B_{jl}|$.
By hypothesis $|S_{kk}|\le1+\frac12\Psi$ and $|S_{kl}|\le\frac12\Psi$ for $k\ne l$ in $P$, so $\|S_P\|_\infty\le1+\Psi$; and $\|C\|_\infty=|h|$, because each row
of $C$ has exactly one entry that may be nonzero, of modulus $|h|$. Hence, with \textup{(D2)},
\[
\|S_PC\|_\infty\le\tau:=(1+\Psi)|h|\le\tfrac98\cdot\tfrac1{320}=\tfrac{9}{2560}<1 .
\]
Taking norms in \eqref{eq:b1:closed}: $\|R_P\|_\infty\le\|S_P\|_\infty+\tau\|R_P\|_\infty$, so $\|R_P\|_\infty\le(1+\Psi)/(1-\tau)$, and then, again by \eqref{eq:b1:closed},
\[
\|R_P-S_P\|_\infty=\|S_PC\,R_P\|_\infty\le\frac{\tau\,(1+\Psi)}{1-\tau}=\frac{(1+\Psi)^2}{1-\tau}\,|h|\le\frac{81}{64}\cdot\frac{2560}{2551}\,|h|\le\frac43|h| ,
\]
because $81\cdot2560\cdot3=622080\le4\cdot64\cdot2551=653056$. This is the first bound in \eqref{eq:b1:RS}; $\frac43|h|\le\frac43t\le\frac1{30}\Psi$ by \textup{(D2)}.
Finally $|R_{kl}+\delta_{kl}|\le|S_{kl}+\delta_{kl}|+|R_{kl}-S_{kl}|\le(\frac12+\frac1{30})\Psi\le\Psi$.
\end{proof}
\noindent\emph{Remark.} The exact identity \eqref{eq:b1:closed} replaces the tenth-order expansion with the trivial bound $\eta^{-1}$ of the original, (6.45);
it is the only place where information is transferred from $H$ to $Q=H^{[p]}$. No inverse of $1-S_PC$ and no Neumann series is needed.
The statement is pointwise in $y$; in particular the endpoint condition \eqref{eq:b1:end} for $Q$ follows from $|h_{ab}|\le t$ and the bound
$|G_H(z_i)_{kl}+\delta_{kl}|\le\frac12\Psi$ at the two points $z_1,z_2$ alone.

\begin{lemma}[expansion of the functional; deterministic]\label{lem:expand}
Assume \textup{(D1)--(D3)}. Let $p=(a,b)$, $Q\in\cG_p$ \textup{(or, more generally, let $Q\in M_N(\C)$ be Hermitian and satisfy the endpoint condition \eqref{eq:b1:end})},
$R=G_Q$, $x_0:=X(Q)$, and $u,w,g$ as in Lemma~\ref{lem:det}. Let $c_1:=4$.

\smallskip\noindent\textup{\textbf{Case $a\ne b$.}} For $v\in\K$ define the real numbers
\begin{gather*}
\Delta_1(v):=-\tfrac1\pi\im\br{u(v;\cdot)},\qquad\Delta_2(v):=\tfrac1\pi\im\br{w(v;\cdot)+\tfrac12u(v;\cdot)^2},\\
\Delta_3(v):=-\tfrac1\pi\im\br{u(v;\cdot)\,w(v;\cdot)},
\end{gather*}
\[
P_1(v):=F'(x_0)\Delta_1(v),\qquad P_2(v):=F'(x_0)\Delta_2(v)+\tfrac12F''(x_0)\Delta_1(v)^2 .
\]
\textup{(a)} $\Delta_k$ is a homogeneous real polynomial of degree $k$ in $(v,\bar v)$ whose coefficients are polynomials
in $R_{kl}(z_i)$, $k,l\in\{a,b\}$, $i=1,2$, and their complex conjugates; explicitly,
\begin{align}
\Delta_1(v)&=-\tfrac1\pi\im\Big(v\,\br{R_{ba}}+\bar v\,\br{R_{ab}}\Big),\label{eq:b1:D1}\\
\Delta_2(v)&=\tfrac1\pi\im\Big(|v|^2\,\br{R_{aa}R_{bb}}+\tfrac12v^2\,\br{R_{ba}^2}+\tfrac12\bar v^2\,\br{R_{ab}^2}\Big),\label{eq:b1:D2}\\
\Delta_3(v)&=-\tfrac1\pi\im\Big(|v|^2v\,\br{R_{ba}\,g}+|v|^2\bar v\,\br{R_{ab}\,g}\Big).\label{eq:b1:D3}
\end{align}
For all $v\in\K$: $|\Delta_1(v)|\le2|v|\Psi$, $|\Delta_2(v)|\le3|v|^2\Psi$, $|\Delta_3(v)|\le3|v|^3\Psi$. Moreover $|g(z_i)-1|\le\frac94\Psi$ and $|g(z_i)|\le\frac43$ for $i=1,2$.\\
\textup{(b)} For all $v\in\K$ with $|v|\le\Psi/40$:
\[
\Big|F\big(X(Q+V_p(v))\big)-F(x_0)-P_1(v)-P_2(v)-F'(x_0)\Delta_3(v)\Big|\le c_1C_F\big(|v|^3\Psi^2+|v|^4\Psi\big),
\]
and $|X(Q+V_p(v))-x_0|\le\frac32|v|\Psi$.

\smallskip\noindent
\textup{\textbf{Case $a=b$}} ($v\in\R$). With $\Delta_1(v):=-\frac v\pi\im\br{R_{aa}}$, $\Delta_2(v):=\frac{v^2}{2\pi}\im\br{R_{aa}^2}$ and the same
$P_1,P_2$: $|\Delta_1(v)|\le|v|\Psi$ and $|\Delta_2(v)|\le|v|^2\Psi$ for all $v\in\R$, and for $|v|\le\frac14$:
\[
\big|F(X(Q+V_p(v)))-F(x_0)-P_1(v)-P_2(v)\big|\le c_1C_F\big(|v|^3\Psi+|v|^4\big),
\]
and $|X(Q+V_p(v))-x_0|\le|v|\Psi+|v|^4$.
\end{lemma}

\begin{proof}
Throughout, $z\in\{z_1,z_2\}$, all resolvent entries are evaluated at $z$, and we use \eqref{eq:b1:end}, $\Psi\le\frac18$, $\frac1\pi<\frac13$ and \eqref{eq:b1:br}.

\medskip\noindent\textbf{Case $a\ne b$. Step 0: sizes of $u$, $g$, $w$ at $z_1,z_2$.} Write $R_{aa}=-1+\alpha$, $R_{bb}=-1+\beta$. By \eqref{eq:b1:end},
$|\alpha|,|\beta|,|R_{ab}|,|R_{ba}|\le\Psi$. Hence
\begin{align}
&|u|\le|v||R_{ba}|+|v||R_{ab}|\le2|v|\Psi,\label{eq:b1:u}\\
&g-1=(-1+\alpha)(-1+\beta)-R_{ab}R_{ba}-1=-\alpha-\beta+\alpha\beta-R_{ab}R_{ba},\qquad|g-1|\le2\Psi+2\Psi^2\le\tfrac94\Psi,\label{eq:b1:g}\\
&|g|\le1+\tfrac94\cdot\tfrac18=\tfrac{41}{32}\le\tfrac43,\qquad|w|=|v|^2|g|\le\tfrac43|v|^2,\label{eq:b1:w}\\
&|\im w|=|v|^2|\im(g-1)|\le\tfrac94|v|^2\Psi,\label{eq:b1:imw}\\
&|\im w^2|=|v|^4|\im(g^2-1)|\le|v|^4|g-1||g+1|\le|v|^4\cdot\tfrac94\Psi\cdot\big(2+\tfrac9{32}\big)=\tfrac{657}{128}|v|^4\Psi\le6|v|^4\Psi .\label{eq:b1:imw2}
\end{align}
In \eqref{eq:b1:imw} and \eqref{eq:b1:imw2} we used that $|v|^2$ and $|v|^4$ are real: \emph{the leading parts of the even-order terms $w$, $w^2$ are real and drop
out under $\im$.}

\medskip\noindent\textbf{Step 1: part (a).} By definition $u=vR_{ba}+\bar vR_{ab}$, $u^2=v^2R_{ba}^2+2|v|^2R_{ab}R_{ba}+\bar v^2R_{ab}^2$ and
$w=|v|^2(R_{aa}R_{bb}-R_{ab}R_{ba})$, so
\[
w+\tfrac12u^2=|v|^2R_{aa}R_{bb}+\tfrac12v^2R_{ba}^2+\tfrac12\bar v^2R_{ab}^2,\qquad uw=|v|^2v\,R_{ba}g+|v|^2\bar v\,R_{ab}g .
\]
Since $v$ does not depend on the spectral parameter, taking $\br{\cdot}$ gives \eqref{eq:b1:D1}--\eqref{eq:b1:D3}. For a complex number $c$ and integers
$j,l\ge0$, $\im(c\,v^j\bar v^l)=\frac1{2i}(c\,v^j\bar v^l-\bar c\,\bar v^jv^l)$; hence each $\Delta_k$ is a real valued homogeneous polynomial of degree $k$ in
$(v,\bar v)$, with coefficients as stated. The bounds follow from \eqref{eq:b1:br} and Step~0:
\begin{align}
|\Delta_1|&\le\tfrac1\pi\big(|u(z_1)|+|u(z_2)|\big)\le\tfrac4\pi|v|\Psi\le\tfrac43|v|\Psi,\label{eq:b1:bD1}\\
|\Delta_2|&\le\tfrac1\pi\sum_{i=1,2}\big(|\im w(z_i)|+\tfrac12|u(z_i)|^2\big)\le\tfrac2\pi\big(\tfrac94+2\Psi\big)|v|^2\Psi\le\tfrac2\pi\cdot\tfrac52|v|^2\Psi\le\tfrac53|v|^2\Psi,\label{eq:b1:bD2}\\
|\Delta_3|&\le\tfrac1\pi\sum_{i=1,2}|u(z_i)||w(z_i)|\le\tfrac2\pi\cdot2|v|\Psi\cdot\tfrac43|v|^2\le\tfrac{16}9|v|^3\Psi .\label{eq:b1:bD3}
\end{align}
These are stronger than the bounds $2|v|\Psi$, $3|v|^2\Psi$, $3|v|^3\Psi$ claimed in (a). No restriction on $|v|$ was used.

\medskip\noindent\textbf{Step 2: the exact formula.} Let $|v|\le\Psi/40$, so $|v|\le\frac1{320}$. Put $\zeta(z):=u(v;z)-w(v;z)=D_p(v;z)-1$. By \eqref{eq:b1:u}, \eqref{eq:b1:w}
and $|v|\le\Psi/40$,
\begin{equation}\label{eq:b1:zeta}
|\zeta(z_i)|\le2|v|\Psi+\tfrac43|v|\cdot\tfrac{\Psi}{40}=\tfrac{61}{30}|v|\Psi\ \le\ \tfrac{61}{30}\cdot\tfrac1{320}\cdot\tfrac18<\tfrac12\qquad(i=1,2).
\end{equation}
By Corollary~\ref{cor:Xchange}(i),
\begin{equation}\label{eq:b1:Delta}
\Delta:=X(Q+V_p(v))-x_0=-\frac1\pi\im\br{\Log(1+\zeta)} .
\end{equation}
(If $Q\in\cG_p$, the bound \eqref{eq:b1:zeta} holds at $y+i\eta$ for every $y\in[E_1,E_2]$ by the same computation, so that $|D_p(v;y+i\eta)-1|\le\frac12$ on the whole
segment; by \eqref{eq:b1:nocut} this information is not needed.)

\medskip\noindent\textbf{Step 3: expansion of the logarithm.} Let $r_3(\zeta):=\Log(1+\zeta)-\zeta+\frac12\zeta^2$. By \eqref{eq:b1:L2} and \eqref{eq:b1:zeta},
\[
|r_3(\zeta(z_i))|\le\tfrac23\big(\tfrac{61}{30}\big)^3|v|^3\Psi^3\le6|v|^3\Psi^3\qquad\big(\tfrac23\cdot\tfrac{226981}{27000}<5.61\big).
\]
Algebraically, $\zeta-\frac12\zeta^2=u-w-\frac12(u-w)^2=u-\big(w+\frac12u^2\big)+uw-\frac12w^2$. Inserting $\Log(1+\zeta)=\zeta-\frac12\zeta^2+r_3(\zeta)$ into
\eqref{eq:b1:Delta} and comparing with the definitions of $\Delta_1,\Delta_2,\Delta_3$,
\begin{equation}\label{eq:b1:split}
\Delta=\Delta_1+\Delta_2+\Delta_3+\Delta_4+\Delta_r,\qquad\Delta_4:=\tfrac1{2\pi}\im\br{w^2},\qquad\Delta_r:=-\tfrac1\pi\im\br{r_3(\zeta)},
\end{equation}
an exact identity, and by \eqref{eq:b1:imw2} and \eqref{eq:b1:br},
\begin{equation}\label{eq:b1:bD4}
|\Delta_4|\le\tfrac1{2\pi}\cdot2\cdot6|v|^4\Psi\le2|v|^4\Psi,\qquad|\Delta_r|\le\tfrac2\pi\cdot6|v|^3\Psi^3\le4|v|^3\Psi^3 .
\end{equation}
\textbf{Step 4: sizes.} By \eqref{eq:b1:split}, \eqref{eq:b1:bD1}--\eqref{eq:b1:bD3}, \eqref{eq:b1:bD4}, and $|v|\le\frac1{320}$, $\Psi\le\frac18$:
\begin{align}
|\Delta-\Delta_1|&\le|\Delta_2|+|\Delta_3|+|\Delta_4|+|\Delta_r|\le|v|^2\Psi\Big(\tfrac53+\tfrac{16}9|v|+2|v|^2+4|v|\Psi^2\Big)\le2|v|^2\Psi,\label{eq:b1:s1}\\
|\Delta+\Delta_1|&\le2|\Delta_1|+|\Delta-\Delta_1|\le\big(\tfrac83+2|v|\big)|v|\Psi\le3|v|\Psi,\label{eq:b1:s2}\\
|\Delta|&\le|\Delta_1|+|\Delta-\Delta_1|\le\big(\tfrac43+2|v|\big)|v|\Psi\le\tfrac32|v|\Psi .\label{eq:b1:s3}
\end{align}
(In \eqref{eq:b1:s1}: $\frac53+\frac{16}{9\cdot320}+\frac2{320^2}+\frac4{320\cdot64}<1.68$.)

\medskip\noindent\textbf{Step 5: Taylor expansion of $F$.} By Taylor's formula with Lagrange remainder and \textup{(D3)},
\[
F(x_0+\Delta)=F(x_0)+F'(x_0)\Delta+\tfrac12F''(x_0)\Delta^2+\rho_3,\qquad|\rho_3|\le\tfrac16C_F|\Delta|^3 .
\]
Subtract $P_1+P_2+F'(x_0)\Delta_3=F'(x_0)(\Delta_1+\Delta_2+\Delta_3)+\frac12F''(x_0)\Delta_1^2$ and use \eqref{eq:b1:split} and $\Delta^2-\Delta_1^2=(\Delta-\Delta_1)(\Delta+\Delta_1)$:
\begin{equation}\label{eq:b1:err}
F(x_0+\Delta)-F(x_0)-P_1-P_2-F'(x_0)\Delta_3=F'(x_0)\big(\Delta_4+\Delta_r\big)+\tfrac12F''(x_0)(\Delta-\Delta_1)(\Delta+\Delta_1)+\rho_3 .
\end{equation}
By \eqref{eq:b1:bD4}, \eqref{eq:b1:s1}--\eqref{eq:b1:s3} the modulus of the right side is at most
\begin{align*}
C_F\Big(2|v|^4\Psi+4|v|^3\Psi^3+\tfrac12\cdot2\cdot3\,|v|^3\Psi^2+\tfrac16\cdot\tfrac{27}8|v|^3\Psi^3\Big)
&=C_F\Big(\big(3+\tfrac{73}{16}\Psi\big)|v|^3\Psi^2+2|v|^4\Psi\Big)\\
&\le C_F\Big(\tfrac{457}{128}|v|^3\Psi^2+2|v|^4\Psi\Big),
\end{align*}
and $\frac{457}{128}<3.58\le c_1$. Since $X(Q+V_p(v))=x_0+\Delta$, this is the first claim of (b); the second is \eqref{eq:b1:s3}.
The bounds on $g$ stated in (a) are \eqref{eq:b1:g}, \eqref{eq:b1:w}.

\medskip\noindent\textbf{Case $a=b$} ($v\in\R$). Write $R_{aa}=-1+\alpha$ with $|\alpha|\le\Psi$ at $z_1,z_2$ by \eqref{eq:b1:end}. Then, at $z_1$ and $z_2$,
\begin{equation}\label{eq:b1:diag0}
\begin{aligned}
&|R_{aa}|\le1+\Psi\le\tfrac98,\qquad|\im R_{aa}|=|\im\alpha|\le\Psi,\\
&|\im R_{aa}^2|=|\im(R_{aa}^2-1)|=|\im(\alpha(\alpha-2))|\le(2+\Psi)\Psi\le\tfrac{17}8\Psi,\\
&|\im R_{aa}^3|=|\im(R_{aa}^3+1)|=|\im(\alpha(3-3\alpha+\alpha^2))|\le(3+3\Psi+\Psi^2)\Psi\le\tfrac{217}{64}\Psi ,
\end{aligned}
\end{equation}
where we used $R^2-1=(R+1)(R-1)$ and $R^3+1=(R+1)(R^2-R+1)$ with $R=-1+\alpha$. Since $v$ is real, \eqref{eq:b1:br} gives, for all $v\in\R$,
\begin{equation}\label{eq:b1:diag1}
|\Delta_1|\le\tfrac{|v|}\pi\cdot2\Psi\le\tfrac23|v|\Psi,\qquad|\Delta_2|\le\tfrac{v^2}{2\pi}\cdot2\cdot\tfrac{17}8\Psi\le\tfrac{17}{24}v^2\Psi ,
\end{equation}
which are stronger than the claimed bounds $|v|\Psi$, $|v|^2\Psi$.

Let now $|v|\le\frac14$ and $\zeta(z):=vR_{aa}(z)=D_p(v;z)-1$. Then $|\zeta(z_i)|\le\frac98|v|\le\frac9{32}<\frac12$, and by Corollary~\ref{cor:Xchange}(i),
$\Delta:=X(Q+V_p(v))-x_0=-\frac1\pi\im\br{\Log(1+\zeta)}$. Let $r_4(\zeta):=\Log(1+\zeta)-\zeta+\frac12\zeta^2-\frac13\zeta^3$; by \eqref{eq:b1:L3},
$|r_4(\zeta(z_i))|\le\frac12\big(\frac98\big)^4v^4=\frac{6561}{8192}v^4$. As $v$ is real, $-\frac1\pi\im\br{\zeta}=\Delta_1$ and $-\frac1\pi\im\br{-\frac12\zeta^2}=\Delta_2$, so
\begin{equation}\label{eq:b1:diagsplit}
\Delta=\Delta_1+\Delta_2+\Delta_3'+\Delta_r',\qquad\Delta_3':=-\tfrac{v^3}{3\pi}\im\br{R_{aa}^3},\qquad\Delta_r':=-\tfrac1\pi\im\br{r_4(\zeta)},
\end{equation}
and by \eqref{eq:b1:diag0}, \eqref{eq:b1:br},
\begin{equation}\label{eq:b1:diag2}
|\Delta_3'|\le\tfrac{|v|^3}{3\pi}\cdot2\cdot\tfrac{217}{64}\Psi\le\tfrac{217}{288}|v|^3\Psi\le|v|^3\Psi,\qquad|\Delta_r'|\le\tfrac2\pi\cdot\tfrac{6561}{8192}v^4\le v^4 .
\end{equation}
Using $|v|\le\frac14$ and $\Psi\le\frac18$ we get from \eqref{eq:b1:diag1}--\eqref{eq:b1:diag2}
\begin{align}
|\Delta-\Delta_1|&\le v^2\Psi+|v|^3\Psi+v^4\le\tfrac54v^2\Psi+v^4,\label{eq:b1:ds1}\\
|\Delta|&\le\tfrac23|v|\Psi+\tfrac54v^2\Psi+v^4\le\big(\tfrac23+\tfrac5{16}\big)|v|\Psi+v^4\le|v|\Psi+v^4\le|v|\big(\tfrac18+\tfrac1{64}\big)\le\tfrac15|v|,\label{eq:b1:ds2}\\
|\Delta+\Delta_1|&\le|\Delta|+|\Delta_1|\le\tfrac9{64}|v|+\tfrac1{12}|v|\le\tfrac13|v| .\label{eq:b1:ds3}
\end{align}
Exactly as in Step~5 (Taylor's formula to second order with remainder $|\rho_3|\le\frac16C_F|\Delta|^3$, and \eqref{eq:b1:diagsplit}),
\[
F(x_0+\Delta)-F(x_0)-P_1-P_2=F'(x_0)\big(\Delta_3'+\Delta_r'\big)+\tfrac12F''(x_0)(\Delta-\Delta_1)(\Delta+\Delta_1)+\rho_3 ,
\]
and, bounding $|\Delta|^3\le(\frac15|v|)^2(|v|\Psi+v^4)$ by \eqref{eq:b1:ds2}, the modulus of the right side is at most
\begin{align*}
&C_F\Big(|v|^3\Psi+v^4+\tfrac12\big(\tfrac54v^2\Psi+v^4\big)\tfrac{|v|}3+\tfrac1{150}\big(|v|^3\Psi+|v|^6\big)\Big)\\
&\qquad\le C_F\Big(\big(1+\tfrac5{24}+\tfrac1{150}\big)|v|^3\Psi+\big(1+\tfrac1{24}+\tfrac1{2400}\big)v^4\Big)\le2C_F\big(|v|^3\Psi+v^4\big),
\end{align*}
where $|v|^5\le\frac14v^4$ and $|v|^6\le\frac1{16}v^4$ were used. Since $2\le c_1$, this is the claimed expansion; the bound
$|X(Q+V_p(v))-x_0|=|\Delta|\le|v|\Psi+v^4$ is contained in \eqref{eq:b1:ds2}.
\end{proof}

\noindent\emph{Remarks on Lemma~\ref{lem:expand}.}
(1) \emph{Where the factors $\Psi$ come from.} Every term of the expansion carries at least one factor $\Psi$: the odd-order terms $u$, $uw$ through the
off-diagonal entries $R_{ab},R_{ba}$ \eqref{eq:b1:u}; the even-order terms $w$, $w^2$ through $\im$, because $g=1+O(\Psi)$ and the leading parts
$|v|^2$, $|v|^4$ are real \eqref{eq:b1:imw}, \eqref{eq:b1:imw2}. For a diagonal entry the same mechanism gives one factor $\Psi$ up to third order \eqref{eq:b1:diag0},
because $R_{aa}^k=(-1)^k+O(\Psi)$.

(2) \emph{Size table.} Table~1 lists every term that occurs in the proof for $a\ne b$, with the constant proved above. The third column gives the order of the
term after integration against a law $\nu$ with $\int|v|^k\,d\nu\le\mathfrak m_kC_0^{k/2}N^{-k/2}$ \eqref{eq:b3:M}, for $\Psi=4N^{-1/3+2\eg}$; it is meant as orientation only
(the integration, including the region $|v|>t$, is carried out in Corollary~\ref{cor:entry}).

\begin{center}\footnotesize
\begin{tabular}{l|l|l|l}
term & bound under \eqref{eq:b1:end}, $|v|\le\Psi/40$ & order after $\int\!\cdot\,d\nu$ & fate\\\hline
$F'\Delta_1$ & $\frac43C_F|v|\Psi$ \eqref{eq:b1:bD1} & $=0$ exactly (mean zero) & $P_1$\\[1pt]
$F'\Delta_2$ & $\frac53C_F|v|^2\Psi$ \eqref{eq:b1:bD2} & second moments only & $P_2$, matched\\[1pt]
$\frac12F''\Delta_1^2$ & $\frac89C_F|v|^2\Psi^2$ & second moments only & $P_2$, matched\\[1pt]
$F'\Delta_3$ & $\frac{16}9C_F|v|^3\Psi$ \eqref{eq:b1:bD3} & $N^{-3/2}\Psi\sim N^{-11/6+2\eg}$ & \textbf{key term}, Lemma~\ref{lem:decor}\\[1pt]
$\frac12F''(\Delta-\Delta_1)(\Delta+\Delta_1)$ & $3C_F|v|^3\Psi^2$ \eqref{eq:b1:s1}, \eqref{eq:b1:s2} & $N^{-3/2}\Psi^2\sim N^{-13/6+4\eg}$ & error (leading)\\[1pt]
$F'\Delta_r$ & $4C_F|v|^3\Psi^3$ \eqref{eq:b1:bD4} & $N^{-5/2+6\eg}$ & error\\[1pt]
$\rho_3$ & $\frac9{16}C_F|v|^3\Psi^3$ \eqref{eq:b1:s3} & $N^{-5/2+6\eg}$ & error\\[1pt]
$F'\Delta_4$ & $2C_F|v|^4\Psi$ \eqref{eq:b1:bD4} & $N^{-2}\Psi\sim N^{-7/3+2\eg}$ & error
\end{tabular}\\[3pt]
Table 1: the terms of the expansion for $a\ne b$.
\end{center}

\noindent The leading part of the fifth row is $F''\Delta_1\Delta_2$; the fifth row also contains $F''(\Delta-\Delta_1)^2$ and $F''\Delta_1\Delta_3$
(both of size $|v|^4\Psi^2$), and the terms $F'''\Delta_1^3$ and $\rho_4$ of a fourth-order expansion are contained in $\rho_3$. The four error rows add up to at most
$C_F\big(\frac{457}{128}|v|^3\Psi^2+2|v|^4\Psi\big)$. For $a=b$ and $|v|\le\frac14$ the corresponding list is: $F'\Delta_1$: $\frac23C_F|v|\Psi$; $F'\Delta_2$: $\frac{17}{24}C_Fv^2\Psi$;
$\frac12F''\Delta_1^2$: $\frac29C_Fv^2\Psi^2$ (these form $P_1$, $P_2$); errors: $F'\Delta_3'$: $C_F|v|^3\Psi$ (order $N^{-3/2}\Psi\sim N^{-11/6+2\eg}$ after integration);
$F'\Delta_r'$: $C_Fv^4$ (order $N^{-2}$); $\frac12F''(\Delta-\Delta_1)(\Delta+\Delta_1)$: $C_F(\frac5{24}|v|^3\Psi+\frac16|v|^5)$; $\rho_3$: $\frac1{150}C_F(|v|^3\Psi+|v|^6)$.

\textbf{Exponent check} (for the standing assumptions, $\Psi=4N^{-1/3+2\eg}$): off-diagonal, $N^{-3/2}\cdot\Psi^2=16N^{-3/2-2/3+4\eg}=16N^{-13/6+4\eg}$ and
$N^{-2}\cdot\Psi=4N^{-7/3+2\eg}\le4N^{-13/6+4\eg}$; diagonal, $N^{-3/2}\cdot\Psi=4N^{-11/6+2\eg}$ and $N^{-2}\le N^{-11/6+2\eg}$.

(3) \emph{The key term.} By \eqref{eq:b1:D3}, for a probability measure $\nu$ on $\K$ with $\int|v|^3d\nu<\infty$ and $\mu_3:=\int|v|^2v\,d\nu$,
\[
\int\Delta_3(v)\,\nu(dv)=-\frac1\pi\im\br{(\mu_3R_{ba}+\overline{\mu_3}R_{ab})\,g} ,
\]
which is the quantity $\mathfrak T_3(Q;\nu)$ of Corollary~\ref{cor:entry}; and by \eqref{eq:b1:D1}, \eqref{eq:b1:D2}, $\int P_1\,d\nu=0$ if $\int v\,d\nu=0$, while $\int P_2\,d\nu$ depends on
$\nu$ only through $\int v^2d\nu$, $\int|v|^2d\nu$, $\int\bar v^2d\nu$.

(4) \emph{Endpoints only; consequence for the use of Theorem~\ref{thm:21}.} Lemma~\ref{lem:expand} and Corollary~\ref{cor:Xchange} use the resolvent of $Q$ only at the two points
$z_1,z_2$, and Lemma~\ref{lem:RfromS}(a) is pointwise in $y$. Hence all deterministic statements of \S\ref{sec:X} and of this subsection remain valid if the good sets
($\cG_p$ in Definition~\ref{def:goodB}, $\Omega_{\rm edge}^c$ in Lemma~\ref{lem:edge}) are defined by conditions at $y\in\{E_1,E_2\}$ only. With such definitions the proof plans of \S\ref{sec:step}--\ref{sec:gfc} would use
the bound of Theorem~\ref{thm:21}(ii) at the two fixed spectral parameters $z_1,z_2$ only, and not its simultaneity in $z$. We do \emph{not} use this simplification here;
it is recorded as an optional simplification for the formalization.

\paragraph{Moments and tail of one entry.} Put $\hat\vartheta:=\min(\vartheta,1)$ and
\[
\mathfrak m_k=\mathfrak m_k(\vartheta):=\Big(\frac{90\,k}{\hat\vartheta}\Big)^{k/\hat\vartheta}\qquad(k=3,4,6).
\]
These are the moment constants of the standing assumptions (the letters $C_2,C_3,C_4$ denote other constants below).
Let $\nu$ be a Borel probability measure on $\K$ such that
\begin{equation}\label{eq:b3:nu}
\sigma^2:=\int|y|^2\,\nu(dy)\le\frac{C_0}N\qquad\text{and}\qquad\nu\big(\{y:|y|>x\sigma\}\big)\le\vartheta^{-1}e^{-x^\vartheta}\quad\text{for all real }x\ge1 .
\end{equation}
Every law $\nu_p$ of an ensemble of size $N$ with $\fA(\Par)$ satisfies \eqref{eq:b3:nu}, by (C) and \eqref{eq:subexp}. Under \eqref{eq:b3:nu}:
\begin{align}
&\int|y|^k\,\nu(dy)\le\mathfrak m_k\,\sigma^k\le\mathfrak m_k\,C_0^{k/2}N^{-k/2}\qquad(k=3,4,6),\label{eq:b3:M}\\
&\nu\big(\{y:|y|>t\}\big)\le N^{-20}.\label{eq:b3:T}
\end{align}
\emph{Proof of \eqref{eq:b3:M}.} For $x\ge1$ we have $\vartheta^{-1}e^{-x^\vartheta}\le\hat\vartheta^{-1}e^{-x^{\hat\vartheta}}$: this is an equality if $\vartheta\le1$, and if $\vartheta>1$ then
$\vartheta^{-1}<1=\hat\vartheta^{-1}$ and $x^\vartheta\ge x$. So $\nu$, regarded as a measure on $\C$, satisfies the hypothesis of Lemma~\ref{lem:moments} with the parameter
$\hat\vartheta\in(0,1]$ and $s=\sigma$, and Lemma~\ref{lem:moments}(a) with $q=k$ gives $\int|y|^k d\nu\le(90k/\hat\vartheta)^{k/\hat\vartheta}\sigma^k$.
\emph{Proof of \eqref{eq:b3:T}.} By \eqref{eq:b3:nu} and \textup{(N4)}, $N^{\eg}\sigma\le N^{\eg}C_0^{1/2}N^{-1/2}\le N^{-1/2+2\eg}=t$, so
$\{|y|>t\}\subset\{|y|>N^{\eg}\sigma\}$. Since $N^{\eg}\ge1$, \eqref{eq:b3:nu} with $x=N^{\eg}$ and \textup{(N3)} give
$\nu(|y|>t)\le\vartheta^{-1}\exp(-N^{\eg\vartheta})\le N^{-20}$. \qed

\begin{proof}[Proof of Lemma~\ref{lem:RfromS}\textup{(b)}]
Let $M$ be an ensemble of size $N$ over $\K$ with $\fA(\Par)$ and $p=(a,b)\in\mathcal P_N$. We show
\begin{equation}\label{eq:b3:incl}
\big\{x\in\Omega:\ H(x)^{[p]}\notin\cG_p\big\}\ \subset\ \Omega_{\rm edge}\cup\big\{x\in\Omega:\ |x_p|>t\big\}.
\end{equation}
Let $x\notin\Omega_{\rm edge}$ with $|x_p|\le t$, and $H=H(x)$. By Definition~\ref{def:ens}, $h_{ab}=x_p$ if $a<b$ and $h_{aa}=\re x_p$ if $a=b$; in both cases $|h_{ab}|\le|x_p|\le t$.
By the standing assumptions $|E_i-2|\le N^{-2/3+\eg}$ for $i=1,2$, so $[E_1,E_2]\subset I_{\rm edge}$, and since $x\in\Omega^c_{\rm edge}$ (Lemma~\ref{lem:edge}),
$|G_H(y+i\eta)_{kl}+\delta_{kl}|\le\frac12\Psi$ for all $k,l\in\{a,b\}$ and all $y\in[E_1,E_2]$. By Lemma~\ref{lem:RfromS}(a) (its hypothesis \textup{(D2)} is \textup{(N4)}),
$H^{[p]}\in\cG_p$. This proves \eqref{eq:b3:incl}.

The three sets in \eqref{eq:b3:incl} are Borel: the first is open (the remark after Definition~\ref{def:goodB}), $\Omega_{\rm edge}$ is open (Lemma~\ref{lem:edge}), and
$\{|x_p|>t\}$ is open. By Lemma~\ref{lem:edge}, $\P_M(\Omega_{\rm edge})\le N^{-20}$. Under $\P_M$ the coordinate $x_p$ has the law $\nu_p$, with $\sigma_p^2\le C_0/N$ by (C)
and with \eqref{eq:subexp}; hence $\P_M(|x_p|>t)=\nu_p(\{|y|>t\})\le N^{-20}$ by the tail bound \eqref{eq:b3:T} above (whose proof uses only \textup{(N3)}, \textup{(N4)}).
Adding the two bounds gives $\P_M(H^{[p]}\notin\cG_p)\le2N^{-20}$.
\end{proof}

\begin{corollary}[integration over one entry]\label{cor:entry}
Assume the standing assumptions of \S\ref{sec:step} \textup{(}in particular $\eg\in(0,\frac1{10}]$ and $N\ge N_1(\Par,\eg)$\textup{)}.
Let $p=(a,b)\in\mathcal P_N$, let $Q\in\cG_p$, $R=G_Q$, $x_0=X(Q)$, and let $u,w,g,\Delta_1,\Delta_2,\Delta_3,P_1,P_2$ be as in Lemmas~\ref{lem:det} and~\ref{lem:expand}.
Let $\nu$ be a Borel probability measure on $\K$, with $\nu(\R)=1$ if $a=b$, which satisfies \eqref{eq:b3:nu} and $\int y\,\nu(dy)=0$. Put
\[
m_{11}:=\int|y|^2\nu(dy)=\sigma^2,\quad m_{20}:=\int y^2\,\nu(dy),\quad\mu_3:=\int|y|^2y\,\nu(dy)\quad\big(\text{so }|\mu_3|\le\mathfrak m_3C_0^{3/2}N^{-3/2}\big),
\]
and
\begin{equation}\label{eq:b3:C2}
C_2=C_2(\vartheta,C_0):=64\,\mathfrak m_3C_0^{3/2}+16\,\mathfrak m_4C_0^{2}+6\big(1+\mathfrak m_6^{1/2}C_0^{3/2}\big).
\end{equation}
Then $v\mapsto F(X(Q+V_p(v)))$ is a bounded continuous function on $\K$, and with
\[
\mathfrak D(Q;\nu):=\int_\K\big[F(X(Q+V_p(v)))-F(X(Q))\big]\,\nu(dv)
\]
we have
\begin{align}
a\ne b:\quad&\Big|\mathfrak D(Q;\nu)-\mathfrak Q_2(Q;\nu)-F'(X(Q))\,\mathfrak T_3(Q;\nu)\Big|\le C_2C_F\,N^{-13/6+4\eg},\label{eq:b3:entry}\\
a=b:\quad&\Big|\mathfrak D(Q;\nu)-\mathfrak Q_2(Q;\nu)\Big|\le C_2C_F\,N^{-11/6+2\eg},\label{eq:b3:entryd}
\end{align}
where $\mathfrak Q_2(Q;\nu):=\int P_2\,d\nu$ and $\mathfrak T_3(Q;\nu):=\int\Delta_3\,d\nu$ are given by the following formulas. For $a\ne b$, with
$\varkappa:=\br{R_{ba}}-\overline{\br{R_{ab}}}$,
\begin{align}
\mathfrak Q_2(Q;\nu)&=\frac{F'(x_0)}\pi\,\im\Big(m_{11}\br{R_{aa}R_{bb}}+\tfrac12m_{20}\br{R_{ba}^2}+\tfrac12\overline{m_{20}}\,\br{R_{ab}^2}\Big)\notag\\
&\qquad+\frac{F''(x_0)}{4\pi^2}\Big(|\varkappa|^2m_{11}-\re\big(\varkappa^2m_{20}\big)\Big),\label{eq:b3:Q2}\\
\mathfrak T_3(Q;\nu)&=-\frac1\pi\,\im\br{\big(\mu_3R_{ba}+\overline{\mu_3}\,R_{ab}\big)\,g};\label{eq:b3:T3}
\end{align}
for $a=b$,
\begin{equation}\label{eq:b3:Q2d}
\mathfrak Q_2(Q;\nu)=m_{11}\Big(\frac{F'(x_0)}{2\pi}\,\im\br{R_{aa}^2}+\frac{F''(x_0)}{2\pi^2}\big(\im\br{R_{aa}}\big)^2\Big).
\end{equation}
In particular $\mathfrak Q_2(Q;\nu)$ depends on $\nu$ only through the second moments $m_{11}=\int|y|^2d\nu$ and $m_{20}=\int y^2d\nu$ \textup{(}and $\int\bar y^2d\nu=\overline{m_{20}}$\textup{)},
and $\mathfrak T_3(Q;\nu)$ only through $\mu_3$. For fixed $\nu$, $\mathfrak Q_2(\cdot\,;\nu)$ and $\mathfrak T_3(\cdot\,;\nu)$ are continuous functions on the Hermitian matrices.
\end{corollary}
\begin{proof}
Throughout, resolvent entries without argument are evaluated at $z_1$ or $z_2$, and $t\le\Psi/40\le\frac1{320}$, $\Psi\le\frac18$ by \textup{(N4)}.
If $a=b$ then $\nu(\R)=1$ and $V_p(v)=vE^{(aa)}$ for $v\in\R$ (Definition~\ref{def:hyb}), so all integrals below are integrals over $\R$.

\emph{Step 0: integrability; the integrals of $P_1$, $P_2$, $\Delta_3$.} The map $v\mapsto Q+V_p(v)$ is $\R$-linear, hence continuous, into the Hermitian matrices, and $X$ is continuous there
(Lemma~\ref{lem:Xspec}(c)); so $v\mapsto F(X(Q+V_p(v)))$ is continuous, and bounded by $C_F$. By \eqref{eq:b3:M}, $\int|y|^3d\nu<\infty$, hence all polynomials of degree $\le3$ in
$(v,\bar v)$ are $\nu$-integrable; in particular $P_1,P_2,\Delta_3$ are, and $|\mu_3|\le\int|y|^3d\nu\le\mathfrak m_3C_0^{3/2}N^{-3/2}$. For a constant $c\in\C$ and a $\nu$-integrable function $\xi$
we have $\int\im(c\,\xi)\,d\nu=\im\big(c\int\xi\,d\nu\big)$, and $\int\bar\xi\,d\nu=\overline{\int\xi\,d\nu}$.

\emph{Case $a\ne b$.} By \eqref{eq:b1:D1} and $\int v\,d\nu=0=\int\bar v\,d\nu$: $\int\Delta_1d\nu=0$, hence $\int P_1d\nu=F'(x_0)\int\Delta_1d\nu=0$.
By \eqref{eq:b1:D2}, $\int\Delta_2\,d\nu=\frac1\pi\im\big(m_{11}\br{R_{aa}R_{bb}}+\frac12m_{20}\br{R_{ba}^2}+\frac12\overline{m_{20}}\br{R_{ab}^2}\big)$.
Next, for complex numbers $c_1,c_2$ and $v$: $\im(c_1v+c_2\bar v)=\im(c_1v)-\im(\overline{c_2}\,v)=\im\big((c_1-\overline{c_2})v\big)$. With $c_1=\br{R_{ba}}$, $c_2=\br{R_{ab}}$ this gives
$\Delta_1(v)=-\frac1\pi\im(\varkappa v)$ by \eqref{eq:b1:D1}. Since $(\im\xi)^2=\frac12\big(|\xi|^2-\re(\xi^2)\big)$ for $\xi\in\C$,
\[
\Delta_1(v)^2=\frac1{2\pi^2}\Big(|\varkappa|^2|v|^2-\re\big(\varkappa^2v^2\big)\Big),\qquad\int\Delta_1^2\,d\nu=\frac1{2\pi^2}\Big(|\varkappa|^2m_{11}-\re\big(\varkappa^2m_{20}\big)\Big).
\]
As $P_2=F'(x_0)\Delta_2+\frac12F''(x_0)\Delta_1^2$, this proves \eqref{eq:b3:Q2}. By \eqref{eq:b1:D3} and $\int|v|^2\bar v\,d\nu=\overline{\mu_3}$,
$\int\Delta_3\,d\nu=-\frac1\pi\im\big(\mu_3\br{R_{ba}g}+\overline{\mu_3}\br{R_{ab}g}\big)$, which is \eqref{eq:b3:T3}.

\emph{Case $a=b$.} Here $\Delta_1(v)=-\frac v\pi\im\br{R_{aa}}$ and $\Delta_2(v)=\frac{v^2}{2\pi}\im\br{R_{aa}^2}$ for $v\in\R$. So $\int P_1d\nu=0$ and, as $\int v^2d\nu=m_{11}$,
$\int P_2\,d\nu=F'(x_0)\frac{m_{11}}{2\pi}\im\br{R_{aa}^2}+\frac12F''(x_0)\frac{m_{11}}{\pi^2}(\im\br{R_{aa}})^2$, which is \eqref{eq:b3:Q2d}.

\emph{Continuity in $Q$.} $Q\mapsto X(Q)$ and $Q\mapsto R_{kl}(z_i)=G_Q(z_i)_{kl}$ are continuous on the Hermitian matrices (Lemma~\ref{lem:Xspec}(c), Lemma~\ref{lem:basic}(h)),
and $F',F''$ are continuous; so the right sides of \eqref{eq:b3:Q2}--\eqref{eq:b3:Q2d} are continuous in $Q$.

\emph{The error function.} Define, for $v\in\K$ ($v\in\R$ if $a=b$),
\[
e(v):=F\big(X(Q+V_p(v))\big)-F(x_0)-P_1(v)-P_2(v)-\mathbf 1(a\ne b)\,F'(x_0)\Delta_3(v).
\]
It is continuous and $\nu$-integrable, and by the above $\int e\,d\nu$ is the quantity inside the modulus in \eqref{eq:b3:entry}, resp.\ in \eqref{eq:b3:entryd}
(recall $F(x_0)=F(X(Q))$ and $\nu(\K)=1$). We bound
$\int|e|\,d\nu$ by splitting $\K=\{|v|\le t\}\cup\{|v|>t\}$.

\emph{Step 1: $a\ne b$, the region $|v|\le t$.} Since $t\le\Psi/40$, Lemma~\ref{lem:expand}(b) (with $c_1=4$) gives $|e(v)|\le4C_F(|v|^3\Psi^2+|v|^4\Psi)$ for $|v|\le t$. By \eqref{eq:b3:M}
and $\Psi=4N^{-1/3+2\eg}$,
\begin{align*}
\int_{\{|v|\le t\}}|e|\,d\nu&\le4C_F\Big(\mathfrak m_3C_0^{3/2}N^{-3/2}\cdot16N^{-2/3+4\eg}+\mathfrak m_4C_0^2N^{-2}\cdot4N^{-1/3+2\eg}\Big)\\
&=C_F\Big(64\,\mathfrak m_3C_0^{3/2}N^{-13/6+4\eg}+16\,\mathfrak m_4C_0^2N^{-7/3+2\eg}\Big)\\
&\le C_F\big(64\,\mathfrak m_3C_0^{3/2}+16\,\mathfrak m_4C_0^2\big)N^{-13/6+4\eg},
\end{align*}
because $-\frac32-\frac23=-\frac{13}6$ and $-\frac73+2\eg\le-\frac{13}6+4\eg$.

\emph{Step 2: $a\ne b$, the region $|v|>t$.} For \emph{all} $v\in\K$, $|F|\le C_F$ and Lemma~\ref{lem:expand}(a) give
\[
|e(v)|\le C_F\Big(2+2|v|\Psi+3|v|^2\Psi+\tfrac12(2|v|\Psi)^2+3|v|^3\Psi\Big)\le C_F\big(2+|v|+|v|^2+|v|^3\big)\le4C_F\big(1+|v|^3\big),
\]
where we used $\Psi\le\frac18$ (so $2\Psi\le1$, $3\Psi+2\Psi^2\le1$, $3\Psi\le1$) and $|v|,|v|^2\le1+|v|^3$. By the Cauchy--Schwarz inequality, $(1+s)^2\le2+2s^2$, \eqref{eq:b3:T} and \eqref{eq:b3:M} with $k=6$ (note $N^{-3}\le1$),
\begin{align*}
\int_{\{|v|>t\}}|e|\,d\nu&\le4C_F\,\nu(|v|>t)^{1/2}\Big(2+2\!\int|v|^6d\nu\Big)^{1/2}\le4\sqrt2\,C_F\,N^{-10}\big(1+\mathfrak m_6C_0^3\big)^{1/2}\\
&\le6\,C_F\big(1+\mathfrak m_6^{1/2}C_0^{3/2}\big)N^{-10}.
\end{align*}
Since $N^{-10}\le N^{-13/6+4\eg}$, Steps 1 and 2 give $\int|e|\,d\nu\le C_2C_FN^{-13/6+4\eg}$, which is \eqref{eq:b3:entry}.

\emph{Step 3: $a=b$.} For $v\in\R$ with $|v|\le t$ $(\le\frac14)$, Lemma~\ref{lem:expand} (case $a=b$) gives $|e(v)|\le4C_F(|v|^3\Psi+v^4)$, so by \eqref{eq:b3:M}
\begin{align*}
\int_{\{|v|\le t\}}|e|\,d\nu&\le4C_F\Big(\mathfrak m_3C_0^{3/2}N^{-3/2}\cdot4N^{-1/3+2\eg}+\mathfrak m_4C_0^2N^{-2}\Big)\\
&\le C_F\big(16\,\mathfrak m_3C_0^{3/2}+4\,\mathfrak m_4C_0^2\big)N^{-11/6+2\eg},
\end{align*}
because $-\frac32-\frac13=-\frac{11}6$ and $-2\le-\frac{11}6+2\eg$. For all $v\in\R$, by $|\Delta_1|\le|v|\Psi$ and $|\Delta_2|\le v^2\Psi$:
$|e(v)|\le C_F(2+|v|\Psi+v^2\Psi+\frac12v^2\Psi^2)\le C_F(2+|v|+v^2)\le4C_F(1+|v|^3)$, and exactly as in Step 2,
$\int_{\{|v|>t\}}|e|\,d\nu\le6C_F(1+\mathfrak m_6^{1/2}C_0^{3/2})N^{-10}\le6C_F(1+\mathfrak m_6^{1/2}C_0^{3/2})N^{-11/6+2\eg}$. The sum is at most $C_2C_FN^{-11/6+2\eg}$, which is \eqref{eq:b3:entryd}.
\end{proof}

\noindent\emph{Remarks.} (1) \emph{Exponent check.} Off-diagonal: third moment $N^{-3/2}$ times $\Psi^2=16N^{-2/3+4\eg}$ gives $N^{-13/6+4\eg}$; the fourth-order term $N^{-2}\Psi$
is smaller by $N^{-1/6-2\eg}$. Diagonal: $N^{-3/2}\cdot\Psi=4N^{-11/6+2\eg}$. The region $|v|>t$ costs $N^{-10}$.
(2) \emph{What remains.} After Corollary~\ref{cor:entry}, the difference of the two integrals for $\nu^v_p$ and $\nu^w_p$ is, up to the errors \eqref{eq:b3:entry},
$\mathfrak Q_2(Q;\nu^v_p)-\mathfrak Q_2(Q;\nu^w_p)=0$ (matching second moments) plus $F'(X(Q))\big(\mathfrak T_3(Q;\nu^v_p)-\mathfrak T_3(Q;\nu^w_p)\big)$. By Lemma~\ref{lem:expand}(a),
$|\mathfrak T_3(Q;\nu)|\le3\mathfrak m_3C_0^{3/2}N^{-3/2}\Psi$ on $\cG_p$, which is of order $N^{-11/6+2\eg}$: too large by a factor $N^{1/6}$ for $N^2/2$ steps. The missing factor $\Psi$
is gained in the next subsection, in expectation over $Q$.
(3) \emph{Vanishing variance.} If $\sigma=0$ then $\nu=\delta_0$ (Lemma~\ref{lem:moments}(a)), all quantities in \eqref{eq:b3:entry}, \eqref{eq:b3:entryd} vanish, and the statement is trivial; the proof does not treat this case separately and never divides by $\sigma$.

\subsection{The third-order term: decorrelation}\label{sec:decor}
The standing assumptions of \S\ref{sec:step} remain in force. With only second moments matched, the term $F'(X(Q))\,\mathfrak T_3(Q;\nu)$ of Corollary~\ref{cor:entry} does not
cancel between $\nu^v_p$ and $\nu^w_p$. Its trivial size is $N^{-3/2}\Psi\approx N^{-11/6+2\eg}$, too large by $N^{1/6}$ for $N^2/2$ steps. The next lemma gains one more
factor $\Psi$, in expectation. It plays the role of Lemma~6.6 of the original, but it concerns a \emph{single} off-diagonal resolvent entry at a fixed spectral parameter.
The mechanism: by the Schur complement formula, $R_{ab}=-R_{aa}\,Y$, where $Y$ is \emph{linear} in the entries of the row $a$ of $Q$, with coefficients that depend only
on the minor $Q^{(a)}$; so $Y$ has conditional mean zero given the minor. The two factors that multiply $Y$, namely $-R_{aa}$ and $\Phi(X(Q))$, differ from quantities
that depend only on the minor (the number $1$, and $\Phi(X(Q^{(a)}))$) by $O(\Psi)$, and $Y$ itself is $O(\Psi)$.

\begin{lemma}[decorrelation of one off-diagonal entry]\label{lem:decor}
Assume the standing assumptions of \S\ref{sec:step} \textup{(}in particular $\eg\in(0,\frac1{10}]$ and $N\ge N_1(\Par,\eg)$\textup{)}.
Let $M$ be any ensemble of size $N$ over $\K$ with $\fA(\Par)$, $p=(a,b)\in\mathcal P_N$ with $a\ne b$, $Q:=H^{[p]}$, $R:=G_Q$, and let
$\Phi\in C^1(\R)$ with $\sup_\R|\Phi|\le C_\Phi$ and $\sup_\R|\Phi'|\le C_\Phi$. Then for $z\in\{z_1,z_2\}$ and $(c,d)\in\{(a,b),(b,a)\}$, the function
$\mathbf 1(Q\in\cG_p)\,\Phi(X(Q))\,R_{cd}(z)$ is bounded and Borel on $\Omega$, and
\[
\Big|\E_M\big[\mathbf 1(Q\in\cG_p)\,\Phi(X(Q))\,R_{cd}(z)\big]\Big|\le7\,C_\Phi\Psi^2\qquad(=112\,C_\Phi N^{-2/3+4\eg}).
\]
\textup{(The proof gives the constant $\frac{11}3$ instead of $7$.)}
\end{lemma}
\begin{proof}
Fix $z\in\{z_1,z_2\}$ and $(c,d)\in\{(a,b),(b,a)\}$; we prove the bound for the entry $R_{cd}(z)$. The index $c$ is the one whose row will be integrated out. Recall $a<b$, so
$N\ge2$. All matrices below are functions of $x\in\Omega$.

\smallskip\noindent\emph{Step 0: notation, dependence on the coordinates, measurability.}
Let $\pi_p:\Omega\to\Omega$ be the linear map which replaces the coordinate $x_p$ by $0$ and leaves the other coordinates unchanged. By Definitions~\ref{def:ens} and~\ref{def:hyb},
\[
Q(x)=H(x)^{[p]}=H(\pi_px)\qquad(x\in\Omega),
\]
since $H(x)^{[p]}$ is $H(x)$ with the entries $(a,b),(b,a)$, i.e.\ $x_p,\overline{x_p}$, replaced by $0$. Let $Q^{(c)}(x)$ be the minor of $Q(x)$ with row and column $c$ removed; it is a
Hermitian matrix indexed by $\{1,\dots,N\}\setminus\{c\}$ (\S\ref{sec:resid}), of size $N-1\ge1$. Put
\begin{gather*}
R(x):=G_{Q(x)}(z),\qquad R^{(c)}(x):=\big(Q^{(c)}(x)-z\big)^{-1},\qquad r_j(x):=R^{(c)}_{jd}(x)\ \ (j\ne c),\\
T(x):=\Phi\big(X(Q^{(c)}(x))\big),
\end{gather*}
where $X(Q^{(c)})=X_{E_1,E_2,\eta}(Q^{(c)})$ is the functional of Definition~\ref{def:X} for the $(N-1)\times(N-1)$ matrix $Q^{(c)}$.

For $j\in\{1,\dots,N\}\setminus\{c\}$ let $q(j):=(\min(c,j),\max(c,j))\in\mathcal P_N$; the map $j\mapsto q(j)$ is injective, and $q(d)=p$. Put
\[
\mathcal J:=\big\{q(j):\ j\in\{1,\dots,N\}\setminus\{c,d\}\big\}\subset\mathcal P_N,\qquad\mathcal J^c:=\mathcal P_N\setminus\mathcal J .
\]
So $\mathcal J$ consists of the $N-2$ pairs which carry the entries of the row $c$ of $H$ other than $h_{cc}$ and $h_{cd}$; note $p\notin\mathcal J$ and $(c,c)\notin\mathcal J$.
As in Lemma~\ref{lem:split} we write $x=(x',y)$ with $x'=x_{\mathcal J^c}\in\K^{\mathcal J^c}$ and $y=x_{\mathcal J}\in\K^{\mathcal J}$, and $\P_{\mathcal J}=\bigotimes_{q\in\mathcal J}\nu_q$.

\emph{\textup{(i)} The minor does not depend on the coordinates in $\mathcal J$.} For $i,j\ne c$ the entry $Q_{ij}(x)$ is a function of the single coordinate $x_{(\min(i,j),\max(i,j))}$ (or is identically $0$),
and the pair $(\min(i,j),\max(i,j))$ does not contain $c$, whereas every pair in $\mathcal J$ contains $c$. Hence $Q^{(c)}(x)$, and therefore $R^{(c)}(x)$, $r_j(x)$ and $T(x)$, depend only on $x'$.
We write $r_j(x')$, $T(x')$.

\emph{\textup{(ii)} The row $c$.} For $j\notin\{c,d\}$: $Q_{cj}(x)=h_{cj}(x)=\xi_j(x_{q(j)})$, where $\xi_j(s):=s$ if $c<j$ and $\xi_j(s):=\bar s$ if $j<c$ $(s\in\K)$; and $Q_{cd}(x)=0$.

\emph{\textup{(iii)} Continuity and bounds.} $x\mapsto Q(x)$ and $x\mapsto Q^{(c)}(x)$ are linear. By Lemma~\ref{lem:basic}(a),(h) (for Hermitian matrices of size $N$ and $N-1$), the entries of $R$ and of $R^{(c)}$
are continuous functions of $x$, $R_{cc}\ne0$, and
\begin{equation}\label{eq:b3:rbound}
|r_j(x')|\le\|R^{(c)}\|\le\eta^{-1}\qquad(j\ne c).
\end{equation}
By Lemma~\ref{lem:Xspec}(c), $x\mapsto X(Q(x))$ and $x\mapsto X(Q^{(c)}(x))$ are continuous; so $T$ is continuous with $|T|\le C_\Phi$. The set $\{x:Q(x)\in\cG_p\}$ is closed (remark after Definition~\ref{def:goodB}).
Hence $\mathbf 1(Q\in\cG_p)\Phi(X(Q))R_{cd}(z)$ is a Borel function on $\Omega$, and it is bounded by $C_\Phi\Psi$, because $|R_{cd}(z)|\le\Psi$ on $\cG_p$.

\smallskip\noindent\emph{Step 1: Schur complement (exact).} By Lemma~\ref{lem:schur}\,\textup{(I2)}, applied to the Hermitian matrix $Q(x)$ with $i=c$, $j=d$ (first formula),
\begin{equation}\label{eq:b3:schur}
R_{cd}=-R_{cc}\,Y,\qquad Y(x):=\sum_{j\ne c}Q_{cj}(x)\,R^{(c)}_{jd}(x)=\sum_{j\notin\{c,d\}}\xi_j\big(x_{q(j)}\big)\,r_j(x')
\end{equation}
at every $x\in\Omega$, by (i) and (ii). Since $Y=-R_{cd}/R_{cc}$, the function $Y$ is continuous on $\Omega$.

\smallskip\noindent\emph{Step 2: $Y$ is centred given the minor, and $\E_M|Y|^2\le\eta^{-2}$.} Fix $x'\in\K^{\mathcal J^c}$. By Lemma~\ref{lem:split}(e), on the probability space $(\K^{\mathcal J},\P_{\mathcal J})$
the coordinate functions $y\mapsto y_q$, $q\in\mathcal J$, are independent with laws $\nu_q$. Hence the random variables $\Xi_j(y):=\xi_j(y_{q(j)})$, $j\notin\{c,d\}$ (functions of distinct coordinates),
are independent, and by the centring condition (Z) of Definition~\ref{def:ass} and the definition \eqref{eq:aver} of the variances,
\[
\int\Xi_j\,d\P_{\mathcal J}=\int\xi_j(s)\,\nu_{q(j)}(ds)=0,\qquad\int|\Xi_j|^2d\P_{\mathcal J}=\int|s|^2\nu_{q(j)}(ds)=\sigma_{cj}^2<\infty
\]
(for $\xi_j(s)=\bar s$ the first integral is the complex conjugate of $\int s\,\nu_{q(j)}(ds)=0$). For $j\ne j'$ the product $\Xi_j\overline{\Xi_{j'}}$ is integrable (Cauchy--Schwarz) and, by independence,
$\int\Xi_j\overline{\Xi_{j'}}\,d\P_{\mathcal J}=\int\Xi_j\,d\P_{\mathcal J}\cdot\overline{\int\Xi_{j'}\,d\P_{\mathcal J}}=0$. By \eqref{eq:b3:schur}, $Y(x',y)=\sum_{j\notin\{c,d\}}r_j(x')\,\Xi_j(y)$ with coefficients that do not depend on $y$. Therefore,
for \emph{every} $x'\in\K^{\mathcal J^c}$,
\begin{align}
\int Y(x',y)\,\P_{\mathcal J}(dy)&=\sum_{j\notin\{c,d\}}r_j(x')\int\Xi_j\,d\P_{\mathcal J}=0,\label{eq:b3:Ymean}\\
\int|Y(x',y)|^2\,\P_{\mathcal J}(dy)&=\sum_{j,j'\notin\{c,d\}}r_j(x')\overline{r_{j'}(x')}\int\Xi_j\overline{\Xi_{j'}}\,d\P_{\mathcal J}=\sum_{j\notin\{c,d\}}|r_j(x')|^2\sigma_{cj}^2\notag\\
&\le\ \eta^{-2}\sum_{j=1}^N\sigma_{cj}^2=\eta^{-2}.\label{eq:b3:Yvar}
\end{align}
In \eqref{eq:b3:Yvar} we used \eqref{eq:b3:rbound}, added the nonnegative terms $j=c,d$, and used (A) (with $\sigma_{cj}=\sigma_{jc}$). \emph{This is an inequality, not an identity:} the summands $j\in\{c,d\}$ are absent from $Y$,
and $|r_j|\le\eta^{-1}$ is only a bound.

By Lemma~\ref{lem:split}(b) (Tonelli) applied to $|Y|^2$ and by \eqref{eq:b3:Yvar},
\begin{equation}\label{eq:b3:Y2}
\E_M|Y|^2=\int P_{\mathcal J}\big(|Y|^2\big)\,d\P_M\le\eta^{-2}.
\end{equation}
In particular $Y\in L^2(\P_M)\subset L^1(\P_M)$. By \eqref{eq:b3:Yvar} and the Cauchy--Schwarz inequality, $\int|Y(x',y)|\P_{\mathcal J}(dy)\le\eta^{-1}<\infty$ for every $x'$, so the set $\mathcal D_Y$ of Lemma~\ref{lem:split}(c) is all of $\Omega$,
and $P_{\mathcal J}Y=0$ at every point by \eqref{eq:b3:Ymean}. The function $T$ is bounded, Borel and independent of the coordinates in $\mathcal J$ (Step~0). By Lemma~\ref{lem:split}(d) with $f=Y$, $g=T$,
\begin{equation}\label{eq:b3:TY}
\E_M\big[T\,Y\big]=\int T\cdot P_{\mathcal J}Y\,d\P_M=0 .
\end{equation}

\smallskip\noindent\emph{Step 3: on the good set (deterministic).} Let $x\in\Omega$ with $Q=Q(x)\in\cG_p$. By \eqref{eq:b1:end}, $|R_{cc}(z_i)+1|\le\Psi$ for $i=1,2$ and $|R_{cd}(z)|\le\Psi$, where $\Psi\le\frac18$ by \textup{(N4)}.

\emph{\textup{(3a)}} Corollary~\ref{cor:Xchange}(ii), applied to the matrix $Q$ and the pair $(c,c)$ (the statement involves only $Q$ and the index whose row and column are removed), gives, since $\frac1\pi<\frac13$,
\[
\big|X(Q)-X(Q^{(c)})\big|\le\tfrac2\pi\big(|R_{cc}(z_1)+1|+|R_{cc}(z_2)+1|\big)\le\tfrac4\pi\Psi\le\tfrac43\Psi .
\]
\emph{\textup{(3b)}} $|R_{cc}(z)|\ge1-|R_{cc}(z)+1|\ge1-\Psi\ge\frac78$, hence by \eqref{eq:b3:schur}: $|Y|=|R_{cd}(z)|/|R_{cc}(z)|\le\frac87\Psi$. \emph{No large deviation bound is used: the size of $Y$ is read off from the identity.}

\emph{\textup{(3c)}} By \eqref{eq:b3:schur},
\[
\Phi(X(Q))\,R_{cd}(z)-T\,Y=\rho\,Y,\qquad\rho:=\Phi(X(Q))\big(-R_{cc}(z)-1\big)+\Big(\Phi(X(Q))-\Phi\big(X(Q^{(c)})\big)\Big).
\]
By the mean value theorem $|\Phi(s)-\Phi(s')|\le C_\Phi|s-s'|$, so by (3a): $|\rho|\le C_\Phi\Psi+\frac43C_\Phi\Psi=\frac73C_\Phi\Psi$. With (3b),
\begin{equation}\label{eq:b3:good}
\big|\Phi(X(Q))\,R_{cd}(z)-T\,Y\big|\le\tfrac73\cdot\tfrac87\,C_\Phi\Psi^2=\tfrac83\,C_\Phi\Psi^2\qquad\text{on }\{x:\ Q(x)\in\cG_p\}.
\end{equation}

\smallskip\noindent\emph{Step 4: the bad set.} Write $\mathbf 1_{\cG}:=\mathbf 1(Q\in\cG_p)$ and $\mathbf 1_{\cG^c}:=1-\mathbf 1_{\cG}$. All functions below are $\P_M$-integrable (Step 0 and \eqref{eq:b3:Y2}), and
\[
\E_M\big[\mathbf 1_{\cG}\,\Phi(X(Q))R_{cd}(z)\big]=\E_M\big[\mathbf 1_{\cG}\big(\Phi(X(Q))R_{cd}(z)-TY\big)\big]+\E_M[TY]-\E_M\big[\mathbf 1_{\cG^c}\,TY\big].
\]
The first term is at most $\frac83C_\Phi\Psi^2$ in modulus by \eqref{eq:b3:good}; the second vanishes by \eqref{eq:b3:TY}. For the third, the Cauchy--Schwarz inequality, $|T|\le C_\Phi$, Lemma~\ref{lem:RfromS}(b) (for the ensemble $M$) and \eqref{eq:b3:Y2} give
\begin{align*}
\big|\E_M\big[\mathbf 1_{\cG^c}\,TY\big]\big|&\le C_\Phi\,\P_M\big(H^{[p]}\notin\cG_p\big)^{1/2}\big(\E_M|Y|^2\big)^{1/2}\le C_\Phi\,(2N^{-20})^{1/2}\,\eta^{-1}\\
&=\sqrt2\,C_\Phi\,N^{-10+\frac23+\eg}\le C_\Phi\Psi^2,
\end{align*}
because $-10+\frac23+\eg\le-9$ and $\sqrt2N^{-9}\le16N^{-2/3}\le16N^{-2/3+4\eg}=\Psi^2$. Altogether
\[
\Big|\E_M\big[\mathbf 1_{\cG}\,\Phi(X(Q))R_{cd}(z)\big]\Big|\le\big(\tfrac83+1\big)C_\Phi\Psi^2=\tfrac{11}3C_\Phi\Psi^2\le7C_\Phi\Psi^2.\qedhere
\]
\end{proof}

\noindent\emph{Remarks.} (1) \emph{Why Step 4 is needed.} The indicator $\mathbf 1(Q\in\cG_p)$ depends on the row $c$ of $Q$, so one cannot integrate out the row $c$ in $\E_M[\mathbf 1_{\cG}TY]$ directly;
the identity \eqref{eq:b3:TY} holds without the indicator, and the indicator is removed at the cost $\P(\cG^c)^{1/2}\|Y\|_{L^2}\le\sqrt2N^{-10}\eta^{-1}$. This is the only place of the comparison branch where an
unbounded function of the matrix is integrated over $\Omega$, and the only place where the trivial bound $\eta^{-1}$ on a resolvent is used.

(2) \emph{What is not used.} No large deviation estimate: the size of $Y$ on the good set is read off from $Y=-R_{cd}/R_{cc}$. No local law for minors: the resolvent $R^{(c)}$ of the minor enters only through
$\|R^{(c)}\|\le\eta^{-1}$ and through the exact identity of Corollary~\ref{cor:Xchange}(ii). No bound along the segment $[E_1,E_2]+i\eta$: only the endpoint condition \eqref{eq:b1:end}.

(3) \emph{Vanishing variances} are harmless: if $\sigma_{cj}=0$ then $\Xi_j=0$ almost surely and the summand disappears from \eqref{eq:b3:Yvar}. The step ``$S$ is the Green function of a generalized Wigner matrix
with all variances positive'' of the original (TeX 4259--4269) has no counterpart here.

(4) \emph{The gain is exactly one factor $\Psi$.} Both corrections in $\rho$ --- the diagonal entry $R_{cc}\ne-1$, and $X$ of the matrix versus $X$ of the minor --- are of order $\Psi$ and are not centred. The argument
therefore gains exactly one factor $\Psi$ over the trivial bound $C_\Phi\Psi$, which is what the power counting of Proposition~\ref{prop:step} needs.

(5) \emph{Real and complex case.} For $\K=\R$ all $\xi_j$ are the identity; for $\K=\C$ the conjugations in $\xi_j$ play no role, since only $\int\Xi_j=0$ and $\int|\Xi_j|^2=\sigma_{cj}^2$ are used.
The second Schur formula of \textup{(I2)} is not needed: the entry $R_{ba}$ is treated by removing the index $b$.

(6) \emph{Numerical check.} The identities $R_{cd}=-R_{cc}Y$, $X(Q)-X(Q^{(c)})=\frac1\pi\im\br{\Log(-R_{cc})}$ (for $c=a$ and $c=b$), the conditional statements \eqref{eq:b3:Ymean}, \eqref{eq:b3:Yvar}
(by exact enumeration over two-point laws on a small matrix) and the deterministic bounds of Step~3 (on synthetic Hermitian matrices which satisfy \eqref{eq:b1:end}) were checked numerically.

\begin{proposition}[one Lindeberg step]\label{prop:step}
Assume the standing assumptions of \S\ref{sec:step} \textup{(}in particular $\eg\in(0,\frac1{10}]$, $N\ge N_1(\Par,\eg)$, $\eta=N^{-2/3-\eg}$, $|E_i-2|\le N^{-2/3+\eg}$ and $F\in C^3(\R)$ with the constant $C_F$\textup{)}.
Let $M^v,M^w$ be ensembles of size $N$ over $\K$ which satisfy $\fA(\Par)$ and have matching second moments \eqref{eq:2m}; let $M_0,\dots,M_\Gamma$ be the hybrids of Definition~\ref{def:hyb}
(for any bijection $\varphi$), $1\le k\le\Gamma$ and $p_k=(a,b)$. Then, with
\begin{equation}\label{eq:b3:C3}
C_3=C_3(\vartheta,C_0):=2C_2+396\,\mathfrak m_3C_0^{3/2}+8
\end{equation}
\textup{(}$C_2$ from \eqref{eq:b3:C2}\textup{)},
\[
\big|\E_{M_{k-1}}F(X(H))-\E_{M_k}F(X(H))\big|\le\begin{cases}C_3\,C_F\,N^{-13/6+4\eg}=C_3C_F\,N^{-2}\cdot N^{-1/6+4\eg},&a\ne b,\\[2pt] C_3\,C_F\,N^{-11/6+2\eg},&a=b.\end{cases}
\]
\end{proposition}
\begin{proof}
Write $p:=p_k$. By Lemma~\ref{lem:hyb}(ii), $M_{k-1}$ and $M_k$ satisfy $\fA(\Par)$, and we use its notation: $x=(x',x_p)$, $\P'=\bigotimes_{q\ne p}\nu^{(k)}_q$ on $\K^{\mathcal P_N\setminus\{p\}}$, $Q(x')=H(x)^{[p]}$.
In this proof $(c,d)$ denotes an element of $\{(a,b),(b,a)\}$.

\smallskip\noindent\emph{Step 1: reduction to one entry.} The function $A\mapsto F(X(A))$ is continuous (Lemma~\ref{lem:Xspec}(c)) and bounded by $C_F$ on the Hermitian $N\times N$ matrices. By \eqref{eq:B2step},
\begin{gather*}
\E_{M_{k-1}}F(X(H))-\E_{M_k}F(X(H))=\int\big[I^v(x')-I^w(x')\big]\,\P'(dx'),\\
I^u(x'):=\int_\K F\big(X(Q(x')+V_p(s))\big)\,\nu^u_p(ds)\qquad(u\in\{v,w\}),
\end{gather*}
and $I^v,I^w$ are Borel functions of $x'$. Since $\nu^v_p,\nu^w_p$ are probability measures, subtracting the same number $F(X(Q(x')))$ from both gives
\begin{equation}\label{eq:b3:DD}
I^v-I^w=D^v-D^w,\qquad D^u(x'):=\mathfrak D\big(Q(x');\nu^u_p\big)=\int_\K\Big[F\big(X(Q(x')+V_p(s))\big)-F\big(X(Q(x'))\big)\Big]\nu^u_p(ds),
\end{equation}
and $|D^u|\le2C_F$.
$D^u$ is Borel, because $I^u$ is and $x'\mapsto F(X(Q(x')))$ is continuous. Let $\cG':=\{x':\ Q(x')\in\cG_p\}$; it is closed, because $x'\mapsto Q(x')$ is linear and $\cG_p$ is closed (Definition~\ref{def:goodB}).

\smallskip\noindent\emph{Step 2: the bad set.} By the last statement of Lemma~\ref{lem:hyb}(ii) (with $\Theta=\mathbf 1(\,\cdot\notin\cG_p)$) and Lemma~\ref{lem:RfromS}(b) for the ensemble $M_{k-1}$,
\begin{equation}\label{eq:b3:bad}
\P'\big((\cG')^c\big)=\P_{M_{k-1}}\big(H^{[p]}\notin\cG_p\big)\le2N^{-20},\quad\text{hence}\quad\Big|\int_{(\cG')^c}\big(D^v-D^w\big)\,d\P'\Big|\le8C_FN^{-20},
\end{equation}
because $|D^v-D^w|\le4C_F$.

\smallskip\noindent\emph{Step 3: the good set, second order.} Each of the laws $\nu\in\{\nu^v_p,\nu^w_p\}$ satisfies the hypotheses of Corollary~\ref{cor:entry}: it is a Borel probability measure on $\K$, with $\nu(\R)=1$ if $a=b$ (R),
$\int y\,d\nu=0$ (centring), and \eqref{eq:b3:nu} holds by (C) and \eqref{eq:subexp} (for $M^v$, resp.\ $M^w$). By \eqref{eq:2m} with $(l,u)=(1,1)$ and $(0,2)$, the numbers $m_{11}$ and $m_{20}$ are the same for $\nu^v_p$ and
$\nu^w_p$; hence, by \eqref{eq:b3:Q2}, \eqref{eq:b3:Q2d},
\begin{equation}\label{eq:b3:Q2eq}
\mathfrak Q_2(Q;\nu^v_p)=\mathfrak Q_2(Q;\nu^w_p)\qquad\text{for every }Q\in\cG_p .
\end{equation}

\emph{Case $a=b$.} For $x'\in\cG'$, \eqref{eq:b3:entryd} and \eqref{eq:b3:Q2eq} give $|D^v(x')-D^w(x')|\le2C_2C_FN^{-11/6+2\eg}$. With \eqref{eq:b3:bad}, $\P'(\cG')\le1$ and $N^{-20}\le N^{-11/6+2\eg}$,
\[
\big|\E_{M_{k-1}}F(X(H))-\E_{M_k}F(X(H))\big|\le(2C_2+8)\,C_FN^{-11/6+2\eg}\le C_3C_FN^{-11/6+2\eg}.
\]

\emph{Case $a\ne b$.} For $u\in\{v,w\}$ and $x'\in\K^{\mathcal P_N\setminus\{p\}}$ put $e^u(x'):=D^u(x')-\mathfrak Q_2(Q(x');\nu^u_p)-F'(X(Q(x')))\,\mathfrak T_3(Q(x');\nu^u_p)$, where $\mathfrak Q_2,\mathfrak T_3$ are \emph{defined} by the right sides of
\eqref{eq:b3:Q2}, \eqref{eq:b3:T3} for every Hermitian $Q$. These are Borel functions (Corollary~\ref{cor:entry}, last sentence), and $|e^u|\le C_2C_FN^{-13/6+4\eg}$ on $\cG'$ by \eqref{eq:b3:entry}. By \eqref{eq:b3:Q2eq}, on $\cG'$,
\[
D^v-D^w=e^v-e^w+F'(X(Q))\,\mathfrak T_3(Q;\nu^v_p)-F'(X(Q))\,\mathfrak T_3(Q;\nu^w_p),
\]
so that
\begin{equation}\label{eq:b3:good2}
\Big|\int_{\cG'}\big(D^v-D^w\big)\,d\P'\Big|\le2C_2C_FN^{-13/6+4\eg}+\sum_{u\in\{v,w\}}\Big|\int_{\cG'}F'(X(Q))\,\mathfrak T_3(Q;\nu^u_p)\,d\P'\Big| .
\end{equation}

\smallskip\noindent\emph{Step 4: the third-order term ($a\ne b$).} Fix $\nu\in\{\nu^v_p,\nu^w_p\}$ and let $\mu_3=\int|y|^2y\,d\nu$, $|\mu_3|\le\mathfrak m_3C_0^{3/2}N^{-3/2}$. For $(c,d)\in\{(a,b),(b,a)\}$ and $i\in\{1,2\}$ put
\[
\mathfrak E_{cd,i}:=\int\mathbf 1_{\cG'}(x')\,F'(X(Q))\,R_{cd}(z_i)\,g(z_i)\,\P'(dx'),\qquad R=G_{Q(x')} .
\]
The integrand is Borel, and bounded by $C_F\cdot\Psi\cdot\frac43$ (on $\cG'$: $|R_{cd}(z_i)|\le\Psi$ by \eqref{eq:b1:end} and $|g(z_i)|\le\frac43$ by Lemma~\ref{lem:expand}(a)). By \eqref{eq:b3:T3}, written out,
$\mathfrak T_3(Q;\nu)=-\frac1\pi\im\sum_{i=1,2}(-1)^i\big(\mu_3R_{ba}(z_i)g(z_i)+\overline{\mu_3}R_{ab}(z_i)g(z_i)\big)$. Since $\mathbf 1_{\cG'}F'(X(Q))$ is real and $\im$ is $\R$-linear,
\begin{equation}\label{eq:b3:T3int}
\int_{\cG'}F'(X(Q))\,\mathfrak T_3(Q;\nu)\,d\P'=-\frac1\pi\,\im\sum_{i=1,2}(-1)^i\Big(\mu_3\,\mathfrak E_{ba,i}+\overline{\mu_3}\,\mathfrak E_{ab,i}\Big).
\end{equation}
Write $g=1+(g-1)$. By the last statement of Lemma~\ref{lem:hyb}(ii) (applied to the real and imaginary parts of the bounded Borel function $\Theta(Q):=\mathbf 1(Q\in\cG_p)F'(X(Q))G_Q(z_i)_{cd}$),
\[
\mathfrak E_{cd,i}=\E_{M_{k-1}}\big[\mathbf 1(Q\in\cG_p)\,F'(X(Q))\,R_{cd}(z_i)\big]+\int\mathbf 1_{\cG'}\,F'(X(Q))\,R_{cd}(z_i)\,\big(g(z_i)-1\big)\,d\P' .
\]
The first term is at most $7C_F\Psi^2$ in modulus by Lemma~\ref{lem:decor}, applied to the ensemble $M=M_{k-1}$ and to $\Phi:=F'$ (which is $C^1$ with $\sup|F'|,\sup|F''|\le C_F$, so $C_\Phi=C_F$). The second term is at most
$C_F\cdot\Psi\cdot\frac94\Psi$, by $|R_{cd}(z_i)|\le\Psi$ and $|g(z_i)-1|\le\frac94\Psi$ on $\cG'$ (Lemma~\ref{lem:expand}(a)). Hence $|\mathfrak E_{cd,i}|\le\frac{37}4C_F\Psi^2$ for the four pairs $(cd,i)$, and by \eqref{eq:b3:T3int},
$\frac1\pi<\frac13$ and $\Psi^2=16N^{-2/3+4\eg}$,
\begin{equation}\label{eq:b3:third}
\begin{aligned}
\Big|\int_{\cG'}F'(X(Q))\,\mathfrak T_3(Q;\nu)\,d\P'\Big|&\le\frac1\pi\cdot4|\mu_3|\cdot\frac{37}4\,C_F\Psi^2\le\frac{37}3\,\mathfrak m_3C_0^{3/2}N^{-3/2}\cdot16N^{-2/3+4\eg}\,C_F\\
&\le198\,\mathfrak m_3C_0^{3/2}\,C_F\,N^{-13/6+4\eg}.
\end{aligned}
\end{equation}

\smallskip\noindent\emph{Step 5: conclusion ($a\ne b$).} By \eqref{eq:b3:DD}, \eqref{eq:b3:bad}, \eqref{eq:b3:good2}, \eqref{eq:b3:third} (for $\nu^v_p$ and for $\nu^w_p$) and $N^{-20}\le N^{-13/6+4\eg}$,
\[
\big|\E_{M_{k-1}}F(X(H))-\E_{M_k}F(X(H))\big|\le\big(8+2C_2+2\cdot198\,\mathfrak m_3C_0^{3/2}\big)C_FN^{-13/6+4\eg}=C_3C_FN^{-13/6+4\eg}.\qedhere
\]
\end{proof}

\noindent\emph{Remarks.} (1) \textbf{Power counting per off-diagonal step:} $N^{-3/2}$ (third moment of the entry) $\times$ $\Psi\approx N^{-1/3+2\eg}$ (one off-diagonal resolvent entry in $\mathfrak T_3$) $\times$
$\Psi\approx N^{-1/3+2\eg}$ (decorrelation gain, Lemma~\ref{lem:decor}; or the second factor $\Psi$ of the error terms of Lemma~\ref{lem:expand}) $=N^{-13/6+4\eg}=N^{-2}\cdot N^{-1/6+4\eg}$. This is the exponent of
Lemma~6.5 of the original, with the explicit constant $4$ in front of $\eg$. For a diagonal step no decorrelation is needed: $N^{-3/2}\cdot\Psi\approx N^{-11/6+2\eg}$, and there are only $N$ such steps.

(2) \emph{Constants.} $C_3$ depends on $(\vartheta,C_0)$ only, through $\mathfrak m_3,\mathfrak m_4,\mathfrak m_6$ and $C_0$; the parameters $\delta_\pm$ enter only through the threshold $N_1(\Par,\eg)$ of the standing assumptions.

(3) \emph{Which hybrid.} The step $k$ uses Lemma~\ref{lem:RfromS}(b) and Lemma~\ref{lem:decor} for the single ensemble $M_{k-1}$ (through the common law $\P'$ of $Q$; one could equally use $M_k$). The laws $\nu^v_{p}$, $\nu^w_{p}$
enter only through Corollary~\ref{cor:entry}.

(4) \emph{Trivial steps.} If $\sigma_{p_k}=0$ then $\nu^v_{p_k}=\nu^w_{p_k}=\delta_0$, $M_{k-1}=M_k$, and the left side is $0$.

\subsection{Green function comparison at the edge}\label{sec:gfc}
\begin{theorem}[edge comparison; the form (6.25) of Theorem~6.3 of the original]\label{thm:gfc}
For every $\Par=(\vartheta,\delta_-,\delta_+,C_0)$ and every $\eg\in(0,\frac1{10}]$ there is $C_4=C_4(\vartheta,C_0)$ such that for all
$N\ge N_1(\Par,\eg)$ \textup{(}the threshold of Lemma~\ref{lem:N0}\textup{)}, $\K\in\{\R,\C\}$, all ensembles $M^v,M^w$ of size $N$ over $\K$ with $\fA(\Par)$ and matching second moments \eqref{eq:2m}, all real $E_1\le E_2$ with
$|E_i-2|\le N^{-2/3+\eg}$ $(i=1,2)$, $\eta=N^{-2/3-\eg}$, and all $F\in C^3(\R)$ with $C_F=\max_{0\le k\le3}\sup_\R|F^{(k)}|<\infty$:
\begin{equation}\label{eq:b3:gfc}
\Big|\E^vF\Big(\frac N\pi\int_{E_1}^{E_2}\im m(y+i\eta)\,dy\Big)-\E^wF\Big(\frac N\pi\int_{E_1}^{E_2}\im m(y+i\eta)\,dy\Big)\Big|\le C_4\,C_F\,N^{-1/6+4\eg}.
\end{equation}
One can take
\begin{equation}\label{eq:b3:C4}
C_4:=\tfrac32\,C_3=3C_2+594\,\mathfrak m_3C_0^{3/2}+12,
\end{equation}
with $C_2$, $C_3$ from \eqref{eq:b3:C2}, \eqref{eq:b3:C3}.
\end{theorem}
\begin{proof}
Let $N\ge N_1(\Par,\eg)$ and let $\K,M^v,M^w,E_1,E_2,\eta,F$ be as in the statement. Then the standing assumptions of \S\ref{sec:step} hold (with $\Psi=4N^{-1/3+2\eg}$, $t=N^{-1/2+2\eg}$).
By Definition~\ref{def:X} and \eqref{eq:green} ($\tr G_H=Nm$),
\[
X(H)=X_{E_1,E_2,\eta}(H)=\frac1\pi\int_{E_1}^{E_2}\im\tr G_H(y+i\eta)\,dy=\frac N\pi\int_{E_1}^{E_2}\im m(y+i\eta)\,dy ,
\]
so the left side of \eqref{eq:b3:gfc} is $|\E^vF(X(H))-\E^wF(X(H))|$; the function $x\mapsto F(X(H(x)))$ is continuous and bounded on $\Omega$ (Lemma~\ref{lem:Xspec}(c)), so all expectations exist.

Fix a bijection $\varphi:\mathcal P_N\to\{1,\dots,\Gamma\}$, $\Gamma=N(N+1)/2$, and let $M_0=M^v,M_1,\dots,M_\Gamma=M^w$ be the hybrids of Definition~\ref{def:hyb}. Telescoping,
\[
\E^vF(X(H))-\E^wF(X(H))=\sum_{k=1}^{\Gamma}\Big(\E_{M_{k-1}}F(X(H))-\E_{M_k}F(X(H))\Big).
\]
Since $\varphi$ is a bijection, exactly $N$ of the pairs $p_1,\dots,p_\Gamma$ are diagonal $(a=b)$ and exactly $N(N-1)/2$ are off-diagonal. By Proposition~\ref{prop:step}, applied to each $k$,
\begin{align*}
\big|\E^vF(X(H))-\E^wF(X(H))\big|&\le\frac{N(N-1)}2\cdot C_3C_FN^{-13/6+4\eg}+N\cdot C_3C_FN^{-11/6+2\eg}\\
&\le\tfrac12C_3C_FN^{-1/6+4\eg}+C_3C_FN^{-5/6+2\eg}\le\tfrac32\,C_3C_F\,N^{-1/6+4\eg},
\end{align*}
because $2-\frac{13}6=-\frac16$, $1-\frac{11}6=-\frac56$ and $N^{-5/6+2\eg}\le N^{-1/6+4\eg}$. The constant $C_4=\frac32C_3$ depends only on $(\vartheta,C_0)$ by \eqref{eq:b3:C2}, \eqref{eq:b3:C3}, and $N_1(\Par,\eg)$ depends only on
$(\Par,\eg)$ (Lemma~\ref{lem:N0}). Only the bounds $\sup|F^{(k)}|\le C_F$, $k=0,1,2,3$, are used: in Lemma~\ref{lem:expand} (hypothesis \textup{(D3)}), Corollary~\ref{cor:entry}, Lemma~\ref{lem:decor} (with $\Phi=F'$) and
Proposition~\ref{prop:step}.
\end{proof}

\noindent\emph{Remarks.} (1) \emph{Uniformity over the $N(N+1)/2$ hybrids.} Proposition~\ref{prop:step} holds for each $k$ with the \emph{same} constant $C_3$ and the same threshold $N_1(\Par,\eg)$, because every
hybrid $M_k$ satisfies $\fA(\Par)$ with the same $\Par$ (Lemma~\ref{lem:hyb}(ii)) and the constants of Theorem~\ref{thm:21} depend only on $\Par$. No union bound over the hybrids is taken, and no supremum over $k$
appears inside a probability: step $k$ bounds a difference of two expectations, using the event $\{Q\in\cG_{p_k}\}$ under the single law $\P_{M_{k-1}}$, whose complement costs $O(N^{-20})$ in that step. The sum of
these costs over all steps is at most $8C_F\,\Gamma\,N^{-20}\le8C_FN^{-18}$, which is included in the bound.

(2) \emph{Exponents and constants.} The exponent $-\frac16+4\eg$ is absolute: the number $4$ in front of $\eg$ and the constant $\frac16$ do not depend on $\Par$. The constant $C_4$ depends on $(\vartheta,C_0)$ only;
$\delta_\pm$ enter only through $N_1(\Par,\eg)$ (via the constants $A_0,C,c,\phi,N_0$ of Theorem~\ref{thm:21}). This is what the quantifier order of Theorem~\ref{thm:24} requires: in \S\ref{sec:proof24} the theorem is used with
$\eg=\frac1{50}$, where $-\frac16+4\eg=-\frac{13}{150}<-\frac1{20}$. Total: $N^2\cdot N^{-3/2}\cdot\Psi^2\approx N^{-1/6+4\eg}$ from the off-diagonal steps, $N\cdot N^{-3/2}\cdot\Psi\approx N^{-5/6+2\eg}$ from the diagonal ones.

(3) \emph{The order of replacement} $\varphi$ is arbitrary; the bound does not depend on it.

(4) \emph{Scope.} Only the integrated form (6.25) of Theorem~6.3 of the original is proved, for bounded $F$ with bounded derivatives; this is what \S\ref{sec:proof24} uses (with $F=q$). The pointwise form (6.24)
and polynomially growing $F$ are not needed for Theorem~\ref{thm:24} and are not treated (see Appendix~\ref{app:diffB}).

\medskip\noindent\textbf{Dependencies of the comparison branch} (``$\Leftarrow$'' means ``uses''; nothing else is used).
Lemma~\ref{lem:edge} $\Leftarrow$ Theorem~\ref{thm:21}(ii) in the form \textup{(BB1)}, Lemma~\ref{lem:mscB}(b), Lemma~\ref{lem:N0}.
Lemma~\ref{lem:RfromS}(b) $\Leftarrow$ Lemma~\ref{lem:edge}, Lemma~\ref{lem:RfromS}(a), \eqref{eq:subexp}, (C).
Corollary~\ref{cor:entry} $\Leftarrow$ Lemma~\ref{lem:expand}, Lemma~\ref{lem:moments}(a), \eqref{eq:subexp}.
Lemma~\ref{lem:decor} $\Leftarrow$ Lemma~\ref{lem:schur}\,(I2), Lemma~\ref{lem:split}, Corollary~\ref{cor:Xchange}(ii), Lemma~\ref{lem:RfromS}(b), centring, (A).
Proposition~\ref{prop:step} $\Leftarrow$ Lemma~\ref{lem:hyb}(ii), Corollary~\ref{cor:entry}, Lemma~\ref{lem:decor}, Lemma~\ref{lem:RfromS}(b), Lemma~\ref{lem:expand}(a), \eqref{eq:2m}.
Theorem~\ref{thm:gfc} $\Leftarrow$ Proposition~\ref{prop:step}.
The comparison branch uses Theorem~\ref{thm:21}(ii) only (through Lemma~\ref{lem:edge}) and never Theorem~\ref{thm:22}; it meets the rigidity branch (\S\ref{sec:sandwich}) only in Proposition~\ref{prop:oneside}.
Every function that is integrated over $\Omega$ (or over the law $\P'$ of $Q$) is bounded, except $Y$ in Lemma~\ref{lem:decor}, which is in $L^2$; the integrals over a single entry law $\nu$
in Corollary~\ref{cor:entry} involve polynomials of degree $\le3$ in the entry, and its sixth moment in the tail estimate. Every event is an open or a closed subset of $\Omega$.

\subsection{Proof of Theorem~\ref{thm:24B}}\label{sec:proof24}
Throughout this subsection
\[
\eg:=\tfrac1{50},\qquad\e=\tfrac12\eg=\tfrac1{100},\qquad\delta:=\tfrac1{20},
\]
and $L_0,W,E_L,\eta,\ell$ are the parameters \eqref{eq:B2par}, \eqref{eq:B2par2} for this $\eg$; in particular $\eta=N^{-2/3-\eg}$ is the same in
Proposition~\ref{prop:sandwich} and in Theorem~\ref{thm:gfc}.

\paragraph{The cutoff.} Let $\mathsf S_4:\R\to\R$ be the polynomial step
\[
\mathsf S_4(t):=0\ \ (t\le0),\quad\mathsf S_4(t):=t^5\big(126-420t+540t^2-315t^3+70t^4\big)\ \ (0\le t\le1),\quad\mathsf S_4(t):=1\ \ (t\ge1).
\]
On $[0,1]$ one has $\mathsf S_4'(t)=630\,t^4(1-t)^4\ge0$ (expand both sides) and $\mathsf S_4(1)=126-420+540-315+70=1$; all derivatives of order $1,\dots,4$ of the polynomial vanish at $t=0$ and at $t=1$
(because $t^4(1-t)^4$ has zeros of order $4$ there), so $\mathsf S_4\in C^4(\R)$, $\mathsf S_4$ is nondecreasing and $0\le\mathsf S_4\le1$. Fix once and for all
\[
q(x):=1-\mathsf S_4(9x-1)\qquad(x\in\R),\qquad C_q:=\max_{0\le k\le3}\sup_{x\in\R}|q^{(k)}(x)|.
\]
Then $q\in C^4(\R)$, $0\le q\le1$, $q(x)=1$ for $x\le\frac19$ (as $9x-1\le0$) and $q(x)=0$ for $x\ge\frac29$ (as $9x-1\ge1$). $C_q$ is an explicit absolute constant:
$|q|\le1$, and for $1\le k\le3$, $q^{(k)}(x)=-9^k\mathsf S_4^{(k)}(9x-1)$, where $\mathsf S_4^{(k)}$ vanishes outside $[0,1]$ and on $[0,1]$ is a polynomial whose modulus is at most the sum of the moduli of its coefficients;
for $k=3$ this sum is $7560+50400+113400+105840+35280=312480$, and it is smaller for $k=1,2$, so $C_q\le9^3\cdot312480<2.3\cdot10^8$ (numerically $C_q=9^3\cdot\frac{315}4\approx5.7\cdot10^4$, the value of $9^3|\mathsf S_4'''|$ at $t=\frac12$).
Thus $q$ is an admissible function in
Proposition~\ref{prop:sandwich} (it is continuous, hence Borel) and an admissible $F$ in Theorem~\ref{thm:gfc}, with $C_F=C_q$. Monotonicity of $q$ is not used.
(Any $q\in C^3(\R)$ with the same three properties would do, for instance $q(x)=1-\mathsf S(9x-1)$ with the smooth transition function $\mathsf S$ of Mathlib; then $C_q$ is finite by compactness, but not explicit.)

\paragraph{The threshold.} Let $N_1=N_1(\Par,\frac1{50})$ be the threshold of Lemma~\ref{lem:N0} (it is also the threshold of Theorem~\ref{thm:gfc}), and let $C_4=C_4(\vartheta,C_0)$ be the
constant of Theorem~\ref{thm:gfc} for $\eg=\frac1{50}$. Define
\begin{equation}\label{eq:B2N0}
N_0^{\rm edge}(\Par):=\max\Big\{N_1,\ (2C_4C_q)^{300/11}\Big\}.
\end{equation}
It depends only on $\Par$ (through the constants $A_0,C,c,\phi,N_0$ of Part~I, through $\vartheta,C_0$, and through $N_1,C_4$), and not on $s$, $\K$ or the ensembles.
For $N\ge N_0^{\rm edge}(\Par)$ the conditions (N1)--(N5) of Lemma~\ref{lem:N0} hold with $\eg=\frac1{50}$, and in addition
\begin{itemize}
\item[\textup{(N6)}] $C_4C_qN^{-1/6+4\eg}+2N^{-20}\le N^{-\delta}$.
\end{itemize}
Indeed $-\frac16+4\eg=-\frac{13}{150}=-\frac1{20}-\frac{11}{300}$, so $C_4C_qN^{-1/6+4\eg}=C_4C_qN^{-11/300}\cdot N^{-\delta}\le\frac12N^{-\delta}$ because
$N^{11/300}\ge2C_4C_q$; and $2N^{-20}\le\frac12N^{-\delta}$ because $N^{20-\delta}\ge4$ (as $N\ge N_1\ge e^{27}$).

\begin{proposition}[one-sided comparison]\label{prop:oneside}
Let $N\ge N_0^{\rm edge}(\Par)$, $\K\in\{\R,\C\}$, and let $M^v,M^w$ be ensembles of size $N$ over $\K$ which satisfy $\fA(\Par)$ and have matching second moments
\eqref{eq:2m}. Then for every $s\in\R$:
\[
\P^v\big(N^{2/3}(\lambda_N-2)\le s-N^{-\e}\big)\le\P^w\big(N^{2/3}(\lambda_N-2)\le s\big)+N^{-\delta}.
\]
\end{proposition}
\begin{proof}
For $r\in\R$ the event $\{x:N^{2/3}(\lambda_N(H(x))-2)\le r\}$ is closed by Lemma~\ref{lem:weyl}(d), so all probabilities below are defined.
Let $\cR^v,\cR^w$ be the event $\cR$ of the proof of Proposition~\ref{prop:sandwich} (it does not depend on the ensemble), considered under $\P^v$ and under $\P^w$;
by \eqref{eq:B2Rc}, $\P^v(\Omega\setminus\cR)\le N^{-20}$ and $\P^w(\Omega\setminus\cR)\le N^{-20}$, and on $\cR$, by (R1), $N^{2/3}|\lambda_N-2|<(\log N)^{L_0}$.
Note $N^{-20}\le N^{-\delta}$.

\emph{Case A: $s>(\log N)^{L_0}$.} On $\cR$, $N^{2/3}(\lambda_N-2)<(\log N)^{L_0}<s$. Hence
$\P^w(N^{2/3}(\lambda_N-2)\le s)+N^{-\delta}\ge\P^w(\cR)+N^{-20}\ge1$, which is at least the left side.

\emph{Case B: $s<-(\log N)^{L_0}$.} On $\cR$, $N^{2/3}(\lambda_N-2)>-(\log N)^{L_0}>s>s-N^{-\e}$. Hence the left side is at most
$\P^v(\Omega\setminus\cR)\le N^{-20}\le N^{-\delta}$, which is at most the right side.

\emph{Case C: $|s|\le(\log N)^{L_0}$.} Put $E:=2+sN^{-2/3}$; then $|E-2|\le W$. By \eqref{eq:B2etaell}, $3\ell=N^{-2/3-\e}$, so
$\{N^{2/3}(\lambda_N-2)\le s-N^{-\e}\}=\{\lambda_N\le E-3\ell\}\subset\{\lambda_N<E-2\ell\}$. The point $E':=E-2\ell$ satisfies $2-2W\le2-W-2\ell\le E'<E\le2+W$ (by (N5), $2\ell\le W$),
and $E'+\ell=E-\ell$, so $X^+_{E'}=X_{E-\ell,E_L,\eta}=X^-_E$. Put $E_1:=E-\ell$, $E_2:=E_L$. Then $E_1<E_2$ (Proposition~\ref{prop:sandwich}),
$|E_1-2|\le W+\ell\le2W\le N^{-2/3+\eg}$ and $|E_2-2|=2W\le N^{-2/3+\eg}$ by (N2) and (N5). Moreover, by Definition~\ref{def:X} and \eqref{eq:green},
$X_{E_1,E_2,\eta}(H)=\frac1\pi\int_{E_1}^{E_2}\im\tr G_H(y+i\eta)\,dy=\frac N\pi\int_{E_1}^{E_2}\im m(y+i\eta)\,dy$, which is the argument of $F$ in Theorem~\ref{thm:gfc}. Therefore
\begin{align*}
\P^v\big(N^{2/3}(\lambda_N-2)\le s-N^{-\e}\big)&\le\P^v(\lambda_N<E')\\
&\le\E^v\,q\big(X^+_{E'}(H)\big)+N^{-20}&&\text{(a)}\\
&=\E^v\,q\big(X_{E_1,E_2,\eta}(H)\big)+N^{-20}&&\text{(b)}\\
&\le\E^w\,q\big(X_{E_1,E_2,\eta}(H)\big)+C_4C_qN^{-1/6+4\eg}+N^{-20}&&\text{(c)}\\
&=\E^w\,q\big(X^-_{E}(H)\big)+C_4C_qN^{-1/6+4\eg}+N^{-20}&&\text{(b)}\\
&\le\P^w(\lambda_N<E)+C_4C_qN^{-1/6+4\eg}+2N^{-20}&&\text{(d)}\\
&\le\P^w\big(N^{2/3}(\lambda_N-2)\le s\big)+N^{-\delta}&&\text{(e)}
\end{align*}
Here (a) is the upper bound of Proposition~\ref{prop:sandwich} for $M^v$ at the point $E'$; (b) is $X^+_{E'}=X_{E_1,E_2,\eta}=X^-_E$; (c) is Theorem~\ref{thm:gfc} with
$F=q$; (d) is the lower bound of Proposition~\ref{prop:sandwich} for $M^w$ at the point $E$; (e) is (N6) together with $\{\lambda_N<E\}\subset\{\lambda_N\le E\}$.
Theorem~\ref{thm:gfc} applies because $N\ge N_1$, $M^v,M^w$ satisfy $\fA(\Par)$ and \eqref{eq:2m}, $\eta=N^{-2/3-\eg}$, $E_1\le E_2$, $|E_i-2|\le N^{-2/3+\eg}$, and $q\in C^3(\R)$
with $C_F=C_q$. Proposition~\ref{prop:sandwich} applies because $N\ge N_1$, $E,E'\in[2-2W,2+W]$, and $M^v$, $M^w$ satisfy $\fA(\Par)$.
\end{proof}

\begin{proof}[Proof of Theorem~\ref{thm:24B}]
Let $\Par$ be given, $N\ge N_0^{\rm edge}(\Par)$, $\K\in\{\R,\C\}$, and let $M^v,M^w$ be as in the theorem; let $s\in\R$.

\emph{The first inequality of \eqref{eq:tw}} is Proposition~\ref{prop:oneside}.

\emph{The second inequality of \eqref{eq:tw}.} The hypotheses of Proposition~\ref{prop:oneside} are symmetric in $(M^v,M^w)$ (\eqref{eq:2m} is an equality). Applied to the pair
$(M^w,M^v)$ at the point $s+N^{-\e}$ it gives $\P^w(N^{2/3}(\lambda_N-2)\le s)\le\P^v(N^{2/3}(\lambda_N-2)\le s+N^{-\e})+N^{-\delta}$.

\emph{The inequalities \eqref{eq:tw1} for $\lambda_1$.} By Lemma~\ref{lem:hyb}(iii), the pair $(-M^v,-M^w)$ satisfies the hypotheses of the theorem (same $\Par$, same $N$, same $\K$),
so \eqref{eq:tw} holds for it, at every point; we use it at the point $-s$. For $r\in\R$ let $S_r$ be the set of Hermitian $N\times N$ matrices $A$ with
$N^{2/3}(\lambda_N(A)-2)\le r$; it is closed, because $|\lambda_N(A)-\lambda_N(A')|\le\|A-A'\|$ (Lemma~\ref{lem:weyl}(b)). By Lemma~\ref{lem:weyl}(c), $\lambda_N(-A)=-\lambda_1(A)$, so
\[
-A\in S_r\iff N^{2/3}\big(-\lambda_1(A)-2\big)\le r\iff N^{2/3}\big(\lambda_1(A)+2\big)\ge-r .
\]
Since the law of $H$ under $\P_{-M}$ is the law of $-H$ under $\P_M$ (Lemma~\ref{lem:hyb}(iii)), for $M\in\{M^v,M^w\}$ and every $r\in\R$:
\begin{equation}\label{eq:B2refl}
\P_{-M}\big(N^{2/3}(\lambda_N-2)\le r\big)=\P_{-M}(H\in S_r)=\P_M(-H\in S_r)=\P_M\big(N^{2/3}(\lambda_1+2)\ge-r\big).
\end{equation}
The first inequality of \eqref{eq:tw} for $(-M^v,-M^w)$ at $-s$ reads
$\P_{-M^v}(N^{2/3}(\lambda_N-2)\le-s-N^{-\e})-N^{-\delta}\le\P_{-M^w}(N^{2/3}(\lambda_N-2)\le-s)$; by \eqref{eq:B2refl} with $r=-s-N^{-\e}$ and $r=-s$ this is
\[
\P^v\big(N^{2/3}(\lambda_1+2)\ge s+N^{-\e}\big)-N^{-\delta}\le\P^w\big(N^{2/3}(\lambda_1+2)\ge s\big),
\]
the first inequality of \eqref{eq:tw1}. The second inequality of \eqref{eq:tw} for $(-M^v,-M^w)$ at $-s$ reads
$\P_{-M^w}(N^{2/3}(\lambda_N-2)\le-s)\le\P_{-M^v}(N^{2/3}(\lambda_N-2)\le-s+N^{-\e})+N^{-\delta}$; by \eqref{eq:B2refl} with $r=-s$ and $r=-s+N^{-\e}$ this is
\[
\P^w\big(N^{2/3}(\lambda_1+2)\ge s\big)\le\P^v\big(N^{2/3}(\lambda_1+2)\ge s-N^{-\e}\big)+N^{-\delta},
\]
the second inequality of \eqref{eq:tw1}.

\emph{Quantifiers.} $\e=\frac1{100}$ and $\delta=\frac1{20}$ are absolute; $N_0^{\rm edge}$ \eqref{eq:B2N0} depends only on $\Par$. The parameters $\delta_\pm$ enter only through the
thresholds $N_0$ (of Part~I) and $N_1$; the exponents and the constant $C_4(\vartheta,C_0)$ do not depend on them.
\end{proof}

\begin{remark}[other values of $\e,\delta$]\label{rem:B2eps}
The same proof works for every $\e\in(0,\frac1{48})$ and every $\delta\in(0,\frac16-8\e)$, with $\eg:=2\e$: then $\eg<\frac1{24}\le\frac1{10}$, so Lemma~\ref{lem:N0},
Proposition~\ref{prop:sandwich} and Theorem~\ref{thm:gfc} apply, and (N6) holds for large $N$ because $-\frac16+4\eg=-\frac16+8\e<-\delta$. The threshold becomes
$\max\{N_1,(2C_4C_q)^{1/(\frac16-8\e-\delta)},4^{1/(20-\delta)}\}$.
\end{remark}


% ================= PART III =================
\clearpage
\part{Strong local ergodicity of Dyson Brownian motion (Theorem 2.3)}
\section{Statement of Theorem 2.3 and overview of the route}\label{sec:C1}

Part~III proves Theorem~\ref{thm:23} of NP. The only results of Parts~I and~II used are those listed in \S\ref{sec:reqA}; Theorem~\ref{thm:21} is not used.
Statements numbered with a letter (such as Lemma~\ref{c2:transfer}) are auxiliary statements.
References of the form ``Part~I, Lemma~$x$'' are to the numbered statements of Part~I.
\emph{Local notation.} The constants $c_0$, $c_1$, $C_1'$, $C_1''$, $C_\sharp$, $C_\star$, the thresholds $N_*$, $N_K$, $N_2$, $N_3$, $N_4$, $N_6$, $N_c$, $N_{\rm q}$, and the functions $\Psi$, $\Lambda$ of Part~III are defined
where they first occur and are used only in Part~III; they are unrelated to the quantities with the same names in Parts~I and~II (for instance the exponent $c_0$ and the
constant $C_1$ of Appendix~\ref{app:ledger}, or $\Psi$ in Section~\ref{sec:edge}). The letter $\delta$ without index denotes only the exponent of Theorem~\ref{thm:23}.

\subsection{Setting and notation}\label{ss:notation}

Throughout, $\K\in\{\R,\C\}$ and $\beta:=\dim_\R\K\in\{1,2\}$; $N\ge 1$.
\begin{itemize}[leftmargin=1.6em,itemsep=2pt]
\item $\cH=\cH_N(\K)$ is the real vector space of self-adjoint $N\times N$ matrices over $\K$ with the inner product
$\la X,Y\ra:=\Tr(XY)$, norm $|X|:=(\Tr X^2)^{1/2}=(\sum_{i,j}|X_{ij}|^2)^{1/2}$, dimension $d=N+\beta N(N-1)/2$.
Orthonormal basis: $E_{ii}$; $X_{ij}:=(E_{ij}+E_{ji})/\sqrt2$ $(i<j)$; and, if $\beta=2$, $Y_{ij}:=\mathrm{i}(E_{ij}-E_{ji})/\sqrt2$ $(i<j)$ (Lemma \ref{c3:GE} and the paragraph before it).
Lebesgue measure $\rd H$ on $\cH$ is the one normalised by this basis; gradients $\nabla$ and the Laplacian $\Delta_\cH$ refer to $\la\cdot,\cdot\ra$. $\|A\|$ is the operator norm.
\item $U_\beta(N)$ is the orthogonal ($\beta=1$) or unitary ($\beta=2$) group. $\Phi:\cH\to\R$ is \emph{invariant} if $\Phi(UHU^*)=\Phi(H)$ for all $U\in U_\beta(N)$, $H\in\cH$.
\item $\Sig:=\{x\in\R^N:x_1<\dots<x_N\}$, $\cSig$ its closure; $\lambda:\cH\to\cSig$ is the vector of ordered eigenvalues
$\lambda_1(H)\le\dots\le\lambda_N(H)$, counted with multiplicity. By (RA4) of \S\ref{sec:reqA}, $\lambda$ is $1$-Lipschitz from $(\cH,|\cdot|)$ to $(\R^N,|\cdot|_\infty)$, hence continuous and Borel. For $x\in\R^N$, $\diag x\in\cH$.
\item \emph{Ensembles.} An ensemble of size $N$ over $\K$ is a family $\nu=(\nu_{ij})_{1\le i\le j\le N}$ of Borel probability measures on $\K$ (Part~I, Definition~\ref{def:ens}); its matrix is $H(x)$, $x\sim\PP_\nu=\bigotimes_{i\le j}\nu_{ij}$, and its variances are $\sigma_{ij}^2=\int|y|^2\nu_{ij}(\rd y)$.
$\mathcal E_N(\vartheta,\delta_-,\delta_+,C_0)$ is the class of ensembles satisfying (R), (D), (Z), (A), (B), (C) (Part~I, Definition~\ref{def:ass}); we write $\delta_\pm$ for the pair $(\delta_-,\delta_+)$.
\item \emph{Gaussian ensemble.} $\sigma^2:=2/(\beta N)$ and $\bG=\bG_{N,\beta}$ is the centred Gaussian measure on $\cH$ with covariance $\sigma^2\,\mathrm{Id}$ (Definition \ref{c3:def:gauss} with $\mathcal E=\cH$). By Lemma \ref{c3:GE}, $\bG(\rd H)=(2\pi\sigma^2)^{-d/2}\exp\big(-\tfrac{\beta N}{4}\Tr H^2\big)\rd H$;
a random $V\in\cH$ has law $\bG$ iff its entries $v_{ij}$, $i\le j$, are independent, $v_{ii}\sim N(0,2/(\beta N))$ real,
and for $i<j$: $v_{ij}\sim N(0,1/N)$ real if $\beta=1$, while $\sqrt2\,\mathrm{Re}\,v_{ij},\sqrt2\,\mathrm{Im}\,v_{ij}$ are i.i.d.\ $N(0,1/N)$ if $\beta=2$
(so $\E|v_{ij}|^2=1/N$, and $\E v_{ij}^2=0$ if $\beta=2$). This is the GOE ($\beta=1$, diagonal variance $2/N$; PD-06) resp.\ the GUE ($\beta=2$), exactly the matrix $V$ of the setting of Theorem~\ref{thm:23} in Part~I.
$V$ always denotes a $\bG$-distributed matrix. We put $\tau_{ij}^2:=\E|v_{ij}|^2$, i.e.\ $\tau_{ij}^2=1/N$ ($i\ne j$), $\tau_{ii}^2=2/(\beta N)$.
\item \emph{The flow.} Let $\nu$ be an ensemble and let $H_0$ and $V$ be random elements of $\cH$, defined on one probability space $(\widehat\Omega,\widehat{\mathcal F},\P)$, independent, with $H_0\overset{d}=H(x)$ under $\PP_\nu$ and $V\sim\bG$ (e.g.\ $\widehat\Omega=\Omega\times\cH$ with $\PP_\nu\otimes\bG$). $\E$ is the expectation on $\widehat\Omega$. For $t\in[0,\infty]$ put
\begin{equation}\label{Ht}
a_t:=e^{-t/2},\qquad s_t:=1-e^{-t},\qquad H_t:=a_tH_0+s_t^{1/2}V\qquad(e^{-\infty}:=0,\ \text{so }H_\infty=V).
\end{equation}
NP takes \eqref{Ht} (the original's (2.29)) as the \emph{definition}; the matrix SDE (2.28) is not used.
\item $\mu=\mu_{N,\beta}:=\lambda_*\bG$ (law of the ordered eigenvalues of $V$) and $\nu_t:=$ law of $\lambda(H_t)$; both are Borel probability measures on $\cSig$, and $\nu_\infty=\mu$.
\item \emph{Semicircle} (Part~I, \eqref{eq:sc}): $\rhosc(x)=\frac1{2\pi}\sqrt{(4-x^2)_+}$, $\nsc(E)=\int_{-\infty}^E\rhosc$; the classical locations $\gamma_j$ satisfy $N\nsc(\gamma_j)=j$, $-2<\gamma_1<\dots<\gamma_N=2$; $\gamma:=(\gamma_1,\dots,\gamma_N)\in\Sig$.
For $x\in\R^N$ the empirical counting function is $\fn_x(E):=N^{-1}\#\{j:x_j\le E\}$; thus $\fn_{\lambda(H)}$ is the function $\fn$ of Part~I for the matrix $H$.
For $0<c\le2$ we put $\varrho_*(c):=\frac1{2\pi}\sqrt{c-c^2/16}$.
\item \emph{Test functions.} For $O:\R^n\to\R$: $\|O\|_\infty:=\sup|O|$ and $\Lip(O):=\sup_{u\ne v}|O(u)-O(v)|/|u-v|$, Euclidean norm on $\R^n$.
\item \emph{Density-free correlation functionals} (PD-05). For a bounded Borel $O:\R^n\to\R$, $1\le n\le N$, $E'\in\R$, $\varrho>0$ and $x\in\R^N$,
\begin{equation}\label{defX}
\cX^{O}_{E',\varrho}(x):=\sum_{\substack{i_1,\dots,i_n\in\{1,\dots,N\}\\ \text{pairwise distinct}}}O\big(N\varrho(x_{i_1}-E'),\dots,N\varrho(x_{i_n}-E')\big),\qquad C_{N,n}:=\frac{N^n(N-n)!}{N!}.
\end{equation}
If $p^{(n)}[P]$ denotes the $n$-point marginal of the symmetrisation of a probability measure $P$ on $\R^N$, then $N^n\int F\,\rd p^{(n)}[P]=C_{N,n}\int\sum_{i_1,\dots,i_n\text{ distinct}}F(x_{i_1},\dots,x_{i_n})\,P(\rd x)$ for bounded Borel $F$
(Part~I, comment (a) after Theorem~\ref{thm:23} there); so $C_{N,n}\int\cX^O_{E',\varrho}\rd P$ is the $n$-point correlation function of $P$ tested against the rescaled observable, and when densities exist the left side of \eqref{eq:abstrthm2} below is the left side of (2.36) of the original (with $\varrho=\rhosc(E)$). The formulation needs no densities.
\end{itemize}

\subsection{The theorem}

The theorem is stated exactly as in Part~I (Theorem~\ref{thm:23} there); it is repeated here so that Part~III can be read on its own.

\renewcommand{\restatedname}{Theorem~\ref*{thm:23}}
\begin{restated}[strong local ergodicity of Dyson Brownian motion; restated from Section~\ref*{sec:main}]
Let $\vartheta,\delta_-,\delta_+,C_0>0$, $\eps'>0$, $\delta>0$, $0<c\le2$, $n\ge1$ and $\rho_O>0$.
There is a constant $C$ depending only on $(\eps',\delta,c,n,\rho_O,\vartheta,\delta_-,\delta_+,C_0)$ such that for all $N\ge\max(n,3)$,
$\K\in\{\R,\C\}$, $\nu\in\mathcal E_N(\vartheta,\delta_-,\delta_+,C_0)$, every Lipschitz function $O:\R^n\to\R$ with support in $[-\rho_O,\rho_O]^n$, and all
$E\in[-2+c,2-c]$, $b\in(0,c/2)$ and $t\ge N^{-1+\delta+\eps'}$, with $\varrho:=\rhosc(E)$:
\begin{multline}\tag{\ref*{eq:abstrthm2}}
 \Big|\int_{E-b}^{E+b}\frac{\rd E'}{2b}\;C_{N,n}\Big(\E\,\cX^{O}_{E',\varrho}\big(\lambda(H_t)\big)-\E\,\cX^{O}_{E',\varrho}\big(\lambda(V)\big)\Big)\Big|\\
 \le C\,\big(\|O\|_\infty+\Lip(O)\big)\,N^{2\eps'}\Big[b^{-1}N^{-1+\eps'}+b^{-1/2}N^{-\delta/2}\Big].
\end{multline}
\end{restated}

\noindent Here $H_t$ is \eqref{Ht} with $H_0$ the matrix of $\nu$. The left side depends only on the law of $H_0$, so the choice of the probability space $\widehat\Omega$ is immaterial.
Since $O$ is Lipschitz with compact support it is bounded and continuous; $|\cX^O_{E',\varrho}|\le N^n\|O\|_\infty$, and $(E',\hat\omega)\mapsto\cX^O_{E',\varrho}(\lambda(H_t(\hat\omega)))$ is continuous in $E'$ and Borel in $\hat\omega$, hence jointly measurable;
so all expectations and the $E'$-integral in \eqref{eq:abstrthm2} exist. Because the left side is linear in $O$, the theorem is equivalent to its special case $\|O\|_\infty+\Lip(O)\le1$.
Theorem \ref{thm:23} is a corollary of the following sharper statement, which is what Part~III proves.

\begin{theorem}[Quantitative form]\label{thm:sharp}
Let $\vartheta,\delta_-,\delta_+,C_0>0$, $0<c\le2$, $n\ge1$ and $\rho_O>0$. There are $A_1\ge1$, $N_{\rm q}\ge3$ and $C_1'>0$, depending only on
$(c,n,\rho_O,\vartheta,\delta_-,\delta_+,C_0)$, such that for all $N\ge\max(N_{\rm q},n)$, $\K\in\{\R,\C\}$, $\nu\in\mathcal E_N(\vartheta,\delta_-,\delta_+,C_0)$, every Lipschitz $O:\R^n\to\R$ with support in $[-\rho_O,\rho_O]^n$, all $t>0$ with $Nt\ge (\log N)^{2}$,
all $E\in[-2+c,2-c]$ and all $b\in(0,c/2)$:
\begin{equation}\label{sharp}
\text{l.h.s.\ of \eqref{eq:abstrthm2}}\;\le\;C_1'\big(\|O\|_\infty+\Lip(O)\big)\,e^{A_1(\log\log N)^2}\Big[\,b^{-1}N^{-1}+b^{-1/2}\,(Nt)^{-1/2}\Big].
\end{equation}
\end{theorem}

\begin{auxremark}[What Section~\ref{sec:C8} proves]\label{c1:oblig}
Theorems \ref{thm:sharp} and \ref{thm:23} are proved in \S\ref{sec:C8}, where the following is established:
\begin{enumerate}[label=(O\arabic*),leftmargin=2.4em,itemsep=1pt]
\item Theorem \ref{thm:sharp} with $A_1$, $N_{\rm q}$, $C_1'$ depending only on $(c,n,\rho_O,\vartheta,\delta_\pm,C_0)$, the dependence on $O$ being only through $\rho_O$ and the factor $\|O\|_\infty+\Lip(O)$ (Remark \ref{rem:Odep}); $\Lip$ is with respect to the Euclidean norm (coordinatewise Lipschitz constants are bounded by it).
\item The deduction of Theorem \ref{thm:23}. The elementary argument is: let $N_2\ge\max(N_{\rm q},n,3)$ depend on $(\eps',\delta,A_1,N_{\rm q},n)$ and be such that $e^{A_1(\log\log N)^2}\le N^{\eps'}$ and $N^{\delta+\eps'}\ge(\log N)^2$ for $N\ge N_2$.
For $N\ge N_2$ and $t\ge N^{-1+\delta+\eps'}$ one has $Nt\ge(\log N)^2$, $(Nt)^{-1/2}\le N^{-\delta/2}$ and $N^{-1}\le N^{-1+\eps'}$, so \eqref{sharp} gives \eqref{eq:abstrthm2} with $C=C_1'$.
For $\max(n,3)\le N<N_2$ the left side of \eqref{eq:abstrthm2} is at most $2C_{N,n}\frac{N!}{(N-n)!}\|O\|_\infty=2N^n\|O\|_\infty\le2N_2^n\|O\|_\infty$, and the right side is at least $C\|O\|_\infty b^{-1/2}N^{-\delta/2}\ge C\|O\|_\infty(c/2)^{-1/2}N_2^{-\delta/2}$;
so $C:=\max\big(C_1',\,2(c/2)^{1/2}N_2^{\,n+\delta/2}\big)$ works. Every quantity here depends only on $(\eps',\delta,c,n,\rho_O,\vartheta,\delta_\pm,C_0)$.
\item The H\"older corollary (Corollary \ref{cor:cont}) and the remark that merely continuous $O$ is not covered (Remark \ref{rem:contO}).
\end{enumerate}
The inputs prepared in the present section and in \S\ref{sec:C2} are: the notation \eqref{Ht}, \eqref{defX}; the constants $A_0,L,D$ of \eqref{defL}; the good set $\Xi$ \eqref{good} with Lemmas \ref{lem:goodC}, \ref{lem:Q}, \ref{lem:countC}.
\end{auxremark}

\begin{remark}[Differences from the original]\label{rem:diff}
The differences between Theorem~\ref{thm:23} and (2.36) of the original, and between the route of Part~III and that of the original, are listed in Appendix~\ref{app:diff}
(items PD-04, PD-05, PD-06, PD-16, PD-21 of \S\ref{app:diffPD}, and \S\ref{sec:diff}).
\end{remark}

\subsection{Overview of the proof and why this route}\label{ss:overview}

The original derives (2.36) from [ESYY, Thm 2.1], whose proof needs: the eigenvalue density $f_t$ as a solution of $\partial_tf=Lf$ on $\Sig$ (stochastic calculus, or the Weyl integration formula and second-order perturbation theory; a Markov semigroup on $\Sig$ with singular drift and uniqueness at the walls $x_i=x_{i+1}$);
the logarithmic Sobolev inequality for the local relaxation measure $\omega$ (Bakry--\'Emery); and an auxiliary $\omega$-reversible dynamics with a $\Gamma_2$ computation (the ``Dirichlet form inequality'', [ESYY, Thm 4.1--4.3]).
None of this is available in Mathlib. NP keeps the \emph{idea} (relative entropy with respect to a local relaxation measure) but replaces every analytic tool:

\begin{enumerate}[leftmargin=1.8em,itemsep=2pt,label=(\Roman*)]
\item \emph{Flow at matrix level, restricted to invariant densities} (\S\ref{sec:C3}--\S\ref{sec:C4}). The matrix Ornstein--Uhlenbeck semigroup $T_u$ is explicit (Mehler formula \eqref{Mehler}). If $f_s=\rd\nu_s/\rd\mu$, then $\Phi_t:=T_{t-s}(f_s\circ\lambda)$ is $C^\infty$, positive and invariant, and $f_t(x):=\Phi_t(\diag x)$ is a version of $\rd\nu_t/\rd\mu$ which is a \emph{smooth symmetric function on all of $\R^N$} (Proposition \ref{prop:U9}).
A $2\times2$ Givens-rotation identity (Lemma \ref{lem:givens}) gives $(A\Phi)(\diag x)=(L\varphi)(x)$ for every invariant $C^2$ function, so $\partial_tf_t=Lf_t$ holds classically off the walls. No eigenvalue perturbation theory, no Haar measure, no Weyl integration formula, no boundary analysis.
\item \emph{Joint eigenvalue density of G$\beta$E by ergodicity} (\S\ref{sec:C5}). The explicit measure $\mu^*\propto\Delta^\beta e^{-\beta N|x|^2/4}$ on $\Sig$ is invariant under the radial flow (integration by parts on $\Sig$ along coordinate lines; the weight vanishes on the walls because $\beta\ge1$), and $T_u\Phi\to\int\Phi\,\rd\bG$ as $u\to\infty$. Hence $\mu=\mu^*$ (Theorem \ref{thm:JD}).
\item \emph{Pr\'ekopa--Leindler does everything else} (\S\ref{sec:C6}): (a) the LSI for $\omega$ on the convex set $\Sig$ (Theorem \ref{thm:LSI}; Bobkov--Ledoux argument, no semigroup); (b) a sub-Gaussian bound for gap observables under \emph{any} convex log-gas (Theorem \ref{thm:conc}), which with the entropy inequality gives
$|\int F\,\rd\nu-\int F\,\rd\omega|\le\frac2N\sqrt{C_\star(n-1)|J|\,S(\nu|\omega)}$ (Corollary \ref{cor:gap}). This replaces [ESYY, Thm 4.1] and the auxiliary dynamics entirely, and needs only the relative entropy, not the Dirichlet form.
\item \emph{Entropy bound} (\S\ref{sec:C7}): $\frac{\rd}{\rd t}S(\nu_t|\omega)\le-\frac{1}{\beta R^2}S(\nu_t|\omega)+\frac{NQ}{2\beta R^4}$ (Proposition \ref{prop:entprod} and (III)(a)); with the a priori bound $S(\nu_s|\mu)\le N^3$ (Lemma \ref{lem:U8}: Gaussian relative entropy and data processing) this gives $S(\nu_t|\omega)\le\frac1N+\frac{32\beta D^2\log N}{t}$ (Theorem \ref{thm:entropy}).
\item \emph{From gaps to correlation functions} (\S\ref{sec:C8}): deterministic bookkeeping on the good set $\Xi$ of \eqref{good}, on which at most $2\ell/\pi+2D$ eigenvalues lie in a window of size $2\ell/N$ (Lemma \ref{lem:countC}).
\end{enumerate}
Here $R>0$ is the confinement scale, $W(x)=|x-\gamma|^2/(2R^2)$, $\omega\propto e^{-NW}\mu$ (Definition \ref{def:omega}), and $Q:=7D^2N^{-1}$ bounds $\E|\lambda(H_r)-\gamma|^2$ uniformly in $r\in[0,\infty]$ by rigidity (Lemma \ref{lem:Q}). The choice $R^2=t/(8\beta\log N)$ yields \eqref{sharp}.
The only input from the random-matrix part of NP is Theorem~\ref{thm:22} (rigidity), applied at each fixed time to the rescaled matrix $\widetilde H_t$, which belongs to one class $\mathcal E_N$ for all $t$ (Lemma \ref{lem:U11}).
\section{Inputs from Part~I; $H_t$ and G$\beta$E as ensembles}\label{sec:C2}

\subsection{Required from Part~I}\label{sec:reqA}

Part~III uses exactly the following statements of Part~I, as they are stated there, and Lemma \ref{c3:GE} of \S\ref{sec:C3} (the Gaussian ensemble on $\cH$). Nothing else from Parts~I and~II is used.

\begin{enumerate}[leftmargin=2.4em,itemsep=2pt,label=(RA\arabic*)]
\item \textbf{Theorem \ref{thm:22} of Part~I (rigidity)}: \emph{for all $\vartheta,\delta_-,\delta_+,C_0>0$ there are $A_0>1$, $C,c>0$, $\phi\in(0,1)$ and $N_0$, depending only on $(\vartheta,\delta_-,\delta_+,C_0)$, such that for every $N\ge N_0$, $\K\in\{\R,\C\}$, $\nu\in\mathcal E_N(\vartheta,\delta_-,\delta_+,C_0)$ and every real $L$ with $A_0\log\log N\le L\le\log(10N)/(10\log\log N)$ \textup{(}Part~I, \eqref{eq:Lbound}\textup{)}:}
\begin{gather*}
\PP_\nu\Big(\exists j:\ |\lambda_j-\gamma_j|\ge(\log N)^L[\min(j,N-j+1)]^{-1/3}N^{-2/3}\Big)\le Ce^{-c(\log N)^{\phi L}},\\
\PP_\nu\Big(\sup_{|E|\le5}|\fn(E)-\nsc(E)|\ge\frac{(\log N)^L}{N}\Big)\le Ce^{-c(\log N)^{\phi L}},
\end{gather*}
\emph{where $\lambda_j=\lambda_j(H(x))$ and $\fn=\fn_{\lambda(H(x))}$.} The class $\mathcal E_N$ (Part~I, Definition \ref{def:ass}) imposes only (R), (D), (Z), (A), (B), (C) on \emph{arbitrary} Borel probability entry laws; in particular it contains Gaussian convolutions and rescaled Gaussian ensembles once these conditions are verified (Lemma \ref{lem:U11}).
\item \textbf{Monotonicity} (Part~I, Lemma \ref{lem:mono}(ii)): $\mathcal E_N(\vartheta,\delta_-,\delta_+,C_0)\subset\mathcal E_N(\vartheta',\delta'_-,\delta'_+,C_0')$ if $0<\vartheta'\le\vartheta$, $0<\delta'_\pm\le\delta_\pm$, $C_0'\ge C_0$.
\item \textbf{Scaling} (Part~I, Lemma \ref{lem:mono}(v)): if a Borel probability measure $\mu_0$ on $\K$ with $s^2=\int|y|^2\rd\mu_0<\infty$ satisfies $\mu_0(\{|y|>xs\})\le\vartheta^{-1}e^{-x^\vartheta}$ for all $x\ge1$, then for every real $a$ its image under $y\mapsto ay$ has second moment $a^2s^2$ and satisfies the same bound, with the same $\vartheta$, relative to $|a|s$; centring and $\mu_0(\R)=1$ are preserved.
\item \textbf{Weyl} (Part~I, Lemma \ref{lem:weyl}(a),(b)): for Hermitian $A,A'$, $\|A\|=\max_k|\lambda_k(A)|$ and $|\lambda_k(A)-\lambda_k(A')|\le\|A-A'\|$ for all $k$.
Consequently $\|A\|^2=\max_k\lambda_k(A)^2\le\sum_k\lambda_k(A)^2=\Tr A^2=|A|^2$ (spectral theorem), so $\lambda:(\cH,|\cdot|)\to(\R^N,|\cdot|_\infty)$ is $1$-Lipschitz; $\lambda$ is continuous and Borel, and so is every function of the form $\Theta\circ\lambda$ with $\Theta$ Borel.
\item \textbf{Semicircle facts} (Part~I, Lemma \ref{lem:nsc}(a),(c),(d),(e)): $0\le\rhosc\le1/\pi$; $\nsc$ is nondecreasing, $\frac1\pi$-Lipschitz, strictly increasing on $[-2,2]$, $\nsc=0$ on $(-\infty,-2]$, $\nsc=1$ on $[2,\infty)$;
$\gamma_j\in(-2,2]$, $N\nsc(\gamma_j)=j$, $\gamma_N=2$; $\rhosc(x)\ge\varrho_*(c)$ for $|x|\le2-c/4$ $(0<c\le2)$; and $\#\{j:\gamma_j\in I\}\le N|I|/\pi+1$ for every bounded interval $I$.
\item \textbf{Moments} (Part~I, Lemma \ref{lem:moments}(a)): if $0<\vartheta\le1$ and $\mu_0$ is a Borel probability measure on $\C$ with $\mu_0(\{|y|>xs\})\le\vartheta^{-1}e^{-x^\vartheta}$ for all $x\ge1$, then $\int|y|^q\rd\mu_0\le(C_\vartheta q)^{q/\vartheta}s^q$ for every real $q\ge1$, with $C_\vartheta=90/\vartheta$.
\end{enumerate}
Theorem~\ref{thm:21} is not needed: the local density bound (``Assumption IV'' of [ESYY]) follows from the counting-function estimate in Theorem~\ref{thm:22}, see Lemma \ref{lem:countC}.
Note also that every ensemble has a vanishing entry $h_{ij}=0$ a.s.\ whenever $\sigma_{ij}=0$ (Part~I, remark after Definition~\ref{def:ass}).

\subsection{The ensembles $\widetilde H_t$ and $\widetilde V$}

For $t\in[0,\infty]$ let $a_t,s_t$ be as in \eqref{Ht} and
\begin{equation}\label{tildeH}
\kappa_t:=\Big(1+\frac{(2-\beta)\,s_t}{\beta N}\Big)^{-1/2}=\begin{cases}(1+s_t/N)^{-1/2},&\beta=1,\\ 1,&\beta=2,\end{cases}\qquad \widetilde H_t:=\kappa_tH_t,\qquad \widetilde V:=\widetilde H_\infty=\kappa_\infty V .
\end{equation}
For $\beta=2$ nothing is rescaled. For $\beta=1$ the rescaling compensates the diagonal variance $2/N$ of the GOE (PD-06), see Lemma \ref{lem:U11}(a).

\begin{auxlemma}[Matrices with independent entries are ensembles]\label{c2:transfer}
Let $M$ be a random element of $\cH$ on some probability space $(\Omega_1,\mathcal F_1,\P_1)$ whose entries $(m_{ij})_{1\le i\le j\le N}$ are independent, and let $\nu_{ij}$ be the law of $m_{ij}$. Then $\nu=(\nu_{ij})_{i\le j}$ is an ensemble with $\nu_{ii}(\R)=1$, and for every Borel set $S\subset\cH$: $\P_1(M\in S)=\PP_\nu(H(x)\in S)$.
In particular, for every Borel $S'\subset\R^N$, $\P_1(\lambda(M)\in S')=\PP_\nu(\lambda(H(x))\in S')$, and $\E_1\Theta(M)=\int\Theta(H(x))\PP_\nu(\rd x)$ for every Borel $\Theta\ge0$ on $\cH$.
\end{auxlemma}

\begin{proof}
Let $X:=(m_{ij})_{i\le j}$, a random element of $\Omega=\K^{\mathcal P_N}$. Since finitely many random variables are independent iff their joint law is the product of their laws, $X_*\P_1=\bigotimes_{i\le j}\nu_{ij}=\PP_\nu$. The diagonal entries of $M\in\cH$ are real, so $\nu_{ii}(\R)=1$.
At every point of $\Omega_1$, $M=H(X)$: for $i<j$, $H(X)_{ij}=m_{ij}$ and $H(X)_{ji}=\overline{m_{ij}}=m_{ji}$ because $M$ is self-adjoint; $H(X)_{ii}=\mathrm{Re}\,m_{ii}=m_{ii}$.
Since $x\mapsto H(x)$ is continuous (it is $\R$-linear), $H^{-1}(S)$ is Borel and $\P_1(M\in S)=\P_1(X\in H^{-1}(S))=\PP_\nu(H^{-1}(S))$. The other two statements follow with $S:=\lambda^{-1}(S')$ (Borel by (RA4)) and by the transformation formula for image measures.
\end{proof}

\begin{auxlemma}[Gaussian tails]\label{c2:gausstail}
(a) If $\xi$ is a real centred Gaussian variable with variance $\tau^2>0$, then $\P(|\xi|>y\tau)\le2e^{-y^2/2}$ for all $y\ge0$.
(b) If $V\sim\bG$, then $\P(|v_{ij}|>y\tau_{ij})\le4e^{-y^2/2}$ for all $i\le j$ and $y\ge0$.
\end{auxlemma}

\begin{proof}
(a) For $\lambda\in\R$, completing the square gives $\E e^{\lambda\xi}=\int\frac{1}{\sqrt{2\pi}\tau}e^{\lambda u-u^2/(2\tau^2)}\rd u=e^{\lambda^2\tau^2/2}\int\frac{1}{\sqrt{2\pi}\tau}e^{-(u-\lambda\tau^2)^2/(2\tau^2)}\rd u=e^{\lambda^2\tau^2/2}$.
By Markov's inequality applied to $e^{\lambda\xi}$ with $\lambda:=y/\tau$: $\P(\xi>y\tau)\le e^{-\lambda y\tau}\E e^{\lambda\xi}=e^{-y^2+y^2/2}=e^{-y^2/2}$. The same bound holds for $-\xi$, which has the same law; add the two.

(b) By Lemma \ref{c3:GE}(ii), $v_{ii}$ is real centred Gaussian with variance $\tau_{ii}^2$, and so is $v_{ij}$ ($i<j$) if $\beta=1$; (a) gives the bound $2e^{-y^2/2}$. If $\beta=2$ and $i<j$, then $v_{ij}=(\xi+\mathrm i\eta)/\sqrt2$ with $\xi,\eta$ real centred Gaussian of variance $1/N=\tau_{ij}^2$.
If $|v_{ij}|>y\tau_{ij}$ then $\xi^2+\eta^2>2y^2\tau_{ij}^2$, so $|\xi|>y\tau_{ij}$ or $|\eta|>y\tau_{ij}$; by (a) the probability is at most $4e^{-y^2/2}$.
\end{proof}

\begin{lemma}[$\widetilde H_t$ and $\widetilde V$ are ensembles with uniform parameters]\label{lem:U11}
Let $N\ge3$, $\vartheta,\delta_-,\delta_+,C_0>0$, $\nu\in\mathcal E_N(\vartheta,\delta_-,\delta_+,C_0)$ over $\K$, and let $H_0,V,H_t$ be as in \eqref{Ht}. Put
\begin{equation*}
\tilde\vartheta:=\tfrac1{14}\min(\vartheta,1),\qquad\tilde\delta_\pm:=\tfrac34\min(\delta_\pm,1),\qquad\tilde C_0:=\max(C_0,2).
\end{equation*}
Then for every $t\in[0,\infty]$:
\begin{enumerate}[label=(\alph*),itemsep=1pt]
\item the entries $h^t_{ij}:=a_th_{ij}+s_t^{1/2}v_{ij}$ $(i\le j)$ of $H_t$ are independent and centred, with
$\sigma_{t,ij}^2:=\E|h^t_{ij}|^2=e^{-t}\sigma_{ij}^2+s_t\tau_{ij}^2$, and $\sum_{i=1}^N\sigma^2_{t,ij}=\kappa_t^{-2}$ for every $j$;
\item the laws $\tilde\nu^t_{ij}$ of the entries $\kappa_th^t_{ij}$ of $\widetilde H_t$ form an ensemble $\tilde\nu^t\in\mathcal E_N(\tilde\vartheta,\tilde\delta_-,\tilde\delta_+,\tilde C_0)$ over $\K$, and $\widetilde H_t$ has the law of $H(x)$ under $\PP_{\tilde\nu^t}$.
In particular (case $t=\infty$) $\widetilde V$ is the matrix of an ensemble in $\mathcal E_N(\tilde\vartheta,\tilde\delta_-,\tilde\delta_+,\tilde C_0)$; this ensemble does not depend on $\nu$.
\item \emph{(deterministic)} $1\ge\kappa_t\ge(1+1/N)^{-1/2}$, $0\le\kappa_t^{-1}-1\le\frac1{2N}$, and at every sample point
\[
\lambda(H_t)=\kappa_t^{-1}\lambda(\widetilde H_t),\qquad |\lambda_j(H_t)-\lambda_j(\widetilde H_t)|\le\frac{|\lambda_j(\widetilde H_t)|}{2N},\qquad \fn_{\lambda(H_t)}(e)=\fn_{\lambda(\widetilde H_t)}(\kappa_te)\quad(e\in\R).
\]
\end{enumerate}
\end{lemma}

\begin{proof}
Write $a=a_t$, $s=s_t$, $\kappa=\kappa_t$ and $r:=(2-\beta)s/(\beta N)\in[0,1/N]$, so that $\kappa^{-2}=1+r$ and $a^2+s=e^{-t}+s=1$.

\emph{(a)} $H_0=H(X)$ in law, with $X\sim\PP_\nu$, so the entries $(h_{ij})_{i\le j}$ of $H_0$ are independent (functions of distinct coordinates of $X$); the entries $(v_{ij})_{i\le j}$ of $V$ are independent (Lemma \ref{c3:GE}(ii)); and $H_0$, $V$ are independent.
Hence the joint law of the $2|\mathcal P_N|$ variables $(h_{ij},v_{ij})_{i\le j}$ is the product of their laws, and the variables $h^t_{ij}$, being functions of the disjoint blocks $\{h_{ij},v_{ij}\}$, are independent.
All of them are square integrable ($\sigma_{ij}^2\le C_0/N$ by (C)); $\E h_{ij}=0$ by (Z) and $\E v_{ij}=0$, so $\E h^t_{ij}=0$ and, by independence, $\E[\overline{h_{ij}}v_{ij}]=0$, whence $\E|h^t_{ij}|^2=a^2\E|h_{ij}|^2+s\E|v_{ij}|^2=e^{-t}\sigma_{ij}^2+s\tau_{ij}^2$.
By (A) and the values of $\tau_{ij}^2$: $\sum_i\sigma^2_{t,ij}=e^{-t}+s\big(\frac{N-1}N+\frac2{\beta N}\big)=e^{-t}+s+s\frac{2-\beta}{\beta N}=1+r=\kappa^{-2}$.

\emph{(b)} $\widetilde H_t=\kappa H_t\in\cH$ has the independent entries $\kappa h^t_{ij}$, $i\le j$, by (a); Lemma \ref{c2:transfer} shows that $\tilde\nu^t$ is an ensemble with real diagonal (R) and that $\widetilde H_t$ has the law of $H(x)$ under $\PP_{\tilde\nu^t}$. Its variances are $\tilde\sigma_{ij}^2:=\kappa^2\sigma^2_{t,ij}$. We verify the conditions of Part~I, Definition \ref{def:ass}.

\emph{(Z)} and \emph{(A)}: $\E\kappa h^t_{ij}=0$ by (a), and $\sum_i\tilde\sigma_{ij}^2=\kappa^2\kappa^{-2}=1$ by (a).

\emph{(C)}: $\kappa\le1$, $\sigma_{ij}^2\le C_0/N$ and $\tau_{ij}^2\le2/N$ give $\tilde\sigma_{ij}^2\le e^{-t}C_0/N+2s/N\le\max(C_0,2)/N$.

\emph{(B)}: let $B=(\sigma_{ij}^2)$, $\widetilde B=(\tilde\sigma_{ij}^2)$, and $v\in\R^N$ with $\sum_iv_i=0$. The variance matrix of $V$ is $(\tau_{ij}^2)=N^{-1}\mathbf 1\mathbf 1^{\mathsf T}+\frac{2-\beta}{\beta N}I$, and $\la v,\mathbf 1\mathbf 1^{\mathsf T}v\ra=(\sum_iv_i)^2=0$. Hence
\[
\la v,\widetilde Bv\ra=\kappa^2\big(e^{-t}\la v,Bv\ra+r|v|^2\big)=\frac{e^{-t}\la v,Bv\ra+r|v|^2}{1+r}.
\]
Put $\bar\delta_\pm:=\min(\delta_\pm,1)$. Using (B) for $\nu$ (if $t=\infty$ the term $e^{-t}\la v,Bv\ra$ vanishes and (B) is not needed):
\[
|v|^2-\la v,\widetilde Bv\ra\ge\frac{1+r-e^{-t}(1-\delta_+)-r}{1+r}|v|^2=\frac{s+e^{-t}\delta_+}{1+r}|v|^2\ge\frac{\bar\delta_+(s+e^{-t})}{1+r}|v|^2=\frac{\bar\delta_+}{1+r}|v|^2,
\]
\[
|v|^2+\la v,\widetilde Bv\ra\ge\frac{1+r+e^{-t}(-1+\delta_-)+r}{1+r}|v|^2=\frac{s+e^{-t}\delta_-+2r}{1+r}|v|^2\ge\frac{\bar\delta_-}{1+r}|v|^2 .
\]
Since $1+r\le1+1/N\le4/3$ ($N\ge3$), $\frac{\bar\delta_\pm}{1+r}\ge\tilde\delta_\pm$, i.e.\ $(-1+\tilde\delta_-)|v|^2\le\la v,\widetilde Bv\ra\le(1-\tilde\delta_+)|v|^2$.

\emph{(D)} with $\tilde\vartheta$: by (RA3) applied with the factor $\kappa$ it suffices to show, for fixed $(i,j)$, $\sigma_t:=\sigma_{t,ij}$ and all real $x\ge1$,
\begin{equation}\tag{T}\label{c2:tail}
\P\big(|h^t_{ij}|>x\sigma_t\big)\le\tilde\vartheta^{-1}e^{-x^{\tilde\vartheta}}.
\end{equation}
Put $\bar\vartheta:=\min(\vartheta,1)$, so $\tilde\vartheta=\bar\vartheta/14$ and, by (RA2), $\nu_{ij}(\{|y|>x\sigma_{ij}\})\le\bar\vartheta^{-1}e^{-x^{\bar\vartheta}}$ for $x\ge1$.
\emph{Case $1\le x<2$.} Then $x^{\tilde\vartheta}\le2^{1/14}<1.06$ and $\tilde\vartheta^{-1}e^{-x^{\tilde\vartheta}}\ge14e^{-1.06}>4$, which exceeds the left side of \eqref{c2:tail}.
\emph{Case $x\ge2$.} If $|h^t_{ij}|>x\sigma_t$ then $a|h_{ij}|>x\sigma_t/2$ or $s^{1/2}|v_{ij}|>x\sigma_t/2$.
For the first event: if $a\sigma_{ij}=0$ then $ah_{ij}=0$ a.s.\ (if $\sigma_{ij}=0$ then $h_{ij}=0$ a.s.), and the event has probability $0$ (strict inequality); otherwise $a\sigma_{ij}\le\sigma_t$ and the event is contained in $\{|h_{ij}|>\frac x2\sigma_{ij}\}$, of probability $\le\bar\vartheta^{-1}e^{-(x/2)^{\bar\vartheta}}$ since $x/2\ge1$.
For the second event: if $s=0$ it is empty; otherwise $s^{1/2}\tau_{ij}\le\sigma_t$ and the event is contained in $\{|v_{ij}|>\frac x2\tau_{ij}\}$, of probability $\le4e^{-x^2/8}$ by Lemma \ref{c2:gausstail}(b).
Now let $y:=x^{\bar\vartheta/2}\ge1$. Then $(x/2)^{\bar\vartheta}=2^{-\bar\vartheta}y^2\ge y^2/2\ge y-1$, because $y^2-2y+2=(y-1)^2+1>0$.
Further, with $u:=x^{1/2}$, $x^2/8-x^{1/2}+1=u^4/8-u+1\ge1-\frac34 2^{1/3}>0$ (the minimum over $u\ge0$ is attained at $u^3=2$, and $(\frac34)^3\cdot2=\frac{27}{32}<1$); since $\bar\vartheta/2\le\frac12$ and $x\ge1$, $x^2/8\ge x^{1/2}-1\ge y-1$.
Hence the left side of \eqref{c2:tail} is at most $(\bar\vartheta^{-1}+4)e\,e^{-y}\le5e\,\bar\vartheta^{-1}e^{-x^{\bar\vartheta/2}}\le14\bar\vartheta^{-1}e^{-x^{\tilde\vartheta}}=\tilde\vartheta^{-1}e^{-x^{\tilde\vartheta}}$, where we used $\bar\vartheta\le1$, $5e<14$, and $x^{\tilde\vartheta}\le x^{\bar\vartheta/2}$ (as $x\ge1$, $\tilde\vartheta\le\bar\vartheta/2$).
This proves \eqref{c2:tail}, so $\tilde\nu^t\in\mathcal E_N(\tilde\vartheta,\tilde\delta_-,\tilde\delta_+,\tilde C_0)$. For $t=\infty$, $h^\infty_{ij}=v_{ij}$ and nothing depends on $\nu$.

\emph{(c)} $\kappa^{-2}=1+r\in[1,1+1/N]$ gives the bounds on $\kappa$, and $\kappa^{-1}-1=\sqrt{1+r}-1\le r/2\le1/(2N)$.
For $c>0$ and Hermitian $A$, the eigenvalues of $cA$ counted with multiplicity are $c\lambda_1(A)\le\dots\le c\lambda_N(A)$, so $\lambda(cA)=c\lambda(A)$; with $A=\widetilde H_t$, $c=\kappa^{-1}$ this is the first identity, and the second follows: $|\lambda_j(H_t)-\lambda_j(\widetilde H_t)|=(\kappa^{-1}-1)|\lambda_j(\widetilde H_t)|$.
Finally $\fn_{\lambda(H_t)}(e)=N^{-1}\#\{j:\kappa^{-1}\lambda_j(\widetilde H_t)\le e\}=N^{-1}\#\{j:\lambda_j(\widetilde H_t)\le\kappa e\}$.
\end{proof}

\subsection{Standing constants, the good set, and the quantity $Q$}

\noindent\textbf{Standing choice of constants.} Fix $(\vartheta,\delta_-,\delta_+,C_0)$ and let $(\tilde\vartheta,\tilde\delta_-,\tilde\delta_+,\tilde C_0)$ be as in Lemma \ref{lem:U11}.
Let $(A_0,C_R,c_R,\phi,N_0)$ be the constants $(A_0,C,c,\phi,N_0)$ of Theorem \ref{thm:22} of Part~I (RA1) for the parameters $(\tilde\vartheta,\tilde\delta_-,\tilde\delta_+,\tilde C_0)$; they depend only on $(\vartheta,\delta_\pm,C_0)$. For $N\ge3$ set
\begin{equation}\label{defL}
L:=A_0\log\log N,\qquad D:=(\log N)^{L}+2=e^{A_0(\log\log N)^2}+2 .
\end{equation}
Since $(\log\log N)^2/\log N\to0$ and $\log\log N\to\infty$, there is $N_*=N_*(\vartheta,\delta_\pm,C_0)\ge\max(N_0,16)$ such that for all $N\ge N_*$
\begin{equation}\tag{N$_*$}\label{c2:Nstar}
10A_0(\log\log N)^2\le\log(10N)\qquad\text{and}\qquad\phi A_0\log\log N\ge2 .
\end{equation}
For $N\ge N_*$, $L$ satisfies Part~I, \eqref{eq:Lbound}, so Theorem \ref{thm:22} applies with this $L$; moreover $A_0(\log\log N)^2\le\frac1{10}\log(10N)\le\frac15\log N$, i.e.
\begin{equation}\tag{D$_*$}\label{c2:Dsmall}
(\log N)^L\le N^{1/5},\qquad D\le N^{1/5}+2 .
\end{equation}
The \emph{good set} is
\begin{equation}\label{good}
\Xi:=\Big\{x\in\cSig:\ \sup_{|e|\le4}|\fn_x(e)-\nsc(e)|\le\frac DN\ \text{ and }\ |x_j-\gamma_j|\le D\,[\min(j,N+1-j)]^{-1/3}N^{-2/3}\ \text{ for all }j\Big\}.
\end{equation}

\begin{auxlemma}[Measurability]\label{c2:Xi}
For every $a>0$ and every $\Delta\ge0$, the function $x\mapsto\sup_{|e|\le a}|\fn_x(e)-\nsc(e)|$ is Borel on $\R^N$, and $\{x\in\cSig:\sup_{|e|\le a}|\fn_x(e)-\nsc(e)|\le\Delta\}$ is Borel. In particular $\Xi$ is a Borel subset of $\cSig$.
\end{auxlemma}

\begin{proof}
For fixed $e$, $x\mapsto\fn_x(e)=N^{-1}\sum_j\mathbf 1\{x_j\le e\}$ is Borel. For fixed $x$, $e\mapsto\fn_x(e)$ is right-continuous: if $e_k\downarrow e$, then for $k$ large no $x_j$ lies in $(e,e_k]$, so $\fn_x(e_k)=\fn_x(e)$; and $\nsc$ is continuous (RA5).
Let $J_0:=(\mathbb Q\cap[-a,a])\cup\{a\}$. Every $e\in[-a,a)$ is the limit of a decreasing sequence in $J_0\cap(e,a]$, so $|\fn_x(e)-\nsc(e)|\le\sup_{e'\in J_0}|\fn_x(e')-\nsc(e')|$; hence the supremum over $[-a,a]$ equals the supremum over the countable set $J_0$, which is Borel.
The second condition in \eqref{good} defines a closed set. Intersections of Borel sets are Borel.
\end{proof}

\begin{lemma}[Good event, uniformly in $t$]\label{lem:goodC}
For every $K>0$ there is $N_K=N_K(K,\vartheta,\delta_\pm,C_0)\ge N_*$ such that for all $N\ge N_K$, $\K\in\{\R,\C\}$, $\nu\in\mathcal E_N(\vartheta,\delta_-,\delta_+,C_0)$ and all $t\in[0,\infty]$:
\[
\P\big(\lambda(H_t)\notin\Xi\big)\le N^{-K};\qquad\text{in particular (}t=\infty\text{)}\quad\mu(\Xi^c)=\P(\lambda(V)\notin\Xi)\le N^{-K}.
\]
\end{lemma}

\begin{proof}
Let $N_K\ge N_*$ be such that for $N\ge N_K$: $c_R\log N\ge K+\max(0,\log(2C_R))$ (note $\log N\ge1$ as $N\ge16$). Fix $N\ge N_K$, $\nu$, $t$, and write $m_j:=\min(j,N+1-j)$.

\emph{Step 1 (the event from Theorem~\ref{thm:22}).} Let
\[
S:=\Big\{y\in\cSig:\ |y_j-\gamma_j|<(\log N)^Lm_j^{-1/3}N^{-2/3}\ \forall j,\ \ \sup_{|e|\le5}|\fn_y(e)-\nsc(e)|<\frac{(\log N)^L}N\Big\},
\]
a Borel set by Lemma \ref{c2:Xi} (the first condition defines an open set). By Lemma \ref{lem:U11}(b), $\tilde\nu^t\in\mathcal E_N(\tilde\vartheta,\tilde\delta_-,\tilde\delta_+,\tilde C_0)$ and $\widetilde H_t$ has the law of $H(x)$ under $\PP_{\tilde\nu^t}$.
The complement $\cSig\setminus S$ is the union of the two sets whose $\PP_{\tilde\nu^t}$-probabilities are bounded in (RA1). Since $N\ge N_*\ge N_0$ and $L$ is admissible,
\[
\P\big(\lambda(\widetilde H_t)\notin S\big)=\PP_{\tilde\nu^t}\big(\lambda(H(x))\notin S\big)\le2C_R\exp\big(-c_R(\log N)^{\phi L}\big).
\]

\emph{Step 2 (deterministic transfer).} Let $y\in S$ and $x:=\kappa_t^{-1}y$. We show $x\in\Xi$. By Lemma \ref{lem:U11}(c) this gives $\{\lambda(\widetilde H_t)\in S\}\subset\{\lambda(H_t)\in\Xi\}$.
By \eqref{c2:Dsmall}, $(\log N)^LN^{-2/3}\le N^{1/5-2/3}\le1$, so $|y_j|\le|\gamma_j|+1\le3$. With $0\le\kappa_t^{-1}-1\le\frac1{2N}$ and $N^{-1}=N^{-1/3}N^{-2/3}\le m_j^{-1/3}N^{-2/3}$ (as $m_j\le N$):
\begin{align*}
|x_j-\gamma_j|&\le|y_j-\gamma_j|+(\kappa_t^{-1}-1)|y_j|\le(\log N)^Lm_j^{-1/3}N^{-2/3}+\tfrac3{2N}\\
&\le\big((\log N)^L+2\big)m_j^{-1/3}N^{-2/3}=D\,m_j^{-1/3}N^{-2/3}.
\end{align*}
For $|e|\le4$: $\fn_x(e)=\fn_y(\kappa_te)$ (as $\kappa_t>0$), $|\kappa_te|\le4\le5$, and $0\le1-\kappa_t=\kappa_t(\kappa_t^{-1}-1)\le\frac1{2N}$; since $\nsc$ is $\frac1\pi$-Lipschitz (RA5),
\[
|\fn_x(e)-\nsc(e)|\le|\fn_y(\kappa_te)-\nsc(\kappa_te)|+\frac{(1-\kappa_t)|e|}{\pi}<\frac{(\log N)^L}N+\frac{4}{2\pi N}<\frac DN .
\]
So $x\in\Xi$.

\emph{Step 3 (probability).} By \eqref{c2:Nstar}, $(\log N)^{\phi L}=(\log N)^{\phi A_0\log\log N}\ge(\log N)^2$. Hence, using $\log N\ge1$ and the choice of $N_K$,
\begin{align*}
2C_R\exp\big(-c_R(\log N)^{\phi L}\big)&\le\exp\big(\log(2C_R)-c_R(\log N)^2\big)\\
&\le\exp\big(\log(2C_R)-(K+\max(0,\log(2C_R)))\log N\big)\le N^{-K}.
\end{align*}
With Steps 1 and 2, $\P(\lambda(H_t)\notin\Xi)\le\P(\lambda(\widetilde H_t)\notin S)\le N^{-K}$. For $t=\infty$, $\lambda(H_\infty)=\lambda(V)$ has law $\mu$.
(The events are measurable: $\Xi$ and $S$ are Borel and $\lambda$ is continuous.)
\end{proof}

\begin{lemma}[Assumption III: the quantity $Q$]\label{lem:Q}
There is $N_3=N_3(\vartheta,\delta_\pm,C_0)$ such that for all $N\ge N_3$, $\K$, $\nu\in\mathcal E_N(\vartheta,\delta_-,\delta_+,C_0)$ and all $t\in[0,\infty]$,
\begin{equation}\label{Qbound}
\E|\lambda(H_t)-\gamma|^2\le 7D^2N^{-1},\qquad\text{in particular}\qquad \int|x-\gamma|^2\mu(\rd x)\le7D^2N^{-1}.
\end{equation}
We write $Q:=7D^2N^{-1}$.
\end{lemma}

\begin{proof}
Let $M:=(360/\tilde\vartheta)^{2/\tilde\vartheta}$ (depends only on $\vartheta$) and $N_3:=\max(N_6,\lceil\frac83M+8\rceil)$, with $N_6$ from Lemma \ref{lem:goodC} for $K=6$. Fix $N\ge N_3$, $\nu$, $t$; write $m_j:=\min(j,N+1-j)$.

\emph{On $\Xi$.} Each value $k\ge1$ of $m_j$ is taken for at most two indices $j$, and $m_j\le\lceil N/2\rceil$. For $m\ge1$, $\sum_{k=1}^mk^{-2/3}\le1+\int_1^mu^{-2/3}\rd u=3m^{1/3}-2\le3m^{1/3}$ (as $k^{-2/3}\le\int_{k-1}^ku^{-2/3}\rd u$ for $k\ge2$). Hence for $x\in\Xi$:
\[
|x-\gamma|^2\le D^2N^{-4/3}\sum_{j=1}^Nm_j^{-2/3}\le D^2N^{-4/3}\cdot2\cdot3\lceil N/2\rceil^{1/3}\le6D^2N^{-1}.
\]

\emph{Off $\Xi$.} By the spectral theorem $\sum_j\lambda_j(H_t)^2=\Tr H_t^2$; with $|\gamma_j|\le2$ and $\kappa_t^{-2}\le1+\frac1N\le\frac43$,
\[
|\lambda(H_t)-\gamma|^2\le2|\lambda(H_t)|^2+2|\gamma|^2\le2\kappa_t^{-2}\Tr\widetilde H_t^2+8N\le Y:=\tfrac83\Tr\widetilde H_t^2+8N .
\]
Write $\tilde h_{ij}$ for the entries of $\widetilde H_t$ and $\tilde\sigma_{ij}^2=\E|\tilde h_{ij}|^2$. By Lemma \ref{lem:U11}(b), for $i\le j$ the law of $\tilde h_{ij}$ satisfies (D) with $\tilde\vartheta\le1$, so (RA6) with $q=4$ gives $\E|\tilde h_{ij}|^4\le(4C_{\tilde\vartheta})^{4/\tilde\vartheta}\tilde\sigma_{ij}^4=M^2\tilde\sigma_{ij}^4$; the same holds for $i>j$ since $|\tilde h_{ij}|=|\tilde h_{ji}|$.
By Minkowski's inequality in $L^2(\P)$ and (A) for $\tilde\nu^t$,
\[
\big(\E(\Tr\widetilde H_t^2)^2\big)^{1/2}\le\sum_{i,j=1}^N\big(\E|\tilde h_{ij}|^4\big)^{1/2}\le M\sum_{i,j}\tilde\sigma_{ij}^2=MN,\qquad\text{so}\qquad(\E Y^2)^{1/2}\le\big(\tfrac83M+8\big)N .
\]
By the Cauchy--Schwarz inequality and Lemma \ref{lem:goodC} with $K=6$,
\[
\E\big[|\lambda(H_t)-\gamma|^2\mathbf 1\{\lambda(H_t)\notin\Xi\}\big]\le(\E Y^2)^{1/2}\,\P(\lambda(H_t)\notin\Xi)^{1/2}\le\big(\tfrac83M+8\big)N\cdot N^{-3}\le N^{-1}\le D^2N^{-1},
\]
since $N\ge\frac83M+8$ and $D\ge1$. Adding the two parts gives $\E|\lambda(H_t)-\gamma|^2\le7D^2N^{-1}$. The case $t=\infty$ is the statement for $\mu$.
\end{proof}

\begin{lemma}[Deterministic consequences of $x\in\Xi$]\label{lem:countC}
Let $N\ge3$, let $D\ge0$ be any number and $\Xi$ be defined by \eqref{good} with this $D$. Let $0<c\le2$, $\varrho_*:=\varrho_*(c)$, $\ell>0$, and assume
\begin{equation}\tag{H$_\ell$}\label{c2:counthyp}
\frac{D+\ell}{N\varrho_*}\le\frac c8 .
\end{equation}
Let $x\in\Xi$.
\begin{enumerate}[label=(\alph*),itemsep=1pt]
\item (local density) For every $e\in[-3,3]$: $\#\{i:|x_i-e|\le\ell/N\}\le 2\ell/\pi+2D=:K_\ell$.
\item (bulk localisation) If $x_i\in[-2+\frac38c,2-\frac38c]$ or $\gamma_i\in[-2+\frac38c,2-\frac38c]$, then $|x_i-\gamma_i|\le\xi_0:=D/(N\varrho_*)$.
\item (index sets) For $E\in[-2+c,2-c]$ and $b\in(0,c/2)$, the set $J:=\{i:\gamma_i\in[E-b,E+b]\}$ has $|J|\le 2Nb/\pi+1$, and
$\mathcal B:=\{i:\operatorname{dist}(\gamma_i,\{E-b,E+b\})\le\xi_0+\ell/N\}$ has $|\mathcal B|\le 4(D/\varrho_*+\ell)/\pi+2$.
\end{enumerate}
\end{lemma}

\begin{proof}
Note $0<\varrho_*\le\frac{\sqrt2}{2\pi}<\frac1\pi$ (since $c-c^2/16\le c\le2$), so \eqref{c2:counthyp} gives $\ell/N\le\varrho_*c/8<\frac1{4\pi}<1$ and $D/N<\varrho_*c/8$.
We use repeatedly: for $u<w$, $\#\{i:x_i\in(u,w]\}=N(\fn_x(w)-\fn_x(u))$, and, for $|u|,|w|\le4$, $|\fn_x(\cdot)-\nsc(\cdot)|\le D/N$ at $u$ and at $w$ (definition of $\Xi$).

(a) Let $0<\eta<1-\ell/N$, $u:=e-\ell/N-\eta$ and $w:=e+\ell/N$; then $-4\le u<w\le4$. Since $[e-\ell/N,e+\ell/N]\subset(u,w]$ and $\nsc$ is $\frac1\pi$-Lipschitz (RA5),
\[
\#\{i:|x_i-e|\le\ell/N\}\le N(\fn_x(w)-\fn_x(u))\le N(\nsc(w)-\nsc(u))+2D\le\frac{2\ell+N\eta}{\pi}+2D .
\]
Letting $\eta\downarrow0$ gives the claim.

(b) \emph{Step 1: $|x_i|\le4$.} If $x_i\in[-2+\frac38c,2-\frac38c]$ this is clear. Let $\gamma_i\in[-2+\frac38c,2-\frac38c]$. The interval $[-2+\frac c4,-2+\frac38c]$ has length $c/8$ and lies in $\{|y|\le2-c/4\}$, where $\rhosc\ge\varrho_*$ (RA5); hence
$\nsc(-2+\frac38c)\ge\varrho_*c/8>D/N$, and by symmetry ($\nsc(-y)=1-\nsc(y)$) $1-\nsc(2-\frac38c)>D/N$. As $\nsc$ is nondecreasing and $\nsc(\gamma_i)=i/N$:
\[
\frac iN>\frac DN\qquad\text{and}\qquad1-\frac iN>\frac DN .
\]
If $x_i>4$, then $\#\{j:x_j\le4\}\le i-1$, so $\fn_x(4)\le\frac{i-1}N$ while $\nsc(4)=1$; the definition of $\Xi$ gives $1-\frac{i-1}N\le\frac DN$, contradicting $1-\frac iN>\frac DN$.
If $x_i<-4$, then $\fn_x(-4)\ge\frac iN$ while $\nsc(-4)=0$; so $\frac iN\le\frac DN$, a contradiction. Hence $|x_i|\le4$.

\emph{Step 2: $|\nsc(x_i)-\nsc(\gamma_i)|\le D/N$.} Since $x_1\le\dots\le x_i$, $\fn_x(x_i)\ge\frac iN$, so $\nsc(x_i)\ge\fn_x(x_i)-\frac DN\ge\frac iN-\frac DN=\nsc(\gamma_i)-\frac DN$.
For the other inequality: if $x_i=-4$ then $\nsc(x_i)=0\le\nsc(\gamma_i)+\frac DN$. Otherwise, for $e\in[-4,x_i)$, $\fn_x(e)\le\frac{i-1}N$ (only $x_j$ with $j<i$ can be $\le e$), so $\nsc(e)\le\frac{i-1}N+\frac DN$; letting $e\uparrow x_i$ and using continuity, $\nsc(x_i)\le\nsc(\gamma_i)+\frac DN$.

\emph{Step 3: conclusion.} Let $p\in\{x_i,\gamma_i\}$ be a point of $[-2+\frac38c,2-\frac38c]$ and $q$ the other point. The interval $[p-\frac c8,p+\frac c8]$ lies in $\{|y|\le2-\frac c4\}$, where $\rhosc\ge\varrho_*$.
If $|q-p|>c/8$, then the segment between $p$ and $q$ contains a subinterval of length $c/8$ of $[p-\frac c8,p+\frac c8]$, so $|\nsc(q)-\nsc(p)|\ge\varrho_*c/8>D/N$ (monotonicity of $\nsc$), contradicting Step 2.
Hence $|q-p|\le c/8$, the segment between $p$ and $q$ lies in $\{|y|\le2-\frac c4\}$, and $D/N\ge|\nsc(q)-\nsc(p)|=|\int_p^q\rhosc|\ge\varrho_*|q-p|$, i.e.\ $|x_i-\gamma_i|\le D/(N\varrho_*)=\xi_0$.

(c) By (RA5), $|J|=\#\{i:\gamma_i\in[E-b,E+b]\}\le N\cdot2b/\pi+1$. With $\rho:=\xi_0+\ell/N$, $\{y:\operatorname{dist}(y,\{E-b,E+b\})\le\rho\}=[E-b-\rho,E-b+\rho]\cup[E+b-\rho,E+b+\rho]$ is a union of two intervals of length $2\rho$; by (RA5),
$|\mathcal B|\le2(2N\rho/\pi+1)=4(N\xi_0+\ell)/\pi+2=4(D/\varrho_*+\ell)/\pi+2$.
\end{proof}

\begin{auxremark}[Use of the constants]\label{c2:rem}
In \S\ref{sec:C8}, Lemma \ref{lem:countC} is applied with $D$ from \eqref{defL} and $\ell=\rho_O/\varrho_*(c)$ (enlarged to $\ell\ge1$ if desired); \eqref{c2:counthyp} then holds for $N\ge N(c,\rho_O,\vartheta,\delta_\pm,C_0)$ by \eqref{c2:Dsmall}.
All thresholds of this section ($N_*$, $N_K$, $N_3$) and the constants $A_0$, $D$, $Q$ depend only on $(\vartheta,\delta_\pm,C_0)$ (and on $K$ for $N_K$); none depends on $t\in[0,\infty]$, on $\K$, on the ensemble, or on $O$.
\end{auxremark}

\section{\INFRA\ The Ornstein--Uhlenbeck (Mehler) semigroup on a Euclidean space}\label{sec:C3}

This section is pure Gaussian calculus on a finite-dimensional Euclidean space. It will be applied with $\mathcal E=\cH$, $\sigma^2=2/(\beta N)$ and $g=\bG$.
Apart from \S\ref{c3:ss:GE}, where the Gaussian ensemble on $\cH$ is described, nothing from the other sections of Part~III is used.
Every differentiation under an integral sign is justified by an explicit dominating function; the constants in these dominating functions are explicit, but they depend on $d$, $\sigma$ and a lower bound $u_1$ for the time variable, and they are used only qualitatively (Remark \ref{rem:M4}).

\medskip\noindent\emph{Conventions.}
$(\mathcal E,\la\cdot,\cdot\ra)$ is a real inner product space of dimension $d\ge1$ with norm $|\cdot|$ and its Borel $\sigma$-algebra. Integrals of $[0,\infty]$-valued Borel functions are Lebesgue integrals with values in $[0,\infty]$; ``Borel $\Phi\ge0$'' means Borel $\Phi:\mathcal E\to[0,\infty]$.
For $f\in C^k(\mathcal E)$, $y\in\mathcal E$ and $e_1,\dots,e_k\in\mathcal E$ we write $\partial_{e_1}\cdots\partial_{e_k}f(y)$ for the iterated directional derivative, and
\[
\mathsf N_kf(y):=\sup\{|\partial_{e_1}\cdots\partial_{e_k}f(y)|:\ |e_1|,\dots,|e_k|\le1\}
\]
(the norm of the $k$-th derivative as a $k$-linear form; $\mathsf N_0f=|f|$). If $(f_i)$ is an orthonormal basis then $\Delta f=\sum_i\partial_{f_i}\partial_{f_i}f$ (independent of the basis), so $|\Delta f|\le d\,\mathsf N_2f$; and $|\la y,\nabla f(y)\ra|=|\partial_yf(y)|\le|y|\,\mathsf N_1f(y)$.
A function of $(u,y)$ is differentiated in $y$ by $\partial_e,\nabla,\Delta$ and in $u$ by $\partial_u$.

\subsection{Lebesgue and Gaussian measures on a Euclidean space}\label{c3:ss:gauss}

Let $\mathbf e=(e_1,\dots,e_d)$ be an orthonormal basis of $\mathcal E$ and $J_{\mathbf e}:\R^d\to\mathcal E$, $J_{\mathbf e}(x):=\sum_kx_ke_k$. Then $J_{\mathbf e}$ is a linear bijection with $|J_{\mathbf e}x|=|x|$ (Euclidean norm of $\R^d$), and $J_{\mathbf e}^{-1}y=(\la y,e_k\ra)_{k\le d}$.
\emph{Lebesgue measure} on $\mathcal E$ is $\lambda_{\mathcal E}:=(J_{\mathbf e})_*\mathrm{Leb}^d$. It does not depend on $\mathbf e$: for another orthonormal basis $\mathbf e'$, $M:=J_{\mathbf e}^{-1}J_{\mathbf e'}$ is a linear isometry of $\R^d$, i.e.\ an orthogonal matrix, so $|\det M|=1$ and $M_*\mathrm{Leb}^d=|\det M|^{-1}\mathrm{Leb}^d=\mathrm{Leb}^d$; hence $(J_{\mathbf e'})_*\mathrm{Leb}^d=(J_{\mathbf e})_*M_*\mathrm{Leb}^d=(J_{\mathbf e})_*\mathrm{Leb}^d$. Consequently:
\begin{enumerate}[label=(L\arabic*),leftmargin=2.4em,itemsep=1pt]
\item\label{c3:L1} $\mathcal O_*\lambda_{\mathcal E}=\lambda_{\mathcal E}$ for every linear isometry $\mathcal O$ of $\mathcal E$ (as $\mathcal O\circ J_{\mathbf e}=J_{\mathcal O\mathbf e}$ and $\mathcal O\mathbf e$ is an orthonormal basis; $\mathcal O$ is bijective because it is injective and $\dim\mathcal E<\infty$);
\item\label{c3:L2} for $m\in\mathcal E$ and $c\in\R\setminus\{0\}$ the map $y\mapsto m+cy$ maps $\lambda_{\mathcal E}$ to $|c|^{-d}\lambda_{\mathcal E}$ (translation invariance and scaling of $\mathrm{Leb}^d$);
\item\label{c3:L3} on $\mathcal E^2:=\mathcal E\times\mathcal E$ with the inner product $\la(y,z),(y',z')\ra:=\la y,y'\ra+\la z,z'\ra$ one has $\lambda_{\mathcal E^2}=\lambda_{\mathcal E}\otimes\lambda_{\mathcal E}$ (use the orthonormal basis $((e_k,0))_k\cup((0,e_k))_k$, whose coordinate map is $J_{\mathbf e}\times J_{\mathbf e}$, and $\mathrm{Leb}^{2d}=\mathrm{Leb}^d\otimes\mathrm{Leb}^d$).
\end{enumerate}

\begin{auxdefinition}[Gaussian measures]\label{c3:def:gauss}
For $s>0$ let $\gamma_s(y):=(2\pi s)^{-d/2}\exp(-|y|^2/(2s))$, $y\in\mathcal E$, and $g_s:=\gamma_s\,\lambda_{\mathcal E}$; let $g_0:=\delta_0$.
For $m\in\mathcal E$ and $s\ge0$, $N(m,s)$ denotes the image of $g_s$ under $y\mapsto m+y$. (Thus $g_s=N(0,s)$ is the centred Gaussian measure with covariance $s\,\mathrm{Id}$; $g_1$ is the standard Gaussian measure of $\mathcal E$, which is Mathlib's \texttt{stdGaussian}.)
When the space has to be indicated we write $\gamma^{\mathcal E}_s$, $g^{\mathcal E}_s$.
\end{auxdefinition}

\begin{auxlemma}[Gaussian measures on $\mathcal E$]\label{c3:gauss}
Let $s>0$.
\begin{enumerate}[label=(\alph*),itemsep=1pt]
\item $g_s$ is a probability measure, and for every orthonormal basis $\mathbf e$, $(J_{\mathbf e}^{-1})_*g_s=\bigotimes_{k=1}^d\mathrm{gaussianReal}(0,s)$. Equivalently: a random vector $Y$ in $\mathcal E$ has law $g_s$ if and only if its coordinates $\la Y,e_1\ra,\dots,\la Y,e_d\ra$ are independent with law $N(0,s)$ each.
\item $\mathcal O_*g_s=g_s$ for every linear isometry $\mathcal O$ of $\mathcal E$.
\item For $m\in\mathcal E$ and $c\ne0$ the image of $g_s$ under $y\mapsto m+cy$ is $\gamma_{c^2s}(\cdot-m)\,\lambda_{\mathcal E}$; in particular $N(m,s)=\gamma_s(\cdot-m)\lambda_{\mathcal E}$. For every $s'\ge0$ the image of $g_1$ under $y\mapsto\sqrt{s'}\,y$ is $g_{s'}$.
\item $g^{\mathcal E}_s\otimes g^{\mathcal E}_s=g^{\mathcal E^2}_s$. For $\alpha,\beta\in\R$ with $\alpha^2+\beta^2=1$ the map $R_{\alpha,\beta}(y,z):=(\alpha y+\beta z,\ \beta y-\alpha z)$ is a linear isometry of $\mathcal E^2$ with $R_{\alpha,\beta}\circ R_{\alpha,\beta}=\mathrm{id}$, and $(R_{\alpha,\beta})_*(g_s\otimes g_s)=g_s\otimes g_s$.
\item For $\alpha,\beta\in\R$ and $s'\ge0$ the image of $g_{s'}\otimes g_{s'}$ under $(y,z)\mapsto\alpha y+\beta z$ is $g_{(\alpha^2+\beta^2)s'}$.
\item If $Y\sim g_s$ then $\E Y=0$ (coordinatewise), $\E|Y|^2=ds$, and $\E|Y|^p<\infty$ for every $p\ge0$; more precisely $Y\overset{d}=\sqrt s\,Z$ with $Z\sim g_1$, and $m_p:=\int(1+|z|)^p\,g_1(\rd z)<\infty$ depends only on $(d,p)$.
\end{enumerate}
\end{auxlemma}

\begin{proof}
(a) Write $\varphi_s(t):=(2\pi s)^{-1/2}e^{-t^2/(2s)}$ for the density of $\mathrm{gaussianReal}(0,s)$ on $\R$. Since $|J_{\mathbf e}x|^2=\sum_kx_k^2$, we have $\gamma_s(J_{\mathbf e}x)=\prod_k\varphi_s(x_k)$, and therefore, for Borel $B\subset\R^d$,
$(J_{\mathbf e}^{-1})_*g_s(B)=\int_{\mathcal E}\mathbf 1_B(J_{\mathbf e}^{-1}y)\gamma_s(y)\,\lambda_{\mathcal E}(\rd y)=\int_{\R^d}\mathbf 1_B(x)\prod_k\varphi_s(x_k)\,\rd x$.
For a box $B=\prod_kB_k$ Tonelli gives $\prod_k\int_{B_k}\varphi_s$, which is the value of the product measure $\bigotimes_k\mathrm{gaussianReal}(0,s)$ on $B$; two finite measures on $\R^d$ that agree on boxes agree (uniqueness of product measures). With $B=\R^d$ the total mass is $1$.
The law of $Y$ is $g_s$ iff the law of $J_{\mathbf e}^{-1}Y=(\la Y,e_k\ra)_k$ is $(J_{\mathbf e}^{-1})_*g_s$ (as $J_{\mathbf e}$ is a measurable bijection with measurable inverse), i.e.\ the product measure; and the joint law of finitely many real random variables is the product of their laws iff they are independent.

(b) $\gamma_s(\mathcal Oy)=\gamma_s(y)$ since $|\mathcal Oy|=|y|$. For Borel $B$, using \ref{c3:L1},
$g_s(\mathcal O^{-1}B)=\int\mathbf 1_B(\mathcal Oy)\gamma_s(\mathcal Oy)\,\lambda_{\mathcal E}(\rd y)=\int\mathbf 1_B(y')\gamma_s(y')\,(\mathcal O_*\lambda_{\mathcal E})(\rd y')=g_s(B)$.

(c) Let $T(y)=m+cy$. By \ref{c3:L2}, for Borel $\Theta\ge0$: $\int\Theta(T y)\gamma_s(y)\lambda_{\mathcal E}(\rd y)=\int\Theta(y')\gamma_s(T^{-1}y')|c|^{-d}\lambda_{\mathcal E}(\rd y')$, and
$|c|^{-d}\gamma_s((y'-m)/c)=|c|^{-d}(2\pi s)^{-d/2}e^{-|y'-m|^2/(2c^2s)}=\gamma_{c^2s}(y'-m)$. For the last claim take $m=0$, $c=\sqrt{s'}$, $s=1$ if $s'>0$; if $s'=0$ the image is $\delta_0=g_0$.

(d) $\gamma_s(y)\gamma_s(z)=(2\pi s)^{-d}e^{-(|y|^2+|z|^2)/(2s)}=\gamma_s^{\mathcal E^2}(y,z)$, and $\lambda_{\mathcal E^2}=\lambda_{\mathcal E}\otimes\lambda_{\mathcal E}$ by \ref{c3:L3}; so $g_s\otimes g_s=\gamma_s^{\mathcal E^2}\lambda_{\mathcal E^2}=g^{\mathcal E^2}_s$ (again by uniqueness of product measures, checking on rectangles with Tonelli).
$R=R_{\alpha,\beta}$ is linear, and
$|\alpha y+\beta z|^2+|\beta y-\alpha z|^2=(\alpha^2+\beta^2)(|y|^2+|z|^2)+2\alpha\beta\la y,z\ra-2\alpha\beta\la y,z\ra=|y|^2+|z|^2$, so $R$ is a linear isometry of $\mathcal E^2$.
$R(R(y,z))=(\alpha(\alpha y+\beta z)+\beta(\beta y-\alpha z),\ \beta(\alpha y+\beta z)-\alpha(\beta y-\alpha z))=(y,z)$. Apply (b) on $\mathcal E^2$.

(e) If $\alpha=\beta=0$ the image is $\delta_0=g_0$. If $s'=0$ the image of $\delta_{(0,0)}$ is $\delta_0=g_0$. Otherwise let $r:=(\alpha^2+\beta^2)^{1/2}>0$ and $\pi_1(y,z):=y$. Then $\alpha y+\beta z=r\,\pi_1(R_{\alpha/r,\beta/r}(y,z))$. By (d), $(\pi_1\circ R_{\alpha/r,\beta/r})_*(g_{s'}\otimes g_{s'})=(\pi_1)_*(g_{s'}\otimes g_{s'})=g_{s'}$, and by (c) the image of $g_{s'}$ under $y\mapsto ry$ is $\gamma_{r^2s'}\lambda_{\mathcal E}=g_{r^2s'}$.

(f) By (a) the coordinates $\la Y,e_k\ra$ are independent $N(0,s)$, so $\E\la Y,e_k\ra=0$ and $\E|Y|^2=\sum_k\E\la Y,e_k\ra^2=ds$. By (c), $Y\overset d=\sqrt sZ$ with $Z\sim g_1$. Finally $(1+|z|)^p\le(1+\sum_k|z_k|)^p\le(d+1)^{p}\max(1,\max_k|z_k|)^p\le(d+1)^p(1+\sum_k|z_k|^p)$ in the coordinates $z_k=\la z,e_k\ra$, and the one-dimensional Gaussian has finite absolute moments of every order; so $m_p<\infty$.
\end{proof}

\subsection{The Gaussian ensemble on $\cH$}\label{c3:ss:GE}

We use the notation of \S\ref{ss:notation}: $\K\in\{\R,\C\}$, $\beta=\dim_\R\K$, $\cH=\cH_N(\K)$ with $\la X,Y\ra=\Tr(XY)$, $|X|=(\Tr X^2)^{1/2}$.
The form $\Tr(XY)$ is real on $\cH$ ($\overline{\Tr(XY)}=\Tr((XY)^*)=\Tr(YX)=\Tr(XY)$), bilinear, symmetric, and $\Tr(X^2)=\sum_{i,j}|x_{ij}|^2>0$ for $X\ne0$; so $(\cH,\la\cdot,\cdot\ra)$ is a real inner product space.
For $H=(h_{ij})\in\cH$ and $i<j$, using $\Tr(HE_{ij})=h_{ji}=\overline{h_{ij}}$:
\begin{equation}\label{c3:coord}
\la H,E_{ii}\ra=h_{ii},\qquad \la H,X_{ij}\ra=\tfrac{h_{ji}+h_{ij}}{\sqrt2}=\sqrt2\,\mathrm{Re}\,h_{ij},\qquad \la H,Y_{ij}\ra=\tfrac{\mathrm i(h_{ji}-h_{ij})}{\sqrt2}=\sqrt2\,\mathrm{Im}\,h_{ij}\quad(\beta=2).
\end{equation}
The family $\mathbf b:=(E_{ii})_{i}\cup(X_{ij})_{i<j}\,[\cup(Y_{ij})_{i<j}$ if $\beta=2]$ consists of elements of $\cH$; it is orthonormal (by $\Tr(E_{ab}E_{cd})=\delta_{bc}\delta_{ad}$, products of two elements supported on different index pairs have zero trace; $\Tr X_{ij}^2=\Tr(E_{ii}+E_{jj})/2=1$, $\Tr Y_{ij}^2=-\Tr(-E_{ii}-E_{jj})/2=1$, $\Tr(X_{ij}Y_{ij})=\frac{\mathrm i}2\Tr(E_{jj}-E_{ii})=0$), and it spans $\cH$, because by \eqref{c3:coord}
$H=\sum_ih_{ii}E_{ii}+\sum_{i<j}\big(\sqrt2\,\mathrm{Re}\,h_{ij}\,X_{ij}+\sqrt2\,\mathrm{Im}\,h_{ij}\,Y_{ij}\big)$ (the $Y$-terms only for $\beta=2$; the $(i,j)$ entry of the right side is $\mathrm{Re}\,h_{ij}+\mathrm i\,\mathrm{Im}\,h_{ij}$). Hence $d=\dim\cH=N+\beta N(N-1)/2$, and $\lambda_\cH$ (\S\ref{c3:ss:gauss}) is the Lebesgue measure $\rd H$ of \S\ref{ss:notation}.

\begin{auxlemma}[The Gaussian ensemble]\label{c3:GE}
Let $\sigma^2:=2/(\beta N)$ and $\bG:=g^{\cH}_{\sigma^2}$ (Definition \ref{c3:def:gauss} with $\mathcal E=\cH$).
\begin{enumerate}[label=(\roman*),itemsep=1pt]
\item $\bG(\rd H)=(2\pi\sigma^2)^{-d/2}\exp\big(-\frac{\beta N}{4}\Tr H^2\big)\,\rd H$.
\item A random matrix $V=(v_{ij})$ with values in $\cH$ has law $\bG$ if and only if the $\K$-valued random variables $(v_{ij})_{i\le j}$ are independent with the following laws:
$v_{ii}\sim N(0,2/(\beta N))$ (real); for $i<j$: $v_{ij}\sim N(0,1/N)$ if $\beta=1$, and, if $\beta=2$, the pair $(\sqrt2\,\mathrm{Re}\,v_{ij},\sqrt2\,\mathrm{Im}\,v_{ij})$ has law $N(0,1/N)\otimes N(0,1/N)$.
In particular $\E v_{ij}=0$, $\E|v_{ij}|^2=1/N$ for $i<j$, $\E v_{ii}^2=2/(\beta N)$, and $\E v_{ij}^2=0$ for $i<j$ if $\beta=2$.
\item For every $U\in U_\beta(N)$ the map $\mathcal O_U:H\mapsto UHU^*$ is a linear isometry of $\cH$ onto itself, and $(\mathcal O_U)_*\bG=\bG$.
\item If $V\sim\bG$, then $V$ is an ensemble in the sense of NP (entries $v_{ij}$, $i\le j$, independent, centred, $v_{ji}=\overline{v_{ij}}$, $v_{ii}\in\R$) with variances $\sigma_{ij}^2:=\E|v_{ij}|^2=1/N$ for $i\ne j$ and $\sigma_{ii}^2=2/(\beta N)$. Moreover $\E|V|^2=\E\Tr V^2=d\sigma^2\le N+1$, and $\E|V|^p<\infty$ for all $p\ge0$.
\end{enumerate}
\end{auxlemma}

\begin{proof}
(i) is Definition \ref{c3:def:gauss} with $|H|^2/(2\sigma^2)=\frac{\beta N}4\Tr H^2$.

(ii) By Lemma \ref{c3:gauss}(a) applied to the orthonormal basis $\mathbf b$, $V\sim\bG$ iff the $d$ coordinates $\la V,b\ra$, $b\in\mathbf b$, are independent $N(0,\sigma^2)$. By \eqref{c3:coord} these coordinates are $v_{ii}$, $\xi_{ij}:=\sqrt2\,\mathrm{Re}\,v_{ij}$ and (for $\beta=2$) $\eta_{ij}:=\sqrt2\,\mathrm{Im}\,v_{ij}$; and $v_{ij}=(\xi_{ij}+\mathrm i\eta_{ij})/\sqrt2$ (with $\eta_{ij}:=0$ if $\beta=1$).
``$\Rightarrow$'': the blocks $\{v_{ii}\}$, $\{\xi_{ij},\eta_{ij}\}$ are disjoint subfamilies of an independent family, so measurable functions of different blocks are independent; this gives the independence of $(v_{ij})_{i\le j}$. The laws: $v_{ii}\sim N(0,\sigma^2)=N(0,2/(\beta N))$; for $\beta=1$, $v_{ij}=\xi_{ij}/\sqrt2\sim N(0,\sigma^2/2)=N(0,1/N)$; for $\beta=2$, $(\xi_{ij},\eta_{ij})\sim N(0,\sigma^2)^{\otimes2}=N(0,1/N)^{\otimes2}$.
``$\Leftarrow$'': if the $v_{ij}$, $i\le j$, are independent with these laws, then the blocks $\{v_{ii}\}$, $\{\xi_{ij},\eta_{ij}\}$ are independent (each block is a measurable function of one $v_{ij}$), and inside a block the joint law is $N(0,\sigma^2)$ resp.\ $N(0,\sigma^2)^{\otimes2}$ (for $\beta=1$: $\xi_{ij}=\sqrt2v_{ij}\sim N(0,2/N)=N(0,\sigma^2)$). So the joint law of all coordinates is $N(0,\sigma^2)^{\otimes d}$, i.e.\ they are independent $N(0,\sigma^2)$.
Moments: $\E|v_{ij}|^2=(\E\xi_{ij}^2+\E\eta_{ij}^2)/2=\beta\sigma^2/2=1/N$; for $\beta=2$, $\E v_{ij}^2=(\E\xi_{ij}^2-\E\eta_{ij}^2)/2+\mathrm i\,\E\xi_{ij}\eta_{ij}=0$.

(iii) $\mathcal O_U$ is $\R$-linear and maps $\cH$ into $\cH$: $(UHU^*)^*=UH^*U^*=UHU^*$, and for $\beta=1$ ($U$ real orthogonal) $UHU^{T}$ is real. It preserves the inner product: $\Tr(UXU^*\,UYU^*)=\Tr(UXYU^*)=\Tr(XYU^*U)=\Tr(XY)$. It is bijective with inverse $C_{U^*}$. Apply Lemma \ref{c3:gauss}(b).

(iv) Independence, centring and the variances are (ii); $v_{ji}=\overline{v_{ij}}$ and $v_{ii}\in\R$ hold because $V\in\cH$.
$\E|V|^2=d\sigma^2$ by Lemma \ref{c3:gauss}(f), and $d\sigma^2=\frac{2}{\beta N}(N+\frac{\beta N(N-1)}2)=\frac2\beta+N-1\le N+1$. Finite moments: Lemma \ref{c3:gauss}(f).
\end{proof}

\begin{auxremark}\label{c3:rem:I1}
In Lemma \ref{c3:GE}, (ii) is the coordinate description of \S\ref{ss:notation}, (iii) is the conjugation invariance, (iv) says that $V$ is an ensemble. No complex Gaussian measure is used: for $\beta=2$ the off-diagonal entry is the image of the pair $(\xi_{ij},\eta_{ij})$ of real Gaussians.
\end{auxremark}

\subsection{The Mehler semigroup}\label{c3:ss:mehler}

From now on $\sigma>0$ is fixed, $g:=g_{\sigma^2}$, so $g(\rd y)=\gamma_{\sigma^2}(y)\,\lambda_{\mathcal E}(\rd y)$ with $\gamma_{\sigma^2}(y)=(2\pi\sigma^2)^{-d/2}e^{-|y|^2/2\sigma^2}$, and for $u\ge0$
\[
a_u:=e^{-u/2},\qquad s_u:=1-e^{-u},\qquad a_u^2+s_u=1,\qquad a_ua_v=a_{u+v},\qquad a_v^2s_u+s_v=s_{u+v}.
\]
For $u\ge0$ and Borel $\Phi:\mathcal E\to[0,\infty]$, and also for Borel $\Phi:\mathcal E\to\R$ whenever the integral converges absolutely (e.g.\ for bounded $\Phi$), define
\begin{equation}\label{Mehler}
(T_u\Phi)(y):=\int\Phi\big(a_uy+s_u^{1/2}z\big)\,g(\rd z)=\int\Phi(w)\,\gamma_{\sigma^2s_u}(w-a_uy)\,\lambda_{\mathcal E}(\rd w)\quad(u>0),
\qquad A:=\frac{\sigma^2}{2}\Delta-\frac12\la y,\nabla\ra .
\end{equation}
The second expression follows from Lemma \ref{c3:gauss}(c): for $u>0$ the image of $g$ under $z\mapsto a_uy+s_u^{1/2}z$ is $\gamma_{\sigma^2s_u}(\cdot-a_uy)\lambda_{\mathcal E}$. For $u=0$, $T_0\Phi=\Phi$.
Since $(y,z)\mapsto\Phi(a_uy+s_u^{1/2}z)$ is Borel on $\mathcal E^2$, $y\mapsto T_u\Phi(y)$ is Borel (Tonelli; Fubini for bounded $\Phi$), and $|T_u\Phi|\le T_u|\Phi|$.
For a Borel probability measure $\rho$ on $\mathcal E$ and $u\ge0$ let
\[
\rho_u:=(m_u)_*(\rho\otimes g),\qquad m_u(y,z):=a_uy+s_u^{1/2}z ,
\]
i.e.\ $\rho_u=\mathrm{law}(a_uY+s_u^{1/2}Z)$ with $Y\sim\rho$, $Z\sim g$ independent. In the application $\rho=\mathrm{law}(H_0)$ and $\rho_t=\mathrm{law}(H_t)$, by \eqref{Ht}.

\begin{auxlemma}[The Mehler kernel]\label{c3:kernel}
For $u>0$ and $y,w\in\mathcal E$ let
\begin{equation}\label{c3:K}
K(u,y,w):=s_u^{-d/2}\exp\Big(\frac{2a_u\la y,w\ra-a_u^2(|y|^2+|w|^2)}{2\sigma^2s_u}\Big).
\end{equation}
\begin{enumerate}[label=(\alph*),itemsep=1pt]
\item $K>0$, $K(u,y,w)=K(u,w,y)$, and $\gamma_{\sigma^2}(w)K(u,y,w)=\gamma_{\sigma^2s_u}(w-a_uy)$, $\gamma_{\sigma^2}(y)K(u,y,w)=\gamma_{\sigma^2s_u}(y-a_uw)$.
\item For $u>0$, $y\in\mathcal E$ and Borel $\Phi\ge0$ (or Borel $\Phi$ with $\int|\Phi|K(u,y,\cdot)\rd g<\infty$): $T_u\Phi(y)=\int\Phi(w)K(u,y,w)\,g(\rd w)$.
\item $K\in C^\infty((0,\infty)\times\mathcal E\times\mathcal E)$. With $v=v(u,y,w):=a_u(w-a_uy)/(\sigma^2s_u)$ and $c_u:=a_u^2/(\sigma^2s_u)$:
$\ \partial_eK=\la e,v\ra K$ and $\partial_{e'}\la e,v\ra=-c_u\la e,e'\ra$ for $e,e'\in\mathcal E$ (derivatives in $y$); and $\partial_uK=A_yK$, where $A_y$ is $A$ acting on the variable $y$.
\item $K(u,y,w)\le s_u^{-d/2}e^{|y|^2/(2\sigma^2)}$ and $\gamma_{\sigma^2}(y)K(u,y,w)\le(2\pi\sigma^2s_u)^{-d/2}\exp\big(-(|y|-|w|)_+^2/(2\sigma^2)\big)$.
\end{enumerate}
\end{auxlemma}

\begin{proof}
Write $a=a_u$, $s=s_u$. (a) $\gamma_{\sigma^2s}(w-ay)/\gamma_{\sigma^2}(w)=s^{-d/2}\exp\big(-\frac{|w-ay|^2}{2\sigma^2s}+\frac{|w|^2}{2\sigma^2}\big)$ and
$-|w-ay|^2+s|w|^2=-|w|^2+2a\la y,w\ra-a^2|y|^2+s|w|^2=2a\la y,w\ra-a^2(|y|^2+|w|^2)$ because $s-1=-a^2$. The exponent in \eqref{c3:K} is symmetric in $(y,w)$, which gives the symmetry and the second identity.
(b) By the second expression in \eqref{Mehler} and (a): $T_u\Phi(y)=\int\Phi(w)K(u,y,w)\gamma_{\sigma^2}(w)\lambda_{\mathcal E}(\rd w)$.
(c) $u\mapsto a_u,s_u,s_u^{-1},s_u^{-d/2}$ are $C^\infty$ on $(0,\infty)$ (as $s_u>0$), the exponent in \eqref{c3:K} is a polynomial in $(y,w)$ with these coefficients, and $\exp$ is $C^\infty$. Let $F:=\log K=-\frac d2\log s+\frac{P}{2\sigma^2s}$, $P:=2a\la y,w\ra-a^2(|y|^2+|w|^2)$.
Then $\partial_eF=\frac{2a\la e,w\ra-2a^2\la e,y\ra}{2\sigma^2s}=\la e,v\ra$, and $\partial_{e'}\la e,v\ra=-\frac{a^2}{\sigma^2s}\la e,e'\ra$. Hence $\Delta K=(|v|^2-dc_u)K$ and $\la y,\nabla K\ra=\la y,v\ra K$, so
$A_yK/K=\frac{\sigma^2}2|v|^2-\frac{\sigma^2d}{2}c_u-\frac12\la y,v\ra$.
On the other hand $\frac{\rd a}{\rd u}=-\frac a2$ and $\frac{\rd s}{\rd u}=a^2$, so $\partial_u(-\frac d2\log s)=-\frac{da^2}{2s}=-\frac{\sigma^2d}{2}c_u$, and with $\partial_uP=-a\la y,w\ra+a^2(|y|^2+|w|^2)$:
\[
\partial_u\frac{P}{2\sigma^2s}=\frac{s\,\partial_uP-a^2P}{2\sigma^2s^2}=\frac{-a(s+2a^2)\la y,w\ra+a^2(s+a^2)(|y|^2+|w|^2)}{2\sigma^2s^2}=\frac{a^2|y|^2+a^2|w|^2-a(1+a^2)\la y,w\ra}{2\sigma^2s^2}.
\]
Also $\frac{\sigma^2}2|v|^2-\frac12\la y,v\ra=\frac{a^2(|w|^2-2a\la y,w\ra+a^2|y|^2)-as(\la y,w\ra-a|y|^2)}{2\sigma^2s^2}=\frac{a^2|w|^2+a^2(a^2+s)|y|^2-a(2a^2+s)\la y,w\ra}{2\sigma^2s^2}$, which is the same expression since $a^2+s=1$.
Hence $\partial_uF=A_yK/K$, i.e.\ $\partial_uK=A_yK$.
(d) By (a), $K=s^{-d/2}\exp\big(\frac{|y|^2}{2\sigma^2}-\frac{|y-aw|^2}{2\sigma^2s}\big)\le s^{-d/2}e^{|y|^2/(2\sigma^2)}$. Since $0<a\le1$, $|y-aw|\ge|y|-|w|$, so $|y-aw|\ge(|y|-|w|)_+$; with $s\le1$, $\gamma_{\sigma^2s}(y-aw)\le(2\pi\sigma^2s)^{-d/2}e^{-(|y|-|w|)_+^2/(2\sigma^2)}$.
\end{proof}

\begin{lemma}[Algebraic properties]\label{lem:M1}
For $u,v\ge0$:
\begin{enumerate}[label=(\roman*),itemsep=1pt]
\item $T_u1=1$, $T_u\Phi\ge0$ for Borel $\Phi\ge0$, $T_uT_v\Phi=T_{u+v}\Phi$ for Borel $\Phi\ge0$ and for bounded Borel $\Phi$, and $(\rho_u)_v=\rho_{u+v}$ for every probability measure $\rho$.
\item (symmetry) $\int\Phi\,T_u\Theta\,\rd g=\int\Theta\,T_u\Phi\,\rd g$ in $[0,\infty]$ for Borel $\Phi,\Theta\ge0$; and for real Borel $\Phi,\Theta$ with $\int|\Phi|\,T_u|\Theta|\,\rd g<\infty$ (e.g.\ $\Theta$ bounded and $\Phi\in L^1(g)$), in which case all integrals converge absolutely (here $T_u\Theta$ is finite $g$-a.e.\ on $\{\Phi\ne0\}$, and $\Phi\,T_u\Theta:=0$ where $\Phi=0$; symmetrically for $\Theta\,T_u\Phi$).
\item $T_u(\Phi\circ\mathcal O)=(T_u\Phi)\circ\mathcal O$ for every linear isometry $\mathcal O$ of $\mathcal E$ and every Borel $\Phi\ge0$ or bounded Borel $\Phi$.
\item If $\rho=\Phi g$ with Borel $\Phi\ge0$, $\int\Phi\,\rd g=1$, then $\rho_u=(T_u\Phi)\,g$.
In general, for $u>0$, $\rho_u$ has the Lebesgue density $p(u,y)=\int\gamma_{\sigma^2s_u}(y-a_uw)\,\rho(\rd w)\in(0,(2\pi\sigma^2s_u)^{-d/2}]$, and
$\rho_u=\Psi(u,\cdot)\,g$ with $\Psi(u,y):=p(u,y)/\gamma_{\sigma^2}(y)=\int K(u,y,w)\,\rho(\rd w)\in(0,\infty)$; in particular $\rho_u\ll g$ and $g\ll\rho_u$.
If $\rho=\Phi g$ as above, then $\Psi(u,y)=T_u\Phi(y)$ for \emph{every} $y\in\mathcal E$ and $u>0$.
\item $\int F\,\rd\rho_u=\int T_uF\,\rd\rho$ for Borel $F\ge0$ and for bounded Borel $F$.
\item If $\int|w|^q\rho(\rd w)<\infty$ for some $q\ge0$, then $\int|y|^q\rho_u(\rd y)\le2^q\big(\int|w|^q\rho(\rd w)+\sigma^qs_u^{q/2}\int|z|^qg_1(\rd z)\big)<\infty$; and $\int|y|^2\rho_u(\rd y)=e^{-u}\int|w|^2\rho(\rd w)+s_u\sigma^2d$ whenever $\int|w|^2\rho(\rd w)<\infty$.
\end{enumerate}
\end{lemma}

\begin{proof}
Write $a=a_u$, $b=s_u^{1/2}$.
(v) For Borel $F\ge0$, by definition of the image measure and Tonelli, $\int F\rd\rho_u=\int F\circ m_u\,\rd(\rho\otimes g)=\int\big(\int F(ay+bz)g(\rd z)\big)\rho(\rd y)=\int T_uF\,\rd\rho$. For bounded $F$ apply this to $F^\pm$ (all terms finite).

(i) $T_u1=1$ as $g$ is a probability measure; $T_u\Phi\ge0$ is clear. Let $\Phi\ge0$ be Borel and $y\in\mathcal E$. By Tonelli,
$T_u(T_v\Phi)(y)=\int\!\!\int\Phi\big(a_v(ay+bz)+s_v^{1/2}z'\big)\,g(\rd z')g(\rd z)=\int\Phi(a_{u+v}y+w)\,\pi(\rd w)$,
where $\pi$ is the image of $g\otimes g$ under $(z,z')\mapsto a_vb\,z+s_v^{1/2}z'$. By Lemma \ref{c3:gauss}(e) (with $s'=\sigma^2$), $\pi=g_{\sigma^2(a_v^2s_u+s_v)}=g_{\sigma^2s_{u+v}}$, which by Lemma \ref{c3:gauss}(c) is the image of $g=g_{\sigma^2}$ under $z\mapsto s_{u+v}^{1/2}z$. Hence $T_uT_v\Phi(y)=\int\Phi(a_{u+v}y+s_{u+v}^{1/2}z)g(\rd z)=T_{u+v}\Phi(y)$.
For bounded $\Phi$ apply this to $\Phi^\pm$ (all values finite, and $T_v$ is linear on bounded functions).
For Borel $B$, by (v) twice and the semigroup property for $\mathbf 1_B$: $(\rho_u)_v(B)=\int T_v\mathbf 1_B\,\rd\rho_u=\int T_u(T_v\mathbf 1_B)\,\rd\rho=\int T_{u+v}\mathbf 1_B\,\rd\rho=\rho_{u+v}(B)$.

(ii) Let $\Phi,\Theta\ge0$ be Borel. As $a^2+b^2=1$, $R:=R_{a,b}$ of Lemma \ref{c3:gauss}(d) preserves $g\otimes g$ and $R\circ R=\mathrm{id}$. Put $G(y,z):=\Phi(y)\Theta(\pi_1R(y,z))$, $\pi_1(y,z)=y$, so that $\pi_1R(y,z)=ay+bz$. Then, by Tonelli and the invariance,
\[
\int\Phi\,T_u\Theta\,\rd g=\int G\,\rd(g\otimes g)=\int G\circ R\,\rd(g\otimes g)=\int\Phi(ay+bz)\Theta(y)\,(g\otimes g)(\rd y,\rd z)=\int\Theta\,T_u\Phi\,\rd g ,
\]
using $G(R(y,z))=\Phi(\pi_1R(y,z))\Theta(\pi_1R(R(y,z)))=\Phi(ay+bz)\Theta(y)$.
For real $\Phi,\Theta$ with $\int|\Phi|T_u|\Theta|\rd g<\infty$: the nonnegative identity gives $\int|\Theta|T_u|\Phi|\rd g<\infty$ as well, so the four identities for $\Phi^\pm,\Theta^\pm$ are between finite numbers, and subtracting them gives the claim. If $\Theta$ is bounded and $\Phi\in L^1(g)$, then $\int|\Phi|T_u|\Theta|\rd g\le\|\Theta\|_\infty\int|\Phi|\rd g$.

(iii) By Lemma \ref{c3:gauss}(b), $\mathcal O_*g=g$, so $T_u(\Phi\circ\mathcal O)(y)=\int\Phi(a\mathcal Oy+b\mathcal Oz)g(\rd z)=\int\Phi(a\mathcal Oy+bz')(\mathcal O_*g)(\rd z')=T_u\Phi(\mathcal Oy)$.

(iv) If $\rho=\Phi g$: for Borel $B$, by (v) and (ii), $\rho_u(B)=\int T_u\mathbf 1_B\,\Phi\,\rd g=\int\mathbf 1_B\,T_u\Phi\,\rd g$.
General $\rho$, $u>0$: by (v), the second expression in \eqref{Mehler} and Tonelli, $\rho_u(B)=\int\big(\int\mathbf 1_B(y)\gamma_{\sigma^2s_u}(y-aw)\lambda_{\mathcal E}(\rd y)\big)\rho(\rd w)=\int_Bp(u,y)\lambda_{\mathcal E}(\rd y)$.
The bounds on $p$ hold since $0<\gamma_{\sigma^2s_u}\le(2\pi\sigma^2s_u)^{-d/2}$ and $\rho$ is a probability measure.
By Lemma \ref{c3:kernel}(a), $p(u,y)/\gamma_{\sigma^2}(y)=\int K(u,y,w)\rho(\rd w)$, and $\rho_u(B)=\int_B\Psi(u,y)\gamma_{\sigma^2}(y)\lambda_{\mathcal E}(\rd y)=\int_B\Psi(u,\cdot)\rd g$. Since $0<\Psi<\infty$, $\rho_u$ and $g$ are mutually absolutely continuous.
If $\rho=\Phi g$, then $\Psi(u,y)=\int K(u,y,w)\Phi(w)g(\rd w)=T_u\Phi(y)$ for every $y$, by Lemma \ref{c3:kernel}(b).

(vi) By (v) with $F=|\cdot|^q$: $\int|y|^q\rd\rho_u=\iint|ay+bz|^qg(\rd z)\rho(\rd y)$, and $|ay+bz|^q\le(|y|+b|z|)^q\le2^q(|y|^q+b^q|z|^q)$; finally $g$ is the image of $g_1$ under $z\mapsto\sigma z$ (Lemma \ref{c3:gauss}(c)), and $\int|z|^qg_1(\rd z)<\infty$ (Lemma \ref{c3:gauss}(f)).
For $q=2$: $\int|ay+bz|^2g(\rd z)=a^2|y|^2+2ab\la y,\int z\,g(\rd z)\ra+b^2\int|z|^2g(\rd z)=a^2|y|^2+s_u\sigma^2d$ by Lemma \ref{c3:gauss}(f); integrate in $\rho(\rd y)$ and use $a^2=e^{-u}$.
\end{proof}

\subsection{Derivatives of the Mehler kernel; differentiation under the integral sign}\label{c3:ss:deriv}

\emph{Constants.} Let $t_0=t_1=1$ and $t_{k+1}:=t_k+k\,t_{k-1}$ ($k\ge1$); $t_k$ is the number of partial matchings of $\{1,\dots,k\}$ (sets of pairwise disjoint two-element subsets), and $t_k\le k!$ by induction ($t_{k+1}\le k!+k(k-1)!\le(k+1)!$).
Let
\[
c_{0,k}:=t_k,\qquad c_{j,k}:=\frac{\sigma^2d}{2}\,c_{j-1,k+2}+\frac12\,c_{j-1,k+1}+\frac k2\,c_{j-1,k}\quad(j\ge1,\ k\ge0),\qquad \hat c_n:=\max_{0\le j\le n}c_{j,n-j}.
\]
For $u_1>0$ let
\[
\kappa=\kappa(u_1):=1+\frac{1}{\sigma^2s_{u_1}}+\frac{1}{\sigma s_{u_1}^{1/2}},\qquad \Lambda(y,w):=\kappa\,(1+|y|+|w|)\ \ (\ge1).
\]
These numbers depend only on $(j,k,n,d,\sigma,u_1)$. On $\R\times\mathcal E$ we use the norm $\|(h,e)\|:=\max(|h|,|e|)$; $D^n_x$ denotes the $n$-th Fr\'echet derivative in $x=(u,y)\in(0,\infty)\times\mathcal E$, and $\|\cdot\|$ the norm of $n$-linear forms.

\begin{auxlemma}[Explicit derivative bounds for the kernel]\label{c3:kbounds}
Let $u_1>0$. For all $u\ge u_1$, $y,w\in\mathcal E$, $j,k,n\ge0$ and $e_1,\dots,e_k\in\mathcal E$ with $|e_i|\le1$:
\begin{enumerate}[label=(\alph*),itemsep=1pt]
\item $\partial_u^j\partial_{e_1}\cdots\partial_{e_k}K(u,y,w)=\partial_{e_1}\cdots\partial_{e_k}(A_y^jK)(u,y,w)$;
\item $|\partial_u^j\partial_{e_1}\cdots\partial_{e_k}K(u,y,w)|\le c_{j,k}\,\Lambda(y,w)^{k+2j}\,K(u,y,w)$;
\item $\|D^n_xK(u,y,w)\|\le2^n\hat c_n\,\Lambda(y,w)^{2n}\,K(u,y,w)$.
\end{enumerate}
\end{auxlemma}

\begin{proof}
Fix $w$; all derivatives below are in $(u,y)$, and $K\in C^\infty$ (Lemma \ref{c3:kernel}(c)), so partial derivatives commute.

\emph{Step 1 ($j=0$; a matching formula).} Let $\mathsf M_k$ be the set of partial matchings of $\{1,\dots,k\}$; for $\pi\in\mathsf M_k$ let $U(\pi)$ be the set of unmatched indices. We claim
\begin{equation}\label{c3:match}
\partial_{e_1}\cdots\partial_{e_k}K=K\sum_{\pi\in\mathsf M_k}(-c_u)^{|\pi|}\prod_{\{i,l\}\in\pi}\la e_i,e_l\ra\prod_{i\in U(\pi)}\la e_i,v\ra .
\end{equation}
For $k=0$ this is trivial. Assume it for $k$ and apply $\partial_{e_{k+1}}$ with the product rule and Lemma \ref{c3:kernel}(c): differentiating $K$ produces the factor $\la e_{k+1},v\ra$, i.e.\ the matching $\pi$ of $\{1,..,k+1\}$ with $k+1$ unmatched; differentiating a factor $\la e_i,v\ra$, $i\in U(\pi)$, produces $-c_u\la e_i,e_{k+1}\ra$, i.e.\ the matching $\pi\cup\{\{i,k+1\}\}$. Every partial matching of $\{1,\dots,k+1\}$ arises exactly once in this way (remove $k+1$, or the pair containing it), which proves \eqref{c3:match} for $k+1$; it also shows $|\mathsf M_{k+1}|=|\mathsf M_k|+k|\mathsf M_{k-1}|$, i.e.\ $|\mathsf M_k|=t_k$.
For $u\ge u_1$: $a_u\le1$ and $s_u\ge s_{u_1}$, so $|v|\le\frac{|w|+|y|}{\sigma^2s_{u_1}}$ and $c_u^{1/2}\le\frac1{\sigma s_{u_1}^{1/2}}$, whence $|v|+c_u^{1/2}\le\Lambda$. Each term of \eqref{c3:match} is bounded by $c_u^{|\pi|}|v|^{k-2|\pi|}\le(c_u^{1/2}+|v|)^k\le\Lambda^k$ (as $|e_i|\le1$), so $|\partial_{e_1}\cdots\partial_{e_k}K|\le t_k\Lambda^kK$, i.e.\ $\mathsf N_k(K(u,\cdot,w))(y)\le c_{0,k}\Lambda^kK$.

\emph{Step 2 (proof of (a)).} $A_y$ has coefficients independent of $u$. By induction on $j$, using Lemma \ref{c3:kernel}(c) and the commutation of $\partial_u$ with $y$-derivatives: $\partial_u^jK=\partial_u(A_y^{j-1}K)=A_y^{j-1}\partial_uK=A_y^jK$. Then $\partial_u^j\partial_{e_1}\cdots\partial_{e_k}K=\partial_{e_1}\cdots\partial_{e_k}\partial_u^jK=\partial_{e_1}\cdots\partial_{e_k}A_y^jK$.

\emph{Step 3 (a commutator).} For $f\in C^\infty(\mathcal E)$ and $e\in\mathcal E$: $\partial_e\la y,\nabla f(y)\ra=\partial_e\big(Df(y)[y]\big)=D^2f(y)[y,e]+Df(y)[e]=\la y,\nabla\partial_ef(y)\ra+\partial_ef(y)$, and $\partial_e\Delta f=\Delta\partial_ef$. Hence $\partial_e(Af)=A\partial_ef-\frac12\partial_ef$, and by induction on $k$
\[
\partial_{e_1}\cdots\partial_{e_k}(Af)=A\,\partial_{e_1}\cdots\partial_{e_k}f-\frac k2\,\partial_{e_1}\cdots\partial_{e_k}f .
\]
With $|\Delta F|\le d\,\mathsf N_2F$ and $|\la y,\nabla F(y)\ra|\le|y|\,\mathsf N_1F(y)$ (Conventions) applied to $F=\partial_{e_1}\cdots\partial_{e_k}f$ this gives, for $|e_i|\le1$,
\begin{equation}\label{c3:NA}
\mathsf N_k(Af)(y)\le\frac{\sigma^2d}2\,\mathsf N_{k+2}f(y)+\frac{|y|}2\,\mathsf N_{k+1}f(y)+\frac k2\,\mathsf N_kf(y).
\end{equation}

\emph{Step 4 (proof of (b)).} We show $\mathsf N_k(A_y^jK)\le c_{j,k}\Lambda^{k+2j}K$ for all $k$, by induction on $j$; by (a) this is (b). The case $j=0$ is Step 1. For $j\ge1$ apply \eqref{c3:NA} to $f=A_y^{j-1}K(u,\cdot,w)$ and the induction hypothesis:
$\mathsf N_k(A_y^jK)\le\big[\frac{\sigma^2d}2c_{j-1,k+2}\Lambda^{k+2j}+\frac{|y|}2c_{j-1,k+1}\Lambda^{k+2j-1}+\frac k2c_{j-1,k}\Lambda^{k+2j-2}\big]K\le c_{j,k}\Lambda^{k+2j}K$, using $|y|\le\Lambda$ and $\Lambda\ge1$.

\emph{Step 5 (proof of (c)).} For $\xi_i=(h_i,e_i)\in\R\times\mathcal E$, $\xi_i=h_i(1,0)+(0,e_i)$, multilinearity and symmetry of $D^n_xK$ give
$D^n_xK[\xi_1,\dots,\xi_n]=\sum_{S\subset\{1,\dots,n\}}\big(\prod_{i\in S}h_i\big)\,\partial_u^{|S|}\prod_{i\notin S}\partial_{e_i}K$.
If $\|\xi_i\|\le1$ for all $i$, (b) bounds the term of $S$ by $c_{|S|,n-|S|}\Lambda^{n+|S|}K\le\hat c_n\Lambda^{2n}K$; there are $2^n$ terms.
\end{proof}

\begin{auxlemma}[Smoothing of kernel integrals, with explicit domination]\label{c3:smooth}
Let $\rho$ be a finite Borel measure on $\mathcal E$ with $M_p(\rho):=\int(1+|w|)^p\rho(\rd w)<\infty$ for every $p\ge0$, and let $h:\mathcal E\to\R$ be bounded and Borel. Put $\Omega:=(0,\infty)\times\mathcal E$ and
\[
\Psi_h(u,y):=\int K(u,y,w)\,h(w)\,\rho(\rd w),\qquad (u,y)\in\Omega .
\]
\begin{enumerate}[label=(\alph*),itemsep=1pt]
\item (Domination.) For $u_1>0$, $r\ge0$, $n\ge0$ and all $(u,y)\in\Omega$ with $u\ge u_1$, $|y|\le r$:
\begin{equation}\label{c3:dom}
\|D^n_xK(u,y,w)\,h(w)\|\le G_{n,u_1,r}(w):=\|h\|_\infty\,2^n\hat c_n\,\kappa(u_1)^{2n}\,s_{u_1}^{-d/2}e^{r^2/(2\sigma^2)}\,(1+r+|w|)^{2n},
\end{equation}
and $\int G_{n,u_1,r}\,\rd\rho\le\|h\|_\infty2^n\hat c_n\kappa(u_1)^{2n}s_{u_1}^{-d/2}e^{r^2/(2\sigma^2)}(1+r)^{2n}M_{2n}(\rho)<\infty$.
\item (Smoothness.) $\Psi_h\in C^\infty(\Omega)$ and $D^n\Psi_h(x)=\int D^n_xK(x,w)h(w)\rho(\rd w)$ (absolutely convergent Bochner integral of $n$-linear forms) for all $n$ and $x\in\Omega$. In particular $\partial_u^j\partial_{e_1}\cdots\partial_{e_k}\Psi_h=\int\partial_u^j\partial_{e_1}\cdots\partial_{e_k}K\,h\,\rd\rho$, and $\partial_u\Psi_h=A\Psi_h$ on $\Omega$.
\item (Gaussian-weighted bound.) For $u\ge u_1>0$, $y\in\mathcal E$, $|e_i|\le1$:
\[
\gamma_{\sigma^2}(y)\,|\partial_u^j\partial_{e_1}\cdots\partial_{e_k}\Psi_h(u,y)|\le\|h\|_\infty\,c_{j,k}\,\kappa(u_1)^{k+2j}(2\pi\sigma^2s_{u_1})^{-d/2}\int(1+|y|+|w|)^{k+2j}e^{-(|y|-|w|)_+^2/(2\sigma^2)}\rho(\rd w).
\]
\item (An elementary integral.) For $p,r\ge0$:
$\displaystyle\int_{\mathcal E}(1+|y|+r)^pe^{-(|y|-r)_+^2/(2\sigma^2)}\lambda_{\mathcal E}(\rd y)\le B_p(1+r)^{p+d}$, with
$B_p:=4^d3^p+(8\pi\sigma^2)^{d/2}(1+4\sigma)^pm_p$ ($m_p$ from Lemma \ref{c3:gauss}(f)).
\end{enumerate}
\end{auxlemma}

\begin{proof}
(a) By Lemma \ref{c3:kbounds}(c), $\|D^n_xK\|\le2^n\hat c_n\Lambda^{2n}K$; here $\Lambda\le\kappa(u_1)(1+r+|w|)$, and by Lemma \ref{c3:kernel}(d) $K(u,y,w)\le s_u^{-d/2}e^{|y|^2/(2\sigma^2)}\le s_{u_1}^{-d/2}e^{r^2/(2\sigma^2)}$, because $s_u\ge s_{u_1}$. The integral bound uses $(1+r+|w|)\le(1+r)(1+|w|)$.

(b) For $n\ge0$ and $x\in\Omega$ put $I_n(x):=\int D^n_xK(x,w)h(w)\rho(\rd w)$. The integrand is measurable in $w$ (continuous in $w$ times Borel $h$) and, by (a), dominated by an integrable function locally uniformly in $x$; so $I_n$ is well defined, and $I_0=\Psi_h$.
Fix $x_0=(u_0,y_0)\in\Omega$ and let $B_0:=\{x:\|x-x_0\|<\eps_0\}$, $\eps_0:=\min(u_0/2,1)$. For $x=(u,y)\in B_0$: $u\ge u_0/2=:u_1$ and $|y|\le|y_0|+1=:r$. For each $w$, $x\mapsto D^n_xK(x,w)h(w)$ is differentiable on $B_0$ with derivative $D^{n+1}_xK(x,w)h(w)$, whose norm is $\le G_{n+1,u_1,r}(w)$ by (a), and $G_{n+1,u_1,r}\in L^1(\rho)$.
By the theorem on differentiation under the integral sign (Mathlib, see Appendix~\ref{app:mathlib}), $I_n$ is differentiable at $x_0$ with $DI_n(x_0)=I_{n+1}(x_0)$. By dominated convergence with the dominating function $G_{n,u_1,r}$, $I_n$ is continuous on $B_0$.
Hence by induction $D^n\Psi_h=I_n$ exists and is continuous on $\Omega$ for every $n$, i.e.\ $\Psi_h\in C^\infty(\Omega)$.
Evaluating $D^{j+k}\Psi_h(x)$ at $((1,0)^{\times j},(0,e_1),\dots,(0,e_k))$ (evaluation at fixed arguments is a continuous linear map and commutes with the Bochner integral) gives the formula for $\partial_u^j\partial_{e_1}\cdots\partial_{e_k}\Psi_h$.
Finally, with an orthonormal basis $(f_i)$ of $\mathcal E$ and Lemma \ref{c3:kernel}(c):
$\partial_u\Psi_h(u,y)=\int A_yK\,h\,\rd\rho=\frac{\sigma^2}2\sum_i\int\partial_{f_i}^2K\,h\,\rd\rho-\frac12\int\partial_yK\,h\,\rd\rho=\frac{\sigma^2}2\Delta\Psi_h-\frac12\la y,\nabla\Psi_h\ra=A\Psi_h$
(here $\partial_y$ is the directional derivative in the direction $y$, and $\int\partial_yK\,h\,\rd\rho=\partial_y\Psi_h=\la y,\nabla\Psi_h\ra$).

(c) By (b) and Lemma \ref{c3:kbounds}(b), $|\partial_u^j\partial_{e_1}\cdots\partial_{e_k}\Psi_h(u,y)|\le\|h\|_\infty c_{j,k}\int\Lambda^{k+2j}K\,\rd\rho$. Multiply by $\gamma_{\sigma^2}(y)$ and use Lemma \ref{c3:kernel}(d) together with $s_u\ge s_{u_1}$ and $\Lambda\le\kappa(u_1)(1+|y|+|w|)$.

(d) On $\{|y|\le2r\}$ the integrand is $\le(1+3r)^p\le3^p(1+r)^p$, and $\{|y|\le2r\}\subset J_{\mathbf e}([-2r,2r]^d)$ (as $|\la y,e_k\ra|\le|y|$) has $\lambda_{\mathcal E}$-measure $\le(4r)^d\le4^d(1+r)^d$.
On $\{|y|>2r\}$: $|y|-r>|y|/2$, so $(|y|-r)_+^2\ge|y|^2/4$, and $1+|y|+r\le1+2|y|$; hence this part is at most
$\int(1+2|y|)^pe^{-|y|^2/(8\sigma^2)}\lambda_{\mathcal E}(\rd y)=(8\pi\sigma^2)^{d/2}\int(1+2|y|)^pg_{4\sigma^2}(\rd y)=(8\pi\sigma^2)^{d/2}\int(1+4\sigma|z|)^pg_1(\rd z)$
(Lemma \ref{c3:gauss}(c): $g_{4\sigma^2}$ is the image of $g_1$ under $z\mapsto2\sigma z$), and $(1+4\sigma|z|)\le(1+4\sigma)(1+|z|)$. Add the two parts and use $(1+r)^{p+d}\ge\max\{(1+r)^p(1+r)^d,1\}$.
\end{proof}

\subsection{Smoothing, the equation, and bounds}\label{c3:ss:M2}

\begin{lemma}[Smoothing, the equation, and bounds]\label{lem:M2}
Let $\rho$ be a probability measure on $\mathcal E$ with $\int|w|^m\rho(\rd w)<\infty$ for all $m$, write $M_p:=\int(1+|w|)^p\rho(\rd w)$, $M_2':=\int|w|^2\rho(\rd w)$, and let $\Psi(u,y):=p(u,y)/\gamma_{\sigma^2}(y)=\int K(u,y,w)\rho(\rd w)$ for $u>0$ (Lemma \ref{lem:M1}(iv); so $\rho_u=\Psi(u,\cdot)g$, and $\Psi(u,\cdot)=T_u\Phi$ everywhere if $\rho=\Phi g$). Then:
\begin{enumerate}[label=(\roman*),itemsep=1pt]
\item $\Psi\in C^\infty((0,\infty)\times\mathcal E)$, $\Psi>0$, and $\partial_u\Psi=A\Psi$. All derivatives may be taken under the integral sign: $\partial_u^j\partial_{e_1}\cdots\partial_{e_k}\Psi=\int\partial_u^j\partial_{e_1}\cdots\partial_{e_k}K(\cdot,\cdot,w)\rho(\rd w)$, with the dominating functions \eqref{c3:dom} ($h\equiv1$).
\item For $0<u_1<u_2<\infty$ (in fact for all $u\ge u_1$) and $y\in\mathcal E$: $|\log\Psi(u,y)|\le C_{\log}(1+|y|^2)$ with
$C_{\log}:=\log2+\frac d2\log\frac1{s_{u_1}}+\frac{1+2M_2'}{\sigma^2s_{u_1}}$.
\item With $\mathfrak M(y):=\sup_{u\in[u_1,u_2]}\sum_{j\le1,\,|\alpha|\le3}|\partial_u^j\partial_y^\alpha\Psi(u,y)|$ (coordinate derivatives in an orthonormal basis): $\mathfrak M$ is Borel and $\int\mathfrak M(y)(1+|y|)^m\,g(\rd y)<\infty$ for every $m\ge0$. Quantitatively, for all $j,k,m\ge0$,
\begin{equation}\label{c3:Mbound}
\int\sup_{u\ge u_1}\mathsf N_k\big(\partial_u^j\Psi(u,\cdot)\big)(y)\,(1+|y|)^m\,g(\rd y)\le c_{j,k}\,\kappa(u_1)^{k+2j}\,(2\pi\sigma^2s_{u_1})^{-d/2}\,B_{m+k+2j}\,M_{m+k+2j+d},
\end{equation}
and $\mathfrak M\le\binom{d+3}{3}\sum_{j\le1,\,k\le3}\sup_{u\ge u_1}\mathsf N_k(\partial_u^j\Psi(u,\cdot))$.
\item (Fisher information) For $u>0$: $\displaystyle\int\frac{|\nabla\Psi(u,\cdot)|^2}{\Psi(u,\cdot)}\rd g\le\frac{2d}{\sigma^2s_u}+\frac{2}{\sigma^4}\int|y|^2\rho_u(\rd y)=\frac{2d}{\sigma^2s_u}+\frac{2(e^{-u}M_2'+s_u\sigma^2d)}{\sigma^4}<\infty$.
\end{enumerate}
\end{lemma}

\begin{proof}
(i) Lemma \ref{c3:smooth}(a),(b) with $h\equiv1$ ($\rho$ is finite and $M_p<\infty$ for all $p$); $\Psi>0$ as $K>0$.

(ii) Let $u\ge u_1$, $a=a_u$, $s=s_u\in[s_{u_1},1)$. Upper bound: by Lemma \ref{c3:kernel}(d), $\Psi(u,y)\le s_{u_1}^{-d/2}e^{|y|^2/(2\sigma^2)}$, so $\log\Psi\le\frac d2\log\frac1{s_{u_1}}+\frac{|y|^2}{2\sigma^2}\le C_{\log}(1+|y|^2)$ (as $\frac1{2\sigma^2}\le\frac1{\sigma^2s_{u_1}}$).
Lower bound: let $R_0:=(2M_2')^{1/2}$. By Chebyshev, $\rho(|w|>R_0)\le\frac12$ (if $M_2'=0$ then $\rho=\delta_0$ and $\rho(|w|>0)=0$). For $|w|\le R_0$, by Lemma \ref{c3:kernel}(a) and $s^{-d/2}\ge1$:
$K(u,y,w)=s^{-d/2}\exp\big(\frac{|y|^2}{2\sigma^2}-\frac{|y-aw|^2}{2\sigma^2s}\big)\ge\exp\big(-\frac{2|y|^2+2|w|^2}{2\sigma^2s}\big)\ge\exp\big(-\frac{|y|^2+2M_2'}{\sigma^2s_{u_1}}\big)$.
Hence $\Psi(u,y)\ge\int_{|w|\le R_0}K\,\rd\rho\ge\frac12\exp\big(-\frac{|y|^2+2M_2'}{\sigma^2s_{u_1}}\big)$, i.e.\ $\log\Psi\ge-\log2-\frac{|y|^2+2M_2'}{\sigma^2s_{u_1}}\ge-C_{\log}(1+|y|^2)$.

(iii) \emph{Measurability.} For fixed $y$, $u\mapsto\partial_u^j\partial^\alpha_y\Psi(u,y)$ and $u\mapsto\mathsf N_k(\partial_u^j\Psi(u,\cdot))(y)$ are continuous ($\Psi\in C^\infty$; $\mathsf N_k$ is the norm of the continuous $k$-th derivative), and they are continuous in $y$ for fixed $u$. So the suprema over $u$ in an interval equal the suprema over the countable set of rational points of the interval together with its endpoints, and are therefore Borel (even lower semicontinuous) functions of $y$.
\emph{Bound \eqref{c3:Mbound}.} By Lemma \ref{c3:smooth}(c) with $h\equiv1$, for every $u\ge u_1$ and $|e_i|\le1$, $\gamma_{\sigma^2}(y)|\partial_u^j\partial_{e_1}\cdots\partial_{e_k}\Psi(u,y)|$ is bounded by a quantity that does not depend on $u$ or on the $e_i$; hence the same bound holds for $\gamma_{\sigma^2}(y)\sup_{u\ge u_1}\mathsf N_k(\partial_u^j\Psi(u,\cdot))(y)$. Multiply by $(1+|y|)^m\le(1+|y|+|w|)^m$ inside the $w$-integral, integrate over $\lambda_{\mathcal E}(\rd y)$, exchange the integrals (Tonelli), and apply Lemma \ref{c3:smooth}(d) with $r=|w|$ and $p=m+k+2j$:
the left side of \eqref{c3:Mbound} is at most $c_{j,k}\kappa^{k+2j}(2\pi\sigma^2s_{u_1})^{-d/2}\int B_{m+k+2j}(1+|w|)^{m+k+2j+d}\rho(\rd w)$.
\emph{$\mathfrak M$.} A coordinate derivative $\partial_y^\alpha$ is $\partial_{f_{i_1}}\cdots\partial_{f_{i_{|\alpha|}}}$ with basis vectors $|f_i|=1$, so $|\partial_u^j\partial_y^\alpha\Psi|\le\mathsf N_{|\alpha|}(\partial_u^j\Psi)$, and there are $\binom{d+3}3$ multi-indices with $|\alpha|\le3$. Integrability of $\mathfrak M(1+|y|)^m$ follows from \eqref{c3:Mbound}.

(iv) Fix $u>0$, $a=a_u$, $s=s_u$, $t:=\sigma^2s$. Then $p(u,\cdot)=\gamma_{\sigma^2}\Psi(u,\cdot)\in C^\infty$, $p>0$, and by (i) and the product rule
$\nabla p=\Psi\nabla\gamma_{\sigma^2}+\gamma_{\sigma^2}\nabla\Psi=\int\nabla_y\big(\gamma_{\sigma^2}(y)K(u,y,w)\big)\rho(\rd w)=\int(\nabla\gamma_t)(y-aw)\,\rho(\rd w)$, by Lemma \ref{c3:kernel}(a); here $\nabla\gamma_t(x)=-\frac{x}{t}\gamma_t(x)$.
By Cauchy--Schwarz in $L^2(\rho)$,
$|\nabla p(y)|^2\le\Big(\int\frac{|\nabla\gamma_t|}{\gamma_t^{1/2}}(y-aw)\,\gamma_t^{1/2}(y-aw)\,\rho(\rd w)\Big)^2\le p(y)\int\frac{|\nabla\gamma_t|^2}{\gamma_t}(y-aw)\,\rho(\rd w)$.
Since $\nabla(1/\gamma_{\sigma^2})(y)=\frac{y}{\sigma^2}\frac1{\gamma_{\sigma^2}(y)}$, we have $\nabla\Psi=\big(\nabla p+p\,y/\sigma^2\big)/\gamma_{\sigma^2}$ and
\[
\frac{|\nabla\Psi|^2}{\Psi}\,\gamma_{\sigma^2}=\frac{|\nabla p+p\,y/\sigma^2|^2}{p}\le\frac{2|\nabla p|^2}{p}+\frac{2|y|^2p}{\sigma^4}\le2\int\frac{|\nabla\gamma_t|^2}{\gamma_t}(y-aw)\rho(\rd w)+\frac{2|y|^2p(y)}{\sigma^4}.
\]
Integrate over $\lambda_{\mathcal E}(\rd y)$ and use Tonelli and translation invariance of $\lambda_{\mathcal E}$:
$\int\!\!\int\frac{|\nabla\gamma_t|^2}{\gamma_t}(y-aw)\lambda_{\mathcal E}(\rd y)\rho(\rd w)=\int\frac{|x|^2}{t^2}\gamma_t(x)\lambda_{\mathcal E}(\rd x)=\frac{dt}{t^2}=\frac{d}{\sigma^2s}$ (Lemma \ref{c3:gauss}(f)), and $\int|y|^2p(u,y)\lambda_{\mathcal E}(\rd y)=\int|y|^2\rho_u(\rd y)$, which is computed in Lemma \ref{lem:M1}(vi).
\end{proof}

\begin{lemma}[Bounded data; mixing]\label{lem:M3}
Let $\Phi$ be bounded and Borel. Then $(u,y)\mapsto T_u\Phi(y)$ is $C^\infty$ on $(0,\infty)\times\mathcal E$, $|T_u\Phi|\le\|\Phi\|_\infty$, $\partial_uT_u\Phi=AT_u\Phi$, and for $u_1>0$, all $u\ge u_1$ (in particular on $[u_1,u_2]$), $y\in\mathcal E$, $j,k\ge0$ and $|e_i|\le1$:
\begin{equation}\label{c3:M3bound}
|\partial_u^j\partial_{e_1}\cdots\partial_{e_k}T_u\Phi(y)|\le C_{j,k}(1+|y|)^{2j+k}\|\Phi\|_\infty,\qquad C_{j,k}:=c_{j,k}\,\kappa(u_1)^{k+2j}\int(2+|z|)^{k+2j}g(\rd z).
\end{equation}
In coordinates: $\sup_{u\in[u_1,u_2]}|\partial_u^j\partial_y^\alpha T_u\Phi(y)|\le C_{j,|\alpha|}(1+|y|)^{2j+|\alpha|}\|\Phi\|_\infty$.
If $\Phi$ is moreover continuous, then for every $y$: $T_u\Phi(y)\to\Phi(y)$ as $u\downarrow0$ and $T_u\Phi(y)\to\int\Phi\,\rd g$ as $u\to\infty$ (and $|T_u\Phi(y)|\le\|\Phi\|_\infty$ throughout).
\end{lemma}

\begin{proof}
By Lemma \ref{c3:kernel}(b), $T_u\Phi=\Psi_h$ of Lemma \ref{c3:smooth} with $\rho:=g$ and $h:=\Phi$; $g$ has finite moments of all orders (Lemma \ref{c3:gauss}(f)). So Lemma \ref{c3:smooth}(b) gives smoothness, the equation, and the derivative formula. $|T_u\Phi|\le T_u|\Phi|\le\|\Phi\|_\infty$ since $g$ is a probability measure.
Bound: by Lemma \ref{c3:kbounds}(b) and Lemma \ref{c3:kernel}(a),
\[
|\partial_u^j\partial_{e_1}\cdots\partial_{e_k}T_u\Phi(y)|\le\|\Phi\|_\infty c_{j,k}\kappa^{k+2j}\int(1+|y|+|w|)^{k+2j}\gamma_{\sigma^2s_u}(w-a_uy)\lambda_{\mathcal E}(\rd w)
=\|\Phi\|_\infty c_{j,k}\kappa^{k+2j}\int(1+|y|+|a_uy+s_u^{1/2}z|)^{k+2j}g(\rd z),
\]
using that $\gamma_{\sigma^2s_u}(\cdot-a_uy)\lambda_{\mathcal E}$ is the image of $g$ under $z\mapsto a_uy+s_u^{1/2}z$ (Lemma \ref{c3:gauss}(c)). Since $a_u,s_u\le1$: $1+|y|+|a_uy+s_u^{1/2}z|\le1+2|y|+|z|\le(1+|y|)(2+|z|)$. The coordinate form follows as in Lemma \ref{lem:M2}(iii).
Limits: $T_u\Phi(y)=\int\Phi(a_uy+s_u^{1/2}z)g(\rd z)$. As $u\downarrow0$, $a_u\to1$ and $s_u\to0$, so $\Phi(a_uy+s_u^{1/2}z)\to\Phi(y)$ for every $z$ by continuity; as $u\to\infty$, $a_u\to0$, $s_u\to1$, so $\Phi(a_uy+s_u^{1/2}z)\to\Phi(z)$. The integrands are bounded by $\|\Phi\|_\infty$ and $g$ is a probability measure, so dominated convergence (along every sequence $u_n\downarrow0$, resp.\ $u_n\to\infty$) gives the limits.
\end{proof}

\subsection{Gaussian relative entropy}\label{c3:ss:KL}

For probability measures $\nu,\nu'$ on a measurable space, $S(\nu|\nu'):=\int\mathfrak h\big(\frac{\rd\nu}{\rd\nu'}\big)\rd\nu'$ if $\nu\ll\nu'$ and $S(\nu|\nu'):=+\infty$ otherwise, where $\mathfrak h(t)=t\log t$ ($\mathfrak h(0)=0$); this is the convention of Lemma \ref{lem:DV} and agrees with Mathlib's \texttt{klDiv} for probability measures.
With $\tilde{\mathfrak h}(t):=t\log t-t+1\ge0$ one has $S(\nu|\nu')=\int\tilde{\mathfrak h}(\frac{\rd\nu}{\rd\nu'})\rd\nu'\in[0,\infty]$ when $\nu\ll\nu'$ (the negative part of $\mathfrak h(f)$ is bounded by $1/e$, and $\int(f-1)\rd\nu'=0$).

\begin{auxlemma}[Gaussian relative entropy]\label{c3:KL}
\begin{enumerate}[label=(\alph*),itemsep=1pt]
\item For $m\in\mathcal E$ and $s>0$: $\ S\big(N(m,\sigma^2s)\,\big|\,g\big)=\frac d2\,(s-1-\log s)+\frac{|m|^2}{2\sigma^2}$.
\item For every probability measure $\rho$ on $\mathcal E$ with $M_2':=\int|w|^2\rho(\rd w)<\infty$ and every $u>0$:
\[
S(\rho_u|g)\le\int S\big(N(a_uw,\sigma^2s_u)\,\big|\,g\big)\rho(\rd w)=\frac d2\,(s_u-1-\log s_u)+\frac{e^{-u}M_2'}{2\sigma^2}\le\frac d2\log\frac1{1-e^{-u}}+\frac{e^{-u}M_2'}{2\sigma^2}.
\]
\end{enumerate}
\end{auxlemma}

\begin{proof}
(a) Let $\varkappa:=N(m,\sigma^2s)=\gamma_{\sigma^2s}(\cdot-m)\lambda_{\mathcal E}$ (Lemma \ref{c3:gauss}(c)). Then $\varkappa=f\,g$ with $f(y):=\gamma_{\sigma^2s}(y-m)/\gamma_{\sigma^2}(y)>0$ and
$\log f(y)=-\frac d2\log s+\frac{|y|^2}{2\sigma^2}-\frac{|y-m|^2}{2\sigma^2s}$, a polynomial of degree $2$ in $y$. If $Y\sim\varkappa$ then $Y=m+\sigma s^{1/2}Z$ in law with $Z\sim g_1$, $\E Z=0$, $\E|Z|^2=d$ (Lemma \ref{c3:gauss}(c),(f)); so $\log f\in L^1(\varkappa)$, $\E|Y|^2=|m|^2+\sigma^2sd$ and $\E|Y-m|^2=\sigma^2sd$. Hence
$S(\varkappa|g)=\int f\log f\,\rd g=\int\log f\,\rd\varkappa=-\frac d2\log s+\frac{|m|^2+\sigma^2sd}{2\sigma^2}-\frac d2$, which is the claim.

(b) Let $u>0$, $k_w:=K(u,\cdot,w)$. By Lemma \ref{c3:kernel}(a) and Lemma \ref{c3:gauss}(c), $k_w\,g=\gamma_{\sigma^2s_u}(\cdot-a_uw)\lambda_{\mathcal E}=N(a_uw,\sigma^2s_u)$, so by (a)
$\int\tilde{\mathfrak h}(k_w)\rd g=S(N(a_uw,\sigma^2s_u)|g)=\frac d2(s_u-1-\log s_u)+\frac{e^{-u}|w|^2}{2\sigma^2}$.
By Lemma \ref{lem:M1}(iv), $\rho_u=\Psi(u,\cdot)g$ with $\Psi(u,y)=\int k_w(y)\rho(\rd w)\in(0,\infty)$.
$\tilde{\mathfrak h}$ is convex and differentiable on $(0,\infty)$ with $\tilde{\mathfrak h}'(t_0)=\log t_0$, so $\tilde{\mathfrak h}(t)\ge\tilde{\mathfrak h}(t_0)+\log t_0\,(t-t_0)$ for $t_0>0$, $t\ge0$ (for $t=0$ it reads $1\ge1-t_0$).
Take $t_0=\Psi(u,y)$, $t=k_w(y)$ and integrate in $\rho(\rd w)$ (the right side is $\rho$-integrable with integral $\tilde{\mathfrak h}(t_0)$): $\int\tilde{\mathfrak h}(k_w(y))\rho(\rd w)\ge\tilde{\mathfrak h}(\Psi(u,y))$.
Integrate in $g(\rd y)$ and use Tonelli ($\tilde{\mathfrak h}\ge0$):
$S(\rho_u|g)=\int\tilde{\mathfrak h}(\Psi(u,\cdot))\rd g\le\int\big(\int\tilde{\mathfrak h}(k_w)\rd g\big)\rho(\rd w)$, which is the first inequality and the equality. The last inequality holds because $s_u-1\le0$ and $s_u=1-e^{-u}$.
\end{proof}

\begin{auxremark}[Use in Lemma \ref{lem:U8}]\label{c3:rem:U8}
With $\mathcal E=\cH$, $\sigma^2=2/(\beta N)$, $g=\bG$ and $\rho=\mathrm{law}(H_0)$, Lemma \ref{c3:KL}(b) gives $S(\mathrm{law}(H_s)|\bG)\le\frac d2\log\frac1{1-e^{-s}}+\frac{\beta N}{4}e^{-s}\,\E|H_0|^2$, and $\E|H_0|^2=\E\Tr H_0^2=\sum_{i,j}\sigma_{ij}^2=N$ under (A). This is the second inequality of Lemma \ref{lem:U8}; no kernel/chain-rule argument (\texttt{klDiv\_compProd\_eq\_add}) is needed. Only the first inequality of Lemma \ref{lem:U8} (data processing under $\lambda$, \texttt{klDiv\_map\_le}) uses Mathlib's relative entropy.
\end{auxremark}

\begin{remark}\label{rem:M4}
Lemmas \ref{lem:M1}--\ref{lem:M3} (with the auxiliary Lemmas \ref{c3:gauss}--\ref{c3:KL}) are the \emph{only} facts about ``Dyson Brownian motion'' used in NP. No stochastic process is constructed; $t\mapsto\mathrm{law}(H_t)$ is the curve $\rho_t$ with $\rho=\mathrm{law}(H_0)$, and $\rho_t$ is computed by Lemma \ref{lem:M1}(iv).
The constants in Lemma \ref{lem:M2}(ii)--(iii), Lemma \ref{lem:M3} and in the dominating functions \eqref{c3:dom} are explicit, but they depend on $d$, $\sigma$, $u_1$ and the moments of $\rho$ (hence on $N$ and on the ensemble): they are used only to justify differentiation under integrals and integrations by parts, never quantitatively. The only quantitative statements are Lemma \ref{lem:M2}(iv) (finiteness only is used) and Lemma \ref{c3:KL}.
\end{remark}

\section{Invariant functions and the eigenvalue density of $H_t$}\label{sec:C4}

\noindent\emph{Conventions of Sections \ref{sec:C4}--\ref{sec:C5}.}
We use the notation of \S\ref{ss:notation}: $\K\in\{\R,\C\}$, $\beta=\dim_\R\K$, the Euclidean space $(\cH,\la\cdot,\cdot\ra)$ of dimension $d$ with the orthonormal basis
\[\mathfrak B:=\{E_{ii}:1\le i\le N\}\cup\{X_{ij}:i<j\}\cup\{Y_{ij}:i<j\}\quad(\text{the }Y_{ij}\text{ only if }\beta=2),\]
the group $U_\beta(N)$, invariance, $\Sig$, $\lambda$, $\bG$ with $\sigma^2=2/(\beta N)$, the flow \eqref{Ht}, $\mu=\lambda_*\bG$ and $\nu_t=\mathrm{law}(\lambda(H_t))$.
From \S\ref{sec:C3} we use $T_u$, $A$, $a_u=e^{-u/2}$, $s_u=1-e^{-u}$, $\rho_u$ and Lemmas \ref{lem:M1}--\ref{lem:M3}, always with $\mathcal E=\cH$ and $g=\bG$; the coordinates on $\cH$ are those of the basis $\mathfrak B$, and the partial derivatives $\partial_y^\alpha$ in Lemmas \ref{lem:M2}, \ref{lem:M3} refer to them.

\begin{itemize}[leftmargin=1.6em,itemsep=1pt]
\item \emph{Derivatives.} For $\Phi\in C^k(\cH)$ (as a function on the finite-dimensional real space $\cH$), $D^j\Phi(H)$ ($j\le k$) is the $j$-th derivative, a $j$-linear form on $\cH$;
$\partial_Z\Phi(H):=D\Phi(H)[Z]$ and $\partial_Z^2\Phi(H):=D^2\Phi(H)[Z,Z]=\frac{\rd^2}{\rd t^2}\Phi(H+tZ)|_{t=0}$ for $Z\in\cH$; $\nabla\Phi(H)\in\cH$ is defined by $\la\nabla\Phi(H),Z\ra=D\Phi(H)[Z]$, and $\Delta_\cH\Phi(H):=\sum_{e\in\mathfrak B}D^2\Phi(H)[e,e]$.
The Hilbert--Schmidt norm of a $j$-linear form $T$ on $\cH$ is $|T|:=\big(\sum_{e_1,\dots,e_j\in\mathfrak B}T[e_1,\dots,e_j]^2\big)^{1/2}$; thus $|D\Phi(H)|=|\nabla\Phi(H)|$.
On $\R^N$ the same notation refers to the standard basis $(e_i)$; $|\nabla^2f|$ is the Frobenius norm of the Hessian, $\Delta=\sum_i\partial_i^2$, and $\partial^\alpha$ for a multi-index $\alpha\in\mathbb N^N$ is the usual partial derivative.
\item \emph{Independence of the basis.} $|T|$ does not change if $\mathfrak B$ is replaced by another orthonormal basis $\mathfrak B'$ of $\cH$, in any one slot and hence in all slots:
for fixed $h_2,\dots,h_j$ let $v\in\cH$ be defined by $\la v,h\ra=T[h,h_2,\dots,h_j]$; then $\sum_{e\in\mathfrak B}T[e,h_2,\dots,h_j]^2=\sum_{e\in\mathfrak B}\la v,e\ra^2=|v|^2=\sum_{e\in\mathfrak B'}T[e,h_2,\dots,h_j]^2$ (Parseval), and similarly in every other slot.
\item \emph{Conjugations.} For $U\in U_\beta(N)$ let $\mathcal O_U(H):=UHU^*$. Then $\mathcal O_U$ maps $\cH$ onto $\cH$ ($(UHU^*)^*=UHU^*$, and $UHU^*$ is real if $\beta=1$), is $\R$-linear with inverse $\mathcal O_{U^*}$, and
$\la\mathcal O_UX,\mathcal O_UY\ra=\Tr(UXU^*UYU^*)=\Tr(XY)$; so $\mathcal O_U$ is a linear isometry of $\cH$. $\Phi$ is invariant iff $\Phi\circ\mathcal O_U=\Phi$ for all $U\in U_\beta(N)$.
Permutation matrices, real diagonal matrices with entries $\pm1$, the rotations $R_\theta$ below, and (if $\beta=2$) diagonal matrices with entries of modulus $1$ belong to $U_\beta(N)$.
For a permutation $\pi$ of $\{1,\dots,N\}$ let $P_\pi e_k:=e_{\pi(k)}$; then $P_\pi\diag(x)P_\pi^{T}=\diag(x\circ\pi^{-1})$, so every rearrangement of $x$ is obtained in this way.
\item \emph{The cone.} $\cSig$, the closure of $\Sig$, equals $\{x\in\R^N:x_1\le\dots\le x_N\}$: this set is closed and contains $\Sig$, and each of its points $x$ is the limit of $x+\eps(1,2,\dots,N)\in\Sig$ as $\eps\downarrow0$.
We write $\partial\Sig:=\cSig\setminus\Sig=\{x\in\cSig:\ x_k=x_{k+1}\text{ for some }k\}$ (the \emph{walls}). Measures on $\cSig$ (such as $\mu$, $\nu_t$) are identified with Borel measures on $\R^N$ carried by $\cSig$.
\end{itemize}

\subsection{Radial calculus by Givens rotations}

\begin{lemma}[Invariant functions]\label{lem:inv}
Let $\Phi:\cH\to\R$ be invariant and $\varphi(x):=\Phi(\diag x)$, $x\in\R^N$. Then $\varphi$ is symmetric, $\Phi=\varphi\circ\lambda$, and $\Phi\in C^k(\cH)\Rightarrow\varphi\in C^k(\R^N)$.
Conversely $\theta\circ\lambda$ is invariant for every $\theta:\cSig\to\R$. For every $u\ge0$, $T_u$ maps invariant (nonnegative or bounded) Borel functions to invariant functions.
Moreover:
\begin{enumerate}[label=\textup{(\alph*)},leftmargin=2em,itemsep=1pt]
\item $\lambda(UHU^*)=\lambda(H)$ for $U\in U_\beta(N)$; for every $H\in\cH$ there is $U\in U_\beta(N)$ with $H=U\diag(\lambda(H))U^*$; $\lambda(\diag x)=x$ for $x\in\cSig$; $|\lambda(H)|=|H|$; and $\max_k|\lambda_k(H)-\lambda_k(H')|\le|H-H'|$ for $H,H'\in\cH$. In particular $\lambda:\cH\to\cSig$ is continuous, hence Borel.
\item If $\Phi\in C^k(\cH)$ is invariant and $0\le j\le k$, then $H\mapsto|D^j\Phi(H)|$ is invariant, and for $x\in\R^N$ and every multi-index $\alpha$ with $|\alpha|=j$: $|\partial^\alpha\varphi(x)|\le|D^j\Phi(\diag x)|$. Moreover $|\nabla\varphi(x)|\le|\nabla\Phi(\diag x)|$ and $|\nabla^2\varphi(x)|\le|D^2\Phi(\diag x)|$.
\end{enumerate}
\end{lemma}

\begin{proof}
\emph{(a)} For $U\in U_\beta(N)$ and $t\in\R$, $\det(t-UHU^*)=\det(U)\det(t-H)\det(U^*)=\det(t-H)$ since $\det U\,\det U^*=|\det U|^2=1$. As $\lambda(H)$ is the nondecreasing list of the roots of $\det(t-H)$ counted with multiplicity (Part~I, before Lemma \ref{lem:weyl}), $\lambda(UHU^*)=\lambda(H)$.
By the spectral theorem (for real symmetric matrices if $\beta=1$, for Hermitian matrices if $\beta=2$) there are $V\in U_\beta(N)$ and $y\in\R^N$ with $H=V\diag(y)V^*$; by the remark just quoted, $y$ is a rearrangement of $\lambda(H)$, so $\diag(y)=P\diag(\lambda(H))P^T$ for a permutation matrix $P$, and $U:=VP\in U_\beta(N)$ does the job.
The characteristic polynomial of $\diag x$ is $\prod_k(t-x_k)$, so $\lambda(\diag x)$ is the nondecreasing rearrangement of $x$, which is $x$ itself if $x\in\cSig$.
With $H=U\diag(\lambda(H))U^*$: $|H|^2=\Tr H^2=\Tr\big(U\diag(\lambda(H))^2U^*\big)=\sum_k\lambda_k(H)^2=|\lambda(H)|^2$.
Finally, by Lemma \ref{lem:weyl}(b),(a) of Part~I, $|\lambda_k(H)-\lambda_k(H')|\le\|H-H'\|=\max_l|\lambda_l(H-H')|\le|\lambda(H-H')|=|H-H'|$.

\emph{Symmetry of $\varphi$; $\Phi=\varphi\circ\lambda$.} For a permutation $\pi$, $\varphi(x\circ\pi^{-1})=\Phi(P_\pi\diag(x)P_\pi^T)=\Phi(\diag x)=\varphi(x)$ by invariance; every rearrangement of $x$ is of the form $x\circ\pi^{-1}$.
For $H\in\cH$ choose $U$ as in (a): $\Phi(H)=\Phi(U\diag(\lambda(H))U^*)=\Phi(\diag\lambda(H))=\varphi(\lambda(H))$.

\emph{Regularity.} $\iota:\R^N\to\cH$, $x\mapsto\diag x$, is linear, so $\varphi=\Phi\circ\iota\in C^k$ if $\Phi\in C^k$, with
$D^j\varphi(x)[v_1,\dots,v_j]=D^j\Phi(\diag x)[\diag v_1,\dots,\diag v_j]$. In particular, if $\alpha$ lists the indices $i_1,\dots,i_j$, then $\partial^\alpha\varphi(x)=D^j\Phi(\diag x)[E_{i_1i_1},\dots,E_{i_ji_j}]$. \hfill$(\ast)$

\emph{Converse.} $\theta(\lambda(UHU^*))=\theta(\lambda(H))$ by (a).

\emph{$T_u$.} $T_0$ is the identity ($a_0=1$, $s_0=0$). Let $u>0$ and let $\Phi\ge0$ be Borel and invariant. By Lemma \ref{lem:M1}(iii) with the linear isometry $\mathcal O=\mathcal O_U$:
$T_u\Phi=T_u(\Phi\circ\mathcal O_U)=(T_u\Phi)\circ\mathcal O_U$; so $T_u\Phi$ is invariant (and Borel, by Tonelli's theorem in \eqref{Mehler}). If $\Phi$ is bounded and invariant, so are $\Phi^\pm:=\max(\pm\Phi,0)$; $T_u\Phi^\pm\le\|\Phi\|_\infty$ is finite and $T_u\Phi=T_u\Phi^+-T_u\Phi^-$ is invariant.

\emph{(b)} Since $\mathcal O_U$ is linear, differentiating $\Phi\circ\mathcal O_U=\Phi$ $j$ times gives $D^j\Phi(\mathcal O_UH)[\mathcal O_Uh_1,\dots,\mathcal O_Uh_j]=D^j\Phi(H)[h_1,\dots,h_j]$.
As $\mathcal O_U\mathfrak B$ is an orthonormal basis of $\cH$ and $|\cdot|$ does not depend on the basis,
\[|D^j\Phi(\mathcal O_UH)|^2=\sum_{e_1,\dots,e_j\in\mathfrak B}D^j\Phi(\mathcal O_UH)[\mathcal O_Ue_1,\dots,\mathcal O_Ue_j]^2=\sum_{e_1,\dots,e_j\in\mathfrak B}D^j\Phi(H)[e_1,\dots,e_j]^2=|D^j\Phi(H)|^2 .\]
By $(\ast)$, $\partial^\alpha\varphi(x)$ is one of the entries $D^j\Phi(\diag x)[e_1,\dots,e_j]$, $e_l\in\mathfrak B$, so $|\partial^\alpha\varphi(x)|\le|D^j\Phi(\diag x)|$.
Likewise $|\nabla\varphi(x)|^2=\sum_i D\Phi(\diag x)[E_{ii}]^2\le\sum_{e\in\mathfrak B}D\Phi(\diag x)[e]^2=|\nabla\Phi(\diag x)|^2$ and $|\nabla^2\varphi(x)|^2=\sum_{i,l}D^2\Phi(\diag x)[E_{ii},E_{ll}]^2\le|D^2\Phi(\diag x)|^2$.
\end{proof}

\begin{lemma}[Radial part of the gradient and Laplacian]\label{lem:givens}
Let $\Phi\in C^2(\cH)$ be invariant, $\varphi=\Phi\circ\diag$, and let $x\in\R^N$ have pairwise distinct coordinates, $D:=\diag x$. Then for all $i<j$ and $Z\in\{X_{ij},Y_{ij}\}$ ($Y_{ij}$ only if $\beta=2$):
\begin{equation}\label{givens}
\partial_Z\Phi(D)=0,\qquad \partial_Z^2\Phi(D)=\frac{\partial_i\varphi(x)-\partial_j\varphi(x)}{x_i-x_j}.
\end{equation}
Consequently $|\nabla\Phi(D)|^2=|\nabla\varphi(x)|^2$, $\la D,\nabla\Phi(D)\ra=x\cdot\nabla\varphi(x)$,
$\Delta_\cH\Phi(D)=\Delta\varphi(x)+\beta\sum_{i<j}\frac{\partial_i\varphi-\partial_j\varphi}{x_i-x_j}(x)$, and, with $\sigma^2=2/(\beta N)$,
\begin{equation}\label{L}
(A\Phi)(\diag x)=(L\varphi)(x),\qquad L:=\frac1{\beta N}\Delta+\sum_{i=1}^N\Big(-\frac{x_i}{2}+\frac1N\sum_{j\ne i}\frac{1}{x_i-x_j}\Big)\partial_i .
\end{equation}
Moreover, for \emph{every} $x\in\R^N$ and every invariant $\Phi\in C^1(\cH)$: $\partial_Z\Phi(\diag x)=0$ for all off-diagonal $Z\in\mathfrak B$, $\nabla\Phi(\diag x)=\sum_i\partial_i\varphi(x)E_{ii}$, $|\nabla\Phi(\diag x)|=|\nabla\varphi(x)|$ and $\la\diag x,\nabla\Phi(\diag x)\ra=x\cdot\nabla\varphi(x)$.
\end{lemma}

\begin{proof}
Let $x\in\R^N$, $D=\diag x$. By $(\ast)$ in the proof of Lemma \ref{lem:inv}, $D\Phi(D)[E_{ii}]=\partial_i\varphi(x)$, and $D\Phi(D)[D]=\sum_ix_i\,\partial_i\varphi(x)$ since $D=\sum_ix_iE_{ii}$.

\emph{Step 1 (first derivatives; all $x$, $\Phi\in C^1$).} Fix $i<j$ and let $S:=\diag(s)$ with $s_j=-1$ and $s_k=1$ for $k\ne j$; $S\in U_\beta(N)$. Conjugation by $S$ multiplies the $(k,l)$ entry by $s_ks_l$; hence $SDS^*=D$, $SX_{ij}S^*=-X_{ij}$, $SY_{ij}S^*=-Y_{ij}$.
For $Z\in\{X_{ij},Y_{ij}\}$ and $t\in\R$, invariance gives $\Phi(D+tZ)=\Phi(\mathcal O_S(D+tZ))=\Phi(D-tZ)$. The function $t\mapsto\Phi(D+tZ)$ is even and differentiable, so its derivative at $0$, $\partial_Z\Phi(D)$, vanishes.
Therefore $\nabla\Phi(D)=\sum_{e\in\mathfrak B}D\Phi(D)[e]\,e=\sum_i\partial_i\varphi(x)E_{ii}$, whence $|\nabla\Phi(D)|^2=\sum_i\partial_i\varphi(x)^2=|\nabla\varphi(x)|^2$ and $\la D,\nabla\Phi(D)\ra=D\Phi(D)[D]=x\cdot\nabla\varphi(x)$.
This proves the first identity in \eqref{givens} and the additional statement.

\emph{Step 2 (second derivative in the direction $X_{ij}$; $x_i\ne x_j$).} Fix $i<j$, put $\varsigma:=x_j-x_i\ne0$, and for $\theta\in\R$ let $R_\theta$ be the real orthogonal matrix with $R_\theta e_i=\cos\theta\,e_i+\sin\theta\,e_j$, $R_\theta e_j=-\sin\theta\,e_i+\cos\theta\,e_j$, $R_\theta e_k=e_k$ ($k\ne i,j$). Writing $c=\cos\theta$, $s=\sin\theta$ and $D=\sum_kx_ke_ke_k^T$,
\[R_\theta DR_\theta^T=\sum_kx_k(R_\theta e_k)(R_\theta e_k)^T=\sum_{k\ne i,j}x_kE_{kk}+x_i\big(c^2E_{ii}+cs(E_{ij}+E_{ji})+s^2E_{jj}\big)+x_j\big(s^2E_{ii}-cs(E_{ij}+E_{ji})+c^2E_{jj}\big),\]
that is,
\begin{equation}\label{c4:rot}
\mathbf c(\theta):=R_\theta DR_\theta^T=D+\varsigma\sin^2\theta\,(E_{ii}-E_{jj})-\sqrt2\,\varsigma\sin\theta\cos\theta\,X_{ij}.
\end{equation}
$\mathbf c:\R\to\cH$ is $C^\infty$ with $\mathbf c(0)=D$,
$\mathbf c'(\theta)=\varsigma\sin2\theta\,(E_{ii}-E_{jj})-\sqrt2\,\varsigma\cos2\theta\,X_{ij}$ and $\mathbf c''(\theta)=2\varsigma\cos2\theta\,(E_{ii}-E_{jj})+2\sqrt2\,\varsigma\sin2\theta\,X_{ij}$; so
$\mathbf c'(0)=-\sqrt2\,\varsigma X_{ij}$ and $\mathbf c''(0)=2\varsigma(E_{ii}-E_{jj})$.
By invariance $g(\theta):=\Phi(\mathbf c(\theta))=\Phi(D)$ for all $\theta$. Since $\Phi\in C^2$, the chain rule gives $g''(\theta)=D^2\Phi(\mathbf c(\theta))[\mathbf c'(\theta),\mathbf c'(\theta)]+D\Phi(\mathbf c(\theta))[\mathbf c''(\theta)]$, and $g''\equiv0$. At $\theta=0$:
\[0=2\varsigma^2\,\partial^2_{X_{ij}}\Phi(D)+2\varsigma\big(\partial_i\varphi(x)-\partial_j\varphi(x)\big).\]
Dividing by $2\varsigma^2$: $\partial^2_{X_{ij}}\Phi(D)=-(\partial_i\varphi(x)-\partial_j\varphi(x))/(x_j-x_i)$, which is the second identity in \eqref{givens} for $Z=X_{ij}$.

\emph{Step 3 (the direction $Y_{ij}$, $\beta=2$).} Let $P:=\diag(p)$ with $p_j=\mathrm i$ and $p_k=1$ for $k\ne j$; $P\in U_2(N)$. Conjugation by $P$ multiplies the $(k,l)$ entry by $p_k\overline{p_l}$, so $PDP^*=D$ and $PX_{ij}P^*$ has $(i,j)$ entry $-\mathrm i/\sqrt2$ and $(j,i)$ entry $\mathrm i/\sqrt2$, i.e.\ $PX_{ij}P^*=-Y_{ij}$.
Hence $\Phi(D+tY_{ij})=\Phi(\mathcal O_P(D-tX_{ij}))=\Phi(D-tX_{ij})$ for all $t$, and taking the second derivative at $t=0$: $\partial^2_{Y_{ij}}\Phi(D)=\partial^2_{X_{ij}}\Phi(D)$. With Step 2 this proves \eqref{givens} for $Z=Y_{ij}$.

\emph{Step 4 (Laplacian and generator).} By $(\ast)$, $D^2\Phi(D)[E_{ii},E_{ii}]=\partial_i^2\varphi(x)$. Summing over $\mathfrak B$, which contains $\beta$ off-diagonal elements for each pair $i<j$, and using \eqref{givens},
\[\Delta_\cH\Phi(D)=\sum_i\partial_i^2\varphi(x)+\beta\sum_{i<j}\frac{\partial_i\varphi(x)-\partial_j\varphi(x)}{x_i-x_j}.\]
In the last sum, the pair $\{i,j\}$ contributes $\partial_i\varphi/(x_i-x_j)+\partial_j\varphi/(x_j-x_i)$, so the sum equals $\sum_i\partial_i\varphi(x)\sum_{j\ne i}(x_i-x_j)^{-1}$.
By \eqref{Mehler}, $A\Phi(D)=\frac{\sigma^2}2\Delta_\cH\Phi(D)-\frac12\la D,\nabla\Phi(D)\ra$ with $\frac{\sigma^2}{2}=\frac1{\beta N}$; inserting the two formulas above and Step 1 gives \eqref{L}.
\end{proof}

\begin{remark}[The generator; PD-06]\label{rem:gen}
For $\beta=2$, $L$ is exactly (2.32) of the original, $\mathcal L=\sum_i\frac1{2N}\partial_i^2+\sum_i\big(-\frac\beta4x_i+\frac\beta{2N}\sum_{j\ne i}\frac1{x_i-x_j}\big)\partial_i$ with $\beta=2$. For $\beta=1$, $L=2\,\mathcal L_{\beta=1}$: the GOE-normalised flow \eqref{Ht} (diagonal variance $2/N$) runs the $\beta=1$ Dyson generator at double speed. This is harmless: all statements are for $t\ge$ threshold.
Formula \eqref{L} is the content of ``It\^o's formula for the eigenvalues''; here it is an identity between two continuous functions on the open dense set of $x$ with pairwise distinct coordinates, and the left side $x\mapsto(A\Phi)(\diag x)$ is continuous on all of $\R^N$ because $\Phi\in C^2(\cH)$.
In particular, for $\varphi=\Phi\circ\diag$ with $\Phi\in C^2(\cH)$ invariant, the singular terms $\frac{\partial_i\varphi-\partial_j\varphi}{x_i-x_j}$ extend continuously across the walls $\{x_i=x_j\}$. \emph{This is how NP handles collisions: nothing happens there.}
\end{remark}

\subsection{The eigenvalue density at time $t$}

In this subsection $H_0$ is a random element of $\cH$, defined on a probability space that also carries $V\sim\bG$ independent of $H_0$, and $H_t$ is given by \eqref{Ht}. Let $\rho:=\mathrm{law}(H_0)$.
Since $e^{-t/2}=a_t$ and $(1-e^{-t})^{1/2}=s_t^{1/2}$, \eqref{Ht} says that $\mathrm{law}(H_t)=\rho_t$ in the notation of \S\ref{sec:C3}.
``Finite moments of all orders'' means $\E|H_0|^m<\infty$ for every $m\in\mathbb N$; then also $\E|H_0|^m<\infty$ for every real $m\ge0$ (as $|H_0|^m\le1+|H_0|^{\lceil m\rceil}$). For the matrices of the ensembles of Theorem \ref{thm:23} this holds by (RA6).

\begin{proposition}[Smooth symmetric version of $\rd\nu_t/\rd\mu$]\label{prop:U9}
Let $H_0$ be a random element of $\cH$ with $\E|H_0|^m<\infty$ for every $m\in\mathbb N$, and let $s>0$. Then:
\begin{enumerate}[label=(\roman*),itemsep=1pt]
\item $\nu_s\ll\mu$. Fix a Borel version $f_s:=\rd\nu_s/\rd\mu$ on $\cSig$ and put $\Phi_s:=f_s\circ\lambda$; then $\Phi_s\bG$ is a probability measure on $\cH$ with $\int|H|^m\Phi_s\rd\bG=\E|H_s|^m<\infty$.
\item For $t>s$ let $\Phi_t:=T_{t-s}\Phi_s$ and $f_t(x):=\Phi_t(\diag x)$, $x\in\R^N$. Then $(t,H)\mapsto\Phi_t(H)$ is $C^\infty$ and positive on $(s,\infty)\times\cH$, $\Phi_t$ is invariant, $f_t$ is $C^\infty$, positive and symmetric on $\R^N$, and
$\nu_t=f_t\,\mu$ on $\cSig$.
\item $\partial_tf_t(x)=(A\Phi_t)(\diag x)$ for all $x\in\R^N$, and $\partial_tf_t(x)=(Lf_t)(x)$ for all $x$ with pairwise distinct coordinates.
\item For $s<t_1<t_2<\infty$ there is $C<\infty$ with $|\log f_t(x)|\le C(1+|x|^2)$ on $[t_1,t_2]\times\R^N$, and
$\mathfrak m(x):=\sup_{t\in[t_1,t_2]}\big(|f_t|+|\nabla f_t|+|\nabla^2f_t|+|\partial_tf_t|+|\nabla\partial_tf_t|\big)(x)$ satisfies $\int_{\cSig}\mathfrak m(x)(1+|x|)^m\mu(\rd x)<\infty$ for all $m$. Moreover $\int_{\cSig}|\nabla f_t|^2/f_t\,\rd\mu<\infty$.
\end{enumerate}
\end{proposition}

\begin{proof}
\emph{(i)} By Lemma \ref{lem:M1}(iv), $\rho_s$ has a Lebesgue density $p(s,\cdot)$; $\bG$ has the everywhere positive Lebesgue density $\gamma_{\sigma^2}$. Hence $\rho_s=(p(s,\cdot)/\gamma_{\sigma^2})\,\bG\ll\bG$.
If $B\subset\cSig$ is Borel with $\mu(B)=0$, then $\bG(\lambda^{-1}(B))=0$ ($\lambda$ is Borel by Lemma \ref{lem:inv}(a)), hence $\nu_s(B)=\rho_s(\lambda^{-1}(B))=0$. So $\nu_s\ll\mu$, and by the Radon--Nikodym theorem there is a Borel $f_s:\cSig\to[0,\infty)$ with $\nu_s=f_s\mu$; we fix one.
$\Phi_s=f_s\circ\lambda$ is Borel, nonnegative and invariant (Lemma \ref{lem:inv}). For Borel $F:\cSig\to[0,\infty]$, the change of variables formula for the image measure $\mu=\lambda_*\bG$ gives
\begin{equation}\label{c4:push}
\int_\cH F(\lambda(H))\,\Phi_s(H)\,\bG(\rd H)=\int_{\cSig}F\,f_s\,\rd\mu=\int_{\cSig}F\,\rd\nu_s=\E\,F(\lambda(H_s)).
\end{equation}
With $F=1$: $\int\Phi_s\rd\bG=1$. With $F(x)=|x|^m$ and $|\lambda(H)|=|H|$ (Lemma \ref{lem:inv}(a)): $\int|H|^m\Phi_s\rd\bG=\E|H_s|^m$.
Finally, for $m\ge1$, $|H_s|^m\le2^{m-1}(|H_0|^m+|V|^m)$ (as $a_s,s_s^{1/2}\le1$), and $\E|V|^m=\int|H|^m\bG(\rd H)<\infty$ (Gaussian moments), so $\E|H_s|^m<\infty$; for $m=0$ there is nothing to prove.

\emph{(ii), smoothness.} By (i), $\rho':=\Phi_s\bG$ is a probability measure on $\cH$ with finite moments of all orders. Apply Lemma \ref{lem:M2} to $\rho'$: the function $\Psi'(u,y):=p'(u,y)/\gamma_{\sigma^2}(y)$, where $p'(u,\cdot)$ is the Lebesgue density of $\rho'_u$, is $C^\infty$ and positive on $(0,\infty)\times\cH$ and satisfies $\partial_u\Psi'=A\Psi'$.
Moreover $\Psi'(u,\cdot)=T_u\Phi_s$ \emph{everywhere} on $\cH$ for $u>0$. (This is the parenthetical statement of Lemma \ref{lem:M2}; for completeness: by Lemma \ref{lem:M1}(iv), $p'(u,y)=\int\gamma_{\sigma^2s_u}(y-a_uw)\Phi_s(w)\gamma_{\sigma^2}(w)\,\rd w$, and
\[\gamma_{\sigma^2s_u}(y-a_uw)\,\gamma_{\sigma^2}(w)=\gamma_{\sigma^2}(y)\,\gamma_{\sigma^2s_u}(w-a_uy)\qquad(y,w\in\cH),\]
because the normalising constants on both sides are $(2\pi\sigma^2)^{-d}s_u^{-d/2}$ and $s_u(|w|^2-|y|^2)+|y-a_uw|^2=|w-a_uy|^2$ by $a_u^2+s_u=1$; dividing by $\gamma_{\sigma^2}(y)$ and comparing with the second formula in \eqref{Mehler} gives $\Psi'(u,y)=(T_u\Phi_s)(y)$.)
Hence $\Phi_t(H)=\Psi'(t-s,H)$, and $(t,H)\mapsto\Phi_t(H)$ is $C^\infty$ and positive on $(s,\infty)\times\cH$ with
\begin{equation}\label{c4:PhiPDE}
\partial_t\Phi_t=A\Phi_t\qquad\text{on }(s,\infty)\times\cH .
\end{equation}
$\Phi_t=T_{t-s}\Phi_s$ is invariant by Lemma \ref{lem:inv}, since $\Phi_s\ge0$ is Borel and invariant. Consequently $f_t=\Phi_t\circ\diag$ is symmetric (Lemma \ref{lem:inv}) and positive, and $(t,x)\mapsto f_t(x)=\Phi_t(\diag x)$ is $C^\infty$ on $(s,\infty)\times\R^N$ as the composition of $(t,H)\mapsto\Phi_t(H)$ with the linear map $(t,x)\mapsto(t,\diag x)$.

\emph{(ii), identification.} Let $t>s$, $B\subset\cSig$ Borel, and $\Theta:=\mathbf 1_B\circ\lambda$, a Borel invariant function with values in $\{0,1\}$. Then:
\begin{itemize}[leftmargin=1.6em,itemsep=1pt]
\item[(a)] $\nu_t(B)=\int\Theta\,\rd\rho_t=\int\Theta\,\rd(\rho_s)_{t-s}=\int T_{t-s}\Theta\,\rd\rho_s$, by $\mathrm{law}(H_t)=\rho_t$, Lemma \ref{lem:M1}(i) ($(\rho_s)_{t-s}=\rho_t$) and Lemma \ref{lem:M1}(v).
\item[(b)] $T_{t-s}\Theta$ is Borel, invariant (Lemma \ref{lem:inv}) and takes values in $[0,1]$ (Lemma \ref{lem:M1}(i): $T_u$ is positive and $T_u1=1$). By Lemma \ref{lem:inv}, $T_{t-s}\Theta=\vartheta\circ\lambda$ with the Borel function $\vartheta:=(T_{t-s}\Theta)\circ\diag$.
By the image-measure formula for $\nu_s=\lambda_*\rho_s$ and by \eqref{c4:push} with $F=\vartheta$:
$\int T_{t-s}\Theta\,\rd\rho_s=\int_{\cSig}\vartheta\,\rd\nu_s=\int_\cH(T_{t-s}\Theta)\,\Phi_s\,\rd\bG$.
\item[(c)] By the symmetry of $T_{t-s}$, Lemma \ref{lem:M1}(ii) (both $\Theta$ and $\Phi_s$ are Borel and nonnegative): $\int(T_{t-s}\Theta)\,\Phi_s\,\rd\bG=\int\Theta\,T_{t-s}\Phi_s\,\rd\bG=\int\Theta\,\Phi_t\,\rd\bG$.
\item[(d)] $\Phi_t$ is invariant, so $\Phi_t=f_t\circ\lambda$ (Lemma \ref{lem:inv}), and $\int\Theta\,\Phi_t\,\rd\bG=\int(\mathbf 1_Bf_t)\circ\lambda\,\rd\bG=\int_Bf_t\,\rd\mu$.
\end{itemize}
Thus $\nu_t(B)=\int_Bf_t\,\rd\mu$ for every Borel $B\subset\cSig$, i.e.\ $\nu_t=f_t\mu$. Note that the law $\rho_s$ of $H_s$ itself is never claimed to be invariant; only its eigenvalue law $\nu_s$ enters, through \eqref{c4:push}.

\emph{(iii)} By \eqref{c4:PhiPDE}, $\partial_tf_t(x)=(\partial_t\Phi_t)(\diag x)=(A\Phi_t)(\diag x)$ for all $x\in\R^N$. If $x$ has pairwise distinct coordinates, Lemma \ref{lem:givens} applied to the invariant $C^2$ function $\Phi_t$ gives $(A\Phi_t)(\diag x)=(Lf_t)(x)$.

\emph{(iv)} Fix $s<t_1<t_2<\infty$ and put $u_l:=t_l-s$, so $0<u_1<u_2$ and $\Phi_t=\Psi'(t-s,\cdot)$ with $t-s\in[u_1,u_2]$ for $t\in[t_1,t_2]$.

\emph{Logarithm.} Lemma \ref{lem:M2}(ii) for $\rho'$ gives $C$ with $|\log\Psi'(u,y)|\le C(1+|y|^2)$ for $u\in[u_1,u_2]$, $y\in\cH$. With $y=\diag x$ and $|\diag x|=|x|$: $|\log f_t(x)|\le C(1+|x|^2)$ for $t\in[t_1,t_2]$, $x\in\R^N$.

\emph{Derivatives.} For $j\in\{0,1\}$, $H\mapsto\partial_t^j\Phi_t(H)$ is $C^\infty$ and invariant (for $j=1$ it is the pointwise limit of the invariant difference quotients $(\Phi_{t+h}-\Phi_t)/h$), and its restriction to diagonal matrices is $\partial_t^jf_t$. Define
\[\widetilde{\mathfrak M}(H):=\sup_{t\in[t_1,t_2]}\ \sum_{j=0}^1\sum_{k=0}^2\big|D^k\partial_t^j\Phi_t(H)\big|\qquad(H\in\cH).\]
Each summand is invariant by Lemma \ref{lem:inv}(b) (for $k=0$, $|\partial_t^j\Phi_t|$ is trivially invariant), so $\widetilde{\mathfrak M}$ is invariant.
It is finite and Borel: the summands are continuous in $(t,H)$, so the supremum over the compact interval $[t_1,t_2]$ is finite and equals the supremum over the countable set $([t_1,t_2]\cap\mathbb Q)\cup\{t_1,t_2\}$. The same argument shows that $\mathfrak m$ is finite and Borel.
By Lemma \ref{lem:inv}(b) applied to $\partial_t^j\Phi_t$ ($j=0,1$),
\[|f_t(x)|+|\nabla f_t(x)|+|\nabla^2f_t(x)|+|\partial_tf_t(x)|+|\nabla\partial_tf_t(x)|\le\sum_{j=0}^1\sum_{k=0}^2|D^k\partial_t^j\Phi_t(\diag x)|,\]
hence $\mathfrak m(x)\le\widetilde{\mathfrak M}(\diag x)$ for all $x\in\R^N$.
On the other hand, for a $k$-linear form $T$ on $\cH$, $|T|\le\sum_{e_1,\dots,e_k\in\mathfrak B}|T[e_1,\dots,e_k]|$, and each ordered tuple $(e_1,\dots,e_k)$ corresponds to a multi-index $\alpha$ with $|\alpha|=k$ (at most $k!\le2$ tuples per $\alpha$ for $k\le2$) with $D^k\partial_t^j\Phi_t[e_1,\dots,e_k]=\partial_y^\alpha\partial_t^j\Psi'(t-s,\cdot)$. Hence $\widetilde{\mathfrak M}\le2\,\mathfrak M$, where $\mathfrak M$ is the function of Lemma \ref{lem:M2}(iii) for $\rho'$ and $[u_1,u_2]$.
Since $\widetilde{\mathfrak M}$ is invariant, $\widetilde{\mathfrak M}=\widetilde{\mathfrak M}\circ\diag\circ\lambda$ (Lemma \ref{lem:inv}); with $|\lambda(H)|=|H|$ and $\mu=\lambda_*\bG$,
\[\int_{\cSig}\mathfrak m(x)(1+|x|)^m\mu(\rd x)\le\int_{\cSig}\widetilde{\mathfrak M}(\diag x)(1+|x|)^m\mu(\rd x)=\int_\cH\widetilde{\mathfrak M}(H)(1+|H|)^m\,\bG(\rd H)\le2\int_\cH\mathfrak M(1+|H|)^m\,\rd\bG<\infty\]
by Lemma \ref{lem:M2}(iii).

\emph{Fisher information.} Let $t>s$. By Lemma \ref{lem:inv}(b), $|\nabla\Phi_t|$ is invariant, hence so is $|\nabla\Phi_t|^2/\Phi_t$, and $|\nabla f_t(x)|\le|\nabla\Phi_t(\diag x)|$ (with equality, by Lemma \ref{lem:givens}). Therefore, as above,
\[\int_{\cSig}\frac{|\nabla f_t|^2}{f_t}\rd\mu\le\int_{\cSig}\frac{|\nabla\Phi_t|^2}{\Phi_t}(\diag x)\,\mu(\rd x)=\int_\cH\frac{|\nabla\Psi'(t-s,\cdot)|^2}{\Psi'(t-s,\cdot)}\rd\bG<\infty\]
by Lemma \ref{lem:M2}(iv) with $u=t-s$; its right side is finite because $\rho'_{u}=\mathrm{law}(a_uY+s_u^{1/2}Z)$ with $Y\sim\rho'$, $Z\sim\bG$, and both have finite second moments.
\end{proof}

\begin{remark}[What is \emph{not} needed]
No uniqueness theorem for $\partial_tf=Lf$, no semigroup on $\Sig$, no domain questions: $f_t$ is \emph{defined} by the matrix flow, and the equation is used only as a pointwise identity, once, in Proposition \ref{prop:entprod} (and for bounded data in Theorem \ref{thm:JD}).
Neither Haar measure on $U_\beta(N)$ nor the invariance of $\mathrm{law}(H_s)$ is used: in the proof of Proposition \ref{prop:U9}(ii) only $\mu=\lambda_*\bG$, the symmetry of $T_u$ in $L^2(\bG)$ and the commutation of $T_u$ with conjugations enter. In particular $\nu_s\ll\mu$ for every $s>0$ also when $H_0$ has a discrete (e.g.\ Bernoulli) law.
The version $f_t$ does not depend on the choice of the version $f_s$: two versions of $f_s$ differ on a $\mu$-null set, so the corresponding $\Phi_s$ differ on a $\bG$-null set, which is Lebesgue-null, and $T_u$ ($u>0$) integrates against a Lebesgue density (second formula in \eqref{Mehler}).
\end{remark}
\section{\INFRA\ Integration by parts on $\Sig$ and the joint eigenvalue density of G$\beta$E}\label{sec:C5}

Let $\Delta(x):=\prod_{i<j}(x_j-x_i)$ (a polynomial; positive on $\Sig$), and
\begin{equation}\label{mustar}
w(x):=\Delta(x)^\beta e^{-\frac{\beta N}{4}|x|^2},\qquad Z_\beta:=\int_{\Sig}w\,\rd x\in(0,\infty),\qquad \mu^*(\rd x):=Z_\beta^{-1}w(x)\mathbf 1_{\Sig}(x)\,\rd x .
\end{equation}
(The value of $Z_\beta$ is never needed.) Here $\beta\in\{1,2\}$, so $w$ is a polynomial times a Gaussian and $w\in C^\infty(\R^N)$; for $N=1$, $\Delta\equiv1$ and $\Sig=\R$.
For $x\in\R^N$ and $i<j$, $|x_j-x_i|\le|x_i|+|x_j|\le2|x|$, hence $0\le|w(x)|\le(2|x|)^{\beta N(N-1)/2}e^{-\frac{\beta N}4|x|^2}$. Since for $p\ge0$, $c>0$ the function $r\mapsto(1+r)^pe^{-cr^2/2}$ is bounded on $[0,\infty)$ and $\int_{\R^N}e^{-c|x|^2/2}\rd x=(2\pi/c)^{N/2}$,
\begin{equation}\label{c5:mom}
\int_{\R^N}(1+|x|)^m\,|w(x)|\,\rd x<\infty\qquad(m\ge0).
\end{equation}
With $m=0$ this gives $Z_\beta<\infty$; $Z_\beta>0$ because $w>0$ on the nonempty open set $\Sig$. So $\mu^*$ is a Borel probability measure on $\R^N$, carried by $\Sig$, with finite moments of all orders.
Recall $\partial\Sig=\cSig\setminus\Sig=\{x\in\cSig:x_k=x_{k+1}\text{ for some }k\}$ (\S\ref{sec:C4}); $w=0$ on $\partial\Sig$ because $\Delta=0$ there and $\beta\ge1$.

\smallskip\noindent\emph{The divergence identity.} On $\Sig$, for $\varphi\in C^2$: $\ \beta N\,w\,L\varphi=\dv(w\nabla\varphi)$, by $\partial_i\log w=\beta\sum_{j\ne i}\frac1{x_i-x_j}-\frac{\beta N}2x_i$. In detail: for $x\in\Sig$ and $1\le i\le N$,
\begin{equation}\label{c5:dw}
\partial_iw(x)=w(x)\,b_i(x),\qquad b_i(x):=\beta\sum_{j\ne i}\frac1{x_i-x_j}-\frac{\beta N}{2}x_i ,
\end{equation}
and consequently, for every $\varphi\in C^2(\R^N)$ and $x\in\Sig$,
\begin{equation}\label{c5:div}
\dv(w\nabla\varphi)(x)=\beta N\,w(x)\,(L\varphi)(x).
\end{equation}
\emph{Proof.} By the product rule, $\partial_i\Delta=\sum_{k<l}\partial_i(x_l-x_k)\prod_{(k',l')\ne(k,l)}(x_{l'}-x_{k'})$; only the factors with $i\in\{k,l\}$ contribute, with $\partial_i(x_l-x_i)=-1$ and $\partial_i(x_i-x_k)=1$. On $\Sig$, $\Delta\ne0$, so dividing by $\Delta$,
$\partial_i\Delta/\Delta=\sum_{l>i}\frac{-1}{x_l-x_i}+\sum_{k<i}\frac1{x_i-x_k}=\sum_{j\ne i}\frac1{x_i-x_j}$.
Hence $\partial_i(\Delta^\beta)=\beta\Delta^{\beta-1}\partial_i\Delta=\Delta^\beta\beta\sum_{j\ne i}(x_i-x_j)^{-1}$, and $\partial_ie^{-\frac{\beta N}4|x|^2}=-\frac{\beta N}2x_ie^{-\frac{\beta N}4|x|^2}$; this is \eqref{c5:dw}. Then
$\dv(w\nabla\varphi)=\sum_i\partial_i(w\,\partial_i\varphi)=w\sum_i\partial_i^2\varphi+w\sum_ib_i\partial_i\varphi=\beta N\,w\big[\frac1{\beta N}\sum_i\partial_i^2\varphi+\sum_i\frac{b_i}{\beta N}\partial_i\varphi\big]$ (to avoid a clash with the Vandermonde product $\Delta(x)$ we write the Laplacian as $\sum_i\partial_i^2$ here), and $\frac{b_i}{\beta N}=-\frac{x_i}2+\frac1N\sum_{j\ne i}\frac1{x_i-x_j}$ is the coefficient of $\partial_i$ in \eqref{L}. \hfill$\square$

\begin{lemma}[Divergence theorem on $\Sig$ for fields vanishing on the walls]\label{lem:div}
Let $\mathbf Y\in C^1(\R^N;\R^N)$ vanish on every wall $\{x_i=x_{i+1}\}$, $1\le i<N$, and assume $|\mathbf Y|\in L^1(\Sig)$ and $\dv\mathbf Y\in L^1(\Sig)$. Then $\int_{\Sig}\dv\mathbf Y\,\rd x=0$.
The conclusion holds under the weaker hypothesis that $\mathbf Y(x)=0$ for every $x\in\partial\Sig$.
\end{lemma}

\begin{proof}
We prove the lemma under the weaker hypothesis $\mathbf Y=0$ on $\partial\Sig$ (which follows from the first one, since every point of $\partial\Sig$ lies on a wall).
\emph{Integrability.} $\Sig$ is open, hence Borel; $\dv\mathbf Y=\sum_i\partial_i\mathbf Y_i$ is continuous.

\emph{Step 1: $\mathbf Y$ vanishes outside a ball $\{|x|\le r_0\}$.} Then each $\partial_i\mathbf Y_i$ is continuous with compact support, so $\partial_i\mathbf Y_i\mathbf 1_{\Sig}\in L^1(\R^N)$ and $\int_{\Sig}\dv\mathbf Y=\sum_i\int_{\Sig}\partial_i\mathbf Y_i$. Fix $i$.
For $x\in\R^N$ write $\hat x:=(x_1,\dots,x_{i-1},x_{i+1},\dots,x_N)\in\R^{N-1}$; the map $x\mapsto(x_i,\hat x)$ is a coordinate permutation $\R^N\to\R\times\R^{N-1}$ and preserves Lebesgue measure.
Let $\Sigma':=\{\hat x\in\R^{N-1}:\hat x\text{ is strictly increasing}\}$ (for $N=1$, $\R^0$ is a point and $\Sigma'=\R^0$), and put $a(\hat x):=x_{i-1}$ if $i>1$, $a(\hat x):=-\infty$ if $i=1$, and $b(\hat x):=x_{i+1}$ if $i<N$, $b(\hat x):=+\infty$ if $i=N$. Then
\begin{equation}\label{c5:slice}
x\in\Sig\iff\hat x\in\Sigma'\ \text{ and }\ a(\hat x)<x_i<b(\hat x).
\end{equation}
Indeed, ``$\Rightarrow$'': a subsequence of a strictly increasing sequence is strictly increasing. ``$\Leftarrow$'': the consecutive pairs $(x_k,x_{k+1})$ with $i\notin\{k,k+1\}$ are consecutive in $\hat x$, and the pairs $(x_{i-1},x_i)$, $(x_i,x_{i+1})$ are ordered by $a<x_i<b$.
By Fubini's theorem and \eqref{c5:slice},
\[\int_{\Sig}\partial_i\mathbf Y_i\,\rd x=\int_{\Sigma'}\Big(\int_{a(\hat x)}^{b(\hat x)}\partial_i\mathbf Y_i(x)\,\rd x_i\Big)\rd\hat x .\]
Fix $\hat x\in\Sigma'$, write $a=a(\hat x)$, $b=b(\hat x)$ (so $a<b$), and let $h(r):=\mathbf Y_i(x_1,\dots,x_{i-1},r,x_{i+1},\dots,x_N)$. Then $h\in C^1(\R)$, $h'(r)=\partial_i\mathbf Y_i(\dots,r,\dots)$, and $h(r)=0$ for $|r|>r_0$.
If $a>-\infty$, the point with $i$-th coordinate $r=a=x_{i-1}$ satisfies $x_1<\dots<x_{i-1}=r<x_{i+1}<\dots<x_N$, so it lies in $\partial\Sig$ and $h(a)=0$; likewise $h(b)=0$ if $b<\infty$.
Since $h=0$ on the open set $\{|r|>r_0\}$, also $h'=0$ there. Put $a':=a$ if $a>-\infty$; if $a=-\infty$, put $a':=-r_0-1-|b|$ when $b<\infty$ and $a':=-r_0-1$ when $b=\infty$. Symmetrically, $b':=b$ if $b<\infty$; if $b=\infty$, $b':=r_0+1+|a|$ when $a>-\infty$ and $b':=r_0+1$ when $a=-\infty$.
Then $a\le a'<b'\le b$, $a',b'$ are finite, $h(a')=h(b')=0$ (either by the wall condition or because $|a'|,|b'|>r_0$), and $h'=0$ on $(a,a')\cup(b',b)$. By the fundamental theorem of calculus
$\int_a^bh'(r)\,\rd r=\int_{a'}^{b'}h'(r)\,\rd r=h(b')-h(a')=0$.
Hence $\int_{\Sig}\partial_i\mathbf Y_i=0$ for each $i$, and $\int_{\Sig}\dv\mathbf Y=0$.

\emph{Step 2: general $\mathbf Y$.} Let $\chi:\R\to[0,1]$, $\chi(\tau):=1$ for $\tau\le1$, $\chi(\tau):=1-3(\tau-1)^2+2(\tau-1)^3$ for $1\le\tau\le2$, $\chi(\tau):=0$ for $\tau\ge2$. Then $\chi\in C^1(\R)$ (the pieces match with value $1$, slope $0$ at $\tau=1$ and value $0$, slope $0$ at $\tau=2$), $\chi'(\tau)=-6(\tau-1)(2-\tau)$ on $[1,2]$, so $|\chi'|\le\frac32$, and $\chi$ is nonincreasing.
For $r\ge1$ put $\chi_r(x):=\chi(|x|^2/r^2)\in C^1(\R^N)$: $\chi_r=1$ on $\{|x|\le r\}$, $\chi_r=0$ on $\{|x|\ge\sqrt2r\}$, and $\nabla\chi_r(x)=\chi'(|x|^2/r^2)\,2x/r^2$ vanishes unless $|x|\le\sqrt2r$, so
\[|\nabla\chi_r(x)|\le\tfrac32\cdot\frac{2\sqrt2\,r}{r^2}=\frac{3\sqrt2}{r}.\]
The field $\chi_r\mathbf Y$ is $C^1$, vanishes outside $\{|x|\le\sqrt2r\}$ and on $\partial\Sig$. By Step 1 and $\dv(\chi_r\mathbf Y)=\chi_r\dv\mathbf Y+\nabla\chi_r\cdot\mathbf Y$ (both terms are in $L^1(\Sig)$ by the hypotheses),
\[0=\int_{\Sig}\chi_r\,\dv\mathbf Y\,\rd x+\int_{\Sig}\nabla\chi_r\cdot\mathbf Y\,\rd x .\]
As $r\to\infty$ along the integers, $\chi_r\dv\mathbf Y\to\dv\mathbf Y$ pointwise with $|\chi_r\dv\mathbf Y|\le|\dv\mathbf Y|\in L^1(\Sig)$, so the first integral tends to $\int_{\Sig}\dv\mathbf Y$ by dominated convergence; the second is bounded by $\frac{3\sqrt2}r\int_{\Sig}|\mathbf Y|\to0$. Hence $\int_{\Sig}\dv\mathbf Y=0$.
\end{proof}

\begin{corollary}[Integration by parts for $L$]\label{cor:ibp}
Let $\varphi\in C^2(\R^N)$, $k\in C^1(\R^N)$, and assume $k\,|\nabla\varphi|$, $\nabla k\cdot\nabla\varphi$ and $k\,L\varphi$ are in $L^1(\mu^*)$. Then
\begin{equation}\label{ibp}
\int_{\Sig}k\,L\varphi\,\rd\mu^*=-\frac1{\beta N}\int_{\Sig}\nabla k\cdot\nabla\varphi\,\rd\mu^* .
\end{equation}
\end{corollary}
\noindent(Here $L\varphi$ is defined on $\Sig$, which carries $\mu^*$; ``$\in L^1(\mu^*)$'' means $\int_{\Sig}|\cdot|\,w\,\rd x<\infty$.)

\begin{proof}
Let $\mathbf Y:=k\,w\,\nabla\varphi$. Since $k\in C^1$, $w\in C^\infty$ and $\nabla\varphi\in C^1$, $\mathbf Y\in C^1(\R^N;\R^N)$; and $\mathbf Y=0$ on $\partial\Sig$ because $w=0$ there.
$\int_{\Sig}|\mathbf Y|\rd x=Z_\beta\int_{\Sig}|k||\nabla\varphi|\,\rd\mu^*<\infty$.
On $\Sig$, by the product rule and \eqref{c5:div},
$\dv\mathbf Y=w\,\nabla k\cdot\nabla\varphi+k\,\dv(w\nabla\varphi)=w\,\nabla k\cdot\nabla\varphi+\beta N\,k\,w\,L\varphi$,
and both terms are in $L^1(\Sig)$ by hypothesis. Lemma \ref{lem:div} (in the form with $\mathbf Y=0$ on $\partial\Sig$) gives
$0=\int_{\Sig}\dv\mathbf Y\,\rd x=Z_\beta\big(\int_{\Sig}\nabla k\cdot\nabla\varphi\,\rd\mu^*+\beta N\int_{\Sig}k\,L\varphi\,\rd\mu^*\big)$, which is \eqref{ibp}.
\end{proof}

\begin{theorem}[Joint eigenvalue density of GOE/GUE]\label{thm:JD}
For $\beta\in\{1,2\}$ and all $N\ge1$: $\ \mu=\lambda_*\bG_{N,\beta}=\mu^*$. In particular $\mu(\partial\Sig)=0$, i.e.\ $\bG$-a.e.\ matrix has simple spectrum.
Moreover, for every random $H_0\in\cH$ (independent of $V$) and every $t>0$, $\nu_t\ll\mu$ and $\nu_t(\partial\Sig)=0$.
\end{theorem}

\begin{proof}
Both $\mu$ and $\mu^*$ are Borel probability measures on $\R^N$ ($\mu$ is carried by $\cSig$ because $\lambda$ takes values there). Let $\theta:\R^N\to\R$ be bounded and continuous, and put $\Phi:=\theta\circ\lambda$.
By Lemma \ref{lem:inv}, $\Phi$ is invariant, and it is bounded and continuous on $\cH$ ($\lambda$ is continuous, Lemma \ref{lem:inv}(a)). For $u>0$ let $\Phi_u:=T_u\Phi$ and $\varphi_u(x):=\Phi_u(\diag x)$.

\emph{Step 1 (regularity and the equation).} By Lemma \ref{lem:M3}, $(u,H)\mapsto\Phi_u(H)$ is $C^\infty$ on $(0,\infty)\times\cH$, $\partial_u\Phi_u=A\Phi_u$, and for $0<u_1<u_2$ there is $C$ (depending on $u_1,u_2,N,\|\theta\|_\infty$) with
$|\partial_u^j\partial^\alpha_y\Phi_u(H)|\le C(1+|H|)^{2j+|\alpha|}$ for $u\in[u_1,u_2]$, $j\le1$, $|\alpha|\le2$.
Hence $(u,x)\mapsto\varphi_u(x)$ is $C^\infty$ on $(0,\infty)\times\R^N$, and, using $\partial_i\varphi_u(x)=\partial_{E_{ii}}\Phi_u(\diag x)$ and $|\diag x|=|x|$: for $u\in[u_1,u_2]$, $x\in\R^N$,
\begin{equation}\label{c5:JDbounds}
|\varphi_u(x)|\le\|\theta\|_\infty,\qquad|\nabla\varphi_u(x)|\le\sqrt N\,C(1+|x|),\qquad|\partial_u\varphi_u(x)|\le C(1+|x|)^2 .
\end{equation}
The first bound holds because $T_u$ is positive with $T_u1=1$ (Lemma \ref{lem:M1}(i)), applied to $\|\theta\|_\infty\pm\Phi\ge0$.
$\Phi_u$ is invariant (Lemma \ref{lem:inv}) and $C^2$, so by Lemma \ref{lem:givens}, for $x\in\Sig$:
\begin{equation}\label{c5:JDeq}
\partial_u\varphi_u(x)=(A\Phi_u)(\diag x)=(L\varphi_u)(x).
\end{equation}

\emph{Step 2 ($m$ is constant).} Let $m(u):=\int_{\Sig}\varphi_u\,\rd\mu^*$, $u>0$. Fix $0<u_1<u_2$. On $(u_1,u_2)$, $u\mapsto\varphi_u(x)$ is differentiable for every $x$, $x\mapsto\varphi_u(x)$ is continuous, and $|\partial_u\varphi_u(x)|\le C(1+|x|)^2$, which is $\mu^*$-integrable by \eqref{c5:mom}. By differentiation under the integral sign, $m$ is differentiable on $(u_1,u_2)$ with
$m'(u)=\int_{\Sig}\partial_u\varphi_u\,\rd\mu^*=\int_{\Sig}L\varphi_u\,\rd\mu^*$ (by \eqref{c5:JDeq}).
Apply Corollary \ref{cor:ibp} with $k\equiv1$ and $\varphi=\varphi_u\in C^2(\R^N)$: $k|\nabla\varphi_u|\le\sqrt NC(1+|x|)$ and $k\,L\varphi_u=\partial_u\varphi_u$ (on $\Sig$) are $\mu^*$-integrable by \eqref{c5:JDbounds}, \eqref{c5:mom}, and $\nabla k\cdot\nabla\varphi_u=0$. Hence $m'(u)=0$.
As $u_1<u_2$ were arbitrary, $m$ is differentiable on $(0,\infty)$ with $m'\equiv0$, hence constant there.

\emph{Step 3 (the two limits).} By Lemma \ref{lem:M3} (with $\Phi$ bounded and continuous):
as $u\to\infty$, $\varphi_u(x)\to\int\Phi\,\rd\bG=\int\theta\circ\lambda\,\rd\bG=\int\theta\,\rd\mu$ for every $x$; as $u\downarrow0$, $\varphi_u(x)\to\Phi(\diag x)=\theta(\lambda(\diag x))=\theta(x)$ for every $x\in\cSig$ (Lemma \ref{lem:inv}(a)).
Since $|\varphi_u|\le\|\theta\|_\infty$ and $\mu^*$ is a probability measure carried by $\Sig$, dominated convergence (along sequences $u_n\to\infty$, $u_n\downarrow0$) gives $m(u)\to\int\theta\,\rd\mu$ as $u\to\infty$ and $m(u)\to\int_{\Sig}\theta\,\rd\mu^*=\int\theta\,\rd\mu^*$ as $u\downarrow0$. As $m$ is constant,
\[\int_{\R^N}\theta\,\rd\mu=\int_{\R^N}\theta\,\rd\mu^*\qquad\text{for every bounded continuous }\theta:\R^N\to\R .\]

\emph{Step 4 (conclusion).} Two finite Borel measures on the metric space $\R^N$ with the same integrals of all bounded continuous functions are equal. Hence $\mu=\mu^*$.
In particular $\mu(\partial\Sig)=\mu^*(\partial\Sig)=0$, as $\mu^*(\R^N\setminus\Sig)=0$. Since $\lambda(H)\in\Sig$ iff the eigenvalues of $H$ are pairwise distinct, $\bG(\{H:H\text{ has a multiple eigenvalue}\})=\mu(\partial\Sig)=0$.
Finally, for $t>0$ the proof of Proposition \ref{prop:U9}(i) (which uses only Lemma \ref{lem:M1}(iv), not the moment assumption) shows $\nu_t\ll\mu$; hence $\nu_t(\partial\Sig)=0$.
\end{proof}

\begin{remark}
This replaces the Weyl integration formula (not available in Mathlib; a chart-based proof needs a parametrisation of $U(N)/\mathbb T^N$ and a Jacobian computation). The proof uses only Lemmas \ref{lem:M1}(i), \ref{lem:M3}, \ref{lem:inv}, \ref{lem:givens}, \ref{lem:div} and Corollary \ref{cor:ibp} --- all needed anyway --- and the continuity of $\lambda$ (Weyl's inequality).
It does not use Proposition \ref{prop:U9}, so there is no circularity: Proposition \ref{prop:U9} uses $\mu=\lambda_*\bG$ only as a definition.
Only the eigenvalue law of the \emph{Gaussian} ensemble is identified; the disintegration of Lebesgue measure on $\cH$ along conjugation orbits is never used.
For $N=1$ the statement reduces to $\bG_{1,\beta}=N(0,2/\beta)$, which is the definition of $\bG$.
\end{remark}

\begin{definition}[Local relaxation measure]\label{def:omega}
For $R>0$ let $W(x)=W_R(x):=\frac1{2R^2}|x-\gamma|^2$ and
\begin{equation}\label{omega}
\omega=\omega_R:=\frac{e^{-NW}}{\int e^{-NW}\rd\mu}\,\mu=\frac{1}{Z_\omega}e^{-\cV_R(x)}\mathbf 1_{\Sig}(x)\,\rd x,\qquad \cV_R(x):=\frac{\beta N}{4}|x|^2+\frac{N}{2R^2}|x-\gamma|^2-\beta\sum_{i<j}\log(x_j-x_i).
\end{equation}
(The second equality is Theorem \ref{thm:JD}.) $\psi:=\rd\omega/\rd\mu=e^{-NW}/\int e^{-NW}\rd\mu$.
\end{definition}

\noindent\emph{Comments.}
(1) \emph{Well-definedness.} $0<e^{-NW}\le1$, so $\int e^{-NW}\rd\mu\in(0,1]$ and $\omega$ is a Borel probability measure on $\R^N$. The function $\psi(x)=e^{-NW(x)}/\int e^{-NW}\rd\mu$ is $C^\infty$ and positive on all of $\R^N$; it is a version of $\rd\omega/\rd\mu$, and $\omega$ and $\mu$ are mutually absolutely continuous.
(2) \emph{The second equality in \eqref{omega}.} By Theorem \ref{thm:JD}, $\mu=Z_\beta^{-1}w\mathbf 1_{\Sig}\rd x$. On $\Sig$, $\Delta>0$ and $\Delta^\beta=\exp\big(\beta\sum_{i<j}\log(x_j-x_i)\big)$, so $e^{-NW}w=e^{-\cV_R}$ there. Hence $\omega=Z_\omega^{-1}e^{-\cV_R}\mathbf 1_{\Sig}\rd x$ with
$Z_\omega:=\int_{\Sig}e^{-\cV_R}\rd x=Z_\beta\int e^{-NW}\rd\mu\in(0,\infty)$.
(3) \emph{Convexity.} $\omega$ is a convex log-gas in the sense of Definition \ref{def:loggas} with $\cV_0(x)=\frac{\beta N}4|x|^2+\frac N{2R^2}|x-\gamma|^2$, which is $\kappa$-convex with $\kappa=N(\frac\beta2+\frac1{R^2})\ge\frac N{R^2}$ (Comment (2) after Definition \ref{def:loggas}); likewise $\mu^*$ ($R=\infty$, $\kappa=\beta N/2$).
(4) \emph{Derivatives used later.} $\nabla W(x)=(x-\gamma)/R^2$, $|\nabla W|^2=|x-\gamma|^2/R^4$, $\nabla\log\psi=-N\nabla W$, and $NW(x)\le\frac N{R^2}(|x|^2+|\gamma|^2)$.
(5) \emph{Notation.} Because $\psi$ denotes $\rd\omega/\rd\mu$, the function $t\mapsto t\log t$ is written $\mathfrak h$ in \S\ref{sec:C6}. Only $\gamma\in\R^N$ (the vector of classical locations, \S\ref{ss:notation}) enters; the statements of this section hold for every fixed $\gamma\in\R^N$.

\section{\INFRA\ Pr\'ekopa--Leindler; LSI and gap concentration for convex log-gases}\label{sec:C6}

\noindent\emph{Conventions of this section.}
All integrals of nonnegative Borel functions are Lebesgue integrals with values in $[0,\infty]$; $|A|$ denotes the Lebesgue measure of a Borel set $A$.
In $[0,\infty]$ we use $0\cdot\infty=0$, $\infty^p=\infty$ and $0^p=0$ for $p>0$ (these are the conventions of Mathlib's \texttt{ENNReal}); with them $(a,b)\mapsto a^\tau b^{1-\tau}$ is nondecreasing in each variable on $[0,\infty]^2$.
$N\ge1$ is an integer, $\Sig=\{x\in\R^N:x_1<\dots<x_N\}$ (\S\ref{ss:notation}); $\Sig$ is open, nonempty and convex: if $x,y\in\Sig$ and $\tau\in[0,1]$ then $z:=\tau x+(1-\tau)y$ satisfies $z_j-z_i=\tau(x_j-x_i)+(1-\tau)(y_j-y_i)>0$ for $i<j$.
$|\cdot|$ is the Euclidean norm on $\R^N$.
In this section $\beta\ge1$ is an arbitrary real number (in the application $\beta\in\{1,2\}$), and nothing from Sections \ref{sec:C1}--\ref{sec:C5} is used except the notation $\Sig$ (the measures $\mu^*$ and $\omega_R$ appear only as examples).
The section is self-contained relative to Mathlib; the Mathlib facts used are listed in Appendix~\ref{app:mathlib}.
Auxiliary statements are proved as numbered steps inside the proofs.

\subsection{Pr\'ekopa--Leindler}

\begin{theorem}[Pr\'ekopa--Leindler]\label{thm:PL}
Let $N\ge1$, $\tau\in(0,1)$ and let $u,v,w:\R^N\to[0,\infty]$ be Borel with
\begin{equation}\label{c6:PLhyp}
w(\tau x+(1-\tau)y)\ \ge\ u(x)^\tau v(y)^{1-\tau}\qquad\text{for all }x,y\in\R^N .
\end{equation}
Then
\begin{equation}\label{c6:PLconcl}
\int_{\R^N} w\,\rd x\ \ge\ \Big(\int_{\R^N} u\,\rd x\Big)^\tau\Big(\int_{\R^N} v\,\rd x\Big)^{1-\tau}.
\end{equation}
\end{theorem}

\begin{proof}
\emph{Step 0 (reductions valid in every dimension).}
If $\int u=0$ or $\int v=0$, the right side of \eqref{c6:PLconcl} is $0$ and there is nothing to prove. So assume $\int u>0$ and $\int v>0$.
For $k\in\mathbb N$ let $u_k:=\min(u,k)$, $v_k:=\min(v,k)$. These are Borel, bounded, $u_k\le u$, $v_k\le v$, so $(u_k,v_k,w)$ satisfies \eqref{c6:PLhyp} by monotonicity of $(a,b)\mapsto a^\tau b^{1-\tau}$.
By monotone convergence $\int u_k\uparrow\int u$ and $\int v_k\uparrow\int v$; in particular $\int u_k>0$ and $\int v_k>0$ for all large $k$. Moreover $a_k^\tau b_k^{1-\tau}\to a^\tau b^{1-\tau}$ whenever $a_k\uparrow a$, $b_k\uparrow b$ in $[0,\infty]$ with $a,b>0$ (continuity on $(0,\infty)^2$ if $a,b<\infty$; if, say, $a=\infty$, then $a_k^\tau\to\infty$ while $b_k^{1-\tau}$ is eventually bounded below by a positive constant). So if \eqref{c6:PLconcl} holds for $(u_k,v_k,w)$ for all large $k$, it holds for $(u,v,w)$. Hence
\begin{equation}\label{c6:PLred}
\text{it suffices to prove \eqref{c6:PLconcl} when $u,v$ are bounded and $\textstyle\int u>0$, $\int v>0$.}
\end{equation}

\emph{Step 1 (Brunn--Minkowski on the line).} \emph{For nonempty compact $K_1,K_2\subset\R$, the set $K_1+K_2:=\{a+b:a\in K_1,b\in K_2\}$ is compact and $|K_1+K_2|\ge|K_1|+|K_2|$.}
Indeed, $K_1+K_2$ is the image of the compact set $K_1\times K_2$ under the continuous map $(a,b)\mapsto a+b$.
Let $a^*:=\max K_1$, $b_*:=\min K_2$. Then $K_1+b_*\subset K_1+K_2$ and $a^*+K_2\subset K_1+K_2$; moreover $K_1+b_*\subset(-\infty,a^*+b_*]$ and $a^*+K_2\subset[a^*+b_*,\infty)$, so $(K_1+b_*)\cap(a^*+K_2)\subset\{a^*+b_*\}$ is a Lebesgue null set. By additivity and translation invariance of Lebesgue measure,
$|K_1+K_2|\ge|(K_1+b_*)\cup(a^*+K_2)|=|K_1+b_*|+|a^*+K_2|=|K_1|+|K_2|$.

\emph{Step 2 (level sets; no measurability of sumsets is needed).} \emph{Let $\tau\in(0,1)$, let $A,B,C\subset\R$ be Borel, $A,B$ nonempty, with $\tau a+(1-\tau)b\in C$ for all $a\in A$, $b\in B$. Then $|C|\ge\tau|A|+(1-\tau)|B|$.}
Lebesgue measure on $\R$ is inner regular on Borel sets: $|A|=\sup\{|K|:K\subset A\text{ compact}\}$; since $A\neq\emptyset$ and $|K\cup\{a\}|=|K|$, the supremum may be restricted to nonempty compact $K$. The same holds for $B$.
For nonempty compact $K_1\subset A$, $K_2\subset B$, the set $\tau K_1+(1-\tau)K_2$ is compact (hence Borel), contained in $C$, and by Step 1 and the scaling property $|cK|=c|K|$ ($c>0$):
$|C|\ge|\tau K_1+(1-\tau)K_2|\ge\tau|K_1|+(1-\tau)|K_2|$.
Taking the supremum over $K_1$ and over $K_2$ proves the claim (both sides are in $[0,\infty]$).

\emph{Step 3 (the case $N=1$).}
By \eqref{c6:PLred} we may assume $u,v$ bounded with positive integrals. Then $a:=\sup_\R u\in(0,\infty)$ and $b:=\sup_\R v\in(0,\infty)$ (true suprema, not essential suprema).
The functions $\tilde u:=u/a$, $\tilde v:=v/b$, $\tilde w:=w/(a^\tau b^{1-\tau})$ satisfy \eqref{c6:PLhyp}, and \eqref{c6:PLconcl} for $(\tilde u,\tilde v,\tilde w)$ is equivalent to \eqref{c6:PLconcl} for $(u,v,w)$ (both sides are divided by $a^\tau b^{1-\tau}$). So we may assume
\[\sup u=\sup v=1 .\]
For $r\in(0,1)$ put $A_r:=\{u>r\}$, $B_r:=\{v>r\}$, $C_r:=\{w>r\}$. These sets are Borel; $A_r$ and $B_r$ are nonempty because $r<1=\sup u=\sup v$; and if $x\in A_r$, $y\in B_r$ then by \eqref{c6:PLhyp}
$w(\tau x+(1-\tau)y)\ge u(x)^\tau v(y)^{1-\tau}>r^\tau r^{1-\tau}=r$, i.e.\ $\tau x+(1-\tau)y\in C_r$.
By Step 2,
\begin{equation}\label{c6:PLlevel}
|C_r|\ \ge\ \tau|A_r|+(1-\tau)|B_r|\qquad(0<r<1).
\end{equation}
The functions $r\mapsto|A_r|,|B_r|,|C_r|$ are nonincreasing, hence Borel. By the layer-cake formula $\int f\,\rd x=\int_0^\infty|\{f>r\}|\,\rd r$ (valid for Borel $f\ge0$), and since $\{u>r\}=\{v>r\}=\emptyset$ for $r\ge1$,
\begin{align*}\int w\ &\ge\ \int\min(w,1)\ =\ \int_0^1|C_r|\,\rd r\ \ge\ \tau\int_0^1|A_r|\,\rd r+(1-\tau)\int_0^1|B_r|\,\rd r\\
&=\ \tau\int u+(1-\tau)\int v\ \ge\ \Big(\int u\Big)^\tau\Big(\int v\Big)^{1-\tau}.\end{align*}
Here the layer-cake formula is applied to the real-valued functions $\min(w,1)$, $u$ and $v$ (for $r<1$, $\{\min(w,1)>r\}=C_r$), so that no $[0,\infty]$-valued version of it is needed.
The last step is the weighted arithmetic--geometric mean inequality $\tau s+(1-\tau)t\ge s^\tau t^{1-\tau}$ for $s,t\in[0,\infty)$; if $s$ or $t$ is $+\infty$, the left side is $+\infty$.
This proves the theorem for $N=1$ (for functions with values in $[0,\infty]$).

\emph{Step 4 (induction $N-1\to N$, $N\ge2$).}
Identify $\R^N=\R\times\R^{N-1}$, $x=(x_1,x')$; Lebesgue measure on $\R^N$ is the product of the Lebesgue measures on $\R$ and $\R^{N-1}$. Define the marginals
\[\bar u(x_1):=\int_{\R^{N-1}}u(x_1,x')\,\rd x',\qquad\bar v(x_1):=\int_{\R^{N-1}}v(x_1,x')\,\rd x',\qquad\bar w(x_1):=\int_{\R^{N-1}}w(x_1,x')\,\rd x' .\]
By Tonelli's theorem $\bar u,\bar v,\bar w:\R\to[0,\infty]$ are Borel and $\int_\R\bar u=\int_{\R^N}u$, and likewise for $v,w$.
Fix $x_1,y_1\in\R$ and put $z_1:=\tau x_1+(1-\tau)y_1$. The sections $u(x_1,\cdot)$, $v(y_1,\cdot)$, $w(z_1,\cdot)$ are Borel functions $\R^{N-1}\to[0,\infty]$, and for all $x',y'\in\R^{N-1}$, by \eqref{c6:PLhyp} at the points $(x_1,x')$, $(y_1,y')$,
\[w\big(z_1,\tau x'+(1-\tau)y'\big)=w\big(\tau(x_1,x')+(1-\tau)(y_1,y')\big)\ge u(x_1,x')^\tau v(y_1,y')^{1-\tau}.\]
By the induction hypothesis (the theorem in dimension $N-1$), $\bar w(z_1)\ge\bar u(x_1)^\tau\bar v(y_1)^{1-\tau}$.
Since $x_1,y_1$ were arbitrary, $(\bar u,\bar v,\bar w)$ satisfies \eqref{c6:PLhyp} on $\R$, with values in $[0,\infty]$. By Step 3,
$\int_{\R^N}w=\int_\R\bar w\ge(\int_\R\bar u)^\tau(\int_\R\bar v)^{1-\tau}=(\int_{\R^N}u)^\tau(\int_{\R^N}v)^{1-\tau}$.
\end{proof}

\smallskip\noindent\emph{On the formulation.}
The marginals in Step 4 may take the value $+\infty$ even when $u,v,w$ are finite; this is the only reason for allowing values in $[0,\infty]$ in the statement, and it removes the truncation of the marginals. In the applications below $u,v,w$ are finite, bounded on compacts of $\Sig$ and vanish outside $\Sig$.
The sumset $\tau A_r+(1-\tau)B_r$ of two Borel sets need not be Borel; Step 2 avoids it.

\subsection{Convexity of the log-gas Hamiltonian}

For $r>0$ let
\begin{equation}\label{c6:defLambda}
\Lambda(r):=\log\frac{1+r}{2\sqrt r}=\log\cosh\big(\tfrac12\log r\big)\ \ge0 .
\end{equation}
(The two expressions agree since $\cosh(\frac12\log r)=\frac12(\sqrt r+1/\sqrt r)$; $\Lambda\ge0$ since $\cosh\ge1$.)
As $\cosh$ is even and increasing on $[0,\infty)$, $\Lambda(r)=\ell(|\log r|)$ with $\ell(\vartheta):=\log\cosh(\vartheta/2)$ nondecreasing on $[0,\infty)$; in particular $\Lambda(r)=\Lambda(1/r)$.

\begin{definition}\label{def:loggas}
A \emph{convex log-gas} (of size $N$, with parameters $\beta\ge1$, $\kappa\ge0$) is a probability measure on $\R^N$ of the form
\[\omega=Z^{-1}e^{-\cV}\mathbf 1_{\Sig}\,\rd x,\quad \cV(x):=\cV_0(x)-\beta\!\!\sum_{1\le i<j\le N}\!\!\log(x_j-x_i)\ \ (x\in\Sig),\quad Z:=\int_{\Sig}e^{-\cV}\rd x<\infty,\]
where $\cV_0:\R^N\to\R$ is continuous and \emph{$\kappa$-convex}: for all $x,y\in\R^N$ and $\tau\in[0,1]$,
\begin{equation}\label{c6:kconv}
\tau\cV_0(x)+(1-\tau)\cV_0(y)-\cV_0(\tau x+(1-\tau)y)\ \ge\ \frac\kappa2\,\tau(1-\tau)\,|x-y|^2 .
\end{equation}
\end{definition}

\noindent\emph{Comments.}
(1) $\cV$ is continuous on the open set $\Sig$, so $e^{-\cV}\mathbf 1_{\Sig}$ (extended by $0$ outside $\Sig$) is a Borel function on $\R^N$, strictly positive on $\Sig$; hence $Z>0$. The requirement $Z<\infty$ is part of the definition; for $\kappa>0$ it is automatic (Lemma \ref{lem:convex}(c)).
We write $\int\varphi\,\rd\omega=\int_{\Sig}\varphi\,\rd\omega=Z^{-1}\int_{\Sig}\varphi\,e^{-\cV}\rd x$. Since $\omega(\R^N\setminus\Sig)=0$, functions need only be defined on $\Sig$.
On $\Sig$ one has $e^{-\cV}=\Delta^\beta e^{-\cV_0}$ with $\Delta(x)=\prod_{i<j}(x_j-x_i)$.

(2) \emph{Examples.} For $a,b\ge0$ and $\gamma\in\R^N$ the function $\cV_0(x)=a|x|^2+b|x-\gamma|^2$ is $\kappa$-convex with $\kappa=2(a+b)$, with equality in \eqref{c6:kconv}: for every $c\in\R^N$,
$\tau|x-c|^2+(1-\tau)|y-c|^2-|\tau x+(1-\tau)y-c|^2=\tau(1-\tau)|x-y|^2$ (expand the squares).
Hence $\mu^*$ of \S\ref{sec:C5} ($a=\beta N/4$, $b=0$) is a convex log-gas with $\kappa=\beta N/2$, and the local relaxation measure $\omega_R$ of \eqref{omega} ($a=\beta N/4$, $b=N/(2R^2)$) is a convex log-gas with
\[\kappa=N\Big(\frac\beta2+\frac1{R^2}\Big)\ \ge\ \frac N{R^2}.\]

\begin{lemma}\label{lem:convex}
For a convex log-gas, $x,y\in\Sig$ and $\tau\in[0,1]$:
\begin{enumerate}[label=\textup{(\alph*)},leftmargin=2em,itemsep=1pt]
\item $\tau\cV(x)+(1-\tau)\cV(y)-\cV(\tau x+(1-\tau)y)\ge\frac\kappa2\tau(1-\tau)|x-y|^2$;
\item $\cV(x)+\cV(y)-2\cV(\frac{x+y}2)\ge2\beta\sum_{i<j}\Lambda\big(\frac{x_j-x_i}{y_j-y_i}\big)$.
\end{enumerate}
Moreover:
\begin{enumerate}[label=\textup{(\alph*)},leftmargin=2em,itemsep=1pt,start=3]
\item
if $\cV_0:\R^N\to\R$ is continuous and $\kappa$-convex with $\kappa>0$, and $\beta\ge0$, then $\int_{\Sig}e^{-\cV}\rd x<\infty$ for $\cV:=\cV_0-\beta\sum_{i<j}\log(x_j-x_i)$; i.e.\ for $\kappa>0$ the condition $Z<\infty$ in Definition \ref{def:loggas} holds automatically.
\end{enumerate}
\end{lemma}

\begin{proof}
Let $x,y\in\Sig$, $\tau\in[0,1]$, $z:=\tau x+(1-\tau)y\in\Sig$. For $i<j$ put $a_{ij}:=x_j-x_i>0$, $b_{ij}:=y_j-y_i>0$, so $z_j-z_i=\tau a_{ij}+(1-\tau)b_{ij}$.

By the definition of $\cV$,
\begin{align}\label{c6:Vdecomp}
\tau\cV(x)+(1-\tau)\cV(y)-\cV(z)={}&\big[\tau\cV_0(x)+(1-\tau)\cV_0(y)-\cV_0(z)\big]\notag\\
&+\beta\sum_{i<j}\big[\log(\tau a_{ij}+(1-\tau)b_{ij})-\tau\log a_{ij}-(1-\tau)\log b_{ij}\big].
\end{align}

(a) By concavity of $\log$ on $(0,\infty)$, every term of the sum in \eqref{c6:Vdecomp} is $\ge0$; as $\beta\ge0$, the second line of \eqref{c6:Vdecomp} is $\ge0$, and the first is $\ge\frac\kappa2\tau(1-\tau)|x-y|^2$ by \eqref{c6:kconv}.

(b) Take $\tau=\frac12$ in \eqref{c6:Vdecomp} and multiply by $2$. The first bracket is $\ge\frac\kappa4|x-y|^2\ge0$ by \eqref{c6:kconv}. For the second,
\[2\log\frac{a+b}2-\log a-\log b=2\log\frac{a+b}{2\sqrt{ab}}=2\log\frac{1+a/b}{2\sqrt{a/b}}=2\Lambda(a/b)\qquad(a,b>0).\]
Hence $\cV(x)+\cV(y)-2\cV(\frac{x+y}2)\ge\frac\kappa4|x-y|^2+2\beta\sum_{i<j}\Lambda(a_{ij}/b_{ij})$, which is stronger than (b).

(c) Let $c_1:=\max_{|e|=1}|\cV_0(e)-\cV_0(0)|<\infty$ (continuity on the compact unit sphere).
For $|x|\ge1$ apply \eqref{c6:kconv} to the pair $(x,0)$ with $\tau:=1/|x|\in(0,1]$, so that $\tau x=e:=x/|x|$:
$\tau\cV_0(x)+(1-\tau)\cV_0(0)-\cV_0(e)\ge\frac\kappa2\tau(1-\tau)|x|^2$. Dividing by $\tau$,
\[\cV_0(x)\ \ge\ \cV_0(0)+|x|\big(\cV_0(e)-\cV_0(0)\big)+\frac\kappa2\big(|x|^2-|x|\big)\ \ge\ \cV_0(0)-\Big(c_1+\frac\kappa2\Big)|x|+\frac\kappa2|x|^2\qquad(|x|\ge1).\]
On $\Sig$, $0<x_j-x_i\le2|x|$, so $e^{-\cV(x)}=\Delta(x)^\beta e^{-\cV_0(x)}\le(2|x|)^{\beta N(N-1)/2}e^{-\cV_0(x)}$. Therefore
\begin{align*}
\int_{\Sig}e^{-\cV}\rd x\le{}&2^{\beta N(N-1)/2}\,|\{|x|\le1\}|\,e^{-\min_{|x|\le1}\cV_0}\\
&+e^{-\cV_0(0)}\int_{|x|>1}(2|x|)^{\beta N(N-1)/2}e^{(c_1+\kappa/2)|x|-\frac\kappa2|x|^2}\rd x<\infty,
\end{align*}
since a polynomial times $e^{c|x|-\frac\kappa2|x|^2}$ is integrable on $\R^N$ for $\kappa>0$ (for instance $c|x|\le\frac\kappa4|x|^2+\frac{c^2}{\kappa}$, and Gaussian moments are finite).
\end{proof}

\subsection{Logarithmic Sobolev inequality on $\Sig$ without semigroups}

\begin{theorem}[LSI for a uniformly convex log-gas]\label{thm:LSI}
Let $\omega$ be a convex log-gas with $\kappa>0$, and let $q\in C^2(\R^N)$, $q>0$, with $\int_{\Sig}q\,\rd\omega=1$. Then
\begin{equation}\label{LSI}
\Ent_\omega(q):=\int_{\Sig}q\log q\,\rd\omega\ \le\ \frac1{2\kappa}\int_{\Sig}\frac{|\nabla q|^2}{q}\,\rd\omega\qquad(\text{the right side may be }+\infty).
\end{equation}
\end{theorem}

\noindent\emph{Notation for the proof.} Let $\mathfrak h(t):=t\log t$ for $t>0$, $\mathfrak h(0):=0$; $\mathfrak h$ is continuous on $[0,\infty)$ and $\mathfrak h\ge-1/e$.
For a Borel function $F:\Sig\to[0,\infty)$ with $\int F\rd\omega<\infty$ put
\[\Ent_\omega(F):=\int\mathfrak h(F)\,\rd\omega-\mathfrak h\Big(\int F\,\rd\omega\Big)\ \in(-\infty,+\infty];\]
the integral is well defined because $\mathfrak h(F)+1/e\ge0$ and $\omega$ is a probability measure. If $\int F\rd\omega=1$ this is the left side of \eqref{LSI}.
In particular the left side of \eqref{LSI} is a well-defined element of $(-\infty,\infty]$, and the theorem asserts that it is finite when the right side is.

\begin{proof}
\emph{Step 1 (exponentials of bounded $C^2$ functions).} \emph{Let $g\in C^2(\R^N)$ be such that
$G_0:=\sup|g|$, $G_1:=\sup|\nabla g|$ and $M:=\sup_{z}\|\nabla^2g(z)\|_{op}$ are finite. Then}
\begin{equation}\label{c6:LSIexp}
\Ent_\omega(e^g)=\int ge^g\,\rd\omega-\int e^g\rd\omega\,\log\int e^g\rd\omega\ \le\ \frac1{2\kappa}\int|\nabla g|^2e^g\,\rd\omega .
\end{equation}

\emph{1a (pointwise inequality).} Put $s_0:=\kappa/(\kappa+M)\in(0,1]$. For $s\in(0,s_0)$ let $\tau:=1-s$; then $\kappa\tau-Ms=\kappa-(\kappa+M)s>0$. Define
\[g_\tau(z):=g(z)+\frac{s\,|\nabla g(z)|^2}{2(\kappa\tau-Ms)}\qquad(z\in\R^N).\]
Let $x,y\in\Sig$, $z:=\tau x+sy\in\Sig$ and $h:=x-y$, so that $x=z+sh$.
Taylor's formula with Lagrange remainder for $\theta(t):=g(z+th)$ on $[0,s]$, with $\theta'(0)=\la\nabla g(z),h\ra$ and $\theta''(t)=\la\nabla^2g(z+th)h,h\ra\le M|h|^2$, gives
\[g(x)\le g(z)+s\la\nabla g(z),h\ra+\tfrac M2s^2|h|^2 .\]
Lemma \ref{lem:convex}(a) gives $\tau\cV(x)+s\cV(y)\ge\cV(z)+\frac\kappa2\tau s|h|^2$. Subtracting,
\begin{align*}
g(x)-\tau\cV(x)-s\cV(y)&\le g(z)-\cV(z)+s\la\nabla g(z),h\ra-\tfrac s2(\kappa\tau-Ms)|h|^2\\
&\le g(z)-\cV(z)+\frac{s|\nabla g(z)|^2}{2(\kappa\tau-Ms)}=g_\tau(z)-\cV(z),
\end{align*}
where we used $\la p,h\ra-\frac c2|h|^2\le\frac{|p|^2}{2c}$ for $c>0$ (namely $\frac c2|h-p/c|^2\ge0$) with $p=\nabla g(z)$, $c=\kappa\tau-Ms$.

\emph{1b (Pr\'ekopa--Leindler).} Define the Borel functions on $\R^N$
\[u:=e^{g/\tau-\cV}\mathbf 1_{\Sig},\qquad v:=e^{-\cV}\mathbf 1_{\Sig},\qquad w:=e^{g_\tau-\cV}\mathbf 1_{\Sig}.\]
For $x,y\in\Sig$, 1a says $u(x)^\tau v(y)^{s}=e^{g(x)-\tau\cV(x)-s\cV(y)}\le w(\tau x+sy)$; if $x\notin\Sig$ or $y\notin\Sig$ then $u(x)^\tau v(y)^s=0$. So \eqref{c6:PLhyp} holds, and Theorem \ref{thm:PL} gives $\int w\ge(\int u)^\tau(\int v)^s$. Dividing by $Z=\int v=Z^\tau Z^s\in(0,\infty)$:
\begin{equation}\label{c6:BLstep}
\Big(\int e^{g/(1-s)}\rd\omega\Big)^{1-s}\ \le\ \int e^{g+c(s)|\nabla g|^2}\rd\omega\qquad(0<s<s_0),\qquad c(s):=\frac{s}{2(\kappa-(\kappa+M)s)} .
\end{equation}

\emph{1c (differentiation at $s=0$).} Let $I:=(-s_0/2,s_0/2)$. For $s\in I$ we have $1-s\ge\frac12$ and $\kappa-(\kappa+M)s\ge\frac\kappa2$, so the functions
\[a(s):=\int e^{g/(1-s)}\rd\omega,\qquad b(s):=\int e^{g+c(s)|\nabla g|^2}\rd\omega\]
are well defined on $I$ (bounded continuous integrands), positive, and $a(0)=b(0)=m:=\int e^g\rd\omega\in[e^{-G_0},e^{G_0}]$.
The integrands are $C^1$ in $s\in I$ for every fixed $x$, with
\begin{align*}
\Big|\partial_se^{g/(1-s)}\Big|&=\frac{|g|}{(1-s)^2}e^{g/(1-s)}\le4G_0e^{2G_0},\\
\Big|\partial_se^{g+c(s)|\nabla g|^2}\Big|&=c'(s)|\nabla g|^2e^{g+c(s)|\nabla g|^2}\le\frac2\kappa G_1^2e^{G_0+G_1^2/\kappa},
\end{align*}
because $c'(s)=\dfrac{\kappa}{2(\kappa-(\kappa+M)s)^2}\in(0,\frac2\kappa]$ and $|c(s)|\le|s|/\kappa\le\frac1\kappa$ on $I$ (the quotient rule gives $c'=\frac{(\kappa-(\kappa+M)s)+(\kappa+M)s}{2(\kappa-(\kappa+M)s)^2}$).
The bounds are constants, hence $\omega$-integrable; by differentiation under the integral sign, $a$ and $b$ are differentiable on $I$ and
\[a'(0)=\int ge^g\,\rd\omega,\qquad b'(0)=c'(0)\int|\nabla g|^2e^g\,\rd\omega=\frac1{2\kappa}\int|\nabla g|^2e^g\,\rd\omega .\]
Hence $\phi(s):=(1-s)\log a(s)-\log b(s)$ is differentiable at $0$, $\phi(0)=0$, and
\[\phi'(0)=-\log m+\frac{a'(0)}{m}-\frac{b'(0)}{m}.\]
By \eqref{c6:BLstep}, $\phi(s)\le0$ for $s\in(0,s_0/2)$. Therefore $\phi'(0)=\lim_{s\downarrow0}\phi(s)/s\le0$
(only the right-hand difference quotient is used, although $\phi$ is differentiable from both sides), i.e.
\[\int ge^g\,\rd\omega-m\log m-\frac1{2\kappa}\int|\nabla g|^2e^g\,\rd\omega\ \le0,\]
which is \eqref{c6:LSIexp} (note $\mathfrak h(e^g)=ge^g$).

\emph{Step 2 (general $q$).} Let $\mathcal I:=\int_{\Sig}|\nabla q|^2/q\,\rd\omega\in[0,\infty]$ (the integrand is continuous and nonnegative). If $\mathcal I=\infty$ there is nothing to prove; assume $\mathcal I<\infty$.
Let $f:=\sqrt q\in C^2(\R^N)$, $f>0$; then $\int f^2\rd\omega=1$, $\nabla f=\nabla q/(2\sqrt q)$ and
\begin{equation}\label{c6:fq}
4\int|\nabla f|^2\rd\omega=\mathcal I .
\end{equation}
\emph{2a (cutoff).} Fix $\chi\in C^\infty(\R)$ with $0\le\chi\le1$, $\chi=1$ on $(-\infty,1]$, $\chi=0$ on $[4,\infty)$, and put $C_\chi:=4\sup|\chi'|<\infty$.
For $r\ge1$ let $\chi_r(x):=\chi(|x|^2/r^2)$. Then $\chi_r\in C^\infty(\R^N)$, $0\le\chi_r\le1$, $\chi_r=1$ on $\{|x|\le r\}$, $\chi_r=0$ on $\{|x|\ge2r\}$, and
$|\nabla\chi_r(x)|=|\chi'(|x|^2/r^2)|\cdot2|x|/r^2\le C_\chi/r$ (the derivative vanishes unless $|x|\le2r$).
For $k\in\mathbb N$ put
\[q_{k,r}:=q\chi_r^2+\frac1k=(f\chi_r)^2+\frac1k,\qquad g_{k,r}:=\log q_{k,r}.\]
Then $q_{k,r}\in C^2(\R^N)$, $q_{k,r}\ge1/k$, so $g_{k,r}\in C^2(\R^N)$; on the open set $\{|x|>2r\}$, $g_{k,r}\equiv-\log k$, so $\nabla g_{k,r}=0$ and $\nabla^2g_{k,r}=0$ there, while on the compact ball $\{|x|\le2r\}$ the continuous functions $g_{k,r},\nabla g_{k,r},\nabla^2g_{k,r}$ are bounded. Thus Step 1 applies to $g=g_{k,r}$; since $e^{g}=q_{k,r}$ and $|\nabla g|^2e^g=|\nabla q_{k,r}|^2/q_{k,r}$, \eqref{c6:LSIexp} reads
\begin{equation}\label{c6:LSIkr}
\Ent_\omega(q_{k,r})\le\frac1{2\kappa}\int\frac{|\nabla q_{k,r}|^2}{q_{k,r}}\,\rd\omega .
\end{equation}
\emph{2b (the right side).} $\nabla q_{k,r}=2(f\chi_r)\nabla(f\chi_r)$, hence
$\dfrac{|\nabla q_{k,r}|^2}{q_{k,r}}=\dfrac{4(f\chi_r)^2}{(f\chi_r)^2+1/k}|\nabla(f\chi_r)|^2\le4|\nabla(f\chi_r)|^2$.
Moreover $|\nabla(f\chi_r)|\le\chi_r|\nabla f|+f|\nabla\chi_r|\le|\nabla f|+\frac{C_\chi}rf$, and $(a_1+a_2)^2\le(1+\eta)a_1^2+(1+\eta^{-1})a_2^2$ for all $a_1,a_2\ge0$, $\eta>0$. With $\int f^2\rd\omega=1$ and \eqref{c6:fq}, for every $\eta>0$:
\begin{equation}\label{c6:LSIrhs}
\int\frac{|\nabla q_{k,r}|^2}{q_{k,r}}\,\rd\omega\le(1+\eta)\,\mathcal I+4(1+\eta^{-1})\frac{C_\chi^2}{r^2}.
\end{equation}
\emph{2c ($k\to\infty$, $r$ fixed).} Let $B_r:=1+\sup_{|x|\le2r}q<\infty$. For all $k$ and $x$: $0\le q\chi_r^2(x)\le q_{k,r}(x)\le B_r$ and $|q_{k,r}(x)-q\chi_r^2(x)|=1/k$. Since $\mathfrak h$ is uniformly continuous on $[0,B_r]$, $\sup_x|\mathfrak h(q_{k,r}(x))-\mathfrak h(q\chi_r^2(x))|\to0$, hence $\int\mathfrak h(q_{k,r})\rd\omega\to\int\mathfrak h(q\chi_r^2)\rd\omega$; also $\int q_{k,r}\rd\omega=m_r+\frac1k\to m_r:=\int q\chi_r^2\rd\omega\in[0,1]$ and $\mathfrak h$ is continuous. So $\Ent_\omega(q_{k,r})\to\Ent_\omega(q\chi_r^2)$, and by \eqref{c6:LSIkr}, \eqref{c6:LSIrhs}
\begin{equation}\label{c6:LSIr}
\int\mathfrak h(q\chi_r^2)\,\rd\omega\le\mathfrak h(m_r)+\frac1{2\kappa}\Big[(1+\eta)\,\mathcal I+4(1+\eta^{-1})\frac{C_\chi^2}{r^2}\Big].
\end{equation}
\emph{2d ($r\to\infty$ along the integers).} For every $x$, $\chi_r(x)=1$ as soon as $r\ge|x|$, so $q\chi_r^2\to q$ and $\mathfrak h(q\chi_r^2)\to\mathfrak h(q)$ pointwise. By dominated convergence ($0\le q\chi_r^2\le q$, $\int q\,\rd\omega=1$), $m_r\to1$ and $\mathfrak h(m_r)\to0$. By Fatou's lemma applied to the nonnegative functions $\mathfrak h(q\chi_r^2)+1/e$,
\[\int\mathfrak h(q)\,\rd\omega\le\liminf_{r\to\infty}\int\mathfrak h(q\chi_r^2)\,\rd\omega\le\frac{1+\eta}{2\kappa}\,\mathcal I .\]
Letting $\eta\downarrow0$ gives \eqref{LSI}.
\end{proof}

\begin{remark}\label{rem:LSI}
(i) No boundary term appears: the domain enters only through the convexity of $\Sig$ (so that $z\in\Sig$ in 1a) and through $\mathbf 1_{\Sig}$ in $u,v,w$. The singularity of $\cV$ at the walls is harmless because only the convexity inequality of Lemma \ref{lem:convex}(a) is used, never a derivative of $\cV$.
The proof uses only two properties of $\omega$: it is a probability measure of the form $Z^{-1}e^{-\cV}\mathbf 1_{\Sig}\rd x$, and Lemma \ref{lem:convex}(a) holds with $\kappa>0$.

(ii) In the application (Theorem \ref{thm:entropy}) $\omega=\omega_R$, $\kappa\ge N/R^2$, and \eqref{LSI} is used in the form $\int q\log q\,\rd\omega\le\frac{R^2}{2N}\int|\nabla q|^2/q\,\rd\omega$.

(iii) The Herbst argument is \emph{not} used in NP: LSI\,+\,Herbst for the observable $F$ of Theorem \ref{thm:conc} only sees $\|\nabla F\|_\infty^2\cdot R^2/N\approx R^2|J|/N$ and would lose a factor $NR^2$ compared with \eqref{conc}; the gain comes from the logarithmic interaction, via Lemma \ref{lem:convex}(b).
\end{remark}

\subsection{Sub-Gaussian concentration of gap observables}

\begin{lemma}[Two-point inequality]\label{lem:2pt}
Let $A,A'\ge0$, $\rho>0$, and let $g:(0,\infty)\to\R$ satisfy $|g|\le A$, $|g(u)-g(v)|\le A'|u-v|$ for all $u,v>0$, and $g=0$ on $(\rho,\infty)$. Then for all $u,v>0$, $s\in\R$:
\[s\,(g(u)-g(v))\le2\Lambda(u/v)+C_\star s^2,\qquad C_\star:=\max\Big(\frac{A^2}{2c_0},\ \tfrac98(e\rho A')^2\Big),\quad c_0:=\log\cosh\tfrac12 .\]
\end{lemma}

\noindent Numerically $c_0=0.1201\ldots$, so $C_\star=\max\big(4.16\ldots A^2,\ 8.31\ldots(\rho A')^2\big)$. The constant depends only on $(A,A',\rho)$.

\begin{proof}
\emph{Step 1 (an elementary bound).} \emph{For $0\le\vartheta\le1$: $\ \ell(\vartheta)=\log\cosh(\vartheta/2)\ge\vartheta^2/9$.}
For real $t$: $\cosh t=1+2\sinh^2(t/2)\ge1+t^2/2$ because $|\sinh x|\ge|x|$. For $y\ge0$: $\log(1+y)\ge y-y^2/2$ (the difference vanishes at $0$ and has derivative $\frac1{1+y}-1+y=\frac{y^2}{1+y}\ge0$).
For $|t|\le\frac12$ put $y:=t^2/2\in[0,\frac18]$: $\log\cosh t\ge\log(1+y)\ge y(1-\frac y2)\ge\frac{15}{16}y=\frac{15}{32}t^2$.
With $t=\vartheta/2$: $\ell(\vartheta)\ge\frac{15}{128}\vartheta^2\ge\frac19\vartheta^2$.

\emph{Step 2.} Fix $u,v>0$, $s\in\R$ and put $\vartheta:=|\log(u/v)|$, so that $\Lambda(u/v)=\ell(\vartheta)$ (see \eqref{c6:defLambda}). Recall $2\alpha\gamma\le\alpha^2+\gamma^2$ for real $\alpha,\gamma$.

\emph{Case 1: $\vartheta\ge1$.} Then $2\Lambda(u/v)\ge2\ell(1)=2c_0$ by monotonicity of $\ell$, and, using $|g|\le A$,
\[s(g(u)-g(v))\le2|s|A=2\cdot\sqrt{2c_0}\cdot\frac{|s|A}{\sqrt{2c_0}}\le2c_0+\frac{A^2}{2c_0}s^2\le2\Lambda(u/v)+C_\star s^2 .\]

\emph{Case 2: $\vartheta<1$ and $\min(u,v)>\rho$.} Then $g(u)=g(v)=0$ and the left side is $0\le2\Lambda(u/v)+C_\star s^2$.

\emph{Case 3: $\vartheta<1$ and $\min(u,v)\le\rho$.} Then $\max(u,v)=e^{\vartheta}\min(u,v)<e\rho$ and
$|u-v|=\max(u,v)\,(1-e^{-\vartheta})\le\max(u,v)\,\vartheta\le e\rho\,\vartheta$. By the Lipschitz bound and Step 1,
\[s(g(u)-g(v))\le|s|A'e\rho\,\vartheta=2\cdot\frac{\sqrt2}3\vartheta\cdot\frac{3}{2\sqrt2}A'e\rho|s|\le\frac29\vartheta^2+\frac98(A'e\rho)^2s^2\le2\ell(\vartheta)+C_\star s^2 .\]
The three cases are exhaustive.
\end{proof}

\begin{theorem}[Gap concentration]\label{thm:conc}
Let $\omega$ be a convex log-gas of size $N$ ($\beta\ge1$, $\kappa\ge0$). Let $n\ge2$, $1\le m_2<\dots<m_n$ integers, $J\subset\{1,\dots,N-m_n\}$, $A,A'\ge0$, $\rho>0$, and let
$G:(0,\infty)^{n-1}\to\R$ be a Borel function, defined on all of $(0,\infty)^{n-1}$, with $|G|\le A$, $G(u)=0$ unless $u\in(0,\rho]^{n-1}$, and $|G(u)-G(u')|\le A'|u_k-u'_k|$ whenever $u,u'\in(0,\infty)^{n-1}$ differ only in the $k$-th coordinate. Put
\begin{equation}\label{defF}
F(x):=\frac1N\sum_{i\in J}G\big(N(x_{i+m_2}-x_i),\dots,N(x_{i+m_n}-x_i)\big),\qquad x\in\Sig .
\end{equation}
Then for all $\lambda\in\R$, with $C_\star=C_\star(A,A',\rho)$ from Lemma \ref{lem:2pt},
\begin{equation}\label{conc}
\int e^{\lambda(F-\int F\rd\omega)}\rd\omega\ \le\ \exp\Big(C_\star(n-1)\frac{|J|}{N^2}\lambda^2\Big).
\end{equation}
\end{theorem}

\noindent\emph{Comments on the hypotheses.}
The coordinates of $u\in(0,\infty)^{n-1}$ are indexed by $k=2,\dots,n$. If $m_n\ge N$ then $J=\emptyset$ and $F\equiv0$.
The coordinatewise Lipschitz condition implies $|G(u)-G(u')|\le A'\sum_k|u_k-u_k'|$ for all $u,u'\in(0,\infty)^{n-1}$ (change one coordinate at a time; the intermediate points stay in the product set $(0,\infty)^{n-1}$), so $G$ is continuous, and the Borel hypothesis is automatic; it is stated for emphasis, as is the fact that $G$ must be defined (and satisfy the hypotheses) on all of $(0,\infty)^{n-1}$, not only near the diagonal or on a set of gaps of full measure.
For $x\in\Sig$ and $i\in J$ all arguments $N(x_{i+m_k}-x_i)$ are positive, so $F$ is well defined on $\Sig$; $F$ is continuous on $\Sig$ and $|F|\le A|J|/N\le A$. In particular $\int F\rd\omega$ and $\int e^{\lambda F}\rd\omega$ are finite.
$C_\star$ does not depend on $N$, $n$, $m$, $J$, $\beta$, $\kappa$, $\cV_0$.

\begin{proof}
Fix $\lambda\in\R$ and put $c_\lambda:=C_\star(n-1)|J|\lambda^2/N^2$.

\emph{Step 1 (telescoping in the coordinates).} \emph{For all $a,b\in(0,\infty)^{n-1}$ and $s\in\R$,}
\begin{equation}\label{c6:telesc}
s\,\big(G(a)-G(b)\big)\ \le\ \sum_{k=2}^n\Big[2\Lambda(a_k/b_k)+C_\star s^2\Big].
\end{equation}
For $k=1,\dots,n$ let $c^{(k)}:=(b_2,\dots,b_k,a_{k+1},\dots,a_n)\in(0,\infty)^{n-1}$, so $c^{(1)}=a$, $c^{(n)}=b$ and $G(a)-G(b)=\sum_{k=2}^n[G(c^{(k-1)})-G(c^{(k)})]$.
For $2\le k\le n$ define $g_k:(0,\infty)\to\R$, $g_k(t):=G(b_2,\dots,b_{k-1},t,a_{k+1},\dots,a_n)$. Then $G(c^{(k-1)})-G(c^{(k)})=g_k(a_k)-g_k(b_k)$, and $g_k$ satisfies the hypotheses of Lemma \ref{lem:2pt} with the same $(A,A',\rho)$: $|g_k|\le A$; $|g_k(t)-g_k(t')|\le A'|t-t'|$ since the two arguments of $G$ differ only in the $k$-th coordinate; and $g_k(t)=0$ for $t>\rho$ since then the argument of $G$ is not in $(0,\rho]^{n-1}$.
Lemma \ref{lem:2pt} gives $s(g_k(a_k)-g_k(b_k))\le2\Lambda(a_k/b_k)+C_\star s^2$; summing over $k$ gives \eqref{c6:telesc}.

\emph{Step 2 (pointwise inequality).} \emph{For all $x,y\in\Sig$,}
\begin{equation}\label{c6:concpt}
\lambda\big(F(x)-F(y)\big)\ \le\ \cV(x)+\cV(y)-2\cV\big(\tfrac{x+y}2\big)+c_\lambda .
\end{equation}
For $i\in J$ and $2\le k\le n$ let $a_{i,k}:=N(x_{i+m_k}-x_i)>0$, $b_{i,k}:=N(y_{i+m_k}-y_i)>0$. By \eqref{defF} and \eqref{c6:telesc} with $s=\lambda/N$,
\[\lambda\big(F(x)-F(y)\big)=\sum_{i\in J}\frac\lambda N\big(G(a_{i,\cdot})-G(b_{i,\cdot})\big)\le2\sum_{i\in J}\sum_{k=2}^n\Lambda\Big(\frac{x_{i+m_k}-x_i}{y_{i+m_k}-y_i}\Big)+(n-1)|J|\,C_\star\frac{\lambda^2}{N^2}.\]
The map $(i,k)\mapsto(i,i+m_k)$ from $J\times\{2,\dots,n\}$ to $\{(i,j):1\le i<j\le N\}$ is well defined ($m_k\ge1$, $i+m_k\le i+m_n\le N$) and injective ($m_2<\dots<m_n$). Since $\Lambda\ge0$ and $\beta\ge1$, Lemma \ref{lem:convex}(b) gives
\[2\sum_{i\in J}\sum_{k=2}^n\Lambda\Big(\frac{x_{i+m_k}-x_i}{y_{i+m_k}-y_i}\Big)\le2\sum_{1\le i<j\le N}\Lambda\Big(\frac{x_j-x_i}{y_j-y_i}\Big)\le2\beta\sum_{i<j}\Lambda\Big(\frac{x_j-x_i}{y_j-y_i}\Big)\le\cV(x)+\cV(y)-2\cV\big(\tfrac{x+y}2\big).\]

\emph{Step 3 (Pr\'ekopa--Leindler with $\tau=\frac12$).} Define the Borel functions on $\R^N$ (with $F$ extended by $0$ outside $\Sig$)
\[u:=e^{\lambda F-c_\lambda-\cV}\mathbf 1_{\Sig},\qquad v:=e^{-\lambda F-\cV}\mathbf 1_{\Sig},\qquad w:=e^{-\cV}\mathbf 1_{\Sig}.\]
For $x,y\in\Sig$ the midpoint $\frac{x+y}2$ is in $\Sig$ and, by \eqref{c6:concpt},
$u(x)^{1/2}v(y)^{1/2}=\exp\big(\tfrac12[\lambda F(x)-\lambda F(y)-c_\lambda-\cV(x)-\cV(y)]\big)\le e^{-\cV(\frac{x+y}2)}=w\big(\tfrac{x+y}2\big)$;
if $x\notin\Sig$ or $y\notin\Sig$ the left side is $0$. Theorem \ref{thm:PL} gives $\int w\ge(\int u)^{1/2}(\int v)^{1/2}$, i.e., after squaring and dividing by $Z^2$ ($\int w=Z$):
\begin{equation}\label{c6:concprod}
\int e^{\lambda F}\rd\omega\cdot\int e^{-\lambda F}\rd\omega\ \le\ e^{c_\lambda}.
\end{equation}

\emph{Step 4 (Jensen).} Let $\bar F:=\int F\rd\omega$. From $e^t\ge1+t$: $\int e^{-\lambda(F-\bar F)}\rd\omega\ge1-\lambda\int(F-\bar F)\rd\omega=1$, i.e.\ $\int e^{-\lambda F}\rd\omega\ge e^{-\lambda\bar F}$. Inserting this into \eqref{c6:concprod}:
$\int e^{\lambda F}\rd\omega\le e^{c_\lambda}e^{\lambda\bar F}$, which is \eqref{conc}.
\end{proof}

\begin{lemma}[Entropy inequality]\label{lem:DV}
Let $\nu\ll\omega$ be probability measures on a measurable space, $F$ a bounded measurable function, $\lambda>0$. Then
\[\int F\rd\nu\le\frac1\lambda\Big[S(\nu|\omega)+\log\int e^{\lambda F}\rd\omega\Big],\]
where $S(\nu|\omega):=\int\log\frac{\rd\nu}{\rd\omega}\rd\nu\in[0,\infty]$ (Mathlib: \texttt{klDiv}).
\end{lemma}

\noindent\emph{Precisely:} let $\varrho:=\rd\nu/\rd\omega$ be a measurable version of the Radon--Nikodym density, with values in $[0,\infty)$, $\int\varrho\,\rd\omega=1$, and $\mathfrak h(t)=t\log t$, $\mathfrak h(0)=0$ as in the proof of Theorem \ref{thm:LSI}. Then
\[S(\nu|\omega)=\int\mathfrak h(\varrho)\,\rd\omega\in[0,\infty],\]
which is well defined since $\mathfrak h\ge-1/e$, and nonnegative since $\mathfrak h(t)\ge t-1$ and $\int(\varrho-1)\rd\omega=0$. (Equality with $\int\log\varrho\,\rd\nu$: $\nu=\varrho\,\omega$ and $\varrho>0$ $\nu$-a.e. For probability measures with $\nu\ll\omega$ this is Mathlib's \texttt{klDiv}$(\nu,\omega)$, which equals $\int\texttt{llr}\,\rd\nu$ if the log-likelihood ratio is $\nu$-integrable and $+\infty$ otherwise; $\texttt{llr}$ is $\nu$-integrable iff $\mathfrak h(\varrho)$ is $\omega$-integrable iff $\int\mathfrak h(\varrho)\rd\omega<\infty$, the negative part of $\mathfrak h$ being bounded.)

\begin{proof}
If $S(\nu|\omega)=\infty$ there is nothing to prove. Assume $S(\nu|\omega)<\infty$, so $\mathfrak h(\varrho)\in L^1(\omega)$.
\emph{Young's inequality:} for $a\ge0$ and $b\in\R$,
\begin{equation}\label{c6:young}
ab\ \le\ \mathfrak h(a)-a+e^b .
\end{equation}
For $a=0$ this is $0\le e^b$. For $a>0$ apply $\log t\le t-1$ to $t:=e^b/a$: $b-\log a\le e^b/a-1$, and multiply by $a$.
Let $m:=\int e^{\lambda F}\rd\omega\in(0,\infty)$ ($F$ is bounded). Apply \eqref{c6:young} pointwise with $a=\varrho(x)$, $b=\lambda F(x)-\log m$:
\[\varrho\,(\lambda F-\log m)\le\mathfrak h(\varrho)-\varrho+\frac{e^{\lambda F}}m .\]
All four terms are $\omega$-integrable ($\varrho\in L^1(\omega)$, $F$ bounded, $\mathfrak h(\varrho)\in L^1(\omega)$). Integrating against $\omega$, with $\int F\varrho\,\rd\omega=\int F\rd\nu$ and $\int\varrho\,\rd\omega=1$:
$\lambda\int F\rd\nu-\log m\le S(\nu|\omega)-1+1$. Divide by $\lambda>0$.
\end{proof}

\begin{corollary}[Gap observables see only the relative entropy]\label{cor:gap}
In the setting of Theorem \ref{thm:conc}, for every probability measure $\nu\ll\omega$ on $\Sig$:
\begin{equation}\label{gapent}
\Big|\int F\,\rd\nu-\int F\,\rd\omega\Big|\le\frac{2}{N}\sqrt{C_\star(n-1)\,|J|\,S(\nu|\omega)} .
\end{equation}
\end{corollary}

\noindent Here $\sqrt{\,\cdot\,}$ of $+\infty$ is $+\infty$, and $0\cdot\infty=0$ (if $C_\star|J|=0$ then $F\equiv0$, see the proof).

\begin{proof}
Let $\bar F:=\int F\rd\omega$, $K:=C_\star(n-1)|J|/N^2$ and $S:=S(\nu|\omega)$. The functions $\pm(F-\bar F)$ are bounded and Borel on $\Sig$, and both $\nu$ and $\omega$ are carried by $\Sig$.
By Theorem \ref{thm:conc} applied with $\lambda$ and with $-\lambda$, $\log\int e^{\pm\lambda(F-\bar F)}\rd\omega\le K\lambda^2$ for every $\lambda>0$. By Lemma \ref{lem:DV} applied to $\pm(F-\bar F)$,
\begin{equation}\label{c6:gapopt}
\Big|\int F\rd\nu-\bar F\Big|\ \le\ \frac S\lambda+K\lambda\qquad\text{for every }\lambda>0 .
\end{equation}
If $K=0$, then $J=\emptyset$ or $C_\star=0$; in the latter case $A=0$; in both cases $F\equiv0$ and \eqref{gapent} is trivial. Let $K>0$.
If $S=\infty$, the right side of \eqref{gapent} is $+\infty$. If $S=0$, let $\lambda\downarrow0$ in \eqref{c6:gapopt}: the left side of \eqref{gapent} is $0$. If $0<S<\infty$, choose $\lambda:=\sqrt{S/K}$: the right side of \eqref{c6:gapopt} is $2\sqrt{SK}=\frac2N\sqrt{C_\star(n-1)|J|S}$.
\end{proof}

\begin{remark}[Comparison with {[ESYY, Thm 4.1]}; not used in any proof]
With $S(\nu|\omega)\le\frac{R^2}{2N}\int|\nabla\log\frac{\rd\nu}{\rd\omega}|^2\rd\nu=4R^2D_\omega(\sqrt q)$ (Theorem \ref{thm:LSI} with $\kappa\ge N/R^2$; $q=\rd\nu/\rd\omega$, and $D_\omega(f)=\frac1{2N}\int|\nabla f|^2\rd\omega$ is the Dirichlet form in the normalisation of [ESYY]) the bound \eqref{gapent} reads $C\sqrt{|J|R^2D_\omega(\sqrt q)}/N$, which is [ESYY, (4.3)] with $\tau\leftrightarrow R^2$, without the additional term $Ce^{-cN^{\eps}}\sqrt{S_\omega(q)}$ --- and without the auxiliary dynamics $\partial_tq=\widetilde Lq$.
In NP the relative entropy $S(\nu_t|\omega)$ is estimated directly (Theorem \ref{thm:entropy}), and Corollary \ref{cor:gap} is applied as it stands.
\end{remark}

\subsection{What the later sections use}\label{ss:useC6}
\begin{itemize}[leftmargin=1.6em,itemsep=1pt]
\item Theorem \ref{thm:LSI} for $\omega=\omega_R$ (Definition \ref{def:loggas}, Comment (2): $\kappa\ge N/R^2$) and $q=f_r/\psi$ (with $\psi=\rd\omega/\rd\mu$ of Definition \ref{def:omega}) smooth and positive on $\R^N$: used once, in Theorem \ref{thm:entropy}.
\item Corollary \ref{cor:gap} for $\omega=\omega_R$ and $\nu\in\{\nu_t,\mu\}$: used once, in Proposition \ref{prop:gapuniv}. The constant is $C_\star(A,A',\rho)=\max\big(A^2/(2\log\cosh\frac12),\ \frac98(e\rho A')^2\big)$; it enters only through $\sqrt{C_\star(n-1)|J|}$.
\item Theorem \ref{thm:PL}, Lemmas \ref{lem:convex}, \ref{lem:2pt}, \ref{lem:DV} and Theorem \ref{thm:conc} are used only inside this section.
\end{itemize}
The numerical sanity checks for this section (Lemma \ref{lem:2pt} on $1.2\cdot10^8$ random triples for 300 random functions $g$; the pointwise inequalities \eqref{c6:concpt} and Step 1a of Theorem \ref{thm:LSI} on random configurations for $\beta\in\{1,2\}$, $\kappa=0$ and $\kappa>0$, $n\in\{2,3,4\}$; \eqref{LSI} and \eqref{conc} for $N=2$ by quadrature) found no violation.

\section{The entropy bound relative to the local relaxation measure}\label{sec:C7}

\noindent\emph{Standing notation of Sections \ref{sec:C7}--\ref{sec:C8}.}
We use the notation of \S\ref{ss:notation}. The parameters $(\vartheta,\delta_-,\delta_+,C_0)$ are fixed, $\nu\in\mathcal E_N(\vartheta,\delta_-,\delta_+,C_0)$ is an ensemble over $\K$, and $H_0$, $V$, $H_t$ ($t\in[0,\infty]$) are as in \eqref{Ht}, on the probability space $(\widehat\Omega,\widehat{\mathcal F},\P)$ of \S\ref{ss:notation}; $\nu_t=\mathrm{law}(\lambda(H_t))$, $\mu=\lambda_*\bG=\nu_\infty$.
All measures on $\cSig$ are regarded as Borel measures on $\R^N$ carried by $\cSig$ (\S\ref{sec:C4}).
For $R>0$, $W=W_R$, $\omega=\omega_R$ and $\psi=\rd\omega/\rd\mu$ are as in Definition \ref{def:omega}, and
\[
Z_W:=\int e^{-NW}\rd\mu\in(0,1],\qquad \psi=e^{-NW}/Z_W,\qquad -\log\psi=NW+\log Z_W .
\]
$S(\cdot|\cdot)$ is the relative entropy of \S\ref{c3:ss:KL} (equivalently of Lemma \ref{lem:DV}): $S(\nu'|\nu'')=\int\mathfrak h(\rd\nu'/\rd\nu'')\,\rd\nu''\in[0,\infty]$ if $\nu'\ll\nu''$ and $+\infty$ otherwise, with $\mathfrak h(\tau)=\tau\log\tau$, $\mathfrak h(0)=0$; for probability measures this is Mathlib's \texttt{klDiv}.
$D$, $L$ and $Q=7D^2N^{-1}$ are the quantities of \eqref{defL} and Lemma \ref{lem:Q}.

\smallskip\noindent\emph{Moments of $H_0$.} The matrix $H_0$ of $\nu$ has finite moments of all orders: $\E|H_0|^m<\infty$ for every real $m\ge0$.
Indeed, by (RA2) $\nu\in\mathcal E_N(\vartheta',\delta_-,\delta_+,C_0)$ with $\vartheta':=\min(\vartheta,1)$, so by (RA6) $\E|h_{ij}|^q\le(C_{\vartheta'}q)^{q/\vartheta'}\sigma_{ij}^q<\infty$ for all $i\le j$ and $q\ge1$ (and $|h_{ji}|=|h_{ij}|$).
For $m\ge2$, Minkowski's inequality in $L^{m/2}(\P)$ gives $(\E|H_0|^m)^{2/m}=\big\|\sum_{i,j}|h_{ij}|^2\big\|_{L^{m/2}}\le\sum_{i,j}(\E|h_{ij}|^m)^{2/m}<\infty$; for $0\le m<2$, $\E|H_0|^m\le1+\E|H_0|^2$.
Hence Proposition \ref{prop:U9} applies to $H_0$. Moreover $\E|H_0|^2=\E\Tr H_0^2=\sum_{i,j}\E|h_{ij}|^2=\sum_j\sum_i\sigma_{ij}^2=N$ by (A) (Part~I, Definition \ref{def:ass}).

\begin{lemma}[A priori entropy]\label{lem:U8}
For every $s>0$:
\[
S(\nu_s|\mu)\ \le\ S(\mathrm{law}(H_s)\,|\,\bG)\ \le\ \frac d2\log\frac1{1-e^{-s}}+\frac{\beta N^2}{4}e^{-s}.
\]
In particular $S(\nu_s|\mu)\le N^3$ whenever $N\ge1$ and $s\ge N^{-2}$.
\end{lemma}

\begin{proof}
Let $\rho:=\mathrm{law}(H_0)$, a Borel probability measure on $\cH$. By \eqref{Ht}, $\mathrm{law}(H_s)=\rho_s$ in the notation of \S\ref{c3:ss:mehler} with $\mathcal E=\cH$, $\sigma^2=2/(\beta N)$, $g=\bG$.

\emph{First inequality.} $\lambda:\cH\to\R^N$ is Borel (RA4), $\nu_s=\lambda_*\rho_s$ and $\mu=\lambda_*\bG$. The data-processing inequality for the relative entropy under a measurable map (Mathlib \texttt{klDiv\_map\_le}) gives $S(\lambda_*\rho_s|\lambda_*\bG)\le S(\rho_s|\bG)$.

\emph{Second inequality.} Lemma \ref{c3:KL}(b) with $u=s$ and $M_2'=\int|w|^2\rho(\rd w)=\E|H_0|^2=N$ (see above) gives
$S(\rho_s|\bG)\le\frac d2\log\frac1{1-e^{-s}}+\frac{e^{-s}N}{2\sigma^2}$, and $\frac{N}{2\sigma^2}=\frac{\beta N^2}4$ (Remark \ref{c3:rem:U8}).

\emph{In particular.} $d=N+\beta N(N-1)/2\le N+N(N-1)=N^2$ as $\beta\le2$. Since $e^{s}\ge1+s$, $1-e^{-s}\ge1-\frac1{1+s}=\frac s{1+s}$, which is increasing in $s$; so for $s\ge N^{-2}$, $1-e^{-s}\ge\frac{N^{-2}}{1+N^{-2}}=\frac1{N^2+1}\ge\frac1{2N^2}$. With $\beta e^{-s}\le2$,
\[
S(\nu_s|\mu)\le\frac{N^2}2\log(2N^2)+\frac{N^2}2=\frac{N^2}{2}\big(\log2+1+2\log N\big)\le\frac{N^2}2\big(\log 2+1+2(N-1)\big)<N^3,
\]
using $\log N\le N-1$ and $\log2<1$.
\end{proof}

\begin{lemma}[Changing the reference measure]\label{lem:ref}
Let $R>0$ and let $\nu'$ be a Borel probability measure on $\R^N$ with $\nu'\ll\mu$ and $\int|x|^2\nu'(\rd x)<\infty$. Then $\nu'\ll\omega$, $\ \rd\nu'/\rd\omega=(\rd\nu'/\rd\mu)/\psi$, and in $[0,\infty]$
\[
S(\nu'|\omega)=S(\nu'|\mu)+N\int W\rd\nu'+\log Z_W\ \le\ S(\nu'|\mu)+\frac{N}{2R^2}\int|x-\gamma|^2\,\nu'(\rd x).
\]
In particular $S(\mu|\omega)\le\frac{N}{2R^2}\int|x-\gamma|^2\mu(\rd x)<\infty$.
\end{lemma}

\begin{proof}
Let $f:=\rd\nu'/\rd\mu$ (a Borel version with values in $[0,\infty)$). Since $\omega=\psi\mu$ with $\psi>0$ everywhere (Definition \ref{def:omega}, Comment (1)), $\mu=\psi^{-1}\omega$ and $\nu'=f\mu=(f/\psi)\,\omega$; so $\nu'\ll\omega$ with $q:=\rd\nu'/\rd\omega=f/\psi$.
Pointwise, $\mathfrak h(q)\psi=f\log f-f\log\psi=\mathfrak h(f)+f\,(NW+\log Z_W)$ (both sides vanish where $f=0$). Now $0\le NW=\frac N{2R^2}|x-\gamma|^2$ and $\int fNW\rd\mu=N\int W\rd\nu'\le\frac N{2R^2}\int|x-\gamma|^2\rd\nu'\le\frac N{R^2}\big(\int|x|^2\rd\nu'+|\gamma|^2\big)<\infty$, and $\int f\log Z_W\rd\mu=\log Z_W$; so $f(NW+\log Z_W)\in L^1(\mu)$.
The negative part of $\mathfrak h(f)$ is bounded by $1/e$, so $\int\mathfrak h(f)\rd\mu\in(-1/e,\infty]$ is well defined, and the integral of the sum is the sum of the integrals:
\[
S(\nu'|\omega)=\int\mathfrak h(q)\,\rd\omega=\int\mathfrak h(q)\psi\,\rd\mu=\int\mathfrak h(f)\rd\mu+N\int W\rd\nu'+\log Z_W=S(\nu'|\mu)+N\int W\rd\nu'+\log Z_W .
\]
Since $Z_W\le1$, $\log Z_W\le0$, which gives the inequality. For $\nu'=\mu$: $f\equiv1$, $S(\mu|\mu)=0$, and $\int|x|^2\rd\mu=\int|\lambda(H)|^2\bG(\rd H)=\int|H|^2\rd\bG=d\sigma^2<\infty$ (Lemma \ref{lem:inv}(a): $|\lambda(H)|=|H|$; Lemma \ref{lem:M1}(vi) or Lemma \ref{c3:GE}).
\end{proof}

\begin{proposition}[Entropy production]\label{prop:entprod}
Let $H_0$ be a random element of $\cH$ with $\E|H_0|^m<\infty$ for all $m$ (e.g.\ the matrix of $\nu\in\mathcal E_N$, see above), let $s>0$, let $(f_t)_{t>s}$ be as in Proposition \ref{prop:U9}, and let $R>0$. Then $\cS(t):=S(\nu_t|\omega)$ is finite and continuously differentiable on $(s,\infty)$, and for every $t>s$
\begin{equation}\label{entprod}
\cS'(t)=-\frac1{\beta N}\int_{\Sig}\nabla f_t\cdot\nabla\big(\log f_t+NW\big)\,\rd\mu\ \le\ -\frac1{2\beta N}I_t+\frac{N}{2\beta R^4}\,\E|\lambda(H_t)-\gamma|^2,
\end{equation}
where $q_t:=f_t/\psi$ is a version of $\rd\nu_t/\rd\omega$ and
\[
I_t:=\int_{\Sig}\Big|\nabla\log\frac{f_t}{\psi}\Big|^2f_t\,\rd\mu=\int_{\Sig}\frac{|\nabla q_t|^2}{q_t}\,\rd\omega\ \in[0,\infty).
\]
Equivalently, $\cS'(t)=-\frac1{\beta N}\big[\int\frac{|\nabla q_t|^2}{q_t}\rd\omega+\int\nabla\log\psi\cdot\nabla q_t\,\rd\omega\big]$.
\end{proposition}

\begin{proof}
Throughout, $f(t,x):=f_t(x)$; by Proposition \ref{prop:U9}(ii), $f$ is $C^\infty$ and positive on $(s,\infty)\times\R^N$, symmetric in $x$, and $\nu_t=f_t\mu$. Since $\mu=\mu^*$ (Theorem \ref{thm:JD}) and $\mu^*(\R^N\setminus\Sig)=0$, all $\mu$-integrals below may be taken over $\Sig$.
By Definition \ref{def:omega}, Comment (4): $\nabla W=(x-\gamma)/R^2$, $|\nabla W|=|x-\gamma|/R^2\le(|x|+|\gamma|)/R^2$, and $0\le NW\le\frac N{R^2}(|x|^2+|\gamma|^2)$.

\emph{Step 0 (finiteness; the formula for $\cS$).} Let $t>s$. By Lemma \ref{lem:inv}(a), $\int|x|^2\rd\nu_t=\E|\lambda(H_t)|^2=\E|H_t|^2<\infty$ (Proposition \ref{prop:U9}(i) applied with $s$ replaced by $t$). Lemma \ref{lem:ref} with $\nu'=\nu_t$ and $f=f_t$ gives
\begin{equation}\label{c7:Sj}
\cS(t)=\int j(t,x)\,\mu(\rd x)+\log Z_W,\qquad j(t,x):=f_t(x)\big(\log f_t(x)+NW(x)\big),
\end{equation}
provided $j(t,\cdot)\in L^1(\mu)$; this holds because, for $t$ in any compact interval $[t_1,t_2]\subset(s,\infty)$, Proposition \ref{prop:U9}(iv) gives a constant $C<\infty$ with $|\log f_t(x)|\le C(1+|x|^2)$, so that
\[
|j(t,x)|\le\mathfrak m(x)\,\big(C(1+|x|^2)+\tfrac N{R^2}(|x|^2+|\gamma|^2)\big)\le C_1\,\mathfrak m(x)(1+|x|)^2,\qquad C_1:=C+\tfrac N{R^2}(1+|\gamma|^2),
\]
with $\mathfrak m$ the function of Proposition \ref{prop:U9}(iv) for $[t_1,t_2]$, and $\int\mathfrak m(1+|x|)^2\rd\mu<\infty$. In particular $\cS(t)<\infty$.

\emph{Step 1 (differentiation in $t$).} Fix $s<t_1<t_2$, with $C$, $C_1$, $\mathfrak m$ as in Step 0. For every $x$, $t\mapsto j(t,x)$ is differentiable on $(t_1,t_2)$ with
\[
\partial_tj(t,x)=\partial_tf_t(x)\,\big(\log f_t(x)+1+NW(x)\big),\qquad |\partial_tj(t,x)|\le\mathfrak m(x)\big(C_1(1+|x|)^2+1\big)\le(C_1+1)\,\mathfrak m(x)(1+|x|)^2,
\]
and $x\mapsto j(t,x)$, $x\mapsto\partial_tj(t,x)$ are continuous, hence Borel. The dominating function $(C_1+1)\mathfrak m(1+|x|)^2$ is $\mu$-integrable and independent of $t\in[t_1,t_2]$.
By differentiation under the integral sign (Mathlib \texttt{hasDerivAt\_integral\_of\_dominated\_loc\_of\_deriv\_le}), $t\mapsto\int j(t,\cdot)\rd\mu$ is differentiable on $(t_1,t_2)$ with derivative $\int\partial_tj(t,\cdot)\rd\mu$.
Moreover $t\mapsto\partial_tj(t,x)$ is continuous for each $x$ ($f\in C^\infty$, $f>0$), so by dominated convergence $t\mapsto\int\partial_tj(t,\cdot)\rd\mu$ is continuous on $(t_1,t_2)$. As $s<t_1<t_2$ are arbitrary, $\cS\in C^1((s,\infty))$ and
\[
\cS'(t)=\int_{\Sig}\partial_tf_t\,k_t\,\rd\mu,\qquad k_t:=\log f_t+1+NW\in C^\infty(\R^N).
\]

\emph{Step 2 (the equation).} Every $x\in\Sig$ has pairwise distinct coordinates, so $\partial_tf_t(x)=(Lf_t)(x)$ by Proposition \ref{prop:U9}(iii). Hence, with $\mu=\mu^*$,
$\cS'(t)=\int_{\Sig}k_t\,Lf_t\,\rd\mu^*$.

\emph{Step 3 (integration by parts).} Apply Corollary \ref{cor:ibp} with $\varphi:=f_t\in C^2(\R^N)$ and $k:=k_t\in C^1(\R^N)$. Its three hypotheses hold, with $\mathfrak m$ for an interval $[t_1,t_2]\ni t$:
\begin{itemize}[leftmargin=1.6em,itemsep=1pt]
\item $|k_t|\,|\nabla f_t|\le(C_1+1)(1+|x|)^2\,\mathfrak m(x)\in L^1(\mu^*)$;
\item $k_t\,Lf_t=k_t\,\partial_tf_t$ on $\Sig$, and $|k_t\,\partial_tf_t|\le(C_1+1)(1+|x|)^2\mathfrak m(x)\in L^1(\mu^*)$;
\item $\nabla k_t\cdot\nabla f_t=\frac{|\nabla f_t|^2}{f_t}+N\nabla W\cdot\nabla f_t$, where $\int|\nabla f_t|^2/f_t\,\rd\mu<\infty$ (Proposition \ref{prop:U9}(iv)) and $|N\nabla W\cdot\nabla f_t|\le\frac N{R^2}(|x|+|\gamma|)\mathfrak m(x)\in L^1(\mu^*)$.
\end{itemize}
Hence \eqref{ibp} gives $\cS'(t)=-\frac1{\beta N}\int_{\Sig}\nabla k_t\cdot\nabla f_t\,\rd\mu^*$, which is the equality in \eqref{entprod} since $\nabla k_t=\nabla(\log f_t+NW)$.

\emph{Step 4 (completing the square).} Put $\mathbf a:=\nabla\log f_t+N\nabla W$. Since $\log\psi=-NW-\log Z_W$, $\mathbf a=\nabla\log(f_t/\psi)=\nabla\log q_t$. Pointwise, using $\nabla f_t=f_t\nabla\log f_t=f_t(\mathbf a-N\nabla W)$ and $N\nabla W\cdot\mathbf a\le\frac12|\mathbf a|^2+\frac12N^2|\nabla W|^2$,
\[
\nabla f_t\cdot\nabla(\log f_t+NW)=f_t\,(\mathbf a-N\nabla W)\cdot\mathbf a\ \ge\ f_t\Big(\frac12|\mathbf a|^2-\frac12N^2|\nabla W|^2\Big).
\]
\emph{Finiteness of $I_t$:} $f_t|\mathbf a|^2\le2\frac{|\nabla f_t|^2}{f_t}+2N^2f_t|\nabla W|^2$, and $\int f_t|\nabla W|^2\rd\mu=R^{-4}\int|x-\gamma|^2\rd\nu_t<\infty$ (Step 0); so $I_t=\int f_t|\mathbf a|^2\rd\mu<\infty$.
Integrating the pointwise inequality against $\mu$ (all terms integrable) and inserting it in the equality of Step 3:
\[
\cS'(t)\le-\frac1{\beta N}\Big[\frac12I_t-\frac{N^2}{2R^4}\int|x-\gamma|^2f_t\,\rd\mu\Big]=-\frac{I_t}{2\beta N}+\frac{N}{2\beta R^4}\E|\lambda(H_t)-\gamma|^2,
\]
since $\int|x-\gamma|^2f_t\rd\mu=\int|x-\gamma|^2\rd\nu_t=\E|\lambda(H_t)-\gamma|^2$.

\emph{The two forms of $I_t$ and of $\cS'$.} $q_t=f_t/\psi$ is $C^\infty$ and positive, $\nu_t=q_t\,\omega$ (Lemma \ref{lem:ref}), $\nabla\log q_t=\mathbf a$, so $\frac{|\nabla q_t|^2}{q_t}=q_t|\mathbf a|^2$ and $\int\frac{|\nabla q_t|^2}{q_t}\rd\omega=\int q_t\psi|\mathbf a|^2\rd\mu=\int f_t|\mathbf a|^2\rd\mu=I_t$.
Finally, with $f_t=q_t\psi$: $\nabla f_t\cdot\nabla\log q_t=\psi\big(\frac{|\nabla q_t|^2}{q_t}+\nabla\log\psi\cdot\nabla q_t\big)$, and integrating against $\mu$ (i.e.\ the bracket against $\omega=\psi\mu$) gives the last formula.
\end{proof}

\begin{theorem}[Entropy bound]\label{thm:entropy}
Let $N_4:=N_3$ with $N_3=N_3(\vartheta,\delta_\pm,C_0)$ from Lemma \ref{lem:Q}. For all $N\ge N_4$, $\K\in\{\R,\C\}$, $\nu\in\mathcal E_N(\vartheta,\delta_-,\delta_+,C_0)$ and all $t>0$ with $Nt\ge(\log N)^2$, with the choice
\begin{equation}\label{Rchoice}
R^2:=\frac{t}{8\beta\log N},\qquad \omega:=\omega_R,
\end{equation}
we have, with $D$ from \eqref{defL},
\begin{equation}\label{entbound}
S(\nu_t|\omega)\le\frac1N+\frac{32\beta D^2\log N}{t},\qquad S(\mu|\omega)\le\frac{28\beta D^2\log N}{t}.
\end{equation}
\end{theorem}
\noindent The hypothesis $Nt\ge(\log N)^2$ is used in the proof only through $t/2\ge N^{-2}$ (step (4) below); so \eqref{entbound}, and with it Proposition~\ref{prop:gapuniv}, holds for every $t\ge2N^{-2}$.
For $t\ll1/N$ the bound is true but weaker than trivial bounds; the hypothesis is kept because it is the one needed in \S\ref{sec:C8}.

\begin{proof}
Note $N\ge N_3\ge16$ (Lemma \ref{lem:Q} and \S\ref{sec:C2}: $N_3\ge N_6\ge N_*\ge16$), so $\log N>0$ and $R$ is well defined. Put $Q:=7D^2N^{-1}$; by Lemma \ref{lem:Q}, $\E|\lambda(H_r)-\gamma|^2\le Q$ for every $r\in[0,\infty]$, in particular $\int|x-\gamma|^2\rd\mu\le Q$.
Note the identity
\begin{equation}\label{c7:NQR}
\frac{NQ}{2R^2}=\frac{7D^2}{2R^2}=\frac{28\beta D^2\log N}{t}.
\end{equation}

\emph{The bound on $S(\mu|\omega)$} is Lemma \ref{lem:ref} with \eqref{c7:NQR}.

\emph{The bound on $S(\nu_t|\omega)$.} Apply Proposition \ref{prop:U9} to $H_0$ with $s:=t/4$ (the moment hypothesis holds, see the beginning of this section), and Proposition \ref{prop:entprod} with this $s$ and $R$ from \eqref{Rchoice}. So $\cS(r)=S(\nu_r|\omega)$ is finite and $C^1$ on $(t/4,\infty)\supset[t/2,t]$.

(1) \emph{LSI.} By Definition \ref{def:omega}, Comment (3) (equivalently Definition \ref{def:loggas}, Comment (2)), $\omega$ is a convex log-gas with $\kappa=N(\frac\beta2+\frac1{R^2})\ge\frac N{R^2}>0$. For $r>t/4$, $q_r=f_r/\psi\in C^\infty(\R^N)$ is positive and $\int_{\Sig}q_r\rd\omega=\int f_r\rd\mu=\nu_r(\R^N)=1$. Theorem \ref{thm:LSI} gives
\[
\cS(r)=\int_{\Sig}\mathfrak h(q_r)\,\rd\omega=\Ent_\omega(q_r)\le\frac1{2\kappa}\int_{\Sig}\frac{|\nabla q_r|^2}{q_r}\rd\omega\le\frac{R^2}{2N}I_r .
\]
(The first equality: $q_r$ is a version of $\rd\nu_r/\rd\omega$, Lemma \ref{lem:ref}, and $\omega$ is carried by $\Sig$.)

(2) \emph{Differential inequality.} By \eqref{entprod}, (1) in the form $-I_r\le-\frac{2N}{R^2}\cS(r)$, and $\E|\lambda(H_r)-\gamma|^2\le Q$: for every $r\in[t/2,t]$,
\[
\cS'(r)\le-a\,\cS(r)+b',\qquad a:=\frac1{\beta R^2},\quad b':=\frac{NQ}{2\beta R^4},\quad \frac{b'}a=\frac{NQ}{2R^2}.
\]

(3) \emph{Gr\"onwall.} The function $h(r):=e^{ar}\big(\cS(r)-b'/a\big)$ is $C^1$ on $(t/4,\infty)$ with $h'(r)=e^{ar}\big(\cS'(r)+a\cS(r)-b'\big)\le0$ for $r\in[t/2,t]$. By the mean value theorem $h(t)\le h(t/2)$, i.e.
\[
\cS(t)\le\frac{b'}a+e^{-at/2}\Big(\cS(t/2)-\frac{b'}a\Big)\le e^{-at/2}\,\cS(t/2)+\frac{NQ}{2R^2}.
\]

(4) \emph{Initial value.} Since $Nt\ge(\log N)^2$ and $N\ge16$: $t/2\ge(\log N)^2/(2N)\ge N^{-2}$ (as $N(\log N)^2\ge2$). By Lemma \ref{lem:ref} (with $\nu'=\nu_{t/2}\ll\mu$, Proposition \ref{prop:U9}(i), and $\int|x|^2\rd\nu_{t/2}<\infty$), Lemma \ref{lem:U8} and Lemma \ref{lem:Q}:
\[
\cS(t/2)\le S(\nu_{t/2}|\mu)+\frac N{2R^2}\E|\lambda(H_{t/2})-\gamma|^2\le N^3+\frac{NQ}{2R^2}.
\]

(5) \emph{Arithmetic.} $\frac{at}2=\frac{t}{2\beta R^2}=\frac{t\cdot8\beta\log N}{2\beta t}=4\log N$, so $e^{-at/2}=N^{-4}$. By (3), (4) and \eqref{c7:NQR},
\[
\cS(t)\le N^{-4}\Big(N^3+\frac{NQ}{2R^2}\Big)+\frac{NQ}{2R^2}=\frac1N+(1+N^{-4})\frac{28\beta D^2\log N}{t}\le\frac1N+\frac{32\beta D^2\log N}{t},
\]
since $28(1+N^{-4})\le32$.
\end{proof}

\begin{remark}[Where the dynamics is used]\label{rem:dyn}
Only in \eqref{entprod}. Its two ingredients are (a) the pointwise identity $\partial_tf_t=Lf_t$ on $\Sig$ for the smooth symmetric density of Proposition \ref{prop:U9}, and (b) Corollary \ref{cor:ibp}.
The cross term $N\nabla W\cdot\nabla f_t$ is what distinguishes $\omega$ from $\mu$; it is the reason why a purely static argument (matrix de Bruijn identity plus projection of the Fisher information onto eigenvalues, which only yields $\int|\nabla f_t|^2/f_t\,\rd\mu\le CN^3/t$) does not suffice.
The only quantitative inputs of Theorem \ref{thm:entropy} are Lemma \ref{lem:U8} (through $S(\nu_{t/2}|\mu)\le N^3$), Lemma \ref{lem:Q} (rigidity, through $Q$), Theorem \ref{thm:LSI} (through $\kappa\ge N/R^2$) and the choice \eqref{Rchoice}; all constants in Proposition \ref{prop:U9}(iv) are used only qualitatively (Remark \ref{rem:M4}).
\end{remark}
\section[From gap observables to correlation functions; proof of Theorems \ref*{thm:sharp} and \ref*{thm:23}]{From gap observables to correlation functions; proof of Theorems \texorpdfstring{\ref{thm:sharp}}{12.1} and \ref{thm:23}}\label{sec:C8}

The notation of \S\ref{sec:C7} remains in force.

\subsection{Gap universality}

\begin{proposition}[Universality of gap observables]\label{prop:gapuniv}
Let $N\ge N_4$, $\K\in\{\R,\C\}$, $\nu\in\mathcal E_N(\vartheta,\delta_-,\delta_+,C_0)$, and $t>0$ with $Nt\ge(\log N)^2$. Let $n\ge2$, integers $1\le m_2<\dots<m_n$, $J\subset\{1,\dots,N-m_n\}$, $A,A'\ge0$, $\rho>0$, and let $G$ and $F$ be as in Theorem \ref{thm:conc} and \eqref{defF}; extend $F$ to $\R^N$ by $F:=0$ on $\R^N\setminus\Sig$. Then
\begin{equation}\label{gapuniv}
\big|\E F(\lambda(H_t))-\E F(\lambda(V))\big|\le\frac{4}{N}\sqrt{C_\star(n-1)|J|}\;\Big(\frac1N+\frac{32\beta D^2\log N}{t}\Big)^{1/2},
\end{equation}
with $C_\star=C_\star(A,A',\rho)$ from Lemma \ref{lem:2pt}.
\end{proposition}

\begin{proof}
$G$ is continuous on $(0,\infty)^{n-1}$ (comment after Theorem \ref{thm:conc}), so $F$ is continuous on the open set $\Sig$; the extended $F$ is Borel on $\R^N$ and $|F|\le A|J|/N$. Hence both expectations exist, and $\E F(\lambda(H_t))=\int F\rd\nu_t$, $\E F(\lambda(V))=\int F\rd\mu$.
By Theorem \ref{thm:JD}, $\nu_t(\partial\Sig)=\mu(\partial\Sig)=0$; as $\nu_t,\mu$ are carried by $\cSig$, they are probability measures carried by $\Sig$, and the integrals do not depend on the extension of $F$ off $\Sig$.
Let $R$, $\omega=\omega_R$ be as in \eqref{Rchoice}. By Theorem \ref{thm:JD}, $\nu_t\ll\mu$, and $\mu=\psi^{-1}\omega\ll\omega$ (Definition \ref{def:omega}, Comment (1)); so $\nu_t\ll\omega$ and $\mu\ll\omega$.
Corollary \ref{cor:gap} (with the convex log-gas $\omega$, Definition \ref{def:omega}, Comment (3)), applied with $\nu=\nu_t$ and with $\nu=\mu$, and the triangle inequality give
\[
\Big|\int F\rd\nu_t-\int F\rd\mu\Big|\le\frac2N\sqrt{C_\star(n-1)|J|}\,\Big(\sqrt{S(\nu_t|\omega)}+\sqrt{S(\mu|\omega)}\Big).
\]
By Theorem \ref{thm:entropy} both entropies are at most $\frac1N+\frac{32\beta D^2\log N}t$ (as $28\le32$). This is \eqref{gapuniv}.
\end{proof}

\noindent This is the NP version of [ESYY, Theorem 4.1]: in the original's variables it reads $C\sqrt{|J|Q/(N\tau)}$ up to $(\log N)^{1/2}$, with $\tau=t$. Note that \eqref{gapuniv} holds for \emph{every} $t$ with $Nt\ge(\log N)^2$; no time average and no auxiliary dynamics are needed.

\subsection{Deterministic reduction}\label{ss:reduction}

\noindent\emph{Fixed data.} Throughout this subsection $n\ge1$, $\rho_O>0$ and $0<c\le2$ are fixed, $\varrho_*:=\varrho_*(c)=\frac1{2\pi}\sqrt{c-c^2/16}$, and
\[
\ell:=\max\Big(1,\frac{\rho_O}{\varrho_*}\Big)\ \ge1 .
\]
$O:\R^n\to\R$ is Lipschitz with support in $[-\rho_O,\rho_O]^n$ (so $O$ is continuous and bounded). $E\in[-2+c,2-c]$, $b\in(0,c/2)$ and $\varrho:=\rhosc(E)$. Since $|E|\le2-c\le2-c/4$, (RA5) gives
\begin{equation}\label{c8:rho}
\varrho_*\le\varrho\le\frac1\pi,\qquad\frac{\rho_O}{\varrho}\le\frac{\rho_O}{\varrho_*}\le\ell .
\end{equation}
$J:=\{i:\gamma_i\in[E-b,E+b]\}$ and $\mathcal B:=\{i:\operatorname{dist}(\gamma_i,\{E-b,E+b\})\le\xi_0+\ell/N\}$, $\xi_0=D/(N\varrho_*)$, are the deterministic index sets of Lemma \ref{lem:countC}(c).
Note $[E-b,E+b]\subset[-2+\frac c2,2-\frac c2]\subset[-2,2]$.

\noindent\emph{Symmetrisation.} $O^{s}(\alpha):=\frac1{n!}\sum_{\pi\in S_n}O(\alpha_{\pi(1)},\dots,\alpha_{\pi(n)})$. Since a permutation of the coordinates is an isometry of $(\R^n,|\cdot|)$ preserving $[-\rho_O,\rho_O]^n$: $\|O^s\|_\infty\le\|O\|_\infty$, $\Lip(O^s)\le\Lip(O)$, and $O^s=0$ outside $[-\rho_O,\rho_O]^n$.

\noindent\emph{Index tuples.} $\mathfrak S_n:=\{m=(m_2,\dots,m_n)\in\mathbb Z^{n-1}:1\le m_2<\dots<m_n\}$ ($\mathfrak S_1=\{\emptyset\}$); always $m_1:=0$, and $m_n:=0$ for $n=1$. For $x\in\R^N$, $E'\in\R$, $1\le i\le N$, $m\in\mathfrak S_n$, $u\in\R^{n-1}$:
\begin{align*}
Y_{i,m}(E',x)&:=\begin{cases}O^{s}\big(N\varrho(x_i-E'),N\varrho(x_{i+m_2}-E'),\dots,N\varrho(x_{i+m_n}-E')\big)&\text{if }i+m_n\le N,\\ 0&\text{else;}\end{cases}\\
G_\varrho(u)&:=\int_\R O^{s}\big(\varrho y,\varrho(y+u_2),\dots,\varrho(y+u_n)\big)\,\rd y\qquad(u=(u_2,\dots,u_n));\\
F_m(x)&:=\frac1N\sum_{i\in J,\ i+m_n\le N}G_\varrho\big(N(x_{i+m_2}-x_i),\dots,N(x_{i+m_n}-x_i)\big).
\end{align*}
For $n=1$, $G_\varrho=\int O(\varrho y)\rd y$ is a number and $F_\emptyset=|J|G_\varrho/N$ is constant.

\begin{lemma}[Algebraic identities and trivial bounds]\label{lem:alg}
For all $x\in\cSig$ and $E'\in\R$:
\begin{enumerate}[label=(\alph*),itemsep=1pt]
\item $\cX^O_{E',\varrho}(x)=n!\sum_{i=1}^N\sum_{m\in\mathfrak S_n}Y_{i,m}(E',x)$ (only the finitely many $m$ with $m_n\le N-1$ contribute);
\item if $i+m_n\le N$: $E'\mapsto Y_{i,m}(E',x)$ is continuous; $Y_{i,m}(E',x)=0$ unless $|x_{i+m_k}-E'|\le\ell/N$ for all $k=1,\dots,n$; $\int_\R|Y_{i,m}(E',x)|\rd E'\le2\ell\|O\|_\infty/N$; and
$\int_\R Y_{i,m}(E',x)\,\rd E'=\frac1NG_\varrho\big(N(x_{i+m_2}-x_i),\dots,N(x_{i+m_n}-x_i)\big)$;
\item $G_\varrho:\R^{n-1}\to\R$ is well defined, $|G_\varrho|\le A:=2\ell\|O\|_\infty$, $G_\varrho(u)=0$ unless $|u_k|\le2\ell$ for all $k$, and $|G_\varrho(u)-G_\varrho(u')|\le A'|u_k-u_k'|$ whenever $u,u'$ differ only in the $k$-th coordinate, with $A':=2\rho_O\Lip(O)$; in particular $G_\varrho$ is continuous.
Consequently, for $n\ge2$, the restriction of $G_\varrho$ to $(0,\infty)^{n-1}$ satisfies the hypotheses of Theorem \ref{thm:conc} with $(A,A',\rho)=(2\ell\|O\|_\infty,\,2\rho_O\Lip(O),\,2\ell)$, and the constant of Lemma \ref{lem:2pt} satisfies
\begin{equation}\label{c8:Cstar}
\sqrt{C_\star}\le C_\sharp\,\big(\|O\|_\infty+\Lip(O)\big),\qquad C_\sharp:=\ell\Big(\sqrt{2/c_0}+3\sqrt2\,e\,\rho_O\Big),\quad c_0=\log\cosh\tfrac12 ;
\end{equation}
\item $|\cX^O_{E',\varrho}(x)|\le N^n\|O\|_\infty$ for all $x\in\R^N$; $F_m$ is continuous on $\R^N$ and $|F_m|\le A|J|/N$; $F_m\equiv0$ if $m_n\ge N$; and for $n\ge2$ and $m_n\le N-1$, the restriction of $F_m$ to $\Sig$ is the function $F$ of \eqref{defF} built with $G:=G_\varrho|_{(0,\infty)^{n-1}}$, the tuple $m$ and the index set $J_m:=\{i\in J:i\le N-m_n\}\subset\{1,\dots,N-m_n\}$, $|J_m|\le|J|$.
\end{enumerate}
\end{lemma}

\begin{proof}
(a) Put $\alpha_j:=N\varrho(x_j-E')$. An ordered $n$-tuple $(i_1,\dots,i_n)$ of pairwise distinct indices is the same as a pair $(\mathcal J,\pi)$ of an $n$-subset $\mathcal J=\{j_1<\dots<j_n\}$ of $\{1,\dots,N\}$ and a permutation $\pi\in S_n$, via $i_k=j_{\pi(k)}$. Hence
$\cX^O_{E',\varrho}(x)=\sum_{\mathcal J}\sum_{\pi}O(\alpha_{j_{\pi(1)}},\dots,\alpha_{j_{\pi(n)}})=n!\sum_{\mathcal J}O^s(\alpha_{j_1},\dots,\alpha_{j_n})$.
The $n$-subsets $\mathcal J$ are in bijection with the pairs $(i,m)$, $m\in\mathfrak S_n$, $1\le i$, $i+m_n\le N$, via $i=j_1$, $m_k=j_k-j_1$; and $O^s(\alpha_{j_1},\dots,\alpha_{j_n})=Y_{i,m}(E',x)$. The terms with $i+m_n>N$ are $0$ by definition.

(b) Continuity: $O^s$ is continuous. If $Y_{i,m}(E',x)\ne0$ then each argument lies in $[-\rho_O,\rho_O]$, i.e.\ $|x_{i+m_k}-E'|\le\rho_O/(N\varrho)\le\ell/N$ by \eqref{c8:rho}. Thus $Y_{i,m}(\cdot,x)$ vanishes outside $[x_i-\ell/N,x_i+\ell/N]$ and $|Y_{i,m}|\le\|O\|_\infty$, which gives the $L^1$ bound.
For the integral substitute $y:=N(x_i-E')$ ($\rd E'=\rd y/N$ up to orientation): $N\varrho(x_{i+m_k}-E')=\varrho\big(y+N(x_{i+m_k}-x_i)\big)$, which gives the formula with $u_k=N(x_{i+m_k}-x_i)$.

(c) For fixed $u$, $y\mapsto O^s(\varrho y,\dots)$ is continuous and vanishes unless $|\varrho y|\le\rho_O$, i.e.\ outside an interval of length $2\rho_O/\varrho\le2\ell$; so $G_\varrho(u)$ is well defined and $|G_\varrho(u)|\le2\ell\|O\|_\infty$.
If $G_\varrho(u)\ne0$, some $y$ has $|\varrho y|\le\rho_O$ and $|\varrho(y+u_k)|\le\rho_O$ for all $k$, so $|u_k|\le2\rho_O/\varrho\le2\ell$.
If $u,u'$ differ only in the $k$-th coordinate, then for every $y$ the two integrands differ only in one argument, by $\varrho|u_k-u_k'|$, and both vanish unless $|\varrho y|\le\rho_O$; hence
$|G_\varrho(u)-G_\varrho(u')|\le\frac{2\rho_O}{\varrho}\cdot\Lip(O)\,\varrho|u_k-u_k'|=2\rho_O\Lip(O)|u_k-u_k'|$.
Changing one coordinate at a time, $|G_\varrho(u)-G_\varrho(u')|\le A'\sum_k|u_k-u'_k|$, so $G_\varrho$ is continuous, hence Borel.
For $n\ge2$ the restriction to $(0,\infty)^{n-1}$ therefore satisfies: Borel, $|G|\le A$, $G(u)=0$ unless $u\in(0,2\ell]^{n-1}$, coordinatewise $A'$-Lipschitz — the hypotheses of Theorem \ref{thm:conc} with $\rho=2\ell$.
By Lemma \ref{lem:2pt}, $C_\star=\max\big(\frac{A^2}{2c_0},\frac98(e\rho A')^2\big)$, so $\sqrt{C_\star}\le\frac{A}{\sqrt{2c_0}}+\frac{3}{2\sqrt2}e\rho A'=\ell\sqrt{2/c_0}\,\|O\|_\infty+3\sqrt2\,e\,\ell\rho_O\Lip(O)$, which gives \eqref{c8:Cstar}.

(d) The number of ordered $n$-tuples of distinct indices is $N!/(N-n)!\le N^n$. $F_m$ is a finite sum of continuous functions, each bounded by $A/N$, with at most $|J|$ terms. If $m_n\ge N$ the sum in $F_m$ is empty. If $m_n\le N-1$, the summation set is $J_m$, and for $x\in\Sig$, $i\in J_m$ all arguments $N(x_{i+m_k}-x_i)$ are positive; so $F_m|_{\Sig}$ is \eqref{defF}.
\end{proof}

\begin{lemma}[Localisation on the good set]\label{lem:local}
Assume \eqref{c2:counthyp}, i.e.\ $(D+\ell)/(N\varrho_*)\le c/8$, with $D$ from \eqref{defL}; let $x\in\Xi$ (see \eqref{good}) and $K_\ell:=2\ell/\pi+2D$ (Lemma \ref{lem:countC}(a)). Then:
\begin{enumerate}[label=(\alph*),itemsep=1pt]
\item $Y_{i,m}(E',x)=0$ for all $E'\in[-3,3]$, all $i$ and all $m\in\mathfrak S_n$ with $m_n\ge K_\ell$; for $i\in J$, the $i$-th summand of $F_m(x)$ vanishes whenever $m_n\ge K_\ell$.
Moreover, for every bounded Borel $P:\R^n\to\R$ vanishing outside $[-\rho_O,\rho_O]^n$ and every $E'\in[-3,3]$: $|\cX^P_{E',\varrho}(x)|\le K_\ell^n\|P\|_\infty$.
\item The number of $m\in\mathfrak S_n$ with $m_n<K_\ell$ is at most $K_\ell^{n-1}$, and
\[
\Big|\int_{E-b}^{E+b}\cX^O_{E',\varrho}(x)\,\rd E'-n!\sum_{m\in\mathfrak S_n,\,m_n<K_\ell}F_m(x)\Big|\le n!\,|\mathcal B|\,K_\ell^{n-1}\,\frac{2\ell\|O\|_\infty}{N}.
\]
\end{enumerate}
\end{lemma}

\begin{proof}
By \eqref{c2:counthyp}, $\xi_0=D/(N\varrho_*)\le c/8$ and $\ell/N\le\varrho_*c/8<c/8$ (as $\varrho_*\le\frac{\sqrt2}{2\pi}<1$).

(a) If $Y_{i,m}(E',x)\ne0$, then by Lemma \ref{lem:alg}(b) $x_i$ and $x_{i+m_n}$ lie in $[E'-\ell/N,E'+\ell/N]$, hence so do the $m_n+1$ points $x_i\le x_{i+1}\le\dots\le x_{i+m_n}$ ($x\in\cSig$). Lemma \ref{lem:countC}(a) with $e=E'\in[-3,3]$ gives $m_n+1\le K_\ell$; so $m_n\ge K_\ell$ is impossible.
Let $i\in J$ and suppose the $i$-th summand of $F_m(x)$ is nonzero. By Lemma \ref{lem:alg}(c), $0\le N(x_{i+m_n}-x_i)\le2\ell$, so the $m_n+1$ points $x_i,\dots,x_{i+m_n}$ lie in $[e-\ell/N,e+\ell/N]$ with $e:=x_i+\ell/N$.
Since $\gamma_i\in[E-b,E+b]\subset[-2+\frac38c,2-\frac38c]$, Lemma \ref{lem:countC}(b) gives $|x_i-\gamma_i|\le\xi_0\le c/8$, so $|x_i|\le2-\frac c2+\frac c8\le2$ and $|e|\le2+\frac c8\le3$. Lemma \ref{lem:countC}(a) gives $m_n+1\le K_\ell$ again.
For the last claim: a nonzero term of $\cX^P_{E',\varrho}(x)$ (see \eqref{defX}) requires $|N\varrho(x_{i_k}-E')|\le\rho_O$, i.e.\ $|x_{i_k}-E'|\le\ell/N$, for $k=1,\dots,n$. By Lemma \ref{lem:countC}(a) the set $\{j:|x_j-E'|\le\ell/N\}$ has at most $K_\ell$ elements, so there are at most $K_\ell^n$ such ordered tuples, each term bounded by $\|P\|_\infty$.

(b) \emph{Counting.} Let $M$ be the largest integer $<K_\ell$ ($K_\ell\ge2D\ge4$). The $m\in\mathfrak S_n$ with $m_n<K_\ell$ correspond to the $(n-1)$-subsets of $\{1,\dots,M\}$; there are $\binom M{n-1}\le M^{n-1}\le K_\ell^{n-1}$ of them (for $n=1$: one).

\emph{Reduction.} By Lemma \ref{lem:alg}(a) and Fubini (finite sum), $\int_{E-b}^{E+b}\cX^O_{E',\varrho}(x)\rd E'=n!\sum_{i,m}\int_{E-b}^{E+b}Y_{i,m}(E',x)\rd E'$; by (a) (note $[E-b,E+b]\subset[-3,3]$) only $m$ with $m_n<K_\ell$ contribute.
By Lemma \ref{lem:alg}(b), $F_m(x)=\sum_{i\in J}\int_\R Y_{i,m}(E',x)\rd E'$ (terms with $i+m_n>N$ vanish on both sides). Hence the quantity to be bounded is $n!\big|\sum_{i}\sum_{m:\,m_n<K_\ell}\Delta_{i,m}\big|$ with
\[
\Delta_{i,m}:=\int_{E-b}^{E+b}Y_{i,m}(E',x)\,\rd E'-\mathbf 1_{i\in J}\int_\R Y_{i,m}(E',x)\,\rd E' .
\]
\emph{Claim: $\Delta_{i,m}=0$ unless $i\in\mathcal B$.}
Case $i\notin J$, $\Delta_{i,m}\ne0$: then $Y_{i,m}(E',x)\ne0$ for some $E'\in[E-b,E+b]$, so $|x_i-E'|\le\ell/N$ and $x_i\in[E-b-\frac\ell N,E+b+\frac\ell N]\subset[-2+\frac c2-\frac c8,2-\frac c2+\frac c8]=[-2+\frac38c,2-\frac38c]$. By Lemma \ref{lem:countC}(b), $|x_i-\gamma_i|\le\xi_0$, so $\gamma_i\in[E-b-\xi_0-\frac\ell N,E+b+\xi_0+\frac\ell N]$; as $\gamma_i\notin[E-b,E+b]$, $\gamma_i$ is within $\xi_0+\ell/N$ of $E-b$ or of $E+b$: $i\in\mathcal B$.
Case $i\in J\setminus\mathcal B$: $\gamma_i\in[E-b,E+b]$ is in the bulk, so $|x_i-\gamma_i|\le\xi_0$ (Lemma \ref{lem:countC}(b)). If $Y_{i,m}(E',x)\ne0$ then $|E'-\gamma_i|\le|E'-x_i|+|x_i-\gamma_i|\le\ell/N+\xi_0<\operatorname{dist}(\gamma_i,\{E-b,E+b\})$; a point $E'\notin[E-b,E+b]$ would be separated from $\gamma_i\in[E-b,E+b]$ by an endpoint, hence at distance $\ge\operatorname{dist}(\gamma_i,\{E-b,E+b\})$ from $\gamma_i$. So the support of $Y_{i,m}(\cdot,x)$ lies in $[E-b,E+b]$ and $\Delta_{i,m}=0$.
\emph{Bound.} For $i\in\mathcal B$: if $i\in J$, $\Delta_{i,m}=-\int_{\R\setminus[E-b,E+b]}Y_{i,m}\rd E'$; if $i\notin J$, $\Delta_{i,m}=\int_{E-b}^{E+b}Y_{i,m}\rd E'$. In both cases $|\Delta_{i,m}|\le\int_\R|Y_{i,m}|\rd E'\le2\ell\|O\|_\infty/N$ (Lemma \ref{lem:alg}(b)). Summing over $i\in\mathcal B$ and the at most $K_\ell^{n-1}$ values of $m$ gives the claim.
\end{proof}

\subsection{Proof of Theorem \ref{thm:sharp}}\label{ss:proofsharp}

\noindent\emph{Constants.} Let $(\vartheta,\delta_\pm,C_0)$, $c$, $n$, $\rho_O$ be given, and $\varrho_*$, $\ell$ as in \S\ref{ss:reduction}. Let $A_0$ be the constant of \eqref{defL}, and $N_*$, $N_K$ ($K=n+2$), $N_3$ the thresholds of \eqref{c2:Nstar}, Lemma \ref{lem:goodC}, Lemma \ref{lem:Q}; $N_4=N_3$. Let $N_c$ be the least integer $\ge16$ with
$\big(N_c^{1/5}+2+\ell\big)/(N_c\varrho_*)\le c/8$ (it exists since the left side, as a function of $N_c\ge1$, decreases to $0$). Put
\[
N_{\rm q}:=\max\big(N_3,\,N_{n+2},\,N_c,\,n^2\big),\qquad A_1:=nA_0+1 .
\]
Then $N_{\rm q}\ge16$, and $N_{\rm q}$, $A_1$ depend only on $(c,n,\rho_O,\vartheta,\delta_\pm,C_0)$. With $C_\sharp$ from \eqref{c8:Cstar} and $C_{\mathcal B}:=\frac{4}{\pi\varrho_*}+\frac{4\ell}\pi+1$ put
\[
C_1'':=8(\ell+2)^n+32\,n!\sqrt{2n}\,(\ell+2)^{n-1}C_\sharp+4\,n!\,\ell\,C_{\mathcal B}(\ell+2)^{n-1}+4+8\,n!\,\ell(\ell+2)^{n-1},\qquad C_1':=3^nC_1'' ;
\]
$C_1'$ depends only on $(c,n,\rho_O)$.

\noindent\emph{Consequences of $N\ge N_{\rm q}$.} Let $N\ge\max(N_{\rm q},n)=N_{\rm q}$.
\begin{itemize}[leftmargin=1.6em,itemsep=1pt]
\item[(i)] $N\ge N_*$, so \eqref{c2:Dsmall} holds: $D\le N^{1/5}+2$. Hence $(D+\ell)/(N\varrho_*)\le(N^{1/5}+2+\ell)/(N\varrho_*)\le c/8$ (monotonicity in $N$), i.e.\ \eqref{c2:counthyp} holds and Lemma \ref{lem:local} applies. Also $D\le N$ and $D\ge2$.
\item[(ii)] $K_\ell=\frac{2\ell}\pi+2D\le(\ell+2)D$ (as $D\ge1$), and by Lemma \ref{lem:countC}(c), $|\mathcal B|\le\frac4\pi\big(\frac D{\varrho_*}+\ell\big)+2\le C_{\mathcal B}D$ (as $D\ge2$).
\item[(iii)] $C_{N,n}=\prod_{k=0}^{n-1}\frac N{N-k}\le2$: indeed $\prod_{k<n}(1-\frac kN)\ge1-\sum_{k<n}\frac kN=1-\frac{n(n-1)}{2N}\ge\frac12$ for $N\ge n^2$.
\item[(iv)] By Lemma \ref{lem:goodC} with $K=n+2$: $\P(\lambda(H_t)\notin\Xi)\le N^{-n-2}$ for every $t\in[0,\infty]$, and $\mu(\Xi^c)\le N^{-n-2}$.
\item[(v)] $D^n(\log N)^{1/2}\le3^ne^{A_1(\log\log N)^2}$: $D=e^{A_0(\log\log N)^2}+2\le3e^{A_0(\log\log N)^2}$, and $(\log N)^{1/2}=e^{\frac12\log\log N}\le e^{(\log\log N)^2}$ because $\log\log N\ge\log\log16>\frac12$.
\end{itemize}

\begin{proof}[Proof of Theorem \ref{thm:sharp}]
Let $N\ge N_{\rm q}$, $\K$, $\nu\in\mathcal E_N(\vartheta,\delta_-,\delta_+,C_0)$, $O$, $t>0$ with $Nt\ge(\log N)^2$, $E\in[-2+c,2-c]$, $b\in(0,c/2)$, and $\varrho=\rhosc(E)$. Write $\|O\|:=\|O\|_\infty$ and $\Theta$ for the left side of \eqref{eq:abstrthm2}.

\emph{Step 1 (a deterministic functional).} $(E',x)\mapsto\cX^O_{E',\varrho}(x)$ is continuous on $\R\times\R^N$ ($O$ is continuous; finite sum), so
\[
\mathcal Z(x):=\int_{E-b}^{E+b}\cX^O_{E',\varrho}(x)\,\rd E'
\]
is continuous in $x$, hence Borel, and $|\mathcal Z|\le2bN^n\|O\|$ (Lemma \ref{lem:alg}(d)). The map $(E',\hat\omega)\mapsto\cX^O_{E',\varrho}(\lambda(H_t(\hat\omega)))$ is jointly measurable and bounded, so by Fubini
$\int_{E-b}^{E+b}\E\,\cX^O_{E',\varrho}(\lambda(H_t))\rd E'=\E\,\mathcal Z(\lambda(H_t))$, and likewise for $V$. Therefore
\[
\Theta=\frac{C_{N,n}}{2b}\big|\E\,\mathcal Z(\lambda(H_t))-\E\,\mathcal Z(\lambda(V))\big|\le\frac1b\big|\E\,\mathcal Z(\lambda(H_t))-\E\,\mathcal Z(\lambda(V))\big| .
\]
For $x\in\Xi$, Lemma \ref{lem:local}(a) with $P=O$ (all $E'\in[E-b,E+b]\subset[-3,3]$) gives $|\mathcal Z(x)|\le2bK_\ell^n\|O\|$.

\emph{Step 2 (case $b<1/N$).} Split each expectation according to $\lambda(\cdot)\in\Xi$ or not and use (iv):
$|\E\,\mathcal Z(\lambda(H_t))|\le2bK_\ell^n\|O\|+2bN^n\|O\|N^{-n-2}$, and the same for $V$. Hence, by (ii) and $b^{-1}N^{-1}>1$,
\[
\Theta\le4\|O\|\big(K_\ell^n+N^{-2}\big)\le8(\ell+2)^nD^n\|O\|\le8(\ell+2)^nD^n\|O\|\,b^{-1}N^{-1}.
\]

\emph{Step 3 (case $b\ge1/N$: decomposition).} By Lemma \ref{lem:countC}(c), $|J|\le\frac{2Nb}\pi+1\le2Nb$ (as $Nb\ge1$). Define, for $x\in\cSig$,
\[
\mathcal F(x):=n!\sum_{m\in\mathfrak S_n,\ m_n<K_\ell}F_m(x),\qquad\mathrm{Err}(x):=\mathcal Z(x)-\mathcal F(x).
\]
By Lemma \ref{lem:local}(b) and (ii), $|\mathrm{Err}(x)|\le n!\,C_{\mathcal B}D\,(\ell+2)^{n-1}D^{n-1}\frac{2\ell\|O\|}N$ for $x\in\Xi$; and for all $x$, by Lemma \ref{lem:alg}(d) and Lemma \ref{lem:local}(b) (counting), $|\mathrm{Err}(x)|\le2bN^n\|O\|+n!K_\ell^{n-1}\frac{A|J|}N\le2bN^n\|O\|+4\,n!\,\ell\,b\,K_\ell^{n-1}\|O\|$.
Since $\E\,\mathcal Z(\lambda(\cdot))=\E\,\mathcal F(\lambda(\cdot))+\E\,\mathrm{Err}(\lambda(\cdot))$ for $H_t$ and for $V$,
\[
\Theta\le(T1)+(T2)+(T3),\qquad (T1):=\frac{n!}b\sum_{m:\,m_n<K_\ell}\big|\E F_m(\lambda(H_t))-\E F_m(\lambda(V))\big|,
\]
\[
(T2):=\frac1b\sum_{X\in\{H_t,V\}}\E\big[|\mathrm{Err}(\lambda(X))|\mathbf 1_{\Xi}(\lambda(X))\big],\qquad(T3):=\frac1b\sum_{X\in\{H_t,V\}}\E\big[|\mathrm{Err}(\lambda(X))|\mathbf 1_{\Xi^c}(\lambda(X))\big].
\]

\emph{Step 4 (T3: the bad event).} By (iv), (ii) and $D\le N$:
\[
\begin{aligned}
(T3)&\le\frac2bN^{-n-2}\big(2bN^n+4\,n!\,\ell\,b\,K_\ell^{n-1}\big)\|O\|\le\big(4N^{-2}+8\,n!\,\ell(\ell+2)^{n-1}N^{n-1}N^{-n-2}\big)\|O\|\\
&\le\big(4+8\,n!\,\ell(\ell+2)^{n-1}\big)\|O\|\,b^{-1}N^{-1},
\end{aligned}
\]
using $N^{-2}\le N^{-1}\le b^{-1}N^{-1}$ (as $b<c/2\le1$).

\emph{Step 5 (T2: boundary indices).} By Step 3, $(T2)\le\frac2b\cdot n!\,C_{\mathcal B}(\ell+2)^{n-1}D^n\frac{2\ell\|O\|}N=4\,n!\,\ell\,C_{\mathcal B}(\ell+2)^{n-1}D^n\|O\|\,b^{-1}N^{-1}$.

\emph{Step 6 (T1: gap universality).} If $n=1$, $F_\emptyset$ is constant and $(T1)=0$. Let $n\ge2$ and $m\in\mathfrak S_n$ with $m_n<K_\ell$. If $m_n\ge N$, $F_m\equiv0$. Otherwise, by Lemma \ref{lem:alg}(c),(d), $F_m|_{\Sig}$ is a gap observable \eqref{defF} with index set $J_m\subset\{1,\dots,N-m_n\}$, $|J_m|\le|J|\le2Nb$, and $G=G_\varrho$ satisfying the hypotheses of Theorem \ref{thm:conc}; $F_m$ is continuous on $\R^N$ and $\nu_t(\partial\Sig)=\mu(\partial\Sig)=0$ (Theorem \ref{thm:JD}), so $\E F_m(\lambda(\cdot))$ equals the expectation of the extension used in Proposition \ref{prop:gapuniv}. By Proposition \ref{prop:gapuniv} ($N\ge N_4$, $Nt\ge(\log N)^2$), $\beta\le2$ and $\sqrt{\alpha+\alpha'}\le\sqrt\alpha+\sqrt{\alpha'}$:
\[
\begin{aligned}
\big|\E F_m(\lambda(H_t))-\E F_m(\lambda(V))\big|&\le\frac4N\sqrt{2C_\star(n-1)Nb}\Big(N^{-1/2}+8D(\log N)^{1/2}t^{-1/2}\Big)\\&=4\sqrt{2C_\star(n-1)}\,b^{1/2}\Big(N^{-1}+8D(\log N)^{1/2}(Nt)^{-1/2}\Big).
\end{aligned}
\]
There are at most $K_\ell^{n-1}\le(\ell+2)^{n-1}D^{n-1}$ such $m$ (Lemma \ref{lem:local}(b), (ii)). With \eqref{c8:Cstar}, $b^{-1/2}N^{-1}\le b^{-1}N^{-1}$ and $1\le D(\log N)^{1/2}$:
\[
\begin{aligned}
(T1)&\le4\,n!\sqrt{2(n-1)}\,(\ell+2)^{n-1}D^{n-1}C_\sharp\big(\|O\|+\Lip(O)\big)\cdot8D(\log N)^{1/2}\Big[b^{-1}N^{-1}+b^{-1/2}(Nt)^{-1/2}\Big]\\
&\le32\,n!\sqrt{2n}\,(\ell+2)^{n-1}C_\sharp\big(\|O\|+\Lip(O)\big)D^n(\log N)^{1/2}\Big[b^{-1}N^{-1}+b^{-1/2}(Nt)^{-1/2}\Big].
\end{aligned}
\]

\emph{Step 7 (conclusion).} In both cases, by Steps 2, 4, 5, 6 and $(\log N)^{1/2}\ge1$, $D\ge1$:
\[
\Theta\le C_1''\big(\|O\|_\infty+\Lip(O)\big)\,D^n(\log N)^{1/2}\Big[b^{-1}N^{-1}+b^{-1/2}(Nt)^{-1/2}\Big],
\]
and by (v) this is at most $C_1'(\|O\|_\infty+\Lip(O))e^{A_1(\log\log N)^2}[b^{-1}N^{-1}+b^{-1/2}(Nt)^{-1/2}]$, which is \eqref{sharp}.
\end{proof}

\begin{remark}[Dependence of the constants]\label{rem:Odep}
In the proof, $N_{\rm q}$ depends only on $(c,n,\rho_O,\vartheta,\delta_\pm,C_0)$, $A_1=nA_0+1$ only on $(n,\vartheta,\delta_\pm,C_0)$, and $C_1'$ only on $(c,n,\rho_O)$; none depends on $\|O\|_\infty$, $\Lip(O)$, $E$, $b$, $t$, $\K$, the ensemble $\nu$ or $N$. The test function enters only through $\rho_O$ (via $\ell$ and the support of $G_\varrho$) and linearly through $\|O\|_\infty$ (Steps 2, 4, 5) and $\sqrt{C_\star}\le C_\sharp(\|O\|_\infty+\Lip(O))$ (Step 6).
The Lipschitz constant is used only in Lemma \ref{lem:alg}(c), where only the coordinatewise Lipschitz constants of $O$ enter; these are bounded by $\Lip(O)$ (Euclidean norm). This is obligation (O1) of Remark \ref{c1:oblig}.
\end{remark}

\subsection{Proof of Theorem \ref{thm:23}}\label{ss:proof23}

This is item (O2) of Remark \ref{c1:oblig}. Let $(\eps',\delta,c,n,\rho_O,\vartheta,\delta_\pm,C_0)$ be given, and let $A_1,N_{\rm q},C_1'$ be the constants of Theorem \ref{thm:sharp} (proved in \S\ref{ss:proofsharp}). Since $(\log\log N)^2/\log N\to0$ and $(\log N)^2/N^{\delta+\eps'}\to0$ as $N\to\infty$, there is a least integer
$N_2\ge\max(N_{\rm q},n,3)$, depending only on $(\eps',\delta,A_1,N_{\rm q},n)$, such that for all $N\ge N_2$
\begin{equation}\label{c8:N2}
e^{A_1(\log\log N)^2}\le N^{\eps'}\qquad\text{and}\qquad N^{\delta+\eps'}\ge(\log N)^2 .
\end{equation}
Put $C:=\max\big(C_1',\,2(c/2)^{1/2}N_2^{\,n+\delta/2}\big)$; it depends only on $(\eps',\delta,c,n,\rho_O,\vartheta,\delta_\pm,C_0)$.

\begin{proof}[Proof of Theorem \ref{thm:23}]
Let $N\ge\max(n,3)$, $\K$, $\nu$, $O$, $E$, $b$ and $t\ge N^{-1+\delta+\eps'}$ be as in Theorem \ref{thm:23}, and let $\Theta$ be the left side of \eqref{eq:abstrthm2}.

\emph{Case $N\ge N_2$.} By \eqref{c8:N2}, $Nt\ge N^{\delta+\eps'}\ge(\log N)^2$, so Theorem \ref{thm:sharp} applies ($N\ge N_{\rm q}$). Moreover $(Nt)^{-1/2}\le N^{-(\delta+\eps')/2}\le N^{-\delta/2}$, $N^{-1}\le N^{-1+\eps'}$ and $e^{A_1(\log\log N)^2}\le N^{\eps'}\le N^{2\eps'}$. Hence
$\Theta\le C_1'(\|O\|_\infty+\Lip(O))N^{2\eps'}[b^{-1}N^{-1+\eps'}+b^{-1/2}N^{-\delta/2}]$, and $C_1'\le C$.

\emph{Case $\max(n,3)\le N<N_2$.} By Lemma \ref{lem:alg}(d) (which needs only that $O$ is bounded) $|\cX^O_{E',\varrho}|\le\frac{N!}{(N-n)!}\|O\|_\infty$, so
$\Theta\le\frac1{2b}\int_{E-b}^{E+b}C_{N,n}\cdot2\frac{N!}{(N-n)!}\|O\|_\infty\,\rd E'=2N^n\|O\|_\infty\le2N_2^n\|O\|_\infty$.
On the other hand $b<c/2$, $N<N_2$ and $N^{2\eps'}\ge1$ give
$C(\|O\|_\infty+\Lip(O))N^{2\eps'}\big[b^{-1}N^{-1+\eps'}+b^{-1/2}N^{-\delta/2}\big]\ge C\|O\|_\infty(c/2)^{-1/2}N_2^{-\delta/2}\ge2N_2^n\|O\|_\infty$
by the choice of $C$. So \eqref{eq:abstrthm2} holds in both cases.
\end{proof}

\noindent\emph{Room.} For $N\ge N_2$ the proof shows the stronger bound $C_1'(\|O\|_\infty+\Lip(O))\big[N^{\eps'}b^{-1}N^{-1}+N^{\eps'/2}b^{-1/2}N^{-\delta/2}\big]$: the first term of \eqref{eq:abstrthm2} is beaten by $N^{2\eps'}$, the second by $N^{3\eps'/2}$.

\begin{corollary}[H\"older test functions]\label{cor:cont}
Let $0<\alpha\le1$ and let $O:\R^n\to\R$ vanish outside $[-\rho_O,\rho_O]^n$ and satisfy $|O(u)-O(v)|\le[O]_\alpha|u-v|^\alpha$ for all $u,v$. In the setting of Theorem \ref{thm:sharp} ($N\ge N_{\rm q}$, $Nt\ge(\log N)^2$, $E\in[-2+c,2-c]$, $b\in(0,c/2)$, $\nu\in\mathcal E_N$), with
\[
\mathcal R:=e^{A_1(\log\log N)^2}\Big[b^{-1}N^{-1}+b^{-1/2}(Nt)^{-1/2}\Big]
\]
(the right side of \eqref{sharp} without the factor $C_1'(\|O\|_\infty+\Lip(O))$):
\[
\text{l.h.s.\ of \eqref{eq:abstrthm2}}\ \le\ C_{\rm H}\,\big(\|O\|_\infty+[O]_\alpha\big)\,D^n\,\min(1,\mathcal R)^{\alpha},\qquad C_{\rm H}:=16(\ell+2)^n+C_1',
\]
$C_{\rm H}$ depending only on $(c,n,\rho_O)$. In particular, for $t\ge N^{-1+\delta+\eps'}$ and $N\ge N_2$ (see \eqref{c8:N2}) the left side of \eqref{eq:abstrthm2} is at most $3^nC_{\rm H}(\|O\|_\infty+[O]_\alpha)N^{2\eps'}\big[b^{-1}N^{-1}+b^{-1/2}N^{-\delta/2}\big]^\alpha$ (using $D^n\le3^nN^{\eps'}$, which follows from (v) of \S\ref{ss:proofsharp} and \eqref{c8:N2}).
\end{corollary}

\begin{proof}
For a bounded Borel $P:\R^n\to\R$ vanishing outside $[-\rho_O,\rho_O]^n$ let $\mathscr L(P)$ denote the expression inside the absolute value on the left of \eqref{eq:abstrthm2} (well defined as in Step 1 of \S\ref{ss:proofsharp} if $P$ is continuous; we only use continuous $P$). $\Lambda$ is linear in $P$.

\emph{Trivial bound.} By Lemma \ref{lem:local}(a) on $\Xi$ (applicable by (i)), $|\cX^P_{E',\varrho}|\le N^n\|P\|_\infty$ off $\Xi$, (iii) and (iv):
$|\mathscr L(P)|\le\frac{C_{N,n}}{2b}\cdot2b\cdot2\big(K_\ell^n+N^nN^{-n-2}\big)\|P\|_\infty\le4(K_\ell^n+1)\|P\|_\infty\le8(\ell+2)^nD^n\|P\|_\infty$.

\emph{Lipschitz approximation.} For $r\in(0,1]$ let $L_r:=[O]_\alpha r^{\alpha-1}$, $\eps_r:=[O]_\alpha r^\alpha$,
\[
\widetilde O_r(u):=\inf_{v\in\R^n}\big(O(v)+L_r|u-v|\big),\qquad O_r:=\phi_{\eps_r}\circ\widetilde O_r,\qquad\phi_\eps(\tau):=\operatorname{sign}(\tau)\,(|\tau|-\eps)_+ .
\]
$O$ is bounded ($|O(u)|=|O(u)-O(v)|\le[O]_\alpha|u-v|^\alpha$ for $v$ outside the cube, and $O$ is continuous on the compact cube), so $\widetilde O_r$ is finite, $\inf O\le\widetilde O_r\le O$, and $\widetilde O_r$ is $L_r$-Lipschitz (an infimum of $L_r$-Lipschitz functions which is finite).
Moreover $O(u)-\widetilde O_r(u)=\sup_v\big(O(u)-O(v)-L_r|u-v|\big)\le\sup_{\tau\ge0}[O]_\alpha(\tau^\alpha-r^{\alpha-1}\tau)\le\eps_r$, since $\tau^\alpha\le r^\alpha$ for $\tau\le r$ and $\tau^\alpha-r^{\alpha-1}\tau=\tau^\alpha(1-(\tau/r)^{1-\alpha})\le0$ for $\tau\ge r$.
$\phi_\eps$ is $1$-Lipschitz, $|\phi_\eps(\tau)-\tau|\le\eps$, $|\phi_\eps(\tau)|\le|\tau|$ and $\phi_\eps=0$ on $[-\eps,\eps]$. Hence: $O_r$ is $L_r$-Lipschitz; $|O_r-O|\le|O_r-\widetilde O_r|+|\widetilde O_r-O|\le2\eps_r$; outside $[-\rho_O,\rho_O]^n$, $O=0$ so $\widetilde O_r\in[-\eps_r,0]$ and $O_r=0$; and $\|O_r\|_\infty\le\|O\|_\infty$ (if $\widetilde O_r(u)\ge0$ then $|O_r(u)|\le\widetilde O_r(u)\le O(u)$; if $\widetilde O_r(u)<0$ then $|O_r(u)|=(|\widetilde O_r(u)|-\eps_r)_+\le|O(u)|$ because $\widetilde O_r(u)\ge O(u)-\eps_r$).

\emph{Conclusion.} By linearity, the trivial bound for $P=O-O_r$ (continuous, vanishing outside the cube, $\|P\|_\infty\le2\eps_r$) and Theorem \ref{thm:sharp} for $O_r$:
\[
|\mathscr L(O)|\le16(\ell+2)^nD^n[O]_\alpha r^\alpha+C_1'\big(\|O\|_\infty+[O]_\alpha r^{\alpha-1}\big)\mathcal R .
\]
If $\mathcal R\ge1$, the trivial bound for $P=O$ gives $|\mathscr L(O)|\le8(\ell+2)^nD^n\|O\|_\infty\le C_{\rm H}\|O\|_\infty D^n\min(1,\mathcal R)^\alpha$. If $\mathcal R<1$, take $r=\mathcal R$: then $\mathcal R\le\mathcal R^\alpha$ and $r^{\alpha-1}\mathcal R=\mathcal R^\alpha$, so $|\mathscr L(O)|\le(16(\ell+2)^n+C_1')(\|O\|_\infty+[O]_\alpha)D^n\mathcal R^\alpha$ (using $D\ge1$).
For the last statement: for $N\ge N_2$ and $t\ge N^{-1+\delta+\eps'}$, $\mathcal R\le N^{\eps'}[b^{-1}N^{-1}+b^{-1/2}N^{-\delta/2}]$ as in \S\ref{ss:proof23}, $\min(1,\mathcal R)^\alpha\le\mathcal R^\alpha$ and $N^{\alpha\eps'}D^n\le3^nN^{2\eps'}$.
\end{proof}

\begin{remark}[Merely continuous $O$: not covered]\label{rem:contO}
For $O\in C_c(\R^n)$ without a modulus of continuity the argument above gives l.h.s.\ $\le C D^n\|O-O'\|_\infty+C(O')\mathcal R$ for every Lipschitz $O'$; because $D^n=e^{nA_0(\log\log N)^2}(1+o(1))\to\infty$, this does \emph{not} even imply convergence to $0$ for arbitrary continuous $O$.
A uniform bound $\E\,\#\{n\text{-tuples of eigenvalues in a window of size }\ell/N\}\le C$ (a Wegner-type estimate, not available in NP for $n\ge2$) would give the qualitative statement, but still no rate $C(O)\times$power.
The original states (2.36) for continuous $O$ with $C=C(O)$; its source [ESYY, Thm 4.1] needs $G$ smooth. Moreover, by the uniform boundedness principle, ``for every continuous $O$ there is $C(O)$'' would be equivalent to a bound with $C(O)=C'\|O\|_\infty$, i.e.\ a total-variation bound at polynomial rate, which neither the original nor [ESYY] proves.
NP therefore states Theorem \ref{thm:23} for Lipschitz $O$ (item PD-21 of Appendix~\ref{app:diff}); H\"older $O$ is covered by Corollary \ref{cor:cont}.
\end{remark}
\section{Dependency graph and constants of Part~III}\label{sec:C9}

This section records the logical structure of Part~III and its constants. It contains no new mathematics.

\subsection{Dependency graph}
{\footnotesize
\begin{alltt}
 Part I: Thm 2.2 (2.25), sup form of (2.26) [RA1]; Lem 2.5(ii),(v) [RA2,RA3]; Lem 3.13 [RA4];
         Lem 3.4 [RA5]; Lem 4.1(a) [RA6]
     |
     v
 Lem \ref*{c2:transfer} (transfer), Lem \ref*{c2:gausstail} (Gaussian tails), Lem \ref*{c3:GE} (G on H, entries)
     '--> Lem \ref*{lem:U11} (H~_t, V~ in one class E_N(theta~, delta~, C~_0), all t in [0,inf])
              '--> Lem \ref*{c2:Xi} (Xi Borel) --> Lem \ref*{lem:goodC} (good event, all t) --+--> Lem \ref*{lem:Q} (Q <= 7 D^2/N)   [Assumption III]
                                                                         '--> used in Sec. \ref*{sec:C8}
              Lem \ref*{lem:countC} (deterministic: counting, bulk localisation, index sets)   [Assumption IV]
 [INFRA] Def \ref*{c3:def:gauss}, Lem \ref*{c3:gauss} (Gaussians on a Euclidean space), Lem \ref*{c3:kernel}-\ref*{c3:smooth} (Mehler kernel, explicit domination)
     '--> Lem \ref*{lem:M1}-\ref*{lem:M3} (Mehler semigroup: symmetry, smoothing, d_u = A, log-bounds, domination, mixing)
     '--> Lem \ref*{c3:KL} (Gaussian relative entropy) --------------------------------------------.
 Lem \ref*{lem:inv} (invariant functions), Lem \ref*{lem:givens} (Givens: (A Phi)(diag x) = (L phi)(x))                |
     |--> Prop \ref*{prop:U9} (smooth symmetric f_t, nu_t = f_t mu, d_t f = L f, domination)          |
     '--> [INFRA] Lem \ref*{lem:div}, Cor \ref*{cor:ibp} (divergence thm / IBP on Sigma_N)                           |
              '--> Thm \ref*{thm:JD} (mu = mu*: joint density of GbetaE; nu_t << mu, nu_t(walls) = 0)    |
                       '--> Def \ref*{def:omega} (omega_R, psi = d omega / d mu)                             |
 [INFRA] Thm \ref*{thm:PL} (Prekopa-Leindler)                                                             |
     |--> Lem \ref*{lem:convex} (convexity) --> Thm \ref*{thm:LSI} (LSI on Sigma_N, constant 1/(2 kappa))               |
     '--> Lem \ref*{lem:2pt} (two-point) --> Thm \ref*{thm:conc} (gap concentration) --+--> Cor \ref*{cor:gap} (gap vs entropy)|
                                    Lem \ref*{lem:DV} (entropy inequality) -'                             |
 Lem \ref*{lem:U8} (a priori entropy) <-- Lem \ref*{c3:KL} + klDiv_map_le <-----------------------------------'
 Lem \ref*{lem:ref} (reference change) <-- Def \ref*{def:omega}
 Prop \ref*{prop:U9} + Thm \ref*{thm:JD} + Cor \ref*{cor:ibp} --> Prop \ref*{prop:entprod} (entropy production)
 Prop \ref*{prop:entprod} + Thm \ref*{thm:LSI} + Lem \ref*{lem:U8} + Lem \ref*{lem:ref} + Lem \ref*{lem:Q} --> Thm \ref*{thm:entropy} (entropy bound)
 Thm \ref*{thm:entropy} + Cor \ref*{cor:gap} + Thm \ref*{thm:JD} --> Prop \ref*{prop:gapuniv} (gap universality)
 Lem \ref*{lem:alg}, \ref*{lem:local} (deterministic reduction; use Lem \ref*{lem:countC}) + Prop \ref*{prop:gapuniv} + Lem \ref*{lem:goodC}
     --> Thm \ref*{thm:sharp} (sharp form, constant C_1'(||O|| + Lip O)) --> Thm 2.3 (Sec. \ref*{ss:proof23})
                                                          '--> Cor \ref*{cor:cont} (Hoelder O)
\end{alltt}
}
The graph is acyclic. Theorem \ref{thm:JD} does not use Proposition \ref{prop:U9}. The three \INFRA\ blocks (Mehler; Pr\'ekopa--Leindler; integration by parts on $\Sig$) are independent of each other and of Part~I.
Only Theorem~\ref{thm:22} of the random-matrix part of NP is used (at each fixed time, for $\widetilde H_t$); Theorem~\ref{thm:21} and Part~II are not used.

\subsection{Exponent and constant checks}
The table collects the bounds that enter the error estimate, each with the place where it is proved in \S\S\ref{sec:C7}--\ref{sec:C8}.

\begin{center}\footnotesize
\renewcommand{\arraystretch}{1.25}
\begin{tabular}{|p{3.1cm}|p{9.1cm}|p{2.7cm}|}
\hline
quantity & bound & source\\\hline
$D$ & $e^{A_0(\log\log N)^2}+2\le N^{1/5}+2$ ($N\ge N_*$) & \eqref{defL}, \eqref{c2:Dsmall}\\
$\P(\lambda(H_t)\notin\Xi)$, $\mu(\Xi^c)$ & $\le N^{-n-2}$ for $N\ge N_{n+2}$, every $t\in[0,\infty]$ & Lemma \ref{lem:goodC}\\
$Q=\sup_r\E|\lambda(H_r)-\gamma|^2$ & $\le7D^2N^{-1}$ & Lemma \ref{lem:Q}\\
$R^2$ & $t/(8\beta\log N)$; then $\frac{NQ}{2R^2}=\frac{28\beta D^2\log N}t$ & \eqref{Rchoice}, \eqref{c7:NQR}\\
a priori entropy & $S(\nu_{t/2}|\mu)\le N^3$ ($t/2\ge N^{-2}$, all $N\ge1$) & Lemma \ref{lem:U8}\\
initial entropy & $S(\nu_{t/2}|\omega)\le N^3+\frac{NQ}{2R^2}$ & Lemma \ref{lem:ref}\\
LSI & $S(\nu_r|\omega)\le\frac{R^2}{2N}I_r$ ($\kappa\ge N/R^2$) & Theorem \ref{thm:LSI}\\
decay on $[t/2,t]$ & $e^{-t/(2\beta R^2)}=N^{-4}$; $\ N^{-4}\cdot N^3=N^{-1}$ & Theorem \ref{thm:entropy}\\
$S(\nu_t|\omega)$, $S(\mu|\omega)$ & $\le N^{-1}+32\beta D^2(\log N)/t$ (resp.\ $28\beta D^2(\log N)/t$) & \eqref{entbound}\\
one gap observable & $\frac4N\sqrt{C_\star(n-1)|J|}\sqrt{S}\le4\sqrt{2C_\star(n-1)}\,b^{1/2}[N^{-1}+8D(\log N)^{\frac12}(Nt)^{-\frac12}]$ ($|J|\le2Nb$) & Prop.\ \ref{prop:gapuniv}\\
$\sqrt{C_\star}$ & $\le C_\sharp(\|O\|_\infty+\Lip O)$, $C_\sharp=\ell(\sqrt{2/c_0}+3\sqrt2e\rho_O)$ & \eqref{c8:Cstar}\\
number of $m$'s & $\le K_\ell^{n-1}\le(\ell+2)^{n-1}D^{n-1}$ (original: $M^{n-1}$, $M^n=N^{\eps'}$) & Lemma \ref{lem:local}\\
(T1) & $\le32\,n!\sqrt{2n}(\ell+2)^{n-1}C_\sharp(\|O\|_\infty+\Lip O)D^n(\log N)^{\frac12}[b^{-1}N^{-1}+b^{-\frac12}(Nt)^{-\frac12}]$ & \S\ref{ss:proofsharp}, Step 6\\
(T2) boundary indices & $\le4\,n!\,\ell\,C_{\mathcal B}(\ell+2)^{n-1}D^n\|O\|_\infty b^{-1}N^{-1}$ (original: $b^{-1}\xi$, $\xi=2N^{-1+A}$) & Step 5, Lemma \ref{lem:local}\\
(T3) bad event & $\le(4+8\,n!\,\ell(\ell+2)^{n-1})\|O\|_\infty N^{-2}$ & Step 4, Lemma \ref{lem:goodC}\\
case $b<1/N$ & $\le8(\ell+2)^nD^n\|O\|_\infty\le8(\ell+2)^nD^n\|O\|_\infty b^{-1}N^{-1}$ & Step 2\\\hline
total & $C_1''(\|O\|_\infty+\Lip O)D^n(\log N)^{\frac12}[b^{-1}N^{-1}+b^{-\frac12}(Nt)^{-\frac12}]$, $\ D^n(\log N)^{\frac12}\le3^ne^{(nA_0+1)(\log\log N)^2}$ & \eqref{sharp}\\
$t\ge N^{-1+\delta+\eps'}$, $N\ge N_2$ & $(Nt)^{-1/2}\le N^{-(\delta+\eps')/2}$; $e^{A_1(\log\log N)^2}\le N^{\eps'}$ & \S\ref{ss:proof23}\\
$\max(n,3)\le N<N_2$ & l.h.s.\ $\le2N_2^n\|O\|_\infty$; $C\ge2(c/2)^{1/2}N_2^{n+\delta/2}$ & \S\ref{ss:proof23}\\\hline
\end{tabular}
\end{center}
\noindent Comparison with \eqref{eq:abstrthm2}: the first term is better by $N^{2\eps'}$ ($N^{\eps'}b^{-1}N^{-1}$ vs.\ $N^{2\eps'}b^{-1}N^{-1+\eps'}$), the second by $N^{3\eps'/2}$.
Example: $n=2$, $\eps'=\frac1{20}$, $\delta=\frac12$, $b=\frac14$, $t=N^{-0.45}$: bound $CN^{o(1)}[4N^{-1}+2N^{-0.275}]$ (original: $C[4N^{-0.85}+2N^{-0.15}]$).
Steps that the original leaves implicit are proved in Lemma \ref{lem:countC}(a); Lemma \ref{lem:U8}; Proposition \ref{prop:U9} and Theorem \ref{thm:JD}; Lemmas \ref{lem:convex}, \ref{lem:Q}, \ref{lem:countC}(b); Lemmas \ref{lem:U11}, \ref{lem:goodC}.
[ESYY] Assumption II is not needed beyond (RA5); Assumption IV is derived from \eqref{good} (Lemma \ref{lem:countC}), not from the local law.

\noindent\emph{Thresholds and constants of Part~III.} $N_*$, $N_K$, $N_3=N_4$ depend only on $(\vartheta,\delta_\pm,C_0)$ (and $K$); $N_c$ on $(c,\rho_O)$; $N_{\rm q}=\max(N_3,N_{n+2},N_c,n^2)$; $A_1=nA_0+1$; $C_1'=3^nC_1''(c,n,\rho_O)$; $N_2=N_2(\eps',\delta,A_1,N_{\rm q},n)$; $C=\max(C_1',2(c/2)^{1/2}N_2^{n+\delta/2})$. Linear dependence on $\|O\|_\infty+\Lip(O)$ throughout (PD-21).

\subsection{Infrastructure results}
The results of Part~III that are not available in Mathlib and are proved here (they are contained in the sections marked \INFRA, see \S\ref{sec:intro}) are: the Gaussian ensemble on $\cH$ (Lemma \ref{c3:GE});
the Mehler semigroup (Lemmas \ref{lem:M1}--\ref{lem:M3}, with Lemma \ref{c3:smooth}); integration by parts on $\Sig$ (Lemma \ref{lem:div}, Corollary \ref{cor:ibp}); the joint eigenvalue density of G$\beta$E (Theorem \ref{thm:JD});
the Pr\'ekopa--Leindler inequality (Theorem \ref{thm:PL}); the logarithmic Sobolev inequality (Theorem \ref{thm:LSI}); gap concentration (Lemma \ref{lem:2pt}, Theorem \ref{thm:conc}, Lemma \ref{lem:DV}); and the Gaussian relative entropy (Lemma \ref{c3:KL}).
The differences between Part~III and the original are listed in Appendix~\ref{app:diff} (\S\ref{sec:diff}).

\subsection{Remarks for the formalization}\label{sec:risks}
\begin{enumerate}[leftmargin=1.8em,itemsep=2pt]
\item Domination for derivatives of Gaussian convolutions: explicit dominating functions are given in Lemma \ref{c3:smooth} and Lemma \ref{lem:M2}(iii); they are used in Proposition \ref{prop:U9}(iv) and in Steps 0--3 of the proof of Proposition \ref{prop:entprod}. This is mathematically routine, but it is the longest analytic bookkeeping of Part~III.
\item Theorem \ref{thm:LSI} uses two-sided differentiability in Step 1 and Fatou's lemma in Step 2; Theorem \ref{thm:PL} is proved via inner regularity, without measurability of sumsets.
\item Lemma \ref{lem:2pt}, Theorem \ref{thm:conc} and Corollary \ref{cor:gap} are new in this form; they were also checked numerically.
\item Theorem \ref{thm:JD} is proved by an ergodicity argument; it was also checked by Monte Carlo simulation for $N=2$.
\item Part~III uses Theorem~\ref{thm:22} in the supremum form over $|E|\le5$, for every ensemble of $\mathcal E_N$, with constants depending only on $(\vartheta,\delta_\pm,C_0)$; it is applied to $\widetilde H_t$ through Lemma \ref{c2:transfer} and Lemma \ref{lem:U11}.
\item Mathlib names marked $^\circ$ in Appendix~\ref{app:mathlib} are expected names that have not been checked against a fixed version of Mathlib.
\end{enumerate}


% ================= APPENDICES =================
\clearpage
\appendix
\section{Exponent ledger}\label{app:ledger}
This appendix is the single place where the exponents $C_1,\dots,C_5$, $c_0$, $b$ of Part~I are fixed and where every relation between them that is used anywhere in Sections~\ref{sec:main}--\ref{sec:rig} is verified.
All exponents are affine functions of $\fa=1/\vth\ge1$ (after the normalisation $\vth\le1$ of Convention~\ref{conv:const}).

\subsection{The final values}\label{app:values}
\begin{equation}\tag{A.1}\label{eq:a5-ledger}
\begin{gathered}
 C_1=18\fa+20,\qquad C_2=12\fa+19,\qquad C_3=24\fa+36,\qquad C_4=30\fa+44,\qquad C_5=18\fa+21,\\
 c_0=\frac1{8C_3}=\frac1{192\fa+288},\qquad b=12\fa+12 .
\end{gathered}
\end{equation}
These are the values \eqref{eq:ledger0} of Section~\ref{sec:main} and the value of $b$ of Section~\ref{sec:strong}. The explicit values of the named (non-exponent) constants are collected in Section~\ref{app:consts}.
The following table records where each quantity is fixed.

\begin{center}
\renewcommand{\arraystretch}{1.15}\small
\begin{tabular}{|p{0.27\textwidth}|p{0.43\textwidth}|p{0.22\textwidth}|}
\hline
quantity & value / bound & fixed in\\\hline
truncation level $T=\kappa^{\fa}$ & $\kappa=K_w=K^2$ & Lemma~\ref{lem:LD}, Def.~\ref{def:levels}\\
linear loss $\ell(\kappa)$ & $5\kappa^{\fa+1/2}$ & \eqref{eq:loss}\\
quadratic loss $\ell_2(\kappa,N)$ & $50\kappa^{2\fa+1}\sqrt{2\lceil\log_2N\rceil+2}$ & \eqref{eq:loss}, decoupling\\
$\fL=3C_0^2\ell_2(K^2,N)$ & $150\le\fL\le K^{4\fa+3}$ & Def.~\ref{def:levels}, Lemma~\ref{lem:levels}\\
$\Emax=25K^{2\fa}\sqrt{C_0\log N}$ & $1\le\Emax\le K^{2\fa+1}$ & Lemmas~\ref{lem:norm0}, \ref{lem:levels}\\
weak law, near & $\Ld+\Lo\le K\fL(N\eta)^{-1/3}\le K^{4\fa+4}(N\eta)^{-1/3}$ & Theorem~\ref{thm:weak}(a)\\
weak law, far & $\Ld+\Lo\le9\fL(N\eta)^{-1/2}$ & Theorem~\ref{thm:weak}(b)\\
$\mathfrak E_0=6\cdot10^6C_s^2\fL^2$ & $\le K^{8\fa+7}$ & Lemma~\ref{lem:apriori}\\
FA exponent $c_F$ & $8\fa+7$ & Theorem~\ref{thm:FA}\\
FA domain & $N\eta\ge p^{12\fa+11}$, $\ \log N\le p\le N^{1/(40\fa+40)}$, $p\ge p_0$ & \eqref{eq:FAhyp}\\
$\mathfrak F=C_Fp_1^{c_F}$ \ ($p_1\le2K^{3/2}$) & $\le K^{12\fa+11}$ & Corollary~\ref{cor:FA}\\
$b$, $\mathfrak E=K^b$, $\mathfrak E_3=K^{b+3}$ & $b=12\fa+12$ \ (needs $b>12\fa+11$) & Section~\ref{sec:strong}\\
$\mathfrak E_0'=3(\Emax+14)\fL^2$ & $\le K^{10\fa+8}$ & Proposition~\ref{prop:S4}\\
first pass & $Y\le42(\varkappa+K^b/(N\eta))$; $\ |\avg X|\le\frac{K^{b+1}}{N\eta}(\varkappa+\frac{K^b}{N\eta})$ & Proposition~\ref{prop:S1}\\
law on the lattice & $\Lambda\le K^{b+4}/(N\eta)$; $\ \Ld+\Lo\le K^{b+5}(\dots)$ & Proposition~\ref{prop:S2}\\
improved law & factor $K^{b+2}$ & Props.~\ref{prop:S3}, \ref{prop:S4}\\
$C_1$ (domain $\eta\ge K^{C_1}/N$) & $18\fa+20$ \ (needs $C_1-1\ge\max\{12\fa+18,\,18\fa+17\}$) & Thm~\ref{thm:weak}, Cor.~\ref{cor:FA}\\
$C_2$ (law) & $12\fa+19=b+7$ \ (needs $C_2\ge b+6$) & end of Section~\ref{sec:strong}\\
$C_3$ (norm) & $24\fa+36=2b+12$ \ (needs $\frac12C_3>b+2$) & Proposition~\ref{prop:norm}\\
$C_5$ (counting) & $18\fa+21=C_1+1$ & Proposition~\ref{prop:rig}\\
$C_4$ (rigidity) & $30\fa+44=C_3+\frac13C_5+1$ & Proposition~\ref{prop:rig}\\
$c_0$ (range $K\le N^{c_0}$) & $1/(8C_3)$ & Lemma~\ref{lem:ledger}\\
$\phi$, $A_0$ & $1/(60\fa+88)$, $\ 180\fa+264$ & Prop.~\ref{prop:reduction}, Lemma~\ref{lem:phi}\\\hline
\end{tabular}
\end{center}

\subsection{The constraints and their verification}\label{app:constraints}
The following list contains every relation between the exponents that is used in Part~I; the right column says where it is used.
\begin{center}
\renewcommand{\arraystretch}{1.15}\small
\begin{tabular}{|l|p{0.40\textwidth}|p{0.42\textwidth}|}
\hline
 & relation & used in\\\hline
(L1) & $C_1-1\ge12\fa+18$ & \eqref{eq:ledgerA3}: Lemmas~\ref{lem:apriori}, \ref{lem:weakdich}, Theorem~\ref{thm:weak}(c); (T4)\\
(L2) & $c_0\,(6\fa+6)\le\frac1{32}$ & \eqref{eq:ledgerA3}: Lemma~\ref{lem:levels}(S4), $|\mathbb L|\le N^{10}$, far regime; (T5)\\
(L3) & $C_1-1\ge18\fa+17$ & Corollary~\ref{cor:FA} ($N\eta\ge p_1^{12\fa+11}$ on $\bD^{\rm near}$)\\
(L4) & $c_0<\dfrac{2}{3(40\fa+40)}$ & Corollary~\ref{cor:FA} ($p_1\le N^{1/(40\fa+40)}$)\\
(L5) & $b\ge12\fa+12$ & Propositions~\ref{prop:S1}--\ref{prop:S4} (via \eqref{eq:a5-T1})\\
(L6) & $C_2\ge b+7$ & \eqref{eq:a5-T3}: Theorem~\ref{thm:K}(i),(ii)\\
(L7) & $C_3\ge2b+12$; in particular $\frac12C_3>b+2$ & \eqref{eq:a5-T3}: Proposition~\ref{prop:norm}, Step 3 (edge exclusion)\\
(L8) & $c_0\,(C_1-1)\le\frac13$ & \eqref{eq:a5-T3}: $\eta_e\approx N^{-2/3}$ lies in $\bD^{\rm near}$\\
(L9) & $C_5\ge\max\big(C_1,\,C_2+\frac13\big)+1$ & \eqref{eq:ledgerrel}: Proposition~\ref{prop:rig}(iv)\\
(L10) & $C_4\ge\max\big(C_5,\,C_3+\frac13C_5\big)+1$ & \eqref{eq:ledgerrel}: Proposition~\ref{prop:rig}(v)\\
(L11) & $c_0\max(C_1,C_3,C_5)\le\frac18$ & \eqref{eq:ledgerrel}: Proposition~\ref{prop:rig}\\
(L12) & $C_1,\dots,C_5\ge1$, $\ c_0>0$ & Proposition~\ref{prop:reduction}\\\hline
\end{tabular}
\end{center}

\begin{lemma}[The ledger closes]\label{lem:ledger}
For every $\fa\ge1$ the values \eqref{eq:a5-ledger} satisfy \textup{(L1)--(L12)}.
\end{lemma}
\begin{proof}
All verifications are comparisons of affine functions of $\fa$ on $\fa\ge1$ (in fact on $\fa\ge0$).

(L1) $C_1-1=18\fa+19\ge12\fa+18$.\quad
(L2) $32(6\fa+6)=192\fa+192\le192\fa+288=1/c_0$.\quad
(L3) $18\fa+19\ge18\fa+17$.

(L4) $3(40\fa+40)=120\fa+120<384\fa+576=2/c_0$.\quad
(L5) holds with equality.\quad
(L6) $b+7=12\fa+19=C_2$.

(L7) $2b+12=24\fa+36=C_3$, and $\frac12C_3=b+6>b+2$.\quad
(L8) $3(C_1-1)=54\fa+57\le192\fa+288=1/c_0$.

(L9) $C_1+1=18\fa+21=C_5$ and $C_2+\frac43=12\fa+20\frac13\le18\fa+21$.

(L10) $C_3+\frac13C_5+1=24\fa+36+6\fa+7+1=30\fa+44=C_4$ and $C_5+1=18\fa+22\le30\fa+44$.

(L11) $C_1<C_5=18\fa+21\le24\fa+36=C_3$, so the maximum is $C_3$, and $c_0C_3=\frac18$.\quad
(L12) is clear.
\end{proof}
\noindent The condition $\frac12C_3>b+2$ needed for the edge exclusion is (L7); the three inequalities of Step~3 of Proposition~\ref{prop:norm} use it in the forms $K^{b+6}\ge16\sqrt2\cdot8$, $K^{b+10}\ge256$ and $K^{4}\ge16\sqrt2$.

\begin{remark}[Robustness]\label{rem:robust}
The proofs of Sections~\ref{sec:prob}--\ref{sec:rig} use the exponents only through (L1)--(L12): any other choice of $(C_1,\dots,C_5,c_0,b)$ satisfying these relations gives Theorem~\ref{thm:K} (with a different $N_0$), and then Proposition~\ref{prop:reduction} gives Theorems~\ref{thm:21} and~\ref{thm:22}
with $\phi$ defined by the same formula. The values \eqref{eq:a5-ledger} are generous; nothing was optimised. The three inputs that determine everything are: the loss $\fL\le K^{4\fa+3}$ of the large deviation bounds at level $K^2$, the exponent $c_F=8\fa+7$ of fluctuation averaging at level $p_1\approx K^{3/2}$
(so $\mathfrak F\le K^{\frac32c_F+\frac12}$), and the domain $N\eta\ge p^{12\fa+11}$ of Theorem~\ref{thm:FA}.
\end{remark}

\subsection{The constant \texorpdfstring{$\phi$}{phi} of Proposition \ref{prop:reduction}}\label{app:phi}
\begin{lemma}\label{lem:phi}
For the values \eqref{eq:a5-ledger} and every $\fa\ge1$,
\[
 \phi=\min\Big\{\frac{10}{C_1},\ \frac{2}{C_2},\ \frac{4}{C_3},\ \frac1{2C_4},\ \frac1{2C_5},\ 5c_0\Big\}=\frac1{2C_4}=\frac1{60\fa+88},\qquad A_0=\frac3\phi=180\fa+264 .
\]
\end{lemma}
\begin{proof}
The six numbers are the reciprocals of $1.8\fa+2$, $\ 6\fa+9.5$, $\ 6\fa+9$, $\ 60\fa+88$, $\ 36\fa+42$, $\ 38.4\fa+57.6$; the largest of these is $60\fa+88$.
\end{proof}
\noindent So Proposition~\ref{prop:reduction} holds as stated, with the explicit constants $\phi=1/(60\fa+88)$, $A_0=180\fa+264$, $C=1$, $c=\frac12$, where $\fa=1/\min(\vth,1)$. With $K=(\log N)^{\phi L}$ the powers of $\log N$ in the final statements are:
domain $K^{C_1}=(\log N)^{\phi C_1L}\le(\log N)^{0.5L}$, far inside the required $(\log N)^{10L}$; errors $K^{C_2}\le(\log N)^{0.4L}$ (required: $4L$), $K^{C_3}\le(\log N)^{0.5L}$ (required: $9L$), $K^{C_4}=(\log N)^{L/2}$ and $K^{C_5}\le(\log N)^{0.4L}$ (required: $L$);
probability $\exp(-\frac12(\log N)^{\phi L})$. (Indeed $\phi C_1=\frac{18\fa+20}{60\fa+88}\le\frac12$, $\phi C_2\le\frac25$, $\phi C_3=\frac{24\fa+36}{60\fa+88}\le\frac12$, $\phi C_5\le\frac25$.) The binding constraint is the rigidity exponent $C_4$.

\subsection{Levels, events and probabilities}\label{app:levels}
\begin{center}
\renewcommand{\arraystretch}{1.15}\small
\begin{tabular}{|p{0.29\textwidth}|p{0.13\textwidth}|p{0.25\textwidth}|p{0.22\textwidth}|}
\hline
event & level & failure probability & defined in\\\hline
$\cG(z)$ (large deviations at $z$) & $K_w=K^2$ & $46\,\fa N^4e^{-K^2}$ & Def.~\ref{def:good}, Lemma~\ref{lem:levels}\\
$\cG_{\rm col}(z)$, $\ \Xi(z)\supset\cG_{\rm col}(z)$ & $K^2$ & $460\,\fa N^8e^{-K^2}$ & Def.~\ref{def:domains}, Thm~\ref{thm:weak}\\
$\{\norm H\le\Emax\}$ & $K^2$ & $(N^2\fa+1)e^{-K^2}$ & Lemma~\ref{lem:norm0}\\
$\cG^1(z)$ within $\cG_{\rm col}(z)$ & $p_1\approx K^{3/2}$ & $e^{-p_1}$ & Proposition~\ref{prop:S1}\\
$\cG^2(z)$ & $p_1$ & $N^{10}e^{-p_1}$ & Proposition~\ref{prop:S2}\\
$\Xi_K^{(1)}$ & $p_1$ & $N^{20}e^{-p_1}$ & end of Section~\ref{sec:strong}\\
$\cG^3(z)$ within $\cG^2(z)$; $\cG^{\rm f}(z)$ within $\cG_{\rm col}(z)$ & $p_2\approx K$ & $e^{-p_2}$ & Props.~\ref{prop:S3}, \ref{prop:S4}\\
$\Xi_K^{(2)}$ & $p_2$ & $N^{21}e^{-p_1}+N^6e^{-p_2}$ & Proposition~\ref{prop:norm}\\
$\Xi_K=\Xi_K^{(1)}\cap\Xi_K^{(2)}$ & --- & $2N^{21}e^{-p_1}+N^6e^{-p_2}\le e^{-K/2}$ & Def.~\ref{def:XiK}, Prop.~\ref{prop:XiK}\\\hline
\end{tabular}
\end{center}
\noindent The relations between the levels are \eqref{eq:levelrel}: $K^2\ge40p_1\log N$, $p_1\ge40p_2\log N$, $p_2\ge K\ge(\log N)^3$. Each passage from a level to the next costs one factor $40\log N$ in the exponent, through a term $2N\,\PP(\cdot)^{1/p}$ in a moment bound ((T6) of Section~\ref{sec:strong}, and the term $N^{2}\PP(\Xi_p^c)^{1/(2p)}$ in Corollary~\ref{cor:FA});
there are exactly two such passages. This is why no lower bound of the form $L\ge A_0\log\log N$ is needed in the proof: it enters only through $K\ge(\log N)^3$. All lattice unions have at most $N^{10}$ members.

\subsection{Thresholds}\label{app:N0}
The threshold $N_0$ of Theorem~\ref{thm:K} is the maximum of the following finitely many numbers, each of which depends only on the normalised model parameters $(\vth,\delta_-,\delta_+,C_0)$:
\begin{enumerate}[label=\textup{(N\arabic*)},leftmargin=2.8em,itemsep=1pt]
\item the threshold of Lemma~\ref{lem:levels} for the constant $K_0:=\max\{K_0^{(5)},\,124800,\,2142\,C_s\}$, where $K_0^{(5)}$ is the maximum of the explicit constants of Section~\ref{sec:weak} (they are named in its proofs:
$\frac54C_c+2072C_s+70$, $(1152C_s)^{1/2}$, $416^2$, $2037$, $8/\lambda_{\rm st}$, $255$, $68$, $9$); this threshold includes $\log N\ge41^2$ (for \eqref{eq:levelrel}), $(\log N)^{5/2}\ge300C_0^2$, $N\ge13$, $N^2\ge C_0$, $(\log N)^6\ge\kappa_\vth+3\log N$ and $K_0\le N^{3/32}$; and the explicit conditions on $N$ alone in the proofs of Section~\ref{sec:weak} (such as $N^{3/2}\ge156/c_*$);
\item the threshold of Theorem~\ref{thm:FA} (conditions \eqref{eq:N0FA}) and the conditions in the proof of Corollary~\ref{cor:FA}: $(\log N)^{3/2}\ge\max\{2^{12\fa+11},\,C_F2^{8\fa+7}\}$, $(\log N)^3\ge p_0$ (the constant of \eqref{eq:FAhyp}) and $C_we^{9/2}\le N^{10}$;
\item the threshold $\max\{13,\sqrt{C_0}\}$ of Lemma~\ref{lem:norm0} (contained in (N1));
\item the explicit conditions of Sections~\ref{sec:strong}--\ref{sec:norm}: $N^{31}\ge\max\{920\,\fa e^8,\,2e^{10}\}$; $\ N^2\ge460\,\fa+1$; $\ N^{15/32}\ge\frac{15600}{\sqrt{10}}$; $\ N^{3/2}\ge23400$; $\ N^{7/32}\ge72$; $\ N^{7/16}\ge4$; $\ N^{13/32}\ge24$; $\ N\ge400$;
$\ N^6\ge462\,\fa+2$ (which gives $N^2\fa+1+460\fa N^{13}+N^{15}\le N^{21}$, $2N^5\le N^6$ and $N^{20}+46\fa N^{14}\le N^{21}$);
\item $N\ge e^{7}$ (Proposition~\ref{prop:XiK}); this implies the conditions $N\ge e^6$ of Proposition~\ref{prop:rig} and $N\ge16$ of Proposition~\ref{prop:reduction}.
\end{enumerate}
Conditions (N3)--(N5) are implied by (N1) (which forces $\log N\ge1681$) except for the dependence on $\fa$ in (N4), which is harmless. We make no attempt to estimate $N_0$: it is astronomically large, because several absorptions are of the form
``$K^{1/2}\ge(\log N)^{3/2}\ge$ a constant'' with constants such as $2^{12\fa+11}$ or $C_F2^{8\fa+7}$, and because of the condition $(\log N)^3\ge p_0$ of (N2)
(for $\vth=1$, $C_0=1$ the constant $p_0$ of \eqref{eq:FAhyp} is of order $10^{19}$). This affects none of the statements, in which $N_0$ is existential, and it is the price of having error terms that are pure powers of $K$ with no numerical constants.

\subsection{Named constants}\label{app:consts}
The following named constants are fixed in the text; none of them affects an exponent of \eqref{eq:a5-ledger}.
\begin{itemize}[leftmargin=2em,itemsep=1pt]
\item $\lambda_{\rm st}=1/(576C_s)$ in \eqref{eq:lambda1}, together with $D\le2\Ld\le2\lambda_{\rm st}$ in Lemma~\ref{lem:stabnear}; the constants of \eqref{eq:quad} are $40C_s^2\Upsilon_{\max}^2+(90+72C_s)\Lambda^3$, and the coefficient of $\Lambda^2$ in \eqref{eq:ap3} is in fact $\frac3{10}$ (the statement keeps $\frac34$, which is what Sections~\ref{sec:strong}--\ref{sec:norm} use).
In Section~\ref{sec:strong} the value of $\lambda_{\rm st}$ enters only through $144C_s\lambda_{\rm st}=\frac14$ (proof of Proposition~\ref{prop:S2}(c)) and through the statements of Section~\ref{sec:weak}.
\item Lemma~\ref{lem:farstep}: $\Ld\le6\fL(N\eta)^{-1/2}$, $\Lo\le3\fL(N\eta)^{-1/2}$; Theorem~\ref{thm:weak}(b): $\Ld+\Lo\le9\fL(N\eta)^{-1/2}$. Proposition~\ref{prop:S4} does not use these numbers: it uses only $\max_i|G_{ii}|\le\frac12$ from Theorem~\ref{thm:weak}(b) and re-derives $\Psi_m\le6(N\eta)^{-1/2}$.
\item Explicit values: $c_*=\frac1{150}$, $C_*=4$ (Lemma~\ref{lem:mscsize}); $\mathfrak E_0=6\cdot10^6C_s^2\fL^2$ (Lemma~\ref{lem:apriori}); $\PP(\cG(z)^c)\le46\,\fa N^4e^{-K^2}$; $c_F=8\fa+7$ and $C_F$ as in \eqref{eq:CF}; $C_6=10$, $C_7=9$ (Lemmas~\ref{lem:count}, \ref{lem:conv}).
\item In Sections~\ref{sec:strong}--\ref{sec:norm}: the constants $42$ and $194$ in \eqref{eq:Y}, \eqref{eq:Xmom}; $7000$ in \eqref{eq:a5-7e}; $\mathfrak E_0'=3(\Emax+14)\fL^2$; the lower bound $K_0\ge\max\{124800,2142C_s\}$ in \eqref{eq:a5-T2}.
\end{itemize}
\section{Dependency graph}\label{app:dep}
The graph below shows the dependencies in Part~I. It is acyclic. An arrow from a box to a lower box means ``is used in''.
{\small
\begin{verbatim}
 Mathlib: spectral theorem, block inverse/determinant, product measures, Fubini,
          Hoeffding, layer cake, Hoelder/Minkowski/Markov, divergence theorem
    |
 [2 model: D2.1, D2.2, L2.5 (class: monotone, reflection, scaling)]
    |
 [3.1 m_sc: L3.1-3.4]     [3.2 resolvent: L3.5-3.7]     [3.4 Weyl L3.13, net L3.14]
    |                          |                               |
 [3.3 stability: L3.8 (exact equations), L3.9 (vector, on e-perp), L3.10 (near;      |
      lambda_st = 1/(576 C_s)), L3.11 (scalar quadratic), L3.12 (far)]                |
    |                                                                                |
 [4 LD: L4.1 moments -> L4.2 tail form -> L4.3 moment form;  L4.7 product splitting] |
    |                                                                                |
 [L4.4 good event G(z)]  [D4.2 levels, L4.6 standing ineq.]  [L4.5 crude norm bound]
    |                                    |                                           |
 [5 weak law: L5.1 a priori bounds, L5.2 dichotomy, L5.3 initial, L5.4 far step     |
              ->  T5.5 (lattice; deterministic on G_col(z); event Xi(z))]            |
    |          \                                                                     |
    |           [6 FA: L6.3-6.7, P6.8 (p = 2), L6.9-6.13 -> T6.1                     |
    |            (uses L3.6 (I1),(I4),(I5), L3.7, L4.1, L4.3, L4.7) -> C6.2          |
    |            (uses T5.5(c) only through P(Xi(z)^c), and (4.13))]                 |
    |          /                                                                     |
 [7 strong law: P7.1 first pass (L5.1, L3.11, C6.2 at p_1)                           |
       -> P7.2 dichotomy + continuity (L5.1, L5.3, L3.5(c), L3.1(f))                 |
       -> P7.3 second pass (C6.2 at p_2);   P7.4 far (T5.5(b), L4.6, L3.12, C6.2);   |
       L7.5 Lipschitz extension (L3.5(c), L3.1(f))]                                  |
    |                         |                                                      |
 T2.6 (i),(ii)          [8 norm: P8.1 (L4.5, P7.3 + L3.2, P7.4 + L3.1(g), L3.5(e)) <-+
    |                         |
    |                    T2.6 (iii);   D8.1 event Xi_K,  P8.2  P(Xi_K^c) <= exp(-K/2)
 [3.5 HS L3.15 -> counting L3.16 (uses L3.3)]
    |                         |
 [9 rigidity: L9.1 conversion (uses L3.4) -> P9.2 ((i),(iii) => (iv),(v); deterministic)]
    |
 T2.6 (iv),(v)   --[Appendix A: L A.1 ledger (L1)-(L12), L A.2 phi]-->
 P2.7 (reduction K -> L)  ->  Theorems 2.1 and 2.2
\end{verbatim}
}
\noindent The weak law does not use Section~\ref{sec:FA}; Section~\ref{sec:FA} uses the weak law only through $\PP(\Xi(z)^c)$; Sections~\ref{sec:strong}--\ref{sec:norm} use Section~\ref{sec:FA} only through Corollary~\ref{cor:FA}; Section~\ref{sec:rig} is deterministic and uses Sections~\ref{sec:strong}--\ref{sec:norm} only through the statements (i), (iii) of Theorem~\ref{thm:K}.
The exponent ledger (Lemma~\ref{lem:ledger}) is used by Lemma~\ref{lem:levels}, Corollary~\ref{cor:FA}, Sections~\ref{sec:strong}--\ref{sec:norm} and Proposition~\ref{prop:rig}, each through the relations listed in Section~\ref{app:constraints}.

\medskip\noindent\textbf{Dependencies of Sections~\ref{sec:strong}--\ref{sec:norm} in detail.}
\begin{center}
\renewcommand{\arraystretch}{1.12}\small
\begin{tabular}{|p{0.20\textwidth}|p{0.73\textwidth}|}
\hline
statement & uses (exactly)\\\hline
(T4)--(T6) & Def.~\ref{def:domains}; Lemma~\ref{lem:mscsize}; Lemma~\ref{lem:stabnear} (size of $w_0$); Lemma~\ref{lem:basic}(a); Lemma~\ref{lem:msc}(b); Lemma~\ref{lem:levels}(S3),(S4); \eqref{eq:levelrel}; (L1), (L2), (L8)\\
Proposition~\ref{prop:S1} & Theorem~\ref{thm:weak}(a); \eqref{eq:ap3}; Lemma~\ref{lem:scalar} (real statement); Corollary~\ref{cor:FA} with $p=p_1$; Lemma~\ref{lem:weyl}(e); (T4), (T6); (L5)\\
Proposition~\ref{prop:S2} & Proposition~\ref{prop:S1}; \eqref{eq:ap2}--\eqref{eq:ap4}; Lemma~\ref{lem:initial}; \eqref{eq:lipnear} (i.e.\ Lemma~\ref{lem:basic}(c), Lemma~\ref{lem:msc}(f)); Lemma~\ref{lem:scalar}; $144C_s\lambda_{\rm st}=\frac14$\\
Proposition~\ref{prop:S3} & Proposition~\ref{prop:S2}(a); \eqref{eq:ap3}; Corollary~\ref{cor:FA} with $p=p_2$; (T4), (T6)\\
Proposition~\ref{prop:S4} & Theorem~\ref{thm:weak}(b) ($\max_i|G_{ii}|\le\frac12$); Lemma~\ref{lem:levels}(S2),(S4),(S5); Lemma~\ref{lem:stabfar}; Corollary~\ref{cor:FA} with $p=p_2$; (T5), (T6)\\
Lemma~\ref{lem:lip} & Lemma~\ref{lem:basic}(c); Lemma~\ref{lem:msc}(f); Def.~\ref{def:domains}\\
Theorem~\ref{thm:K}(i),(ii) & Proposition~\ref{prop:S2}(b),(c); Lemma~\ref{lem:lip}; (L6)\\
Proposition~\ref{prop:norm} & Lemma~\ref{lem:norm0}; Proposition~\ref{prop:S4} with Lemma~\ref{lem:msc}(g); Proposition~\ref{prop:S3} with \eqref{eq:imm}; Lemma~\ref{lem:basic}(e); Lemma~\ref{lem:weyl}(a); (T5); (L7), (L8)\\
Proposition~\ref{prop:XiK}, Theorem~\ref{thm:K} & Propositions~\ref{prop:S2}, \ref{prop:norm}; \eqref{eq:Krange}; Proposition~\ref{prop:rig} (for (iv), (v)); Lemma~\ref{lem:ledger}\\\hline
\end{tabular}
\end{center}



\section{Differences from the original}\label{app:diff}
``The original'' is the paper of Erd\H{o}s, Yau and Yin (arXiv:1007.4652v7); its equation and theorem numbers are quoted in upright type, e.g.\ (2.17), Theorem~2.3.
This appendix is the only list of differences between NP and the original. Section~\ref{app:diffPD} lists every difference between the statements of the original and those of NP,
under the labels PD-01, \dots, PD-21, which are used for cross-reference throughout NP. Sections~\ref{app:diffNP}, \ref{app:diffB} and~\ref{sec:diff} list, for Parts~I, II and~III,
the remaining differences (constants, quantifiers, formulations) and compare the proofs with those of the original; items already contained in Section~\ref{app:diffPD} are not repeated there.
None of the differences weakens a main result, with one exception that is a correction: Theorem~\ref{thm:23} is stated for Lipschitz observables (PD-21).

\emph{Out of scope.} The following statements of the original are not part of NP: the corollary (2.37) on averaged correlation functions and gap distributions; (2.42) and the extension to several extreme eigenvalues;
Theorem~6.4 (Green function comparison in its general form); the quaternion self-dual case; the remarks on weakening (2.17).

\subsection{Differences in the statements: PD-01 -- PD-21}\label{app:diffPD}
\begin{description}[leftmargin=1.5em,itemsep=3pt,font=\normalfont\bfseries]
\item[PD-01] \emph{Place:} remark after (2.18)--(2.20). \emph{Original:} ``the upper bound on $L$ originates from $\mathbf S_L\ne\emptyset$''; but at the upper endpoint of (2.18) the set $\mathbf S_L$ is empty. \emph{NP:} adopted as a remark: the range \eqref{eq:Lbound} is that of the original; at its upper endpoint (i), (ii) of Theorem~\ref{thm:21} are vacuous (Remark~\ref{rem:Kstronger}).
\item[PD-02] \emph{Place:} (2.18), (4.4), Theorem~3.1. \emph{Original:} three different upper bounds for the exponents $L$, $\ell$. \emph{NP:} adopted: one range \eqref{eq:Lbound} for $L$ in the final statements; the proofs use only the level $K$ with the single range \eqref{eq:Krange}; there is no second exponent $\ell$.
\item[PD-03] \emph{Place:} (1.3). \emph{Original:} the classical location $\gamma_N$ is not unique ($N\nsc(\gamma)=N$ for all $\gamma\ge2$). \emph{NP:} adopted: $\gamma_j:=\inf\{x:N\nsc(x)\ge j\}$, so $\gamma_N=2$ (\eqref{eq:sc}, Lemma~\ref{lem:nsc}(c)).
\item[PD-04] \emph{Place:} Theorem~2.3, (2.36). \emph{Original:} $E\in[2-c,2+c]$; $\varrho(E)$ undefined, and $\rhosc(E)=0$ for $E\ge2$. \emph{NP:} adopted: bulk energies $E\in[-2+c,2-c]$, $0<c\le2$, $b\in(0,c/2)$, $\varrho=\rhosc(E)$ (Theorem~\ref{thm:23}).
\item[PD-05] \emph{Place:} (2.33), (2.36). \emph{Original:} correlation functions $p^{(n)}$ written as densities. \emph{NP:} adopted: $p^{(n)}$ is the $n$-point marginal of the symmetrised eigenvalue law, a measure; \eqref{eq:abstrthm2} is in the equivalent density-free form with $\mathcal X^O_{E',\varrho}$ and $C_{N,n}$ (comment (a) after Theorem~\ref{thm:23}).
\item[PD-06] \emph{Place:} (2.28), (2.31)--(2.32), symmetric case. \emph{Original:} diagonal variance of the matrix Brownian motion and the coefficient of the generator are those of the Hermitian case. \emph{NP:} adopted: for $\kk=\R$, $V$ is the GOE with diagonal variance $2/N$; $H_t$ is defined by the formula (2.29); the generator of the eigenvalue flow is then twice (2.32) (a time change), and (A) holds for $H_t$ up to a factor $1+O(1/N)$, removed by a rescaling in Part~III (comment (d) after Theorem~\ref{thm:23}, Lemma~\ref{lem:mono}(v)).
\item[PD-07] \emph{Place:} band matrix example after (2.16). \emph{Original:} $\sigma_{ij}^2=(N/2)^{-1}\bone([i-j]_N\le N/4)$: the absolute value is missing and the row sums are $1+2/N$. \emph{NP:} not used: NP contains no examples. The corrected example is $\sigma_{ij}^2=(N/2+1)^{-1}\bone(|[i-j]_N|\le N/4)$ ($N$ divisible by $4$).
\item[PD-08] \emph{Place:} (2.37). \emph{Original:} ``$\lim_{N\to0}$''. \emph{NP:} out of scope: (2.37) is not a result of NP (read $N\to\infty$).
\item[PD-09] \emph{Place:} Lemma~4.6, (4.44). \emph{Original:} the bound on $\PP(\lambda_N\ge3)$ lacks ``$N\ge N_0$'' and is quoted from an earlier paper. \emph{NP:} adopted and superseded: the a priori norm bound is proved in NP with an explicit threshold (Lemma~\ref{lem:norm0}, Propositions~\ref{prop:S4}, \ref{prop:norm}).
\item[PD-10] \emph{Place:} Lemma~5.1, (5.2). \emph{Original:} the hypothesis of the lemma (quoted from an earlier paper) does not match its use ($E_1=-4$, constant $10$). \emph{NP:} adopted and superseded: Lemma~\ref{lem:count} is a new deterministic statement whose hypothesis \eqref{eq:HShyp} is exactly what Theorem~\ref{thm:K}(i) provides; it is proved from Lemma~\ref{lem:HS}.
\item[PD-11] \emph{Place:} Lemma~3.3 and its uses. \emph{Original:} large deviation bounds stated for equal variances but applied to unequal, possibly vanishing variances (``all variances being positive''). \emph{NP:} adopted: weighted large deviation bounds, valid for vanishing variances (Lemmas~\ref{lem:LD}, \ref{lem:LDmom}).
\item[PD-12] \emph{Place:} Theorem~3.1, Lemmas~3.6, 3.11. \emph{Original:} Theorem~3.1 only for Hermitian matrices; dependence of the constants incomplete. \emph{NP:} adopted: every statement of NP is for both $\kk=\R$ and $\kk=\C$, with constants depending on $(\vth,\delta_-,\delta_+,C_0)$ as specified in Convention~\ref{conv:const} (Theorem~\ref{thm:weak}).
\item[PD-13] \emph{Place:} discussion before Theorem~2.3. \emph{Original:} an assumption of an earlier paper is paraphrased as $U''\ge\delta$. \emph{NP:} not applicable: NP does not quote that paper, and the only confining potential in Part~III is quadratic.
\item[PD-14] \emph{Place:} (2.3). \emph{Original:} (B) as a condition on the spectrum of $B$. \emph{NP:} adopted: (B) is the equivalent quadratic form bound on $\be^\perp$ (Definition~\ref{def:ass}; equivalence in Lemma~\ref{lem:mono}(i)).
\item[PD-15] \emph{Place:} (2.17), (2.40). \emph{Original:} conditions on random variables; (2.40) is stated for the rescaled entries $v_{ij}$, $w_{ij}$. \emph{NP:} adopted: both conditions are stated for the laws $\nu_{ij}$ of $h_{ij}$ (\eqref{eq:subexp}, \eqref{eq:2m}); both are scale invariant (Lemma~\ref{lem:mono}(v)).
\item[PD-16] \emph{Place:} Theorem~2.3. \emph{Original:} dependence of $C$ incomplete; $N\ge n$ implicit. \emph{NP:} adopted: $C$ depends on $(\eps',\delta,c,n,\rho_O,\vth,\delta_-,\delta_+,C_0)$ only; $N\ge\max(n,3)$ (Theorem~\ref{thm:23}).
\item[PD-17] \emph{Place:} proof of Theorem~2.2, (5.11)--(5.15). \emph{Original:} rigidity proved for $j\le N/2$ only. \emph{NP:} superseded: Lemma~\ref{lem:conv} treats both edges directly; the reflection $H\mapsto-H$ (Lemma~\ref{lem:mono}(iv), Lemma~\ref{lem:weyl}(c)) is not needed.
\item[PD-18] \emph{Place:} definition of $\Omega_h$ before (3.32). \emph{Original:} $\{|h_{ij}|\ge(\log N)^{\ell/10}|\sigma_{ij}|\}$: if $\sigma_{ij}=0$ this is the whole space. \emph{NP:} adopted: strict inequalities in \eqref{eq:subexp} and in every event comparing $|h_{ij}|$ with a multiple of $\sigma_{ij}$ (Definition~\ref{def:ass} and the sentence after it; Definition~\ref{def:good}).
\item[PD-19] \emph{Place:} Lemma~4.1. \emph{Original:} constants ``depending on $\vth$''. \emph{NP:} adopted: the dependence of the constants of the fluctuation averaging bound and of its domain of validity on $(\vth,\delta_-,\delta_+,C_0)$ is stated explicitly (Theorem~\ref{thm:FA}, Corollary~\ref{cor:FA}).
\item[PD-20] \emph{Place:} (2.26). \emph{Original:} supremum over $E$ which need not be attained ($\fn$ is right-continuous). \emph{NP:} adopted: \eqref{eq:nn} is stated with the supremum itself, which is a Borel function (Lemma~\ref{lem:weyl}(f)); no ``$\exists E$'' form is used.
\item[PD-21] \emph{Place:} Theorem~2.3, (2.36). \emph{Original:} $O$ continuous with compact support, $C$ depending on $O$. \emph{NP:} adopted: $O$ Lipschitz with support in $[-\rho_O,\rho_O]^n$, and $C$ depends on $O$ only through $\rho_O$ and the factor $\norm O_\infty+\operatorname{Lip}(O)$. The source of the original's proof (an earlier paper) requires a smooth observable in the corresponding step,
and a bound of the form $C(O)\times$(power of $N$) for all continuous $O$ would, by the uniform boundedness principle, be equivalent to a bound with $C(O)=C'\norm O_\infty$, i.e.\ to a total variation bound, which is not what the method gives. H\"older continuous $O$: corollary in Part~III, with a weaker power.
\end{description}

\subsection{Part~I (Theorems~\ref{thm:21} and~\ref{thm:22})}\label{app:diffNP}
\begin{enumerate}[leftmargin=2em,itemsep=2pt]
\item \emph{Level form.} Theorems~\ref{thm:21} and~\ref{thm:22} are deduced from Theorem~\ref{thm:K}, which has one parameter $K\in[(\log N)^3,N^{c_0}]$, error factors $K^{C_i}$ with $C_i$ depending on $\vth$ alone, and probability $1-e^{-K/2}$.
This is legitimate because $\phi$ and $A_0$ are existential constants in the original. Theorem~\ref{thm:K} is slightly stronger (Remark~\ref{rem:Kstronger}).
\item \emph{One event.} Theorem~\ref{thm:K} provides a single measurable event on which all bounds hold for all $z\in\bD_K$, all $E$ and all $j$. In Theorem~\ref{thm:21}(i),(ii) the unions over the uncountable set $\bS_L$ are measured by the outer probability $\PP^*$.
\item \emph{Quantifiers.} In Theorems~\ref{thm:21}, \ref{thm:22}: the constants $(A_0,C,c,\phi,N_0)$ depend only on $(\vth,\delta_-,\delta_+,C_0)$ and are chosen before $N$, $\kk$, the ensemble and $L$; the same constants serve both theorems; $C=1$, $c=\frac12$.
\item \emph{The model.} The probability space is the product space $\kk^{\mathcal P_N}$ with the product of the entry laws; $H$ is Hermitian at every point; the diagonal laws are real (R). The assumptions are conditions on the family of laws (Definition~\ref{def:ass}).
The eigenvalues are ordered, $\lambda_1\le\dots\le\lambda_N$.
\item \emph{The function $\msc$} is defined by \eqref{eq:defmsc}, the product of two principal square roots (which is the branch with $\im\msc>0$ on all of $\C_+$), not by $\sqrt{z^2-4}$ with a principal branch. All its properties are proved (Lemmas~\ref{lem:msc}--\ref{lem:stieltjes}).
\item \emph{Counting function.} Theorem~\ref{thm:K}(iv) holds with the supremum over all $E\in\R$, not only $|E|\le5$.
\item \emph{Explicit constants.} $c_*=\frac1{150}$, $C_*=4$ (Lemma~\ref{lem:mscsize}); $C_6=10$ (Lemma~\ref{lem:count}); $C_7=9$ (Lemma~\ref{lem:conv}); the edge constants $2.8$, $3.6$ (Lemma~\ref{lem:nsc}(c)). The thresholds $N_0$ are not optimised and are very large.
\end{enumerate}

\subsubsection*{Proofs: the original versus Part~I}\label{app:diffproofs}
\begin{center}
\renewcommand{\arraystretch}{1.1}\small
\begin{tabular}{|p{0.36\textwidth}|p{0.58\textwidth}|}
\hline
original & NP (Part~I)\\\hline
Theorems 2.1, 2.2 & Theorem~\ref{thm:K} and Proposition~\ref{prop:reduction}\\
Theorem 3.1 (weak law, fixed $z$) & Theorem~\ref{thm:weak} (lattice, deterministic on good events; includes the far regime)\\
Definition 3.1, Lemma 3.2 (quoted) & Lemma~\ref{lem:schur} (block Schur complement; proved)\\
Lemma 3.3 (quoted, equal variances) & Lemmas~\ref{lem:LD}, \ref{lem:LDmom} (weighted; Hoeffding and decoupling; PD-11)\\
Lemma 3.4 (quoted) & Lemmas~\ref{lem:msc}, \ref{lem:mscsize} (proved, explicit constants)\\
Lemma 3.5 (bounds on minors) & Lemmas~\ref{lem:minor}, \ref{lem:apriori}; the trace identity (I5) of Lemma~\ref{lem:schur} replaces interlacing\\
Lemma 3.6 (events $\Omega$, $\Psi$) & Definition~\ref{def:good}, Lemma~\ref{lem:good} (no a priori event needed)\\
Lemmas 3.7--3.9 (stability, with $\log N$) & Lemmas~\ref{lem:vecstab}, \ref{lem:stabnear} (on $\be^\perp$, no $\log N$)\\
Lemmas 3.10, 3.11, 3.12 & Lemmas~\ref{lem:weakdich}, \ref{lem:initial}, Theorem~\ref{thm:weak}\\
Lemma 4.1 and Section 7 & Theorem~\ref{thm:FA} (conditioning on a minor, finite Neumann expansion)\\
Corollary 4.2, Lemmas 4.3, 4.4 (iteration) & Propositions~\ref{prop:S1}, \ref{prop:S2} (one step)\\
Lemma 4.5 & Proposition~\ref{prop:S3} and Proposition~\ref{prop:norm}\\
Lemma 4.6 (quoted; PD-09) & Lemma~\ref{lem:norm0}, Propositions~\ref{prop:S4}, \ref{prop:norm}\\
Lemma 5.1 (quoted; PD-10) & Lemmas~\ref{lem:HS}, \ref{lem:count} (Helffer--Sj\"ostrand from the divergence theorem on rectangles; explicit polynomial cutoffs)\\
Stieltjes transform of $\rhosc$ (used implicitly) & Lemma~\ref{lem:stieltjes} ($x=2\cos\theta$ and a geometric series)\\
(5.11)--(5.15), $j>N/2$ by symmetry (PD-17) & Lemma~\ref{lem:conv}, Proposition~\ref{prop:rig} (both edges directly)\\\hline
\end{tabular}
\end{center}

\subsection{Part~II (Theorem~\ref{thm:24})}\label{app:diffB}
This section lists every place where the statements or the route of Section~\ref{sec:edge} differ from \S6 of the original. Numbers of the form (6.xx), Lemma~6.x, Theorem~6.x \emph{in this section}
refer to the original unless they are given as references to NP; ``TeX $n$'' is a line of the source of the original. The main results are unchanged: Theorem~\ref{thm:24} is (2.41) of the original
together with the analogous statement for $\lambda_1$.

\begin{enumerate}\itemsep2pt
\item[\textbf{II.1.}] \textbf{Theorem 2.4 (statement).}
(a) The exponents are absolute: $\e=\frac1{100}$, $\delta=\frac1{20}$ for every $\vartheta$ (Theorem~\ref{thm:24B}); the original has ``$\e>0$ and $\delta>0$ depending on $\vartheta$'', which is implied. Only $N_0$ depends on
$(\vartheta,\delta_\pm,C_0)$; it does not depend on $s$.
(b) The statement for the smallest eigenvalue, \eqref{eq:tw1}, is written out and proved (via $H\mapsto-H$); the original only says ``Analogous result holds for the smallest eigenvalue $\lambda_1$''.
The quantifiers are: $\eps,\delta$ depend only on $\vth$ (in fact they are absolute), and $N_0$ depends on $(\vth,\delta_-,\delta_+,C_0)$ but not on $s$.
(c) The moment matching condition and the decay condition are conditions on the laws of the entries (PD-15).
(d) Events of the form $\{|h_{ij}|>x\sigma_{ij}\}$ use strict inequalities, so that vanishing variances are allowed; nothing in \S\ref{sec:edge} divides by a variance.
(e) The joint statement (2.42) for several eigenvalues and Theorem~6.4 are not treated (out of scope, see the beginning of this appendix).

\item[\textbf{II.2.}] \textbf{Lemma 6.1 and Corollary 6.2} (sharp versus smoothed counting function) are replaced by Proposition~\ref{prop:sandwich}, which compares $\P(\lambda_N<E)$ directly with
$\E q(X^\pm_E(H))$ on the rigidity event of \eqref{eq:rigidity}. Not used: (2.19) and (2.26) of the original, the bound of $\tr f(H)$ by $\im m$, the averaging over the window edge in the proof of Lemma~6.1, and the monotonicity
of $q$. The parameters are $\eta=N^{-2/3-2\e}$ and $\ell=\frac13N^{-2/3-\e}$, instead of $\eta=N^{-2/3-9\e}$, $\ell_1=N^{-2/3-3\e}$, $\ell=\frac12N^{-2/3-\e}$; the probability bounds are $N^{-20}$.

\item[\textbf{II.3.}] \textbf{Theorem 6.3} (Green function comparison at the edge). Theorem~\ref{thm:gfc} proves the integrated form (6.25) only, for $F\in C^3(\R)$ with bounded derivatives up to order $3$; this is the
form that the proof of Theorem~\ref{thm:24} needs. Differences: the error is $C_4C_FN^{-1/6+4\eg}$ with the explicit exponent $4\eg$ (original: $CN^{-1/6+C\e}$) and with $C_4$ depending on $(\vartheta,C_0)$ only; it holds for every
$\eg\in(0,\frac1{10}]$ (original: $\e<\e_0(C_1)$); the argument of $F$ is $\frac N\pi\int_{E_1}^{E_2}\im m$, i.e.\ the factor $\frac1\pi$ is included (immaterial: replace $F$ by $F(\pi\,\cdot)$).
\emph{Not proved:} the pointwise form (6.24) for $F(N\eta\im m(z))$, and polynomially growing $F$ as in (6.23). Neither is used for the main results.

\item[\textbf{II.4.}] \textbf{Lemma 6.5} (one-step comparison). In the original the functional is rewritten as $\eta\int\sum_{i\ne j}S_{ij}\overline{S_{ji}}$ and expanded entrywise by the resolvent expansion; the terms are sorted
according to $|\{i,j\}\cap\{a,b\}|$ ((6.53)--(6.69)). In NP the change of the functional under the replacement of one entry is given \emph{exactly} by a $2\times2$ determinant: Lemma~\ref{lem:det}, Corollary~\ref{cor:Xchange}(i),
$X(Q+V)-X(Q)=-\frac1\pi\im\br{\Log D_p}$ with $D_p=1+vR_{ba}+\bar vR_{ab}-|v|^2(R_{aa}R_{bb}-R_{ab}R_{ba})$. Only the four entries $R_{aa},R_{bb},R_{ab},R_{ba}$ at the two points $z_1,z_2$ occur; there is no sum over $i,j$,
no integral over the spectral parameter inside the expansion, and no case distinction. Lemma~\ref{lem:expand}, Corollary~\ref{cor:entry} and Proposition~\ref{prop:step} replace Lemma~6.5; the error exponents are the same,
$N^{-13/6+C\e}$ $(a\ne b)$ and $N^{-11/6+C\e}$ $(a=b)$, now with $C=4$, resp.\ $C=2$. The second-order terms $A_N,B_N$ of the original correspond to $\mathfrak Q_2$ \eqref{eq:b3:Q2}, \eqref{eq:b3:Q2d}, given by explicit formulas.

\item[\textbf{II.5.}] \textbf{Lemma 6.6} (the third-order term) is replaced by Lemma~\ref{lem:decor}: a bound $7C_\Phi\Psi^2$ for $\E[\mathbf 1_{\cG_p}\Phi(X(Q))R_{ab}(z)]$, a statement about \emph{one} off-diagonal resolvent entry at one
spectral parameter (original: $|\E F'(x^R)R_{ij}\overline{R_{ja}R_{bi}}|\le N^{-4/3+C\e}$ for $a,b,i,j$ distinct, (6.70)). The proof is new: $R_{ab}=-R_{aa}Y$ with $Y$ centred given the minor $Q^{(a)}$, and
$X(Q)-X(Q^{(a)})=\frac1\pi\im\br{\Log(-R_{aa})}$ (Corollary~\ref{cor:Xchange}(ii)). Not used: the large deviation bounds of Lemma~3.3 of the original, bounds on the quantities $Z^{(S)}_{ja}$, any local law for minors, and the
passage from $R$ to $S$. In particular the step ``$S$ is the Green function of a usual generalized Wigner matrix with all variances being positive'' (TeX 4259--4269), which is not justified under (A)--(C) (variances may
vanish, and Lemma~3.3 is stated for equal variances), does not occur.

\item[\textbf{II.6.}] \textbf{Bounds for $R=(Q-z)^{-1}$.} The original transfers the bound (6.33) from $S$ to $R$ by a tenth-order resolvent expansion closed with the trivial bound $|R_{ij}|\le\eta^{-1}$ ((6.45)). NP uses the exact closed
identity $R_P=S_P+S_PV_PR_P$ for the $2\times2$ (or $1\times1$) blocks with indices in $P=\{a,b\}$ (Lemma~\ref{lem:RfromS}(a)); only these entries of $R$ are needed. The matrix $Q$, which does not satisfy (A), is never fed to
Theorem~\ref{thm:21}.

\item[\textbf{II.7.}] \textbf{Uniformity in the hybrids.} The original states (6.33) with a maximum over all hybrids $\gamma$ inside the probability, and with error probability $C\exp[-c(\log N)^\xi]$ for any $\xi$. In NP each
Lindeberg step uses an event of its own hybrid only (Lemma~\ref{lem:edge}, Lemma~\ref{lem:RfromS}(b)), with probability bound $2N^{-20}$; the uniformity in the step is the uniformity of the constants of
Theorem~\ref{thm:21} over the class $\fA(\Par)$, which contains every hybrid because the variance matrix is common to all of them (Lemma~\ref{lem:hyb}(ii)). No union bound over the hybrids is taken.

\item[\textbf{II.8.}] \textbf{The edge behaviour of $\msc$.} The exact identity $(\msc+1)^2=(2-z)\msc$ (Lemma~\ref{lem:mscB}) replaces the asymptotic bound (6.34); accordingly the resolvent entries are compared with $-\delta_{kl}$
(the edge value $\msc(2)=-1$) rather than with $\delta_{kl}\msc(z)$, and the control parameter is $\Psi=4N^{-1/3+2\eg}$ (original: $N^{-1/3+2\e}$).

\item[\textbf{II.9.}] \textbf{The branch of the logarithm; simultaneity in $z$.} The determinant formulas involve the principal logarithm. By \eqref{eq:b1:nocut}, $D_p(v;z)$ and $-R_{aa}(z)$ avoid the cut $(-\infty,0]$ for
every $z\in\C_+$, so the exact formulas need no condition along the segment $[E_1,E_2]+i\eta$. The good sets (Lemma~\ref{lem:edge}, Definition~\ref{def:goodB}) are nevertheless defined with ``for all $y$'' in the respective
interval, as in the original's (6.33) (maximum over $E$); this uses the simultaneity in $z$ of Theorem~\ref{thm:21}(ii), stated with the outer probability. Conditions at the two points $z_1,z_2$ would suffice
(remark~(4) after Lemma~\ref{lem:expand}).
\end{enumerate}
\noindent No statement of \S\ref{sec:edge} is weaker than the corresponding statement of the original that is used for Theorem~\ref{thm:24}; Item II.3 records the two forms of Theorem~6.3 that are not reproduced.

\subsection{Part~III (Theorem~\ref{thm:23})}\label{sec:diff}
\begin{enumerate}[leftmargin=1.8em,itemsep=2pt,label=\textbf{III.\arabic*.}]
\item \emph{The statement.} The differences in the statement of Theorem~\ref{thm:23} are PD-04, PD-05, PD-06, PD-16 and PD-21 (Section~\ref{app:diffPD}). In Part~III, PD-06 appears as the rescaling $\kappa_t$ of \eqref{tildeH}
and as the factor $2$ in the eigenvalue generator (Remark~\ref{rem:gen}). The constant is linear in $\norm O_\infty+\operatorname{Lip}(O)$ and depends on $O$ otherwise only through $\rho_O$.
The original states (2.36) for continuous $O$ with $C=C(O)$, citing [ESYY, Thm 2.1]; the proof there applies [ESYY, Thm 4.1] to $G=\int\widetilde O\,\rd y$, which must be smooth, with $C=C(G)$.
For H\"older continuous $O$ NP gets a weaker power rate (Corollary~\ref{cor:cont}); merely continuous $O$ is \emph{not covered} (Remark~\ref{rem:contO}).
\item \emph{The flow.} The flow is defined by the explicit formula (2.29), $H_t=e^{-t/2}H+(1-e^{-t})^{1/2}V$ (see \eqref{Ht}); the SDE (2.28) is not used, and the equation (2.31) for the eigenvalue density appears only as the pointwise identity
of Proposition~\ref{prop:U9}(iii) for a smooth symmetric density. The reference measure is the law of $\lambda(V)$; its identification with the density (2.30) of G$\beta$E is a theorem of NP (Theorem~\ref{thm:JD}), which is not needed for the statement.
\item \emph{A stronger bound.} Part~III proves Theorem~\ref{thm:sharp}, from which \eqref{eq:abstrthm2} follows: the factor $N^{2\eps'}\cdot N^{\eps'}$ is replaced by $e^{A_1(\log\log N)^2}$ and the time threshold
$N^{-1+\delta+\eps'}$ by $(\log N)^2/N$; the bound is uniform in $E$, $b$ and the ensemble.
\item \emph{The route.} [ESYY] Assumption IV is derived from the counting-function form of rigidity \eqref{good} (Lemma~\ref{lem:countC}), not from the local semicircle law; Assumption II is not needed beyond (RA5);
[ESYY, Thm 4.1] (Dirichlet form inequality and auxiliary dynamics) is replaced by Corollary~\ref{cor:gap}, which needs only the relative entropy. The logarithmic Sobolev inequality is proved from the Pr\'ekopa--Leindler
inequality (Theorem~\ref{thm:LSI}) and not by the Bakry--\'Emery argument.
\end{enumerate}

\section{Facts used from Mathlib}\label{app:mathlib}
This appendix lists the standard facts that NP uses without proof, all of which are available in Mathlib; each fact is listed once, followed (in brackets) by the places where it is used.

\emph{Status of the names.} A Mathlib name without mark has been checked to exist; a name marked $^\circ$ is an expected name that must still be checked against the version of Mathlib
used for the formalization. Names that occur only in Section~\ref{sec:edge} are not marked individually and should be checked as well.

\subsection{Linear algebra and matrices}
\begin{enumerate}[leftmargin=2em,itemsep=1pt]\raggedright
\item \emph{Spectral theorem for Hermitian matrices} over $\kk\in\{\R,\C\}$: real eigenvalues and an orthonormal eigenbasis (\texttt{Matrix.IsHermitian.eigenvalues}, \texttt{.eigenvectorBasis},
\texttt{.spectral\_theorem}, \texttt{.eigenvectorUnitary}; for $\kk=\R$ the unitary group is the orthogonal group). Note that \texttt{eigenvalues\textsubscript{0}} is \emph{decreasing}, so NP's $\lambda_k$ is
\texttt{eigenvalues\textsubscript{0}} with the index reversed. Consequences used: the characteristic polynomial is invariant under unitary conjugation;
$\det(A-z)=\prod_j(\lambda_j(A)-z)$ (\texttt{Matrix.IsHermitian.det\_eq\_prod\_eigenvalues}); $\tr A^2=\sum_k\lambda_k(A)^2$ (\texttt{Matrix.IsHermitian.trace\_eq\_sum\_eigenvalues} applied to $A^2$);
the eigenvalues of $cA$, $c>0$. \hfill[Lemma~\ref{lem:mono}(i); Lemma~\ref{lem:basic}(d)--(g), Lemma~\ref{lem:weyl}; \S\ref{sec:edge} (through Lemma~\ref{lem:basic}(d)); \S\S\ref{sec:C1}, \ref{sec:C2}; \S\S\ref{sec:C4}, \ref{sec:C5}]
\item \emph{Finite-dimensional linear algebra}: an injective endomorphism of a finite-dimensional space is bijective (\texttt{LinearMap.injective\_iff\_surjective}$^\circ$); a one-sided inverse of a square
matrix is two-sided (\texttt{Matrix.mul\_eq\_one\_comm}$^\circ$); $\dim(V\cap W)\ge\dim V+\dim W-n$ (\texttt{Submodule.finrank\_sup\_add\_finrank\_inf\_eq}$^\circ$); the inverse of a symmetric matrix is symmetric;
inverse of a diagonal matrix. \hfill[\S\ref{sec:det}; \S\ref{sec:FA}]
\item \emph{Determinants}: multiplicativity; $|\det U|=1$ for unitary $U$, $\det O=\pm1$ for orthogonal $O$; the determinant and the inverse of a $2\times2$ matrix (\texttt{Matrix.det\_fin\_two},
\texttt{Matrix.inv\_def}$^\circ$, \texttt{Matrix.adjugate\_fin\_two}$^\circ$); $\det(1+\mathsf X\mathsf Y)=\det(1+\mathsf Y\mathsf X)$ for rectangular $\mathsf X,\mathsf Y$ (\texttt{Matrix.det\_one\_add\_mul\_comm});
a $2\times2$ complex matrix has two eigenvalues counted with multiplicity whose product is its determinant (algebraic closedness of $\C$; \texttt{Matrix.det\_eq\_prod\_roots\_charpoly}).
\hfill[\S\ref{sec:det}; \S\ref{sec:edge}; \S\ref{sec:C3}; \S\S\ref{sec:C4}, \ref{sec:C5}]
\item \emph{Block matrices}: the block inverse formula \eqref{eq:BI} (\texttt{Matrix.invOf\_fromBlocks\textsubscript{22}\_eq}) and the block determinant formula \eqref{eq:BD}
(\texttt{Matrix.det\_fromBlocks\textsubscript{22}}, \texttt{Matrix.det\_fromBlocks\textsubscript{11}}); invariance of inverse and determinant under reindexing
(\texttt{Matrix.det\_submatrix\_equiv\_self}$^\circ$, \texttt{Matrix.inv\_submatrix\_equiv}$^\circ$). Both block formulas are also proved in Step~0 of the proof of Lemma~\ref{lem:schur} from the determinant
of block triangular matrices. \hfill[\S\ref{sec:det}; \S\ref{sec:FA} (through Lemma~\ref{lem:schur}); \S\ref{sec:edge}]
\item \emph{Matrix powers and geometric sums}: entries of matrix powers as sums over index sequences (induction on \texttt{pow\_succ}); $(1-T)\sum_{n<n_0}T^n=1-T^{n_0}$ (\texttt{mul\_neg\_geom\_sum}$^\circ$);
$(u+v)^p-v^p=u\sum_{i<p}(u+v)^iv^{p-1-i}$ (\texttt{Commute.geom\_sum\textsubscript{2}\_mul}$^\circ$). \hfill[\S\ref{sec:FA}]
\item \emph{Trace, adjoint, unitary groups, permutations}: $\tr(AB)=\tr(BA)$ (\texttt{Matrix.trace\_mul\_comm}$^\circ$), cyclicity of the trace; $(AB)^*=B^*A^*$; the unitary and orthogonal groups
(\texttt{Matrix.unitaryGroup}); permutation matrices (\texttt{Matrix.PEquiv.toMatrix}$^\circ$ / \texttt{Equiv.Perm.permMatrix}$^\circ$); Parseval in a finite-dimensional inner product space
(\texttt{OrthonormalBasis.sum\_inner\_mul\_inner}$^\circ$). \hfill[\S\ref{sec:C3}; \S\S\ref{sec:C4}, \ref{sec:C5}]
\item \emph{Inequalities in $\C^n$}: Cauchy--Schwarz. \hfill[\S\ref{sec:det}]
\end{enumerate}

\subsection{Measure theory and integration}
\begin{enumerate}[leftmargin=2em,itemsep=1pt]\raggedright
\item \emph{Product measures}: \texttt{MeasureTheory.Measure.pi}, \texttt{MeasureTheory.Measure.prod}; uniqueness of a product measure from its values on boxes (\texttt{MeasureTheory.Measure.pi\_eq}$^\circ$);
splitting of a finite product over a subset of the index set (\texttt{MeasureTheory.measurePreserving\_piEquivPiSubtypeProd}$^\circ$; this is Lemma~\ref{lem:split}(a), which is also proved in the text
from the uniqueness of measures on a generating $\pi$-system); splitting off one coordinate of $\R^N$ is measure preserving (\texttt{MeasureTheory.volume\_preserving\_piFinSuccAbove}$^\circ$);
$\mathrm{Leb}^{2d}=\mathrm{Leb}^d\otimes\mathrm{Leb}^d$ (\texttt{MeasureTheory.volume\_pi}, \texttt{measurePreserving\_sumPiEquivProdPi}$^\circ$); measure-preserving coordinate permutations of $\R^N$.
\hfill[\S\S\ref{sec:prob}, \ref{sec:weak}; \S\ref{sec:FA}; \S\ref{sec:C3}; \S\S\ref{sec:C4}, \ref{sec:C5}; \S\ref{sec:C6}]
\item \emph{Tonelli and Fubini} (\texttt{MeasureTheory.lintegral\_prod}$^\circ$, \texttt{MeasureTheory.integral\_prod}$^\circ$), including the measurability of partial integrals and marginals
(\texttt{Measurable.lintegral\_prod\_right'}$^\circ$, \texttt{MeasureTheory.StronglyMeasurable.integral\_prod\_right'}$^\circ$), on $\R^2$, $\R^3$, $\R^N$ and for series.
\hfill[Lemmas~\ref{lem:stieltjes}, \ref{lem:HS}, \ref{lem:count}; \S\S\ref{sec:prob}, \ref{sec:weak}; \S\ref{sec:FA}; \S\ref{sec:C3}; \S\S\ref{sec:C4}, \ref{sec:C5}; \S\ref{sec:C6}; \S\S\ref{sec:C7}--\ref{sec:C9}]
\item \emph{Image measures and change of variables}: $\int g\,\rd(T_*\mu)=\int g\circ T\,\rd\mu$ (\texttt{MeasureTheory.Measure.map}, \texttt{MeasureTheory.Measure.map\_apply}, \texttt{MeasureTheory.lintegral\_map},
\texttt{MeasureTheory.integral\_map}); image of a product measure under a coordinatewise map; measures with density (\texttt{withDensity}$^\circ$).
\hfill[Lemma~\ref{lem:mono}(iv),(v); \S\S\ref{sec:prob}, \ref{sec:weak}; \S\ref{sec:edge}; \S\S\ref{sec:C1}, \ref{sec:C2}; \S\ref{sec:C3}; \S\S\ref{sec:C4}, \ref{sec:C5}; \S\S\ref{sec:C7}--\ref{sec:C9}]
\item \emph{Lebesgue measure on finite-dimensional spaces}: $\mathrm{Leb}^d$ as the product \texttt{MeasureTheory.Measure.pi}; the image of an additive Haar measure under an invertible linear map is
$|\det|^{-1}$ times it (\texttt{MeasureTheory.Measure.map\_linearMap\_addHaar\_eq\_smul\_addHaar}$^\circ$); translation invariance; equivalently, the volume of a finite-dimensional inner product space and its
invariance under linear isometries (\texttt{LinearIsometryEquiv.measurePreserving}$^\circ$, \texttt{OrthonormalBasis.volume\_parallelepiped}$^\circ$); scaling on $\R$ (\texttt{Real.volume\_preimage\_mul\_left}$^\circ$);
inner regularity of Lebesgue measure on Borel subsets of $\R$ by compact sets (\texttt{MeasurableSet.measure\_eq\_iSup\_isCompact}$^\circ$; Theorem~\ref{thm:PL}, Step~2). \hfill[\S\ref{sec:C3}; \S\ref{sec:C6}]
\item \emph{Convergence theorems}: monotone convergence (\texttt{MeasureTheory.lintegral\_iSup}$^\circ$); Fatou's lemma (\texttt{MeasureTheory.lintegral\_liminf\_le}$^\circ$); dominated convergence, also along
filters (\texttt{MeasureTheory.tendsto\_integral\_filter\_of\_dominated\_convergence}$^\circ$) and for series. \hfill[\S\S\ref{sec:main}, \ref{sec:det}, \ref{sec:rig}; \S\S\ref{sec:prob}, \ref{sec:weak}; \S\ref{sec:C3}; \S\S\ref{sec:C4}, \ref{sec:C5}; \S\ref{sec:C6}; \S\S\ref{sec:C7}--\ref{sec:C9}]
\item \emph{Parametric integrals}: continuity (\texttt{MeasureTheory.continuousOn\_of\_dominated}$^\circ$), also over a compact interval; differentiation under the integral sign with a dominating function
(\texttt{hasFDerivAt\_integral\_of\_dominated\_of\_fderiv\_le}$^\circ$, \texttt{hasDerivAt\_integral\_of\_dominated\_loc\_of\_deriv\_le}$^\circ$). \hfill[\S\ref{sec:C3}; \S\S\ref{sec:C4}, \ref{sec:C5}; \S\ref{sec:C6}; \S\S\ref{sec:C7}--\ref{sec:C9}]
\item \emph{Layer cake}: $\int f=\int_0^\infty|\{f>r\}|\,\rd r$ for Borel $f\ge0$ and its power version (\texttt{MeasureTheory.lintegral\_rpow\_eq\_lintegral\_meas\_lt\_mul};
\texttt{MeasureTheory.lintegral\_eq\_lintegral\_meas\_lt}$^\circ$). \hfill[Lemma~\ref{lem:moments}; Theorem~\ref{thm:PL}, Step~3]
\item \emph{Radon--Nikodym}: \texttt{MeasureTheory.Measure.rnDeriv}$^\circ$, \texttt{Measure.withDensity\_rnDeriv\_eq}$^\circ$; $\int F\,\rd\nu=\int F\frac{\rd\nu}{\rd\omega}\rd\omega$ for $\nu\ll\omega$;
$\nu'=f\mu$, $\omega=\psi\mu\Rightarrow\nu'=(f/\psi)\omega$. \hfill[\S\S\ref{sec:C4}, \ref{sec:C5}; \S\ref{sec:C6}; \S\S\ref{sec:C7}--\ref{sec:C9}]
\item \emph{Uniqueness of measures}: finite Borel measures on a metric space that agree on all bounded continuous functions are equal
(\texttt{MeasureTheory.ext\_of\_forall\_integral\_eq\_of\_IsFiniteMeasure}$^\circ$ or \texttt{MeasureTheory.ext\_of\_forall\_lintegral\_eq\_of\_IsFiniteMeasure}$^\circ$). \hfill[\S\S\ref{sec:C4}, \ref{sec:C5}]
\item \emph{Outer measure}: $\PP(\Omega')\le\PP^*(U)$ for $\Omega'\subset U$ (\texttt{MeasureTheory.measure\_mono}; in Mathlib a measure evaluated on an arbitrary set is its outer measure).
\hfill[statement of Theorem~\ref{thm:21}; \S\ref{sec:edge}]
\item \emph{The Bochner integral}: linearity and monotonicity (\texttt{MeasureTheory.integral\_mono}, \texttt{integral\_indicator\_one}); $|\int f|\le\int|f|$ (\texttt{MeasureTheory.norm\_integral\_le\_integral\_norm});
$\int\im(c\,\xi)=\im(c\int\xi)$ and $\int\bar\xi=\overline{\int\xi}$ for integrable complex $\xi$ (\texttt{integral\_im}, \texttt{integral\_conj}, \texttt{integral\_const\_mul}); continuous linear maps commute with
Bochner integrals (\texttt{ContinuousLinearMap.integral\_comp\_comm}$^\circ$); a probability measure with a finite third absolute moment integrates all polynomials of degree $\le3$.
\hfill[\S\ref{sec:edge}; \S\ref{sec:C3}]
\item \emph{$L^p$ inequalities} (on a probability space, $p\ge1$ real): H\"older's inequality for finitely many functions (\texttt{ENNReal.lintegral\_prod\_norm\_pow\_le}$^\circ$); Minkowski's inequality
(\texttt{ENNReal.lintegral\_Lp\_add\_le}$^\circ$, \texttt{MeasureTheory.eLpNorm\_add\_le}$^\circ$); monotonicity of $L^q$-norms (\texttt{MeasureTheory.eLpNorm\_le\_eLpNorm\_of\_exponent\_le}$^\circ$);
Jensen's inequality for $|\cdot|^q$ (\texttt{ConvexOn.map\_integral\_le}$^\circ$); $\norm f_p\le\norm g_p$ for $0\le f\le g$, $\norm{c\,\bone_A}_p=c\,\PP(A)^{1/p}$, $\norm{f\bone_{A^c}}_p\le\sup|f|\cdot\PP(A^c)^{1/p}$,
$\norm{f^{1/2}}_p\le\norm f_p^{1/2}$ for $f\ge0$; the Cauchy--Schwarz inequality $\int|fg|\le(\int|f|^2)^{1/2}(\int|g|^2)^{1/2}$, in particular $|\int\bone_Af|\le\PP(A)^{1/2}(\int|f|^2)^{1/2}$ and
$L^2\subset L^1$ (\texttt{MeasureTheory.integral\_mul\_le\_Lp\_mul\_Lq\_of\_nonneg} with $p=q=2$). In Sections~\ref{sec:strong}--\ref{sec:norm} all functions to which these are applied are bounded
and continuous on $\Omega$ times indicators of closed sets, so integrability is automatic. \hfill[\S\ref{sec:FA}; \S\S\ref{sec:strong}, \ref{sec:norm}; \S\ref{sec:edge}; \S\S\ref{sec:C1}, \ref{sec:C2}; \S\S\ref{sec:C7}--\ref{sec:C9}]
\item \emph{Markov's (Chebyshev's) inequality} $\PP(f>t)\le t^{-p}\EE f^p$ for $f\ge0$, $t>0$ (\texttt{MeasureTheory.mul\_meas\_ge\_le\_lintegral}$^\circ$ applied to $f^p$). \hfill[\S\ref{sec:FA}; \S\S\ref{sec:strong}, \ref{sec:norm}; \S\S\ref{sec:C1}, \ref{sec:C2}; \S\ref{sec:C3}]
\item \emph{Union bound} (\texttt{MeasureTheory.measure\_iUnion\_le}, \texttt{measure\_union\_le}) over finite index sets (in Part~I of cardinality at most $N^{10}$). \hfill[\S\S\ref{sec:prob}, \ref{sec:weak}; \S\S\ref{sec:strong}, \ref{sec:norm}; \S\ref{sec:edge}]
\item \emph{Conditional expectation}: its defining property, only for the reformulation in Lemma~\ref{lem:split}(d). \hfill[\S\S\ref{sec:prob}, \ref{sec:weak}]
\end{enumerate}

\subsection{Probability: independence, Hoeffding, Gaussian measures, relative entropy}
\begin{enumerate}[leftmargin=2em,itemsep=1pt]\raggedright
\item \emph{Independence}: \texttt{ProbabilityTheory.iIndepFun}; finitely many random variables are independent iff their joint law is the product of their laws
(\texttt{iIndepFun\_iff\_map\_fun\_eq\_pi\_map}$^\circ$, \texttt{indepFun\_iff\_map\_prod\_eq\_prod\_map\_map}$^\circ$, \texttt{ProbabilityTheory.iIndepFun\_iff\_map\_prod\_eq\_prod\_map\_map}$^\circ$,
\texttt{ProbabilityTheory.iIndepFun\_iff\_pi\_map\_eq\_map}$^\circ$); measurable functions of disjoint blocks of an independent family are independent (\texttt{iIndepFun.comp}$^\circ$,
\texttt{ProbabilityTheory.iIndepFun.finset\_comp}$^\circ$, \texttt{iIndepFun.indepFun\_finset}$^\circ$, \texttt{IndepFun.comp}); the coordinates of a product measure are independent
(\texttt{ProbabilityTheory.iIndepFun\_pi}$^\circ$); for independent integrable $\xi,\zeta$, $\int\xi\zeta=\int\xi\int\zeta$ (\texttt{ProbabilityTheory.IndepFun.integral\_mul\_of\_integrable}).
(The splitting of the product law itself is Lemma~\ref{lem:split}.) \hfill[\S\S\ref{sec:prob}, \ref{sec:weak}; \S\ref{sec:FA}; \S\ref{sec:edge}; \S\S\ref{sec:C1}, \ref{sec:C2}; \S\ref{sec:C3}]
\item \emph{Hoeffding's lemma and inequality}, i.e.\ statement (H) of Section~\ref{sec:prob}: \texttt{ProbabilityTheory.hasSubgaussianMGF\_of\_mem\_Icc\_of\_integral\_eq\_zero}$^\circ$,
\texttt{ProbabilityTheory.HasSubgaussianMGF.measure\_sum\_ge\_le\_of\_iIndepFun}. \hfill[\S\S\ref{sec:prob}, \ref{sec:weak}]
\item \emph{Real Gaussian measures}: \texttt{ProbabilityTheory.gaussianReal} with its density \texttt{gaussianPDFReal}$^\circ$, mean $0$, variance $v$, finite absolute moments of all orders
(\texttt{memLp\_id\_gaussianReal}$^\circ$), its moment generating function (\texttt{ProbabilityTheory.mgf}, \texttt{mgf\_id\_gaussianReal}$^\circ$); \texttt{ProbabilityTheory.stdGaussian} may be used as the
definition of $g_1$ (Definition~\ref{c3:def:gauss}); the Gaussian integrals $\int e^{-(u-m)^2/(2\tau^2)}\rd u=\sqrt{2\pi}\,\tau$ and $\int_{\R^N}e^{-c|x|^2}\rd x=(\pi/c)^{N/2}$ (\texttt{integral\_gaussian} and products);
finiteness of polynomial moments of Gaussian densities; integrability of $|x|^pe^{-\frac\kappa4|x|^2}$ on $\R^N$. \hfill[\S\S\ref{sec:C1}, \ref{sec:C2}; \S\ref{sec:C3}; \S\S\ref{sec:C4}, \ref{sec:C5}; \S\ref{sec:C6}]
\item \emph{Relative entropy}: \texttt{InformationTheory.klDiv} and \texttt{MeasureTheory.llr}, used only to identify $S(\cdot|\cdot)$ (\S\ref{c3:ss:KL}, Lemma~\ref{lem:DV}); the data-processing inequality
under a measurable map \texttt{klDiv\_map\_le} (Lemma~\ref{lem:U8}). The chain rule for \texttt{klDiv} is \emph{not} used. \hfill[\S\ref{sec:C3}; \S\ref{sec:C6}; \S\S\ref{sec:C7}--\ref{sec:C9}]
\end{enumerate}

\subsection{Topology and measurability}
\begin{enumerate}[leftmargin=2em,itemsep=1pt]\raggedright
\item Continuous functions are Borel (\texttt{Continuous.measurable}); open and closed sets are Borel (\texttt{IsOpen.measurableSet}, \texttt{IsClosed.measurableSet}); intersections of closed sets are closed
(\texttt{isClosed\_iInter}); preimages of closed sets under continuous maps are closed (\texttt{IsClosed.preimage}); sets $\{f\le g\}$ (resp.\ $\{f<g\}$) with $f,g$ continuous are closed (resp.\ open), hence Borel.
\hfill[\S\ref{sec:det}; \S\S\ref{sec:prob}, \ref{sec:weak}; \S\S\ref{sec:strong}, \ref{sec:norm}; \S\ref{sec:edge}; \S\S\ref{sec:C4}, \ref{sec:C5}]
\item The product $\sigma$-algebra on a finite product of second countable spaces is the Borel $\sigma$-algebra (\texttt{Pi.borelSpace}$^\circ$). \hfill[\S\ref{sec:det}]
\item A countable supremum of measurable functions is measurable (\texttt{Measurable.iSup}$^\circ$); a monotone function $\R\to[0,\infty]$ is Borel; density of $\mathbb Q$ in $\R$. \hfill[\S\ref{sec:det}; \S\S\ref{sec:C1}, \ref{sec:C2}; \S\S\ref{sec:C4}, \ref{sec:C5}; \S\ref{sec:C6}]
\item Continuous images of compact sets are compact; a continuous function on a compact set is bounded (\texttt{IsCompact.exists\_bound\_of\_continuousOn}) and uniformly continuous. \hfill[\S\ref{sec:edge}; \S\ref{sec:C6}]
\end{enumerate}

\subsection{Real and complex analysis}
\begin{enumerate}[leftmargin=2em,itemsep=1pt]\raggedright
\item \emph{Complex numbers.} The principal square root on $\C\setminus(-\infty,0]$ (\texttt{Complex.sqrt} or \texttt{cpow} with exponent $\frac12$): $(\sqrt\zeta)^2=\zeta$, $\re\sqrt\zeta>0$, uniqueness with these two
properties, and holomorphy on the slit plane; the geometric series; Euler's formula for $\sin$, $\cos$; holomorphy of rational functions; for a holomorphic $\Delta$ on an open set,
$\partial_x\Delta=\Delta'$, $\partial_y\Delta=\ii\Delta'$. \hfill[\S\S\ref{sec:main}, \ref{sec:det}, \ref{sec:rig}]
\item \emph{The principal logarithm} \texttt{Complex.log}: $\exp(\Log w)=w$ for $w\ne0$; complex differentiability on the slit plane with derivative $1/w$ (\texttt{Complex.hasStrictDerivAt\_log});
the bounds \eqref{eq:b1:L1}--\eqref{eq:b1:L3} are instances of \texttt{Complex.norm\_log\_sub\_logTaylor\_le} (a proof is given in \S\ref{sec:X}). \hfill[\S\ref{sec:edge}]
\item \emph{One-variable calculus.} The fundamental theorem of calculus, also for continuously differentiable functions $[E_1,E_2]\to\C$ (\texttt{intervalIntegral.integral\_eq\_sub\_of\_hasDerivAt});
integration by parts for $C^1$ functions on a compact interval; substitution $\int_{\phi(a)}^{\phi(b)}g=\int_a^b(g\circ\phi)\phi'$ for $\phi\in C^1$, $g$ continuous (\texttt{intervalIntegral.integral\_comp\_mul\_deriv}),
$x=2-u$ (\texttt{intervalIntegral.integral\_comp\_sub\_left}) and affine changes of variables (\texttt{integral\_comp\_mul\_add}$^\circ$); monotonicity of the integral (\texttt{intervalIntegral.integral\_mono\_on});
product and chain rules; the intermediate value theorem; the mean value theorem and monotonicity from the sign of the derivative (\texttt{antitoneOn\_of\_deriv\_nonpos}$^\circ$; a function with zero
derivative on an interval is constant, \texttt{is\_const\_of\_deriv\_eq\_zero}$^\circ$); the mean value inequality $|\Phi(s)-\Phi(s')|\le\sup|\Phi'|\,|s-s'|$ for $\Phi\in C^1(\R)$
(\texttt{Convex.norm\_image\_sub\_le\_of\_norm\_deriv\_le}); a differentiable even function has derivative $0$ at $0$; Taylor's formula with Lagrange remainder (\texttt{taylor\_mean\_remainder\_lagrange}).
\hfill[\S\S\ref{sec:main}, \ref{sec:det}, \ref{sec:rig}; \S\ref{sec:edge}; \S\S\ref{sec:C4}, \ref{sec:C5}; \S\ref{sec:C6}; \S\S\ref{sec:C7}--\ref{sec:C9}]
\item \emph{Specific integrals}: $\int_\R\frac{\rd s}{1+s^2}=\pi$ (\texttt{integral\_univ\_inv\_one\_add\_sq}); $\int_0^\pi\sin^2=\frac\pi2$ and $\int_0^\pi\cos(j\theta)\,\rd\theta=0$ for integers $j\ge1$;
$\int_0^r\sqrt u\,\rd u=\frac23r^{3/2}$ and $\int_{m-1}^mx^{-4/3}\rd x$; the substitution $x=u^{\fa}$ over $[1,\infty)$ and $\int_0^\infty e^{-u/2}\rd u=2$. \hfill[\S\S\ref{sec:main}, \ref{sec:det}, \ref{sec:rig}; \S\S\ref{sec:prob}, \ref{sec:weak}; \S\ref{sec:edge}]
\item \emph{Calculus in several variables.} The divergence theorem on a rectangle for a continuously differentiable vector field with values in $\C$
(\texttt{MeasureTheory.integral2\_divergence\_prod\_of\_hasFDerivWithinAt\_off\_countable}; used in Lemma~\ref{lem:HS} only); $C^k$ functions on finite-dimensional real spaces (\texttt{ContDiff}), composition with
continuous linear maps and of $C^\infty$ maps, the chain rule for $t\mapsto\Phi(\mathbf c(t))$ up to second order, in particular for $t\mapsto g(z+th)$ with $g\in C^2(\R^N)$ (\texttt{iteratedFDeriv},
\texttt{ContDiff.comp}); symmetry of the second derivative of $C^2$ maps (\texttt{second\_derivative\_symmetric}$^\circ$, \texttt{ContDiffAt.isSymmSndFDerivAt}$^\circ$); a function on an open set whose
iterated derivatives all exist and are continuous is $C^\infty$ (\texttt{contDiffOn\_of\_continuousOn\_differentiableOn}$^\circ$). \hfill[\S\S\ref{sec:main}, \ref{sec:det}, \ref{sec:rig}; \S\ref{sec:C3}; \S\S\ref{sec:C4}, \ref{sec:C5}; \S\ref{sec:C6}]
\item \emph{Smooth transition function} \texttt{Real.smoothTransition} with \texttt{.contDiff}, \texttt{.zero\_of\_nonpos}, \texttt{.one\_of\_one\_le}, \texttt{.nonneg}, \texttt{.le\_one} ; used for $\chi(t):=\texttt{smoothTransition}((4-t)/3)$ in Section~\ref{sec:C6}. No bump function is needed in Parts~I and~II: the cutoffs of \S\ref{sec:HS}
and of \S\ref{sec:proof24} are explicit piecewise polynomials. \hfill[\S\ref{sec:C6}]
\item \emph{Elementary functions.} Elementary properties of $\log$, $\exp$ and real powers; monotonicity of $s\mapsto s^a$ on $[0,\infty)$ ($a>0$) and of $a\mapsto N^a$ ($N\ge1$), $N^aN^b=N^{a+b}$,
$(\log N)^{aL_0}=\exp(aL_0\log\log N)$ (\texttt{Real.rpow\_le\_rpow}, \texttt{Real.rpow\_le\_rpow\_of\_exponent\_le}, \texttt{Real.rpow\_def\_of\_pos}, \texttt{Real.rpow\_add});
$x^n/n!\le e^x$ for $x\ge0$ (\texttt{Real.pow\_div\_factorial\_le\_exp}$^\circ$); $1+x\le e^x$ (\texttt{Real.add\_one\_le\_exp}); $\log t\le t-1$ (\texttt{Real.log\_le\_sub\_one\_of\_pos});
convexity of $t\mapsto t\log t$ (\texttt{Real.strictConvexOn\_mul\_log}$^\circ$); concavity of $\log$ (\texttt{strictConcaveOn\_log\_Ioi}$^\circ$); $|\sinh x|\ge|x|$, $\cosh t=1+2\sinh^2(t/2)$, monotonicity of
$\cosh$ on $[0,\infty)$; $\arctan$: derivative, strict monotonicity, range $(-\frac\pi2,\frac\pi2)$, oddness, $\arctan x+\arctan\frac1x=\frac\pi2$ for $x>0$ (\texttt{Real.arctan\_inv\_of\_pos});
$\pi>3$ (\texttt{Real.pi\_gt\_three}). \hfill[\S\S\ref{sec:main}, \ref{sec:det}, \ref{sec:rig}; \S\S\ref{sec:prob}, \ref{sec:weak}; \S\ref{sec:edge}; \S\ref{sec:C3}; \S\ref{sec:C6}; \S\S\ref{sec:C7}--\ref{sec:C9}]
\item \emph{Asymptotics.} $(\log\log N)^2/\log N\to0$ and $(\log N)^2/N^a\to0$ for $a>0$ (\texttt{Real.isLittleO\_log\_rpow\_atTop}$^\circ$). \hfill[\S\S\ref{sec:C1}, \ref{sec:C2}; \S\S\ref{sec:C7}--\ref{sec:C9}]
\item \emph{Elementary inequalities.} Cauchy--Schwarz and Young; the weighted arithmetic--geometric mean inequality $\tau s+(1-\tau)t\ge s^\tau t^{1-\tau}$
(\texttt{Real.geom\_mean\_le\_arith\_mean2\_weighted}$^\circ$); $(a+b)^2\le2a^2+2b^2$, $(a+b)^3\le4(a^3+b^3)$ ($a,b\ge0$), $(\alpha+\beta+\gamma)^2\le3(\alpha^2+\beta^2+\gamma^2)$; $\sqrt{\alpha+\beta}\le\sqrt\alpha+\sqrt\beta$,
$2\sqrt{\alpha\beta}\le t\alpha+\beta/t$, $\sqrt{1+r}\le1+r/2$; $(\im\xi)^2=\frac12(|\xi|^2-\re\xi^2)$ for $\xi\in\C$; monotonicity of $\eta\mapsto\sqrt{\kappa+\eta}\,\eta$ and of $x\mapsto x^p$;
absorption (if $0\le Y<\infty$ and $Y\le\theta Y+B$ with $0\le\theta<1$ then $Y\le B/(1-\theta)$); comparisons of the form $N^Ce^{-K}\le\frac12e^{-K/2}$ when $K\ge(\log N)^3$ and $(\log N)^2\ge2C+2$;
$\sum_{k\le m}k^{-2/3}\le3m^{1/3}$ by comparison with $\int u^{-2/3}$. Numerical inequalities between explicit constants (such as $(3\pi/2)^{2/3}>2.8$ or $2^{1/4}/13>\frac1{150}$) are decidable by
interval arithmetic. \hfill[\S\S\ref{sec:main}, \ref{sec:det}, \ref{sec:rig}; \S\ref{sec:det}; \S\S\ref{sec:prob}, \ref{sec:weak}; \S\S\ref{sec:strong}, \ref{sec:norm}; \S\ref{sec:edge}; \S\S\ref{sec:C1}, \ref{sec:C2}; \S\ref{sec:C6}]
\item \emph{Floor and ceiling}: $x-1<\lfloor x\rfloor\le x\le\lceil x\rceil<x+1$ (\texttt{Int.floor\_le}, \texttt{Int.lt\_floor\_add\_one}, \texttt{Int.le\_ceil}, \texttt{Int.ceil\_lt\_add\_one}); $\lfloor a\rfloor+1>a$
(\texttt{Nat.lt\_floor\_add\_one}). \hfill[\S\S\ref{sec:strong}, \ref{sec:norm}; \S\ref{sec:edge}]
\item \emph{Lipschitz functions}: the infimum of a family of $L$-Lipschitz real functions, if finite, is $L$-Lipschitz (\texttt{LipschitzWith.iInf}$^\circ$, or directly). \hfill[\S\S\ref{sec:C7}--\ref{sec:C9}]
\end{enumerate}

\subsection{Combinatorics and finite sums}
\begin{enumerate}[leftmargin=2em,itemsep=1pt]\raggedright
\item The number of $(e_s)\in\Z_{\ge0}^p$ with $\sum e_s=j$ is $\binom{j+p-1}{p-1}\le2^{j+p-1}$ (\texttt{Nat.multichoose}$^\circ$, \texttt{Sym.card\_sym\_eq\_multichoose}$^\circ$); the binomial theorem
(\texttt{add\_pow}); geometric series. \hfill[\S\S\ref{sec:main}, \ref{sec:det}, \ref{sec:rig}; \S\S\ref{sec:prob}, \ref{sec:weak}; \S\ref{sec:FA}]
\item Telescoping sums (\texttt{Finset.sum\_range\_succ\_sub\_sum}); the numbers of diagonal and off-diagonal pairs in $\mathcal P_N$ are $N$ and $N(N-1)/2$. \hfill[\S\ref{sec:edge}]
\item The bijection between ordered $n$-tuples of distinct indices and pairs (subset, permutation); $\binom Mk\le M^k$; $\prod(1-a_k)\ge1-\sum a_k$ for $a_k\in[0,1]$. \hfill[\S\S\ref{sec:C7}--\ref{sec:C9}]
\item Induction on a finite chain $z_0,\dots,z_{k_*}$ (the column of a lattice point, Sections~\ref{sec:weak} and~\ref{sec:strong}). \hfill[\S\S\ref{sec:strong}, \ref{sec:norm}]
\end{enumerate}

\subsection{What is not used}
The following are \emph{not} used: Courant--Fischer in full, Cauchy interlacing, Neumann series, the operator norm of $B$, derivatives of matrix functions [\S\ref{sec:det}];
no martingale inequality, no moment generating function of an unbounded variable, no regular conditional probability [\S\S\ref{sec:prob}, \ref{sec:weak}]; no conditional expectation and no limit in Sections~\ref{sec:strong}--\ref{sec:norm},
whose only probabilistic inputs are Corollary~\ref{cor:FA} (at $p_1$ and at $p_2$), Lemma~\ref{lem:norm0} and the probabilities of the events $\cG_{\rm col}(z)$ [\S\S\ref{sec:strong}, \ref{sec:norm}]; no bump function [\S\S\ref{sec:main}, \ref{sec:det}, \ref{sec:rig}];
Haar measure (on $U_\beta(N)$ or elsewhere), complex Gaussian measures, any semigroup, Markov process or Dirichlet-form theory, stochastic calculus, the chain rule for \texttt{klDiv} [Part~III];
the Weyl integration formula, change of variables on manifolds, the implicit function theorem, eigenvalue perturbation theory, Rademacher's theorem [\S\S\ref{sec:C4}--\ref{sec:C6}];
Sobolev spaces, measurability of sumsets, the Brunn--Minkowski inequality in dimension $>1$, Pinsker's inequality [\S\ref{sec:C6}].

\noindent\emph{Remarks.} Each fact used in Section~\ref{sec:edge} has a proof of a few lines. Besides Mathlib, Sections~\ref{sec:C4}--\ref{sec:C5} use from other sections only
Lemmas~\ref{lem:M1}--\ref{lem:M3} (Section~\ref{sec:C3}), Lemma~\ref{lem:weyl}(a),(b) and Definition~\ref{def:loggas} with its Comment~(2) (Section~\ref{sec:C6}).


\end{document}
